\part{II. Auslegungen über das zweite Buch Mosis}

\chapter[A. D. Martin Luthers Auslegung über etliche Capitel des]{A. D. Martin Luthers Auslegung über etliche Capitel des zweiten Buchs Mosis.\footnote{Nachdem Luther die Predigten über das erste Buch Mosis beendigt hatte, fing er am 19. Sonntage nach Trinitatis, den 2. October 1524, an, in den Wochengottesdiensten über das zweite Buch Mosis zu predigen. Erst im Jahre 1526 kam er damit zu Ende. Diese Zeitbestimmung, welche Aurifaber im ersten Theile der Eislebischen Sammlung gibt, wird bestätigt durch eine Notiz in einem von Bachs Predigten in der Zwickauer Rathsschulbibliothek (Buchwald, Ungedruckte Predigten, Bd. I, S. XX, Anm. 3.): „Anno 1524 Sonntag nach Michaelis hatt er angefangen zu predigen und erklären das ander buch Mose.“ Diese Predigten sind nicht aus Luthers Feder geflossen, sondern aufgefangen worden. Ein Manuscript derselben war im Besitz des Propsts Georg Buchholzer in Berlin; doch wurde ihm daselbe, wie er an den Grafen Vollrath von Mansfeld schrieb, im Jahre 1564 während seiner Krankheit entwendet. Ein anderes, unvollständiges Exemplar (welches nur die Predigten über die ersten achtzehn Capitel enthielt) ermittelte man unter den Büchern des Philipp Fabricius, gewesenen Pfarrers zu Ungleben; dieses erkaufte Herr Wolf von Schönburg, Herr zu Glauchau und Waldenburg, und überließ es an Joh. Aurifaber zur Herausgabe. Aurifaber veröffentlichte nun im Jahre 1564 diese Predigten im ersten Theile der Eislebenschen Sammlung, wo sie sich Bl. 279 bis 392 b finden. Danach sind sie wieder abgedruckt in der Altenburger Ausgabe, Bd. III, S. 523; in der Leipziger, Bd. III, S. 471, und in der Erlanger, Bd. 35, S. 1. Wir geben den Text nach der Eislebenschen Sammlung.}}

[Vom 1. bis 18. Capitel.]

Gepredigt zu Wittenberg vom 2. October 1524 bis 1526.

\section{Vorrede}

1. Wir haben bis anher das erste Buch Mosis ausgelegt, und gehört von mancherlei schönen, tröstlichen Sprüchen, darinnen die Hauptartikel unseres christlichen Glaubens begriffen sind. Auch sind uns daneben vorgehalten worden allerlei herrliche Exempel und Vorbilder eines göttlichen und christlichen Lebens an den Altvätern und Patriarchen, und haben aus demselbigen Buche auch gelernt, wie GOtt der allmächtige Schöpfer [des] Himmels, [der] Erde und aller Creaturen sei, den wir allein ehren, lieben, fürchten und über alles ihm vertrauen sollen; desgleichen, wie alle Dinge in GOttes Gewalt und Hand stehen, und er alles noch regieren, schützen und erhalten müsse; item, wie das menschliche Geschlecht sich mehren und
\quellseite{344}
machen solle; auch, wie Adam und Eva durch den Fall in Sünde, Tod, des Teufels und Höllen Gewalt gerathen sind, aber durch den verheißenen Samen wiederum Gerechtigkeit, ewiges Leben und Seligkeit erlangt haben.

2. Auch haben wir gehört, wie die beiden Geschlechter der Menschen, als, die Gottesfürchtigen und Gottlosen, das ist, die rechte, wahrhaftige Kirche GOttes, und die falsche Kirche, des Teufels Braut, in der Welt gelebt, gewachsen und es getrieben haben. Und sind uns in diesem Buche tröstliche Exempel göttlicher Gnaden und Wohlthaten auch vorgestellet, wie GOtt die Frommen gesegnet und ihnen Gutes gethan, und dagegen Exempel des Zorns, wie er die Bösen gestraft habe. Ist also dasselbe Buch wohl werth, daß wir kundig und läufftig darinnen werden, auf daß wir daraus zum Reiche Christi und zu einem christlichen, bußfertigen Leben uns zu ziehen und zu schicken wüßten.

3. Nun greifen wir zur Auslegung des andern Buchs Mosis, darinnen gehandelt wird vom Auszuge der Kinder Israel aus Egypten, da sie aus dem Diensthause und harten Gefängniß Pharaonis errettet und erlöset worden sind. Das hält uns auch vor schöne, herrliche Exempel, wie GOtt wahrhaftig seine Verheißungen hält; item, Exempel göttlicher Gnade und Güte gegen betrübten, geplagten und geängstigten Christen; und dagegen Exempel seines Zorns, Strafen und Ungnade gegen den Halsstarrigen und Unbußfertigen. Und, nachdem Abrahams, Isaaks und Jakobs Nachkommen ein eigen Volk und Reich werden sollten, so gibt GOtt in diesem Buche ihnen auch Gesetze und Rechte, bestellt auch den Gottesdienst und Religion, und ordnet Obrigkeiten; als den Mosen, der ein Heerführer und Hauptmann dieses Volks sei, der sie aus Egypten führen solle; und theilt das Volk in zwölf Fürstenthümer oder Stämme, darüber er zwölf Fürsten setzt. [Er] verfasset also dies Volk mit Gesetzen und Rechten, auf daß (weil dieses Volks sehr viel war) ein jeglicher gleichwohl wüßte, was er thun sollte.

\section{Das erste Capitel}

1. Wir hören in diesem Capitel von mancherlei Stücken, sonderlich, wie GOtt anfähet, seinem Volke aus der Dienstbarkeit Pharaonis zu helfen, darinnen sie in Egypten waren, und wie er den Mosen ordnet zum Herzog über die Kinder Israel, und ihn wunderlich zu dem Amte bringt, der doch allerlei Ausflüchte suchte, auf daß er solchem Beruf ausschlüge. Erstlich aber sagt der Text:

\begin{schriftzitat}
V. 1–7. Dies sind die Namen der Kinder Israel, die mit Jakob in Egypten kamen; ein jeglicher kam mit seinem Hause hinein: Ruben, Simeon, Levi, Juda, Isaschar, Sebulon, Benjamin, Dan, Naphtali, Gad, Asser. Und aller Seelen, die aus den Lenden Jakobs gekommen waren, der waren siebenzig. Joseph aber war zuvor in Egypten. Da nun Joseph gestorben war, und alle seine Brüder, und alle, die zu der Zeit gelebt hatten, wuchsen die Kinder Israel und zeugeten Kinder, und mehreten sich, und wurden ihrer sehr viel, daß ihrer das Land voll ward.
\end{schriftzitat}

2. Im ersten Buche Mosis haben wir gehört, wie GOtt eine Linie oder Schnur gezogen und geführt hat von Adam, auch allen andern Altvätern und Patriarchen bis auf Jakob, und hat sich GOtt nicht anders gestellt, denn als sähe er allein auf die Schnur, so von Adam, Seth, Noa, Sem, Abraham und Isaac, bis auf den Patriarchen
\quellseite{345}
Jacob gehe, und thut gleich, als kenne er die andern Geschlechte in der Welt nicht, wisse auch nichts von ihnen; wie er denn die andern alle, so neben dieser Linie und Schnur gelebt haben, mehrlich\footnote{mehrlich = kaum.} berührt. Wiewohl nun dieselbige Schnur des menschlichen Geschlechts sehr klein gewesen (denn er ist damit nicht weiter gekommen, denn bis auf Jacob und seine zwölf Söhne, und die von ihnen gebornen, als nämlich siebenzig Seelen, welche in Egypten gezogen waren, 1 Mos. 46, 27.), so wird doch im Anfange dieses ersten Capitels beschrieben, wie dies Volk Israel in Egypten gewachsen sei, andere Kinder gezeugt habe, sich gemehret und also zugenommen, daß das Land ihrer gar voll worden sei, und sie hernach eigene Regimenten und Reiche, geistlich und weltlich bekommen.

3. Daran sehen wir unseres Herrn Gottes Meisterstücke, wie wunderbarlich er regiere, seine Verheißungen und Zusagungen erfülle und halte. Dünne und schmal greift es GOtt an, da er seine Verheißung, dem Abraham und seinem Samen gegeben, in das Werk setzen will. Denn er ihm [1 Mos. 15, 5.] verheißen hatte, sein Same sollte also ein groß, unzählig Volk werden, als der Sand am Meer, und die Sterne am Himmel wären; auch [1 Mos. 17, 6.] sollten Könige und Fürsten von seinen Lenden kommen; aber die Verheißung läßt sich ansehen, als wollte sie gar zu Wasser werden. Denn Abraham und Sara, dergleichen Isaak und Rebekka, sind einzelne Leute, und sterben dahin, haben wenig Kinder. Wo bleibt da die Verheißung? Dünne und geringe fähet es Gott an; allein mit dem Patriarchen Jacob geht es ein wenig an, dem GOtt zwölf Söhne gibt, davon bald siebenzig Seelen gezeuget werden, und will nun die Verheißung etwas blühen und läßt sich sehen.

4. Diese sind das große Volk, das der Patriarch Abraham haben soll. Er zwar stirbt, und fährt dahin in diesem Glauben, daß GOtt ihm die Verheißung halten werde, ob er gleich nicht mehr mit seiner Sara denn nur den Isaak sieht. Denn GOtt ist wahrhaftig in seinem Worte und Zusagungen. Aber gleichwohl stellt er sich also wunderlich und seltsam, gleich als wollte er die Verheißung nicht halten; und thut es doch. Dadurch wir sollen verursacht werden, GOttes Verheißungen zu glauben, festiglich darauf zu bauen, und in allen Trübsalen und Anfechtungen uns darauf zu verlassen, auch mit Geduld seiner Hülfe zu erwarten, daß GOtt zuletzt noch komme, uns tröste und aus Nöthen errette. Folget weiter:

\begin{schriftzitat}
V. 8. Da kam ein neuer König auf in Egypten, der wußte nichts von Joseph.
\end{schriftzitat}

5. Wie das zugegangen sei, ob vielleicht das Land Egypten mit dem vorigen Könige über einen Haufen gefallen, das kann ich nicht wissen. Aber es läßt sich ansehen, daß der fromme Pharao, der zur Zeit Josephs gelebt und regieret hatte, gestorben, und ein neuer König an seiner Statt erwählet sei worden, der von dem Joseph nichts gewußt habe. Denn Joseph hat bei hundert und zehn Jahren gelebt [1 Mos. 50, 22.] und war bei neun und dreißig Jahre alt, da sein Vater, der Patriarch Jacob, zu ihm in Egypten kam. Darnach ist er noch siebenzig Jahre alt worden, und hat auf das wenigste ein fünfzig Jahre nach seines Vaters Tode regiert. Da hat es sich in fünfzig Jahren gar sehr verändert, daß einer oder zween neue Könige in Egypten gefolgt sind nach Jacobs und Josephs Tode.

6. Egypten ist ein feines Land und Königreich gewesen, und damals in der Blüthe gestanden. Als Joseph, voll des Heiligen Geistes, darinnen regiert hat, da ist es das schönste und fruchtbarlichste Königreich unter der Sonne gewesen; aber da Joseph stirbt, und das rechte Haupt und der Kern hinweg ist, da geht es viel anders zu; wie es denn pflegt mit Regimenten und Reichen zu steigen und zu fallen.

7. Also thut es Gott mit allen Königreichen und Regimenten. Er gibt einem jeden Lande seine Zeit zu wachsen und zu steigen, daß es an Reichthum, Macht, Ehre und Gewalt zunimmt und erweitert wird, grünen, in Ehren und Würden schweben möge. Aber wenn es auch wiederum unserm HErrn GOtt Zeit dünkt, daß solches Land wieder fallen und zu Boden gehen solle, so fällt es auch plötzlich wieder dahin, daß es niemand aufhalten kann. Darum sagt der Prophet Daniel [Cap. 2, 21.] recht: Deus constituit et transfert regna, GOtt ist es, der da Königreiche einsetzet und verändert.

8. Also half er dem römischen Reiche auf die Beine, daß es viel andere Königreiche unter sich zwang und niederdrückte, und ein gewaltiges Kaiserthum ward; aber es ist wieder zu Sumpfe gegangen. Item, das jüdische Reich hat GOtt
\quellseite{346}
auch herrlich und groß gemacht, aber nun liegt es wieder in der Asche. Dergleichen das babylonische Kaiserthum ist ein gewaltiges Regiment auf Erden gewesen; aber es hat auch nicht lange gestanden, darnach ist es wieder über einen Haufen gefallen. Also gibt GOtt auch einer jeglichen Stadt ihre Zeit, da sie im Aufnehmen und in der Blüthe ist; wenn denn die Zeit ihres Zunehmens vorüber ist, so ist denn plötzlich das Regiment aus, und kommt wieder in das Abnehmen und Fallen.

9. Also handelt unser HErr GOtt: Er kann aus Nichts Alles machen; ein kleines Reich kann er wohl weiter und groß machen, und wiederum, ein groß Kaiserthum leichtlich stürzen, daß alle Welt daran genug zu fühlen hat [Jes. 13, 19.]. Und wenn es denn im Wechsel oder im Wenden ist, daß er ein Königreich zu Boden stößt, so hilft da denn kein Wehren noch Aufhalten; ja, wenn es an eisernen Ketten hinge, oder auf eisernen Pfählen stände, so muß es zu Trümmern und zu Boden gehen.

10. Also ist Egypten ein selig, christlich, herrlich Reich und Fürstenthum gewesen, welches GOttes Wort reichlich gehabt, gepredigt durch die hohen Patriarchen, Jakob und Joseph, auch durch die zwölf Söhne Jakobs, welche unserm HErrngott gar eine feine Kirche daselbst gepflanzet haben. So hat GOtt diesem Reiche auch gegeben löbliche Könige und Regenten, die da Herberge, Schutz und Schirm der Kirche mitgetheilt, und ihre Ernährer, Säugammen und Schutzherren gewesen [Jes. 49, 23.], die [die] Patriarchen gehauset und geherbergt haben.

11. Das weltliche Regiment ist auch herrlich und vortrefflich, daß in diesem Reiche Gericht und Gerechtigkeit zu finden war. So gibt ihm GOtt auch großen Reichthum und Nahrung. Es hat hohe Leute im geistlichen und weltlichen Regimente gehabt, und ist dieses Reich auf starken und festen Beinen gestanden. Aber da die Zeit kommt, daß dieses Reich auch soll einen Stoß und Fall nehmen, so liegt es in einem Hui alles über einem Haufen, daß da die Religion und der Gottesdienst untergeht, und Gerechtigkeit, Gericht, Disciplin und Zucht aufhört und ein Ende hat.

12. Wann geschieht aber dieses? Solches meldet der Text: „Da Joseph gestorben war, und alle seine Brüder, und alle, die zu seiner Zeit gelebt hatten, da kam ein neuer König auf in Egypten, der wußte nichts von Joseph“ etc.

Der Untergang der Kirche und Religion, und Zerstörung der weltlichen Regierung findet sich dann, wenn die gottlosen Regenten, Patriarchen, Prediger, Bischöfe und Pfarrherren, auch die christlichen, frommen Könige, Fürsten und tüchtige Regenten weggenommen werden. Eben also steht auch geschrieben im Buche der Richter am 2. Capitel, V. 8. ff.: „Als Josua, der Sohn Nun, gestorben war, der Knecht des HErrn, als er hundert und zehn Jahre alt war; da auch alle, die zu der Zeit gelebt hatten, zu ihren Vätern versammelt worden, kam nach ihnen ein ander Geschlecht auf, das den HErrn nicht kannte, noch die Werke, die er an Israel gethan hatte. Da thaten die Kinder Israel übel vor dem HErrn, und dieneten Baalim, und verließen den HErrn, ihrer Väter GOtt, der sie aus Egyptenland geführet hatte, und folgten andern Göttern nach.“ Also geht es, wenn die Säulen und Stützel hinweg sind, so die Kirche und weltlich Regiment getragen haben, so fällt es über einen Haufen. Wenn der Kern aus der Nuß gefressen ist, so behält man nur die ledigen Schalen in der Hand, es bleiben nur die Hülsen oder der Schaum übrig. Je größer denn zuvor der Segen GOttes gewesen ist, also hoch ist hernach sein Fluch, Vermaledeyung und Ungnade [Matth. 11, 23.]. Und hat man in demselbigen Lande zuvor die heiligen GOttes geehrt, so denkt man hernach, sie umzubringen und zu tödten, wie den Kindern von Israel in Egypten geschah. Also wird das Land umgekehrt, und werden die Leute GOtt undankbar, weichen von seinem Worte, und verlieren den wahrhaftigen Gottesdienst.

\begin{schriftzitat}
V. 9–11. Und Pharao sprach zu seinem Volke:
\end{schriftzitat}

\begin{quote}
Siehe, des Volkes der Kinder Israel ist viel, und mehr denn wir. Wohlan, wir wollen sie mit Listen dämpfen, daß ihrer nicht so viel werden. Denn wo sich ein Krieg wider uns erhöbe, möchten sie sich auch zu unsern Feinden schlagen, und wider uns streiten [und zum Lande ausziehen].\footnote{Die eingeklammerten Worte fehlen in der Eislebenschen Ausgabe.} Und er setzte Frohnvögte über sie, die sie mit Diensten drücken sollten. Denn man bauete dem Pharao die Städte Pithon und Raemses zu Schatzhäusern.
\end{quote}

13. Hier lerne der Welt Undankbarkeit erkennen. Wenn man ihr gleich alles Gutes thut und lange dient, so schlägt sie doch letztlich einen dafür zum Lohne an den Kopf. Die Welt vergilt
\quellseite{347}
Gutthat mit Uebelthat. Es haben die Egypter bis daher von Joseph und den Kindern Israel viel Gutes empfangen, als, GOttes wahrhaftiges Erkenntniß, daß sie der greulichen, heidnischen Abgötterei los worden, und nun den wahren, rechten Gottesdienst wußten, und den Weg zur ewigen Seligkeit treffen konnten. So hatten sie auch ein christlich, wohlgeordnet, weltlich Regiment, das GOtt wohlgefällig war. Item, sie hatten fromme Obrigkeit gehabt, fruchtbare Jahre, und viel Jahre her guten Frieden und andere leibliche Wohlthaten mehr; also, daß GOtt, um der Israeliten willen, das ganze Reich Egypten gesegnet hatte. Der Wirth mußte des Gastes genießen. Die Egypter hatten von dem Volke Israel ihrer Seelen Heil und Seligkeit, gut Regiment, einen gemeinen Landfrieden, Nahrung und Ueberfluß, da die Zehrungen anfielen, und alle geistliche und leibliche Wohlfahrt. Wie denn noch die gottlose Welt der Christen viel genießt; und um der Frommen willen muß es den bösen Buben oft wohl gehen. Aber dieses ist bei den Egyptern alles vergessen, es gedenkt jetzt Pharao und die Egypter, wie sie die Kinder von Israel alle mit List vertilgen und unterdrücken möchten.

14. Solches ist der Welt Trankgeld und Lohn für getreue Dienste. Sie gibt solch böse Trankgeld. Deß müssen wir auch gewöhnen, es geht in der Welt nicht anders. Denn, wie die Heiden gesagt haben: Nihil citius senescit, quam gratia, man vergißt nichts so bald, als Wohlthat; item: Beneficium in pulvere scribitur, in marmore noxa, Wohlthat schlägt man bald in Wind, aber wenn einer dem andern etwas zu Verdrieß gethan hat, das kann man lange gedenken.

15. Nun schlägt es Pharao weislich an: er will erstlich die Kinder von Israel tilgen durch saure, schwere, harte Arbeit, auf daß sie durch Ziegelsteinstreichen, und andere Frohndienste an den Festungen, abgemergelt würden. Er will ihnen das Futter zu hoch hängen, ob er sie täuben\footnote{täuben = dämpfen, unterdrücken.} könnte. Aber was geschieht? Der Text spricht:

\begin{quote}
3. 12. Je mehr sie das Volk drückten, je mehr sich es mehrte und ausbreitete. [Und sie hielten die Kinder Israel wie einen Greuel.]\footnote{Die eingeklammerten Worte stehen nicht in der Eislebischen Ausgabe.}
\end{quote}

16. Alhier stehe abermals unsers Herrn Gottes Wunderwort, wie er sein Wort und Verheißung vollziehen und hinausführen will. Er geht damit also um, daß man nicht sagen kann, daß es Menschen also hinausgeführt hätten. Es hatte GOtt dem Abraham auch verheißen, er wollte ein groß Volk aus seinem Samen machen. Nun geht er mit den Kindern von Israel also um, daß sie in Egypten Dienstboten und Knechte sein müssen, und läßt sich das Gegenspiel ansehen, als würden sie gar ausgerottet werden. Dennoch geschieht es zuletzt, daß sie ein groß Volk werden; denn GOtt führt sie in großer Anzahl aus Egypten [Ps. 136, 11. 12.]. Also, obwohl Pharao allhier wider die Kinder von Israel auch sehr tyrannisirt, tobt und wüthet, er plagt sie hart mit Frohndiensten und andern Beschwerungen, dennoch hilft es alles nicht, er muß dieses Volk bleiben lassen.

17. GOtt spottet der Tyrannen noch dazu, wenn sie wider die Christen sich also mit ihrem Toben und Verfolgungen legen, und thut es ihnen zum Hohn und Verdruß, ja, zum Trotz, daß er die Seinen schützt und erhält. Als sollte er sagen: Lieber, fahret zu und wehret mir’s! Ich will dennoch mein Volk lassen wachsen und groß werden, daß ihr alle selbst bekennen müsset, daß es durch keine menschliche Gewalt geschehe. Denn GOtt läßt die Kinder von Israel in Egypten zwingen, dringen, drücken und schwächen, und macht sie dennoch so fruchtbar. Es geht mit ihnen von statten. Das Kindergebären ist der Juden Weibern leichtlicher angekommen denn sonst; es ist ihnen in der Geburt nicht also sauer worden als den egyptischen Weibern. Denn GOtt hat es im Sinne gehabt, dieses Volk fortzubringen. Je mehr man nun allhier gewehret hat, je mehr das Volk sich gemehrt; je weniger man es in der Welt leiden will, je mehr es GOtt erhält [Sprüchw. 19, 21.]. Und nachdem GOtt verheißen hatte, sie sollten ein groß Volk werden, so waren sie sicher in ihrem Gewissen, GOtt gefiele es wohl, daß ein jeder Mann ein Weib nehme. Darum, je mehr Pharao wehrt, je mehr es geht.

18. Also geht es auch mit dem lieben Evangelio. Wo dasselbige aufgeht, da müssen alle Anschläge der Tyrannen zurückgehen, die sonst steuern und hindern wollen, daß das Evangelium nicht fortkomme. Aber wider GOttes Wort
\quellseite{348}
hilft keine Weisheit, Macht oder Gewalt, sondern es bringt durch, und reißt alles hinweg, was sich in Weg legt. Denn gleichwie ein Wasser in einem Teiche mag eine Zeitlang aufgehalten werden, daß es nicht ausfließen kann, wenn man es mit Reißholz und Erde verdämmt und verschüttet, und ihm den Lauf verstopft; aber es schwellet sich im Teiche, und wächst das Wasser drinnen, mehrt und sammelt sich, bis daß es irgend ein Loch am Damm finde, da es ausbrechen und ausfließen mag; oder, wo der Teich gar voll Wassers wird, so läuft es über den Damm her, denn es will ungehalten sein, und seinen Strom, Fluß und Lauf haben: also legt sich der Teufel und die Welt auch wider das liebe Evangelium, will es aufhalten und stopfen, hindern Prediger und Zuhörer, legen sich mit allerlei Tyrannei und Verfolgung dawider; aber es hilft nicht. Es will das göttliche Wort seinen Lauf haben, und sollte der Teufel und die Welt darüber rasend und unsinnig werden.

19. Was hat aber GOtt daran für eine Freude? Er thut es darum, daß wir nicht gedenken, wir thun es, sondern, daß es von GOtt herkomme, der thue es allein. Denn also geht es auch fort, wenn gleich die Welt darüber toll und thöricht würde. Deß haben wir allhier ein Exempel. Je mehr Pharao das Volk Israel gedenkt zu schwächen, je mehr es sich ausbreitet und wächst. Es ist Pharao allhier gar unsinnig. Er bauet zwo feste Städte, oder zwo Festungen, zwei Schatzhäuser, da müssen die Kinder von Israel an arbeiten als die Esel. Das ist der erste Anschlag des Königes in Egypten und seiner Räthe, damit sie vermeinen, dieses Volks weniger zu machen. Aber GOtt gedenkt mittlerweile: O, meine Hand ist nicht zu schwach, daß ich dieses Volk nicht ernähren, schützen und vertheidigen sollte. Jes. 50, 2. 59, 1., und spottet noch des Pharao dazu.

20. Denn also thut GOtt wider der Welt und des Teufels Zorn und Toben. Sie müssen mit ihrer Tyrannei, Toben und Wüthen nicht allein dem Evangelio keinen Schaden thun, sondern es nur je weiter fördern und ausbreiten, und den Kopf am Evangelio entzweirennen. Wie denn der andere Psalm, V. 1–5., solches uns auch lehrt, da David spricht: „Die Heiden toben, und die Völker reden vergeblich; die Könige im Lande lehnen sich auf wider den HErrn und seinen Gesalbten und sprechen: Lasset uns zerreißen ihre Bande, und von uns werfen ihre Seile. Aber der im Himmel wohnet, lachet ihrer, und der HErr spottet ihrer. Er wird einst mit ihnen reden im Zorn, und mit seinem Grimm wird er sie erschrecken“, daß sie über eine kleine Zeit werden müssen zu Boden gehen.

21. Dieses geschieht alles darum, auf daß wir lernen sollen, GOtt vertrauen und an ihn glauben. Denn er will und kann uns nicht verlassen. Unsere Feinde thun mit ihrem Zorn nicht mehr, denn daß sie GOttes Werk je länger je mehr fördern, und sich nur hindern. GOtt läßt die Feinde sich abtoben, auf daß sie sich an ihm vergreifen, und GOtt sie wieder bezahle und ihnen ihren verdienten Lohn gebe. Aber mittlerzeit läßt Gott das andre Theil übel drücken, und sieht durch die Finger, daß sie hart geplagt werden, auf daß er ihnen bald aushelfe und sie errette.

22. Es geht heutzutage uns auch also. Unsere Lehre wird heftig verfolgt, man ertränkt, henkt und verbrennt hin und wieder die Christen, ich und du werden übel geplagt; und wir sehen wohl, was jetzt der Papst, Cardinäle, Bischöfe und die Fürsten im Sinne haben. Könnten sie uns in einem Löffel alle ersäufen, so nähmen sie keinen Zober\footnote{Zober = Zuber.} dazu. Aber GOtt sitzt droben im Himmel, sieht ihnen zu, und spricht: „Ihr Buben, ihr hättet längst wohl verdient, daß ich euch mit einer scharfen Ruthe stäupete; darum eilet nur mit eurer Verfolgung, seid sehr zornig und böse, auf daß ich desto eher über euch komme und zu Grunde vertilge.

23. Denn die Feinde des Evangelii müssen an unserm HErrn GOtt ihre Fäuste also verbrennen, auf daß sie über ihrer Verfolgung zu Grunde gehen und in der Asche liegen. Denn GOtt ist ihnen zu stark und mächtig, daß sie ihn von seinem Stuhle und Throne nicht stoßen mögen; darum geht [das] Spiel über ihrem Kopf aus. Wie denn auch der HErr Christus in den Geschichten der Apostel [Cap. 9, 4. 5.] zu Paulo sagt, da er noch wider die Christenheit schnaubte mit der Verfolgung: „Saul, Saul, was verfolgst du mich? Es wird dir schwer sein, wider einen Stachel löcken.“ Denn
\quellseite{349}
wer solches thut, der macht sich selbst blutdürstig, verlahmt und verderbt seine Füße. Wie denn St. Paulus mit seiner Verfolgung auch nicht viel mehr ausrichtete, sondern zuletzt mit seinem Toben aufhören mußte.

24. Das ist sehr tröstlich allen Christen, daß sie wissen: schreien sie in ihrer Angst und Noth zu GOtt, so werde er sie erhören und aus der Verfolgung erlösen, und ihren Feinden und Verfolgern ihren Lohn auch geben [Ps. 142, V. 7.].

\begin{schriftzitat}
V. 15–17. Und der König in Egypten sprach zu den Wehmüttern der hebräischen Weiber, deren eine hieß Siphra, und die andere Pua: Wenn ihr den hebräischen Weibern helft, und auf dem Stuhle sehet, daß [es] ein Sohn ist, so tödtet ihn. Ist [es] aber eine Tochter, so lasset sie leben. Aber die Wehmütter fürchteten GOtt, und thaten nicht, wie der König zu Egypten ihnen gesagt hatte, sondern ließen die Kinder leben.
\end{schriftzitat}

25. Da der erste Anschlag nicht helfen will, da fähet der König eine andere Practic an, welche unser HErr GOtt ihm doch auch zu Boden schlägt. Denn GOttes Gewalt ist stark. Sie stellt sich wohl zuweilen schwach und ohnmächtig, aber zuletzt läßt sie sich wieder sehen und bringt durch. So stellt sich auch GOtt zuweilen wohl albern und närrisch, aber seine Weisheit besteht endlich. Denn die menschliche und göttliche Weisheit laufen wider einander, und muß die menschliche Weisheit unten liegen.

26. Es will Pharao allhier der Kinder Israel Söhne, und nicht die Töchter würgen lassen, auf daß des jüdischen Volks nicht mehr würde, und die Egypter aus dem Lande trieben; gedenkt: Wir können der Israeliten Weiber wohl zwingen, sie müssen unsere Mägde sein, und mit der Zeit wollen wir sie auch umbringen. Die Weiber sollen los gehen, aber die Männer nicht; die Weiber müssen sonst mit der Zeit wohl sterben.

27. Also meint Pharao, er habe unserm HErrn GOtt gar das Loch verrannt und verstopft, daß er nicht ein Haar hindurch brächte, und den Kindern Israel nicht würde helfen können; darum so nimmt er auch zu solchem seinem Anschlag die egyptischen, und nicht die hebräischen Wehmütter. Aber GOtt macht des Königs Rath mit seinen Wehmüttern zu Schanden; wie denn Salomo im Buch der Sprüche am 21. Capitel, V. 30., sagt: „Es hilft keine Weisheit, kein Verstand, keine Kunst wider den HErrn“ [Hiob 5, 13.]. Und der Prophet Jesaias [Cap. 8, 10.] sagt: „Beschließet einen Rath, und werde nichts draus; beredet euch, und es bestehe nicht; denn hier ist Immanuel.“ Es muß dem Pharao nicht also hinausgehen, wie er es im Sinne hat. An eigenen Anschlägen und geredetem Tache, da gehet viel ein. Den Tyrannen fehlen oft ihre Anschläge.

28. Warum braucht aber GOtt der egyptischen Wehmütter dazu, hätte er doch wohl andere hierzu nehmen können, daß den Israeliten wäre geholfen worden? Aber GOtt thut es darum, daß die Feinde der christlichen Kirche müssen selbst den Christen, oft wider ihren Willen und Dank, Gutes thun. Denn die egyptischen Weiber waren sonst von Natur den hebräischen Weibern feind und entgegen; [den]noch sollen sie, wider des Königs ernstes Gebot und gestrengen Befehl, der Hebräer Knäblein leben lassen. Es hätte einer allhier wohl sagen mögen: Ei, wie klug ist unser HErr GOtt! Sollten der Egypter Wehmütter helfen den Kindern von Israel, sind sie doch ihre Todfeinde? Aber was thut unser HErr GOtt? Wo bei Menschen kein Rath, Heil, Hülfe noch Trost ist, da will er helfen; wie er denn im 9. Psalm, V. 10., ein Helfer in der Noth genannt wird. Darum fährt er zu, und nimmt den Wehmüttern das Herz, und kehrt ihnen Sinn und Muth um, und neigt ihr Herz zu den israelitischen Kindern. Das ist unserm HErrn GOtt eine geringe Kunst, denn er herrscht mitten unter seinen Feinden, sagt der 110. Psalm, V. 2. Darum macht er oft der Christen Widersacher und Feinde ihnen zu Freunden. Wie denn Salomo im Buche der Sprüche am 16. Capitel, V. 7., gar herrlich und tröstlich auch davon redet, und spricht: „Wenn jemandes Wege dem HErrn wohlgefallen, so macht er auch seine Feinde mit ihm zufrieden.“

29. Also muß dieser Anschlag dem Könige Pharao auch fehlen. Denn GOtt gibt den Wehmüttern eine Furcht und Schrecken in das Herz, daß sie die Kindlein nicht tödten; wie denn GOtt im fünften Gebote solches ernstlich verboten, daß man nicht solle todtschlagen; wollen derhalben wider dies Gebot GOttes nicht handeln, fürchten sich vor GOttes Zorn und Strafen. Also kann GOtt den Eisenfressern, zornigen Junkern und Scharrhansen steuern, und sie zu
\quellseite{350}
Memmen machen. Darum recht der 72. Psalm, V. 9., sagt: „Vor ihm werden sich neigen die in der Wüste, und seine Feinde werden Staub lecken“, das ist, die Feinde Christi müssen nicht also hoffärtig sein, sondern ihr Pochen, Trotz, Scharren und Verfolgen lassen. Den Staub lecken heißt, in Angst und Zittern sein, und nicht wissen, was man vor Furcht und Schrecken thun solle, daß man auch ließe mit Füßen über sich hingehen. Steine küssen, Erde lecken bedeutet, erschrecken vor GOttes Zorn, ablassen von der Verfolgung; wie Paulus, der Verfolger, da er mit dem Blitz vor Damasko niedergeschlagen ward, erschrickt er auch, leckte Staub, steckte seine Nase in die Asche, spricht mit Zittern: „HErr, was soll ich thun?“ Apost. 9, 6., wird ein Apostel und Prediger. Zuvor hat er die Christen dürftig erwürget, nun macht er sie lebendig mit der Lehre des Evangelii, und thut den Christen alles Gutes.

30. Es haben aber diese frommen, ehrlichen Matronen, diese zwo Wehmütter, etwas gewagt, daß sie dieses gewaltigen Königes Mandat übertreten, und die unschuldigen kleinen Kindlein nicht würgen wollen, sondern haben GOtt mehr vor Augen gehabt, denn ihren König. Und ist recht und wohl gethan. Es wollten Sauls Trabanten [1 Sam. 22, 17.] auch nicht die fünf und achtzig Priester zu Nobe\footnote{Erlanger: Niobe.} erschlagen, welche Doeg verrathen, daß sie dem David in der Flucht Schaubrode hatten zu essen gegeben, und der König ihnen gebot, sie sollten mit der Schärfe des Schwerts die Priester erwürgen. Also [1 Kön. 18, 4.] da Ahab, der König in Israel, mit seiner gottlosen Isebel, Eliam und die andern Propheten, um des Baals falschen Gottesdienstes willen, aus dem Königreiche jagte, und sie alle wollte todt haben, da war Obadja, des Königs Ahab Kämmerer, heimlicher des Königs, seines Herrn, Verbot, und versteckte hundert Propheten in Höhlen, speisete und ernährte sie.

31. Diese alle haben GOtt mehr gefürchtet denn ihre Könige. Wie denn solches auch der HErr Christus gebeut, da er Matth. 10, 28. spricht: „Fürchtet nicht die, so den Leib tödten können, sondern den, der Leib und Seele kann in das höllische Feuer werfen.“ Wäre es aber der König inne worden, so hätte er die egyptischen Wehmütter viel härter gestraft, denn alle hebräischen Weiber. Aber GOtt mit seiner Klugheit macht den großen König also zu Schanden, daß zwo alte Betteln müssen einen solchen mächtigen König äffen und närren. Sie fürchten sich vor seiner großen Gewalt nicht, sondern vor GOttes Ungnade und Zorn. Solches thut GOtt allein.

32. Nun wird es der König zuletzt inne, und ist sehr zornig über die Wehmütter. Da erdenken sie eine Lüge, wie denn im Texte folgt:

\begin{schriftzitat}
V. 18. 19. Da rief der König in Egypten den Wehmüttern und sprach: Warum thut ihr das, daß ihr die Kinder leben lasset? Die Wehmütter antworteten Pharao: Die hebräischen Weiber sind nicht wie die egyptischen; denn sie sind harte Weiber. Ehe die Wehmutter zu ihnen kommt, haben sie geboren.
\end{schriftzitat}

33. Sie entschuldigen sich und sprechen: Ehe denn wir zu ihnen kommen, so haben wir es versäumt. Hier ist die Frage: Ob auch die Wehmütter gesündigt haben, daß sie ihrem Könige und Obrigkeit eine Lüge drehen und eine Lüge sagen? Abraham that auch eine solche Lüge, da er sein Weib Sarai vor dem König Abimelech nannte seine Schwester [1 Mos. 12, 13.]. Darauf soll man antworten, daß, ob die Christen wohl den Heiligen Geist haben, so haben sie doch auch noch Fleisch und Blut an ihnen: so hängen und kleben an ihnen auch noch viel Sünden und Gebrechen, darum sie immerdar beten müssen: „HErr, vergib uns unsere Schuld [Matth. 6, 12.]. Und ob die Heiligen GOttes unterweilen zu wenig oder zu viel thun, so geht es doch mit dem Haufen und Menge der Vergebung der Sünden hinweg [Ps. 32, 6.] und gehört unter den Himmel der Gnaden, das ist, unter die Vergebung der Sünden.

34. Lügen aber heißt das, wenn man dem Nächsten will damit Schaden thun. Wenn denn eine Lüge dieser Art ist, daß man einem am Gut, Ehre, Leib und Leben, Leide thun will oder betrügen, so ist es Sünde; denn ich suche damit meinen eigenen Nutzen und anderer Schaden. Das ist mendacium perniciosum. Aber wenn ich also lüge, daß ich einem nicht zu Schaden, sondern zu Dienste und Nutzen lüge, daß ich sein Gutes und Bestes fördere, so nennt man es eine freundliche Lüge, mendacium officiosum, wenn sich einer also stellt, und etwas, andern zu Dienste, erdichtet. Als Lucä am

\quellseite{351}
24. Capitel, V. 28., stellt sich der HErr Christus gegen den zweien Jüngern, die gen Emmaus gingen, als wollte er vorüber gehen, und that es doch nicht; er hatte es nicht im Sinne, und ließ sich doch also bitten. Also ist diese Lüge allhier auch. Die Wehmütter wollen damit den Kindern von Israel dienen, und sie erretten aus des Pharaonis grimmigem Toben; sie wollen den Israeliten damit keinen Schaden thun, sondern vorkommen,\footnote{vorkommen = zuvorkommen, wehren.} daß Pharao nicht unschuldig Blut mit den kleinen Kindlein vergieße; wollen dem Mordteufel wehren, der da gerne ein Blutbad mit den Knäblein der Hebräer hätte angerichtet. So wollen sie, die Wehmütter, mit dieser ihrer Vorsichtigkeit dem Teufel die Schanze brechen, und der Hebräer Knäblein bei dem Leben erhalten.

\begin{schriftzitat}
V. 20. 21. Darum that GOtt den Wehmüttern Gutes, und das Volk mehrete sich, und ward sehr viel. Und weil die Wehmütter Gott fürchteten, bauete er ihnen Häuser.
\end{schriftzitat}

35. Siehe, wie GOtt reichlich vergilt den Wehmüttern in Egypten, daß sie der Hebräer Knäblein in der Geburt nicht erwürget haben, daß ihnen GOtt hinwieder Gutes thut; ihre Gottseligkeit muß reichlich belohnt werden. Wie denn Gottesfurcht die Verheißung hat, daß sie mit geistlichen und leiblichen Gütern soll von GOtt vergolten und bezahlt werden. Denn also sagt auch der HErr Christus Matth. 6, 33.: „Suchet zum ersten das Reich GOttes und seine Gerechtigkeit, so soll euch das andere alles zufallen.“ Und St. Paulus zum Timotheo [1. Ep. 4, 8.] spricht: „Die Gottseligkeit ist zu allen Dingen nütze, sie hat Verheißung des gegenwärtigen und zukünftigen Lebens.“ Und 1 Sam. 2, 30. spricht GOtt: „Die mich ehren, die will ich wiederum ehren.“ Das ist, die da gottselig leben, die müssen von GOtt wiederum gesegnet werden, daß es ihnen hier in diesem Leben und dort im Himmelreich wohl gehe.

36. Ein Haus in der Schrift heißt nicht ein Balken, Sparre oder Dach, sondern eine Hütte, Wohnung oder Behausung, ein Regiment, da Mann und Weib, Eltern und Kinder, Herren, Frauen und Gesinde bei einander wohnen, Essen und Trinken, Futter und Decke haben. Also wird auch eine Kirche genannt ein Gotteshaus, da die Menschen zusammen kommen, Gottes Wort handeln, die Sacramente gebrauchen, und den wahrhaftigen Gottesdienst treiben, ob gleich kein Haus oder Dach allda ist. Also heißt nun ein Haus, ein Regiment oder Haushaltung, daß Gott die Wehmütter gesegnet hat, daß sie Männer, Kinder, Reichthum, Geld und Gut bekommen, und ihr Geschlecht sich gemehret und ausgebreitet habe. Denn Haus auch oft ein Geschlecht heißt; wie wir Teutschen sagen, das Haus zu Sachsen, das Haus Brandenburg, das ist, der Stamm, die Fürsten zu Sachsen und die Markgrafen zu Brandenburg. Diese Belohnung haben die Wehmütter für ihren Dienst und Amt bekommen, daß sie der Hebräer Knäblein in der Geburt haben leben lassen, und nicht für die Lüge, daß sie den König Pharao betrogen und sein böses Vorhaben, Tyrannei und mörderische Anschläge zunichte gemacht haben.

\begin{schriftzitat}
V. 22. Da gebot Pharao alle seinem Volk, und sprach: Alle Söhne, die geboren werden, werfet in das Wasser, und alle Töchter lasset leben.
\end{schriftzitat}

37. Da die ersten zween Anschläge Pharaonis, seine Verfolgung und Tyrannei, als, mit den Frohndiensten, auch daß man die Knäblein in der Geburt hat morden sollen, in Born gefallen, da fährt der König Pharao gleichwohl fort mit seinem blutigen, mörderischen Vorsatz, er wollte gerne das Volk Israel mit Wurzel und allem ausrotten; darum so läßt er ein öffentlich Edict oder Mandat im Reiche ausgehen, daß man alle Kinder in das Wasser werfen und ersäufen solle.

38. Dieses ist ein ernstes Gebot gewesen, damit der Teufel, der Erzunmörder, den Kindern von Israel heftig hat zugesetzt, und gar viel Kinder werden umgebracht worden sein, und mag dies Gebot irgend ein zwanzig Jahre gestanden haben. Es wird Moses auch gleich in diesem Edict geboren. Sein Bruder Aaron ist drei Jahre älter denn er, und ist nicht in dieser Gefahr des Todes gewesen, sondern Moses wird in das Wasser gesetzt, daß man ihn ersäufen sollte.

39. Aber was thut unser HErr GOtt dazu? Dieser dritte Anschlag Pharaonis muß auch zunichte gemacht werden. Denn Verfolgung und Mord muß der Kirche GOttes und den Christen nicht schaden, der Christen muß dadurch nicht weniger werden, wie Pharao durch seine blutigen Anschläge gerne das Volk getilgt hätte, sondern die Kirche nimmt unter dem Kreuze, des
\quellseite{352}
Teufels und [der] Welt Tyrannei, je länger je mehr zu; wie der alte Lehrer der Kirche, Tertullianus, gar fein gesagt hat: Sanguine martyrum rigatur ecclesia, das ist, die Kirche wird durch der Christen Blut besprenget, gefeuchtet und gewässert. Es wachsen aus der Christen Blut immer andere Christen, die an ihrem Glauben, Bekenntniß Christi, Geduld im Kreuz und Beständigkeit sich spiegeln, verwundern, und auch zum christlichen Glauben bekehrt und selig werden.

40. Es hat der Pabst und sein Anhang zu unserer Zeit auch sehr getobt wider das Evangelium und seine Anhänger, und viel Christen hin und wieder in das Elend von Weib und Kind, Hab und Gut verjagt, verbrennen, henken, [er]tränken, köpfen und erwürgen lassen, und damit die Predigt des Evangelii mehren, und die Zahl der Christen geringer machen wollen. Aber das Evangelium ist je länger je weiter ausgebreitet worden unter diesem Toben und Wüthen des Teufels; so sind auch je länger je mehr Christen worden, gleichwie sich auch allhier die Kinder von Israel unter dieser Verfolgung gemehrt haben.

41. Der Teufel meint es wohl, sammt der Welt, böse, und wollte die Christen gerne aufreißen; aber Gott schickt es den Seinen zum Besten. Solche Verfolgung, Noth und Anfechtung muß ihnen eine Ursache sein, zu GOtt ernstlich zu schreien, ihn zu vertrauen, mit Geduld seiner Hülfe zu erwarten. Der kann es denn wohl anders schicken, nämlich, daß die Christen aus der Noth endlich erlöst, und die Gottlosen mit ihrem Wüthen gestürzt werden. Wie er es denn allhier auch also macht, daß Pharao wohl der Israeliter Söhne in das Wasser wirft und ertränkt; aber er muß zuletzt aufhören. Denn GOtt nimmt den König Pharao auch endlich, sammt aller Mannschaft in Egypten, und treibt sie durch Mosen wieder zum Bade in das rothe Meer, darinnen sie alle ersaufen müssen. Das soll das Ende der Tyrannen sein.

42. Also haben wir bis anher im ersten Capitel des andern Buchs Mosis gehört, wie die Kinder von Israel durch den König Pharao in Egypten sind mit allerlei Plagen beschwert und untergedrückt worden. Als erstlich: wie er mit Last und Arbeit sie geschwächt hat, auf daß sie ja sich nicht mehren. Aber es half nicht. Zum andern, wie der König den Wehemüttern geboten, die jungen Kinder in der Geburt zu tödten. Dieses geschah auch nicht. Zum dritten, daß er ernstlich befahl, die Kinder, so Männlein waren, in das Wasser zu werfen und zu ersäufen, die Mägdlein aber lebendig zu lassen. Darauf aber erfolgt auch nichts.

43. Also muß es gehen dem Volke, so da das wahrhaftige Wort GOttes hat, durch welches Wort sie mit Leib und Seele dem Teufel entlaufen. Dieselbigen müssen in der Welt erhalten: diese schlachtet der Teufel als Schlachtschafe [Ps. 44, 23.], diese martert und plagt er, da er sonst läßt falsche Lehrer und andere gottlose Leute wohl sicher und im Frieden und Ruhe bleiben; allein die GOttes Wort haben, die müssen ihm Haare lassen und Verfolgung von ihm leiden. Wie es denn St. Paulus [2 Tim. 3, 12.] auch geweissagt hat, „daß alle, die da wollen gottselig in Christo Jesu leben, die müssen Verfolgung leiden“. Es sagt’s der HErr Christus auch [Matth. 10, 38. 16, 24.]: „Wer da will mein Jünger sein, der nehme sein Kreuz auf sich, und folge mir nach.“ Hat doch der HErr Christus selbst vom Teufel und der Welt Verfolgung müssen leiden; was wollten wir es denn besser haben, denn er es gehabt?

44. Aber wir haben auch von GOttes großen Wunderwerken gehört, die er an den Seinen, so ihm vertrauen, seinem Worte und Verheißung glauben, thue. Denn, ob er sich erstlich wohl stellt gegen sie, als wollte er sie gar verlassen, daß sie müßten untergehen und verderben: doch tröstet er sie wieder und hilft ihnen aus der Verfolgung, greift den Tyrannen in das Spiel, verwirft ihnen die Würfel, zerreißt die Karten, und macht sie mit allen ihren Anschlägen und Verfolgungen zu Schanden.

\subsection{Allegoria oder geistliche Deutung des ersten Capitels}

45. Wir haben eurer Liebe, und sonderlich denen, so die heilige Schrift studiren und Prediger werden wollen, insonderheit gesagt, daß sie sich hüten sollen vor den geistlichen Deutungen oder Allegorieen (wie man es nennt). Denn ich auch bis anher darauf gegangen bin, daß ich die Historien auf heimliche Verständniß geführt; da haben\footnote{So von uns herübergenommen aus der alten Ausgabe. Eislebensche: hat man.} sich wohl vorzusehen, die damit
\quellseite{353}
umgehen, und man muß sie recht unterweisen.

46. Wehren kann man es nicht, daß man nicht sollte Historien vor sich nehmen und heimliche Bedeutungen daraus ziehen, welches Paulus [1 Cor. 13, 2. 14, 2.] nennt Mysteria. Als, daß ich ein grob Creuzel gebe: Isaac wird auf dem Altare geopfert und bleibt doch lebendig, daß es bedeute, Christus müsse sterben und wieder auferstehen und lebendig werden. Das muß man zulassen, daß man es thue. Denn Paulus sagt: Der Heilige Geist redet Mysteria, das ist sein Amt, er kann es auch treffen und deuten die Mysteria. Wie denn auch St. Paulus thut zum Galatern [Cap. 4, 23.] mit den beiden Weibern, der Sara und Hagar, so die beiden Testamente bedeuten. Item, mit den beiden Söhnen, den Isaac und Ismael, welche zweierlei Völker bedeuten. Wenn er es nicht selbst so geführt und gedeutet, so sollten wohl Säue herein fallen, wild und wüst es gedeutet haben.

47. Derhalben, so leidet sich es nicht, daß ein jeglicher mit seinem Kopf in die heilige Schrift falle, und darinnen grübele und mähre,\footnote{mähren = herumrühren.} wie er will. Es soll sich deß niemand unterwinden, er habe denn den Heiligen Geist. Und St. Paulus hat zum Römern am 12. Capitel, V. 7., gesagt: „Wenn jemand Prophezeiungen führt, so sollen sie dem Glauben ähnlich sein“, analoga fidei. Wer heimliche Deutung führen will, der sehe zu, daß er sie deute auf den Glauben, daß sie sich zum Glauben reimen.

48. Das rede ich darum, daß ich selbst viel Zeit im Gregorio, Hieronymo, Cypriano, Augustino, Origene verderbt und verloren habe. Denn die Väter haben zu ihrer Zeit eine sonderliche Lust und Liebe zu den Allegorieen gehabt, sind damit umher spaziert und [haben] alle Bücher vollgetlickt. Origenes ist fast ein Fürst und König über die Allegorieen, und hat die ganze Bibel durchaus voll solcher heimlicher Deutungen gemacht, die denn nicht eines Drecks werth sind. Die Ursache ist diese, daß sie alle ihrem Dünkel, Kopf und Meinung, wie sie es recht angesehen, und nicht St. Paulo gefolgt haben, der da will den Heiligen Geist darinnen lassen handeln, oder daß die Allegorieen dem Glauben ähnlich wären.

49. Und ich wollte auch nicht davon predigen,

wenn ich es nicht darum thäte, daß ich euch gewöhnte, recht die Allegorieen zu führen, und den Auslegern und Lehrern, die also irren mit den Allegorieen, könnte die Gewalt nehmen und ihnen widerreden, und den rechten Verstand behalten. Origenes hat Christo damit viel Schaden gethan. Denn die Juden spotteten derselbigen Lehrer, die durch Allegorieen die Schrift auslegten, und wußten Christo sein Wort also zu deuten, daß alles, so von Isaac geschrieben wäre, solle von Christo verstanden werden. Denn leichtfertige Geister, die des göttlichen Worts voll und satt waren, gaben sich allein darauf, daß sie viel Allegorieen suchten; und da konnte man schier nicht recht GOttes Wort fassen vor denselbigen Deutungen.

50. Das war köstlich Ding, und da sperrte man den Leuten das Maul auf; wenn man etwas Neues hört und etwas Seltsames bringt, da will man es alles wissen, schreibt und flickt alles voll; aber da wird nichts Gutes aus. Denn wenn man in der Historie überdrüssig wird, so sucht man das, so die Welt ansieht, und sucht einer seine Ehre, der andere sonst etwas an der Schrift, und predigen allein darum, daß sie den Leuten Augen, Ohren, Maul und Nasen aufsperren, und um ein Haar: O, ein gelehrter Mann ist das! Aber man verliert darüber den rechten Grund und Verstand der Schrift, und führt die Leute auf eitel Holzwege.

51. Sehet die Bücher St. Hieronymi an, ich habe sie auch gelesen. Aber er meint, er habe es wohl getroffen, führt prächtige Worte, aber es ist schlecht Ding, und sind nur die Schalen von der Nuß, die Hülsen von der Erbsie. Man muß die Nuß aufbeißen und aufbrechen, da findet man denn den Kern; wer das nicht thut, der kriegt einen Dreck in das Maul. Ein junger Lehrer fällt wohl darauf, und meint, es sei recht und gut, ja, es sei ein heimlich Mysterium drinne verborgen, gleichwie ein Kern in der Nuß steckt; achten die Historien als die Schalen, davon nichts gelegen sei, halten es als ein todt Ding. Ich habe die Historien auch also gehandelt, als daß das Beste sei, wie man eigentlich und rechtschaffen ein christlich Leben führen soll im Glauben, Liebe, Geduld, und wie man im Kreuze beten und sich halten soll.

52. Darum will ich gebeten haben, wer da will studiren in der heiligen Schrift, der büße seinen Vorwitz und Lust nicht an der heiligen
\quellseite{354}
Schrift, sondern nimm vor dich Homerum, Ovidium, Virgilium, oder sonst einen Poeten, und versuche da deine Kunst, und plumpe nicht hier herein in die heilige Schrift, du hast denn zuvor die Füße gewaschen, auf daß du den Kern aus der Nuß und nicht die Schalen überkommest; darum wende es um und suche das Beste. Das Hauptstück und den Grund, oder das Beste in der Schrift, nennen sie auch die Schalen ohne die Nuß, als, die Historien obenhin lesen und wissen. Du aber laß auch dein bestes Studiren sein, daß wir wissen, wie die Historien gehen, wie Abraham im Worte GOttes und Glauben gelebt, und wie es ihm ergangen sei.

53. Das bessert die Leute, richtet auf und tröstet die Gewissen, die in großer Angst und Trübsal stecken, und gibt ein Exempel, recht zu leben, auf daß sie auch also bleiben im einfältigen Verstande. Und bringt die Historie die feinen Exempel des Glaubens und der Liebe, ja, sie bringt ihr Leben hervor. Suchst du und grübelst sonst etwas Anderes daraus, das ist nicht gut; denn du bist auf der unrechten Bahn und gehst auf dem Holzwege, und hast schon dem Teufel hinten und vorne die Thüre aufgethan, und hast einen leichtfertigen Geist, so die heilige Schrift verachtet. Da wird denn der Teufelskopf so voller Allegorieen gemacht, daß man mit Säcken voll dieselbigen ausschüttet, und wirft gar zum Narren darüber, da du doch meinst, du seiest sehr klug.

54. Nun, ich will es nicht wehren; aber ich bitte, daß ihr den Herrn, den rechten Schatz, behaltet, und das vornehmste Hauptstück in der heiligen Schrift, nämlich, daß ihr die heilige Schrift nach der Historie Weg wohl lernet. Denn ich warne euch, darum, daß ihr es in den alten Lehrern viel anders finden werdet, sie sind gar voll der Allegorieen. Origenes hat sich und seine Discipel und die ganze Welt damit verführt, so doch St. Hieronymus von Origene gerühmt, daß er, nach den Aposteln, auf keinen seine Lehre setzen wolle, als auf den Origenem. Aber ich wollte mir St. Hieronymi und St. Origenis Kunst nicht hüten. Denn ein jeglicher spitziger Kopf könnte es auch thun. Da sie mir sollen die Historie tractiren, wie man im Wort GOttes leben, auch wie man glauben und die Liebe üben solle, da sind sie gar stumm, gehen steif drüber hin.

55. Diese Meister der Christenheit haben nicht Eine Historie recht ausgeführt. Da gebührt sich, es anzuzeigen, wie das Leben im Glauben daher geht, und wie einer nur am Glauben fest hält, und Weib und Kind darüber läßt. Das sollte das Vornehmste sein, deß sollte einer voll stecken. Darum behaltet die Regel: daß, wer die Propheten handeln will, der handele sie nach dem Glauben, und mäßige sich der Allegorieen. Denn sie haben tolle Gedanken daraus gemacht, und wer da studirt, der sehe sich vor; denn diese Doctores sind ja so fährlich zu lesen, als irgend Poeten, ja, Virgilius ist viel unschädlicher denn sie. Da wissen wir, daß Virgilius und andere Poeten Narren sind und gewesen haben, aber diese lassen sich mit ihrer Lehre ansehen, daß es alles gut sei, was sie vorgeben, und streichen doch einen Dreck in das Maul.

56. Da sehet zu, die ihr studirt, daß ihr nicht zurück lernet, und mehr vergesset, denn ihr gelernt habt. Denn also ist mir es auch gegangen. Es fehlte mir an einem rechten Meister, da ich noch in die Schule ging, der mich solches unterrichtet hätte. Aber euch wird es jetzt gesagt, sehet zu und behaltet es. Man muß aus der Schrift den rechten Schatz, den Kern, Kraft, Macht, Saft und Schmack nehmen, welches ist das Exempel des Glaubens und der Liebe; darauf soll man sehen; wo es GOtt heraus geschrieben hat, da darfst du nicht tief darnach graben. Darnach, wenn du dies vornehmste Stück hast, so kannst du heimliche Deutung mit einführen neben dem öffentlichen Text, den schmücken, und als schöne Spangen dazu heften.

57. Wie St. Paulus auch also thut, da er lange in der Epistel an die Galater gepredigt hat, daß der Glaube gerecht macht und nicht die Werke des Gesetzes, da kommt er darnach zuletzt [Cap. 4, 22. ff.] auch auf die Allegoriam oder Deutung, und läßt es dabei bleiben. Damit er anzeigt, daß man Allegorieen wohl führen mag, aber also, daß sie sich reimen und dienen zum Grunde. Denn ich weiß auch wohl, daß die heimlichen Deutungen halten den Stich nicht, sollen auch nicht der Grund sein, darauf wir uns verlassen. Denn das soll und muß allein das klare, ausdrückliche und öffentliche Wort GOttes thun, als, an Christum glauben, und den Nächsten lieben; also ist man selig. Andere Lehren und Allegorieen lässest du fahren, als die auch ist St. Pauli [Gal. 4, 22. ff.] mit dem Abraham, daß seine zween Söhne bedeuten die
\quellseite{355}
zwei Testamente.\footnote{Hier scheint uns die Relation unklar und mangelhaft zu sein. Es möchte etwa zu lesen sein: Andere Lehren und Allegorieen lasse du fahren, nimm aber allein diejenigen an, die in der heiligen Schrift gegründet sind, als die auch ist etc.} Denn sonst,\footnote{In der Eisleben’schen Ausgabe steht „sonst“ vor „allhier“.} wenn dieselbige allhier nicht Grund hätte, so würde mein Herz wanken und immerdar bedenken, wie ich des gewiß sein möchte. Denn man würde sagen: Wer weiß, ob es dies bedeutet, ob es auch also sei? Das Herz muß da zappeln, und kann nicht sicher sein, darf sich auch auf die Allegorieen nicht gründen, oder verlassen. Ich muß den öffentlichen Text und Blatt der heiligen Schrift haben.

58. Darum so seid gerüstet, daß ihr die geistlichen Bedeutungen für euer geringstes Studiren haltet, und ergreifet den Glauben, die Liebe und das Wort GOttes, das die Patriarchen gehabt haben und dadurch sind erhalten worden. Dieses laß dein höchstes Studiren sein, und dann dasjenige hernach kommen, daß die Doctores sonst für das Größeste gehalten haben; dann wirst du nützlich predigen und studiren. Du mußt es umkehren, und die Allegorieen zurück werfen, und nicht deuten, wie Origenes gethan hat; denn ich sie alle dahin deute und führe, daß sie dem Predigtamte und Glauben dienen. Also thut ihr auch. Wie ihr denn allhier in einem groben Exempel des dritten Capitels [V. 2. ff.] hören werdet, daß Moses den feurigen Busch brennen sahe. Diesen grünen Busch, darinnen die Feuerflamme und GOtt selbst war, hat man gezogen auf die Jungfrau Maria, und im Advent durch die ganze Welt gesungen. Das hat der Welt darnach wohlgefallen, daß es dieses bedeutete; plumpt so bald hinein und muß köstlich Ding sein, daß es die Jungfrau Maria soll treffen.

59. Wenn nun einer gefragt hätte den, der solches vorgegeben und gedeutet hätte: Wie reimt es sich, daß der Busch soll die Jungfrau Maria sein, die da stehe mit ihrer Jungfrauschaft und mit schwangerem Leibe? Gib mir des Ursachen: so stände er als ein Pfeifer. Denn sprichst du: Wo reimt sich’s hin? Führt es auch zu Christo? Schickt es sich auch zum Glauben und christlichem Wesen? Da muß man sagen: Nein, denn es geht allein auf die Jungfrau Maria, des HErrn Christi Mutter, und auf ihre Jungfrauschaft; dort reimt sich’s denn also hin, gleich wie eine Faust auf ein Auge. Darum ist es ein gefährlich Ding, die heilige Schrift also führen und deuten, wie mir wollen, wie uns recht dünkt. Denn wie reimt sich, eine Mutter sein oder schwanger gehen, und brennen? Wie kann der feurige Busch die Mutter Maria sein? Item, daß der Busch nicht verzehrt wird, solle bedeuten, daß Maria Jungfrau bleibt.

60. Der Lust hat zu diesem Narrenwerk, der mag die Historie am Feste conceptionis Mariae, im Advent, ferner lesen, und sehen, was der selbige Narr daselbst mit der Schrift geäußelt habe, da er die Arche Noäh, die hohen Berge, so aus dem Wasser gucken; item, den Baum des Paradieses und anderes mehr, alles auf Mariam führt, alle Buchstaben mit der heimlichen Deutung auf die Jungfrau Mariam zeucht. Und wir haben es auch im Papstthum gethan, daß wir alles auf die Jungfrau Mariam gezogen haben, als daß sie die sei, da man allen Trost und Zuversicht auf setzen solle; item, daß sie uns aus Nöthen erretten könne. Item, nehmet das Fest corporis Christi auch vor euch, und sehet, was da auch für ein Flickwerk ist, daß es siehet, als wenn es die Säue gespieen hätten. Die Noten und der Gesang ist wohl gut; sonst ist hier ein Stück, und dort ein Fleck daran gehängt, und muß das Himmelsbrod das Abendmahl des HErrn sein, daß es\footnote{Eislebenische: er.} siehet, wie ein gestickter und geflickter Mantel. Das haben die höchsten Theologen gethan, und sind Narren in der Haut drinnen gewesen.

61. Sehet Gregorium an, wie derselbige Mann gaukelt. Wenn er das Maul aufthut in seinen Predigten, so vergißt er des Textes der Schrift, und sagt nichts vom Glauben, von der Liebe noch Kreuze, sondern allein von den geistlichen Deutungen; [den]noch haben diese Lehrer müssen die Säulen der Christenheit heißen, die da die Kirche erhalten sollten. Sie halten ein gut Jahr.

62. Dieses sei gesagt zur Vorrede und Unterrichtung, auf daß man sich vorsehe, und nicht leichtfertig und vorwitzig sei, die heilige Schrift also zu handeln und mit ihr zu gaukeln. Man soll und kann es wohl nicht wehren; aber gehet mit den Allegorieen sanft um. Wenn ihr die Historien [zu Grunde]\footnote{Von uns hinzugefügt.} gelegt habt, dann thut
\quellseite{356}
eine jämmerliche Allegoriam dazu. Thut es denn der Geist, so wird es sich finden, ob die Allegorie auf den Glauben gehe oder nicht etc. Darum schließe ich, und sage, daß wenig Allegorieen auf Mariam gehören.

63. Nun wollen wir das erste Capitel vor uns nehmen, und die Geheimnisse, oder verborgenen, heimlichen Deutungen daraus nehmen und handeln. Denn ich will es nicht lassen eine geistliche Auslegung heißen, denn die ganze heilige Schrift ist geistlich, und ohne sie keine anderen Bücher.

64. Das erste Stück ist, daß das Volk Israel, so da GOttes Wort hatte, vom Könige Pharao mit Ziegelmachen und anderer schweren Arbeit ist untergedrückt und geplagt worden, auf daß sie sich nicht mehren könnten. Das ist das Hauptstück in dieser Historie, daß, die GOttes Volk sein wollen, müssen Verfolgung leiden. Denn Christus [Matth. 16, 24.] sagt: „Wer mein Jünger will werden, der nehme sein Kreuz auf sich, und folge mir nach.“ Also St. Paulus [2 Tim. 3, 12.]: „Wer gottselig in Christo Jesu leben will, muß Verfolgung leiden.“ Item, im 44. Psalm, V. 23., werden die Christen Schlachtschafe genannt, die um Christi willen täglich geplagt werden etc. Deß haben wir allhier ein klares, helles Exempel; und zeuch es nun auf das Predigtamt, so kann diese Geschichte recht und wohl leiden eine heimliche Deutung und einen verborgenen Verstand, nämlich, wie die rechten Prediger und falschen Lehrer mit den Leuten umgehen und thun. Das Volk Israel hat GOttes Wort und rechtschaffene Prediger, da kommt Pharao mit seinen Treibern und seiner Tyrannei, drückt das Volk unter, und muß die Unterdrückung vom Leibe auch auf das Gewissen gezogen werden. Die heilige Schrift ist deß gar voll, wie die menschlichen Satzungen die Gewissen beschweren und den Glauben verderben, und machen, daß er muß untergehen; wie denn der Prophet Jeremias und auch andere Propheten gar voller Klagen sind wider die menschlichen Satzungen, die die Gewissen beschweren und unterdrücken mit falscher Lehre von guten Werken.

65. Also ist diese Deutung vom Predigtamte, als von den rechten und falschen Predigern, ohne Fahr, und reimt sich mit der heiligen Schrift und mit dem Glauben. Und ich bin deß sicher, wenn diese Historie sich reimt auf das Predigtamt, daß die Deutung gewiß ist, und daß der Heilige Geist solches selbst auslegt. Denn des Menschen Herz kann nimmermehr Ruhe haben, es habe denn die rechte Lehre oder das göttliche Wort, daß sie an GOtt glauben und fest an ihm halten; und wer das hat, der kann auch ein recht Leben führen. Diese Lehre kann der Teufel nicht leiden, darum richtet er falsche Lehre an, und schickt falsche Prediger, läßt das Volk zwingen mit Gesetzen zu guten Werken, dadurch sie gen Himmel kommen, wie wir Mönche mit unsern Gelübden und andern Werken gethan haben. Diese falschen Lehrer sind des Teufels Frohnvögte, die eine Last über die andere schlagen auf die Gewissen, martern uns mit dem Ziegelstreichen. Denn die Werke haben uns nicht gedient, sind uns auch nicht zugute kommen, wir haben nicht mehr denn Mühe davon gehabt, und sind dadurch beschwert und gemartert worden mit dem Ziegelstreichen; es ist nicht uns zugute kommen, sondern dem Pabste, Bischöfen, Mönchen und Pfaffen, die sind reich davon worden.

66. Also kann man die Historie recht deuten, daß Pharao zwo Städte bauet von der sauren Arbeit der Kinder Israel. Denn also geht es allen, die unter Menschengesetzen und unter dem Pabstthume leben, sie werden daraus nichts gebessert, und ihre Gewissen werden davon nur gedrückt, gemartert und geplagt.

67. Wem kommt es zugute, daß diese Ziegel gestrichen werden? Fette Wänste werden draus, unsere Güter werden uns aus dem Beutel genommen; wir sind des Pabsts Ziegelstreicher, er hat seinen Wanst gemästet und gefüllt von unsern Gütern; wir aber sind hart beschwert worden, auf daß wir ja nicht schwanger würden und Kinder zeugten; das ist, man hat uns GOttes Wort nicht gepredigt, auf daß wir ja nicht an GOtt glaubten und selig würden. Denn wo man GOttes Wort nicht predigt, glaubt, betet und von GOtt lehrt, da weiß man nichts von GOtt, und da sind wir nicht GOttes Kinder, die zum Reiche GOttes gezeugt würden. Gleichwie die Kinder von Israel allhier bemühet wurden und müde gemacht, daß sie nicht konnten sich mehren und leiblich Kinder zeugen.

68. Dieses könnte man recht ziehen wider die guten Werke, und wäre dieser heimliche Sinn und Verstand auch recht, auf daß unsere guten Werke niedergeschlagen, und GOttes Wort und
\quellseite{357}
unser Glaube aufgerichtet würden. Wie Sanct Paulus auch gethan hat, da er [Gal. 4, 22. ff.] sagt: Sara sei die rechte Hausfrau und die Freie, die da GOttes Wort und den Glauben hat; aber Hagar, die andere Frau, sei das Gesetz, die mit Werken umgeht. Also reimt sich denn die Prophezei und ist dem Glauben ähnlich, wenn man die Allegorieen zum Worte und Glauben wider die Werke führt.

69. Dem Pharaoni werden zwo Städte gebauet, und kommt sein Reich zu Gewalt, Ehren und Reichthum, daß er ein großer, starker, mächtiger König wird. Dergestalt sind auch wir vordes durch des Papstes Reich sehr beschwert, aber seine Irrthümer, seine Decretales und Canones dadurch bestätigt worden. Welches der Teufel darum thut, auf daß er GOttes Wort ja hindere, und uns wehre, daß wir uns nicht mehren, oder zum großen Volk gemacht, und zum Reiche GOttes gezeugt würden.

70. Darum hat Pharao auch einen rechten Namen. Denn Pharao heißt, bloß stehen, oder müßig gehen; als, wenn ein Weib keinen Schleier auf dem Haupte, oder ein Mann seinen Hut nicht aufgesetzt hat, und nicht Kleider am Leibe hat, oder seines Schmucks und Deckels müßig geht. Wie denn im fünften Buch Mosis [Cap. 32, 42.] im Liede gesagt wird von Entblößung des Hauptes; das ist, wenn ein Mensch von GOtt gerissen ist, daß GOtt nicht über uns schwebt, denn er sonst unser Haupt ist; wie St. Paulus [1 Cor. 11, 3.] sagt, daß der Mann des Weibes Haupt sei; also ist GOtt des Mannes Haupt, der Mann aber regiert das Weib, und GOtt regiert durch den Mann; GOtt ist des Mannes Regent. Und das ist auch von der Obrigkeit gesagt, daß GOtt durch sie regiert [Röm. 13, 1. ff. Ps. 82, 1. ff.]. Darum soll man sich vor ihr demüthigen, denn sie ist der Deckel, das Haupt und Regent, da wir müssen gehen, wie GOtt durch sie uns regiert.

71. Wer nun einen solchen Regenten hat, der heißt nicht bloß, ist auch nicht barhäuptig, denn er geht unter GOtt, und hat einen Deckel, der ihn regiert, schützet und seiner wartet; gleichwie ein Weib unter dem Mann geht, und einen Regenten am Manne hat. Wer aber ohne Regenten ist, den heißt die Schrift bloß und barhäuptig. So heißt nun Pharao ein solcher Regent oder König, der uns barhäuptig macht, daß wir ohne Regenten und ohne Reich sind; das ist, diese falschen Lehrer reißen uns von GOtt, daß Christus und GOtt nicht unser Haupt, HErr und König sei, daß wir keinen Glauben an Christum haben, sondern wir haben ihn verloren, gehen barhäuptig, sind unter Pharaone, gleichwie Pharao auch barhäuptig heißt.

72. Diese reimt sich fein zusammen. Denn das ist die Art aller derer, die sich mit den guten Werken bläuen, daß, wenn sie viel Werke thun, so wollen sie dadurch selig werden, und erlangen doch die Seligkeit nicht, es wird nichts draus. Also sind wir barhäuptig, und gehen ohne das Haupt, Christum, wenn wir durch Werke wollen selig werden; wie denn Paulus zum Galatern [Cap. 5, 4.] auch sagt: „Ihr seid aus der Gnade gefallen, und habt Christum verloren, dieweil ihr wollt durch die Werke selig werden“, das ist, ihr seid aus dem Reiche Christi gefallen, und geht ohne das Haupt Christum; gleichwie ein unzüchtiges Weib von ihrem Manne läuft, und sich zu einem andern hält. Das Haupt ist wohl da, aber die Hure will ohne Haupt sein und barhäuptig gehen.

73. Sonst wird in diesem andern Buche Mosis [Cap. 32, 25.] auch gesagt, daß Moses habe gesehen, quod populus esset nudatus, daß das Volk barhäuptig und bloß ging. Denn Aaron hatte sie bloß gemacht, dieweil er einen sonderlichen Gottesdienst vorgab, daß sie dem güldenen Kalbe opfern sollten, damit er sie von GOtt risse und barhäuptig machte, denn sie vertraueten auf daselbige gute Werk.

74. Das bedeutet dieser König, darum muß er diesen Namen haben, auf daß er anzeige, wer er sei, und was das für Gesellen sind, die da von Gesetzen und guten Werken lehren, und den Glauben schweigen, daß es Pharaonis Frohnvögte sind, so die Gewissen beschweren mit guten Werken, und niemand zu Ruhe und Frieden helfen. Darum kann man in der Christenheit nicht recht predigen, man predige denn allein vom Glauben wider die guten Werke und Menschengesetz. Ich sage aber von denen, die Christen sind. Denn mit den Bösen ist es schon aus, die gehören unter den Henker, der da heißt; Ziegel streichen und die Gewissen beschwert; aber wir reden von denen, die Christen sind und in der Christenheit bleiben, dieselbigen sind nicht unter dem Gesetz.

75. Ziegelstreichen, Thonarbeiten in Egypten heißt, daß man sich wohl bläuet und martert mit guten Werken, auf daß man fromm werde, und
\quellseite{358}
sich selbst beschweren mit Satzungen, davon man denn ein blöde Gewissen kriegt; denn keiner davon reich oder selig wird.

76. Die andere Allegoria ist auch auf das Predigtamt zu ziehen; wenn es nicht helfen will, daß durch die Arbeit und Last des Volks weniger werde, daß der König gebeut den Wehmüttern in Egypten, daß sie der Hebräer Knäblein, alsbald sie jung werden, in der Geburt tödten sollen. Denn die falschen Propheten, wenn sie hören, daß ein Mensch vom Glauben predigt, so wollen sie straks einen todt haben, und thun die Leute in den Bann. [Die] Mägdlein aber lassen sie leben, das ist, die von der Predigt des Glaubens abweichen und auf die guten Werke fallen, und wollen vor GOtt ein besser Leben führen, denn andere Leute. Diejenigen, so stark im Glauben sind, und als Männer sich halten, und sich auf gute Werke nicht verlassen wollen, dieselbigen schlaget todt, spricht der Teufel; die andern, die bei den guten Werken bleiben, die lasset leben. Aber das sind rechtschaffene Christen, die bei der Lehre des Glaubens bleiben und darüber halten, werden auch aus solcher Lehre gebessert.

77. Daß Pharao auch die Knäblein alle in das Wasser werfen und ersäufen läßt, das ist die Macht und Gewalt, wenn man mit Gewalt und Durst in die Lehre greift, und mit Wasser, das ist, mit Menschensatzungen, die Leute ersäuft. Denn Wasser ist, überflüssige, falsche Lehre, da man zu guten Werken fällt und ohne Glauben lebt, darinnen muß man ersaufen und verderben. Also habt ihr das erste Capitel gehört, darinnen der Glaube und das Predigtamt ist bedeutet worden, und ist also diese Lehre bestätigt, daß der Glaube allein vor GOtt gelte, und die Werke nichts thun, denn daß sie nur dem Nächsten dienen.

\section{Das zweite Capitel}

\begin{schriftzitat}
V. 1–4. Und es ging hin ein Mann vom Hause Levi, und nahm eine Tochter Levi. Und das Weib ward schwanger, und gebar einen Sohn, und da sie sahe, daß es ein fein Kind war, verbarg sie ihn drei Monden. Und da sie ihn nicht länger verbergen konnte, machte sie ein Kästlein von Rohr, und verklebte es mit Thon und Pech, und legte das Kind drein, und legte ihn in den Schilf am Ufer des Wassers. Aber seine Schwester stund von ferne, daß sie erfahren wollte, wie es ihm gehen würde.
\end{schriftzitat}

1. Die Summa dieses Capitels ist, wie Moses geboren, und heimlich verborgen, aber zuletzt in ein Kästlein gelegt und auf das Wasser gesetzt wird, da ihn denn des Königs Tochter findet, und für ihren Sohn an dem königlichen Hofe erzog; und wie Moses nach vierzig Jahren zu den Brüdern seines Volkes ausging, einen Egypter todtschlug, und aus dem Lande fliehen mußte, darnach der Kinder von Israel Herzog, Hauptmann und Heerführer ward, der sie aus Egypten führte, und ein eigen Königreich aus ihnen zurichtete, Religion, Gottesdienst und weltliche Polizei, Gesetze, Zucht und Disciplin ihnen vorschrieb und ordnete.

2. Hieraus lernen wir das große Mirakel GOttes, wie er für seine Gläubigen so väterlich und reichlich sorgt, und ihnen wunderlich aus allen Nöthen hilft, daß seinen Christen auch oft ihre eigenen Feinde und ärgsten Widersacher zu ihren Ehren, Nutzen und Wohlfahrt selbst helfen und dienen müssen. Denn der König Pharao muß Moses in seiner Tochter Schooß auferziehen, welcher doch Moses und alle Knäblein der Hebräer ersäufen und umbringen wollte. Solches merkt niemand. Denn GOtt offenbart seinen Rath und Vorhaben nicht bald oder jedermänniglich, sondern hält es eine Zeitlang heimlich und verborgen. Und dieweil es niemand weiß, so geht es dahin, daß man nicht drauf Achtung gibt. Und richtet denn GOtt seine Sachen so wunderberlich aus, daß es niemand inne wird, fühlt oder verstehet, bis daß es GOtt verrichtet hat.

3. Siehe aber, ob GOtt nicht des zornigen Pharaonis noch dazu in der Faust spottet. Denn, da in dieser Verfolgung der Kinder von Israel die Eysen (wie man pflegt zu sagen) am Berge stehen, und eitel Tod und Untergang dieses Volks vor Augen ist, da wird Moses geboren.

4. Es wird sonder Zweifel den Kindern von
\quellseite{359}
Israel übel in Egypten ergangen sein. Denn die egyptischen Weiber haben ein fleißig Auge auf die hebräischen Weiber gehabt, wenn sie schwanger gegangen sind, und haben es offenbar und verrathen. Aber mitten in dieser Tyrannei, und da das Edict am strengsten und härtesten erfordert und gehalten wird, da läßt GOtt des Königs Feind, und den Hauptmann der Kinder von Israel geboren werden. Da Pharao gar zornig und böse ist, muß Moses geboren werden, der den Pharaonen in das Meer stürzen soll. Ja, das noch mehr ist, GOtt schickt es also, daß diesen Feind, Mosen, des Pharao Tochter selbst ernähren und am königlichen Hofe auferziehen muß. Es ist ein wunderlich Ding, daß dieser König, so der Kinder von Israel abgesagter Feind ist, soll des Volks GOttes Heiland, Helfer und Schutzherrn selbst auferziehen, welcher soll hernach ihr, auch sein Land und Leute verderben. Es zieht Pharao eine Schlange auf in seinem eigenen Busen, die ihn hernach beißen soll; er erzeucht einen Wolf, der ihn zerreißen und fressen soll.

5. Also wunderbarlich sind unsers HErrn GOttes Gerichte [Jes. 40], 14. Röm. 11, 33. Hätten sie es aber an Pharaonis Hofe geahnt, sie hätten Mosi hunderttausend Tode angelegt. Aber GOtt spricht: Ich will dir sagen, Pharao soll bei seinem eigenen Brode und mit seinen Gütern seinen eigenen Feind auferziehen, und eine Schlange im Busen ernähren; seine eigene Tochter muß des Mosis Mutter sein. Das heißt, dem Pharao in das Maul und Nase gegriffen, und alle seine Anschläge und Tyrannei zu Schanden gemacht.

6. Mit diesem herrlichen Exempel göttlicher Gnade und Güte gegen die Gläubigen sollen wir uns auch trösten, und im Kreuz und Trübsal zum Glauben, Anrufung und Hoffnung der göttlichen Hülfe bereiten und schicken. Denn die Hülfe muß doch endlich kommen, denn GOtt verläßt die Seinen nicht. Er spricht selbst im ein und neunzigsten Psalm, V. 15.: „Ich bin bei ihm in der Noth, ich will ihn herausreißen“ oder erretten, „und er soll mich preisen“ [Psalm 50, 15.]. So sagt der vier und dreißigste Psalm auch, V. 19.: „Der HErr ist nahe bei allen denen, die ihn anrufen“, und die ihn in der Wahrheit anrufen, er siehet auf unsere Noth und Elend, und will uns daraus erretten [Ps. 145, 18.].

7. Wer die Eltern Mosis gewesen sind, wird allhier auch vermeldet, daß er aus dem Stamme Levi geboren sei. Denn ein Mann vom Stamme oder Hause Levi sei hingegangen, und habe eine Tochter Levi zur Ehe genommen; des Vaters Name hat geheißen Amram, und der Mutter Name Jochebed, wie es im sechsten Capitel, V. 20., geschrieben steht.

8. Diese Eltern hätten das Kind, Mosen, gerne verborgen, und bei dem Leben geschützt, behalten es heimlich bei sich drei Monden, müssen es aber zuletzt, aus Furcht des königlichen Edicts, lassen in das Wasser setzen in einem Kästlein, mit Thon und Pech verkleibet, auf daß das Kindlein nicht so bald ertrinken möchte. Diese Eltern sind ohne Zweifel Christen gewesen, haben auf GOttes Zusagung gesehen, und gewartet auf die Erlösung aus dem Zwinghause Egypti, welche sich unter der Verfolgung und Tyrannei Pharaonis, gleich als eine Maus in der Falle, haben drücken müssen, und auf GOttes Hülfe gesehen und gehofft.

9. Davon redet auch die Epistel zu den Hebräern am elften Capitel, V. 23., daß durch den Glauben sei Moses drei Monden verborgen von seinen Eltern, da er geboren war, darum, daß sie sahen, wie er ein schönes Kind war, und fürchteten sich nicht vor des Königs Gebot. Sie sehen, daß das Kindlein schön ist, und haben Hoffnung, GOtt werde ihnen nach seinen Verheißungen einmal geben einen Mann, der sie aus Egypten führete, wie er zugesagt hatte dem Patriarchen Abraham, da [1 Mos. 15, 13. ff.] gesagt wird: „Das sollt du wissen, daß dein Same wird fremd sein in einem Lande, das nicht sein ist, und da wird man sie zu dienen zwingen und plagen vierhundert Jahr. Aber ich will Richter sein über das Volk, dem sie dienen müssen, darnach sollen sie ausziehen mit großem Gute. Sie aber sollen nach vier Diemmalen wieder hieher kommen.“ Und 1 Mos. 48, 21. spricht der Patriarch Jakob zu seinen Kindern: „Ich sterbe, und GOtt wird mit euch sein, und wird euch wieder bringen in das Land eurer Väter“; und gebot im 49. Capitel, V. 29., seinen Kindern, daß sie sollten ihn begraben im Lande Canaan, darinnen Abraham und Isaak auch begraben lagen, anzuzeigen, daß die Kinder Israel nicht sollten in Egypten bleiben, sondern wieder heraus in ihr Vaterland, das Land Canaan, geführet werden.

\quellseite{360}
10. Und im 50. Capitel, V. 24. 25., sagt Joseph zu seinen Brüdern: „Ich sterbe, und GOTT wird euch heimsuchen, und aus Egyptenland führen in das Land, das er Abraham, Isaak und Jakob geschworen hat. Darum nahm er einen Eid von ihnen, daß sie im Auszuge aus Egypten seine Gebeine mit in das Land Canaan führen sollten.“ Mit diesen Zusagungen trösteten sie sich, und hofften, es werde mit der Dienstbarkeit und Tyrannei Pharaonis bald ein Ende nehmen. Denn wenn die Saite auf das höchste gespannt ist, so zerspringet sie gerne, und wenn das Wetter am härtesten und kältesten ist, so bricht es. Also haben sie mitten in der äußersten Noth gehofft, jetzt sei die Zeit der Erlösung vor der Thür, GOTT werde nun helfen, und vielleicht durch diesen ihren Sohn (welcher also ein schönes Kindlein war) dem Volke Israel helfen; darum verbergen sie das Kindlein bei drei Monden bei sich.

\begin{schriftzitat}
V. 5. 6. Und die Tochter Pharao ging hernieder, und wollte baden im Wasser, und ihre Jungfrauen gingen an dem Rande des Wassers. Und da sie das Kästlein im Schilfe sahe, sandte sie ihre Magd hin, und ließ es holen. Und da sie es aufthat, sahe sie das Kind, und siehe, das Knäblein weinete. Da jammerte es sie, und sprach: Es ist der hebräischen Kindlein eines.
\end{schriftzitat}

11. Siehe allhier, wie GOTT auf das arme Kindlein, Mosen, so ein gnädiges Auge hat, daß es der Teufel nicht ersäufen muß, sondern muß aus dem Wasser gezogen werden und lebendig bleiben, und da sonst seines Vaters Amram Mägde seiner sollten gewartet haben, so muß Amrams Tochter das Kindlein Mosen selbst im Kästlein zum Wasser tragen und es fließen lassen, sieht gleichwohl von ferne hin nach, wo das Kindlein bleibe. Da hätte es nun im Wasser verderben müssen, und wäre von den Fischen gefressen worden; aber GOTT hat noch also genaue Achtung auf Mosen, daß des Königes Tochter ihn aus dem Wasser nimmt, und befiehlt ihn aufzuziehen.

12. Also seltsam schickt es GOTT, daß auch seine eigene Mutter muß Mosis Säugamme werden; er wird wieder in seines Vaters Haus gebracht, darinnen ernährt und auferzogen. Und ehe Mosis sollte Noth leiden, so müßte ihm noch der Teufel selbst dienen, und des wüthenden Tyrannen Tochter ihn auferziehen, und für ihren Sohn halten. Dieses großen Königes Tochter muß noch unseres Herrn GOttes Magd sein, und Mosen auferziehen, wiewohl unwissend. Moses ist ein hebräischer Findling, und Pharao Tochter wird oft mit ihm gespielt, ihn getäuzelt, auch sehr geliebet haben.

13. Solches ist unseres Herrn GOttes Kunst, also sorgt er für die Seinen, und hilft ihnen wunderbarlich. „Er tödtet, und macht wieder lebendig, er führet in die Hölle, und bringet wieder gen Himmel“, wie Hanna, Samuels Mutter, in ihrem Lobgesange singt [1 Sam. 2, 6.]. Es muß das Kindlein, Mosis, erst in Gefahr des Lebens kommen, und dem Tode in Rachen geworfen sein, ehe denn er soll zum Heerführer des Volks Israel gemacht werden. Wiederum hebt GOTT ihn auch endlich empor. Auf daß Mosis groß würde, wird er an den Hof zu des Königs Tochter gebracht, wird ihr Sohn, da er sollte hernach die Verfolgung in Egypten umstoßen. Aber wir sehen allhier auch, wie GOTT die Welt zu Narren und Thoren macht. Denn wer will seiner Weisheit widerstehen? Sie wird dieses göttlichen Rathes mit Mosis nicht inne. Denn es scheint auch nicht, daß aus diesem Kindlein ein solcher Mann sollte werden, denn GOTT thut es heimlich, bis daß er es ausgerichtet hat. Denn er regieret also, daß er dasjenige, so gar nichts ist, fein kann zu etwas machen, und uns vor die Augen stellen, daß wir es sehen mögen.

\begin{schriftzitat}
V. 11. 12. Zu den Zeiten, da Mose war groß worden, ging er aus zu seinen Brüdern, und sahe ihre Last, und ward gewahr, daß ein Egypter schlug seiner Brüder, der hebräischen, einen. Und er wandte sich hin- und her, und da er sahe, daß kein Mensch da war, erschlug er den Egypter, und begrub ihn in den Sand etc.
\end{schriftzitat}

14. Die Juden haben allhier viel Lügen erdacht, und vorgegeben, der König Pharao sei über Mosen zornig worden, und ihm vorlegen lassen seine königliche Krone, einen Apfel und Kohlen, und daran erkennen wollen, ob er sein Sohn wäre oder nicht. Aber Mosis habe die Krone weggestoßen und in die Kohlen gebissen, daher habe Moses hernach gelispelt. Aber es sind jüdische Märlein. St. Stephanus redet in den Geschichten der Apostel am 7. Capitel, V. 17–22., viel anders davon und spricht: „Da sich nun die Zeit der Verheißung nahete,
\quellseite{361}
die GOtt Abraham geschworen hatte, wuchs das Volk, und mehret sich in Egypten, bis daß ein anderer König aufkam, der nichts wußte von Joseph. Dieser trieb Hinterlist mit unserm Geschlechte, und handelte unsere Väter übel, und schaffete, daß man die jungen Kindlein hinwerfen mußte, daß sie nicht lebendig blieben. Zu der Zeit ward Moses geboren, und war ein feines Kind vor GOtt, und ward drei Monat ernähret in seines Vaters Hause. Als er aber hingeworfen ward, nahm ihn die Tochter Pharaonis auf, und zog ihn auf zu einem Sohn. Und Moses ward gelehret in aller Weisheit der Egypter, und ward mächtig in Werken und Worten etc.

15. Moses ist ein Mann, der in Egypten ist wohl auferzogen gewesen, und hat Pharao Tochter wohl gewußt, daß er ein hebräischer Mann sei; darum hat er allerlei Sprachen, Künste und anderes mehr lernen müssen, und ist vor Pharao mit Rath und That ein ansehnlicher Mann gewesen. Er ist auferzogen worden als ein edler Knabe, den man zum Rath im Regiment gebrauchen wollte, welcher dermaßen tüchtig wäre in das weltliche Regiment.

16. Dieses hat gewährt bis in das vierzigste Jahr. Da denn mittlerweile (wie ich es dafür halte) das greuliche, erschreckliche Edict und Gebot, die hebräischen Knäblein zu ersäufen, wird aufgehoben sein. Nun geht die Erlösung an, und will sich das Spiel machen; jedoch verzeucht sich es noch eine Weile, bis daß Moses ein großer\footnote{Eiselbenische: größer. — Zu unserer Lesart vergleiche Hebr. 11, 24.} Mann wird, da zeucht er vom Hofe, und will besuchen und besuchen seine Brüder, die Kinder von Israel; denn da stand ihm sein Herz hin. Am Hofe war er der Königin Sohn, ein herrlicher Hofjunker und großer Hans; da war Gewalt, Ehre, Reichthum und Pracht vollauf, und hatte die Anwartung,\footnote{Anwartung = Anwartschaft.} daß nach Pharaonis Tode er möge König worden. Aber er mußte viel leiden, und am Hofe, der heidnischen Religion und des sündlichen, ärgerlichen Lebens halben, von Trunkenheit, Ehrgeiz, Hoffahrt, und Finanzerei\footnote{Finanzerei = Ränke, Betrug.} und Untreu viel sehen, daß ihm wehe that. Sollte er nun dazu stille schweigen, so hatte er im Herzen weder Ruhe noch Frieden; sollte er aber reden, das wilde, wüste Hofleben tadeln und strafen, so hatte er von den Hofschranzen auch Ungnade, Feindschaft, Widerstand und Gefährlichkeit zu gewarten, daß er von ihnen verachtet wurde und für einen Pfaffen und Narren gehalten; darum so ist sein Herz voller Schmerzen gewesen.

17. Solches Stücklein preiset die Epistel zum Hebräern am 11. Capitel, daß sich Moses vom königlichen Hofe Egypti abgezogen, und das Hofleben verlassen habe, und spricht, V. 24. 25. 26.: „Durch den Glauben verleugnete Moses, da er groß ward, ein Sohn zu heißen der Tochter Pharao, und erwählete viel lieber, mit dem Volke GOttes Ungemach zu leiden, denn die zeitliche Ergötzung der Sünden zu haben, und achtete die Schmach Christi für größern Reichthum denn die Schätze Egypti; denn er sahe hin auf die Belohnung“ etc. Und es ist auch ein schön Exempel des Glaubens, daß man GOtt zu Ehren und Gehorsam könne ein Königreich, Land und Leute, so groß Geld und Güter, Ehre und Gewalt in die Schanze schlagen und verlassen, und dagegen zu den armen, bedrängten, beschwerten, verfolgten und geplagten Kindern von Israel sich halten, mit denen Armuth, Gefahr, Spott, Verachtung, Hunger und Kummer leiden. Ein Weltmensch hätte das nicht gethan, sondern gesagt: Es soll sich einer eben so mehr zum Haupte, als (mit Züchten zu reden) zum Hintern oder Rücken halten, oder gesprochen: In großen Wassern da sähet man große Fische. Ein solcher bliebe sein Lebenlang am Hofe. Wie denn manchem das Hofleben sehr beliebt, daß er darnach rennt und läuft, wie er gen Hofe komme, und hoch oben am Brette sitzen möge. Denn wer nicht zu Hofe ist, der will gar erfrieren. Es gedenkt jedermann dem Feuer nahe zu kommen, auf daß er reich, gewaltig und groß werde.

18. Und ist denjenigen Gesellen gleich als den Hühnern, die in einem Hofe gehen, und sehen, daß andere Hühner, so da in einem Korbe eingesperret sind, volllauf zu essen und zu trinken haben, da gehen dieselbigen Hühner um den Korb her, und wären auch gerne hinein zu den andern, die man mästet, und bald schlachten und essen soll. Also sehen solche Gesellen auch, daß am Hofe gele Suppen gegessen werden, viel Fressens und Saufens daran ist, schöne Kleidung, Hoffahrt, Pracht gespürt, auch Geld und Gut allda erlangt wird; darüber mancher
\quellseite{362}
gar in Gefahr Leibes, Lebens, Ehre, Güter, auch in Verlust der Seligkeit kommt, und dies Feuer ihn übel brennt, oder der Hochmuth ihn noch wohl gar erstickt. Wie denn an manchem Hofe gar ein gottlos Sauleben geführt wird, und man daran um Leib und Seele kommt. Daher denn auch das gemeine Sprüchwort entstanden: Lange zu Hofe, lange zu Hölle.

19. Und dieweil der egyptische Hof auch über diesen Leisten ist geschlagen gewesen, daß Abgötterei allda eingerissen, Unglaube, Tyrannei, Gewalt, Unrecht, Unterdrückung der Unterthanen überhand genommen, so gibt Moses dem Hofleben Urlaub, denn er mit gutem Gewissen bei solchem gottlosen, sündlichen und ärgerlichen Wesen nicht länger hat bleiben können. Denn Glaube und ein ärgerlich, sündlich Leben können sich nicht mit einander vertragen. So war auch die Zeit vorhanden, da GOtt die Kinder von Israel aus dem Diensthause des Königs Pharaonis erlösen wollte, und solches soll durch Mosen geschehen. Darum so zieht GOtt den Mosen vom Hofleben, und Moses folgt unserm HErrn GOtt, setzt hintan alle weltliche, vergängliche Wollust, Freude und Herrlichkeit, als die da eine Occasion und Anreizung sind zu vielen Sünden, Argern und Bösen, und weiß, daß GOtt diesen Verlust an irdischer Freude und zeitlichen Gütern wohl reichlich ihm erstatten könne, glaubt auch festiglich, daß GOtt die Noth und Verfolgung, so er und die Kinder von Israel von den Egyptern leiden müssen, werde wenden, und sie daraus erlösen und retten.

20. Als nun Moses aus dem Hofe des Königs Pharaonis geht, da sieht er, wie es einem seiner Brüder geht, daß er von einem Egyptischen todgeschlagen wird: da fährt Moses zu, rächt den hebräischen Mann, und schlägt den Egypter wieder todt. Denn sein Sinn ist dahin je und allewege gestanden, daß er allen gerne geholfen hätte.

21. Da erhebt sich nun eine Frage: ob Moses recht gethan habe, daß er den Egypter todgeschlagen hat? Aber ich lasse solche Frage in ihrem Werth und Würden bleiben, und halte es dafür, er habe recht daran gethan. Und Moses weiß es auch selbst nicht anders, und hält es bei sich dafür, er sündige nicht daran. Denn er hat die dreißig Jahre zugesehen, wie sich dies Volk unter den Egyptern leiden mußte, und gar keine Errettung vorhanden war. So hat er gehört von seinen Brüdern, daß GOtt ihrem Großvater, Abraham, verheißen hatte, er wollte das jüdische Volk heimsuchen, erretten, und aus dem Diensthause Egypten führen, dagegen aber die Egypter strafen. An dieser Erlösung hat er nicht gezweifelt, sondern gehofft, dieweil im Sprüchwort gesagt wird: Wenn die Saite auf das höchste gespannt wird, so zerspringt sie gerne; darum der Egypter grausame Tyrannei wider die Kinder Israel brechen, und ihr\footnote{Siebenische: sein.} Ende gewinnen würde, dieweil sie aufs höchste gekommen und getrieben wäre.

22. So hat ihm auch der Heilige Geist das Herz und den Muth gegeben, daß er zufährt, und den Egypter erschlägt. Denn aus Fleisches und Blutes Nachgierigkeit thut er es nicht, sondern aus des Heiligen Geistes Treibung; es ist in ihm ein göttlicher, heroischer impetus. Denn GOtt wollte den Mann, Mosen, dazu gebrauchen, daß er sollte der Kinder von Israel Hauptmann sein, und sie von dem Pharaone erlösen. Zwar an dieser That sollten die Kinder von Israel gleich ein Vorbild oder Vorbild haben, was GOtt durch den Mann, Mosen, am Pharaone und allen Egyptern thun wollte, nämlich, daß er sie auch erwürgen und tödten würde; wie sie denn hernachmals im rothen Meere alle ersoffen.

23. Darum so ist dieser Todtschlag Mosis eine besondere That, gleichwie [1 Kön. 18, 40.] Elias, der Prophet, zufuhr, und die Baals-pfaffen, zur Zeit des Königs Ahab in Israel, alle todtschlug: item [4 Mos. 25, 7. 8.], Pinehas, der Hohepriester, ein Messer nahm, und den israelitischen Mann erstach, der öffentliche Hurerei mit einem midianitischen Weibe trieb. Aber diese sonderlichen Thaten der großen Wunderleute, die aus besonderer Eingebung des Heiligen Geistes geschehen, soll man nicht nachthun, und wir können auch solchen Exempeln nicht nachfolgen. Man muß den Wundlerleuten Gottes ihre heroischen Wundertaten allein lassen, und nicht eine Folge oder Exempel daraus machen, daß jedermann auch nachthun wollte, sondern man soll solche Thaten bleiben lassen, es wäre denn, daß man auch den Heiligen Geist hätte, als sie gehabt, und solchen hohen Beruf von GOtt bekäme, und man auch solche Ursachen hätte, als sie gehabt.

24. Sonst enthält sich jeder vom Todtschlagen
\quellseite{363}
, denn der Mord ist eine Sünde wider das fünfte Gebot, da GOtt spricht: „Du sollst nicht tödten.“ Moses hatte für sich GOttes Wort, daß das jüdische Volk sollte von den Egyptern errettet werden, und daß GOtt solches wollte durch ihn thun, wie denn St. Stephanus in den Geschichten der Apostel am 7. Cap., V. 23–25., solches zeuget, da er spricht: „Als Moses vierzig Jahre alt war, stieg es in sein Herz, zu besuchen seine Brüder, die Kinder von Israel, und sahe einen Unrecht leiden, da überhalf er, und rächte den, dem Leide geschahe, und erschlug den Egypter. Er meinte aber, seine Brüder sollten es vernehmen, daß GOtt durch seine Hand ihnen Heil gebe, aber sie vernahmen es nicht.“ Und es sollten sich die Hebräer, seine Brüder, je daran gespiegelt und getröstet haben, wie sie GOtt von der Egypter Tyrannei auch alle erretten würde.

\begin{schriftzitat}
V. 13–15. Auf einen andern Tag ging er auch aus, und sahe zween hebräische Männer sich mit einander zanken, und sprach zu dem Ungerechten: Warum schlägst du deinen Nächsten? Er aber sprach: Wer hat dich zum Obersten oder Richter über uns gesetzt? Willst du mich auch erwürgen, wie du den Egypter erwürget hast? Da fürchtete sich Mose und sprach: Wie ist das laut worden? Und es kam vor Pharao, der trachtete nach Mose, daß er ihn erwürgete. Aber Mose flohe vor Pharao, und hielt sich im Lande Midian, [und] wohnete bei einem Brunnen.
\end{schriftzitat}

25. Als der erschlagene Egypter in dem Sande verscharrt war, und Moses sahe, daß zween hebräische Männer sich mit einander zankten, da wollte er sie auch vertragen, und sein Amt ausrichten als ein Regent und Herr der Kinder von Israel, der da seinem Volke vorstehen sollte mit Gericht und Gerechtigkeit, sie schützen und beschirmen, daß ein gemeiner Landfriede erhalten würde, damit die Frommen gehegt, die Bösen gestraft werden möchten; ein Scheidemann, will dem Unschuldigen nicht lassen Unrecht geschehen. [Es] fähet Moses derhalben sein Amt zeitlich an, daß er will ein Heerführer des jüdischen Volks sein, unterfängt sich bald seines Amts und Gewalt. Denn, wie man sagt, so frummt es sich bald, was da will ein Hake\footnote{Eislebensche: „eine Hacke.“ Das Wort „eine“ wird falsch sein, denn bei Luther ist „Hacke“ oder „Hacken“ (das ist, Haken) männlich.} werden.

werden. Item: Urit mature, quod vult urtica manere. Er läßt sich sehen, daß er ihr Herr und Helfer sein sollte, welchen GOtt geschickt hätte den Kindern von Israel.

26. Wie geht es nun Mosi über dem, daß er dem unschuldigen Hebräer helfen will? Er will Scheider sein, und dürfte wohl Scheidemanns Theil bekommen; wie man denn pflegt im Sprichwort zu sagen: Dem Scheider wird gemeiniglich das Beste davon, sie werden gerne geschlagen: also wird auch dieser Hebräer, Mosis Bruder, des guten Mosis Verräther. Er spricht zu ihm: „Willst du mich auch erschlagen wie den Egypter?“ und bringt es vor den König Pharao, welcher Mosi nachtrachtet, daß er ihn umbrächte. Pharao hätte dem Mosi den Hals abgehauen, als einem Todtschläger, hätte er ihn bekommen. Muß also Moses fliehen vor dem Pharao, das ist Mosis Lohn und Traungeld, darum, daß er den königlichen Hof verlassen, und sich zu seinen Brüdern gesellt hat, und will sie aus Egypten führen. Da schelten sie ihn einen Mörder, und kommt durch sie in Gefahr Leibes und Lebens. Er wollte bei den abgöttischen Egyptern in großem Reichthum und Herrlichkeit nicht bleiben am königlichen Hofe, und fällt unter die undankbaren, boshaften Israeliten, an denen kein Dank zu erstehen ist.

27. Also geht es in der Welt, daß mancher dem Regen will entlaufen, und fällt gar in das Wasser. Aber diese Undankbarkeit der Welt soll uns nicht abschrecken, andern Gutes zu thun. Es geht Christen nicht anders in der Welt, Wohlthat wird ihnen mit Uebelthat vergolten. So folgt auch gemeiniglich Gefährlichkeit, Verfolgung, Jammer und Noth, Armuth und Verachtung, Verjagung und allerlei Elend alle denen, die da GOtt dienen wollen, und ihres Berufs, Standes und Amtes getreulich warten [2 Tim. 3, 12.]. Denn da Moses sein Amt allhier ausrichtet, und dem Untergedrückten hilft, so muß er darüber Leibes und Lebens vor dem Könige Pharao nicht sicher sein. Aber GOtt ist mit ihm in diesem Exilio, und er enthält sich im Lande Midian, wohnt unter den Unbekannten, und kriegt da Weib und Kind.

\begin{schriftzitat}
V. 16–20. Der Priester aber in Midian hatte sieben Töchter, die kamen Wasser zu schöpfen, und fülleten die Rinnen, daß sie ihres Vaters Schafe tränkten. Da kamen die Hirten und stießen sie davon. Aber Moses machte sich auf und half
\end{schriftzitat}

\quellseite{364}
ihnen, und tränkte ihre Schafe. Und da sie zu ihrem Vater Reguel kamen, sprach er: Wie seid ihr heute so bald gekommen? Sie sprachen: Ein egyptischer Mann errettete uns von den Hirten, und schöpfete uns, und tränkte die Schafe. Er sprach zu seinen Töchtern: Wo ist er? Warum habt ihr den Mann gelassen, daß ihr ihn nicht ludet mit uns zu essen?

28. Die Midianiter haben an Capten gestoßen gegen Morgen, dahin ist Moses geflohen. Denn als er seine Brüder erstlich besucht hatte, zieht er wieder an den königlichen Hof, und hat sich daran gehalten, bis daß sein Todtschlag ausbricht, und der König ihn erwürgen will; da fleucht er in Midian, und bestellet ihm Gott allda eine Herberge und Wirth. Denn die sieben Töchter des Priesters (so da Wasser am Brunnen schöpften, aber von den Hirten abgestoßen wurden, und Moses ihnen half, auch sie errettete von den Hirten) bringen das Geschrei von Mose an ihren Vater Reguel, und machen, daß er ihn zur Herberge aufnimmt, gibt ihm nachmals eine Tochter zur Ehe. Das hat nun fast bis in das vierzigste Jahr gewährt.

29. Also wunderlich stellt sich GOtt. Es hat Moses Hoffnung gehabt, die Erlösung des jüdischen Volks würde nun bald angehen, denn sie waren mit unerträglicher Last und Arbeit beschwert. Aber da Moses anhebt das Werk der Erlösung, und schlägt einen Egyptier zu Tode, da bricht ihm Gott also bald den Anfang, daß er muß in Midian flüchtig werden, und schiebt die Erlösung auf bis ins vierzigste Jahr; läßt mittlerweile den Mosen einen Fremdling werden, in einem andern und unbekannten Lande, auf daß es ja keinen Schein noch Ansehen hätte, daß er der Mann sollte sein, der das Volk Israel erlösen, und Pharao mit den Egyptern im rothen Meer ersäufen sollte.

30. Diesen Rathschlag GOttes hat niemand gesehen, auch Moses selbst nicht, sondern GOtt weiß ihn allein. Aber wir sollen in diese Weise und Verzüge GOttes, daß er allezeit seine Rathschläge also wunderbarlich hinaus führt, uns mit rechtem Glauben und Vertrauen schicken lernen. Denn der leidige Unglaube hat das Unglück, und verderbt es alles, daß wir GOtt nicht können in die Faust sehen, stille halten, und seiner Hülfe mit Geduld erwarten. Darum sagt Sirach auch [Cap. 2, 16.]: Vae illis, qui perdiderunt sustinentiam [Wehe denen, so nicht beharren], und der 106. Psalm, V. 13., spricht: „Sie hielten dem Rath GOttes nicht aus.“

31. So gibt GOtt nun dem Mose die Ziporam zum Weibe, welche ihm Kinder gebiert. Den ersten Sohn nennt er Gersiom, das ist, einen Gast; den andern Elieser, das ist, GOtt meine Hülfe, daß ihm GOtt aus des Pharaonis Händen geholfen und erlöst hätte. Denn wie sich ein Fall zugetragen hat, darnach haben die Juden ihre Kinder genannt. Ja, Jakob wird daher genannt, daß er seinen Bruder Esau mit der Ferse begriffen hatte [1 Mos. 25, 26.]. Also nennt Moses seinen ersten Sohn auch Gersiom, einen Fremdling oder Gast, darnach den andern Elieser, GOtt mein Helfer; denn GOtt hatte ihn erhalten und ihm geholfen in Egypten.

32. Wir sollen nun hieraus auch lernen, daß, werden wir GOtt vertrauen und an ihn glauben, so sollen wir Haus und Hof, Weib und Kind, Herberge, Nahrung und alle Leibes Nothdurft finden, wo wir auch in der Welt hinkommen mögen, wir sind gleich im Exilio oder nicht. Denn hier steht Beweis, wie Gott Mose in seinem Exilio in Midian reichlich gesegnet. Wie denn der 34. Psalm, V. 11., auch sagt: „Die den HErrn fürchten, sollen nicht Mangel haben an irgend einem Gut.“ Und der sieben und dreißigste Psalm, V. 25., sagt: „Ich bin jung gewesen und alt worden, und habe nie gesehen den Gerechten verlassen, noch seinen Samen nach Brod gehen.“ Und sonst sagt ein anderer Psalm [Ps. 112, 1. 3.]: „Wohl dem, der den HErrn fürchtet, Reichthum und die Fülle wird in seinem Hause sein“ etc.

\subsection{Allegoria oder geistliche Deutung dieses andern Capitels}

33. Wir haben bisher gehört, wie man die heilige Schrift soll handeln, wenn man suchen will heimliche Deutung, daß man ja zuvor sehe, wie man bei dem Hauptsinn und schriftlichen Verstande bleibe, welches der Grund und Kern ist der heiligen Schrift. Darnach, wenn man Lust hat, auch solche heimliche, und, wie man es nennt, geistliche Deutung zu suchen, so ziehe man es dahin, daß es auf den Glauben und das Predigtamt gefügt werde; sonst fährt man hin und her, und ist ungewiß, was man thut.

\quellseite{365}
34. Als, im ersten Capitel hat man gehört, wie die Kinder von Israel mit vieler Arbeit und Last sind in Egypten geplagt worden, und wie der König die Knäblein hat heißen im Wasser ersäufen. Das sind Exempel, so da bedeuten diejenigen, so mit Menschensatzungen die Gewissen regieren, und aus der Gnade GOttes entfallen, und nicht von GOtt, sondern von Menschen regiert werden, welches eine verdammliche Lehre ist, und die Leute von GOtt zum Teufel wendet, und macht ihnen viel Mühe und Arbeit.

35. In diesem zweiten Capitel hören wir, wie Moses geboren, und heimlich von seinen Eltern drei Monat lang lebendig verborgen und verhohlen, und in einem Rohrkästlein, mit Pech verkleibet, auf das Wasser gelegt wird. Dazu, wie ihn des Königs Tochter findet, und das Kästlein aufbrechen ließ, und Moses als ihren Sohn am Hofe auferzieht, bis er 51 Jahre alt wird; da geht er hinaus, seine Brüder, die Israeliter, zu besuchen, und erschlägt einen Egypter, und muß entfliehen, nimmt ein Weib, und zeugt zween Söhne. So viel wird in dem andern Capitel gehandelt.

36. Es wird aber in dieser Historie, wie ihr gehört habt, uns der Glaube gelehrt, nämlich, daß uns GOtt auch durch unsere Feinde und Widerwärtiger ernährt und schützt, die uns sonst gar umbringen wollen. Denn er macht durch Tod Leben, und durch Sünde Gerechtigkeit, und thut den Seinen Gutes, auch durch ihre Feinde, welche Kunst er allein kann; wie auch sonst die Schrift [Sprüchw. 16, 7.] sagt: „Wem der HErr wohl will, dem müssen auch alle seine Feinde günstig sein.“

37. Die heimliche Deutung ziehen wir dahin: Das Kästlein ist das göttliche Wort, Moses liegt in dem Kästlein. Denn in diesem Buche ist Moses ein Prediger und Lehrer, der GOttes Wort führt, und sonderlich das Gesetz GOttes predigt; darum bedeutet das Rohrkästlein billig den Verstand göttliches Wortes, und die heilige Schrift, die darinnen liegt, verkleibet und verpicht. Denn durch die heilige Schrift ist es beschrieben, und der göttliche Verstand der Schrift darinnen verfaßt. Solches Kästlein ist nun nicht schön geschmückt, als wäre es mit Gold und Silber überzogen, oder von köstlichem Fladderholz\footnote{Hesek. 27, 5. Fladderholz (Fladernholz) = Maserholz, von Ahornbäumen.} gemacht, sondern ist von Rohr geflochten, und mit Pech vergossen, oder mit Thon verkleibet. Da ist nichts denn Rohr, Schilf und Binzen\footnote{Binze = Binse; eine Flechtpflanze.} zu sehen. Denn es in Egypten groß Rohr gehabt, das wie Papier anzusehen gewesen, und man auch Schiffe daraus hat machen können.

38. Aber es bedeutet die Einfalt der heiligen Schrift und des göttlichen Wortes, welche mit schlechten, einfältigen, und nicht geschmückten noch prächtigen Worten einher geht, also, daß viel sich daran ärgern; wie St. Paulus [1 Cor. 2, 1. Gal. 4, 13.] auch sagt: „Ich bin nicht zu euch kommen mit prächtigen Worten menschlicher Weisheit.“ Denn sie scheint als ein schlechter Rohrkasten, mit Thon und Pech etc. verschmiert; nichtsdestoweniger liegt Moses, der Knabe, sein lebendig darinnen, und hat den wahrhaftigen Verstand des göttlichen Worts, und des Königs Tochter erbarmt sich über ihn, und nimmt ihn für ihren Sohn an.

39. Die heilige Schrift lautet als ein überdrüssig Ding, ist ungeziert und ungeschmückt, daß einer nicht gedächte, daß Menschen dem göttlichen Wort glauben sollten, dieweil es gar nichts scheint, noch geschmückt ist [Jes. 53, 1.]. Aber der Glaube kommt aus dem göttlichen Worte, und wird der göttlichen Kraft allein gegeben, und nicht dem äußerlichen Schmuck, sondern der inwendigen Eingebung des Heiligen Geistes, daß man dem ungestalten Worte GOttes trauet. Darnach, so ist es noch greulicher und scheußlicher, wenn du siehst, daß alles unförmlich scheint, was die heilige Schrift gebeut und erfordert, als, wenn sie lehrt, man soll nicht geizig sein, das Fleisch und den alten Adam tödten; daneben die Vernunft, die schöne Wege, lehrt, wie man Geld und Gut, Weisheit, Kunst und Ehre überkomme.

40. Nun liegt das Rohrkästlein und Moses im Wasser, und schwimmt, das ist, GOttes Wort geht zu Toden\footnote{zu Toden gehen = eingemengt werden.} unter den Menschenlehren; wie denn auch droben im ersten Capitel [§ 77] von dem Wasser gesagt worden ist. Denn GOttes Wort geht für und für ohne Unterlaß unter Menschenlehren um. Ein jeglicher Ketzer will dafür angesehen sein, als habe er GOttes Wort für sich.

41. Aber zuletzt findet man das Kind Moses am Rande oder Ufer des Wassers, das ist, die
\quellseite{366}
menschliche Weisheit oder Vernunft, die schöne Mäde, die mit ihren Jungfrauen spazieren geht, finden das Buch, die heilige Schrift, und das schöne Kind, Moses, gefällt ihnen wohl, thun das Buch auf, und weil Moses klein und jung ist, so nimmt ihn die Vernunft auf. Wie denn die Kinder von Israel sich von Mose führen lassen bis an den Berg Sinai und sagen [2 Mos. 19, 8.]: „Was GOtt mit uns redet, das wollen wir thun.“

42. Aber wenn Moses vierzig Jahre alt wird, und heraus kommt zu seinen Brüdern, und ansähet zu richten, urtheilen, verdammen und todtschlagen, da ist es mit Mose aus, und hat ein Ende mit ihm; dann wollen sie ihn todtschlagen. Das ist, wenn GOttes Wort klein und jung ist, so wird es angenommen; aber wenn es wächst und recht erkannt wird, daß es uns straft, und tadelt unser böses und ärgerliches Leben, zeigt uns auch unsere Sünde an, dann wird man ihm feind. Das heißt, Mosen tödten. Denn das Wort GOttes tödtet, und zeigt die Sünde an. Das Gesetz ist darum gegeben, daß die Sünder daraus erkennen sollen, wo es ihnen fehlt, wie St. Paulus [Röm. 3, 20.] spricht: Per legem cognitio peccati est [durch das Gesetz kommt Erkenntniß der Sünde]. So singt man auch: „Gesetz ist ein Spiegel zart, der uns anzeigt die sündge Art, in unserm Fleisch verborgen.“ Die Leute werden aus dem Gesetze gewahr, daß sie vor GOtt gar todt sind mit ihren guten Werken. Derhalben so thut Moses keinen Schaden, er sei denn zuvor groß und alt worden.

43. Also soll man es allezeit darauf ziehen, daß des Gesetzes Amt und Art sei, Sünde anzuzeigen und tödten; wie St. Paulus [1 Cor. 15, 56.] auch sagt: Virtus peccati lex [die Kraft der Sünde ist das Gesetz]. Es thut das Gesetz anders nichts, denn daß es Sünde erregt; da wird man denn bald dem Gesetze feind und gram. Denn niemand will erkannt werden, daß er unrecht und sträflich sei; welches denn das Gesetz nicht leiden kann. Darüber hebt sich denn der Hader, Widerwille, Groll und Uneinigkeit zwischen den Menschen und dem Gesetz oder den Gesetzpredigern.

\section{Das dritte Capitel}

\begin{schriftzitat}
V. 1. Moses aber hütete der Schafe Jethro, seines Schwähers, des Priesters in Midian.
\end{schriftzitat}

1. Das dritte Capitel lehrt, wie Moses in der Wüste umher gezogen sei, als er Egyptenland verlassen, darinnen habe er einen Busch mit Feuer gesehen, daraus GOtt mit ihm redet, und ihm befiehlt, das Volk Israel aus Egypten zu führen. Daß also das Hauptstück ist in diesem dritten und auch nachfolgendem vierten Capitel vom Berufe Mosis, daß er der Kinder von Israel Herzog und Heerführer sein solle, und sie aus der Tyrannei Pharaonis ledig und los machen. Auf daß wir sehen mögen, was es thue, wenn GOtt einen beruft, oder wenn er einen nicht beruft.

2. Dagegen hören wir auch, wie Moses sich lange zerrt und wehrt, und will den Beruf nicht haben. Denn Moses mit Gott gleichsam einen Zank und Hader anfähet, und will dies Amt nicht annehmen; GOtt fordert und berufet ihn fast zum sechsten Mal, noch sperret und widert er sich immerdar.

3. Moses macht sich gar zu klein und geringe; sieht dagegen die Sache, nämlich die Ausführung der Kinder von Israel aus dem Diensthause Egypti, also groß an, daß es ihm werde unmöglich sein, und daß er hieher zu rechnen gar nichts sei; denn da war keine menschliche Hülfe zu sehen. Und obnohl der tyrannische und böse König in Egypten stirbt, der die Kinder Israel mit Frohndiensten beschwert und ihre Kindlein ersäuft hatte, und der Moses so feind war; da man denn Hoffnung gehabt, es sollte nun besser, oder, wie man sagt, um einen Rock wärmer werden, so wird es noch viel ärger. Denn es kommt ein ärgerer König in das Regiment, denn der vorige gewesen war. So war Moses auch nicht mehr am königlichen Hofe, sondern aus Egypten flüchtig, und lebte in Midian bei vierzig Jahren im Exilio. Derhalben so will Moses nicht hinan an den Beruf; er will nicht fort, ob GOtt ihm wohl verhieß, daß er ihm Wort und Mund dazu geben wollte. [Er] treibt und reizt GOtt
\quellseite{367}
auf sich, daß er auf Mosen gar zornig wird, und spricht: Bringe mich nicht auf, fürchte dich vor meinem Zorn.

4. Dieses, von dem Berufe Mosis, behandelt wohl; denn niemand soll sich in ein öffentlich Amt ohne GOttes Beruf eindringen. Und wisset solches um der neuen Flattergeister willen, die sich eindringen und einschleichen, da sie doch GOtt nicht dazu berufen noch geschickt hat; wollen Prediger sein, und die Leute lehren, ohne Erforderung und Beruf [Jer. 23, 21.]. Moses hatte in seinem Geiste vorlängst seinen Brüdern gewünscht die Erlösung aus dem Diensthause Egypti, wie er denn droben Geistes und Muths genug hatte, da er den Egypter todtschlug. Als sollte er sagen: Mich dünkt es, ich sollte euch helfen. Und GOtt straft es nicht; wie denn Stephanus solchen Mord preist in den Geschichten der Apostel [Cap. 7, 25.] und spricht: „Sie vernahmen es nicht.“ Noch will er solche Ausführung aus Egypten nicht angreifen, GOtt berufe ihn denn zuvor und schicke ihn zu den Kindern von Israel.

5. Denn es ist nicht genug, daß man sich des Geistes rühme, GOtt will es auch nicht haben, daß man denen glaube, die da vorgeben und sagen: Glaubet meinem Geiste. Nem, die da sagen: Der Geist treibt mich, der Geist heißt es mich. Sonst ständen wir alle gleich auf einem Haufen, und keiner hörte den andern.

6. Aber wo GOtt beruft und treibt zum Predigtamt, da geht denn das Werk von statten, und reißt hindurch. Wie denn Gott alhier ein großes Werk anrichtet, daran Moses selbst gar verzweifelt, daß er ein solch groß Volk aus einem so mächtigen Königreiche führen sollte, und spricht: Was bin ich gegen einem solchen König? Ich bin wahrlich ein feiner, hübscher, geschmückter Geselle. Aber es schadet nicht. Mose, spricht GOtt, sei du keck, bete und sei getrost, ich will den Pharao verstöckeln, daß er sich sperren soll; aber ich will der Zeichen desto mehr bei dir thun, auf daß du gewiß seiest, ich habe dich zu dem Werke berufen.

7. Denn wer was Neues anfangen soll, der bringe nicht mit ihm diesen Ruhm: der Geist habe es ihm in das Herz gegeben. Bringt er diesen Trost und Ruhm mit sich, so schlage man ihn nur aus. Denn also soll man erstlich thun und sagen: Will mich der heilige Geist zu einem Amte (es sei nun, was es wolle) gebrauchen, und will GOtt haben, daß ich es anfangen soll, und ich habe es im Sinne, so wird er es also geben, daß die Zeichen hernach folgen, daß er das Senden und Anheben durch himmlische Zeichen bestätige.

8. Denn das ist beschlossen, daß der Geist niemand wird das Herz aufblasen oder reizen, er bestätige es denn zuvor mit Zeichen; äußerliche Beweisung muß er mit sich bringen. Das ist die innerliche Berufung.

9. Die andere Berufung ist auch göttlich, aber sie geschieht durch Menschen. Die erste geschieht allein von GOtt, ohne einiges Mittel; und dieser Beruf muß äußerliche Zeichen und Zeugnisse haben. Der andere Beruf darf keiner Zeichen; als, ich predigte allhier zu Wittenberg nimmermehr, wenn ich von GOtt dazu nicht gezwungen und durch den Churfürsten zu Sachsen erfordert wäre, daß ich es thun müßte. Also ist es mit einem andern auch. Denn wenn mich die Leute zwingen und dringen wollen, und ich kann es thun, oder ich kann es gleich nicht thun, daß man von mir begehrt, so thue ich, so viel als ich kann. Da treibt er durch Menschen, und so steht auch GOttes Gebot da, daß mich der Heilige Geist auch beruft und spricht [3 Mos. 19, 18.]: Du sollst den Nächsten lieben als dich selbst. Es soll kein Mensch ihm selbst allein leben, sondern sollen dem Nächsten auch dienen. Dies Gebot ist über alle geschlagen, über mich und über dich. Wenn mich dasjenige Gebot ergreift und mir vorgehalten wird, so hilft kein Wehren, es wäre denn, daß ich mich so lange wehren wollte, bis ich drüber in GOttes Ungnade käme.

10. Dieser Beruf ist nun durch Menschen, und doch auch von GOtt bestätigt; darum gedenke, und diene dem Nächsten darinnen, sonst kommen andere über quer Feld einhergeplumpst, und dringen sich in Aemter, darein sie nicht berufen sind, auch nicht darum gebeten noch ersucht. Als, wenn ich wollte gen Leipzig gehen, und daselbst auf dem Schloße predigen, und sagen: Der Geist treibt mich, ich muß predigen; das ist unrecht.

11. Wenn er dir nicht Siegel und Zeichen gibt, so enthalte dich deß\footnote{Eislebensche: dichs = dich deß.} und bleibe in der gemeinen Weise, daß einer den andern hin und wieder ersuche. Als, an diesem Orte allhier,
\quellseite{368}
Da ist dieser Predigtstuhl unser aller, wir treten zusammen; wer aber zum Predigtamte hier nicht erwählt ist, der steht ab, und predigt nicht. Wer heißt nun die Rottengeister diese Ordnung brechen?

12. Diese beiderlei Berufe sind nun also gethan. Einer, als der schlecht von GOtt geschieht; und der sich hierauf beruft, den nimm nicht an, er thue dir denn Zeichen; oder glaube ihm nicht, er bringe dir denn zuvor Wahrzeichen, daß er vom Heiligen Geiste berufen sei. Denn solcher Schleicher und Winkelprediger findet man heutiges Tages viel, die da sagen, sie sind von GOtt berufen; aber sie zeigen an GOtt oder Menschen, die sie berufen haben. Ja, sie können wohl hinter den Leuten herkommen und einherschleichen, und so lange machen, daß man sie hernach erwählt und beruft; man kann die Leute mit Worten bald dahin bereden. Aber es sind Diebe, Mörder und Wölfe [Joh. 10, 1.]. Nun so mangelt es ihnen daran, daß man frage: Kommt ihr vom Heiligen Geiste, so weiset Zeichen. Oder, rühmt ihr euch, daß ihr von Menschen berufen seid, so thut es dar und beweiset es, wie recht und der heiligen Schrift gemäß ist. Denn GOtt hat auch die Apostel in die Welt nicht schicken wollen, sie wären denn mit Zeichen zuvor bestätigt.

13. Der andere Beruf, so durch Menschen geschieht, ist zuvor bestätigt durch den Befehl GOttes auf dem Berge Sinai [3 Mos. 19, 18. 5 Mos. 6, 5.]: „Liebe GOtt, und den Nächsten als dich selbst.“ Wenn dich dies Gebot treibt, so bedarfst du keines Zeichens, denn GOtt hat es zuvor befohlen, und ich muß es thun. Nun nehmen die Leute dies Gebot, und halten mir es vor; diesen Spruch haben mir Moses und GOtt im Himmel bestätigt, wenn ich demselbigen folge. Also predige ich ohne alle Zeichen, und ist dennoch der Beruf GOttes; denn er geht aus dem Gebot der Liebe daher, und wird von GOtt gezwungen. Also sind auch viel Propheten und Leviten berufen worden, welcher natürlich Amt gewesen, daß sie predigen sollten, gleichwie unsere Doctores Theologiæ dazu gemacht und geweiht sind. Nun, diese durfte man nicht dazu berufen, denn sie waren Propheten und predigten ohne Zeichen, denn sie thaten es aus ihrem natürlichen Amte.

14. Aber dies ist mein Ruhm und Troß, daß ich weiß, daß ich recht thue, und durch das Gebot der Liebe dazu berufen bin. Das ist mein Ruhm. Ich rühme mich nicht meines Geistes, denn da könnte ich nicht ein Theil meines Amtes thun, wenn ich auch gleich wüßte, daß ich die ganze Welt bekehren sollte, sondern da liegt GOttes Wort und Gebot. Und wenn es nach GOttes Wort geht, so ist es recht; wenn es aber nach meinem Geiste geht, so liegt’s bald in der Asche alles, was ich anfange. Also treibt man mich mit Gewalt hinein, und zwingt mich, daß ich predige, wiewohl ich gerne außerhalb dem Predigtamt wäre geblieben. Deß habe ich mich auch gerühmt, da vor Zeiten der Papst meines Berufs Ursachen von mir wissen wollte.

15. Das ist der Beruf eines öffentlichen Amts unter den Christen. Wenn man aber unter den Haufen käme, da nicht Christen wären, da möchte man thun wie die Apostel, und nicht warten des Berufs. Denn man hat da nicht das Amt zu predigen; und einer spräche: Allhier sind nicht Christen, ich will predigen und sie unterrichten vom Christenthum, und es schlüge sich ein Haufe zusammen, erwählten und beriefen mich zu ihrem Bischofe, da hätte ich einen Beruf.

16. Moses ist allhier achtzig Jahre alt, und hat bei vierzig Jahren einen großen Geist, ist ein trefflicher Mann; noch sperrt und wehrt er sich, das öffentliche Amt zu führen, und will nicht hinan, bis daß er hört seines GOttes Ungnade, ob er wohl treffliche, herrliche Zusagung hatte, daß GOtt wollte bei ihm sein. Denn wo GOttes Wort bestätigt, da ist sein Gebot, und da sind die Zeichen auch bestätigt. GOtt thut kein Zeichen, wo es nicht hoch vonnöthen ist. Christus wollte den Juden kein Zeichen nach ihrem Kopf und Gefallen thun [Matth. 12, 39.], aber wo es vonnöthen war, da konnte er wohl Mirakel und Wunderwerke thun. Denn sonst wäre es ein Gaukelwerk, wie denn Herodes auch von Christo Zeichen haben wollte [Luc. 23, 8. 9.]. Wir könnten auch Zeichen thun, aber es ist ohne Noth, denn unsere Lehre ist zuvor bestätigt, und ist keine neue Lehre mehr.

17. Das ist das erste Stück dieses dritten Capitels, von dem Berufe Mosis, da er das große Volk, die Kinder von Israel, soll wider des gewaltigen Königes Pharaonis Willen aus Egypten führen. Und haben gehört von zweierlei Beruf; einem, der da ohne Mittel von GOtt geschieht, wenn der Heilige Geist beruft und in das Herz gibt. Demselbigen glaube nicht, sie
\quellseite{369}
geben dir denn ein Zeichen, daß ihnen GOtt ein solches Amt befohlen habe, da niemand auch um weiß, denn er allein. Und wen GOtt berufen hat, dem gibt er auch ein Zeichen, wie er alhier Mosi thut. Danach, so ist der andere Beruf, da einer durch das gemeine Gebot berufen wird, nämlich, daß man den Menschen dienen solle. Da hat man in diesem Berufe auch GOttes Wort, und ist da des Zeichens nicht vonnöthen, daß das Evangelium und das Gesetz recht sei. Denn beide Lehren vor vielen Jahren durch Zeichen bestätigt, und auch von der Welt angenommen sind, also, daß man der Welt das Maul wohl könne stopfen; denn jedermann weiß, daß solche Lehre recht ist, und man muß damit zufrieden sein.

\begin{schriftzitat}
V. 1. Moses aber hütete der Schafe Jethro, seines Schwähers, des Priesters in Midian, und trieb die Schafe enthinter in die Wüste, und kam an den Berg GOttes, Horeb.
\end{schriftzitat}

18. Horeb und Sina ist ein Berg in Arabia; und [es] liegt Egypten gegen Abend, das gelobte Land\footnote{In der Gießlebenschen Ausgabe: „dem gelobten Lande“. Dadurch werden die Himmelsrichtungen gerade entgegengesetzt angegeben. Der Fehler wird ohne Zweifel bei dem Abschreiber zu suchen sein.} aber gegen Morgen, und hat daran gestoßen Arabia Petraea. Die Edomiter oder Ismaeliter haben daselbst gewohnt, wiewohl sie die Schrift mit andern Namen nennt. Daselbst ist Moses ein Knecht gewesen, und ist vierzig Jahre alt. Er kommt aus einem königlichen Hoflager, da er herrlich gehalten gewesen ist, aber er will lieber mit seinem hohen Geiste in Armuth mit seinen Brüdern sein, denn in großen Ehren an des Königs Hofe schweben. Zuvor war er des Königs Sohn, alhier ist er vierzig Jahre ein Knecht, und hütet der Schafe. Das ist ein großer Sprung, von eines großen Königes Saale und Hofe in ein Hirtenhaus kommen.

19. Midian ist ein Priester gewesen, und ein rechtschaffener Mann. Denn dieselbigen Länder sind bessere Christen gewesen, denn wir sind; sie haben nicht mit Menschenjagungen genarret, wie wir thun. Midian ist Abrahams Sohn gewesen, von der Ketura, Abrahams drittem Weibe, geboren [1 Mos. 25, 2.], und da hat es nicht gefehlt, sie haben von Abraham viel gelernt, wie er in seinem Hause gewandelt habe. Sie haben alle GOtt recht erkannt, und den rechten Gottesdienst gelernt, wie auch die Ismaeliten. Aber das ist dagegen der Mangel, daß der Glaube nicht allen folgt, die GOttes Wort wissen und haben. „Denn der Glaube ist nicht aller“ [2 Thess. 3, 2.]. Wo denn der Glaube gefallen ist, so bleibt gewiß eine Heuchelei, da man zusammen backt die Vernunft und das göttliche Wort, auf daß Ein Kuche draus werde; wie sie denn im Papstthume das Sacrament des Abendmahls, die Taufe und das göttliche Wort mit Menschensatzungen gemengt haben.

20. Nun ist Midian ein solcher Heide nicht gewesen, der da Abgötter angebetet hätte, sondern er hat den rechten GOtt gelehrt, wiewohl er eine sonderliche Weise erfunden hat, welches wir jetzt Abgötterei nennen; wie denn allezeit die Vernunft thut. Also ist Laban auch rechtschaffen gewesen, und doch darnach zu einem Götzendienst abgerissen worden [1 Mos. 31, 30.]. Also ist nun Moses bei den Heiden geblieben, und hat ihnen gepredigt, wie Joseph in Egypten gepredigt hat [1 Mos. 39, 9. ff.].

\begin{schriftzitat}
V. 2. 3. Und der Engel des HErrn erschien ihm in einer feurigen Flamme aus dem Busche, und er sahe, daß der Busch mit Feuer brannte, und ward doch nicht verzehret, und sprach: Ich will dahin, und besehen dies große Gesicht, warum der Busch nicht verbrennet.
\end{schriftzitat}

21. Allhier ist eine Frage: ob GOtt oder ein Engel mit Mose geredet habe? Denn Apost. 7, 30. sagt Stephanus auch, daß über vierzig Jahre Mosi in der Wüste auf dem Berge Sinai der Engel des HErrn in einer Feuerflamme des Busches erschienen sei. Denn etliche sagen, daß Moses allhier das göttliche Wesen selbst gesehen habe. Nun steht da geschrieben, der Engel des HErrn sei Mosi erschienen, habe mit ihm aus dem Busche geredet. Moses hat GOtt allhier nicht gesehen. Darum spricht auch Lucas und Stephanus, daß das Gesetz durch die Engel gegeben sei etc. Diese Himmelsfürsten macht GOtt zu Predigern, daß es also durch die Handreichung der Engel zu uns kommt; aber wir Christen haben den HErrn selbst gehört. Christus, GOttes Sohn, so über alle Engel ist, hat mit seinem Munde uns gepredigt, und das Evangelium in der Welt auszubreiten angefangen, und es darnach seinen Aposteln befohlen zu predigen. Er hat das Evangelium in seiner Menschheit ihm vorbehalten zu predigen, und das Gesetz durch
\quellseite{370}
die Engel geben lassen [Gal. 3, 19.]. Darum soll man allezeit verstehen, daß GOtt durch Engel redet.

22. Es ist ein Engel in diesem Busche, der führt zu Mose GOttes Wort, und redet als GOtt; und Moses hört auch GOttes Wort, es gehe gleich aus einem Busche oder durch einen Stein. Das Wort soll man ansehen, und nicht auf die Person schauen. Denn GOtt nimmt jetzt einen Engel, bald Petrum oder Magdalenen, oder auch irgendes wohl einen Esel, wie mit dem Bileam geschah [4 Mos. 22, 28. 30.], durch welche er sein Wort redet. Darum so wird recht gesagt: „GOtt redete mit Mose“, wiewohl es nicht allezeit gemeldet wird, daß er durch die Engel geredet habe. Denn wenn GOtt für sich redet sein Wort, so redet er es, daß wir alle Augenblicke selig sind. Wenn er unmündliche Dinge redete, so müßten wir flugs vor der großen Herrlichkeit zerschmelzen und vergehen; aber wenn er nicht leiblich redet, so gebraucht er Creaturen dazu, wer auch dieselbigen Creaturen sein mögen.

23. So hat nun GOtt allhier durch einen Engel geredet. Matthäi am 19. Capitel, V. 5., steht geschrieben, daß GOtt gesprochen habe: „Der Mensch wird Vater und Mutter verlassen, und an seinem Weibe hangen“, da es doch Adam sprach [1 Mos. 2, 24.]. Aber GOttes Sprechen und Adams Sprechen ist einerlei. Denn es ist aus GOtt und in GOtt geredet, darum so hat es GOtt gesprochen.

\begin{schriftzitat}
V. 4. 5. Du aber der HErr sahe, daß er hinging zu sehen, rief ihm GOtt aus dem Busche und sprach: Mose, Mose! Er antwortete: Hier bin ich. Er sprach: Tritt nicht herzu, zieh deine Schuhe aus von deinen Füßen, denn der Ort, da du auf stehest, ist ein heilig Land.
\end{schriftzitat}

24. Also thut nun hier GOtt mit Mose, wie er allezeit mit seinen Heiligen pflegt zu thun, daß er zuvor einen gar zunichte macht, ehe er ihn erhebt. Zuvor zerbricht und zerreißt GOtt alles, ehe denn er es bauet; er wirft einen vor zu Boden, ehe er ihn wieder aufrichtet; er tödtet, ehe er lebendig macht. Solches kehret der Teufel gar um. Aber Moses erschrickt, wird blöde, und fällt dahin alle seine Vermessenheit. Es muß aber solches also geschehen um seines äußerlichen Berufs willen, zu Ehren der Stätte, die da heilig gewesen ist.

25. Woher ist sie heilig? Heilig wird sie daher genannt, nicht, daß sie mit Oel gesalbt und geheiliget wäre, sondern daß GOtt sich allda dem Mose offenbart und mit ihm redet. GOttes Wort ist daselbst, darum so muß auch das Volk, so da wohnt, heilig sein. Ob gleich gottlose Schäfle daselbst sind, so ist doch der Ort selig. Und daher wird auch das Volk selig, und die Ohren und Augen selig, die da hören und sehen das göttliche Wort; wie der HErr Christus sagt [Luc. 10, 23. 24.]. Denn wo GOttes Wort hinkommt, da hört und sieht man heilige Dinge; denn das Wort ist heilig und macht heilige Leute. Darum, wo das Wort hin erschallt, da ist eine heilige Statt.

26. Derhalben vermiß dich nicht, als dürfest du dich nicht demüthigen, sondern mit ganzer Furcht soll man herzugehen. Denn da ist nicht ein Busch oder ein Feuer, wie andere Büsche und Feuer, sondern da ist etwas Neues, nämlich GOttes Wort, das will GOtt in Ehren gehalten haben. Und mit Zittern und Furcht muß man herzugehen, und es nicht achten als ein menschliches Wort. Da schlägt er Mosen gar darnieder, weil Gott mit seinen Engeln und mit seinem Worte allhier gegenwärtig ist. Da muß der Mensch sagen: HErr, ich bin nichts; und bald fällt der Muth. Wenn GOtt also einem das Herz rührt, dem fallen Fittig und Federn bald aus, daß er nicht hoch fliege; wie allhier Moses nicht mehr will in den Busch sehen, denn er ist in eine Furcht gefallen.

27. Darnach kommt GOtt, und tröstet den Mosen wieder. Denn er will an dem Berge Horeb anfangen die Erlösung der Kinder von Israel, und wird Moses abgefertigt als zum Hauptmann solches Werks; redet mit Mose davon. So will er hernachher an diesem Orte sein Gesetz geben, und ein eigen Reich und Volk aus den Kindern von Israel machen. Darum, weil GOttes Wort und Gesetz allhier ist, so muß die Stätte oder der Ort heilig sein.

28. Also nennt der Patriarch Jacob Bethel auch eine heilige Stätte, denn er sahe daselbst im Traume die Leiter, so vom Himmel auf die Erde reichte, und GOtt oben auf der Leiter stand, und die Engel darauf vom Himmel auf die Erde stiegen [1 Mos. 28, 12.]. GOtt redete da mit ihm, der Himmel war da offen. Dergleichen wird Jerusalem auch eine heilige Stadt genannt [Matth. 27, 53.], denn da war der
\quellseite{371}
Tempel, das Gesetz, der Propheten Lehre und Verheißung von Christo. Wie denn noch alle unsere Kirchen auch heilig sind darum, daß GOttes Wort darinnen gepredigt und die Sacramente gereicht werden.

29. Allhier soll man auch wohl merken, daß GOtt in diesem Texte gar reich und überflüssig mit Worten ist, und daß er gar viel mit Mose redet. Denn ich habe es zuvor oft gesagt, und sage es noch, daß es ein theuer und köstlich Ding ist, wenn man GOttes Wort hört. Und man sollte alle Lande durchlaufen, daß man würdig sein möchte, einen Buchstaben von GOttes Wort zu hören: vielmehr sollte man die Hände aufheben, Haus und Hof aufthun, daß GOtt so viel mit uns redet, wie er allhier mit Mose auch thut.

30. Daß aber der Busch voll Feuers ist und nicht verbrennt, und Moses die Schuhe ausziehen muß etc., was das sei, wollen wir sparen zu einer andern Zeit. Es geschieht aber darum, daß Moses gedemüthigt vor GOtt käme, in Erkenntniß seiner Nichtigkeit, und gar ein anderer Mann würde, als durch welch Instrument GOtt große Dinge ausrichten wollte. Darum hat er müssen den Busch brennen sehen und erschrecken, und in eine andere Haut kriechen.

\begin{schriftzitat}
V. 6. Und sprach weiter: Ich bin der GOtt deines Vaters, der GOtt Abrahams, der GOtt Isaaks, und der GOtt Jakobs. Und Mose verhüllte sein Angesicht, denn er fürchtete sich, GOtt anzuschauen.
\end{schriftzitat}

31. GOtt rühmt allhier, daß dies sein Name ewiglich sei, dabei soll man ihn nennen, und wer diesen GOtt nicht haben will, der fehlt des rechten GOttes. Der HErr Christus führt diesen Text ein wider die Sadducäer, Matthäi am 22. Capitel, V. 32. [Marc. 12, 26. Luc. 20, 37.], und legt diesen Text kein aus von der Auferstehung der Todten. Denn sie hielten es dafür, es wäre wohl ein Gott, aber\footnote{Eislebensche: oder.} kein ander Leben nach diesem Leben, und auch kein Engel noch Geist, und verstanden aus Gottes Wort Mosis nur nach dem äußerlichen Wesen, wie es die Juden auch nicht anders ansahen. Aber der HErr Christus sagt [Matth. 22, 29.] zu den Sadducäern: „Ihr irret, und verstehet die Schrift nicht, noch die Kraft GOttes“; ihr wisset nicht, was die Auferstehung der Todten sei.

32. Aber allhier wird es angezeigt, daß GOtt spricht: „Ich bin ein GOtt Abrahams, Isaaks und Jakobs.“ Daraus schließt Christus, daß ein ander Leben nach diesem Leben sei. Und daß die Auferstehung der Todten in diesem Texte begriffen werde, beweist er also, daß GOtt ein GOtt der Lebendigen und nicht der Todten sei; denn dasjenige, so nichts ist, item, zerstoben und zerflogen ist, das hat keinen GOtt. Soll ein GOtt sein, so muß er jemandes GOtt sein, von dem ausfließt ewige Güte gegen diejenigen, in welchen er ist.

33. Und alle Heiden haben sich auch Gutes gegen GOtt versehen, darum hat einer diesen, der andere jenen Gott erdacht, um seiner andern Ursache willen, denn, wenn Unglück ihnen zustände, daß man allhier Rath und Hülfe finde. Also haben die Heiden viel Abgötter gemacht, denn es ist uns natürlich eingepflanzt, daß GOtt ein solches Ding sei, so helfe allen denen, die ihn anrufen. Daran haben die Heiden nicht gefehlt, daß GOtt etwas sei, so da helfe; dies Licht steckt noch in der Natur, sonst sagten die Leute nicht: Hilf mir. Sie müssen eine Erkenntniß von GOtt haben. Aber da fehlt es an, und allhier ist die Vernunft blind, und richtet Abgötterei an, daß sie die Gottheit andern Dingen zuschreibt, so nicht GOtt sind, und den rechten GOtt nicht erkennt.

34. Die Natur weiß, daß GOtt sei, der da helfe; aber wer derselbige GOtt sei, das weiß sie nicht. Darum träumt sie ihr, dieser wird GOtt sein, und also wird es um das göttliche Wesen geschaffen sein. Es war der HErr Christus auf Erden und that GOttes Werk, und redete GOttes Wort; noch konnten ihn die Juden selbst mit ihrer Vernunft nicht erkennen, daß er Gott wäre, denn sie doch vor Augen mit seinen Wunderwerken sahen [Joh. 10, 33.], sondern sprachen: GOtt ist, der alles erschaffen hat, allen hilft, und alles Gutes thut etc.; aber daß dieser GOtt sollte sein, das ist er nicht. Denn wenn es in das Leuten kommt, so ist GOttes bald gefehlt. Man weiß wohl, daß ein GOtt sei; aber wer derselbige GOtt sei, da ist die ganze Welt gar blind.

35. Nun, daß ich es klar mache: Wir Mönche haben hoch von GOtt gepredigt, und wie man GOtt dienen solle, nämlich, daß dieser GOtt sei, der hoch im Himmel sitze und alles geschaffen hat, und seinen Sohn in die Welt geschickt etc. Da hat
\quellseite{372}
man gesagt: Diesem GOtt will ich dienen, dies und jenes thun, und geloben Armuth, Keuschheit und Gehorsam. Dies Werk wird ihm wohlgefallen. Also deuten sie GOtt, und sagen, was sein Wille sei. Wenn ich also GOtt vorschlage, wie ich ihm dienen will, da ist es vergeben, und ich habe des rechten GOttes gefehlt. Denn GOtt ist nicht also gesinnt, wie ich gedenke, sondern was GOtt gebeut und heißt, das soll man thun. Also will er getroffen sein, und nicht von meinem Gutdünken. Derhalben läßt er diese alle fehlen, irren und blind bleiben, die GOtt also gedeutet haben.

36. Es ist eine gemeine Erkenntniß GOttes (wie Röm. 1, 19. ff. und Cap. 2, 1. ff. gesagt wird) in der Welt geblieben. Aber wenn es zum Treffen kommt, so fehlen sie alle GOttes. Darum gebührt niemand, GOtt zu deuten oder zu nennen, denn wie er sich selbst nennt, sich selbst deutet und sein Werk vorschlägt. Da sein Deuten währet, da ist es, sonst fehlen wir alle. So ist nun besser, von allen Gottesdiensten abstehen, wo man nicht gewiß ist, daß GOttes Wort da ist, und da sich GOtt nicht selbst deutet. Denn, daß eine Nonne spricht, sie sei Christi Braut, daran ist nichts, denn sie ist des Teufels Braut, darum, daß sie Dinge vor sich nimmt aus eigener Andacht und Gutdünkel, und deutet ihr Ding auf GOttes Werk. Sondern also sagt man recht: GOtt, der Himmel und Erde geschaffen hat, spricht: Wer glaubt an Christum, soll das ewige Leben haben, und dem sollen die Sünden vergeben sein und soll selig werden. Da treffen wir GOtt. Denn ich habe das Deuten (als, an Christum glauben) nicht gemacht. Denn von mir hätte ich also gesagt: Ich will so lange beten, fasten, mich martern und die Heiligen anrufen, bis daß ich selig werde. Solch Deuten ist, unsern HErrn GOtt nennen, wie ich will; das heißt denn nichts.

37. GOtt wird wohl allezeit also erkannt, daß er ein solch Ding sei, so allen helfe; aber niemand kann ihn nennen, wenn er sich nicht selbst deutet und Zeichen steckt, und spricht: An dem Namen, Zeichen oder Werke sollst du mich treffen und kennen. Gebe ich aber GOtt einen Namen, und stecke ihm ein Zeichen, so habe ich sein gefehlt; denn GOtt hat uns seinen Namen und Zeichen gesteckt an Christo. Christus soll sein Name sein. Wer durch den Glauben an Christum sich hängt, und getauft wird, der hat GOtt gefunden. Nimmst du aber etwas Anderes vor, so fehlst du GOttes. Als, wenn du sagen wolltest: Lieber HErr GOtt, der du sitzest im Himmel unter den Engeln, ich komme zu dir, und will dir dienen und ein Kloster bauen, eine Mönchskappe anziehen, Wallfahrt laufen, dies und jenes thun etc. Hiermit deute ich ihn, und nenne ihn ohne seinen Befehl und Wort. Da habe ich gefehlt.

38. Also wird der Artikel von [der] Auferstehung der Todten allhier auch angezeigt; denn GOtt spricht: „Ich bin ein GOtt Abrahams, Isaaks und Jakobs.“ So er denn nun ein GOtt ist des Abraham, so muß Abraham leben. Ist er auch ein GOtt Isaaks, so muß Isaak leben. Derhalben so muß man alles, was da Abraham heißt, bleiben lassen, nicht allein die Seele, sondern auch den ganzen Mann mit Leib und Seele, der Sara, die alte Rachel,\footnote{So wird Sara wahrscheinlich genannt, weil sie, ebenso wie Rahel (\hebraeisch{רחל}) [1 Mos. 29, 31.], unfruchtbar war.} zum Weibe gehabt; derselbige leibliche Mann, so einen Sohn mit Sara gezeuget hat. Demnach muß der Mann Isaak auch leben. Dies sind starke, treffliche Worte, daß Isaak, Abraham, Jakob wieder leben und auferstehen sollen, nicht allein mit dem Leibe, sondern der ganze Mensch mit Leib und Seele.

39. Wie kann denn das möglich sein? ist doch Abraham, Isaak, Jakob gestorben, verfault und verwest, und scheint gar nichts an ihnen? Darauf antwortet Christus, und spricht: daß sein Wort wahr sei; und der Text der heiligen Schrift lüge nicht, daß die Todten leben müssen. Denn, ob sie gleich gestorben sind, so leben sie doch noch; denn im Tode leben sie auch. Denn unserm HErrn GOtt leben sie alle, uns Menschen aber sind sie nur allein aus den Augen entzogen. Gleichwohl leben sie noch, und ist das noch dahinten, daß sie am jüngsten Tage sollen wieder hervorgebracht und uns gezeiget werden. Sie leben aber in GOtt, und was vor ihm lebt, das ist ein bereitetes Leben; da ist es um einen Augenblick zu thun, so werden wir auch wieder leben. Daher spricht St. Paulus [1 Cor. 15, 54.] aus dem Propheten Hosea\footnote{Statt „Hosea“ sollte wohl „Jesaia“ gelesen werden, denn die Worte sind, wie sie hier stehen, 1 Cor. 15, 55. aus Jesaia angezogen.} [Cap. 13, 14.]: der Tod sei im Siege verschlungen etc. [Jes. 25, 8. 1 Cor. 15, 55.]

40. Es sind die Patriarchen darum gestorben,
\quellseite{373}
daß ihr sterblicher Leib nur verändert werde. Es gilt nur eine Veränderung. GOtt macht jetzt dran im Tode an dem sterblichen Leibe, daß er wieder so helle und schön werden soll, als jetzt die Sonne am Himmel ist. Es ist kein Tod vor unserm HErrn GOtt, und vor unsern fleischlichen Augen soll das Sterben nur ein Wegziehen sein. Denn gleichwie es mit unserer Geburt zugeht, also ist es auch um unser Sterben gethan. Was bist du und ich vor hundert Jahren gewesen? noch hat Gott uns hernach geschaffen und gemacht. Also kann er derer Leib, Fleisch und Blut auch wohl wieder hervorbringen, so vor hundert Jahren gestorben sind. Denn vor ihm lebt alles, als, diejenigen, so noch geboren werden sollen, und die, so gewesen sind.

41. Also beschließet er, daß eine Auferstehung der Todten sei. Denn Abraham, Isaac, Jakob sind nicht gestorben, sondern leben. Sind sie doch begraben? Ja, sie sind hinweg; aber es soll nur eine Veränderung sein, da die Verstorbenen aus unsern Augen weggenommen sind, und GOtt macht sie jetzt anders, denn sie sollen vor ihm leben.

\begin{schriftzitat}
V. 7–10. Und der HErr sprach: Ich habe gesehen das Elend meines Volks in Egypten, und habe ihr Geschrei gehöret über die, so sie treiben; ich habe ihr Leid erfahren, und bin hernieder gefahren, daß ich sie errette von der Egypter Gewalt, und sie ausführe aus diesem Lande, in ein gut und weit Land, in ein Land, darinnen Milch und Honig fleußt, nämlich an den Ort der Canaaniter, Hethiter, Amoriter, Pheresiter, Heviter und Jebusiter. Weil denn nun das Geschrei der Kinder Israel vor mich gekommen ist, und ich habe ihre Angst auch dazu gesehen, wie sie die Egypter ängsten; so gehe nun hin, ich will dich zu Pharao senden, daß du mein Volk, die Kinder Israel, aus Egypten führest.
\end{schriftzitat}

42. Dieses ist der Beruf Mosis, da GOtt ihn in der Wüste und hinter den Schafen in Midian holt, und zum Herrn oder Herzog über sein Volk setzt, gleichwie er den Saul von den Eseln nimmt, und zum König über Israel ordnet [1 Sam. 9, 2. ff.], ja, wie er den David hinter den Lämmern findet, und an Sauls Statt zum Könige salben und bestätigen läßt [1 Sam. 16, 11. ff.]. Denn GOtt weiß die Seinen wohl zu finden, die er zu großen Aemtern brauchen will; sie dürfen nicht darnach rennen und laufen, oder viel Practiken und Anschläge gebrauchen, GOtt sieht sie selbst und stößt sie aus in seinen göttlichen Beruf.

43. Lernet auch aus diesem Texte, wie GOtt auf seiner Heiligen Elend, Noth, Anfechtung, Trübsal oder Kreuz ein Auge habe, darauf sehe, ihr Geschrei, Seufzen und Gebet erhöre, und sie erretten wolle. Denn es jammert ihn der Seinen Trübsal, Angst und Herzeleid, da sie innen stecken; er hat ein väterlich Mitleiden und Erbarmen darüber, und will ihnen daraus helfen, wie der 103. Psalm, V. 13., auch sagt: „Gleichwie sich ein Vater erbarmet über seine Kinder, also will sich der HErr erbarmen über die, so ihn fürchten.“ So sagt sonst ein anderer Psalm [Ps. 56, 9.]: „GOtt zähle alle die Flucht der Seinen, und sammle ihre Thränen in seinen Sack“ etc. Denn allhier hilft er endlich seinem Volke aus der Tyrannei Pharaonis, und bringt sie in das gelobte Land, das er ihren Großvätern, Abraham, Isaac und Jakob, verheißen hatte. GOtt ist wahrhaftig in seinen Zusagungen; er hat Hülfe verheißen den Kindern Israel, und dieselbige setzt er nun ins Werk. Nun, wie hält sich Mosis gegen diesem göttlichen Beruf?

\begin{schriftzitat}
V. 11. 12. Mose aber sprach zu GOtt: Wer bin ich, daß ich zu Pharao gehe, und führe die Kinder Israel aus Egypten? Er sprach: Ich will mit dir sein; und das soll dir das Zeichen sein, daß ich dich gesandt habe: Wenn du mein Volk aus Egypten geführet hast, werdet ihr GOtt opfern auf diesem Berge.
\end{schriftzitat}

44. GOtt hat ein seltsam Mirakel und Wunderwerk allhier gethan, daß er Mosi in einem feurigen Busche erscheint, und der Busch doch nicht versehret noch verbrannt wird, damit er Mosen versichere, er wolle bei ihm sein, und sein Beruf sei ein göttlicher Beruf. Darum so gibt er ihm hiermit eine Verheißung, und spricht: er wolle bei ihm sein. Denn wo ein göttlicher, ordentlicher Beruf ist, da ist auch GOttes Zusagung, daß Glück und Segen solle bei demselbigen Berufe sein. GOtt will da auch sein mit seiner Hülfe, auf daß sich der göttlichen Gnade und Hülfe diejenigen trösten und erfreuen mögen, so da in ihrem Berufe allerlei Gefährlichkeit und Anstöße ausstehen müssen.

45. Als GOtt den Propheten Jeremiam [Jer. 1, 6–10.] in die Stadt Jerusalem schickte, und er sich entschuldigte, daß er jung wäre und nicht predigen könnte, gab ihm GOtt auch dergleichen
\quellseite{374}
Verheißung und sprach: „Du sollst nicht sagen, ich bin zu jung, sondern du sollst gehen, wohin ich dich sende, und predigen, was ich dich heiße. Fürchte dich nicht vor ihnen, denn ich bin bei dir, und will dich erretten. Ich lege meine Worte in deinen Mund. Ich setze dich heute dieses Tages über Völker und Königreiche, daß du ausreißen, zerbrechen, verstören und verderben sollst, und bauen und pflanzen.“

46. Wie hält sich nun Moses gegen diesem Beruf? Er sperrt sich und will nicht hinan, sieht erstlich seine Person an, und spricht: Wer bin ich, daß ich der Kinder Israel Hauptmann und Herzog sein sollte? Er bekennt seine Schwachheit und Unvermögen, spricht, er sei ein Hirte oder Schäfer, der vor dem Wolfe schwerlich Schafe könne vertheidigen; wie sollte er denn die Kinder Israel vor dem zornigen Tyrannen Pharaone können schützen? So habe er keine Kriegserrüstung, weder Spieß noch Schwert, Reiter oder Wagen, damit er sie könnte losmachen aus Egypten.

47. Darum so wehrt sich Moses. Wie wir denn hernach mehr hören werden, daß er zu GOtt spricht: Was soll ich von dir sagen? Wer bist du? Darnach wendet er noch mehr vor, und spricht: Sie werden mir es nicht glauben, und klagt sehr über des Volks Härigkeit oder Halsstarrigkeit; item, er spricht: er habe nicht eine gute Sprache; ist derhalben der Mangel an seiner Person, an unserm HErrn GOtt, an dem Könige, und an dem jüdischen Volke. Zuletzt, da er sieht, daß GOtt ihm ungnädig wird, spricht er [4, 13.]: „Sende, wen du willst“, und laß mich unverworren. [Er] will nicht fort.

48. Dies alles ist zu unserer Lehre und Trost geschrieben, daß an eine Sache, die GOtt betrifft, niemand eher gehen soll; er sei denn dazu berufen. Denn wer sich selbst eindringt, und hat kein Zeichen oder Wort von GOtt, daß GOtt ihn haben will, der wird wenig Gutes ausrichten.

\begin{schriftzitat}
V. 13. 14. Mose sprach zu GOtt: Siehe, wenn ich zu den Kindern Israel komme, und spreche zu ihnen: Der GOtt eurer Väter hat mich zu euch gesandt, und sie mir sagen werden: Wie heißt sein Name? Was soll ich ihnen sagen? GOtt sprach zu Mose: Ich werde sein, der ich sein werde. Und sprach: Also sollst du zu den Kindern Israel sagen: Ich werde es sein, der hat mich zu euch gesandt.
\end{schriftzitat}

49. Zuletzt spricht Moses: HErr, wie soll ich den Kindern von Israel sagen, daß du mich zu ihnen gesandt habest? Er soll GOtt deuten, wie er heiße, und er kann es nicht thun; wenn er sich es auch unterstände, so thäte er närrisch. Darum so fährt GOtt hervor und deutet seinen Namen selbst, und ist ein wunderlicher Name GOttes. Er deutet sich, und spricht: „Ich werde sein, der ich sein werde.“ Also heißt mein Name. Das ist blind Ding, daß er spricht: Ich bin, der ich bin, oder werde sein. Warum führt er diesen Namen? Ich habe zuvor gesagt, daß diesen Namen ihm niemand geben könne, wenn man GOtt deuten soll; denn da fehlen wir GOttes. Er rückt mit dem Worte über alle Creaturen, so nicht GOtt sind, noch das ewige Leben geben können; denn da ist sonst keiner, der da ewiglich wirkt, als GOtt allein; der sagt allhier: „Ich bin es.“ Das kann sonst niemand sagen, denn es vergeht alles.

50. All unser Leben ist eine fahrende, lose Habe. Ich kann wohl sagen: Ich fahre dahin; aber sagen kann ich nicht: Ich bin wesend, vergehe oder wanke nicht, ich bin und bleibe ewiglich. Mit welchem Titel reißt GOtt unsere Herzen und Augen von allen Creaturen, und zeucht es allein auf sich. Denn GOtt will allhier sagen: Es bleibt nichts auf Erden, denn ich allein. Ich habe das Wesen allein; wer andern Dingen anhängt, der fährt dahin. Heute ist der Mensch stark, morgen aber krank; heute bin ich in Ehren, morgen in Schanden; heute jung, morgen alt; heute reich, morgen arm. Darum spricht allhier GOtt: Ich bin der GOtt, an dem du hangen sollst, und sonst keiner Creatur vertrauen. Das ist das erste Gebot; wer dasselbige also führt und deutet, der deutet es recht.

51. Die Vernunft aber spricht also: Ich hänge an diesem Werke, als, daß ich so viel Messen habe gelesen, so lange im Kloster gelebt, gefastet, gebetet und so einen harten Orden geführt habe, ergreift also etwas Zeitliches und Vergängliches, damit wir uns trösten sollen, und daß wir durch Armuth gleich todbleich werden und sind vermattet. Gewöhnlich aber meinen wir, GOtt sehe auf uns, wenn es uns wohl geht. Und hiermit muß der edele, theure Name GOttes zu meinem Werke gesetzt werden; [ich] hänge mich an vergängliche Dinge, und laß den fahren, der unvergänglich und ewig ist. Derhalben so will mich GOtt gar bloß ausziehen
\quellseite{375}
, auf daß ich nicht vertraue auf meine Verdienste, Ehre, gute Werke, Fürsten oder Obrigkeit, oder wie es sonst einen Namen haben möge, sondern mich allein an den halte, der da ewig ist, und der allhier spricht: Hältst du dich an mich, so fehlt dir es nicht.

52. Wer hätte ihn also können deuten, wenn er sich nicht selbst also erklärt hätte, da er spricht: An mich muß man sich hängen durch den Glauben. Wenn man in der Vernunft steckt, so kommt man nicht zu GOtt, es sei denn, daß ich ihm einen Haufen guter Werke brächte, tröste zu ihm zu, und spräche: HErr, dies und das habe ich gethan. Das heißt denn, GOtt einen andern Namen geben, und GOttes Namen den Creaturen zueignen, und trotzen auf Creaturen, da wir doch auf ihn sollten trotzen. Also sollte man GOttes Namen deuten.

53. Ferner spricht er zu Mose: Ich will dir ein äußerlich Zeichen geben, daran du mich treffest: „Ich bin der GOtt Abrahams, Isaaks und Jakobs“ etc., davon soll man deuten, reden und predigen etc. Wie ist er denn Abrahams GOtt gewesen? Er gab ihm das Zeichen, die Beschneidung, und also glaubte Abraham GOtt; er gab ihm eine äußerliche Anleittung, wie man ihn finden sollte. Und das ist allein den Juden gesagt, und nicht den Heiden. Abraham hat das Zeugniß in der heiligen Schrift [1 Mos. 15, 6. Röm. 4, 3. Gal. 3, 6. Jac. 2, 23.], daß er GOtt geglaubt hat, und ist durch den Glauben gerechtfertigt worden; also ist GOtt des Abraham GOtt. Isaak hat diesem Glauben nachgefolgt, und zum Wahrzeichen die Beschneidung an seinem Leibe auch getragen. Also hat Isaak GOtt auch für seinen GOtt gehabt, und ist also GOtt durch den Glauben Abrahams GOtt worden. Denn er hat geglaubt, daß in seinem Samen gesegnet werden sollten alle Geschlechte auf Erden [1 Mos. 22, 18.]. Daher ist Abraham durch den Glauben GOttes Sohn worden [Joh. 1, 12.], und GOtt ist sein Vater worden. Daher spricht er allhier: Willst du mich recht deuten, lehren und predigen, so schlage kein Werk vor, sondern wie ich Abrahams GOtt worden bin, also heiße ich noch, und will dein GOtt auch sein. Wirst du nun eine andere Brücke zu mir machen, so fehlst du. Denn, ist er Abrahams GOtt, so soll ich Abrahams Wort, Predigt, Lehre, Religion, Glauben und Zeichen auch haben, oder komme nicht zu GOtt.

54. Die Juden sprechen: Wir glauben an GOtt, der Himmel und Erde geschaffen hat etc., und glauben an den GOtt Abrahams, Isaaks und Jakobs. Das geschieht aber allein mit dem bloßen Namen, und malen uns GOtt ab mit Buchstaben, und gehen mit Mosis Gesetzen und guten Werken um. Das ist weit gefehlt. Denn das ist mein Name, spricht GOtt allhier, auch ehe denn Moses ist geboren worden; aber ihr wollt mit Mosis guten Werken euch zu mir schicken, und den Himmel verdienen.

55. Nun spricht er allhier nicht: Ich bin Mosis oder der Leviten GOtt, sondern ein GOtt Abrahams, Isaaks und Jakobs; hebt also mit diesem Namen auf alle Gesetze und Werke Mosis, und bleibt allein dabei, daß niemand GOtt erlange, auch niemand zu ihm kommen könne, noch ihn\footnote{Eißlebenische: ihnen.} erkennen, er halte ihn denn für den GOtt, an welchen auch Abraham geglaubt hat. GOtt ist Abrahams GOtt gewesen ohne das Gesetz, und ehe denn das Gesetz gegeben ward; aber also ist er Abrahams GOtt worden, daß er geglaubt hat, daß GOtt durch seinen Samen der Welt helfen würde. In den Glauben mußt du auch kommen, und darnach das äußerliche Zeichen daran geheftet haben. Also will GOtt gepredigt und genannt sein in der Welt, das soll man ewiglich von ihm sagen [Röm. 1, 16.]. Wer GOtt anders sucht, der fehlt sein. Also soll ewiglich sein Name bleiben und gepredigt werden.

56. Nun, hier wird sich aber ein Spiel heben. Das Wort GOttes und der Glaube ist geblieben, und durch das Wort und den Glauben kommt man allein zu GOtt; aber wo nun hinaus? Wir nennen GOtt nicht einen GOtt Abrahams, so werden wir auch nicht beschnitten, und halten nicht, was Abraham hält, sondern kommen zu Christo, und halten uns an die Taufe; und dennoch wird allhier gesagt: Dies soll ewiglich mein Name sein? So werden wir Heiden müssen Juden werden, und uns beschneiden lassen, und unsere Taufe fahren lassen. Die Juden müssen wohl, daß ewig allhier nicht eine Zeit sei, die nimmermehr aufhört; wie denn auch gesagt wird: Kinder und Kindes Kind, das auch nicht eine Zeit ist, so nimmermehr aufhöret, sondern ist von einem zeitlichen Laufe geredet, von der Zeit, die auf Erden hier geht.

\quellseite{376}
57. Es ist ein Unterschied unter dem gepredigten Evangelio in der Welt, und dem Evangelio, so zuvor verheißen war [Röm. 1, 2. ff.]. Ein Ding, so verheißen worden ist, und so in das Werk gesetzt worden, ist zweierlei. Der Glaube, und das ewige Leben, sind auch zweierlei. Ich habe das ewige Leben in der Verheißung, und nicht im Wesen; jetzt habe ich es im Dunkel, und sehe es nicht, sondern glaube es; hernach werde ich es wohl fühlen [1 Cor. 13, 12. 1 Joh. 3, 2.]. Das Evangelium ist in die Welt jetzt erschollen und allenthalben verkündigt; aber in diesem Worte: „In deinem Samen sollen gesegnet werden“ etc., da ist es allein verheißen gewesen. Er spricht: In dieser Stunde soll das Evangelium gehen; und es geht. Wiederum auch: Es soll werden, es soll kommen, daß man von dem Segen und Benedeiung predigen wird, so weit die Welt ist; der gebenedeite Same wird von dir, Abraham, kommen, so da Christus heißt [Gal. 3, 29.]. Aber Christus ist damals noch nicht gekommen gewesen, sondern allein verheißen, wie denn auch das Evangelium verheißen gewesen.

58. Was nun Juden sind, aus Abrahams Samen, die werden an diesen Text gebunden, daß sie an den GOtt Abrahä glauben, der Abrahä die Verheißung gegeben hat. Aber es mangelt ihnen daran, daß sie es nicht dafür halten wollen, daß GOtt die Verheißung, dem Abrahä gethan, ausgerichtet und vollbracht habe. Aber GOtt spricht: Man soll mich ergreifen und halten in der Verheißung, Abrahä gethan. Das gilt nun nicht weiter, denn die Verheißung währt; darnach, wenn die Verheißung erfüllt ist, so muß es aus sein. Darum muß man nicht jetzt glauben, daß das Evangelium und Christus verheißen ist, sondern daß es nun alles geschehen und erfüllt ist, und daß GOtt wahrhaftig Abrahams GOtt ist, und diese Verheißung, Abrahä geschehen, nun erfüllt worden sei. Ich wollte denn ganz zurück springen und sagen, Christus wäre nicht gekommen.

59. Es ist es nun auf Christum gezielt, wenn der kommt, der soll den Namen ausrichten, wie Abraham versprochen ist [2 Cor. 5, 19.]. So geht es allein auf Christum; als sollte er sagen: Ich bin ein GOtt, der Abrahä den zukünftigen Christum mit seinem Evangelio verheißen hat. Wer mich nun haben will, der muß mich jetzt [als] einen solchen GOtt erkennen, als der da hat Abraham verheißen Christum und sein Evangelium. Hat er nun verheißen Christum und das Evangelium, so muß es also kommen; und wenn es kommt, so heißt denn GOtt derjenige, so da gibt und hält, was er zugesagt, und der sich uns verbindet, daß er will GOtt sein, und geben, was er verheißt. Darum so müssen alle seine Verheißungen erfüllt werden, und wenn die Verheißung geschieht, dann so ist es aus [2 Cor. 1, 20.]. Darum, nach Christi Geburt läuft seine Zeit mehr; man sagt da nicht: Kind, Kindeskind, also soll man mich fassen etc. Es kann nicht länger währen, denn bis daß die Verheißung kommt und erfüllt wird.

60. Also weist uns GOtt auf sich und auf sein Wort, und zeigt an, daß niemand GOtt treffe, denn allein durch sein Wort, und daß man das Wort nicht fassen könne, denn durch den Glauben. Darum so soll man kein Wort, Weise, noch Werk Gott vorziehen, denn, wie er mir [be]stimmt, wie ich glauben soll; damit sind alle Werke Mosis und unserer guten Werke aufgehoben. Und dann wissen wir, wer GOtt sei, und können ihn deuten. Denn allhier sagt er: Ich heiße, der es ist, der ich bin, und der ich sein werde. Wodurch? Du hast mein Wort, und hast das Zeichen Abrahä; an dem hange. Wenn nun die Verheißung erfüllt ist, so wirst du ein ander Wort und ein ander Zeichen haben.

61. Also hält dieser Text viel schöner, herrlicher Lehren in sich, als, Glauben an GOtt, item, von dem Reiche Christi, und von der Auferstehung der Todten. Hernacher werden nun etliche Wunderzeichen kommen im vierten Capitel, von dem Stabe, so zur Schlange wird, und von den Händen, die er in den Busen steckt und aussätzig werden etc. Und wenn wir nach der Historie diese Wunderzeichen ansehen, so will GOtt durch dieselbigen Mosen gewöhnen, fest und kühn zu werden, die Sache der Ausführung der Kinder Israel aus Egypten getrost anzufahren. Als sollte GOtt sagen: Mose, du hast ein schrecklich Werk vor dir, aber greif drein. Es scheint jetzt unmöglich, daß du den mächtigen König Pharao angreifen und stürzen solltest, aber wage es kurz und getrost. Die Schlange hat ein schrecklich Angesicht und ist giftig; aber ergreif sie bei dem Schwanze. Die Hand ist aussätzig; aber wage es nur auf mein Wort, so wird es ausgerichtet sein.

\quellseite{377}
\subsection{Allegoria oder geistliche Deutung des dritten Capitels}

62. Im dritten Capitel haben wir gehört, wie Moses die Schafe hütet, und er einen feurigen Busch sieht, daraus GOtt ihm erscheint, und befiehlt, daß er das Volk Israel aus Egypten führen soll; auch wie sich Moses wehrt und vorgibt, er habe eine schwere Sprache und Zunge, er könne nicht wohl reden etc. Das ist das dritte Capitel.

63. Daß aber der Busch brennt und doch nicht vom Feuer verzehrt wird, solches hat man in der ganzen Welt auf die Jungfrau Maria gezogen, daß ihre Jungfrauschaft verwahrt und behalten worden sei. Also geht es, wenn man in die heilige Schrift plumpt, daß man durch menschliche Vernunft die heilige Schrift anderswohin zieht. Wenn man heimliche Deutung auf Werke oder Personen stellt, so ist es falsch und unrecht. Sonst kann man Mariam wohl loben, aber es ist nicht vonnöthen, daß man Christo die Federn ausziehe, und stecke sie Marien auf, oder daß man Christo das Kleid ausziehe, und lege es der Jungfrau Marien an.

64. Es ist aber der HErr Christus derselbige Busch. Denn er ist der rechte Engel, von GOtt aus dem Himmel gesandt, wie zum Römern am fünfzehnten Capitel, V. 8., geschrieben steht: „Christus ist der Juden Apostel, auf daß er erfüllete die Verheißungen, den Patriarchen gegeben.“ Er ist ein Apostel, das ist, der rechte Erz- und höchste Bote, der in die Welt gesandt ist. Dieser HErr Christus wird dem Mosi gezeigt, daß ist, allen denen, die unter dem Gesetze waren. Denn Moses trägt das Bildniß und Figur des Gesetzes, und bedeutet des Gesetzes Amt; und das Volk, so unter dem Gesetze liegt, ist Moses; denn er sollte das Gesetz geben, und das Volk ausführen. Daß aber Christus brennt, und nicht verbrennt wird, das wird auf die Gottheit gedeutet; und nicht unrecht.

65. Aber ich will nicht zu hoch fahren, sondern diesmal vornehmen, das noch näher ist, nämlich, daß dies Christum, den wahrhaftigen Menschen, bedeute, der da erscheint in gleicher Gestalt als andere Menschen; gleichwie die eherne Schlange [4 Mos. 21, 9.] in der Wüste von Mose aufgerichtet, anderer Schlangen Gestalt und Bildniß hatte, aber da war kein Gift. Wie auch zum Römern am 8. Capitel, V. 3., gesagt wird: „Christus hat die Gestalt sündliches Fleisches, doch ohne Sünde“; und zum Philippern am 2. Capitel, V. 6–8. Christus hat auf Erden im Fleische gepredigt, wie ich und du, daß sein und mein Fleisch man nicht kann von einander scheiden. Jedoch, dieweil alle unser Fleisch vergiftet, und er hat auch unser Fleisch, darum so ist es eine solche Gleichniß, und hat das Fleisch, Sünde und Gift, aber doch ist nicht an seiner Person Sünde und Gift da, ist also in einem Fleische geboren, das sündig ist, und ist doch an ihm ohne Sünde; das ist, Christus ist wahrhaftig vom Fleische geboren. Also ist der Busch allhier voller Feuer, daß er brennt, und wird doch nicht verzehrt. Es bewegt auch Mosen nicht; als sollte er sagen: Ist Feuer da, so muß es verzehren, denn das ist seine Natur. Auf das Brennen und Verzehren sieht Moses; da ist das Fleisch, das brennt, und scheint wie ein ander sündlich Fleisch, und ist doch da keine Sünde. Christus ist der grüne Busch ohne alle Sünde, und dennoch ist er dem Fleische gleich in unsere Gestalt gekommen und fühlt eben das Fleisch an seinem Leibe, das wir haben, jedoch ohne Sünde.

66. Das ist der erste Blick\footnote{Blick = Anblick, Erscheinung. — Diese Bedeutung fehlt bei Diez.} und der erste Trost, daß wir einen solchen Engel oder Boten haben, welcher ohne alle Sünde ist, und doch gleich unserm Fleisch und Blute. Wir sind der feurige Busch, so verbrennen sollte, sind eitel Feuer und Gift; er aber bleibt unverbrannt. Darum hat er unser Feuer an sich genommen, aber er ist darob nicht verzehrt oder verbrannt, sondern wir sind darob verzehrt. Also hat er Christum gekrieget im Blick, der da gleich ist andern Menschen, und die Sünde schadet ihm nicht, sondern er hilft uns.

67. Also soll man die Allegoriam auf den wahrhaftigen Menschen Christum ziehen, und nicht auf die Gottheit, es sei denn, daß es nöthlich die heilige Schrift zwinge. Denn die heilige Schrift treibt am allermeisten und vornehmlich, daß Christus allein erkannt werde, daß er GOtt und Mensch sei. Denn, daß er GOtt ist, das hat er für sich, und hilft mich nicht, und deshalb ist er nicht mein Christus oder Heiland; sondern, daß er auch wahrhaftiger Mensch ist, das gilt mir.

\quellseite{378}
68. So viel hat er daran gewendet, daß ich es sauge und ertriege. Er gibt es dahin, und läßt es einen\footnote{Eislebensche: eines.} Schatz und Kost werden, was an uns gewandt wird; da stellt er sich als ein Heiland, Mittler und Helfer, gibt uns Hülfe und Arznei. Und, kürzlich zu reden, so will er, daß wir sollen seiner brauchen zu unserm Nutzen.

69. Also soll er Mosi erscheinen als ein Heiland, auf daß unsere Sünde auf ihn gelegt wird, und er am Kreuze gebraten würde, als wäre er voller Sünde, und ist doch ohne alle Sünde. Mosis und Christus müssen einander begegnen und erscheinen, aber Mosis verwundert sich darüber. Denn das ist die fröhliche Botschaft und das tröstliche Erkenntniß, daß Christus unsere Sünde auf sich nimmt, und die Sünden schaden ihm doch nicht.

70. Diese geistliche Deutung stärkt unsern Glauben, daß wir Christum recht und wohl erkennen. Denn, wenn ich es gleich lange auf Mariam und ihre Jungfrauschaft deute, so bin ich ichs gar nichts gebessert; denn ihre Jungfrauschaft ist eine Gabe, ihr verliehen und ihr Eigenthum. Aber diese Allegorie soll im gemeinen Nutzen bleiben, und jedermann soll davon Verstand und Besserung gewinnen, nämlich, daß Mosis Christum, den Heiland, zum ersten sieht, welcher die Sünde auf sich nimmt; und sie schaden ihm doch nicht. Damit angezeigt wird, daß man Christum nirgends abgemalt und beschrieben findet, wie er für uns hat leiden und sterben sollen, denn allein im Gesetze und im Mose. Wie denn der HErr Christus in dem Evangelisten Luca [Cap. 24, 44.] selbst bezeugt, daß Mosis, die Psalmen und Propheten von seinem Leiden, Sterben und Auferstehung geweissagt haben. Mosis hat Christum herrlich beschrieben, wie er Mensch werden und die Sünde wegnehmen, und durch seinen Tod und Auferstehung in sein ewiges Reich kommen sollte, und uns von Sünden, Tod, Teufel und Hölle helfen. Das ist Mosis Prophezei und Weissagung, wie es im ersten Buch Mosis und im Levitico gefunden wird.

71. Das erste Buch und Wort GOttes, darinnen wir Christum erkennen, ist Mosis Buch. Also erscheint Christus Mosi erst im feurigen Busche; als sollte er sagen: Mit diesem Bilde oder Figur zeige ich an, daß, wer Christum erkennen will, wer er sei, der findet es erstlich im Mose, daß er von Abrahams Samen kommen werde und die Welt erlösen, die Sünde auf sich nehmen und tödten, und dennoch dieser grüne Busch unverzehrt und unverletzt bleiben; die Sünden können ihn nicht unterdrücken und tilgen. Also ist in dieser Figur begriffen die Verfolgung, Kreuz, Leiden und Tod des HErrn Christi, auch sein ganzes Reich, und was wir Gutes und Wohlthaten von ihm haben, daß dieses alles im Gesetze Mosis steht.

72. Aber diese Predigt ist also fremd und seltsam, daß sie auch die Heiligen nicht bald glauben: Mosis muß sich selbst verwundern über den herrlichen, großen Gütern und Reichthum, so uns geschenkt worden sind durch Christum. Darum will er zum feurigen Busche treten, und sehen, wie es komme, daß er nicht verbrenne. Aber GOtt sagt ihm: „Beug deine Schuhe aus, denn die Stätte ist heilig.“ Aus dem Busche ist sie heilig gemacht. Man muß erstlich beichten und die Sünde bekennen, ablegen und sich rein machen, ehe man zu GOtt kommt [Jes. 1, 16. 17.]. Unser HErr GOtt hat nicht mit Sündern zu thun [Ps. 5, 5.], wir sollen rein und heilig zu ihm kommen, auf daß wir los werden. Unsere guten Werke, so wir an uns haben, das ist, die das Gesetz halten, und wollen damit kommen, wo Christus ist, die richten nichts aus. Denn GOttes Gesetz und Menschen Gesetz (wo sie also verstanden werden, daß man es mit guten Werken will ausrichten und hinan zu GOtt kommen) thut eines wie das andere, das ist, es wird nichts damit ausgerichtet, wenn man nicht sieht, daß das Gesetz geistlich sei. Denn es heißt allhier: Mose, tritt nicht herzu, du seiest denn heilig und rein; du Volk Israel kommst aus Vermessenheit, als seiest du heilig, und [ihr] wähnet, daß Christus euch komme aus eurem Verdienste und guten Werken.

73. Darum so heißen Schuhe allhier nicht Sünde oder gute Werke, sondern die Lehre der Werke, wie zum Ephesern am 6. Capitel, V. 15., auch gesagt wird, und im Hohenliede Salomonis Cap. 7, 1. steht, daß die Braut in blauen Schuhen gehe; das ist, das Evangelium oder GOttes Wort ist ein Schuh, darinnen die Gewissen gehen sollen, und dieses der Ursachen halben, daß man daraus einen Vorrath und Bereitung habe, daß man sicher gehe mit dem
\quellseite{379}
Gewissen. Aber diese Schuhe Mosis sind von Fellen und Häuten oder Leder gemacht, das ist, es sind Lehren, die aus GOttes Wort genommen und gezogen sind, und machen die Leute doch fleischlich. Also bin ich auch im Papstthum ein Schuhster worden, und habe mein Leben zum Paar Schuhe gemacht, und die Lehre angezogen, und gemeint, wenn ich dies oder jenes Werk thue, so sei ich selig.

74. Aber GOtt spricht: „Zeuch die Schuhe aus, denn diese Stätte ist heilig“, das ist, du wirst dadurch vielmehr einen Sünder machen, denn einen solchen Heiligen. St. Paulus und andere haben auch diesem Volke die Schuhe wollen ausziehen; aber welche bekehrt worden sind, denen sind die Schuhe ausgezogen. Durch das Gesetz werden wir nicht gerechtfertigt [Gal. 2, 16.]. Diese Schuhe müssen wir ausziehen und barfuß gehen, auf daß man wisse, deine guten Werke helfen dich nichts, wenn du zu diesem feurigen Busche gehen willst, sondern du mußt lernen, daß es mit dir alles aus und verloren sei. Wie denn allhier Moses es auch nicht verdient hat, daß ihm der Busch erscheint, sondern unversehens; und Moses hat darum auch nicht die Schuhe angezogen. Die Gnade und der HErr Christus kommen zu uns, ehe denn wir daran gedenken, oder darum bitten. Darum muß man alle Vermessenheit auf das Gesetz, Menschenlehre und gute Werke ausziehen, und frisch hinan auf Christum sich wagen, als auf unsern Heiland.

75. Das Feuer ist seine Marter und Leiden, da er am Kreuz gebraten, und doch nicht gar verzehrt wird. Das scheint, als sei er nicht mehr, und als sei es gar aus mit ihm und zu Pulver worden. Aber er steht auf von den Todten, und kommt als ein gewaltiger König. Wer nun hierzu kommen soll, der muß die Schuhe ausziehen. Denn die Apostel selbst verdienten es nicht, das Evangelium zu hören, das aus diesem Busche allhier gepredigt wird.

\section[Ein andere Predigt über 2 Mos. 3, 1–6]{Ein andere Predigt D. M. Luthers über den Text 2 Mos. 3, 1–6., da Moses den feurigen Busch sahe etc.}

Gethan zu Wittenberg in der Pfarrkirche, am Osterdienstage, Nachmittags, den 3. April 1526.\footnote{Diese Zeitbestimmung ergibt sich mit Sicherheit aus der Zusammengehörigkeit dieser Predigt mit den beiden Predigten über 1 Mos. 3 und 22, welche wir in diesem Bande Col. 650 und 662 mitgetheilt haben, denn der erste Paragraph dieser Predigt gibt eine kurze Recapitulation dessen, wovon Luther am Ostermontage geredet hatte. Außerdem ist in dem Index Voachs in der Zwickauer Rathsschulbibliothek das von uns gesetzte Datum für diese Predigt angegeben, deren Inhalt so bezeichnet ist: „figurae quibus resurrectio significata“. Endlich wird das Datum auch durch § 14 dieser Predigt bestätigt. Den Nachweis, wo diese Predigt zuerst gedruckt worden ist, findet man in diesem Bande Col. 650.}

1. Dieweil ihr, lieben Freunde, die beiden Tage aus dem Neuen und Alten Testamente genugsam gehört habt, wie unser HErr Jesus Christus erstanden ist, und was er mit seiner Auferstehung ausgerichtet, als nämlich, daß er der Schlange, dem Teufel, den Kopf zertreten, vom Fluche erlöst, und als ein rechter GOtt und Mensch den reichen Segen über alle Menschen gebracht, daß sie haben Gerechtigkeit, Vergebung der Sünde, Leben und Seligkeit, Ueberwindung des Todes, Hölle und Teufels, daß sie mögen sprechen mit gläubigen Herzen: Christi Sterben und Auferstehen ist meine Gerechtigkeit, dadurch ich gerecht bin gemacht, wie St. Paul spricht zum Römern am 4. Capitel, V. 25.: „Christus ist um unserer Sünden willen dahin gegeben, und um unserer Gerechtigkeit willen auferwecket“, dieweil wir aber heute noch das fröhliche Osterfest feiern, so wollen wir aus dem heiligen Mose, wie in seinem andern Buche am 3. Capitel geschrieben, diesen Text, der also lautet, vor uns nehmen:

\quellseite{380}
3. 1–6. Moses aber hütete der Schafe Jethro, seines Schwähers, des Priesters in Midian, und trieb die Schafe enhinter in die Wüste, und kam an den Berg GOttes, Horeb. Und der Engel des HErrn erschien ihm in einer feurigen Flamme aus dem Busche, und er sahe, daß der Busch mit Feuer brannte, und ward doch nicht verzehrt. Und sprach: Ich will dahin und sehen dies Gesicht, warum der Busch nicht verbrennt. Da aber der HErr sahe, daß er hinging zu sehen, rief ihm GOtt aus dem Busch und sprach: Mose, Mose! Er antwortete: Hier bin ich. Er sprach: Tritt nicht herzu, zeuch deine Schuhe aus von deinen Füßen, denn der Ort, da du auf stehest, ist ein heilig Land. Und sprach weiter: Ich bin der GOtt deines Vaters, der GOtt Abrahams, der GOtt Isaaks, und der GOtt Jakobs. Und Mose verhüllete sein Angesicht, denn er fürchtete sich, GOtt anzusehen.

2. Das ist der Text aus Mose. Nun, lieben Freunde, diese Worte des heiligen Mosis haben die Väter der Kirche auf die zarte, reine, keusche, unbefleckte Jungfrau Maria, die Mutter Christi, gezogen. Wie man auch in der Kirche gesungen hat: Rubum, quem viderat Moses, daß es soll bedeuten die Jungfrauschaft Mariä. Gleichwie der Busch brennt und nicht verzehrt wird, also hat die Mutter Christi auch geboren ihren lieben Sohn, ohne Versehrung ihrer Jungfrauschaft.

3. Wiewohl es nun eine gute Meinung oder Deutung ist (denn diesen Artikel müssen wir wahr sein und bleiben lassen, daß Maria ist eine reine Jungfrau gewesen vor der Geburt, in der Geburt, und nach der Geburt, und bleibt ewiglich eine reine Jungfrau), so ist doch diese Meinung oder Deutung hier zu schwach, und hält den Stich nicht, denn der Heilige Geist etwas Höheres will vorbilden. Denn da stehen die klaren und hellen Worte: Ich bin der GOtt deines Vaters, Abrahams, Isaaks und Jakobs, welche Worte auf Marien nicht können gezogen werden. Denn Maria ist nicht GOtt, darum kann sie nicht der Busch sein, sondern GOtt ist selber der Busch, der da brennt und nicht verzehrt wird, bleibt ganz und sein lustig grün mitten im Feuer.

4. Darum können wir es nicht besser treffen, denn wir bleiben bei den Worten Christi, daß Moses von ihm geschrieben habe; wie er denn sagt Joh. 5, 46.: „Wenn ihr Mosi glaubtet, so glaubtet\footnote{Erlanger: glaubet.} ihr auch mir, denn er hat von mir geschrieben.“ Denn Christus ist es allein, auf den Moses und alle Propheten weisen [Matth. 11, 13. Apost. 10, 43.]. Derwegen ist dieser feurige und brennende Busch eine Figur Christi, was er für eine Person sein wird, und wie er zwo Naturen in Einer Person haben soll, und wie er soll leiden, sterben und wieder auferstehen von dem Tode. Das ist es, das man hier suchen soll, das Moses lehrt, und das wir auch glauben und lernen sollen, wollen wir anders selig werden.

5. Nun, das erste Stück. Was soll denn Christus für eine Person sein, und was für zwo Naturen soll er haben in Einer Person? Da hören wir es im Texte: „Ich bin der GOtt Abrahams, Isaaks und Jakobs.“ Das ist Eine Natur, daß er ist ewiger, allmächtiger GOtt, gleich dem Vater und dem Heiligen Geiste, gleicher Natur, Substanz und Wesens, doch eine andere Person, denn der Vater und Heilige Geist. Wie denn der Prophet Micha (welcher 760 Jahre vor der Geburt Christi gelebt) auch so lehrt, da er Cap. 5, 1. spricht: „Und du Bethlehem Ephrata, die du klein bist gegen den Tausenden in Juda, aus dir soll der kommen, der in Israel HErr sei.“ Ja, was ist es denn für ein HErr? So sagt der Prophet: „Welches Ausgang von Anfang und von ewig her gewesen ist.“

6. Da beschreibt der Prophet klärlich, daß Christus noch einen andern Ausgang habe denn zu Bethlehem, als nämlich, aus dem väterlichen Herzen GOttes, ehe alle Creaturen gewesen und ihren Anfang gehabt. Denn ehe alle Creaturen gewesen und geschaffen worden, ist nichts gewesen denn GOtt. Dieser Sohn, der zu Bethlehem soll ausgehen, das ist, geboren werden, sagt Micha, ist von Ewigkeit gewesen, ehe alle Creaturen ihren Anfang gehabt; so muß der Sohn zu Bethlehem sein ewiger, allmächtiger GOtt. Das ist nun die erste, göttliche Natur, die Christus hat, daß er rechter, ewiger und allmächtiger GOtt sei, gleich dem Vater und Heiligen Geiste, der mit Mose aus dem brennenden Busche geredet hat; wie denn gestern\footnote{In der zweiten Predigt am Ostermontage, Col. 662 f.} genug gesagt worden.

7. Die andere Natur ist die menschliche Natur, daß er ist ein feiner, grüner, lieblicher Busch,
\quellseite{381}
wie ihn auch Jesaias [Cap. 11, 1.] nennt, eine feine, liebliche, grüne Ruthe oder frischer Zweig aus dem alten Stamme oder Wurzel Jesse. Wie sich denn der HErr Christus selber heißt, als man ihn zum Tode führte, da er [Luc. 23, 31.] zu den Weibern, die ihn beweinen, spricht: „Thut man das am grünen Holz, was will am dürren werden?“ Denn Christus ist allein das grüne Holz oder der grüne Busch, das ist, ein rechter, reiner, heiliger Mensch. Aber alle Menschen, niemand ausgenommen, die werden in Sünden empfangen, geboren, und sind des Todes und verdammt ewiglich; darum sind sie dürre Hölzer oder Klötzer, die nur zum Feuer dienen, sonst nirgend zu.

8. Aber Christi Empfängniß ist rein und heilig. Denn da kommt der Heilige Geist, und wirkt aus den reinen Blutstropfen der reinen Jungfrau Marien, daß aus GOtt oder GOttes Sohne ein wahrhaftiger, rechter Mensch wird; darum ist auch seine Geburt rein und heilig. Denn wo die Empfängniß rein und heilig ist, da ist auch die Geburt rein und heilig. Derwegen ist er allein der grüne Busch, denn in ihm ist keine Sünde befunden worden, wie denn sonst alle Menschen haben. Das ist nun die andere Natur, die der Mann hat, daß er wahrhaftiger Mensch ist, wie ich bin, aber ohne alle Sünde.

9. Das sind die zwo Naturen, die in dieser einigen Person Christi sind, daß er zugleich wahrer GOtt und Mensch ist, und sind die zwo Naturen mit einander so vereinigt, daß man keine von der andern scheiden kann, sondern [sie] müssen zusammen bleiben. Wie nun Leib und Seele in einem Menschen vereinigt sind, und ist doch nur Eine Person und Ein Mensch; also auch hier, in der Person Christi, muß die göttliche und menschliche Natur, das ist, GOtt und Mensch in Einer Person zusammen bleiben unzertrennt. Das ist eine hohe Lehre, von der weiß kein Jude, Türke oder der Papst nichts, allein Moses und die Propheten, und die Christen. Und das ist nun von der Person und beiden Naturen gesagt, daß Christus rechter, wahrhaftiger, ewiger, allmächtiger GOtt und Mensch ist, als das liebe grüne Büschlein oder Rüthlein, rein, lieblich, ohne alle Makel und Sünde.

10. [Wir] wollen jetzt aber hören, wie es dem HErrn Christo gehen soll. Davon sagt Moses, wie der Busch brennt und wird nicht verzehrt, wie sehr und gewaltig auch die Flammen umher brennen. Da ist nun bedeutet Christi Leiden und Auferstehen; denn „brennen“ in der Schrift heißt leiden. Also soll nun dieser GOtt, der ein grüner Strauch ist, das ist, ein rechter, reiner, wahrhaftiger Mensch, von der Jungfrau geboren, leiden und sterben.

11. Dieses ist auch die Ursache seiner Menschwerdung. Denn wenn GOtt nicht wäre Mensch geworden, so könnte er nicht leiden oder sterben. Denn „GOtt ist ein Geist“ [Joh. 4, 24.], der nicht kann leiden oder sterben. Soll aber GOtt leiden und sterben, so muß er Mensch werden. Dieweil aber GOtt Mensch wird, so leidet er und stirbt in dieser Person Christi, die zugleich GOtt und Mensch ist, und GOtt bleibt nicht außer dieser Person Christi, wie etliche Ketzer und Schwärmer trennen und scheiden die Gottheit von der Menschheit, und sagen, die Menschheit Christi habe allein gelitten und uns erlöst; das ist falsch. Denn die Menschheit hat wohl allein gelitten, dieweil die göttliche Natur nicht kann leiden oder sterben, hat uns aber nicht allein erlöst, sondern auch die Gottheit, das ist, GOttes Sohn, wie St. Paulus sagt zum Römern am achten Capitel, V. 32., „daß GOtt auch seines einigen Sohnes nicht hat verschonet, sondern hat ihn für uns alle dahin gegeben“.

12. Darum hat uns nicht allein Marien Sohn, nach der Menschheit, wie die Papisten und etliche Schwärmer narren, erlöst, ist für uns gestorben und gekreuzigt, sondern der ganze Christus, der da ist GOtt und Mensch, hat uns erlöst, und der ist gekreuzigt, gestorben und auferstanden nach dem Fleisch; in der ersten Epistel Petri am dritten Capitel, V. 18. Denn die Wunden, die sie Christo in seine Hände und Füße geschlagen, haben sie nicht allein Marien Sohne in seine Hände und Füße geschlagen, sondern GOttes Sohn, oder GOtt selber; haben auch nicht allein des Menschen-Sohns Fleisch gekreuzigt und gemartert, sondern [es] ist also GOttes Sohn und GOtt selbst ermordet und erwürgt, denn GOtt und Mensch ist Eine Person. Darum ist der GOtt gekreuzigt und gestorben, der Mensch worden; nicht der abgesonderte GOtt, sondern der vereinigte GOtt mit der Menschheit; nicht nach der Gottheit, sondern nach der menschlichen Natur, die er angenommen hat. Wie auch St. Paulus [1 Cor. 2, 8.] sagt: Dominum gloriae crucifixerunt.

\quellseite{382}
13. Denn von dem abgesonderten GOtt ist es beides falsch, als, daß Christus GOtt ist, und GOtt gestorben ist. Beides, sage ich, ist falsch. Denn da ist GOtt nicht Mensch; darum kann er auch nicht sterben. Derhalben sagen und reden wir es von dem GOtt, der zugleich GOtt und Mensch ist in Einer Person, als von Jesu Christo; wie denn auch die lieben Väter, als Marentius und Fulgentius, sagen: Nos dicimus Deum natum, passum, crucifixum, mortuum, propter unitatem et conjunctionem duarum naturarum in Christo. Das können aber unsere Papisten und Schwärmer nicht verstehen oder vernehmen.

14. Darum theilen und trennen sie die Person Christi. Das muß aber nicht sein. Denn wir ehren und beten nicht allein die schlechte, bloße Menschheit in Christo an, wie uns die Schwärmer Schuld geben, sondern die Gottheit und Menschheit, das ist, GOtt und Mensch zugleich, als den rechten Schöpfer Himmels und der Erde, in Einer Person zusammen vereinigt. Wie denn das Concilium zu Epheso, vor 1095\footnote{„1095“ von uns gesetzt statt „1115“ in den Ausgaben. Das Concil zu Ephesus fand im Jahre 431 statt. Da nun diese Predigt am Osterdienstage 1526 gehalten ward, muß „1095“ gelesen werden. Der Abschreiber hat aus MXCV die Zahl MCXV gemacht. (Genau dasselbe Versehen (110 [CX] statt 90 [XC]) findet sich in der alten Ausgabe Walchs zweimal auf der letzten Seite des 9. Bandes, wo wir dieselbe Veränderung vorgenommen hatten.} Jahren wider den Ketzer Nestorium, auch schließt und setzt: Propter assumptem veneror assumptum, et propter invisibilem adoro visibilem. Das ist recht gelehrt und gesagt.

15. So spricht auch der alte Lehrer der Kirche, Johannes Damascenus, lib. 4. cap. 3 von Christo: Caro secundum sui naturam non est adorabilis; adoratur autem in incarnato Dei verbo, non propter se ipsam, sed propter conjunctum ipsi secundum hypostasin Deum verbum. Et non dicimus, quod carnem adoramus nudam, sed carnem Dei, id est, incarnatum Deum. Das ist wahrlich ein schöner, gewaltiger Spruch, den wir zu dieser jährlichen Zeit wohl behalten und merken sollten. Denn die beiden Naturen in Christo wollen unzertrennt und ungeschieden sein. Aber darein können sich unsere neuen Sacramentierer nicht richten, wir Christen aber müssen es lehren und lernen. Denn das wäre gar eine schlechte, ja, gar keine Erlösung gewesen, die uns wider Sünde, Tod, Teufel und Hölle nicht viel geschützt hätte, wenn nur allein der Mensch Christus, und nicht zugleich GOtt, oder GOttes Sohn, in dieser Person vereinigt, gekreuzigt und gestorben wäre. Davon wir ein ander Mal, mit GOttes Hülfe, mehr sagen wollen, und ist diese Tage genug getrieben.

16. Was ist aber das, daß der Busch nicht verzehrt oder verbrannt wird, sondern mitten im Feuer grünt und ganz bleibt? Das ist es: Ob Christus gleich wird sterben und leiden, soll er doch nicht im Tode bleiben, sondern wieder erstehen, grünen und lebendig werden. Warum denn das? Darum, denn er ist nicht ein schlechter Mensch, sondern wahrhaftiger GOtt. Wenn er ein schlechter Mensch wäre gewesen, wie wir, so wäre er dem Tode und Teufel viel zu schwach gewesen; dieweil er aber GOtt ist, und GOtt das Leben ist [Joh. 1. 4.], so kann das Leben\footnote{So von uns herübergenommen aus der alten Ausgabe Walchs. Erlanger: und das Leben kann.} nach keiner ewigen göttlichen Natur nicht sterben. Denn ob es gleich nach dem Fleische stirbt, kann es im Tode nicht bleiben. Denn Leben muß Leben bleiben; wie auch die Kirche gesungen: Vita in ligno moritur. Derhalben so muß dieser GOtt, der Mensch ist, vom Tode wieder erstehen, grünen und lebendig werden, damit er bringe Vergebung der Sünde, Gerechtigkeit, Leben und Seligkeit allen, die an ihn glauben. Das heißt denn: Infernus ex morsu despoliatur, die Hölle hat die Zähne an diesem Christo, der das Leben ist, stumpf gebissen, und hat ihre Kraft darüber verloren.

17. Denn der Christus grünet wieder aufs schönste mitten im Leiden, und geht ein zu seiner Herrlichkeit; wie er selber spricht [Luc. 24, 26.], das ist, er erstehet vom Tode, wird ein HErr des Himmels, der Erde und der Hölle, daß alle Creaturen im Himmel, in der Hölle und auf Erden müssen unter seinen Füßen liegen, und ihm unterthan sein, und für einen HErrn erkennen [Phil. 2, 9. 10.].

18. Das ist nun die rechte Figur oder Bedeutung dieses brennenden Busches, den Moses sieht, daß der GOtt, der Abrahams, Isaaks und Jakobs GOtt ist, soll Mensch werden, leiden, sterben und auferstehen vom Tode, und allen Gläubigen bringen Gerechtigkeit, Vergebung der Sünden, Leben und Seligkeit, wie die Christenheit gesungen: „Also heilig ist der Tag, daß den kein Mann mit Lob erfüllen mag, alleine
\quellseite{383}
der heilige GOttes-Sohn, der die Hölle zerbrach, und den leidigen Teufel darinne band, damit erlöset der HErr die Christenheit, das war GOtt selber, Kyrieleis.“ Da steht die ganze Frucht der Auferstehung Christi, als, daß er die Hölle zerbrochen, den Teufel gebunden, und seine Macht genommen, und seine liebe Christenheit erlöset. Ja, wer ist es aber gewesen? GOtt, GOtt selber, oder GOttes Sohn. Das ist je köstlich, lieblich, tröstlich und wohl gesungen, wenn wir es nur auch betrachteten.

19. Aber wie geht es nun Mosi? Wie er nun hin will und ziehen, wie dieser Busch brennt, und nicht verzehrt wird, da sagt GOtt zu Mose: „Tritt nicht herzu, zeuch vor deine Schuhe aus, denn der Ort, da du auf stehest, ist ein heilig Land.“ Was will da werden? Liegt denn GOtt am Lande so groß, und an den Schuhen? Ist denn nicht ein Ort so wohl heilig als der andere? Was thun denn die Schuhe dazu? Ja, der schlechte Ort Landes, oder die Schuhe thun es nicht; sondern das ist es, das GOtt will haben, daß, wo die Lehre geht und gepredigt wird, daß GOtt Mensch worden, gestorben und wieder erstanden ist, da ist das heilige Land, und die christliche Kirche, und das rechte auserwählte Volk, da GOtt ist und wohnet, wie Christus Joh. 14, 23. spricht. Da muß Moses die Augen, die Adam im Paradiese zu weit aufgethan, wieder zuthun, den alten Adam und die alten Schuhe ablegen und ausziehen, die Augen der Vernunft zuthun, und die Augen des Herzens aufthun.

20. Diese Lehre von Jesu Christo muß allein durch das Wort im Glauben gefaßt werden, denn sie ist dem alten Adam und der menschlichen Vernunft zu hoch, und GOtt hat sie allein in das Wort gelegt, da kann man sie auch fassen, und sonst nicht. Darum die Juden, desgleichen Türken und der Papist, dieweil sie das Wort von Christo fahren lassen, und gaffen dieweil anderswohin, und speculiren und reden von Christo nach der menschlichen Vernunft, kommen sie immer weiter von Christo, und werden zu Narren darüber, lästern, spotten den lieben HErrn Christum. Das macht es, wenn man das Wort fahren läßt von diesem Manne.

21. Derhalben müssen wir allein auf das Wort sehen und sagen: Da steht mein Artikel des Glaubens einfältig: Ich glaube an Jesum Christum, GOttes Sohn, unsern HErrn, der geboren, gelitten, gestorben, begraben und wieder auferstanden ist. Das habe ich nicht erdacht, sondern GOtt hat mir es durch den Heiligen Geist schreiben und predigen lassen. Darauf will ich bleiben, dieweil ich höre, wie GOtt, oder GOttes Sohn, ist Mensch worden, gestorben und auferstanden. Kann ich es nicht verstehen oder begreifen, so will ich es schlecht glauben.

22. Wie es aber zugegangen, daß GOtt stirbt, und ist doch nicht GOtt der Vater, oder GOtt der heilige Geist, sondern GOtt der Sohn, die andere Person in der Gottheit, der ist gestorben und auferstanden, und sind doch nicht drei Götter, sondern Ein allmächtiger GOtt, aber drei unterschiedliche Personen in dem ewigen, allmächtigen göttlichen Wesen, das wollen wir sparen bis in jenes Leben, da werden wir es in dem Spiegel der heiligen Dreifaltigkeit oder Gottheit von Angesicht zu Angesicht sehen. Hier aber wollen wir schlecht, recht und einfältig dem Wort glauben und trauen; weiter können wir es nicht bringen.

23. So hat Moses seine Schuhe, den alten Adam oder Menschen, müssen ausziehen und ablegen, und seine Augen oder Angesicht verhüllen, und allein durch das Wort, mit festem Glauben und Vertrauen, diesen Christum, sein Leiden, Sterben und Auferstehen, mit den Augen des Herzens ansehen und fassen müssen. Thun wir das auch, wie Moses gethan hat, so werden wir auch zu GOtt treten. Denn ohne diesen GOtt, der gestorben ist und auferstanden, fürchte und scheue, oder unterstehe sich kein Mensch, wie heilig, fromm, oder wie viel guter Werke er voll sei, zu GOtt zu treten, oder zu ihm zu kommen.

24. Denn GOtt der Vater kann keinen leiden, der zu ihm gehen oder treten will, er bringe denn diesen seinen geliebten Sohn Christum mit sich. Wie auch der Sohn selber spricht Johannis am vierzehnten Capitel, V. 6.: „Ich bin der Weg, die Wahrheit, und das Leben, niemand kommt zum Vater, denn durch Mich, Mich, Mich.“ Hier werden alle Heiligen und alle ihre Verdienste ganz und gar ausgeschlossen, daß nichts vor dem Vater gelten soll, es heiße denn Christus.

25. Derhalben sind verdammt Juden, Türken und der Papist, die diesen Sohn GOttes mit seinem Leiden, Sterben und Auferstehen verachten, und einen andern Weg zu GOtt zu kommen vornehmen, als durch die todten Heiligen, verdammten Ablaß, Seelmessen, laufige Mönchskappen und Platten, und andere erdichtete Men-.
\quellseite{384}
schenwerke, die GOtt nicht geboten oder geheißen. Denn durch diesen Sohn und sein Wort soll man ihn allein suchen und antreffen. Selig aber sind die alle ewiglich, die in dem geliebten Sohne GOttes, der uns durch sein Sterben und Auferstehen gerecht gemacht (wie Moses und die Propheten gethan), mit festem Glauben an sein Wort, im Herzen beständig verharren und bleiben.

26. Joseph in Egypten, wie im ersten Buch Mosis am ein und vierzigsten Capitel steht, ist auch gewesen eine Figur Christi. Denn Joseph wird von seinen eigenen Brüdern verkauft, darnach durch die Hure, des Potiphar des Hofmeisters Weib, als er ihr nicht folgen will, so belogen, daß er in das Gefängniß geworfen wird. Darnach, als Joseph lange im Gefängniß gelegen, kommt GOtt, und bringt ihn zu solchen Ehren, daß er muß Herr und Fürst werden über ganz Egyptenland.

27. Also ist es auch gegangen unserm HErrn Christo, der wird auch von seinen Brüdern, das ist, von Abrahams Fleische und Blute, den Juden, verkauft. Darnach, da er mit der jüdischen Synagoge nicht huren will, und ihre gottlose, teuflische Lehre zu Jerusalem helfen bestätigen und loben, da werden sie zornig, fangen und tödten ihn, und meinen, es sei nun gar mit ihm aus. Aber ehe drei Tage um sind, bringt ihn GOtt zu solchen Ehren, daß er durch den Tod zu seiner Herrlichkeit eingeht, ersteht vom Tode, wird ein HErr über Himmel, Erde, Hölle und alle Creaturen; das alles muß ihm unter seinen Füßen liegen, und ihn für einen HErrn erkennen und halten.

28. Damit wollen wir nun das fröhliche, liebliche Osterfest beschlossen haben. Bittet aber von Herzen, daß GOtt, der Vater unsers lieben HErrn JEsu Christi, sammt dem Heiligen Geiste, uns alle in solchem Glauben und Erkenntniß wolle erhalten. Thut er es nicht, so ist es mit uns verloren. Wenn wir gleich täglich schreien mit dem Munde: Ich glaube, glaube, so ist es doch nur ein Gewäsche. Denn St. Paulus spricht in der andern Epistel zum Thessalonichern am dritten Capitel, V. 2.: „Der Glaube ist nicht jedermanns Ding.“ Derhalben lasset uns schreien mit dem Vater des besessenen Männchen [Marc. 9, 24.]: „O lieber HErr, ich glaube, hilf aber meinem Unglauben“; und mit den lieben Aposteln aus Grunde unsers Herzens sagen [Luc. 17, 5.]: „O HErr, stärke uns den Glauben.“ Thun wir das, so will uns GOtt den stärken, mehren und geben; sonst nicht.

29. Darnach müssen und sollen wir, als die lieben Kinder GOttes, den zehn Geboten GOttes gehorsamlich nachleben, damit die Leute unsere guten Werke sehen und unsern Vater im Himmel preisen; wie Christus spricht, Matthäi am 5. Capitel, V. 16. Amen.

\section{Das vierte Capitel}

\begin{schriftzitat}
V. 1–9. Mose antwortete und sprach: Siehe, GOtt, sie werden mir nicht glauben, noch meine Stimme hören, sondern werden sagen: Der HErr ist dir nicht erschienen. Der HErr sprach zu ihm: Was ist es, das du in deiner Hand hast? Er sprach: Ein Stab. Er sprach: Wirf ihn von dir auf die Erde. Und er warf ihn von sich, da ward er zur Schlange; und Mose floh vor ihr. Aber der HErr sprach zu ihm: Strecke deine Hand aus, und erhasche sie bei dem Schwanz. Da streckte er seine Hand aus, und hielt sie, und sie ward zum Stabe in seiner Hand. Darum werden sie glauben, daß dir erschienen sei der HErr, der GOtt ihrer Väter, der GOtt Abrahams, der GOtt Isaaks, der GOtt Jakobs. Und der HErr sprach weiter zu ihm: Stecke deine Hand in deinen Busen; und er steckte sie in seinen Busen, und zog sie heraus; siehe, da war sie aussätzig wie Schnee. Und er sprach: Thue sie wieder in den Busen; und er that sie wieder in den Busen, und zog sie heraus; siehe, da ward sie wieder wie ander Fleisch. Wenn sie dir nun nicht werden glauben, noch hören, was du sagst, bei einem Zeichen, so werden sie doch glauben, wenn du sagst vom andern Zeichen. Wenn sie aber diesen zweien Zeichen nicht glauben werden, noch deine Stimme
\end{schriftzitat}

\quellseite{385}
hören, so nimm des Wassers aus dem Strom, und geuß es auf das trockne Land, so wird dasselbe Wasser, das du aus dem Strome genommen hast, Blut werden auf dem trocknen Lande.

1. Wir haben bisher gehört in den dreien Capiteln des andern Buchs Mosis, wie GOtt Mosen berufen hat zu einem großen, herrlichen Werk, in dem, daß er die Kinder von Israel hat aus Egypten führen sollen, und den gewaltigen, mächtigen König in Egypten zu Boden stürzen. Auch haben wir gehört von GOttes Verheißung und Zusagung, daß er wolle bei Mose sein, ihm helfen, daß er dies große Werk hinaus führen könne. Aber Moses hat sich gesperrt und geweigert, solch Amt und Werk anzunehmen und auszurichten.

2. Nun fährt GOtt in diesem vierten Capitel zu, und gibt Mosi dreierlei Zeichen oder Mirakel, daran er und auch die Kinder Israel seine göttliche Gegenwärtigkeit merken mögen. Erstlich, ein Stab oder Stange, die er aus seiner Hand wirft auf die Erde, welcher zur Schlange wird. Zum andern, daß er seine Hand in den Busen steckt, und sie aussätzig wird, wie Schnee; da er sie aber wieder in den Busen that, und zog sie wieder heraus, da ward sie wieder wie ander Fleisch. Zum dritten: „Wenn sie diesen beiden Zeichen nicht glauben werden, noch deine Stimme hören“, spricht GOtt, „so nimm das Wasser aus dem Strom, und geuß es auf das trockne Land, so wird dasjenige Wasser, das du aus dem Strom genommen hast, Blut werden auf dem trocknen Lande.“

3. Diese Zeichen that Moses zur Bestätigung seines Berufs, auf daß er in seinem Amte und Berufe getrost fortführe, und nicht daran zweifelte, GOtt wäre bei ihm, der mit seinem Munde lehrete und redete, was Moses thun und reden sollte, item, die Kinder Israel durch diese Mirakel auch vergewissert und versichert würden, daß Moses sollte der Mann sein, von GOtt dazu sonderlich bestellt, daß er sollte sie aus Egypten erretten. Sonst hatte Moses für sich den feurigen Busch und GOttes Zusagung, damit er sich in seinem Berufe getröstet. Aber diese drei Wunderzeichen sollen um der Kinder Israel willen geschehen, daß sie glauben, die Zeit ihrer Erlösung wäre auch vorhanden. Was aber der Stab oder Knüttel, die aussätzige Hand, und das Wasser sei, das wollen wir in der Allegorie hören.

\begin{schriftzitat}
V. 10–12. Mose aber sprach zu dem HErrn: Ach mein HErr! ich bin je und je nicht wohlberedet gewesen, sint der Zeit du mit deinem Knechte geredet hast; denn ich habe eine schwere Sprache, und eine schwere Zunge. Der HErr sprach zu ihm: Wer hat dem Menschen den Mund geschaffen? oder wer hat den Stummen, oder Tauben, oder Sehenden, oder Blinden gemacht? Habe ich es nicht gethan, der HErr? So gehe nun hin. Ich will mit deinem Munde sein, und dich lehren, was du sagen sollst.
\end{schriftzitat}

4. Jetzt sind wir in der Historie und wollen daraus sehen, wie viel Mühe es kostet, ehe denn ein Mensch GOtt vertraue, und sein Wort annehme, sonderlich, wenn es etwas Hohes und Großes ist, als hier, in diesem Capitel, da sich GOtt also sehr bemüht und erniedriget, und so freundlich die Leute lockt und treibt, bis sie glauben. Moses bleibt gleichwohl bei seiner letzten Entschuldigung, dieweil er nicht anders kann, und hat eine andere Ausflucht oder Loch gesucht, da er hinaus will, auf daß er GOtt nicht Gehorsam leiste. Denn er spricht: Ob ich gleich alles thue, so kann ich doch nicht reden. Er bekennt frei, daß er kein Schwätzer sei. Die Juden legen es aus, es sei das Lispeln gewesen; ich aber glaube, Moses habe eine schwere Sprache und Zunge gehabt, daß er gestammelt und gestottert habe, und ein Wort bis in zwei- oder dreimal geredet, und doch nicht recht herausgebracht habe. Diesen Gebrechen hat er gehabt, und daß er nicht gelispelt.

5. Die Juden geben es also vor, daß er feurige Kohlen in das Maul gelegt, damit er die Zunge verbrannt, und daher habe er gelispelt. Aber der Text spricht, er habe gestottert. Es sind jüdische Märlein. Es ist der letzte Auszug oder Ausflucht, dadurch Moses dies Amt will ablegen. Aber GOtt antwortet ihm recht darauf, ist noch nicht zornig, und spricht: „Wer hat dem Menschen den Mund geschaffen? oder wer hat den Stummen oder Tauben, Sehenden oder Blinden gemacht? Habe ich es nicht gethan?“ Wenn GOtt so klug wäre, als wir sind, so hätte er alle Dinge besser angefangen, denn sonst geschehen.

6. Denn allhier nimmt er zu diesem schweren, hohen Werke einen, der nicht wohl reden kann, wie es denn Moses selbst bekennt; noch spricht GOtt zu ihm: Gehe hin und richte es wohl aus.

\quellseite{386}
Welches eben also lautet, als wenn ich zum Blinden sagte, daß er wohl sehen, und zu einem Lahmen, daß er wohl laufen, und zu einem Stummen, daß er wohl reden sollte. Konnte GOtt nicht einen andern finden, den er zu diesem Werke gebrauchte? Aber es ist darum geschrieben, daß wir lernen sollen, wie GOtt gesinnt sei. Was da gilt vor der Welt, das achtet er nicht; er verwirft und verstößt, was andere zu sich reißen; was andere lieben und aufheben, das wirft er weg; und was der Teufel nicht mag, das nimmt er an. Er gibt dem Mose eine rechte Antwort und spricht: Du bist klug und ein feiner Gesell; es sind Sticherlinge,\footnote{Sticherlinge ist wohl so viel als spitzige Reden.} als sollte er sagen: Meinst du, daß ich nicht wisse, daß du stammelst, und als hörte ich es nicht? Also dünkt es uns. Denn wir immerdar GOtt meistern in seinen Werken, gleich als wenn wir zum allerersten die Fehler, Mängel und Gebrechen sähen, GOtt aber sie nicht sähe. Was liegt daran, will GOtt sagen, ob du taub, blind oder stumm seiest? Wie denn, wenn ich dich's heiße, und gebiete dir etwas, kann ich dich nicht alsdann sehend, hörend und redend machen? Wer ist, der mit dir redet? Es ist nicht Jung Schuster, sondern der, welcher den Blinden die Augen und Gesicht, den Tauben die Ohren und Gehör gibt, und die Unberedten beredt machen kann; wiederum, die großen Schwätzer zu Stummen macht. Und du wolltest mir Ziel und Maß setzen, der du nicht reden kannst; da ich dich doch darum erwählt habe, daß du nicht reden kannst? Wenn du wohl reden könntest, so solltest du dich deß noch wohl überheben. Auf daß sie nun sehen, daß ich der Mann sei, der solches thue, und nicht du, so gebrauche ich dich Stammelnden zu diesem Werke. Denn wenn einer also geschickt wäre, als Gabriel und alle Engel, und ihn GOtt nicht berufe, so würde er doch nichts ausrichten. GOtt macht Beredte und Stumme. Wiederum, ist einer unberedt, und GOtt beruft ihn, so führt er es hinaus, er sei, wie er wolle, auf daß die Welt sehe, wir sind es nicht, die es treiben, sondern GOtt thut es.

7. Derhalben so hat er allhier seine höchste Weisheit an Tag geben wollen, daß dieser Stotterer, Moses, so ein trefflicher Mann wird, der alle Könige überredet. Es möchte einer sagen: Ei, greif es flüglich an, lieber Herr Gott. Ja, antwortet er, ich bin ja der Meister; habe ich es nicht gethan? Ist nicht Moses gefangen mit seiner Zunge, daß er stottert und stammelt? Aber was er reden soll, das hat er von mir. Wenn ich nicht will, so wird er gar stumm; wenn ich will, so muß er wieder reden. Darauf sehen wir nicht, und denken nicht dran, daß GOtt so nahe sei, und sollte wohl daran zweifeln, daß mir GOtt einen Mund gemacht habe; sondern meinen, es sei die Sprache uns angeboren. Aber es ist niemand auf Erden, der ein Wort reden könnte, wenn es nicht GOtt gäbe. Wir schlagen es in Wind, und meinen, wir haben es pur plumpsweise.

8. Nun Moses ist von GOtt gefangen und auf allen Seiten beschlossen. Denn GOtt sagt zu ihm: Fürchtest du den König? Fürchtest du dich so sehr? Wenn dich die Unmöglichkeit abschreckt zu glauben, so hast du die Wunderzeichen; kannst du nicht reden, so will ich in dir reden. Also kann Moses nicht weiter, denn er sieht, daß GOtt seine Entschuldigung nicht annimmt; darum spricht er: Du hast mich beschlossen, daß ich gar rede; denn du willst mir Stärke, Kräfte und alles Vermögen selbst geben. Noch spricht er: Ich mag es nicht thun; und zeugt nichts mehr an, denn daß er spricht: „Nein Herr, sende, welchen du senden willst“; als sollte er sagen: Sende einen andern, wen du willst, ich bin es zufrieden; allein sende mich nur nicht.

9. Als nun Moses aus seinem eigenen Sinne und Willen diesen Beruf wegwerfen will, da wird der HErr sehr zornig über Mosen, und sprach:

\begin{schriftzitat}
V. 14–17. Ei, weiß ich denn nicht, daß dein Bruder Aaron, aus dem Stamme Levi, beredt ist? Und siehe, er wird heraus gehen dir entgegen, und wenn er dich sieht, wird er sich von Herzen freuen. Du sollst zu ihm reden, und die Worte in seinen Mund legen, und ich will mit deinem und seinem Munde sein, und euch lehren, was ihr thun sollt. Und er soll für dich zum Volke reden; er soll dein Mund sein, und du sollst sein Gott sein. Und diesen Stab nimm in deine Hand, damit du Zeichen thun sollst.
\end{schriftzitat}

10. Da muß Moses weichen. Zuvor hat GOtt freundlich mit ihm gehandelt, aber allhier ist er nun zornig. Und wenn GOtt zürnen will, so muß einer bleiben. Also hat Moses einen stattlichen Beruf, auf daß wir sehen, wie GOtt so wunderlich sein Werk hinausführt mit
\quellseite{387}
den Menschen, da sichs niemand zu versieht, sie auch selbst nicht bedenken dürfen. Denn alhier muß dieser Stammler und arme Bettler ausrichten, das sonst vier Könige nicht vermöchten zu thun. Und ist also Moses gefangen, daß er spricht zu GOtt: Ich will dir weichen. GOtt hat mit vielen Worten mit Mose gehandelt, daß es schier Sünde und Schande ist.

11. GOtt spricht zu ihm: „Du sollst sein Gott sein, und er soll dein Mund sein.“ Wie kommt Moses zu diesen Ehren? GOtt spricht zu ihm: Du hast die Worte, du sollst Aaron sagen, wie er predigen soll, und Aaron soll sich nach deinen Reden richten, daß er nicht länger, weiter, noch kürzer rede, denn du ihm vorschreibst, und wie du es von GOtt empfangen hast.

12. Alhier erhebt sich die Frage: warum Moses Aarons Gott heiße? In der heiligen Schrift werden die Christen und sonderlich Obrigkeit, Fürsten, Richter, oft Götter genannt, als im 82. Psalm, V. 6.: „Ich habe gesagt, ihr seid Götter und Kinder des Allerhöchsten“; und im 21. Capitel des andern Buchs Mosis, V. 5. 6., wird gesagt: „Wenn eine Sache vorfällt, daß ein erkaufter Knecht ein Weib genommen, und Kinder gezeuget hat, und nicht will frei werden, so bringe ihn sein Herr vor die Götter.“ Ein Mensch, so GOttes Wort hat, wird Gott genannt, und die es hören, und von ihm annehmen, die sind Aaron; und der Mund, der es ihnen gibt, ist ihr Gott.

13. Dieses wird um aller derer willen geschrieben, die da GOttes Wort führen und reden. Denn Moses ist nicht allein Gott, sondern die sind auch Götter, so da GOttes Wort haben. Denn wir sind GOttes Kinder [Joh. 1, 12. Röm. 8, 15.], und welche GOttes Wort haben und GOtt glauben, die haben GOttes Geist und Kraft, auch die göttliche Wahrheit, Weisheit, Herz, Sinn und Muth, und alles, was GOtt ist. Wenn er denn dieses alles hat, so GOttes ist, so muß er auch GOttes Namen haben. Derhalben ist solcher Mensch theilhaftig worden aller Güter GOttes, und auch des göttlichen Namens, wie es St. Petrus auch in seiner [zweiten] Epistel [Cap. 1, 4.] sagt. Denn, wie GOtt ein HErr aller Dinge ist, also werden sie, die Christen, auch genannt, daß sie sind weise, gerecht und lebendig etc. Gleichwie ein Vater seinem Sohne mittheilt und auserbt Namen und Gut, wie der Vater heißt, also hat der Sohn auch einen Namen. Der Sohn bekommt und erbt des Vaters Namen und Gut; also werden auch Alle Götter geheißen, die GOttes Wort haben.

14. Der lateinische Text hat wollen klug sein, und diese Worte also gegeben: Tu eris ei in his, quae ad Deum pertinent. Aber es heißt also: Du, Moses, sollst Aarons Gott sein, ich setze dich zum Gott über Aaronen. Denn Moses führt GOttes Wort. Wer Mosen verachtet, der verachtet GOtt, denn er führt GOttes Wort. Wie der HErr Christus [Luc. 10, 16.] auch sagt: „Wer euch höret, der höret mich.“ GOtt kann nicht anders reden, denn er will durch Menschen reden, durch mich und dich. Und wenn er auch gleich durch eine steinerne oder hölzerne Säule redete, dennoch so ist es GOttes Wort. Wird es nun verachtet, so wird GOtt selbst verachtet, und wer Mosen, den Mund Gottes verachtet, der da GOttes Wort führt und GOttes Zeuge ist, der verachtet GOtt selbst. Also regiert GOtt seine Christen, also hoch preist und ehrt er sie. Das ist nun die Berufung Mosis. Und er geht dahin. Man hat ihn nicht auf behangen Wagen geführt, denn er ist nicht reich gewesen. Was er für Befehl und Worte an den König habe, das habt ihr gehört.

\begin{schriftzitat}
V. 18–26. Mose ging hin, und kam wieder zu Jethro, seinem Schwäher, und sprach zu ihm: Lieber, laß mich gehen, daß ich wieder zu meinen Brüdern komme, die in Egypten sind, und sehe, ob sie noch leben. Jethro sprach zu ihm: Gehe hin mit Frieden. Auch sprach der HErr zu ihm in Midian: Gehe hin und zeuch wieder in Egypten; denn die Leute sind todt, die nach deinem Leben stunden. Also nahm Mose sein Weib und seine Söhne, und führete sie auf einem Esel, zog wieder in Egyptenland, und nahm den Stab GOttes in seine Hand. Und der HErr sprach zu Mose: Siehe zu, wenn du wieder in Egypten kommst, daß du alle die Wunder thust vor Pharao, die ich unter deine Hand gegeben habe. Ich aber will sein Herz verstocken, daß er das Volk nicht lassen wird. Und sollst zu ihm sagen: So sagt der HErr: Israel ist mein erstgeborner Sohn; und ich gebiete dir, daß du meinen Sohn ziehen lässest, daß er mir diene. Wirst du dich deß\footnote{Im Original: daß.} weigern, so will ich deinen erstgebornen
\end{schriftzitat}

\quellseite{388}
Sohn erwürgen. Und als er unterwegs in der Herberge war, kam ihm der HErr entgegen, und wollte ihn tödten. Da nahm Zipora einen Stein, und beschnitt ihrem Sohne die Vorhaut, und rührte ihm seine Füße an, und sprach: Du bist mir ein Blutbräutigam. Da ließ er von ihm ab. Sie sprach aber Blutbräutigam, um der Beschneidung willen.

15. Da kommt nun ein Stücklein drein: Als Moses abgeschieden ist von seinem Schwäher Jethro, dem Priester in Midian, da zieht er in Egypten und hat den Stab, damit er Zeichen thun soll, und dem König sagen: Laß mein Volk gehen; ich gebiete dir, daß du meinen erstgebornen Sohn, Israel, ziehen lässest. Wie geht es nun dem armen Mann unterwegs? Da fällt eine neue Historie drein. Denn da Moses in Egypten zieht, begegnet ihm ein Unfall oder Todesfall, daß er geschwind und plötzlich krank wird, und seine Frau, Zipora, nicht anders meint, er würde sterben; wie es denn der Text meldet, wiewohl nicht mit ausdrücklichen Worten. Und das war die Ursache, daß Moses, vom Stamme Abrahä, seinen Sohn nicht beschnitten hatte.

16. Da dies sein Weib nun spürt, daß solches die Ursache sei, daß ihr Mann sterben soll, daß der Sohn nicht beschnitten wäre, da nimmt sie einen scharfen Stein für ein Messer, und beschneidet den Sohn, wird kleinmüthig und zornig, wie die Weiber pflegen, spricht: „Du bist mir ein Blutbräutigam“; als sollte sie sagen: Es kostet Blut, daß du mein Mann worden bist, und ich dein Weib bin, deßhalb man muß Blut vergießen. Darum nennt sie ihn einen Blutbräutigam, um der Beschneidung willen ihres Kindes, welche unter den Heiden eine Schande war.

17. Warum beschreibt Moses dies Gaukelwerk? Warum mengt der Heilige Geist dies närrische, tolle Ding unter so große, wichtige und hohe Händel? Ist Moses irgend trunken gewesen? Ihr habt gehört, daß Abraham der erste Patriarch ist, dem befohlen worden, die Knäblein zu beschneiden. Das war ein schändlich, schmählich, toll und unsinnig Gebot vor dem Vernunft, daß sie sich selbst beschneiden sollten, und an dem Orte, den man nicht gerne nennen und hören mag. Aber Abraham ist GOtt gehorsam, und läßt sich beschneiden, ob es wohl ein schmählich und schändlich Ding ist. Aber es ist darum geschehen, daß GOtt unsern Sinn brechen, und die menschliche Weißheit zur Thorheit machen wollte. Deshalb haben die Juden viel Hohn und Spott unter den Heiden hören müssen. Man hat spöttlich und höhnisch ihnen nachgeredet, und die Beschneidung für ein närrisch Ding geachtet. Darum wollten viele die böse Nachrede nicht dulden noch leiden, und ließen sich nicht beschneiden.

18. Wie man denn auch unter uns Christen viel findet, die der Welt Lästerung, da man uns heißt Rotten, Secten und Ketzer, nicht ertragen können, und lassen eher das Christenthum darüber fahren. Also hat Mosis Weib auch gethan; da er mit ihr ein Gast im fremden Lande gewesen, hätte er für seine Person den Sohn wohl gerne beschneiden lassen, aber er hat nicht gekonnt, sein Weib hat nicht gewollt in diese Schmach treten, daß ihr Sohn beschnitten würde, sondern er sollte leben, wie man unter den Heiden lebte, sie wollte diesen Kleck und Schandfleck nicht haben. Denn die Beschneidung war eine Schande unter den Heiden. Da muß Moses leiden, was zu leiden ist, und das Kind unbeschnitten mit sich nehmen. Aber auf dem Wege da greift GOtt Mosen an, und will ihn tödten, darum, daß er des Weibes Halsstarrigkeit gefolgt und seinen Sohn nicht hat beschneiden lassen. Da wird das Weib zornig, und will lieber den Sohn sterben lassen, denn daß ihr der Mann getödtet würde, und spricht: Es kostet mich nun Blut, daß du mein Mann sein sollst. [Sie] hat sich also geschämt, daß sie nicht will fort ziehen, sondern kehrt wieder zurück anheim mit den Kindern; wie hernach im 18. Capitel, V. 1. ff., geschrieben steht, daß Jethro, Mosis Schwäher, gehört, was GOtt durch Mosen gethan; darum nimmt er Ziporam, Mosis Weib, so zurück gezogen, und seine zween Söhne, und kommt zu Mose in die Wüste. Daß sie aber hier beschneidet das Kind, thut sie allein darum, daß sie Mosen errettet.

19. Warum ist dies nun geschrieben? Wenn unser HErr GOtt seine Sache will hinausführen, so fähet er es also an, daß die Welt darüber zu Narren wird, und daß man es nicht glauben kann, und weder ich noch du also gedenken mögen, daß es solle hinausgehen. Denn wir sehen immerdar in diesen Gedanken, GOtt werde es machen, wie wir es wollen, und wenn es nicht
\quellseite{389}
also geht, so sei es falsch und unrecht; wir können uns nicht darinne brechen. Ich meine auch oft, daß ich viel wisse; aber es ist nicht wahr. Und viel Geister rühmen sich des noch heutigen Tages also, aber es ist vergeblich. GOtt hat Mosen mit gewaltigem Zeichen und trefflichen Worten bisher berufen, daß er nicht mehr thun könnte, will dazu auch in seinem Munde sein; dennoch, da er will die Sachen anfangen, da will GOtt ihn tödten. Das geht eben, als wenn GOtt mich hieße den Kaiser todtschlagen, und ehe denn ich gen Prata,\footnote{Prata oder Pratau, auch Brate, ist das nächste Dorf bei Wittenberg, jenseits der Elbe.} auf das Dorf käme, so wollte mich GOtt tödten. Ist das nicht närrisch angefangen?

20. Also heißt er hier Mosen ein trefflich Werk angreifen, und will ihn doch zuvor tödten. So that GOtt mit Abraham auch [1 Mos. 22, 1. ff.]. Er gab ihm die Verheißung, daß sein Same sich mehren sollte, als der Sand am Meere; da war Abraham ganz greis, kein Sohn Isaak würde lange leben, und viel Kinder haben. Aber nicht lange hernach heißt Gott Abraham, diesen Sohn Isaak opfern und schlachten. Wie kann man allhier diese zwei Dinge zusammenbringen? Wenn die Vernunft dieses ansieht, so spricht sie: GOtt sei toll; der Teufel rede also, und nicht GOtt; denn Gott sei zu wahrhaftig und halte über seinen Verheißungen. Und was dergleichen geboten wird, das ist nicht ein gering, schlecht Ding, nämlich Gut, Ehre und Freunde nehmen lassen, und auch zuweilen den Leib dazu. Zwar das ginge noch hin; aber wenn GOtt uns sein Wort und Verheißung gibt, und nimmt es wieder hinweg, das ist ein Großes.

21. Es ist nicht eine geringe Tentation, sondern die auch die allerhöchsten und gelehrtesten Geister betrüget, die da meinen, es stehe wohl um sie, und gehen sicher herdurch, wie sie gedenken. Aber GOtt macht es bald, daß sie zu Narren werden und fliehen müssen. GOtt macht sie kühn, auf daß er von ihnen gefürchtet werde, und daß sie wissen, daß man von seinem Worte nicht weichen soll; und dennoch versucht er die Seinen so hoch, daß sie oft müssen weichen. Mosi ist allhier zu Sinne gewesen, als müßte er sterben, und hat gedacht: Wo sind nun die Worte, die GOtt zuvor zu mir geredet hat? Item: Wo bleiben die Wunderzeichen, die ich empfangen habe auszurichten? Noch hat es müssen fest stehen bleiben, daß er der Mann sein würde, der solches ausrichten sollte, und müßte ehe dieser Moises von den Todten wieder auferstehen. Darum muß er also schließen: GOtt hat mir verheißen, ich soll sein Volk aus Egypten führen: sterbe ich nun, so muß ich doch aus dem Grabe wieder herauskommen, und doch wieder lebendig werden. Denn ehe Gott lügen sollte, so müßten eher hundert Mosen an dieses einigen Statt hervorkommen.

22. Diese Tentation kommt Mosi in der Todesnoth, daß er gedenkt: Ich soll sterben, und habe doch GOttes Wort und seine Sacramente, daß ich solle leben. In diese Klemme bringt GOtt den Mosen, daß er in einem Augenblick sterben und wieder leben muß. Wer kann das fassen? Aber man soll allezeit auf der ersten\footnote{Eislebenische: auf die ersten.} Verheißung stehen, und auf den ersten Befehl GOttes sehen, ob gleich darnach tausenderlei Tode kämen, so soll man doch sagen: Tod hin, Tod her, GOtt leugt mir nicht, ich werde leben.

23. Und solches ist darum geschehen, auf daß Moses ein Ausbund von einem christlichen Manne würde, der da eigentlich wüßte, GOtt könne nicht lügen, und daß die ersten Worte GOttes bei ihm zu Kräften kämen, und er ein Held und starker Mann im Glauben und voller Geistes gemacht würde.

24. Also handelt GOtt mit uns, wenn er sein Wort zu Kräften in uns bringen will, daß uns dennoch wehe thut. Wer nun also auf die erste Verheißung und Wort unsers HErrn GOttes trauen und sich verlassen kann, wohl demjenigen! Denn, was hernach folgt, das kann er gehen lassen, wie es geht. Er hat uns verheißen das ewige Leben und seine Gnade. Wenn nun solche Gedanken uns vorkämen und einfielen: Hörst du, daß du sollst sterben und verdammt werden? Dann soll ich gehen zu seiner ersten Verheißung, und durch sein Wort und Werk dringen und kämpfen wider GOttes Wort und Werk. Denn das Wort und Werk, so hernach folgt, geschieht darum, daß GOtt versucht unsern Glauben, ob wir fest daran halten wollen.

25. Denn Abraham hat GOtt verheißen den Isaak, und so viel Nachkömmlinge und Kinder,
\quellseite{390}
als Sand am Meere; noch spricht er zu ihm [1 Mos. 22, 2.]: Tödte mir den Sohn und schlachte ihn etc. Dieses ist der ersten Verheißung gar entgegen. Aber es ist eine Tentation, da GOtt versuchen will, ob Abraham auch fest im Glauben an seinem Worte halten wolle. Also ist Moses auch gewiß, daß er den König in Egypten stürzen, und das Volk Israel ausführen solle; das sollte und müßte gewiß geschehen. Nun versucht GOtt den Mosen, ob er halten wolle oder nicht. Sein Weib hat es nicht erwehren können, darum zog sie gerne wieder zurück. Denn sie gedenkt: Ach! ich hoffte, ich wollte eine große Frau worden sein, Freude, Ehre und Herrlichkeit in meines Mannes Ehestande erlebt haben; nun ist es eitel Tod und Noth, [ich] ziehe mit ihm aus meinem Vaterlande in das Elend. O trolle dich nur wieder zurück zum Vater in Midian.

26. Also wird Moses probirt und versucht mit diesem Stücklein, auf daß er fest würde im Glauben. Und das ist auch die Ursache, daß dieser Text ist vom Heiligen Geist beschrieben worden. Es lautet wohl ungereimt und seltsam genug, daß GOtt also mit uns umgehe, aber wir müssen es gewöhnen. Denn Gott thut es darum, daß wir ihm allein in seine Hand sehen sollen.

27. Also ist Moses hingegangen mit Aaron, und that alle Zeichen vor dem Volke, die ihm GOtt befohlen hat. Und die Kinder Israel glaubten, wiewohl schwächlich, und wurden froh, dankten GOtt, daß ihre Erlösung nun angehen sollte.

\subsection{Allegoria oder geistliche Deutung des vierten Capitels}

28. Im vierten Capitel wird Mosi befohlen, daß er das Volk ausführen solle; darauf ist er zuvor vergewissert mit dem Gesichte des feurigen Busches, davon wir gesagt, daß es den HErrn Christum bedeute. Nun, wenn er den Busch, das ist, diesen HErrn Christum bei sich hat, so ist das Gesetz erfüllt. Das Gesetz und der HErr Christus müssen zusammen kommen.

29. Da folgen weiter andere Figuren. Mose soll Wunderzeichen haben, auf daß man ihm glaube; so hat er einen Stab in seiner Hand, den wirft er weg, daraus wird eine Schlange. Item, seine Hand wird aussätzig, und das Wasser wird zu Blut. Das sind drei Wunderzeichen, dadurch er das Volk überreden soll, und sie versichern, daß sie ihm glauben.

30. Ich weiß nicht, ob ich es werde treffen oder nicht mit diesen Allegorieen. Denn ich habe zuvor gesagt, daß es nicht eine vornehmlich Kunst sei, alles auf heimliche Deutung führen; darum wollte ich gerne damit den schändlichen Deutungen Origenis, Hieronymi und anderer gewehrt haben, die von den guten Werken in das Herz und Gewissen hinein gedeutet werden, auf daß dagegen alle Predigten von dem göttlichen, gewissen Worte bestätigt würden. Wenn man das nicht thun kann, so ist es viel besser, man lasse die Allegorieen anstehen. Denn Origenis und anderer mehr Allegorieen sind ungereimt und heidnisch.

31. Daß er nun den Stab wegwirft, ist vom Predigtamte zu verstehen, wenn man GOttes Wort nicht in der Faust hat, daß es nicht im Schwange geht; da ist denn solche Lehre anders nicht, denn wie ein greuliches Gift, das eine Schlange bei ihr hat. Da muß Moses selbst vor ihr fliehen. Also geht es zwischen dem Gesetze und menschlichem Herzen, wenn die Leute GOttes Gnade und den Heiligen Geist nicht haben, das Gesetz zu halten. Dann so ist das Gesetz ein Tod und eine Schlange. Ich habe dann kein Herz, Lust noch Begierde zum Gesetze, es schmeckt mir nicht, und je härter darauf gedrungen, und je mehr es gepredigt wird, je feinder ich ihm werde [Röm. 7, 23.]. Denn mein Herz ist viel anders gesinnt; je mehr man das Gesetz treibt, je weniger man schafft. Das Herz spricht: Ach! daß du nicht das Gesetz: Du sollst nicht ehebrechen, gegeben wäre, dann wollte ich frei sein, ginge hin und her, und triebe meine Buhlerei. Item, wenn das Gesetz hinweg wäre: Du sollst nicht stehlen, ei, wie wollte ich geizen, wuchern und scharren.

32. Aber wenn das Gesetz, als der Stab aus der Faust, kommt, so schrecket es, denn es sagt: Stipendium peccati mors [Röm. 6, 23.]. Nun kommt Moses auch ins Evangelium, daß man nicht allein das Gesetz nach der Schrift, sondern auch das Evangelium höre; und wenn das Evangelium auch trifft, so wird man wieder fröhlich. Das Herz muß eigentlich erschrecken und fliehen; aber wenn das Evangelium auch recht trifft, denn wird es lieblich. Denn es tröstet den Sünder und verkündigt Vergebung der Sünde, wie Christus zum Gichtbrüchigen sagt [Matth. 9, 2.]:

\quellseite{391}
\begin{quote}
„Mein Sohn, sei getrost, deine Sünden sind dir vergeben.“
\end{quote}

33. Die aussätzige Hand bedeutet das heuchlerische Leben, so man in der Welt führt. Denn Aussatz ist falsche Lehre und Menschentand. Wenn die Menschen nicht recht glauben, und entweder durch ihre Kräfte und Vermögen zum ewigen Leben geführt, oder durch GOttes und Menschen Gesetze auf ihre eigenen guten Werke gewiesen werden, da ist die Hand aussätzig; und mit den Menschen sollte Moses umgehen. Aber wenn man den Aussatz geheilt hat, und die Hand wieder in Busen steckt, so wird sie wieder rein; das ist, wenn das Gesetz erstlich falsch verstanden ist, daß der Heilige Geist kommt und es recht predigt, nämlich, daß wir nichts vermögen und Sünder sind, und uns allein durch Christum geholfen müsse werden, da werden wir rein und heilig.

34. Daß auch das Wasser zu Blut wird, bedeutet Menschenlehre und das Gesetz; denn die tödten und bluten. Sintemal das Gesetz nichts thut, denn würgt und mordet. Das hat unsere Lehre von Menschensatzungen in den Klöstern wohl beweiset, die ist anders nicht gewesen, denn das Wasser Egypti, so zu Blut ist worden, da die unschuldigen Gewissen dadurch ermordet sind. Denn je mehr man drin hangt, je mehr man stirbt. Damit will Moses anzeigen, es sei zu thun um reine, rechtschaffene Lehre im Volke GOttes, auf daß das Gewissen auf keine andere Lehre oder Werk baue oder vertraue, denn allein auf Christum. Das sind die vier Capitel.

\section{Das fünfte Capitel}

\begin{schriftzitat}
V. 1. 2. Darnach gingen Mose und Aaron hinein und sprachen zu Pharao: So sagt der HErr, der GOtt Israel: Laß mein Volk ziehen, daß mir's feire in der Wüste. Pharao antwortete: Wer ist der HErr, deß Stimme ich hören müsse, und Israel ziehen lassen? Ich weiß nichts von dem HErrn, will auch Israel nicht lassen ziehen.
\end{schriftzitat}

1. Bisher haben wir gehört, wie GOtt Mosen zum Hauptmann und Fürsten ordnet, der sein Volk aus Egypten führen soll, und hat ihn bestätiget mit Wundern, Verheißungen und Wunderzeichen, daß er als ein Herzog gerüstet und geharnischt ist. Nun wollen wir hören, wie er zum König Pharao gehen wird, und sein Amt der Erlösung angreifen, auch was er darinnen ausgerichtet, und es mit ihm für einen Fortgang haben wird.

\begin{schriftzitat}
V. 3–9. Sie sprachen: Der Hebräer GOtt hat uns gerufen, so laß uns nun hinziehen drei Tagreisen in die Wüste, und dem HErrn, unserm GOtt, opfern, daß uns nicht widerfahre Pestilenz oder Schwert. Da sprach der König in Egypten zu ihnen: Warum macht ihr (du Mose und Aaron), daß das Volk sein Werk läßt anstehen? Gehet hin an eure Dienste. Weiter sprach Pharao: Siehe, des Volks ist schon zu viel im Lande, und ihr wollet sie noch feiern heißen von ihrem Dienste. Darum befahl Pharao desselben Tages den Vögten des Volks und ihren Amtleuten, und sprach: Ihr sollt dem Volke nicht mehr Stroh sammeln und geben, daß sie Ziegel brennen, wie bis anher; lasset sie selbst hingehen, und Stroh zusammen lesen. Und die Zahl der Ziegel, die sie bisher gemacht haben, sollet ihr ihnen gleichwohl auflegen, und nichts mindern. Denn sie gehen müßig, darum schreien sie, und sprechen: Wir wollen hinziehen, und unserm GOtt opfern. Man drücke die Leute mit Arbeit, daß sie zu schaffen haben, und sich nicht kehren an solche Worte.
\end{schriftzitat}

2. Da geht nun die Erlösung an. Ihr habt aber droben [Cap. 1, § 52 ff.] gehört, daß man in der heiligen Schrift, oder in diesen Historien, sonderlich auf GOttes Wort sehen und fleißig Achtung haben solle, als auf das vornehmste Hauptstück, darauf man viel mehr merken müsse denn auf gute Werke. Wiewohl die guten Werke, welche von Menschen gethan werden, viel einen größeren Schein haben, und den Leuten das Maul mehr aufsperren, denn das göttliche Wort, welches einfältig ist, schlecht und recht vor sich geht [Ezech. 1, 12.].

3. Das ist aber die Art des göttlichen Worts, und also geht es ihm. Wenn es anfangen will, seine Kraft und Gewalt zu erzeigen, so wird es
\quellseite{392}
zuvor geschwächt. Denn unsers HErrn GOttes Gewalt muß erst vor der Welt Schwachheit werden, und seine Weisheit muß eine Thorheit sein, wie St. Paulus in der 1. [Epistel] zum Corinth[ern] am ersten, V. 21. 25., klagt. Wiewohl GOttes Thorheit die größte Weisheit ist, über alle Weisheit auf Erden, und GOttes Schwachheit, daß er sich schwach stellt, ist stärker denn alle Menschen sind, und die Armuth GOttes ist viel köstlicher denn aller Welt Reichthum. Aber es gehört Glaube dazu und ein sonderlicher Mensch, der es verstehe, daß es dem göttlichen Worte also gehe.

4. Der Apostel St. Paulus ist sehr kühn, daß er unserm HErrn GOtt zuschreibt, er sei schwach, krank und arm, und daß es seinem Worte auch also geht in der Welt. GOtt redet von großen Sachen, als, erretten, erlösen, selig machen, und helfen von allen Sünden etc. Es sind große, mächtige Worte; aber wenn es angeht und in Schwang und zum Treffen kommen soll, so stellt sich’s also, gleich als sei es der Kleist, der über ein Jahr gefallen sollte.

5. Also geschieht es allhier auch, wenn Moses zu Pharao spricht: Laß dies Volk ziehen; bald fällt dies Wort in die Asche und in den Koth. Denn Pharao spricht: „Wer ist der HErr? Ich weiß von keinem HErrn, ich will das Volk Israel nicht ziehen lassen“; und setzt mehr Frohnvögte über sie, wird sehrer erzürnt, und ordnet andere Amtleute, die sie härter drängen denn zuvor, werden nunmals viel ärger geplagt.

6. Darnach fährt der Tyrann Pharao fort, legt sich an Mosis Hals, schilt Mosen, und hält ihn für einen Aufrührer und Meutmacher.\footnote{Meutmacher = Anstifter von Meuterei.} Also kommt Moses in große Gefahr, bekommt einen groben Filz\footnote{Filz = Verweis.} und dürres Capitel, da er spricht: Ihr seid müßig, habt nichts zu thun, des Volks ist viel, beschweret sie noch mit mehr Frohndiensten. Und steht das Volk in großer Noth und Arbeit; Pharao aber treibt nur seine Spötterei draus. Er vermeint das göttliche Wort unterzudrücken, und macht es schwach und hält es für eine Thorheit, das muß seine Posserei\footnote{Posserei = Possenspiel.} und Gelächter sein. Also, Moses, der mit den Seinen GOttes Wort für sich hat, muß drüber für einen Narren von Pharaone und den Seinen, ja, auch wohl sonst in der Welt geachtet werden, und richtet nichts aus.

7. Ueber das werden sie selbst uneins drüber, das Volk und Moses, da das Volk spricht: Hättet ihr uns lassen arbeiten, wie zuvor; nun sinken wir gar vor den Egyptern. Ihr habt ihnen das Schwert in die Hand gegeben, damit sie uns nun schlagen. Ist das erretten und erlösen? Ihr führt uns erst recht zur Marter, daß wir desto mehr geplagt werden; ihr wollt das Licht anzünden, und löscht es gar aus. Und er, Moses, spricht auch: Lieber HErr GOtt, warum heißt du mich etwas thun, und ziehst darnach deine Hand ab, und lässest mich drinnen stecken? Wie denn in diesem Capitel weiter im Texte folgt:

\begin{schriftzitat}
V. 22. 23. Mose aber kam wieder zu dem HErrn, und sprach: HErr, warum thust du so übel an diesem Volke? Warum hast du mich hergesandt? Denn, seit dem, daß ich bin eingegangen zu Pharao, mit ihm zu reden in deinem Namen, hat er das Volk noch härter geplagt und du hast dein Volk nicht errettet.
\end{schriftzitat}

8. Also klagt der Prophet Jeremias auch, und spricht zu GOtt [Cap. 20, 7.]: Du hast mich betrogen und verführet. Ei, wie sollte GOtt einen verführen? Sagt er doch [Jer. 1, 18.]: „Ich will dir eine eherne Stirne, und eiserne Mauer sein“, und läßt dennoch den Propheten Jeremiam in das Gefängniß werfen.

9. Aber das ist GOttes Art, daß er sein Wort wunderbarlich führt, und wenn er sein Werk auch angreift, so läßt sich es ansehen, als wollte nichts daraus werden, ja, das Gegenspiel geschieht wohl. Aber warum thut dieses GOtt? Darum, er will sein Werk desto wunderbarer ausrichten, und das göttliche Wort soll seine Kraft und Macht desto mehr in der Schwachheit sehen lassen und beweisen.

10. Denn durch Thorheit macht er Weisheit, durch Krankheit macht er Stärke und Gesundheit, durch Nichts macht er Alles, Jes. 63, 1. Also läßt er auch sein Wort krank und schwach werden. Alsdann sind die Widersacher des göttlichen Worts viel stärker und klüger, und werden hoffärtig. Aber GOtt gedenkt also: Harret, ich will euch recht setzen, daß ihr anlaufen sollt. Und wenn sie denn am allerstärksten sind, so bricht GOtt herein und schlägt zu Boden der Welt und aller seiner Feinde Stärke, Weisheit, Kunst und alles.

\quellseite{393}
11. Das ist uns zum Troste geschrieben, auf daß wir lernen in unsern Berufen und Aemtern an GOttes Wort fest halten, und von uns selbst abtreten, aber auf GOtt unsere Hoffnung und Vertrauen setzen. Denn um unsertwillen läßt GOtt sein Wort auch zur Thorheit und schwach werden, und bricht doch gleichwohl hindurch; wie denn diese Erlösung noch folgen mußte mit den Kindern von Israel, ob es sich gleich nicht ließ dafür ansehen, und ob es wohl schwerlich und unmöglich zugehet. Denn also sagt GOtt im Ende dieses fünften Capitels:

\begin{quote}
Cap. 6, 1. Nun sollst du sehen, was ich Pharao thun werde; denn durch eine starke Hand muß er sie lassen ziehen, er muß sie noch durch eine starke Hand aus seinem Lande von sich treiben.
\end{quote}

12. Also sind auch alle Christen GOttes Schwachheit und Thorheit, aber GOtt macht sie wieder stark und weise; wie auch im Propheten Joel am dritten Capitel, V. 15., gesagt wird: „Der Schwache achte sich stark.“ St. Petrus, der Laie und arme Fischer, hängt sich an alle heiligen, gelehrtesten und klugen Leute in der Welt, ja, auch an das ganze römische Reich; wiewohl es ein geringes Ansehen hat, daß St. Petrus mit der Predigt des Evangelii das römische Reich meistern sollte. Und einer hätte mögen sagen: O wäre GOtt klug, so griffe er es also närrisch nicht an. Aber St. Petrus war (gleichwie auch Moses allhier ist) GOttes Schwachheit und Thorheit, und mußte den Kopf und sein Leben herstrecken. Aber gleichwohl hat er mit dem Evangelio das römische Reich ersäuft.

13. Denn dies ist die Natur des göttlichen Worts, so wir wohl lernen sollen, daß wenn Pharao stolz ist, pocht, tobt und müthet, so stellt sich GOtt und sein Wort gar schwach, matt und krank, gleich als läge es gar zu Boden. Aber es wird wieder stark, reißt hindurch und überwindet alles; wie denn auch folgen wird, daß die Kinder von Israel werden durch das rothe Meer geführt, und Pharao darinnen ersäuft.

14. Das andere Stück in diesem Capitel ist von des Teufels Art und Eigenschaft, davon auch Christus im Evangelio Matthäi am zwölften Capitel, V. 29., sagt: „Wenn der starke Gewappnete seinen Palast bewahrt, so ist das Seine zufrieden“ etc. [Luc. 11, 21.] Der Teufel kann das wohl leiden, daß man sein Reich läßt gehen und bleiben, und er die Gewissen und Seelen gefangen behalten mag. Aber wenn GOttes Wort kommt, und sein Reich stürmt und angreift, da wird er toll und unsinnig und kann es nicht leiden; und wenn er geduldig dazu wäre, so würde er viel mehr gewinnen, ausrichten und schaffen, denn also. Aber das thut er nicht; alsbald man ihn nur ein wenig anrühret, so schmerzet es ihn und will es nicht leiden. Dieses sehen wir in den evangelischen Historien [Matth. 8, 32. Marc. 9, 26.], wenn der HErr Christus den Teufel austreibt, so schäumt er, sprühet um sich, und treibt seltsame Geberde, er bebt und reißt um sich, als wäre er unsinnig und thöricht; also ungerne weicht er. Und wie er nun in den Menschen thut, also handelt er auch in seinem Reiche und an seinem ganzen Leibe.

15. Der Papst ist auch angegriffen, sein Ablaß, Wallfahrt, Messe, Heiligen-Anrufen und anderer Lügentand sind offenbar und zu Schanden gemacht. Da siehe, wie hier sein Reich tobt und müthet (denn der Teufel muß ausfahren); siehe, wie der Papst uns lästert, schändet, vermaledieet und verdammt, und wie er sprühet und flucht.

16. Also geht es auch allhier in diesem Capitel zu, da GOttes Wort und Verheißung kommt, und dem Volke Israel helfen will, da wird dies Volk härter gehalten, denn zuvor nie. Denn Pharao nimmt dem Volk Israel auch das Stroh. Dennoch, je härter er das Volk hält, je weniger er gewinnt. Denn GOttes Wort ist GOttes Kraft und Gewalt wider des Teufels Macht und Gewalt [Röm. 1, 16. 2 Cor. 12, 9.]. Darum so ist es vergeblich, daß sich der Teufel also sperrt, zerrt und wehrt, und also feindlich hält. Denn dadurch hilft der Teufel selbst, daß das Volk Israel los wird; wie denn allhier Pharao geschieht, welcher über seiner Tyrannei und Toben kommt um Leib und Gut, und muß im rothen Meere ersaufen. Wer doch dies allein glauben könnte, daß, je härter der Teufel in Todesnöthen, und dergleichen in andern Anliegen, uns zwingt, je gewisser GOtt da bei uns sei, und je steifer über uns halte.

\quellseite{394}
\section{Das sechste Capitel}

\begin{schriftzitat}
V. 2–5. Und GOtt redete mit Mose, und sprach zu ihm: Ich bin der HErr, und bin erschienen Abraham, Isaak und Jakob, daß ich ihr allmächtiger GOtt sein wollte, aber meinen Namen, HErr, habe ich ihnen nicht kund gethan. Auch habe ich meinen Bund mit ihnen aufgerichtet, daß ich ihnen geben will das Land Canaan, das Land ihrer Wallfahrt, darinnen sie Fremdlinge gewesen sind. Auch habe ich gehöret die Wehklage der Kinder Israel, die die Egypter mit Fröhnen beschweren, und habe an meinen Bund gedacht.
\end{schriftzitat}

1. In diesem Capitel wiederholt GOtt seine Verheißung, daß er die Kinder von Israel wolle aus Egypten führen, und das gelobte Land ihnen geben, will mit dieser neuen Zusagung des Mosis und der Kinder Israel Glauben stärken, daß sie an der Erlösung nicht zweifeln. Aber diese Worte sind ihnen umsonst gepredigt worden etc.

2. Die heilige Schrift und der Heilige Geist verwahren sich mit höchstem Fleiß, wenn sie reden von den Wohlthaten, so GOtt uns erzeigen wolle, daß sie zuvor hinwegnehmen allen Ruhm und Verdienst, und uns dahin weisen, daß alles, was GOtt in uns thut und wirkt, er darum thue, daß er seine Wohlthat und Freude darinnen gehabt, und er von Ewigkeit es also verordnet.

3. Also, damit Moses allhier, und GOttes Wort auch geschwächt und zur Thorheit gemacht werde, so geht es blödiglich und geringe an. GOtt predigt allhier, daß er erschienen sei Abraham, Isaak und Jakob, und gedenkt des Bundes, den er ihnen verheißen habe, als, das gelobte Land zu geben, und spricht: Um dieses Bundes willen soll euch das widerfahren, und nicht um eures Seufzens, Schreiens und Klagens willen, auf daß sie ja nichts hätten vorzuwenden.

4. Denn kein Wehklagen soll nicht satt und genug sein, und die guten Werke und das Leiden soll nicht gelten, man gehe denn zuvor in GOttes Wort daher. Darum spricht GOtt allhier: Dieweil ich euch verheißen habe dies Land, und ihr seid der Same, dem dies Land zugesagt ist, um dieser Verheißung willen, so sehe ich an dies Wehklagen. Denn diese Zeit, darinnen ihr also gemartert und geplagt werdet, gibt mir Ursache, zu gedenken eures Wehklagens; verfaßt sie also, daß sie nicht ein Haar breit sich rühmen könnten einiges Verdienstes. Denn die Kinder von Israel auch wenig gute Werke und eigene Gerechtigkeit und Frömmigkeit gehabt: das Volk steckte voller Unglauben, Ungeduld, Murren, und Moses voller Schwachheit und Ungehorsam. Sondern, was ihnen gegeben würde, daß es geschähe durch seine ewige Güte, um seiner Wahrheit und Verheißung willen, damit sie GOtt Danksagung für seine Gnade und Barmherzigkeit thäten, und sie nicht sagen dürften, sie wären die Leute, die GOtt um irgend eines Werks willen ansehen müßte. Darum spricht GOtt allhier: Sehet darauf, ich will euch erretten, nicht um eurer gelben Haare willen, sondern daß ich es vierhundert Jahre zuvor euren Vätern verheißen habe, und also beschlossen, daß ich es thun wolle; auf daß ihr wisset, es sei euch aus Gnaden gegeben, und euren Vätern aus Gnaden übergeben und verheißen.

\begin{schriftzitat}
V. 3. Aber meinen Namen, HErr, habe ich ihnen nicht kund gethan.
\end{schriftzitat}

5. Nun wollen wir die Worte besehen. Die Juden sprechen, daß der Name GOttes genannt sei Tetragrammaton ADONAI, das muß man lassen bleiben. Er spricht: Ich bin ihnen erschienen, also, daß ich ihr GOtt sei. Darum so haben sie GOtt erkannt und von GOtt gewußt zu reden, daß er allmächtig, wahrhaftig, gnädig und barmherzig sei. GOtt ist ihnen oft erschienen, hat mit ihnen geredet und große Wunder viel bei ihnen gethan. Und auf diese Weise haben sie ihn besser erkannt denn wir. Warum spricht er denn, daß sie ihn nicht erkannt haben? Wir müssen also deuten den Namen GOttes, daß er nicht sei so hell und klar gepredigt und gepreiset worden, als zur Zeit des HErrn Christi, der lieben Apostel, und hernach eine Weile. Derhalben Gottes Name zur selbigen Zeit kein groß Geschrei hatte in der Welt, daß er allenthalben wohl wäre erkannt worden. Ich, spricht
\quellseite{395}
er, bin ihnen erschienen; aber ich habe keinen Namen bei ihnen gehabt, das ist, die Patriarchen sind nicht die gewesen, so die Predigt und das Geschrei von mir hätten geführt, als nun du, Moses, und die Kinder von Israel thun sollen, daß mein Name in der ganzen Welt bekannt würde. Nun aber soll die Erlösung der Kinder Israel ins Werk gesetzt, und die Verheißungen erfüllt, und dem ganzen Volke Israel geholfen werden, da er zuvor nur einzelnen Personen, als Abraham, Isaak und Jakob, geholfen hat.

6. Also deuten wir es, und kommen hindurch; wer es besser weiß, dem gönnen wir es wohl. Moses hat ein Amt gehabt; von dem Amt wird allhier geredet. Sonst hat Abraham in seinem Hause auch gepredigt den Namen GOttes [1 Mos. 12, 8. 13, 4.], aber einen solchen Befehl hat er nicht gehabt, den Namen GOttes zu predigen, als Moses gethan hat, oder wie wir thun, oder wie St. Paulus, die lieben Apostel, und der HErr Christus gethan hat. Jes. 41, 25. wird gesagt: „Er soll meinen Namen predigen“; also macht er seinen Namen kund und bekannt den Patriarchen.

7. Aber allhier hat man viel lahmer Boten der Juden, von den Wirckeln dieses Namens Adonai, die wir jetzt wollen fahren lassen. Es will aber Moses sagen: Jetzt sollt ihr anheben und predigen diesen Namen GOttes, und eitel treffliche, schöne, liebliche Predigten davon thun. Aber das ist auch schwach, es gilt nicht vor dem Volk; um der harten Arbeit und großen Beschwerung willen glauben sie nicht. Jedoch hat GOtt ein Mitleiden mit ihnen und hält es ihnen zugute, hält an, predigt fort, läßt nicht ab.

\begin{schriftzitat}
V. 9. Moses sagte solches den Kindern Israel; aber sie hörten ihn nicht vor Seufzen und Angst, und vor harter Arbeit.
\end{schriftzitat}

8. Da Moses durch GOttes Wort und Zusagung gestärkt worden, ist er GOtt gehorsam, predigt den Kindern von Israel abermals von der Erlösung, vermahnt sie zum Glauben, Geduld im Kreuz, und daß sie GOttes Beruf folgen sollen.

9. Aber hier spürt man den Ungehorsam und Unglauben der Kinder Israel. Denn, da es nicht bald fort will mit der Erlösung, so werden sie Gott den Sack gar vor die Thüre, wollen seinem Worte nicht mehr glauben, auch keinen Trost zulassen. Denn, wenn es uns übel geht, so gedenkt Fleisch und Blut nicht daran, daß es solle wiederum einmal gut werden, sondern meint, nun müsse man gar verderben, so doch auf einen Platzregen und Ungewitter bald pflegt wieder ein Sonnenschein und helle, schön Wetter zu folgen [Tob. 3, 23.]. Also pflegt GOtt nach langem Kreuz und Trübsal auch wieder zu helfen, und hilft viel herrlicher, je schwerer das Leiden gewesen ist.

\begin{schriftzitat}
V. 14–20. Dies sind die Häupter des Hauses ihrer Väter. Die Kinder Ruben, des ersten Sohns Israel, sind diese: Hanoch, Pallu, Hezron, Charmi. Das sind die Geschlechter von Ruben. Die Kinder Simeon sind diese: Jemuel, Jamin, Ohad, Jachin, Zohar und Saul, der Sohn des canaanäischen Weibes. Das sind Simeons Geschlechte. Dies sind die Namen der Kinder Levi in ihren Geschlechtern: Gerson, Kahath, Merari. Aber Levi ward hundert und sieben und dreißig Jahr alt. Die Kinder Gerson sind diese: Libni und Simei, in ihren Geschlechten. Die Kinder Kahath sind diese: Amram, Jesehar, Hebron, Ussiel. Kahath aber ward hundert und drei und dreißig Jahr alt. Die Kinder Merari sind diese: Mahali und Musi. Das sind diese Geschlechte Levi, in ihren Freundschaften. Und Amram nahm seine Muhme Jochebed zum Weibe, die gebar ihm Aaron und Mose. Aber Amram ward hundert und sieben und dreißig Jahr alt etc.
\end{schriftzitat}

10. Nun folgt weiter, welches in zweien Blättern beschrieben wird, daß ein anderer mit zweien Worten begriffen hätte, und werden erzählt Mosis und Aarons Geschlechte und Stammregister. Denn, soll Moses ein Lehrer und Herzog werden, der das Volk Israel ausführe aus Egypten, so muß man auch wissen, wer er sei, aus was Stamme er geboren sei, daß man seine Ankunft wisse; und um des dritten Stammes, als Levi willen, daraus Moses und Aaron kommen, müssen die zween Stämme, als Ruben und Simeon, zuvor auch genannt werden. Denn diese beiden, Ruben und Simeon, waren älter denn Levi. Moses ist aber auch der Lea Sohn, der armen Aschenbrödelein, denn Levi ist von Jakob, und der Lea Sohn. Nun hat Levi wieder einen Sohn, als Kahath [1 Mos. 46, 11.], darnach Kahaths Sohn ist Amram, welcher Mosis Vater ist gewesen. Das ist Mosis Geschlecht. Diese alle drei, Levi, Kahath und Amram, sind
\quellseite{396}
in Egypten gestorben, und haben alle in Egypten gelebt. So ist Jochebed Mosis Mutter gewesen. Und droben [Cap. 2, 1.] wird gesagt, daß ein Mann aus dem Stamme Levi habe ein Weib genommen, als seine Muhme; da wird Jochebed Levi Tochter genannt. Aber ich glaube vielmehr, daß sie Gersons oder Merari Schwester sei gewesen, und Amrams Weibe, daß sie im andern Gliede ihm angehört und seines Geblüts gewesen sei. Jeseer ist Amrams Bruder; wie denn Elisabeth Zacharia Weib ist, welche auch Aarons Tochter ist [Luc. 1, 5.]. Aber Aarons Weib wird auch Elisiba geheißen, von dem Stamme Juda; denn diese beiden Geschlechte griffen in einander mit Heirath. Wenn die Töchter Güter hatten, mußten sie sich nicht vergreifen; wenn sie aber nicht Güter hatten, so mochten sie aus dem Geschlechte greifen, und in ein ander Geschlecht freien.

\begin{schriftzitat}
V. 26. 27. Das ist der Aaron und Mose, zu denen der HErr sprach: Führt die Kinder Israel aus Egyptenland, mit ihrem Heer. Sie sind es, die mit Pharao, dem Könige in Egypten, redeten, daß sie die Kinder Israel aus Egypten führten, nämlich Mose und Aaron.
\end{schriftzitat}

11. Aus dem Stamme Levi werden nun große Leute geboren; wie denn Moses und Aaron auch daraus gekommen sind. Und wird Moses alhier beschrieben als ein Heerführer und Hauptmann des Volks Israel, daß er aus diesem Stamme gekommen sey, und GOtt so und so mit ihm geredet habe.

\begin{schriftzitat}
V. 28–30. Und des Tages redete der HErr mit Mose in Egyptenlande, und sprach: Ich bin der HErr, rede mit Pharao, dem Könige in Egypten, alles, was ich mit dir rede. Und er antwortete vor dem HErrn: Siehe, ich bin von unbeschnittenen Lippen, wie wird mich denn Pharao hören?
\end{schriftzitat}

12. GOtt wiederholt abermal seinen Befehl und Gebot zu Mose, daß er solle die Kinder Israel aus Egypten führen; [er] will ihn getrost und muthig durch sein Wort in diesem Werke und Berufe machen. Aber Moses läßt seinen Ungehorsam und Schwachheit abermals sehen, entschuldigt sich, will nicht fort, fürchtet der Haut. Daran wir aber ein Exempel der Gebrechlichkeit und Schwachheit der Heiligen sehen, wie sie den Hund auch hinten lassen, wie Sünde immerdar in ihnen sich regt, und sie um Vergebung deßhalben täglich bitten müssen.

\section{Das siebente Capitel}

1. Wir haben gehört im fünften und sechsten Capitel dieses Buchs, wie GOtt der HErr diese zwei, Mosen und Aaron, stärkt und tröstet, daß sie das Amt, mit Pharao zu reden und das Volk Israel aus Egypten zu führen, annehmen, und gibt ihnen Wunderzeichen, als daß der Staub zur Schlange, und das Wasser in Egypten zu Blut, und die ausätzige Hand wieder rein wird. Welches alles darum geschah, daß sie muthig und kühn würden, dies mächtige, große Werk anzugreifen, und nicht gedachten, daß es durch ihre Kraft geschähe, sondern daß es GOttes selbst eigenes Werk wäre, da er wollte bei ihnen gegenwärtig sein. Auch wurde durch diese Mirakel, so Moses that in Egypten, jedermann kund, daß Moses von GOtt gesandt wäre, und daß die Ausführung des Volks Israel aus Egypten durch GOttes Geheiß und Beistand geschähe, und nicht aus freventlicher Leichtfertigkeit, oder vorwitziger und aufrührischer Weise. Diese Mirakel führten und wiesen jedermann auf GOtt, als den Stifter, Urheber und Anfänger dieser Ausführung, daß er sein Volk wollte erlösen von dem Diensthause Egypti.

2. Denn darum läßt GOtt in der Kirche seine Heiligen und Werkzeuge oft Mirakel oder Wunder thun, daß er Ehre, Lob und Preis wil davon haben. Wie denn in diesem Capitel [V. 5.] GOtt spricht: „Die Egypt[e]r sollen es inne werden, daß ich der HErr bin, wenn ich nun meine Hand ausstrecke über Egypten, und die Kinder Israel von ihnen wegführen werde.“ Die Mirakel sollen dienen, daß wir den wahrhaftigen GOtt erkennen. Wie denn auch die Wunderwerke, so der HErr Christus that, offenbarten seine Herrlichkeit, daß man an ihn glaubete [Joh. 2, 11.]. Denn es waren göttliche Mirakel, wie Nicodemus solches bekennt [Joh. 3, 2.] da
\quellseite{397}
er sagt: „Meister, wir wissen, daß du ein Lehrer von GOtt gekommen bist, denn niemand kann die Zeichen thun, die du thust, es sei denn GOtt mit ihm.“

3. Item, durch die Mirakel der Heiligen GOttes wird die christliche, wahrhaftige Lehre bestätigt. Als, die Apostel thun im neuen Testamente große Mirakel; wie denn Petrus Apost. 3, 8. den Lahmen gesund macht, im Namen und Kraft des HErrn Christi. Dadurch wird überzeugt, daß die Lehre von Christo wahrhaftig sei, denn sie werde nicht allein mit Worten von den Aposteln gelehrt, sondern Christus bekräftige sie auch mit Wunderwerken [Marc. 16, 17.], so nicht aus menschlichen Kräften gethan werden.

4. Letzlich, so wird durch Mirakel der wahrhaftige Gottesdienst von GOtt selbst gestiftet, vergewissert und bestätigt. Als, durch Elia und anderer Propheten Mirakel wird das Volk im rechten Gottesdienste versichert und bekräftigt, und vom abgöttischen, baalitischen Gottesdienste abgeführt.

5. Nun, im siebenten Capitel, hören wir, wie Pharao halsstarriger und verstockter wird, antwortet ihnen, daß er das Volk nicht wolle lassen, sondern er beschweret es noch mehr mit Arbeit und Ziegelstreichen, daß auch das Volk Israel anfängt Zeter über Mosi'en zu schreien.

\begin{schriftzitat}
V. 1. 2. Der HErr sprach zu Mose: Siehe zu, ich habe dich einen Gott gesetzt über Pharao; und Aaron, dein Bruder, soll dein Prophet sein. Du sollst reden alles, was ich dir gebieten werde; aber Aaron, dein Bruder, soll vor Pharao reden, daß er die Kinder Israel aus seinem Lande lasse.
\end{schriftzitat}

6. Mit diesen Worten wird angezeigt, wie wunderbarlich GOtt seine Sachen führt, wenn er dieselbigen will in das Werk setzen und ausführen, damit unser Glaube gespeiset und ernähret werde, ihm zu vertrauen, und festiglich zu glauben, daß er es hinausführen kann [Matth. 4, 4. 5 Mos. 8, 3.]. Denn GOtt spricht alhier zu Mose, er solle ein Gott sein über Pharao, und Aaron solle Mosis Prophet sein. Der lateinische Text hat sich gescheut, also zu reden; aber es steht dürre allda. Und auf daß ihr euch auch nicht daran stoßet, wenn solche Art zu reden vorfällt, so wollen wir etwas davon handeln. Denn es werden noch solcher scharfen Stücke mehr kommen.

7. Eure Liebe hat oft gehört, daß unter allen Werken GOttes, so GOtt mit uns und allen Creaturen hat, das herrlichste und größte Werk sei sein Wort, davon alle Dinge geheiliget und nach GOtt genannt werden. Wie man denn gewöhnlich die Kirche ein Gotteshaus heißt, nicht, daß da GOtt wäre, sondern daß da GOttes Wort gehört und gepredigt wird. Und wenn es auch gleich unter einer grünen Linde oder Weide gepredigt würde, so hieße doch derselbige Ort GOttes Wohnung und Stätte; denn GOttes Wort regiert daselbst [1 Tim. 4, 5.]. Nicht, daß wir den Ort dazu geweiht haben, geziert oder schön ausgestrichen hätten, wie unsere Bischöfen ihren Güterchen thun, und mit ihrem Oele und anderer Zugehöre sie schmieren, sondern von dem göttlichen Worte, das ich nicht dahin bringe, sondern es kommt aus göttlicher Ordnung dahin. Wo nun das ist, daselbst ist auch sein Haus, und da will GOtt geehrt sein. Denn GOttes Wort macht nicht allein die Stätte heilig und zu GOttes Stätte und Hause; es macht auch eben dasselbige Wort aus den Leuten Götter.

8. Und dieses ist auch die Ursache (wie vormals [Cap. 3, § 25 ff.] auch angezeigt), daß der HErr zu Mose sagt, da er zu dem feurigen Busch wollte treten: „Zeuch deine Schuhe aus, denn dieser Ort ist heilig“, denn GOttes Wort ist da gewesen, GOtt hat daselbst geredet. Darum denn auch die da GOttes Wort haben, und es ihnen gegeben wird, daß sie GOttes Wort treiben und handeln, dieselbigen werden genannt Götter. Daher der HErr Christus selbst aus dem 82. Psalm, V. 6., einen Spruch, Joh. 10, 34–36., einführt, und schließt also wider die Juden: „Ihr sagt, ich lästere GOtt, daß ich gesagt habe, ich sei GOttes Sohn; nennt doch die Schrift Götter diejenigen, so da GOttes Wort hören, und zu denen GOttes Wort geredet wird“; wie viel mehr bin ich GOttes Sohn, der ich zu euch vom Vater gesandt bin? So muß nun der auch GOtt sein und GOtt angehören, der GOttes Wort redet.

9. GOtt ist unser Vater; darum, wie er Vater ist, also sind wir auch Söhne. Also gibt die heilige Schrift den Heiligen einen Namen, die GOttes Wort haben, daß sie auch GOttes Namen erben. Darum werden wir auch in der [ersten] Epistel Johannis [Cap. 3, 1. 2.] GOttes Söhne genannt; und gleichwie wir Erben aller Güter GOttes sind, also haben wir auch den göttlichen Namen, daß wir herrschen über
\quellseite{398}
Tod, Teufel und alles Unglück, und seiner Güter theilhaftig seien; aber durch den Glauben. Davon haben wir den tröstlichen, vortrefflichen Titel, daß wir Götter genannt werden.

10. Solche Kraft und Macht hat das göttliche Wort, und also hoch achtet die heilige Schrift die Schriften. Dies merket, daß es ein solch groß Ding sei, GOttes Wort rein und rechtschaffen haben, daß GOtt selbst einen nennt „Gott“; wie an diesem Orte Moses Pharaonis Gott sein sollte, denn er hatte GOttes Wort über Pharaonem. Also sind und sollen auch alle Christen genannt werden, die an Christum wahrhaftig glauben. Denn das ist GOttes Sohn sein, wenn wir GOttes Wort haben und dabei bleiben. Darum kann man die Schriften nicht genugsam preisen dieser Ehre halben, so wir haben, wir können sie auch nicht genugsam begreifen [Ps. 139, 6.]. Aber durch den Glauben wird es geschmeckt ein wenig [Ps. 34, 9.], und in dem andern Leben werden wir es sehen. Darum sollen wir das göttliche Wort lieben und hoch ehren, und dasselbige nicht verachten, als sei es ein Menschenwort, darauf nichts erfolgen würde [zech. 12, 22.]. Denn aus GOtt folgt alles Gutes. So sollen wir unsere Augen aufthun, und sehen, was wir doch empfahen und haben, wenn wir GOttes Wort haben.

11. Wenn bei uns GOttes Wort nicht ist, so sind wir im Reiche des Teufels, und sind junge Teufel und Teufels Kinder. Also sagt der HErr Christus [Matth. 16, 23.] auch zu Petro, da er widerret, daß er nicht in Jerusalem ziehen sollte: Hinderet mich nicht, du Teufel. Aber wer GOttes Wort hat, der ist ein junger Gott. Wer das Wort GOttes nicht hat, sondern hat des Teufels Sinn und Wort, das ist, falsche Lehre und Irrthum, das sind alles junge Teufel.

12. Also redet die Schrift gewaltiglich, preist, ehrt und rühmt das göttliche Wort; denn daraus rühmen wir uns, daß wir GOttes Söhne sind. Und drunten im 21. Capitel dieses Buchs, V. 6., heißt Moses die Richter „Götter“, das ist, die an GOttes Statt sitzen, Mosis Gesetz in der Hand haben, und richten nach dem göttlichen Worte, regieren nicht nach der Vernunft, sondern nach GOttes Wort; wie denn GOtt solch Volk darum geordnet hat.

13. Derhalben ist es nun eine ungewöhnliche Rede und Sprache, daß allhier gesagt wird: Ich will dich zum Gott setzen über andere Leute. Aber man muß sich an solche Weise der heiligen Schrift gewöhnen, auf daß wir lernen, wer Gott heiße. Es wird gesagt: Ueber Pharaonem will ich dich zum Gott setzen. Dies ist abermal ein Trost. Denn zuvor ließ GOtt Mose sinken, und Pharaonis Hohn und Spott sein, und mit den Kindern Israel alles ärger werden, und muß Moses gar bloß stehen und verzagt werden. Er hieß Moses sagen, daß das Volk Israel aus Egypten sollte ausziehen, und er, Mose, sollte es ausführen. Aber, je mehr Moses solches heißt, und sagt zum Ausführen, je bitterer und grimmiger Pharao, und je verzagter Moses wird. Aber da Moses also schwach wird, da kommt GOtt und tröstet ihn, und spricht: „Du sollst Pharaonis Gott sein“, und es mit ihm machen, wie du nur willst. Aber das soll geschehen durch mein Wort, das ich dir gebe. Durch dies Wort sollst du über ihn herrschen, er sperre und wehre sich nun so sehr, als er will, so sollst du es doch mit ihm machen, wie du willst.

14. Also geht es dem, der da GOttes Wort hat und GOttes Sohn ist, der hat Gewalt über alles. Denn ein Christ ist ein solcher gewaltiger Mann, daß ihm alle Creaturen müssen gehorsam sein; wiewohl solches nicht scheint, so ist es doch in der Wahrheit also. Was ist auf Erden kräftiger denn der Tod, erschrecklicher denn die Sünde, und bitterer denn das böse Gewissen? [Hohelied 8, 6.] Dennoch spricht ein Christ, daß er über diese alle ein Herr sei. Aber wir schlagen es in Wind, gaffen und sperren das Maul auf gegen den Vikaren der Apostel, so doch der HErr Christus [Joh. 14, 12.] sagt: Die Zeichen, die ich thue, wird ein Christ, der an mich glaubt, auch thun, und noch größere Werke, denn ich, thun. Als wollte er sagen: Es läßt sich für ein groß Werk ansehen, daß ich die Blinden sehend mache, die Tauben hörend, und die Stummen redend; aber er wird es auch thun, und viel mehr denn dieses thun.

15. Sage mir, ist es nicht ein groß Ding, daß ich armer Madensack dem Tode ein Klipplein schlagen, und dem Teufel Trotz bieten, und aller Fürsten Tyrannei in Wind schlagen soll, und über die Gewissen, und über alles herrschen? Wenn ich gleich einen krancken gesund mache, so ist es wohl ein Wunderzeichen, auch wenn ich einen sehend mache; aber er muß doch sterben, und dieses alles wieder lassen. Aber GOttes
\quellseite{399}
Wort zeigt, und spricht: Ich will dir eine solche Gewalt geben, die nicht sonst zu finden.

16. Darum, wenn ein Mensch mit Gottes Wort gesäuget ist, und hat GOttes Wort, so kann er als ein Christ wohl die Blinden sehend heben und die Lahmen gehend machen; denn GOttes Wort ist in ihm, und er ist GOttes Sohn, und wir thun durch GOttes Gaben und Gnade, was GOtt durch sich thut; das ist ein anderes. So ist es nun ein groß Ding, daß ein Christ den Tod und Teufel soll fressen; item, wenn ihm der Kopf abgerissen wird, daß er soll gedenken, der Tod sei anders nichts denn ein Schlag, und er werde seinen Kopf bald wieder kriegen.

17. Ueber dies alles, so werden wir in der heiligen Schrift [Ps. 68, 13. 30.] auch Könige genannt; nicht, daß wir Schwert und eine hohe Krone führen, die voller Smaragden und anderer Edelgesteine wäre, so groß als dieser Kirchspieler, sondern darum, daß wir trotzen können wider den Tod, Teufel und Hölle, daß sie uns nicht schaden müssen [Röm. 8, 1. 1 Cor. 15, 55.]. Diese Krone sollte ja einem gar wohl stehen.

18. So spricht nun allhier GOtt zu Mose: Gehe hin und verzage nicht, ich will mein Wort nicht lassen, das ich gesagt habe: Führe mein Volk aus, ich will dir beistehen. So sich es nun gleich ein wenig verzieht, so hast du doch GOttes Wort; du wirst sie schlagen, und mit ihnen handeln als irgend ein Gott.

19. Also ist Moses nicht nach seiner Person ein Herr über Pharaonem, über den Tod und den Teufel, sondern nach seinem Amte, wie er es denn auch mit seinen Zeichen beweist. Diesen Stab und Gewalt hat GOtt uns gegeben, daß wir Christen solche Herren sein sollen. Aber wir schlagen es in Wind und verachten es. Unser Thun und gute Werke, die bringen uns nicht dahin, so führt uns der Teufel doch immerdar in- das Thun hinein, so wir es doch nicht bekommen, denn allein durch GOttes Wort, dadurch werden wir Götter.

20. Moses ist durch sein Thun nicht zum Gott worden, ja, der HErr spricht zu ihm: O Moses, du bist gefallen und verzagt; aber ich will dich zum Gott machen. Es kommt frei aus GOttes Geheiß, und er will es ohne Werke haben, was vor GOtt soll bestehen. Werke gelten gegen den Leuten, und sollen den Leuten dienen, die sind damit zufrieden; aber GOtt ist nicht zufrieden damit, er will den Glauben und sein Wort haben. GOtt will etwas Höheres haben, und spricht: Predige mir mein Wort, das gilt so viel als ich; was soll ich mit deinen grindigen, krätzigen, schäbigen und stinkenden Werken machen?

\subsection{Und Aaron, dein Bruder, soll dein Prophet sein}

21. Ein Prophet wird genannt, der seinen Verstand von GOtt hat ohne Mittel, dem der Heilige Geist das Wort in [den] Mund legt. Denn er ist die Quelle, und sie haben keinen andern Meister denn GOtt [1 Cor. 14, 1. 2.]. Niemand kann einen Propheten machen durch menschliche Predigt und Lehre [2 Petr. 1, 21.], und ob es gleich GOttes Wort ist, und ich das Wort auf das allerreinste predige, so mache ich doch keinen Propheten; einen weisen und verständigen Mann kann ich machen. Als Matthäi am drei und zwanzigsten Capitel, V. 34., werden Diese genannt, welche ohne den Propheten die Lehre schöpfen, denn Gott durch Leute redet, und nicht ohne Mittel. Aber Propheten sind, die ohne alle Mittel die Lehre von GOtt haben. Also wird allhier gesagt: „Aaron soll dein Prophet sein.“ Gleich wie ich, GOtt, auch Propheten mache, ohne alle Mittel, wenn ich mit ihnen rede; also bist du Gott, und Aaron ist Mosis Prophet, denn er soll ohne Mittel von dir lernen, wie du von mir gelernt hast.

22. Wohlan, so ist nun Moses GOttes Prophet, und Aaron ist Mosis Prophet. Ich wollte aber lieber GOttes Prophet sein, denn Mosis; denn wo GOtt nicht ist, da geht es übel zu. Und dieses Prangen, oder diese Gewalt und Macht haben wir aus GOttes Wort wider die guten Werke, daß man sie verdamme und verwerfe, wenn man damit vor GOtt handeln will. Die Leute können der guten Werke nicht entbehren, oder müßten gar verderben. Denn GOtt hat es also geordnet, daß dieser krank ist, und daß oft ein Auge am menschlichen Leibe, oder die Hand, oder der Mund gebrechlich, und bedarf Hülfe und Rath. So plumpet denn die tolle Hure, die Vernunft, her, und will es hinaustragen vor Gott, die Hände und Zunge abschneiden, Augen ausstechen, und damit etwas verdienen; aber es gilt nichts vor GOtt, denn allein der Glaube.

23. Das ist die Prophezei, daß GOtt sagt: „Ich will dir anzeigen, was du Pharao sagen sollst. Du sollst sein Gott sein.“ GOtt ist
\quellseite{400}
noch über ihn, er redet auch durch uns, wir sehen ihn aber nicht. Daher sagt St. Paulus [2 Cor. 5, 20.], wir führen Christi Botschaft, wir sind seine Mundboten, die wir der Welt antragen, daß sie mit GOtt sich versöhnen solle; das ist unsere Werbung,\footnote{Werbung = Botschaft.} daß wir sagen: „Lasset euch versöhnen mit GOtt.“ GOtt, der im Himmel sitzt, hat es uns befohlen, daß wir das Wort also führen, gleichwie Aaron Mosis Wort führt.

\begin{schriftzitat}
V. 3. Aber ich will Pharao Herz verhärten.
\end{schriftzitat}

24. Mosi ist der Befehl oder die Instruction gegeben, nämlich: „Aaron soll dein Mund sein, und du sollst sein Gott sein.“ Was folgt aber endlich daraus? Wie stellt und hält sich Pharao gegen dieser göttlichen Botschaft und Werbung? Anders nicht, denn wie auch allhier gesagt wird, er will weder sehen noch hören, weder weichen noch\footnote{Siebensche: oder.} gehorchen. Darum so ergeht aus gerechtem Gericht, daß GOtt verhängt, daß das Herz Pharaonis verstockt wird und bleibt. Er will Mosen nicht hören, noch diesem Antragen und trefflichen Wunderzeichen stattgeben.

25. Alhier ist ein tiefes Loch, da man sollte viel von predigen; aber ich thue es ungerne, um der schändlichen, vorwitzigen Natur willen, die allezeit Lust hat, etwas Seltsames und Neues zu wissen, das ich denn gerne vermeide, wo ich es kann umgehen. Diesem Laster ist GOtt sehr feind, und er kann es übel leiden, wenn man rühmen, trotzen und gesehen sein will, außer der heiligen Schrift, wenn wir die Gedanken schöpfen, daß wir die heilige Schrift gar wohl wissen, und den Menschen die Ehre geben, und nicht GOtt; gehen in den Gedanken, daß wir GOttes Wort haben, GOttes Kinder und Güter sind. Aber wenn wir daraus ein Stolzieren machen, so läßt GOtt den Teufel eintreten, sein Wort wegnehmen, und aus dem Gott einen Teufel werden.

26. Ich habe ihrer viel gesehen, und die noch heutiges Tages leben, wenn sie etwas sehen, das wir zuvor nicht gelehrt haben, das schmeckt ihnen als Zucker, und gehen vor: O das ist zuvor nicht gelehrt worden. Diesem ist unser HErr GOtt auch feind, und ein solcher fällt aus einem Gedanken in andere Gedanken, bis ihm der Kopf abgerissen wird, und zu Boden stürzt und gar untergeht. Man gafft darnach, wie man immerdar etwas Seltsames und Neues wissen möge; denn das Alte ist das Manna, oder das tägliche Brod. Aber ist es nicht seltsam und unerhört, daß wir haben, jedoch ist\footnote{„ist“ von uns herübergenommen aus der alten Ausgabe Walchs. Der Sinn ist: Wenn gleich das, was wir lehren, nicht etwas Seltsames ist, so ist es doch eine gesunde Nahrung, deren man überdrüssig ist.} es Fleisch, deß man müde ist zu essen. Dann wird auch GOtt nur Fragen geben, wie er in den hohen Schulen gegeben hat, da die Welt mit Büchern ist überschüttet worden, und die Gewissen sich darinnen wohl gemartert und geplagt haben. Diese Zeit wird wieder kommen. Das ist die Ursache, warum ich solches nicht gerne handele; denn der vorwitzige Teufel und unser Fleisch hört nicht auf in diesem Stücke.

27. Ich lasse mit Allegorieen spielen und handeln, wer da will; allein schauet ihr darauf, daß ihr das Hauptstück nicht verlaßt, sondern treiben und üben möget. Denn wir haben den Teufel zum Feinde, welcher nicht ruht, er reißt einem das Hauptstück bald hinweg, und gibt einem seltsame Fragen ein, daß man grübele und wühle. Darum so schneidet es also, und ziehet die Bibel herzu zu Christo, und folget nicht den Fragen. Denn die Bibel und GOtt ziehen sich fein zu Christo, wie er uns allein selig mache. Andere haben seltsame Gedanken, und führen sich von Christo, wollen etwas Neues haben. Aber die heilige Schrift will von nichts anders wissen, noch uns vorlegen, denn Christum. Und wer die Schrift also führt, oder durch die Schrift also zu Christo geführt wird, der bleibt wohl und geht auf richtiger Bahn.

\begin{schriftzitat}
V. 3–5. Aber ich will Pharao Herz verhärten, daß ich meiner Zeichen und Wunder viel thue in Egyptenland. Und Pharao wird euch nicht hören. Auf daß ich meine Hand in Egypten beweise, und führe mein Heer, mein Volk, die Kinder Israel, aus Egyptenland, durch große Gerichte. Und die Egypter sollen es inne werden, daß ich der HErr bin, wenn ich nun meine Hand ausstrecke über Egypten, und die Kinder Israel von ihnen wegführen werde.
\end{schriftzitat}

28. Man pflegt hier zu fragen: Wie es komme, daß GOtt das Herz Pharaonis will verstocken, warum hat er denn die Zeichen kommen lassen? Moses wird hernach immer also schreiben: Pharao habe ihn nicht gehöret etc., wie denn
\quellseite{401}
Gott also geredet hatte. Und auf diesen Spruch zeucht sich's hernach, wie ihm GOtt allhier zuvor sagt. Es ist gleich also ergangen, wie er gesagt und verheißen hat, spricht Moses. Aber diesen Punkt hat St. Paulus zum Römern am 9., V. 17. 18., und 11. Capitel, V. 7–10., und das folgende Capitel im Exodo ausgelegt, und die Vernunft verwundert sich sehr darüber, und will drinnen grübeln, und GOtt suchen, damit muß sie zu schaffen haben. Da kommt denn eine Frage um die andere, und wenn sie hinein kommt, gedenkt sie: Verspottet Gott, weiß ist die Schuld anders, denn sein? Aber das sind leichtfertige Geister.

29. Also soll man anfangen, daß man spreche: Ich will vom Untersten anfangen, und GOttes Gnade und sein Wort erkennen. Und wenn wir denn in GOttes Wort kommen, so soll uns dasselbige tröstlich und süß sein. Wie denn dies allhier dem Mosi der höchste Trost und Trotz gewesen ist; und er kann keinen größeren Muth fassen und schöpfen, denn daß Gott Pharaonem verspottet. Ein anderer hätte den Trost nicht geschöpft. Moses stand in dem Falle, daß er GOttes Wort führte, das kostet ihm Leib und Leben, und daß er das Volk ausführte, das kostet ihm seine Ehre. Nun ist es ein groß Werk, daß ein Mensch in seinen Sinn nicht nehmen sollte. Aber dieser Rückhalt thut das Beste, es wird's GOtt hinaus führen.

30. Gleichwie wir auch viel toller Fürsten jetzt wider das Evangelium haben, die da schnurren und murren; wie allhier Pharao auch gethan hat. So sind wir dagegen schwach, und unser Trost, Muth und Pochen ist das, daß es gehen wird, wie GOtt es wird haben wollen; ihre und unsere Hälse sind in seinen Händen, und wir haben den Vortheil, daß es GOttes Wort ist, so wir haben; alsdann, wenn sie schnurren und poltern, so stehen wir desto fester.

31. Wie denn auch St. Paulus zu Timotheo sagt, er habe keinen guten Tag gehabt, spricht [2. Ep. 2, 9. 10.]: Ich muß leiden, und alles Leid tragen, um der Gläubigen willen, auf daß sie den Glauben erlangen. Und diesen hochmüthigen Titel, daß er leide um der Auserwählten Glauben willen, führt auch St. Paulus zum Colossern [Cap. 1, 24.], welches er vor seinem Ende geschrieben hat. Er hat manchen Tück und Buben erkennen müssen. Dieser hat ihm gebräuet, und ist sein Widersacher gewesen; jener hat seine Predigt angenommen, ist doch ein falscher Bruder geblieben; daß er [Gal. 2, 4. f.] spricht, er habe keinem vertrauen dürfen, aber nun freue und rühme er sich über dem Glauben der Auserwählten. Als sollte er sagen: Lasset uns solches trösten, ob gleich etliche fallen, und Büberei anrichten [Matth. 12, 14. Luc. 6, 11. Phil. 3, 18. 19.], so werden doch auch etliche beständig bleiben. Denn die den Glauben haben sollen, die haben ihn, wie er zum Timotheo sagt. Demas und die Galater fielen dahin [2 Tim. 4, 10.], also ärgerte man sich, und es machte St. Paulo bange, daß er lieber gestorben wäre. Dergleichen er denn auch [2 Tim. 3, 1–8.] darüber klagt, daß in den letzten Tagen greuliche Zerrüttungen eintreten werden, sündliche Menschen sein werden, die von sich selbst viel halten, geizig, stolz, hoffärtig, Lästerer, Schänder; die da eine Geberde eines göttseligen Wandels haben, aber seine Kraft verleugnen; sie lernen immerdar, und kommen nimmermehr zu der Erkenntniß der Wahrheit. Und nennt mit Namen Moses Widersacher, spricht: „Gleicher Weise, wie Jannes und Jambres Mosi widerstunden, also widerstehen auch diese der Wahrheit.“

32. Aber laß es gehen, GOttes Grund steht fest, und er hat ein Siegel, das heißt also: „GOtt kennet die Seinen“ [2 Tim. 2, 19.]. Nun wird es dennoch gehen, wie GOtt will. Es werden neue Propheten kommen, ich und andere werden dann gestorben sein, so wird sich der Jammer machen. Hin ist hin; laß laufen, was da läuft. GOtt treibt also sein Spiel; noch bleibt, was bleiben soll. Ein Mensch könnte wider diesen Haufen keinen Trost geben; aber GOtt will allhier sagen: Mein Moses, sei du zufrieden, es wird doch gehen, wie ich es will haben, halte fest, ich will verstocken Pharaonis Herz; du weißt es doch, daß ich es thun will, du hast GOttes Wort, item, die Kinder von Israel auch; was nicht gehen will, das bleibe dahinten.

33. Dieses gehört nicht für die müßigen Geister, die mit Fragen drein fallen. GOtt hat diese Art zu reden droben, mit Adam, Eva und Abraham nicht geführt; darum soll man es ansehen, wo es die Schrift also braucht. In Mose werden nicht stolzere noch hochmüthigere Reden geführt, denn allhier. Denn da das Volk Israel in den höchsten Nöthen ist, und Pharao mit
\quellseite{402}
grausamer Gewalt und Tyrannei fährt, so soll Mosis keinen andern Trost haben, denn daß er es lasse gehen, wie es GOtt haben will. Dies Liedlein müssen wir auch noch singen: Mitte vadere, sicut vadit, quia vult vadere, ut vadit. Die zornigen Fürsten, die falschen Propheten und Rottengeister sind verstockt, GOtt hat es also haben wollen. Warum thut denn GOtt also? So wir in seinen Willen unsern Willen wollen stellen, so ist kein anderer GOtt, und hat auch keinen andern Willen; er hat auch über ihm kein ander Reich, sein Wille ist der oberste, und wenn derselbige Wille geht, das will ich auch, so soll man es lassen gehen.

34. Nun ist eine andere Frage: Warum GOtt Mosen hat heißen predigen, da doch GOtt selbst sagt: Pharao wird euch nicht hören? Ist es nicht ein närrisch Ding, wenn einer spräche zu einem: Lieber, predige Pharaoni, und wisse, er wird dich nicht hören, ich will ihn verstocken? Da wollte ich einem die Schlüssel vor die Füße werfen, und sagen: Predige du selber. Aber das ist die Antwort: Uns ist befohlen zu predigen, aber nicht befohlen, die Leute gerechtfertigt und fromm zu machen. Da sollen alle Prediger und Christen getrost sein, und ein jeglicher seinem Berufe folgen, und denselben getreulich ausrichten. Mosi wird allein das Wort GOttes befohlen, und wird ihm nicht auferlegt, Pharaonem weich oder hart zu predigen. Das Wort ist ihm befohlen, das ist GOttes Wille, und das soll er führen, ob ihn auch gleich niemand hören wollte. Es geschieht ihm zu Troste, auf daß er nicht erschrecke, wenn ihm niemand folgen und gehorsam sein will. Wenn ich mich des sollte annehmen, daß mein Wort und Predigt verachtet wird, so sollte ich wohl aufhören zu predigen. Aber fahre du fort, Mose, predige du; wirst du darüber verachtet, das laß mir befohlen sein. Man soll allein darauf sehen, daß es GOtt also haben will, und will also predigen lassen.

35. Das sage ich darum, auf daß man sich vorsehe, und nicht zu hoch in die Fragen fahre, sondern herunter sich lasse, und Christum erkennen lerne, und ein jeder auf seinen Beruf und auf sein Ding sehe, was ihm GOtt befohlen habe, und bekümmere sich nicht mit den Fragen. Es ist dir nicht befohlen, solches zu forschen; du hast die Schuhe noch nicht ausgezogen, du kommst zu frühe, du lässest das anstehen, was du wissen solltest, und achtest das nicht, was du wissen und fragen solltest. Mosis fragt allhier GOtt nicht, und bekümmert sich auch nicht, warum GOtt den Pharaonem verstocke; er disputirt allhier deßhalben mit GOtt nicht.

36. Aber gleichwohl, daß ihn GOtt verstocket will, geschieht nicht ohne Ursache. Denn GOtt will viel Wunderzeichen thun, da Mosis die Kinder Israel soll ausführen. Im neunten Capitel [V. 16.] wird gesagt: Ich will dich verstocken, auf daß ich an dir beweise meine Kraft, und meine Ehre und Preis der Welt bekannt werde. Da dienen sie zu, daß der Tod, der Teufel und die zornigen Fürsten GOttes Ehre hoch machen. Wenn die tollen Fürsten nicht also tobeten, und der Tod und die Sünden das Gewissen das Gewissen nicht also beschwerten, sondern friedlich einher gingen, so würde die Kraft des Glaubens und göttlichen Worts von niemand erkannt. Aber dieweil der Teufel und die Fürsten toben, so scheint und leuchtet der Glaube, und das göttliche Wort tröstet mich und macht mich muthig, und das Gewissen erfährt es, daß ein Nachdruck und Gewalt bei dem göttlichen Worte ist. Daß ein arm Herz, Fleisch und Blut sollte den Muth und die Freudigkeit haben, die rasenden Fürsten, den Teufel und Tod zu trotzen, und verachten; wie an den heiligen Märtyrern, auch an den Jungfräulein, St. Agnes und Agatha, zu sehen: solches geschähe nimmermehr, wenn sich die Welt und der Teufel wider die Christen nicht also gewaltig legten.

37. „Verstockt“ heißt nicht allein, wie wir sagen, verzweifeln, sondern, so Fürsten sind, die sich verlassen auf ihre Gewalt, auf Land und Leute, und einen Muth kriegen, als wollten sie GOtt fressen. Dieser Muth, diese Sicherheit und Kühnheit heißt verhärten oder verstocken. GOtt läßt ein armes Häuflein Christen einhergehen, welches sie gedenken bald unterzudrücken. Denn was ist Mosis und Aaron gegen uns zu rechnen? wir wollen es je wohl in einem Brei fressen, meinen die Egypter, es ist ein geringer Haufe. Also brüsten sie sich, als sei alle Gewalt ihr [Ps. 73, 6–9.], der Muth wächst ihnen, daß sie so aufgeblasen sind, wie die Hopfensäcke. Dieser große Muth heißt verstockt sein. Wiewohl solches nicht die Güter und Reichthümer allein machen, sondern auch inwendig der Muth, daß sie gedenken: Hui, wir wollen sie fressen, oder gar
\quellseite{403}
versenkten. Sie werden einen großen Anhang haben, und nichts wird zwischen euch und den Widersachern liegen, denn GOttes Wort. Auf eurer Seite wird Armuth, auf ihrer Reichthum sein; auf eurer Seite Schwachheit, auf ihrer Gewalt und Stärke; und sollen auf euch fallen mit aller Pracht, und dennoch fehlen. Damit will ich anzeigen die Kraft meines göttlichen Worts mitten in der Verfolgung, mitten im Tode, in der Sünde und in des Teufels Zorn.

\begin{schriftzitat}
V. 10–12. Da gingen Mose und Aaron hinein zu Pharao, und thaten, wie ihnen der HErr geboten hatte. Und Aaron warf seinen Stab vor Pharao, und vor seine Knechte, und er ward zur Schlange. Da forderte Pharao die Weisen und Schwarzkunstiger. Und die egyptischen Zauberer thaten auch also mit ihren Beschwörern, und warf ein jeder seinen Stab von sich, und wurden Schlangen daraus.
\end{schriftzitat}

38. Alhier sehen wir, daß der Teufel durch seine Diener und Werkzeuge, oder falsche Propheten, auch könne Mirakel thun. Denn die Zauberer auch machten aus ihren Stäben Schlangen; item, daß die Wasserströme in Egypten Blut werden; daß sie Frösche aus dem Wasser bringen. Solches erinnert und warnt auch die heilige Schrift. Als, 5 Mos. 13, 1. ff. wird gemeldet, daß auch falsche Propheten können Wunder thun. Und der HErr Christus Matth. 24, 24. spricht: „Es werden aufstehen viel falsche Propheten, und große Zeichen und Wunder thun.“ Item, Paulus 2 Thess. 2, 9. 10. sagt: Die Zukunft des Antichrist geschieht nach der Wirkung des Satans, mit allerlei lügenhaftigen Kräften, Zeichen und Wundern, und mit allerlei Verführung zur Ungerechtigkeit.

39. Dieses sind Teufels Wunder; denn die Zauberer thun solche Mirakel aus GOttes Verhängniß durch des Teufels Kraft und Wirkung, damit sie Moses Mirakel, die da göttliche Mirakel waren, hindern, und Pharao verbleiben und verstocken, daß er die göttliche Wahrheit nicht annehme, sondern im Irrthum stecken bleibe, und wider GOtt und sein Volk sich legte mit seinem Toben und Wüthen. Also thun auch Wunder die beiden Zauberer in den Geschichten der Apostel, Elymas und Simon (Apost. 8, 9–11. 13, 8.). Item (Cap. 16, 16.), das Mägdelein, so einen Wahrsagergeist hat. Aber sie thun es aus des Teufels Kraft, durch GOttes Verhängniß. Also 1 Sam. 28, 11–14. bringt die Wahrsagerin aus dem Grabe ein Gespenst, das dem Propheten Samuel an Gestalt und Kleidern ähnlich war, durch des Teufels Blendung und Betrug. Daher St. Paulus [2 Thess. 2, 10.] des Teufels Zeichen und Wunder auch lügenhaftige Kräfte und Verführung nennt. Denn der König Saul ward durch des Teufels Betrug recht geäfft, daß er nicht anders wußte, es wäre der verstorbene Samuel, da es doch nicht Samuel war, der im HErrn ruhte, sondern ein Gespenst des Teufels, dem Saul zur Strafe von GOtt verhängt, dieweil er sonst Samuels Lehre in seinem Leben verachtet und in Wind geschlagen hatte.

\subsection{Und Aarons Stab verschlang ihre Stäbe}

40. Der Teufel wird mit seinen Wundern von GOttes Mirakeln zu Schanden gemacht. Denn wenn GOttes Finger oder Mirakel und des Teufels Betrügerei zusammenstoßen, so wird der Teufel mit seinen Wundern von GOttes Wundern eingetrieben. Wie denn der HErr Christus im Evangelio den Teufeln gebeut, daß sie schweigen müssen, er macht sie zu Schanden, treibt sie aus, überwindet und bindet sie, ja, er zerbricht die Werke des Teufels; der Teufel verkriecht sich mit seinen Mirakeln vor GOttes Wort.

\begin{schriftzitat}
V. 19–21. Und der HErr sprach zu Mose:
\end{schriftzitat}

\begin{quote}
Sage Aaron: Nimm deinen Stab, und recke deine Hand aus über die Wasser in Egypten, über ihre Bäche, und Ströme, und Seen, und über alle Wasserstrümpfe, daß sie Blut werden, und sei Blut in ganz Egyptenland, beide in hölzernen und steinernen Gefäßen. Mose und Aaron thaten, wie ihnen der HErr geboten hatte, und hub den Stab auf, und schlug in das Wasser, das im Strome war, vor Pharao und seinen Knechten, und alles Wasser im Strome ward in Blut verwandelt. Und die Fische im Strome starben, und der Strom ward stinkend, daß die Egypter nicht trinken konnten des Wassers aus dem Strome, und ward Blut in ganz Egyptenland etc.
\end{quote}

41. So ist nun das erste Zeichen gewesen, daß GOtt Egyptenland voller Bluts machte, daß alle Wasser, so da standen, als die Teiche der Wasserströme waren, voller Blut floßen. Dieses ist ein greulich Zeichen gewesen, dadurch GOtt sie zu seinem Gehorsam hat treiben wollen, sie zur Buße reizen und locken. Wie denn GOtt
\quellseite{404}
seine Plagen und Strafen darum über [die] Gottlosen kommen läßt, nicht, daß er sie damit gar vertilgen und verderben wollte, sondern ihnen helfen möchte; gleichwie ein Vater sein Kind stäupt, nicht, daß er als ein Henker oder Stockmeister es totschlage, sondern durch den Schilling\footnote{Schilling = Ruthe, Strafe.} seiner Bosheit, Unart und Schalkheit steuere, und das Kind fromm werde. Denn die Ruthe nimmt die Thorheit von des Kindes Herzen [Sprichw. 22, 15.]. Dennoch wird es in Sünd geschlagen; denn Pharao sahe, daß seine Zauberer solche Zeichen nachthaten.

42. Diese Plage ist sieben Tage gestanden, auf daß man nicht meinte, es wäre ein zufälliger Handel, sondern man es annehmen müßte als ein Mirakel. Und Pharao mußte andere Wasser suchen; denn sie neben den Strömen Wasser suchten, da sonst zuvor keine Borne waren. Ist das nicht eine Härtigkeit, daß Menschen also hart und verstockt sein sollen, daß ein solch Zeichen sichtlich und greiflich vorgelegt wird, und das menschliche Herz soll noch sagen: Nein, es ist nichts? Aber es geht also, daß GOtt mit seinen Plagen und Schlägen bei den Gottlosen wenig ausrichtet; sie werden je länger je ärger, und sind wie die stätigen Pferde, je schwerer man auf dieselbigen schlägt, je weniger man sie kann fortbringen, sondern weichen nur hinter sich. Wie denn Jesaias auch davon sagt, und spricht [Cap. 1, 5. Jer. 5, 3.] (welches wir uns zur Warnung mögen lassen gesagt sein): Dies Volk hat sich nicht bekehret zu dem, der es schlug.

43. Alhier hat Moses und jedermann lernen müssen, daß es nichts helfe, wie sehr man predige, oder wie herrlich man GOttes Wort habe, wenn GOtt nicht mitwirkt, und das Herz umkehre, daß es lasse das Wort eingehen und fasse es. Denn er muß das Gedeihen dazu geben zu dem, was gepflanzt, gesetzet und begossen ist [1 Cor. 3, 6. 7.]. GOtt ist der, so da verhärtet, und wenn er verhärtet, so kann es niemand weich machen. Wiederum, wenn er etwas weich macht, so kann es niemand hart machen etc.

44. Der Himmel ist zu hoch und die Hölle zu tief. Wenn GOtt mit seinem Worte leuchtet, so geht es gar hell auf am Himmel, Sonne, Mond, Sterne etc., und leuchtet über alle Maß, über alle Gesetzlehre, über alle guten Werke, und bringt uns ein solch Licht, daß wir nicht im Finsterniß bleiben mögen; wie wir denn jetzt auch haben. Wiederum, verkennet Gott sein Wort, so stecket er uns so tief in die Hölle, daß es kein Ende hat. Wie wir denn im Papstthume auch zu Narren worden sind, da wir närrischen Werken angehangen; was wir angerührt haben, das ist Dunkel und Finsterniß gewesen, und sind darinnen immer untergesunken. Also ist es entweder zu hoch erleuchtet, oder zu tief versenkt; darum hat es keine Maß. Damit er anzeigt, es könne sonst nicht sein, daß solche Wunderzeichen das menschliche Herz nicht rühren sollten, aber GOttes Werk und Wunder sei es, der eines Herz also verhärtet und verstockt.

45. So haben wir nun das erste Zeichen, daß alle Wasser in Egypten in Blut verwandelt sind, und ist dennoch keiner, der sich daran kehrt.

\section{Das achte Capitel}

\begin{schriftzitat}
V. 1–7. Der HErr sprach zu Mose: Gehe hinein zu Pharao, und sprich zu ihm: So sagt der HErr: Laß mein Volk, daß mir’s diene. Wo du dich weigerst, siehe, so will ich alle deine Grenze mit Fröschen plagen, daß der Strom soll von Fröschen wimmeln. Die sollen herauf kriechen, und kommen in dein Haus, in deine Kammer, auf dein Lager, auf dein Bette, auch in die Häuser deiner Knechte, unter dein Volk, in deine Backöfen, und in deine Teiche. Und sollen die Frösche auf dich, und auf dein Volk, und auf alle deine Knechte kriechen. Und der HErr sprach zu Mose: Sage Aaron: Recke deine Hand aus mit deinem Stabe über die Bäche und Ströme, und Seen, und laß Frösche über Egyptenland kommen. Und Aaron reckete seine Hand über die Wasser in Egypten, und kamen Frösche herauf, daß Egyptenland bedecket ward. Da thaten die Zauberer auch also mit ihrem Beschwören, und ließen Frösche über Egyptenland kommen.
\end{schriftzitat}

\quellseite{405}
1. Wir haben allhier das andere Wunderzeichen, nämlich, daß alle Ströme und Seen in Egypten voller Frösche werden. Denn nun folgen drei oder vier Plagen nach einander, dadurch GOtt die Egypter zur Buße und Besserung ihres sündlichen Lebens, auch zu seiner Furcht hat reizen wollen, denn dieses alles Strafen der Sünden haben sein sollen. Erstlich haben wir gehabt Blut, darauf folgen die Frösche, hernach Läuse, zum vierten böse Würmer. Und in diesem Capitel sind drei Plagen beschrieben, aber es wird allezeit daneben angezeigt, warum diese Plagen nicht kräftig genug gewesen sind.

2. Moses führt GOttes Wort und die Wunderzeichen, so auf das Wort folgen. Das sollte einen ja bewegen. Wenn ich einen Prediger hätte, der da GOttes Wort handelte, und daneben Zeichen von sich gäbe, daß es die Leute sähen, und wollten sich daran nicht kehren, wäre das nicht böse? Dieses geschieht allhier beides: Das Wort und die Wunderzeichen sind da; und dennoch schafft Moses nichts, er muß Patienz haben, und beides geht zurück, er muß nicht gehört, auch noch dazu verachtet werden.

3. GOtt tröstet aber Mosen in dem, daß er spricht: Kehre dich nicht daran, ich bin es, der ich Pharaonem verhärte. Wenn GOtt mir also sein Wort befehlen würde, und Kraft gäbe Wunderzeichen zu thun, und dennoch spräche, ich sollte damit nichts ausrichten: sollte mich es nicht lüsten, ein Prediger zu sein? Und dennoch muß [es] also ergehen. Und wir wollen uns noch darüber verwundern, wie es komme, daß die Leute wider das Evangelium so handeln und toben! Es ist ein groß Zeichen und Wunderwerk GOttes, daß etliche noch das Evangelium annehmen. Aber daß die Welt voller Secten und Rotengeister wird, das ist sich nicht zu verwundern; es sollte die Welt gar verstockt sein. Denn sie ist Fleisch und Blut. Wem das Evangelium zu Herzen geht, der danke GOtt, denn es ist eine große Gnade, wenn GOttes Wort uns in das Herz geht.

\begin{schriftzitat}
V. 8–14.
\end{schriftzitat}

\begin{quote}
Da forderte Pharao Mosen und Aaron, und sprach: Bittet den HErrn für mich, daß er die Frösche von mir und von meinem Volke nehme, so will ich das Volk lassen, daß es dem HErrn opfere. Mose sprach: Habe du die Ehre, und stimme mir, wann ich für dich, für deine Knechte und für dein Volk bitten soll, daß die Frösche von dir und von deinem Hause vertrieben werden, und allein im Strome bleiben. Er sprach: Morgen. Er sprach: Wie du gesagt hast, auf daß du erfahrest, daß niemand ist, wie der HErr, unser GOtt, so sollen die Frösche von dir, von deinem Hause, von deinen Knechten, und von deinem Volke genommen werden, und allein im Strome bleiben. Also gingen Mose und Aaron von Pharao, und Mose schrie zu dem HErrn der Frösche halben, wie er Pharao hatte zugesagt. Und der HErr that, wie Moses gesagt hatte, und die Frösche starben in den Häusern, in den Höfen, und auf dem Felde. Und sie häuften sie zusammen, hie einen Haufen, und da einen Haufen, und das Land stank davon.
\end{quote}

4. Pharao ist nicht ein schlechter Mann gewesen, er hat seine Fürsten, Räthe und Junker gehabt, die Hof viel Verdienst gethan haben; item seine Zauberer oder Philosophen und weise Leute. So ist Egypten auch gar ein schön, herrliches Königreich gewesen. Moses ist dagegen ein armer Bettler, und aus der Wüste daher gekommen. Und was sind Moses und Aaron gegen den Weisen und Gewaltigen in Egypten? Noch demüthigt sich Pharao gegen sie; welches zu verwundern; dazu demüthigt er sich also sehr, daß er will, man soll GOtt für ihn bitten, stellt sich gar fromm.

5. Moses nimmt es an, gleich als wäre Pharao bekehrt. Denn er dicht und stellt sich, gleich als fürchte er den GOtt, den Moses predigt; er erzeigt sich, als erkenne er die Wunderzeichen, und nehme das Wort GOttes an. Und dennoch ist er ein böser Bube, hat im Sinne, daß er gerne die Kinder Israel alle erwürgen möchte, ob er wohl freundlich mit Mose und Aaron redet, sucht ihre Fürbitte, als fürchte er sich vor GOtt, und als wollte er fromm werden. Aber es ist ihm kein Ernst, er ist ihnen toffeind: könnte er Mosen und Aaron in einem Löffel ersäufen, so nähme er keinen Zober dazu. Davon sagt auch der acht und zwanzigste Psalm, V. 3.: „Die Gottlosen und Uebelthäter reden freundlich mit ihrem Nächsten, und haben Böses im Herzen.“ Das ist die Art und Natur aller Heuchler, welche wir müssen kennen lernen. Aber unser HErr GOtt kartet es auch wunderlich; welches denn die Welt irret. Es ist ihm stracks darum zu thun, daß er dadurch uns an sein Wort binde.

6. Ich habe droben oft gesagt, und sage es
\quellseite{406}
noch, daß man muß von einander scheiden das Leben der Heiligen, und das Wort GOttes, so sie führen; wie man sonst den Himmel von der Erde scheidet. Man kann nicht genug davon predigen. Ich rede jetzt von der heiligen Leute, als St. Petri und Mariä, guten Werken, nicht von des HErrn Christi und der Engel Werk. Denn der Teufel kann anders nichts, denn gute Werke lehren. Und wenn wir gestorben sind, so werdet ihr es sehen, wie er wider das Evangelium toben und wüthen wird. Welchen Griff weder ihr, noch die Prediger sehen werden, ob sie wohl weise sind. Allezeit bringt er Werke herein. Da scheidet denn die Werke oder das Leben von dem Worte, auf daß ihr nicht von dem Worte GOttes zu den guten Werken verführt werdet, denn da ist einer verloren. Darum sprich du: Ich will nicht Werke haben, es wären denn solche Werke, bei welchen GOttes Wort steht. Also sollst du ihm thun. Denn wir sind an die Worte, so da GOtt redet, gebunden, die sollen uns zu GOtt führen. Werk und Beruf müssen sein, aber das Wort GOttes soll uns gewisser sein.

7. Christus hat viel Werke gethan, so ich nicht thue. Er ging auf dem Meer [Matth. 14, 25.], und machte die Blinden sehend [Matth. 9, 29. Joh. 9, 6. 7.]; aber er hat mich es nicht geheißen, daß ich dasselbige auch nachthun sollte. Also hat er auch viel gelassen, so er mich nicht heißt lassen, als, er nahm kein Weib. Item, Johannes der Täufer aß nicht Fleisch, sondern Honigseim [Marc. 1, 6.]. Christus pflügte nicht, noch hat jemals einiges Handwerk getrieben. Wenn einer da sagete: Was er gelassen, sollen wir auch lassen; wie will das geschehen? Ich wollte nicht einen Heller darauf geben, auch auf die Werke Christi. Folgen will ich nicht, es sei denn, daß das Wort da steht. Denn ich will mit dem 119. Psalm, V. 105., singen: Eloquia tua lucerna mea [„dein Wort ist meines Fußes Leuchte“]. Unserm Herrn GOtt gebühret zu reden und zu lehren, mir aber gebühret zu thun; das ist, gute Werke soll ich thun und halten, aber er ist der Meister, er soll durch mich thun.

8. Wieso hat allhier Werke, die gar gleich sind, wie andere thun. Wenn ich einen Todten auferweckte zur Bestätigung des göttlichen Worts, und ein anderer käme, und thäte es auch, was wäre es? Gar zu Boden geschlagen. Wir wollen weder Zeichen noch etwas Anderes thun.

Aber allhier hat er es alles beides, und thut es; und dennoch wird es danieder geschlagen. Denn sie haben gesagt: Da kommt der, und hat das Wort und will Zeichen thun; haben es doch die Unsern auch, und können es auch thun. Damit ist es ergangen, daß Pharao gar verhärtet worden ist, wiewohl er sich fromm stellt; aber es gehört dahin, daß keiner dem andern trauen soll, ob er fromm sei, oder fromm scheine.

9. Denn die Bösen werden doch zu Schanden gemacht, und die verzweifelten bösen Buben meinen oft, und wollen fromm sein, aber es ist nur ein Schein und kein Ernst; wie denn Pharao oft sich also stellt, aber das Herz ist weit davon. Denn wenn die Strafe gegenwärtig ist, so ist er fromm; wenns überhin ist, so geht er wie zuvor; er bessert sich nicht aus Grund seines Herzens, sondern es ist eine gezwungene Andacht; gleichwie wir zur Zeit der Pestilenz und Theuerung fromm werden. Aber GOtt will das Herz haben.

10. Ich vermahne euch, daß ihr GOttes Wort rein behaltet, nämlich also, daß ich euch hinauf führe ohne Werk. Ich weiß wohl, was da wird vonnöthen sein. Daß ihr es sollet für Gott anrichten, [dazu] läßt er es euch nicht in die Faust bringen und predigen: Dies und das sollt du thun.\footnote{Hier ist die Relation, wie es uns scheint, völlig sinnlos, deshalb haben wir geändert. Die Eislebensche bietet: „Ich weiß wohl, was da will vonnöthen sein, daß ihr wollet für Gott anrichten. Läßte (sic) er's euch nicht in die Faust bringen und predigen: Dies und das sollt du thun? Denn“ etc. Der Sinn ist: GOtt läßt euch nicht darum predigen, daß ihr Werke thun sollet, damit ihr einen Gott aus denselben machet. Denn etc.} Denn dadurch werde ich GOttes Sohn, nicht daß ich das thue und jenes lasse, sondern daß Christus mein HErr ist.

11. Der Teufel ist dem göttlichen Worte also feind, daß er auf allerlei Weise sich versucht, ein ander Loch zu finden. Er meint nicht die äußerlichen Werke und die Früchte des Fleisches, als Hurerei, Diebstahl, Wucher, wiewohl sie es auch sind; sondern die der Heilige Geist für Früchte erkennt und urteilt, als ob sich sein Thun und Predigen richte hinauf nach Christo. Thue deine Werk und Thun wirst du ein Christen-mensch. Da richte dich nach. Da findest du Feigen und Trauben. Jene sagen aber also: Wirst du nicht dieses oder das thun, so bist du verdammt; und legen dir dann die heilige Schrift vor. Dann kannst du dich davon nicht erretten,
\quellseite{407}
du bist gefangen; denn GOttes Name liegt dir im Wege, und du hörst, Christus habe es gethan. Wie dünkt dich, wenn du hörst Christum und GOtt nennen? Dann erschrickst du.

12. Alhier wisse, daß dir Christus auf zweierlei Weise vorgemalt wird. Einmal, daß er dich zu GOttes Sohne macht, wie er Mosen zum Gott über Pharaonem machte. Das geschieht nun durch sein Wort, da er dir kein Werk vorschlägt, das du thun solltest, sondern er will dein sein, und du sollst sein werden; durch seinen Tod und Blutvergießen will er dich selig machen.

13. Das wird kein falscher Apostel predigen: Jesum Christum kann niemand einen HErrn heißen, denn aus dem Heiligen Geiste [1 Cor. 12, 3.]. Daß Christus unser Schatz sei, das sei unser Siegel. So weißt du nun ehrlich, daß du durch Christum, ohne deine Werke, dahin kommst durch seine Güte. O wie ist der Teufel der Predigt so feind, und sieht sie so scheel an, und wird ihrer viel stürzen, die da jetzt meinen, sie haben das Evangelium, die Hauptpredigt. Er, der Satan, will diese Predigt stehlen.

14. Es ist nicht genug und damit ausgerichtet, daß man sagt: Christus thut es, darum thue es auch; sondern, Christus ist dein, und du bist sein. Dieser Knote bricht dem Satan den Hals, daß die Engel sagen in den Wiegenachten: „Euch ist geboren Christus der Heiland“, euch ist er geschenkt [Luc. 2, 11.]. Diese Worte zerhören und nehmen dem Teufel sein Reich. Andere Werke, als Bilderstürmen, und anderes, schadet dem Teufel nicht. Also mußt du Christum erstlich haben als deinen Schatz, und als die Hauptpredigt. Darnach sollst du seinem Werke folgen, wenn du sein Wort hast. Dann ist das höchste und neue Gebot, nämlich die Liebe, welche er von dir nicht eher will gethan haben, du hast denn zuvor einen neuen Geist, das ist, du hast denn sein Wort und glaubest demselbigen. Da hörst du nicht von ihm, wie du einem Bilde sollst einen Arm oder Bein entzwei brechen, wie die Schwärmergeister predigen.

\begin{schriftzitat}
V. 16–19. Und der HErr sprach zu Mose:
\end{schriftzitat}

\begin{quote}
Sage Aaron: Reck deinen Stab aus, und schlage in den Staub auf Erden, daß Läuse werden in ganz Egyptenlande. Sie thaten also. Und Aaron reckte seine Hand aus mit seinem Stabe, und schlug in den Staub auf Erden, und es wurden Läuse an den Menschen und an dem Viehe, aller Staub des Landes ward Läuse in ganz Egyptenlande. Die Zauberer thaten auch also mit ihrem Beschwören, daß sie Läuse heraus brächten; aber sie konnten nicht; und die Läuse waren beide an Menschen und Vieh. Da sprachen die Zauberer zu Pharao: Das ist GOttes Finger. Aber das Herz Pharao ward verstockt, und hörte sie nicht; wie denn der HErr gesagt hatte.
\end{quote}

15. Dies ist die dritte Plage, so über das Reich Egypten geht, daß es mit Läusen gestraft wird, ob die Egypter wollten frömmer werden, und sich zu GOtt bekehren.

16. Moses muß um dieser Zeichen willen leiden, daß seine Vocation nicht von statten geht; denn auch diese Zauberer Werke thun können, und ein Theil der Zeichen nachthun, bis auf das dritte Zeichen, da sie Läuse machen sollen. Da sind die Läuse so köstliche Thiere, daß sie der Teufel durch seine Werkzeuge, diese egyptischen Zauberer oder Schwarzkünstler, nicht kann nachmachen; da trägt und regt er den Schwantz und das Haupt umsonst empor. Sonst können die Zauberer alles nachthun, allein Läuse können sie nicht nachmachen. Ei, welch ein Hohn ist es diesem stolzen Geiste!

17. Ob es Mücken oder Wanzen gewesen sind, weiß man nicht; die Hebräi nennen’s Läuse, so da gelebt haben an den Menschen und Viehe; wie denn die Läuse pflegen zu thun. Diese kleinen Thierlein, wie gesagt, können sie nicht nachmachen, Mosi zum Troste, auf daß sie zu Schanden werden. Wie denn St. Paulus [2 Tim. 3, 8. 9.] auch sagt: Jambres und andere widerstanden Mosi; gleichwie jetzt die falschen Ketzer wider die Wahrheit sich auflehnen; aber sie werden es nicht hinausführen, sondern offenbar wird werden ihre Thorheit. Sie bekennen allhier, daß es GOttes Finger und Kraft sei, und werden darüber gar zu Schanden.

18. Warum hat er es aber nicht im ersten Zeichen gethan, sondern er macht sie im dritten Zeichen zu Schanden? Da wisse du, daß GOtt den Secten und Rotten ein wenig verhängt, daß sie einen Fortgang haben. Denn wenn es uns bald gelänge, und wir sie dämpften, so hätten wir nicht Ursache, GOtt anzurufen. Darum müssen die Bösen zunehmen, und mit ihrer falschen Lehre andere verführen, wie auch St. Paulus [1 Cor. 11, 19.] sagt: „Es müssen Irrthümer kommen, daß die Auserwählten bewähret
\quellseite{408}
würden.“ GOtt läßt ihr Gift unter etliche fallen und angenommen werden, daß es einen Schein habe, als wollten sie unser Ding alles unterdrücken; aber zu seiner Zeit fallen sie, und viel sind zu ihrer Zeit dem Irrthume zugefallen; aber hernach werden sie es nicht hinausführen. Eine Zeitlang nehmen sie zu, grünen und blühen; aber darnach gehen sie zu Boden. Ueber den Läusen müssen sie zu Schaden werden. Ein oder zwei Zeichen läßt sie GOtt thun, aber das dritte Zeichen können sie nicht vollbringen.

19. Wohlan, so merket es fleißig und wohl, daß der Satan nicht kann allerlei Wunder und Mirakel thun; wie allhier aus des Teufels Betrug die Teufelskünstler Mosi etliche Wunder nachthun, aber nicht alle, sintemal sie mit ihrem Beschwören keine Läuse heraus bringen. Denn ob seine Gewalt wohl groß ist, so ist sie doch nicht infinita, unmeßlich, als GOttes Gewalt und Macht, sondern sie hat ein Ziel und Maß. Es ist dem Teufel viel Dinges unmöglich, als Todte kann er nicht wieder lebendig machen, wie Christus den Lazarn vom Tode auferweckt [Joh. 11, 43. f.]. Auch kann er alte, verlebte, unfruchtbare Weiber nicht fruchtbar machen c. Der Teufel kann nicht rechtschaffene Wunder und Mirakel thun, sondern es sind betrügliche Mirakel.

20. Warum verhängt aber GOtt dem Teufel und seinen Schuppen, daß sie Mirakel thun? Ei, er will durch der falschen Propheten Mirakel die Gottesfürchtigen probiren und bewähren, wie im fünften Buch Mosis am dreizehnten, V. 1. ff., geschrieben steht: „Wenn ein Prophet oder Träumer unter euch wird aufstehen, und gibt dir ein Zeichen oder Wunder, und das Zeichen oder Wunder kommt, davon er dir gesagt hat, und spricht: Laß uns andern Göttern folgen, die ihr nicht kennet, und ihnen dienen; so sollst du nicht gehorchen den Worten solches Propheten oder Träumers. Denn der HErr euer GOtt versucht euch, daß er erfahre, ob ihr ihn von ganzem Herzen und von ganzer Seele lieb habet.“ Darum geschehen solche Mirakel den Christen zur Warnung.

21. Zum andern verhängt GOtt solches zur Strafe über die Welt, von wegen der Verachtung und Ueberdrus des göttlichen Worts, wie solches St. Paulus 2 Thess. 2, 9–12. bezeugt, da er spricht: „Des Antichrist Zukunft wird geschehen nach Wirkung des Satans, mit allerlei lügenhaftigen Kräften, Zeichen und Wundern, mit allerlei Verführung zur Ungerechtigkeit, unter denen, die verloren werden, dafür, daß sie die Liebe zur Wahrheit nicht haben angenommen, daß sie selig würden. Darum wird ihnen GOtt kräftige Irrthümer senden, daß sie glauben der Lüge, auf daß gerichtet werden alle,\footnote{Siebenbisch: also.} die der Wahrheit nicht glauben, sondern haben Lust an der Ungerechtigkeit.“

22. Sie sprechen aber, es sei der Finger GOttes im dritten Zeichen. Es sind wahrlich gelehrte Leute gewesen, daß sie also reden vom Finger GOttes. Die Propheten haben dahin auch gesehen. Aber was ist das für eine Rede, hat GOtt Finger? Der 19. Psalm, V. 2., sagt: „Du hast den Himmel gemacht, deiner Finger Werk.“ Also hat der HErr Christus durch den Finger GOttes die Teufel ausgetrieben. Der Finger GOttes ist der Heilige Geist, also deutet es GOtt selbst. Denn Matthäus\footnote{In der Siebenbischen steht hier (irrtümlich) „Lucas“ und an der folgenden Stelle „Matthäus“.} [Cap. 12, 28.] spricht: „Durch den Geist GOttes.“ Lucas [Cap. 11, 20.] nennt es: „Durch den Finger GOttes.“ So sie es nun selbst also heißen, so ist es recht, und so habe ich es nicht also ausgelegt. Gleich aber wie ein Mensch alles thut mit den Fingern, und wenn sie abgehauen sind, so kann man nichts fassen, treiben oder machen, also ist der Heilige Geist auch GOttes Finger, der den Menschen in die Herzen Gaben und Geschenke gibt und alles ausrichtet, und GOtt thut alles durch den Heiligen Geist, daß er der sei, der es ausrichtete.

23. Also sollten sie sagen: Bisher haben wir auch Zeichen gethan; aber da können wir nichts, allhier ist GOttes Finger, Moses macht Läuse, mit unserer großen Schande. Sie müssen bekennen, daß ihre Zeichen nicht aus GOttes Finger sind [Weish. 1, 8.]. Sie sagen: O! der Mann hat den Geist, ist voller Geistes, hat GOttes Finger; unser Geist hat nur einen Schein, er kann nicht hindurch dringen, wie Mosis Geist. Das ist, der Heilige Geist ist es, der alles ausrichtet.

24. Also geht es auch: Wenn GOtt gleich viel predigt, und träumet durch das Gesicht, oder verheißt durch das Evangelium, so geht es doch nicht, ist auch nichts; es ist gleich als ein Arm ohne Hand und Finger. Aber wenn der Heilige
\quellseite{409}
Geist kommt, und gibt es in das Herz, das ist, wenn nicht allein der Mund predigt, sondern auch der Heilige Geist in das Herz scheinet, alsdann thut der Mensch alles.

25. Ueber das sehen wir allhier auch, ob wohl Pharao sieht, daß seine Meister zu Schanden werden, so kehret er sich doch nichts dran, sondern er wird je länger verstockter und ärger; wie Salomo auch sagt Sprüchwörter am achtzehnten Capitel, V. 3.: „Wenn der Gottlose herrscht, da kommet Verachtung und Schmach mit Hohn.“ Sie schlagen es alles in Wind. Also thun die Unsern auch, sie meinen, es sei ein Uebergang, es werde wohl anders werden etc.

\begin{schriftzitat}
V. 20. 21. Und der HErr sprach zu Mose: Mache dich morgen früh auf, und tritt vor Pharao, siehe, er wird an das Wasser gehen, und sprich zu ihm: So sagt der HErr: Laß mein Volk, daß mir es diene. Wo nicht, siehe, so will ich allerlei Ungeziefer lassen kommen über dich, deine Knechte, dein Volk, und dein Haus, daß aller Egypter Häuser, und das Feld, und was drauf ist, voll Ungeziefer werden sollen etc.
\end{schriftzitat}

26. Das ist die vierte Plage, allerlei böse Würmer und schädliche Thiere, das ist, allerlei Fliegen; also haben es etliche verdeutscht. Aber unser Deutsch ist besser, Ungeziefer. Diese Plage soll die Egypter auch treiben zur Buße und Besserung, und daß sie die Kinder von Israel ließen aus Egypten ziehen.

27. Solch Ungeziefer kommt nun allein über Egyptenland, und nicht über das Land Gosen, darinnen die Kinder von Israel wohnten, welches ein sonderlich Wunderwerck ist, und GOtt damit seine Allmacht und Herrlichkeit hat beweisen wollen, daß er der Israeliter GOtt sei, und seines Volkes schone, aber die Egypter, als seines Volkes Feinde und Widersacher, strafen könne, daß man ihn, als den rechten, wahrhaftigen GOtt, allein ehrte, lobte und preisete. Wie er denn alle Mirakel und Plagen über Egyptern darum will haben ergehen lassen, daß er die Macht seiner Stärke und Herrlichkeit an den Egyptern sehen ließe, und sein Name in aller Welt gepreiset würde. Die Egypter sollten an diesen Plagen sehen, daß er allein GOtt wäre.

\section{Das neunte Capitel}

1. Wir haben im achten und neunten Capitel sieben Plagen, so über Egyptenland [er]gangen sind. Erstlich, daß die Wasserströme voller Bluts worden; darnach, wie GOtt über das Land geschickt hat Frösche, Läuse, Würmer, Pestilenz, Drüse oder Geschwär an den Leuten, und Hagel; und haben gesagt, daß Moses schier bei allen Plagen dazu gesetzt, daß Pharaonis Herz je länger je seher sei verhärtet und verstockt worden. Wenn er nur ein wenig Luft gewonnen, so hat er wieder angefangen zu tyrannisiren. Es ist mit ihm ergangen, wie man im gemeinen Sprüchwort sagt: Da der Kranke genas, er nie ärger maß. Wenn er aber wieder geplagt worden ist, so hat er sich denn fromm gestellt, und Mosen bei zehnmal mit Worten getäuscht, und es ihm keinen Ernst sein lassen.

\begin{schriftzitat}
V. 1–5. Der HErr sprach zu Mose: Gehe hinein zu Pharao, und sprich zu ihm: Also sagt der HErr HErr, der GOtt der Hebräer: Laß mein Volk, daß sie mir dienen. Wo du dich deß weigerst, und sie weiter aufhälst, siehe, so wird die Hand des HErrn sein über dein Vieh auf dem Felde, über Pferde, über Esel, über Kameele, über Ochsen, über Schafe, mit einer fast schweren Pestilenz. Und der HErr wird ein Besonderes thun, zwischen dem Vieh der Israeliter und der Egypter, daß nichts sterbe aus allem, das die Kinder von Israel haben. Und der HErr bestimmte eine Zeit, und sprach: Morgen wird der HErr solches auf Erden thun etc.
\end{schriftzitat}

2. Das ist die fünfte Plage, da die Pestilenz unter das Vieh kommt. Diese Strafe hat GOtt gedräuet, über die Gottlosen lassen zu ergehen, 5 Mos. 28, 21., und es sollte die Pestilenz ihnen eine Bußpredigt sein, daß sie durch diese Strafe abgeschreckt wären von Sünden. Wie denn Davids Königreich, Land und Leute auch mit Pestilenz drei Tage lang heimgesucht ward, als David sein Volk gezählet und gemustert hatte
\quellseite{410}
[2 Sam. 24, 15.], und straft GOtt des David Hoffahrt mit Pestilenz; wie er denn allhier Pharaonis Tyrannei mit der Pestilenz straft.

\begin{schriftzitat}
V. 6. 7. Und der HErr that solches des Morgens, und starb allerlei Vieh der Egypten; aber des Viehes der Kinder Israel starb nicht eins. Und Pharao sandte darnach, und siehe, es war des Viehes Israel nicht eins gestorben. Aber das Herz Pharaonis ward verstockt, und ließ das Volk nicht.
\end{schriftzitat}

3. Das ist auch ein Wunderwerk, daß GOtt einen Unterschied macht zwischen den Egyptern und Israeliten, und diese Plage nicht über die Kinder Israel kommt, die sonst über die Egypter geht, sondern, wie sonst die heilige Schrift [Ps. 34, 22. f.] sagt: Das Unglück trifft den Gottlosen allein. Darum schützt und errettet GOtt seine Gläubigen. Wie denn sonst die Kinder von Israel von andern Plagen auch frei waren, mit welchen GOtt die Egypter heimsuchte; in ihrem Lande Gosen, da die Kinder von Israel innen wohnten, da hagelte es nicht, da war es nicht finster, sondern licht, und da ward auch kein Erstgeborner ermürgt. Also sorgt GOtt für seine Christen, daß wenn gleich Landsterben oder andere Unfälle gemein kommen, so müssen sie ihnen nicht schaden.

\begin{schriftzitat}
V. 8. 9. Da sprach der HErr zu Mose und Aaron: Nehmet eure Fäuste voll Ruß aus dem Ofen. Und Mose sprengte ihn gegen Himmel vor Pharao, daß über ganz Egyptenland stäube, und Schwären und Drüsen auffahren, beide an Menschen und am Vieh, in ganz Egyptenland.
\end{schriftzitat}

4. Das ist nun die sechste Plage und Strafe GOttes, daß Menschen und Vieh Geschwäre und Drüsen an ihrem Leibe bekommen; welche Plage GOtt 5 Mos. 28, 21. 27. allen Gottlosen auch dräuet; geht also immerdar eine Plage und Strafe nach der andern daher, wie eine Bulge und Wasserwelle die andere treibt, bis daß der Garaus darauf kommt und Pharao mit allen Egyptern im rothen Meere ersäuft. Und also läßt GOtt dem Gottlosen Raum und Zeit zur Buße und Besserung, und zu seiner Bekehrung, wie denn St. Paulus Röm. 2, 4. sagt: „Die Langmüthigkeit GOttes vermahnet uns zur Buße“ etc.

\begin{schriftzitat}
V. 22. Da sprach der HErr zu Mose: Reck deine Hand auf gen Himmel, daß es hagele über ganz Egypten etc.
\end{schriftzitat}

5. Folgt die siebente Plage über die Egypter, als Hagel, mit welcher Plage GOtt auch hinter Pharao und seinen Egyptern her ist, sie Moses lehren will, und die bösen Buben fromm machen. Wie denn im Propheten Haggai, Cap. 2, 18., der Hagel eine Ruthe und Plage GOttes über die Bösen genannt wird, da GOtt spricht: „Ich plagte euch mit Dürre, Brandkorn und Hagel in aller eurer Arbeit, noch kehret ihr euch nicht zu mir, spricht der Herr“ etc. Weiter sagt der Text:

\begin{schriftzitat}
V. 16. Ich habe dich erwecket, daß meine Kraft an dir erscheine, und mein Name verkündiget werde in allen Landen.
\end{schriftzitat}

6. Wir haben oft gehört, daß GOtt mit uns also handelt, daß man es muß greifen, es sei ein göttlich Werk, auf daß man nicht gegen ihm stolz bleibe, und meine, der freie Wille könne auch etwas. Derhalben so führt er die Seinen herunter, auf daß ihnen alle ihre Weisheit, Macht, Gewalt und Witz zu kurz werde und zerrinne, und wir es fühlen müssen im Werke, wenn die Ochsen am Berge stehen, daß wir uns selber nicht helfen können, sondern verzweifeln müssen und allein bei GOtt Hülfe suchen. Wiederum, erhebt und erhöhet GOtt unsere Widersacher auch also sehr, daß sie stolz werden, und meinen, sie haben das Spiel gewonnen und stehe in ihrer Faust. GOtt gibt ihnen Reichthum genug; item, Witz, Anhang, Gunst, Weisheit, Verstand und Kraft, dagegen sie unserm HErrn GOtt nur ein Klipplein schlagen.

7. Aber GOtt will allhier sagen: Du hast mein Volk unter dich gedrückt, daß es gar feucht unter dir, und es scheint, als könnte ich ihm nicht helfen; jedoch sollst du mein Volk gehen lassen, wiewohl du es nicht gedenkst zu thun. Nun, ich habe es gemacht, und dir diesen Muth und Trotz, den du hast von deinem Reichthum und von deiner Weisheit, gegeben, daß du dich wider mich aufbläfest. Aber du weißt nicht, was ich thue. Du hast im Sinne, daß du große Ehre und einen herrlichen Namen daran habest, wenn du mein Volk unterdrückst; aber ich will dadurch gespeiset werden, wenn ich dich zu Boden stürzen und zu Asche machen werde.

8. Derhalben wird unser GOtt ein wunderbarer GOtt genannt, der da hilft aus der Sünde und Tode, und aus dem, das gar verloren war, und schlägt zu Boden alles, was da
\quellseite{411}
etwas vor ihm sein will, auf daß man sage: Dieser GOtt weiß allen zu helfen. Denn die, so verzweifeln wollen, tröstet und richtet er auf, und die Hoffährtigen kann er nicht leiden. Wie auch St. Petrus sagt in seiner [ersten] Epistel [Cap. 5, 5.]: „GOtt widerstrebet den Hoffährtigen, aber den Demüthigen gibt er Gnade.“

9. Solches sieht man erstlich nicht, warum und wie es GOtt thue, bis daß es ausgerichtet ist. Es müssen Pharao und seines Gleichen erweckt werden, die da aufgeblasen und hoffärtig sind, auf daß GOtt etwas Großes habe, daran er seine Weisheit und Gewalt beweise, und daßselbige zerschlage. Wie denn GOtt noch also thut; da er das Evangelium durch Christum, seinen Sohn, Johannem den Täufer und seine lieben Apostel wollte predigen und der Welt offenbaren lassen, da richtete er zuvor auf das gewaltige, mächtige Reich, als das römische Kaiserthum, und ließ kommen die weisesten und heiligsten Leute im Judenthum.

10. Da war Witz, Kunst, Stärke, Gewalt, Reichthum, Wohlust am höchsten, und das konnte die Welt meisterlich aufmutzen, und wenn sie sich ansahen nach ihrer Vernunft, wie mächtig sie wären, und wie weit ihr Arm in der Welt reichte, so fuhren sie empor in ihrem Sinne und stolzierten also, daß sie auch ihren eigenen Namen der Stadt Rom nicht nannten, sie griffen denn zuvor an das Barett und zogen es ab. Hielten also diesen Namen gleich als ein Idolum; daher sie auch im römischen Rathe beschlossen, daß es ein Gott wäre, und haben aller Welt Gut zu sich gezogen, waren prächtig, gewaltig und reich, und verachteten trotziglich GOttes Wort, nahmen keine Bußpredigt, Strafe, Vermahnung und Warnung an, waren gar verstockt, daß sie mit sehenden Augen nichts sahen, und mit hörenden Ohren nichts vernahmen, wie Jesaias sagt, Cap. 6, 10. Nun, GOtt hat Rom also gemacht, sie hat sich selbst nicht also erhöht. Aber hernach spricht GOtt, gleichwie alhier zum Pharaone: Ich habe dir diesen Trotz gegeben, auf daß die ganze Welt davon zu singen und zu sagen habe.

11. Wider diese gewaltige Monarchie ließ GOtt St. Petrum, den Fischer, St. Paulum und andere Apostel predigen, und hängte sie an diese Gewalt der Römer; gleichwie er alhier Mosen an den König Egypti hängt. Da möchte schier einer sagen: Wie ist GOtt so närrisch? Was greift er dies große Regiment an mit armen, elenden Menschen? Die Vernunft spricht: Was sollte ein Fischer ausrichten? Aber GOtt führt es hindurch, ob es wohl viel Blut kostet; es wurden darob viel tausend Märtyrer, die da frisch dahin starben, die man als die Krauthäupte weghieb.

12. Aber dies römische Reich ist zerstoben, daß man nicht weiß, wo es jetzt sei, allein daß der Pabst uns überredet hat, der römische Kaiser habe es. Aber er ziehe dahin und setze sich dahin; das Nest ist zerstört, und die Vögel sind ausgeflogen. Gleichergestalt ist das egyptische Königreich auch zu Grund gegangen und beiseite geräumt.

13. Also ehrt GOtt sein Wort, daß er dies Reich läßt untergehen; wie er auch alhier sagt zum Pharaone: „Ich will dich stürzen, auf daß man erkenne, daß ich der HErr sei.“ Da geht denn Moses in großer Schwachheit daher, wider die große Gewalt, so Pharao hat, und stürzt ihn gleichwohl in das rothe Meer. Also närrisch greift es unser HErr GOtt an.

14. Dies wird uns nun erstlich zum Troste, Licht und Verstand gesagt in GOttes Werken, daß du wissen solltest, wenn du in Nöthen und Aengsten bist, daß GOtt derjenige sei, der es erweckt, wenn der Teufel und Tod sich wider dich legen. Denn so spricht GOtt: Diesen Teufel, der dich fressen will, habe ich in meiner Faust, sein böser Wille und hoher Muth steht in meiner Macht und Gewalt. Ich habe es also geschickt: Du bist mein armes Würmlein [Jes. 41, 14.] und hast mein Wort in deinem Herzen, daran du glaubst; ich aber rufe dem Teufel, daß er dich fresse. Dann brüstet er sich, stellt sich zornig, als wollte er dich auf einen Bissen verschlingen; du aber bist dagegen kleinmüthig, furchtsam und erschrocken; noch, wenn du dich schwingst auf deinen Glauben, so ist GOtt über den Tod, Teufel, Welt, Sünde und Hölle, und über alles, daß dir der Teufel nicht ein Härlein krümmen könnte. Sondern wie ein großer Fisch nach einem Würmlein schnappt, so am Angel steckt, und verschlingt es, wird drüber zu Schanden, also sollte der Teufel und Welt mit ihrer Tyrannei wider die Christen auch zu Schanden werden.

15. GOtt erweckte den Teufel wider den Hiob und hetzte ihn\footnote{„ihnen“ in der Eislebenischen.} an, reizte den Teufel selbst mit
\quellseite{412}
Worten auf den Hiob und sprach [Hiob 1, 8.]: „Siehest du auch meinen Knecht Hiob, wie gerecht und unschuldig er lebet?“ und sagt: Da hast du ihn, gehe hin, nimm ihn, du hast alles, sein Leib und Gut, in deiner Hand, über die Güter und den Leib habe Gewalt; allein schone seines Lebens. Nun, der Teufel fährt hin, und nimmt ihm alle seine Güter und die Gesundheit, daß Hiob nicht anders meint, denn: Nun bin ich des Todes, und gedachte, er würde ihn gar verschlingen und fressen. Aber GOtt hat ihm ein Ziel gesteckt, und spricht: Ich bin auch da, du Teufel, du sollst mir ihn nicht fressen. Daher wurde Hiob wieder gesund, reich und mächtig, bekam wieder Kinder und Güter, und konnte der Teufel wider den Willen Gottes dem Hiob kein Leid thun und gar nichts schaffen.

16. Also, wenn wir auch um des Evangelii willen von den zornigen Fürsten überwältigt, und von Krankheit und Tod geplagt werden, so spricht flugs die Welt: Das hat der Teufel gethan; GOtt ist fromm, er thut es nicht. Dann spricht GOtt dagegen: Lieber, nimm mir diese Ehre nicht hinweg, ich habe es wahrlich gethan, liebes Kind, auf daß du mich erkennest, denn sonst erkennet du mich nicht, und würdest mich nicht anrufen [Pred. 3, 14.]. Denn Noth lehrt beten und rufen zu Gott.

17. Daher sollt du sehen, spricht GOtt, wie ich aus der Sünde, Tode, des Teufels und der Hölle Gewalt, auch aus allem Unglück gewaltiglich erretten könne; denn es gilt allein dazu, daß mein Name, meine Ehre, Macht und Weisheit erkannt werde. Denn wenn ich dir nicht mehr zuschickte, denn du mit deiner Stärke und Gewalt überwinden könntest, und wenn ich dich nur so lange hungern ließe, bis du für Geld Essen und Trinken kaufen und dir selbst helfen möchtest, dann gedächtest du nimmermehr an mich, du lernest meine Kraft nicht erkennen.

18. Wenn ich dich also ließe hinleben, daß du deine Feinde selbst überwändest, und dir nicht mehr Noth unter Augen stießen sollte, denn du mit deiner Vernunft und mit deinen Kräften könntest hinausführen: dann würdest du dich selbst, und nicht mich, erkennen, du würdest meiner vergessen, und über deiner Weisheit dich überheben. Derhalben so will ich es also machen, auf daß ich allein erkannt und geehrt werde, und will dir Noth und Unglück zuschicken, darinnen dir kein Engel noch einige Creatur helfen soll, denn ich allein. Und dann wirst du sehen, wie ich größer sei denn du, und mehr helfen könne denn Menschen.

19. St. Paulus zum Römern am neunten Capitel, V. 17., citirt solches, daß GOtt sich rühme (und man soll ihm lassen diesen Ruhm), daß er uns Krankheit und alles Unglück und Anfechtung zuschicke. Wiewohl es der Teufel und böse Leute thun, dennoch so rühmt er sich, daß er ihren Muth und ihr Herz in seiner Hand habe, und gebrauche des Satans und der gottlosen Leute dazu, die Seinen zu drücken und zu plagen; sie müssen seine Ruthen sein, damit er seine Kinder stäupe. Wie denn allhier auch geschieht. Pharao ist wider GOtt und GOttes Feind, so ist GOtt auch sein Feind, und er drückt und martert auch das Volk Israel übel. Aber GOtt rühmt sich allhier und spricht: Diese Gewalt und Muth hättest du nicht, wenn ich ihn dir nicht gäbe. Pharao muß der Kinder von Israel Ruthe sein.

20. Hier wird nun wieder erregt diese Frage: Ob GOtt derjenige sei, der die Menschen verhärte und zwinge zum Bösen und zur Sünde? Warum verdammt er denn die Menschen? Die Vernunft schließt allhier: Sollte GOtt die Sünde verdammen wollen, so würde er nicht heißen sündigen, oder würde den Ungerechten und Gottlosen nicht geschaffen haben. Nun, die Vernunft will allezeit GOtt hofmeistern, ob er Fug und Recht habe, will GOtt messen nach ihrem Gesetze und Gedanken. GOtt sollte säuberlicher handeln und nicht also erschrecken, sondern auf die und jene Weise es machen, und stellt GOtt also ein Gesetz vor. Aber das mußt du aus deinem Kopfe lassen, wenn du von GOtt reden willst, daß du kein Gesetz oder Maß auf GOtt gebest; denn er ist nicht eine Creatur, er ist unermesslich.

21. Dem Menschen ist ein Maß gesetzt, ich soll so und so thun; mein Leben ist endlich, es kann gefasset werden, und hat eine Regel, Maß, Weise und Gesetz [Weish. 11, 22.]. Da du mit GOtt also auch handeln wolltest\footnote{Eislebensche: wollest.}, so hast du GOttes gefehlt. Denn was da mit GOtt vorgenommen wird nach Gesetz, Maß und Ziel, das trifft nicht zu. Die Vernunft kann nicht höher, denn daß sie gedenkt: Also, und nicht anders sollte GOtt es machen, und urtheilt bald also,
\quellseite{413}
spricht: Ist es doch nicht gut, daß man verstößt? Und macht ihm also ein Maß; sie meint, GOtt sei wie ein Mensch, daß man von GOtt als von Menschen urtheile. Also verstehet es die Vernunft nicht, und will mich klug sein und von GOtt richten [Matth. 11, 25. Ps. 51, 6.]. Aber GOtt gibt dir Gesetze, und nimmt von dir keines; er steckt dir ein Ziel, und du nicht ihm. Darum ist es nicht recht, daß du es also willst, und also für recht und gut ansiehst, sondern wisse, daß er es also will haben und also gebeut; sein Wille ist gesetzt über alle Gesetze. Wenn er spricht: Ich will es also haben, dann so ist es über alle Gesetze, denn er ist ein unendlicher GOtt, und hat es Macht und Fug.

22. Sagt man aber: Ja, ich verstehe es nicht, daß es gut sei, daß er verstößt. Ja, Lieber, vor deinen Augen ist es böse; meinst du, daß du GOtt seiest? GOtt hat kein Maß, Gesetz oder Ziel (wie gesagt), darum so kann er dawider nicht thun, er kann wider Gesetz nicht sündigen, dieweil ihm keines vorgestellet. Derhalben ist es gut alles, was er thut [1 Mos. 1, 31.].

23. Es fließt auch daher eine andere Frage: ob GOtt zur Sünde treibe? Solches macht, daß ich GOtt fasse in einen Ring und Zirkel, oder in ein Glas, darinnen ich ihn will behalten. Er hat mir vorgeschrieben, wie ich leben und ihm dienen solle; da meine ich denn, er solle auch also leben. Er gibt das Gesetz aus, aber er nimmt es nicht wieder hinaus. Es gebühret niemand, denn allein GOtt, Gesetz und Lehre zu geben, wie man leben und fromm sein solle. GOtt aber soll ich kein Gesetz ordnen, wie er die Welt oder Menschen regieren möge. So halte du es, wie du willst, dennoch ist es recht, was GOtt thut; denn es ist sein Wille nicht unrecht noch böse, er hat nicht Maß oder Gesetze, warum er diesen erleuchtet, oder jenen verstößt. Sollte ich hierin GOtt messen und urtheilen nach meiner Vernunft, so ist er ungerecht und hat viel mehr Sünde denn der Teufel, ja, er ist erschrecklicher und greulicher denn der Teufel. Denn er handelt und geht mit uns um mit Gewalt, plagt und martert uns, und achtet unser nicht.

24. Hierüber möchte einer thöricht werden, wenn er nicht seine Vernunft gefangen nimmt, und aus dem Kopfe ihm treiben läßt alle solche Gedanken, und beruht nur darauf, daß GOtt niemand messen, oder ihm Gesetze vorschreiben solle; denn GOtt sei gar exlex,\footnote{das ist, vom Gesetze ausgenommen.} wie man sagt. Aber man kann dies die Vernunft nicht bereden, viel weniger kann man sie es überreden, oder ihr aus den Augen reißen das heillose, verfluchte Grübeln und Forschen in so hohen, unbegreiflichen Sachen, da sie stets spricht: Quare? Cur? Warum? Dann ist das Gesetz da. Wenn GOtt dies oder jenes darnach thäte, dann wäre es recht. Aber mit diesem Messen bringt man sich um Leib, Leben und um unserm Herrn Gott, da heißt es: Miß in des Teufels Namen hin. Aber ein jeglich Herz, das da sagen kann: Lieber GOtt, mache es, wie es dir gefällt, ich bin zufrieden; das kann nicht untergehen, aber die andern müssen zu Boden gehen.

25. Darum spricht St. Paulus zum Römern am 9. Capitel, V. 20.: „Wer bist du, der du mit Gott rechtst?“ Du bist ein Mensch, und willst mit GOtt rechten; wo willst du es nehmen? Du mußt das Gesetz haben, willst du mit GOtt rechten, nämlich, also müsse man thun; und dieweil es also nicht gethan ist, so ist man dem Gesetze nicht nachgekommen. Hui, willst du auch mit GOtt also handeln? Das gebühret sich nicht. Mit dem Nächsten magst du also umgehen, der hat das Gesetz, das soll er und du thun, nicht rauben, stehlen, ehebrechen \&c., aber Gott will thun, wie es ihm gefällt, und muß also thun, denn sein Wille ist das Gesetz, es kann nicht anders sein.

26. Der Mensch wird darum fromm genannt, wenn er handelt und lebt nach dem Gesetze. Mit GOtt kehre es gar um; da heißt ein Werk darum gut, daß es GOtt thut. Mein Werk ist darum nicht gut, daß ich es thue, sondern daß es nach dem Gesetze GOttes gethan ist, darinnen mir vorgeschrieben ist, was ich thun soll. Ich muß aus meinem Sinne in ein Höheres treten, nämlich in das Gesetz GOttes. GOtt ist nicht darum fromm, daß er dies Werk thut, sondern darum ist das Werk recht, gut, heilig und wohlgethan, denn er selbst thut es; und also die Gutheit herkommt von GOtt, und nicht von dem Werke. GOtt ist der Thäter, und nimmt die Gutheit nicht von dem Werke oder Gesetze. Aber wir nehmen die Gutheit nicht davon, daß wir Gutes verbringen,\footnote{verbringen = vollbringen.} und davon bei den Leuten geachtet werden, sondern, daß dem Gesetze recht geschehe, so muß es durch den Heiligen Geist erfüllt werden
\quellseite{414}
; alsdann so bekommen wir auch den Namen, daß wir fromm sind.

27. Unsere Vernunft sieht durch ein gefärbt, roth oder blau Glas, das kann sie nicht von den Augen thun; darum alles, was sie ansieht, das muß auch roth, blau oder grün sein; sie kann diesen Pharaonem nicht wohl auslegen, daß GOtt zum Bösen treibe, und verhärtet oder reizt entweder zum Guten oder Bösen. GOtt thut daran wohl, und nicht Unrecht. Aber der, so getrieben wird, der thut Unrecht; denn er hat GOttes Gebot vor ihm, daß er nicht also thun sollte, und der Teufel treibt ihn doch, daß er also handelt und thut, und nicht lebt, wie GOtt haben will. GOtt will, du sollst sein Gesetz für dich haben; so reizet dich der Teufel, daß du wider das Gesetz handelst. Wollest\footnote{Eislebensche: Wollest.} du denn sagen: Ist denn GOttes Wille wider sich selbst? Das ist zu hoch. GOttes Wille ist da, aber wie das zugeht, das soll ich nicht wissen.

28. Ich soll heruntersehen, was GOtt haben will. Nun hat er mir seinen Willen offenbart durch das Gesetz und Evangelium, und gelehrt, was ich thun soll; damit soll ich umgehen und nicht hinaufflettern, und fragen, warum GOtt dies oder jenes thue? Laß solches anstehen. Wenn du aber zum Glauben und wahrhaftigen Verstande kommen bist, und das Kreuz erfahren hast, so wirst du es verstehen.

29. Die Vernunft fähet allezeit oben am Dache an zu bauen, und nicht unten. Wie man ihrer denn viel findet, die da nimmermehr von Christo haben predigen hören, sind rohe und wilde Leute, martern\footnote{martern = fluchen. „Marter“ war damals ein häufiges Fluchwort, besonders unter den Landsknechten.} und fluchen, als wären sie voller Teufel, und suchen nur zum ersten, warum GOtt dies oder jenes thue, kommen mit den beschissenen Füßen und der blinden Vernunft hinauf an das Licht, und messen GOtt nach der Vernunft. Aber wir sollen vor uns nehmen die Weise, welche GOtt St. Paulo [1 Tim. 6, 19.] gegeben hat, und am Grunde anheben, das Dach wird sich dann wohl finden, laß GOtt mit seinem heimlichen Rathe mit Frieden, und klettere nicht hinauf mit deiner Vernunft ans Dach. Er will dich nicht also hinauf haben, sondern er kommt zu dir, und hat eine Leiter, einen Weg und Brücke zu dir gemacht, und spricht: Ich steige vom Himmel zu dir hinab, und werde Mensch in der Jungfrau Marien Leibe, liege in der Krippe zu Bethlehem, leide und sterbe für dich; da glaube an mich, und wage es auf mich, der ich für dich gekreuzigt bin. Matth. 9, 21. 22.

30. Also steige ich gen Himmel, und da werde ich dann nicht hinauf in die Gottheit klettern und grübeln. Man sollte im Jahre nur Ein Mal von der Gottheit predigen, auf daß man wüßte, daß in Sachen der Seligkeit von unten anzufahen wäre, das ist, wie Christus zu uns käme, daß man predigte, wie dies Kind, Christus, Milch und Butter ist, an der Mutter Brüsten liegt, und zu Bethlehem zu finden sei, und da lernen, warum Christus kommen, was man an ihm habe [Jes. 7, 15. 1 Petr. 2, 2.]. Wenn ich zu GOtt sagen wollte: Warum thust du das? so antwortet er: Ich weiß es wohl, was dahinter ist. Wenn wir das Quare unterlassen könnten, so würde der Teufel nicht hereinkommen mit solchen und dergleichen Fragen, ob wir versehen sind zur Seligkeit oder nicht? Item, wie Christus könne GOtt und Mensch sein? \&c.

31. Sollte man nicht vielmehr predigen vom Glauben und von der Liebe? Ja, sagt man, ich habe das lange wohl gewußt. Aber, Lieber, begib dich nicht auf solche Fragen, handle du mit der Menschheit Christi, da bist du gewiß, daß GOtt seinen Sohn in das Fleisch geschickt hat; laß ihn drinnen stecken, allhier suche ihn, er hat sich in der Jungfrau Marien Leib hinein gesenkt, und uns seine Menschheit vorgelegt, da will er, daß du dieselbige solltest erkennen, anschauen, und dich drinnen üben. „Er ist der Weg, die Wahrheit und das Leben“ \&c. [Joh. 14, 6.]. Dennoch wollen wir höher steigen, und wissen, wie er dies oder jenes versehen hat, diesen verstößt, den andern nicht. Wer klug und weise ist, der bleibe auf dieser vorgesteckten Bahn. Er kommt erst zu uns, und wir steigen nicht eher zu ihm gen Himmel, sondern er wirft den Sohn herunter ins Fleisch, läßt ihn geboren werden; darnach führt er denselben, läßt ihn schlachten und kreuzigen. Dies ist der Zweck, dahin wir sollen sehen und zielen.

32. Wie gibt der HErr Christus dem Apostel Philippo eine Schlappe, der auch wunderliche Gedanken von GOtt hatte, fragte, wo GOtt der himmlische Vater wäre, was er machte, ob er im Himmel Schalken ausnähme? sagte [Joh. 14, 8.] zu Christo: „Zeige uns den Vater, so genüget uns.“ Da antwortete der Herr Chri-
\quellseite{415}
Jesus und sprach: Allhier ist der Vater, weiset auf sich, spricht: „Wer mich siehet, der siehet auch den Vater.“ Willst du durch einen andern Weg gen Himmel zu GOtt klettern? Er spricht: „Siehe, Bruder, der Vater ist in mir, und ich in dem Vater“ [Joh. 14, 10.]. Halte deine Augen fest auf mich, durch meine Menschheit kommt man auf den Vater, der Vater schließet sich in meine Menschheit, und hat sich der Vater durch meine Menschheit der ganzen Welt vorgelegt. Also heftet er ihn an seine Menschheit, und reißt ihn herum von den irrigen Gedanken.

33. Denn, wenn ich sage: Christus, so vom Vater gesandt worden, ist gestorben, und hat mich armen, verdammten Sünder erlöst, alsdann so komme ich bald zum Vater. Frage ich aber: Wer hat es ihn geheißen? so wisse: Er hat es gerne gethan, und hat es für sich gethan aus lauter Liebe, Gnade, Güte und Barmherzigkeit. Da finde ich denn GOtt den Vater recht, und da ist er eitel Güte und Liebe, denn wir sehen seine Gutheit im Sohne [Joh. 3, 16. 16, 15.]. Darum lasse diese Frage fahren, und klimmere nicht hinauf, warum er dieses oder jenes also gemacht habe. Also muß ich meinen Gedanken begegnen, daß sie zu Boden fallen. Denn ich habe einen andern Weg, den ich gehen muß und diese Gedanken fahren lassen. Wie denn der HErr Christus zu St. Philippus sagte, der auch dort hinaus zu hoch sahe: Siehe, Philippe, „wer mich siehet, der siehet auch den Vater“ [Joh. 14, 9.].

34. Also heftet der HErr Christus seiner Jünger Gedanken, Herz und Augen allezeit auf sich, und spricht: Wenn du mich siehest, wie mich der himmlische Vater gesandt hat, daß ich dir predigen solle, und um deinetwillen sterben, da hast du des himmlischen Vaters Willen und Wohlgefallen. Glaubst du nun das, so wirst du selig, und kannst nicht erschreckt werden, sondern lebst ewiglich in diesem Glauben. In diesem Glauben und mit diesem Herzen geht man hinan, und wenn sich ein Mensch also heftet und bindet auf die Menschheit Christi, in welcher alle Schätze und Reichthümer sind [Col. 2, 3.], dann so findet sich eine süße Predigt, warum GOtt den Pharaonem verstößt, und wie er mit der Versehung umgehe. Wenn ich in das Evangelium komme, dann danke ich GOtt. Denn ich habe Christum, der für mich gestorben ist, der da ist ein HErr über alles, und der Vater hat mir ihn vom Himmel gegeben und geschenkt. Darum so hat er auch aller Feinde Herzen in seiner Hand, das weiß ich, und das tröstet mich auch, daß ich meinen Glauben stärken kann, und sagen: Es hat nicht Noth, Christus und der Vater, den ich habe, die haben alles in ihrer Faust und Gewalt [Röm. 8, 33. 34.].

35. So sollen mir nun diese Sprüche sparen bis zur Zeit der Noth, wenn wir in Kerkern liegen, und uns mit dem Schwert gedr äuet wird, daß ich sage: Das Schwert wird nicht schneiden, es will’s denn mein himmlischer Vater haben. Also kann ich dieser Worte gebrauchen zum Nutzen meiner Seligkeit, wenn ich im Leiden und Anfechtung bin. Denn sonst verdrießt mich es auf GOtt, ich werde auf ihn läunisch etc. Ich bin zu grün dazu, daß ich oben will anheben. Es ist gleich, als wenn ich einem kleinen Kindlein wollte Malvasier zu trinken geben, der da gehört für große, starke Leute zu trinken. Wenn ich alt bin und müde von Arbeit, und trinke Malvasier, dann so schmeckt er mir wohl, und stärkt mir das Leben, da sonst ein Kind den Tod sollte daran trinken. Also hätte GOtt diesen Spruch nimmermehr gesagt, wenn Mose nicht in Nöthen wäre gewesen, und in solchen Nöthen, die da scheinen, als wollten sie nimmermehr ein Ende haben. Da will GOtt sagen: Mose, und du Volk Israel, haltet fest, es hat keine Noth, ich habe es also gemacht.

36. Darum so sieht die Vernunft die Zeit und die Person nicht an; aber in der Zeit, da Noth und Angst ist gewesen, da sind diese Sprüche gehandelt worden. Du willst davon reden, wenn du bei der Zeche und im Bierhause sitzest; wenn du da deines Weibes und Biers wartest, und ließest die Frage von der Versehung unterwegen. Du bist noch nicht die Person, noch ich auch nicht, die davon reden soll, es ist zu frühe drum. Man sage zu demselbigen: Weißt du auch, was Christus ist, wie er geboren sei, was er mit seinem Leben und Sterben habe ausgerichtet und gethan? Man frage sie wieder: Bist du auch um des Evangelii willen jemals in Gefahr des Todes gewesen? Spricht er: Nein; so antworte du: Was fragst du denn darnach, das dir nicht müße ist, und auch nicht befohlen ist, daß du es wissen solltest? Und warum willst du diese hohen Dinge wissen, der du noch nie einig Kreuz, Trübsal und Anfechtung erfahren hast, noch Christum verstehst?

\quellseite{416}
37. Also pflege ich sie abzuweisen, die viel von der Versehung fragen und wissen wollen, daß ich sage: Hebet nicht zu hoch an, ihr werdet sonst den Hals abbringen und Mordprüfungen thun; gehet aber erst hin gen Bethlehem, und suchet das Kindlein Christum in der Krippe, und sehet, wie die Mutter Maria mit dem Kindlein Christo umgehet, und wie Christus für euch gestorben sei, und wann er für euch gelitten, und was er eurethalben gethan habe. Item, von diesen Stücken laß dich hören, und gib Bericht, wer du auch seiest; dann will ich dir antworten auf die Frage von der Versehung.

38. So ist nun dies die Summa dieses Capitels, daß man GOtt in seinen Werken nicht messen, urtheilen noch richten solle, sondern er soll alles messen und urtheilen, und sein Messen und Wille ist sein Sinn. Er mache es, wie er wolle, Wo kein Gesetz ist, da ist auch keine Sünde noch Unrecht; wo aber Sünde und Unrecht sollen sein, da muß Gesetz vorhergehen. Die Vernunft urtheilt sich und alle Menschen nach dem Gesetz, und will GOtt auch also achten; darum so fehlt sie. Wer das nicht verstehen kann, der schweige nur stille und laß es andere urtheilen. GOtt hat nicht ein Gesetz; sondern wie er will, so ist es gewollt, sein Wille ist sein Richtigkeit, Maß und Gewicht. Diesen Spruch befehle ich euch, daß ihr ihn brauchet, wenn die Noth herdringt, auf daß ihr lernet, GOtt also erkennen, und auf ihn trotzen, wie auch Moses allhier thut.

\section{Das zehnte Capitel}

\begin{schriftzitat}
V. 1. 2. Und der HErr sprach zu Mose: Gehe hinein zu Pharao, denn ich habe sein und seiner Knechte Herz verhärtet, auf daß ich diese meine Zeichen unter ihnen thue, und daß du verkündigest vor den Ohren deiner Kinder, und deiner Kindeskinder, was ich in Egypten ausgerichtet habe, und wie ich meine Zeichen unter ihnen beweiset habe, daß ihr wisset, ich bin der HErr.
\end{schriftzitat}

1. Wir haben am Ende des neunten Capitels gehört, wie dieser Spruch zu verstehen sei: „Ich will das Herz Pharaonis verstöcken“ etc., nämlich, wenn wir von GOtt reden wollen, daß wir vor allen Dingen unten anfangen sollen, und den Weg erkennen lernen, der uns von GOtt vorgestellt ist, nicht einen eigenen Weg suchen, oder aus eigener Vernunft uns hinan machen; denn sonst wird das Hinterste zuvörderst, und das Vörderste zuhinterst kommen, und einen unglückseligen Gang gewinnen [2 Mos. 33, 23.]. Es ist nicht vergebens geschrieben, ist auch kein Narrenspiel oder Gaukelwerk, daß GOtt Christum, seinen Sohn, hat in das Fleisch kommen lassen, daß er uns Menschen erschiene, und uns erleuchtete [Jes. 43, 11. Joh. 12, 47.]. Wenn es nicht hoch vonnöthen wäre gewesen, so wäre er wohl droben im Himmel geblieben etc. Also toll und thöricht sind wir, daß wir verlassen das Licht, so er in seinem Fleische uns angezündet hat; dann geht es auch also, wenn wir etwas Anderes suchen, daß wir darüber zu Narren werden.

\begin{schriftzitat}
V. 3–11. Also gingen Mose und Aaron hinein zu Pharao, und sprachen zu ihm: So spricht der HErr, der Hebräer GOtt: Wie lange weigerst du dich vor mir zu demüthigen, daß du mein Volk lässest, mir zu dienen? Weigerst du dich mein Volk zu lassen, siehe, so will ich morgen Heuschrecken kommen lassen an allen Orten, daß sie das Land bedecken, also, daß man das Land nicht sehen könnte, und sollen fressen, was euch übrig, und errettet ist vor dem Hagel, und sollen alle eure grünen Bäume fressen auf dem Felde, und sollen erfüllen dein Haus, aller deiner Knechte Häuser, und aller Egypter Häuser, dergleichen nicht gesehen haben deine Väter, und deiner Väter Väter, sint der Zeit sie auf Erden gewesen, bis auf diesen Tag. Und er wandte sich, und ging von Pharao hinaus. Da sprachen die Knechte Pharao zu ihm: Wie lange sollen wir damit geplagt sein? Laß die Leute ziehen, daß sie dem HErrn, ihrem GOtt, dienen. Willst du zuvor erfahren, daß Egypten untergegangen sei? Mose und Aaron wurden wieder zu Pharao gebracht, der sprach zu ihnen: Gehet hin, und dienet dem HErrn, eurem GOtt. Welche sind sie aber, die
\end{schriftzitat}

\quellseite{417}
hinziehen sollen? Mose sprach: Wir wollen ziehen mit Jung und Alt, mit Söhnen und Töchtern, mit Schafen und Rindern; denn wir haben ein Fest des HErrn. Er sprach zu ihnen: Awe ja, der HErr sei mit euch. Sollte ich euch, und eure Kinder dazu ziehen lassen? Sehet da, ob ihr nicht Böses vorhabt? Nicht also, sondern ihr Männer ziehet hin, und dienet dem HErrn; denn das habt ihr auch gesucht. Und man stieß sie heraus von Pharao etc.

2. In diesem zehnten Capitel sehen wir, wie dämisch\footnote{dämisch = dumm, thöricht, schläfrig, langsam, unbeholfen. — Erlanger: „dänisch“.} der Teufel sei, und wie ungerne er weicht. Der König Pharao hat bis anher alle Zeichen in Wind geschlagen, und je mehr deren kommen, je weniger er ihrer achtet. Wiewohl er sich demüthigt, bittet Gnade, und begehrt, daß ihm die Sünde soll vergeben werden, und ist so fromm, daß Moses muß für ihn bitten, und bekennt sich, er sei ein Sünder: noch, wenn es zum Werke und zum Treffen kommt, und es soll sein, wie Moses gesagt hat, daß er die Kinder von Israel aus Egyptenland ziehen lasse, da spricht er zu Mose: Trolle dich etc. Denn Gott ließ den Teufel über sein Herz regieren, und ließ ihn machen, wie es nur dem Teufel gefiel.

3. Was GOtt nicht regiert, das regiert der Teufel. Wenn GOtt regiert, so macht er es also, daß der Mensch lustig und willig wird zu thun, was GOtt gefällt, das ist, er gibt ihm den Heiligen Geist in das Herz. Aber wenn er den Heiligen Geist nicht ausschüttet oder gibt in das Herz, so regiert der Teufel, und gibt dem Menschen ein nach alle seinem Willen. Wie St. Paulus zum Timotheo [2. Ep. 2, 26.] auch sagt, daß die Gottlosen vom Teufel gefangen sind, zu thun nach alle seinem Willen. Er hat sie also gefaßt, daß, was diese Menschen thun, das thun sie ihrem Gott, dem Teufel, zu Dienst. Denn er hat sie eingenommen und gefangen, daß sie seinen Willen thun. Sie gehen mit Lust, Liebe und großem Ernst hinan, auszurichten, was er ihnen eingibt. Darum verstößt sie Gott, und thut seine Hand ab, und gibt sie dem Teufel dahin.

4. Also geht es allhier diesem Könige auch, er stellt sich, als wollte er sich bessern, aber es ist Heuchelei und ein Schein; wenn es der Heilige Geist nicht thut und treibt, so währt es nicht lange mit dieser Heuchelei. Erstlich erlaubt er den Juden, daß sie aus Egyptenland ziehen sollen, aber er hielt es ihnen nicht. Da nun die Heuschrecken kamen, will er allein, was Männer sind, ziehen lassen. Aber Moses spricht: Sie sollen alle ausziehen, Männer und Weiber, Jung und Alt, Groß und Klein. Da spricht Pharao: Awe ja, GOtt ehre euch. Welch eine stolze Antwort gibt der Eselkopf, spricht: Der HErr sei mit euch. Als wollte er sagen: Je, daß euch GOtt ehre, GOtt ehre euch, lieben Gesellen, wollt ihr da hinaus? Nun sehe ich, was ihr im Sinne habt; ihr könntet das wohl ausrichten, und Weib und Kinder allhier lassen; da hätte ich Lust zu, daß ihr mir entziehen wollet; ja, ich wollte es euch bestellen und schenken.

5. Das im Texte gesetzt wird für die achte Plage, muß man dieweil also nennen. Unsere Leute heißen es Raupen, und mag etwas fast gleich sein unsern Heuschrecken, so lange Füße haben und hüpfen, mögen aber etwas größer sein. Es werden auch wohl Heuschrecken genannt ein Thierlein, das man hat pflegen zu essen, rein und gesund; wie auch Johannes der Täufer Heuschrecken gegessen hat [Matth. 3, 4.]. Es hat pflegen wegzufressen alles, was auf den Bäumen und auf dem Felde wächst, ist grün gewesen, wie unsere Raupen noch in unserm Lande thun. Aber wir haben jetzt nicht solche Heuschrecken.

6. Nun gleichwie etliche Lande ihre eigenen Vortheile haben, also haben sie auch ihre eigenen Plagen. Alhier hat man Vier und keinen Wein, denn Wein findet man sonst anderswo. Also haben sie diese Plage mit den Arphephagel\footnote{Dies Wort: „Arphephagel“ ist vielleicht ein aus „Arbe“ und „Hargol“ (3 Mos. 11, 22.) gebildetes Compositum.} in Egypten auch gehabt, daß die Heuschrecken alles aufgefressen haben, und haufenweise gekommen sind, und sich aus der Luft daher geschwungen, gleich als wenn eine Wolke fällt. Es sind nicht eine, zwo oder drei daher-geflogen, sondern mit einem großen Schwarm eingefallen, und [haben] gleich die Sonne bedeckt, und die Luft finster gemacht; sind dahergefallen wie ein Schnee, der alles bedeckt. Ist also dieser Vogel den Egyptern nicht unbekannt, gleichwie uns die Raupen auch nicht fremd sind;
\quellseite{418}
wir kennen sie wohl, allein, daß oft ein Jahr mehr Raupen kommen denn das andere. Also ist dieser Vogel den Egyptiern auch wohl bekannt gewesen. Aber daß sie allhier mit so großen Haufen und Menge kommen, das war ihnen fremd und eine sonderliche Plage; noch schlugen sie es in Wind.

7. Diese Heuschrecken nennen sie den Tod; denn Pharao spricht: „Bittet den HErrn, euren GOtt, daß er diesen Tod von mir wegnehme.“ Denn diese Heuschrecken thaten solchen Schaden im Lande an den Früchten, fraßen alles Kraut und Gemächs hinweg, daß hernachher ein großer Hunger und Theurung darauf folgte; sie verderbten und fraßen alle Früchte im Lande hinweg, daß die Leute Hungers starben.

8. Dies achte Zeichen oder Plage mit den Heuschrecken geht auch hinweg, und wird nichts daraus, da sie doch hieran GOttes Zorn, Ungnade und Strafen gegen den ungehorsamen Menschen sollten erkannt haben und von Sünden abgelassen. Wie denn GOtt der Menschen Sünde dräuet zu strafen mit Heuschrecken, 5 Mos. 28, 38.: „Du wirst viel Samens ausführen auf das Feld, und wenig einsammeln; denn die Heuschrecken werden’s abfressen.“

\begin{schriftzitat}
V. 21–23. Der HErr sprach zu Mose: Recke deine Hand gen Himmel, daß [es] so finster werde in Egyptenland, daß man es greifen mag. Und Mose reckte seine Hand gen Himmel. Da ward eine dicke Finsterniß im ganzen Egyptenlande, drei Tage, daß niemand den andern sahe, noch aufstund von dem Ort, da er war, in dreien Tagen. Aber bei allen Kindern Israel war es licht in ihren Wohnungen.
\end{schriftzitat}

9. Nun folgt die neunte Plage. Es kommen so dicke Finsternisse, daß man sie greifen möchte, wie wir Deutschen pflegen zu reden. Solche Plage der Finsterniß hat GOtt gebräuet zur Strafe allen Gottlosen, die GOttes Wort ungehorsam sind, 5 Mos. 28, 29. Und es sind auch noch alle gottlosen Tyrannen und Ketzer geistliche Blinde. Da bittet Pharao nicht mehr, daß die Plagen aufhören mögen, sondern er spricht truckes: Sie sollen wegziehen. Da aber das Licht wieder in das Land kommt, denkt Pharao: Ach ihr Buben, es wird Gaukelspiel sein, und spricht zu Mose: Trolle dich von meinen Augen. Moses spricht: Ja, ich will es thun, wie du gesagt hast, ich will nicht mehr vor deine Augen kommen. Ich habe für dich gebeten, und Zeichen vor dir gethan; aber es hilft alles nichts. Also haben wir von neun Zeichen oder Plagen gehört.

10. Das ist aber ein herrlich Mirakel und Wunderwerk GOttes, daß [es] im Lande Gosen, da die Kinder Israel wohnen, licht ist. Wie denn sonst die andern Plagen, so über die Egypter [er]gangen, über die Israeliter nicht gekommen sind. Denn GOtt ist bei seinem Volke gegenwärtig, so haben sie auch GOttes Wort, welches das rechte Licht ist. An welchem Orte dasselbige gepredigt wird, da hat man das wahrhaftige Licht; wird aber das Wort GOttes verloren oder verachtet, so hat man denn allda greifliche Finsterniß. Darum sagt der 119. Psalm, V. 105.: Dein Wort sei meinen Füßen eine Leuchte. Und St. Petrus in seiner [zweiten] Epistel [Cap. 1, 19.] nennt das göttliche Wort auch ein Licht, das da leuchtet im Tunkeln.

11. Nun ist noch hinterstellig die letzte Plage. Wir haben aber bisher gehört, wie GOtt ein Spiegelfechten anrichtet, und läßt dem Pharaoni durch Mosen vorgeben, daß sein Volk drei Tage solle reisen,\footnote{Eißlebensche: rein sein.} umb ihren GOtt dienen. Er spricht nicht, daß sie aus dem Lande gar wegziehen sollen, und er sie in ein ander Land führen wolle, wie denn Pharao diesen Braten wohl riecht. Warum sagt es aber GOtt nicht ausdrücklich: Wir wollen kurzum davon, und uns trollen in ein ander Land? Moses spricht allein: Wir wollen in die Wüste ziehen, GOtt hat uns berufen zu einem Fest; und ist doch die Sache endlich, daß sie in ein ander Land ziehen wollen, und darf Moses alhie also bald mitlügen. Geht es nun hin, also ein Affenspiel vor die Nasen machen? Oder, hat er unserm HErrn GOtt nicht vertraut? Wie wenn er dürrer heraus sagte, sie wollten in ein ander Land ziehen, und hätte des Gottesdienstes geschwiegen? Wie denn alhier solches den Pharao hart verdrießt.

12. Da wisse, daß es nicht vonnöthen sei, daß man alleweg dem Teufel sage, was GOtt im Sinne und vorhabe. Es ist wahr, GOtt hat wollen das Volk lassen drei Tage in die Wüste ziehen; aber er sagt nicht, warum. Es ist auch nicht vonnöthen, daß er spreche: Ziehet drei Tage in die Wüste, und darnach fahret weiter, oder was GOtt weiter thun wolle. Darum so
\quellseite{419}
leugt allhier Moses nicht, und wenn er auch gleich gelogen hätte, so wäre Pharao eben recht geschehen. Gleichwie David [1 Sam. 21, 2.] auch log vor dem Priester Abimelech, als er die Schaubrode von ihm forderte, und vorgab, er reiste in Geschäften des Königes Saul. Denn Pharao betrügt und leugt dem Moses viel mehr; und wenn Moses gleich gelogen hätte, so wäre es ohne Gefahr gewesen, denn er thut Pharaoni keinen Schaden.

13. Es ist nicht eine Lüge, wenn sich einer anders stellt, denn es an ihm selbst ist; aber wenn einer anders vorgibt, denn es ist, und Betrug da ist, das ist Lügen. GOtt hat sich oft anders vorgestellet, und leugt dennoch nicht. Gleichwie auch die Eltern sich oft anders stellen gegen den Kindern; oder wenn ich zürne mit meinen Brüdern im Kloster, um ihrer Sünde willen, und habe doch Liebe im Herzen \&c. Also haben sich oft die Heiligen GOttes anders gestellt, denn sie im Sinne gehabt; aber sie haben keinen Schaden gethan, sondern nur das Ihre dadurch fördern wollen. Also ist es Moses nicht zu thun um den Gottesdienst, sondern daß er aus dem Lande käme, und Pharaoni geschieht daran kein Schade.

\begin{schriftzitat}
V. 27–29. Aber der HErr verstockte das Herz Pharao, daß er sie nicht lassen wollte. Und Pharao sprach zu ihm: Gehe von mir, und hüte dich, daß du nicht mehr vor meine Augen kommest. Denn welches Tages du vor meine Augen kommst, sollst du sterben. Mose antwortete: Wie du gesagt hast, ich will nicht mehr vor deine Augen kommen.
\end{schriftzitat}

14. Hiermit haben wir nun das Ende, und hat Moses ausgepredigt, und sein Wort vollzogen; der Befehl ist aus, so ihm von GOtt gegeben war, daß er dem Könige Pharao predigen sollte. Nun ist vonnöthen, daß GOtt daselbst wirke und helfe, wenn Menschen nicht mehr helfen können. Die Macht Pharaonis ist zu groß und stark, Moses muß nun sagen, er könne nicht mehr, sondern: Lieber GOtt, thue du es. Das ist eine häßliche, verdrießliche, unlustige Predigt gewesen, mit der es nicht soll von statten gehen; sondern es soll darum geschehen, daß GOtt Ursache habe, Pharaonen zu stürzen, und seine göttliche Ehre in der Welt zu offenbaren.

15. Also sollen alle Prediger auch das Ihre thun, und den eventum unserm HErrn GOtt befehlen. Moses gibt es und befiehlt es auch unserm HErrn GOtt in seine Hand, daß ihm Pharao verderbt, er soll nicht mehr vor ihn kommen, stößt ihn von sich, ja, er stößt GOtt selbst und sein Wort von sich. Nun laufen GOtt und Pharao auf einander. Moses, der ist schwach, aber GOtt wird nun stark werden.

16. Wenn der Welt lange genug gepredigt ist mit allem höchsten Fleiß und Treuen, und das Wort GOttes hört auf, dann ist es nicht weit dahin, daß die Welt von GOtt mit Strafen heimgesucht und visitirt wird, daß der Teufel die Oberhand bekommen und alles zu Grund und Boden gehen müsse. Denn da Noah aufhört zu predigen, und in die Arche ging, da kam die Sündfluth [1 Mos. 7, 10.]. Item, als Lot nicht mehr zu Sodom und Gomorra predigte, sondern von den Engeln ausgeführt ward nach Zoar, da kam Feuer, Schwefel und Pech, und vertilgte Sodom und Gomorra [1 Mos. 19, 24.]. Eben also geht es auch allhier in diesem Königreiche Egypten; da Moses Predigt ein Ende hat, und Pharao visitirt ist, da gnade denn GOtt, denn er gnaden soll. Denn bald darauf werden alle erste Geburt in Egypten erschlagen, und muß Pharao mit aller Mannschaft der Egypter im rothen Meere ersaufen.

17. Es ist dem jüdischen Volke auch also ergangen. Da die Propheten, und darnach Johannes der Täufer, Christus, der Sohn GOttes, und die lieben Apostel nicht mehr predigten und visitirten, da kamen die Assyrer, Babylonier, und hernach die Römer, und warfen’s alles über einen Haufen. Denn also sagt der HErr Christus [Matth. 10, 14. 15.]: „Wer euch nicht annimmt, und euer Wort nicht hört, aus deß Hause, und aus derselben Stadt gehet, und schlaget den Staub von euren Füßen. Wahrlich, ich sage euch, es wird der\footnote{Eißlebenische: dem.} Sodomer und Gomorrer Lande erträglicher gehen am Tage des Gerichtes, denn derselbigen Stadt.“

\quellseite{420}
\subsection{Erkl. 35, 182–185}

\subsection{Auslegungen über das zweite Buch Mosis}

\section{Das elfte Capitel}

\begin{schriftzitat}
V. 1–3. Und der HErr sprach zu Mose: Ich will noch eine Plage über Pharao und Egypten kommen lassen, darnach wird er euch lassen von hinnen, und wird nicht allein alles lassen, sondern euch auch von hinnen treiben. So sage nun vor dem Volke, daß ein jeglicher von seinem Nächsten, und eine jegliche von ihrer Nächsten, silberne und güldene Gefäße fordere; denn der HErr wird dem Volke Gnade geben vor den Egypten. Und Mose war ein sehr großer Mann in Egyptenlande vor den Knechten Pharao, und vor dem Volke.
\end{schriftzitat}

1. In diesem Capitel dräuet GOtt, daß er will noch eine Strafe über Pharao und die Egypter gehen lassen, welches soll die zehente und letzte Strafe sein, nämlich, daß alle erste Geburt solle erschlagen werden. Und spricht GOtt: er will es also damit machen, daß Pharao soll froh werden, daß er euch ausrichte, und das Volk ziehen lasse, ja, er werde sie mit Gewalt selbst treiben, daß sie sollen ausziehen. Moses kann nun nicht mehr. Denn bisher hat Pharao des Mosis Wort, Predigt und Dräuung verachtet, und nicht folgen wollen, so spricht GOtt: Ich will es selbst ausrichten, was du Moses nicht thun kannst; ich kann es wohl thun [Bi. 12, 6.]. Und ist nun Pharaoni genug gepredigt, und gehört dies Capitel zu dem vorigen, da Moses sagt: Es soll ja sein, ich will nicht mehr vor deine Augen kommen; aber das lasse ich dir zur Letze: Höre ich auf zu predigen, so wird ein anderer nach mir kommen, der wird es besser machen, der wird schlagen alle Erstgeburt, und alsdann wirst du mir Urlaub geben.

2. Alhier möchte man fragen, ob die Heiligen auch zürnen? Moses gibt hiermit dem Lande und dem Könige Pharaoni das Valete und den Urlaub, daß er nicht mehr für den König bitten wolle, sondern hielt ihn für einen verstockten und verhärteten Menschen. Wie denn St. Paulus [Tit. 3, 10. 11.] auch solches gebeut, daß man, nach Einer oder zweier Vermahnung, einen solchen soll fahren lassen, und ihn meiden, als den, der dem Teufel übergeben ist. Diesem lebt Mose auch hier gemäß, und spricht also: Ich will es auch lassen hingehen, und dir nicht mehr predigen, alsbald wird die zehente Plage kommen, daß alle Erstgeburt erschlagen werden soll, so wirst du froh werden, daß du mich und das Volk lässest gehen. Hat also ihm, seinem ganzen Reiche und Volke alle Freundschaft, Gebet und Förderung abgesagt.

3. GOtt gebeut aber Mosi und dem Volke Israel, daß sie nicht leer ausziehen sollen, sondern ein jeglicher von seinem Nachbarn güldene und silberne Gefäße borgen, dieselben viel und lange Jahre her hatten den Egyptern mit großer Mühe und Arbeit gedient, und ohne Lohn und Vergeltung, dazu ihrer Kinder viel waren getödtet worden; man hatte sie übel bedrängt und gezwängt.

4. Nun macht GOtt auch, daß sie Gnade finden vor den Egyptern. Sonst möchte jemand sagen: Wie ist es möglich, weil die Egypter die Hebräer also anfeindeten, daß sie ihnen ihr Silbergeschirr borgen und leihen sollten? Ja, spricht GOtt, ich will es wohl zuwege bringen; ich will den Egyptern Gunst geben gegen den Hebräern, daß sie es gerne thun.

5. Die Israeliten waren Gäste im Lande; wenn sie bei den Bürgern und Bürgerinnen baten um das Silbergeschirr, und sprachen: Morgen haben wir Fest, leihet uns eure silbernen Ketten, oder silbernen Schüsseln und Becher, so haben sie es gethan, und haben ihnen geglaubt. Etliche haben sich über sie erbarmt, und ist ihnen vielleicht der gemeine Mann auch günstig gewesen, ihr Elend und Noth gesehen. Sie werden nicht alle böse gewesen sein, als der König und der Landsfürst mit seinem Knechten. Denn es ist eine Synedoche, wenn man sagt: Das Volk ist entweder böse oder fromm, das ist, das halbe Theil. Denn das ist der Brauch und Art der heiligen Schrift, also zu reden.

6. So hat nun GOtt dem Volke Gnade gegeben, daß sie die Becher und Silbergeschirr der Egypter zu sich brachten auf das österliche Fest, und am Abend, da sie den Schmuck weg haben, ziehen sie hinaus und nehmen es alles mit sich. Was will man allhier sagen? Haben sie dieses nicht gestohlen, und mit Unrecht an sich gebracht,
\quellseite{421}
daß sie auf guten Glauben und Getrauen das Silbergeschirr den Egyptern abborgen, und entwenden es gleichwohl? Der Text spricht: GOtt hat sie es geheißen, das Silbergeschirr zu borgen. Denn GOtt hätte es wohl können machen, daß die Israeliten den Egyptern das Silbergeschirr mit Gewalt genommen hätten, aber er thut es nicht; denn er will sein Volk in der Demuth und unter dem Kreuz führen. Es läßt sich wohl ansehen und scheint, als haben sie es heimlich gethan, aber es ist doch gleichwohl GOttes Befehl gewesen.

7. Derhalben ist es keine Sünde. Denn was GOtt heißt, das ist recht. Wenn ich oder du es gethan hätten, so wäre es unrecht, aber dieweil er es heißt und gebeut, so ist es recht und wohlgethan. Denn der Egypter Güter sind GOttes. Er spricht: Ich habe es ihnen gegeben; wenn ich nun heiße, solches ihnen nehmen, so ist es recht; wer will mich urtheilen, daß ich daran Unrecht thue? Wenn es aber ohne mein Geheiß geschehen wäre, so wäre es unrecht. Aber ich thue es öffentlich, und kann es augenscheinlich alles wegnehmen; denn es ist alles mein, ich mache es damit, wie ich will [Ps. 50, 10.]. [Ps. 24, 1. ff. 1 Cor. 4, 7.]. Darum so ist allhier GOttes Wort und GOttes Gebot, und ist recht gethan.

\begin{schriftzitat}
V. 4. 5. Und Mose sprach: So sagt der HErr: Ich will zu Mitternacht ausgehen in Egyptenland, und alle Erstgeburt in Egyptenland soll sterben, vom dem ersten Sohne Pharao an, der auf seinem Stuhle sitzt, bis an den ersten Sohn der Magd, die hinter der Mühle ist, und alle Erstgeburt unter dem Vieh etc.
\end{schriftzitat}

8. Die letzte und zehute Plage ist diese, daß GOtt alle Erstgeburt schlagen läßt, vom Obersten an, bis auf der Thürmagd oder der geringsten Dienstmagd Sohn. Denn in Egypten hat es alles gedient, Mann und Weib hat man geheißen Knechte und Mägde, es sind arme Leute gewesen, die alle gedient haben. Da sagt Moses: Es soll in keinem Hause abgehen, es soll ein Todter drinnen liegen.

9. Ehe aber diese Plage hergeht, daß die Erstgeburt erschlagen worden ist, so beschreibt Moses im folgenden zwölften Capitel das österliche Fest, das GOtt eingesetzt hat den Israeliten.

\section{Das zwölfste Capitel}

\begin{schriftzitat}
V. 1–5. Der HErr aber sprach zu Mose und Aaron in Egyptenland: Dieser Mond soll bei euch der erste Mond sein, und von ihm sollt ihr die Monden des Jahres anheben. Saget der ganzen Gemeine Israel, und sprecht: Am zehnten Tage dieses Monds nehme ein jeglicher ein Schaf, wo ein Hausvater ist, je\footnote{Gießlebenische: ja.} ein Schaf zu einem Hause. Wo ihrer aber in Einem Hause zum Schaf zu wenig sind, so nehme er es, und sein nächster Nachbar an seinem Hause, bis ihrer so viel wird, daß sie das Schaf aufessen mögen. Ihr sollt aber ein solch Schaf nehmen, da kein Fehl an ist, ein Männlein, und Eines Jahres alt; von den Lämmern und Ziegen sollt ihr es nehmen.
\end{schriftzitat}

1. Wir haben durch die elf Capitel bisher gehandelt, wie Moses sein Amt ausrichtet vor dem Könige Pharaone, und sind mit der Auslegung gekommen bis auf die letzte Plage. Nun werden im 12. und 13. Capitel folgen etliche Gesetze, welche GOtt gibt, ehe denn er die zehnte Plage ausrichtet. Denn zur zehnten Strafe dräuet GOtt, daß er alle Erstgeburt, von des Königes Sohn bis auf der geringsten Magd Sohn, will todtschlagen. Ehe aber dieses in das Werk gesetzt wird, gibt er ein Gebot, wie die Kinder von Israel das Osterlamm essen sollen, wie sie den Abend zuvor, ehe denn sie des Morgens ausziehen, Fest halten sollen, das Silbergeschirr von den Nachbarn entlehnen, und in Egyptenland warten, bis daß die Erstgebornen getödtet wären; alsdann sollten sie ausziehen und Fest halten. Darum mußten sie nun, wie sie dieselbige Nacht ausziehen sollten, und gibt ihnen das Gebot vom Osterlamm. Denn weil sie das Osterlämmlein aßen, ward von GOtt
\quellseite{422}
alle Erstgeburt todgeschlagen. Und führte GOtt also die Kinder Israel aus Egypten im Monat Abib, das ist, im April.\footnote{Hier findet sich in der Eislebenschen Ausgabe die Bemerkung: Alhie hat D. Mart. Luther das 12. 13. 14. und 15. Capitel durchlesen auf der Kanzel.}

2. Also haben wir bisher gehört von neun Zeichen, so Moses wider Pharaonem gethan hat, und letztlich, wie GOtt auch die Erstgeburt erschlagen wolle. Und treibt allhier Moses zwei Lehren: erstlich, wie man das Osterfest halten solle; darnach, wie man die Erstgeburt opfern möge. Das Osterfest sollten sie darum halten, daß sie zur selbigen Zeit ausgezogen waren aus dem Königreiche Egypten; aber die Erstgeburt sollten sie darum opfern, daß GOtt die Erstgeburt der Egypter todgeschlagen hatte. Davon wollen wir zu seiner Zeit weiter hören.

3. Auf daß man aber nicht gedenke, gleich als hätte Moses allhier allen Menschen solch Gesetz gegeben, so müssen wir wissen, daß diese Historie uns zum Exempel der Nachfolge nicht gegeben ist; wir dürfen das Osterlamm nicht schlachten, noch die Erstgeburt opfern. Denn Moses ist nicht uns, sondern allein dem jüdischen Volke zum Lehrer gegeben. Daran ist viel gelegen, daß wir es wohl wissen; denn wenn wir Moses in allen Geboten halten sollten, so müßten wir auch beschnitten werden, und dürften nicht die Taufe annehmen, müßten auch das Osterlamm essen. Aber ihr habt gesehen und gehört, daß Moses nicht ist gegeben der ganzen Welt zum Lehrer, wie sonst der HErr Christus dazu gesandt ist, sondern Moses ist geordnet zum Propheten, Leiter und Heerführer des Volks, so da hieß Abrahä Same, oder so da Abrahams Fleisch und Geblüt war.

4. Denn droben [1 Mos. 12, 7. 26,3. 28,13.] habt ihr gehört, wie GOtt mit Abraham, Isaak und Jakob redet, und verhieß ihnen, daß er ihrem Samen wollte das gelobte Land geben. Diese Geschichte und Worte sind nicht über die ganze Welt gegangen, sondern nur geblieben und gefasset in der Schnur und dem Stamm Abrahams. Darum läßt sich Moses nicht weiter treiben, noch sich bestellen, denn nur auf dies Volk, als die Kinder Israel, denen sagt er, wie sie sich gegen GOtt halten, und auch auf Erden gegen den Menschen leben sollen, und thun solches alles aus GOttes Befehl und Geheiß, nicht länger denn bis auf die Zeit, die Abrahä verheißen [1 Mos. 22, 18.], daß in seinem Samen alle Geschlechter auf Erden sollten gesegnet werden, das ist, bis daß der HErr Christus komme.

5. Dieses sage ich darum, daß ich gerne wollte, daß ihr einen aufrichtigen, sichern, feinen, unterschiedlichen Verstand hätte zwischen Mose und dem HErrn Christo. Item, auch um der Juden willen, die da wollen, daß Moses durch die ganze Welt herrschen, und von jedermann seine Gesetze gehalten werden sollten, die doch den Text nicht recht ansehen, noch bedenken, daß Moses uns nicht ist zum Heiland gegeben, sondern er hat seine Kirche, und er ist Abt und Prior in seinem Kloster.

6. Aber der gebenedeyte Same, dem Abraham verheißen (da gesagt ward: „In deinem Samen sollen gesegnet werden alle Geschlechter auf Erden“), der ist ein anderer Mann gewesen denn Moses. Von dem wird gesagt: Und es wird ein Same gegeben werden, dadurch nicht allein dein Geschlecht, dein Fleisch und Blut, sondern alles, was auf Erden ist, soll gesegnet werden. Alle sollte er sagen: Dein Fleisch und Blut, und diesen Samen, will ich regieren und für mein Volk halten. Aber wenn das rechte Volk kommt, der wahrhaftige Same, da wird eine solche Predigt ausgehen, daß alle Heiden unter dem Himmel sollen gesegnet werden, daß er solle ein Vater aller gläubigen Heiden unter der Sonne sein.

7. Moses geht uns so viel an, daß wir nur seine Prophezeien und Exempel ansehen, als, wie Abraham geglaubt habe, und wie er die Verheißung vom Christo, dem gebenedeyten Samen, empfangen habe. Mehr kann er mir nicht dienen, und Moses kann auch nicht mehr, denn daß er uns Exempel des Glaubens vor die Augen halte, und daß er die Prophezeien von Christo beschrieben hat, und eingesetzt, wie der Mensch inwendig und auswendig leben solle, geistlich und weltlich sich halten möge. Darum folgt es nicht, daß man wollte vorgeben: Moses heißt, man solle sich lassen beschneiden; item, Abraham ist beschnitten worden; item, Noah machte einen Kasten, darein er ging vor der Sündfluth, darum so sollen wir es auch thun. So verbeyt solches auch GOttes Wort. Auf die Verheißung soll man sehen, die Moses führt, daß in Abrahams Samen sollten alle Geschlechter auf Erden gesegnet werden.

\quellseite{423}
8. GOtt hat einer jeglichen Zeit ihren Lehrer gegeben; aber zur selbigen Zeit, da diese Predigt ausging, daß alle Heiden sollten durch die Predigt des gebenedeiten Samens gesegnet werden, da wollte er einen Doctor schicken und geben, nicht Mosen, sondern Christum den HErrn selbst, der Abrahams Sohn und sein Same sein sollte. Da ist denn Jesaias, Jeremias, Moses und andere Propheten ausgestrichen; sie gelten nichts mehr, und soll Christus allein predigen; nicht irgend in einem Winkel, sondern in der ganzen Welt; denn er soll aller HErr sein, und soll nicht allein für sich predigen, sondern auch durch seine verordneten Diener.

9. Darum ist Moses allein in dem Stamme Abrahä gegeben worden; so weit Abrahams Fleisch und Blut reicht, da muß man Mosi gehorsam sein, und ihn hören, bis daß da komme die Verheißung, dem Patriarchen Abraham gegeben. Zu dem Stücke soll Moses sich brauchen lassen, derhalben ist er auch lange nach Abraham geboren; denn GOtt hat Abraham lange zuvor den Messiam verheißen [Gal. 3, 16.]. Abraham wußte von Mose nichts, die Verheißung war lange zuvor da, auf daß sich niemand rühmte, er wäre ein solcher, und Moses selbst mußte wohl bekennen, er wäre es nicht, wie er denn sagt im fünften Buch am 18. Capitel, V. 18.: „Einen andern Propheten wird euch GOtt erwecken aus euren Brüdern, denn wird er sein Wort in seinen Mund legen, den sollt ihr hören.“

10. Nachdem auch Moses ist gestorben, ist keiner kommen, der da wäre Mosi gleich gewesen [5 Mos. 34, 10.], denn der allein, so da von Abrahams Samen kommen sollte, würde viel höher und größer sein denn Moses. Derhalben ist Mosis Same dem Abrahä nicht verheißen; das kann man nicht leugnen. Denn Mosis Lehre ist nicht über die ganze Welt gegangen, sondern ist in ihrem Zirkel geblieben und gegangen. Ein anderer Same Abrahä sollte kommen, der sollte über die ganze Welt gehen, und alle Geschlechter segnen. Das bist du nicht, Moses. Danach, so soll auch seine Predigt anders klingen denn deine, lieber Moses. Das hat GOtt viel Jahre zuvor lassen weissagen und verkündigen, daß der Messias allen Landen Segen und Gnade predigen sollte, nämlich, daß durch Christum wir sollen selig werden; item, daß es durch Christum alles sei ausgerichtet, und uns GOttes Barmherzigkeit und Gnade allein durch Christum gepredigt werde.

11. Das erste und höchste ist ein Hauptspruch von Christo, und eine Predigt von ihm, die gilt mehr, denn sonst hundert Predigten oder Sprüche im Mose. Aber GOttes Wort wird nicht so wunderbärlich angesehen und geachtet, als es ist. Es ist ein einig kurz Wort, das von Christo gesagt wird: „In deinem Samen sollen gesegnet werden alle Geschlechter auf Erden“; das ist, alles, was da wird von Gnaden predigen, das wird von dem HErrn Christo kommen, und auf den Samen Abrahä zeigen. Moses hat ein ander Geschrei und anders gepredigt, und mit vielen Worten hat GOtt mit Mose geredet; aber es ist ein anderer HErr vorhanden gewesen denn Moses, welchen allein die Juden auf sich ziehen, aber diesen Lehrer, Abrahä verheißen, sollen wir Heiden auch auf uns ziehen [Joh. 10, 16. 17.]. Denn er predigt Segen, und nicht das Gesetz; denn er sollte das Gesetz, Tod, Teufel und alles hinwegnehmen. Dieser ist den Christen gegeben.

12. Wir haben in Mose viel Prophezeien von Christo, die wir daraus ziehen und nehmen sollen, auf daß wir unsern Glauben darauf wissen zu gründen. Der Lehrer, Moses, gibt Zeugniß von Christo, und er ist darum von GOtt geschickt, und soll auch um dieser Ursache willen fleißig gelesen werden. Darnach sind in Mosis Büchern seine Exempel, wie GOtt durch das Wort und im Glauben sein Volk geführt habe, und noch uns auch also regieren und führen wolle. Da ist denn Moses ein nützlicher, guter Lehrer.

13. Aber wenn einer Mosen anzeucht, und spricht: Er gehet die Beschneidung, darum so mußt du dich beschneiden lassen, so antworte du: O, Hans Grobianus bist du, Moses ist mein Magister und Lehrer nicht, er wird mich durch das Gesetz nicht lehren selig zu werden; ich habe viel einen bessern Lehrer, nämlich den HErrn Christum, der da [Matth. 11, 29.] spricht: Nimm von mir Gnade und Barmherzigkeit.

14. Moses sagt, als ein Gesetzprediger: Thue dies, thue jenes. Denn das ist des Gesetzes Predigt, die soll nun ein Ende haben; und wenn man uns Mosen mit seinem Gesetze wollte auf den Hals hängen, so sollen wir dem Mosi billig feind werden. Darum so nehmen wir Mosen [als] einen Gesetzgeber an, sonst halten wir uns
\quellseite{424}
an unsern lieben HErrn Christum Jesum. Moses ist aus, er hat mit seinem Amte genug gethan.

15. Danach, so kann ich des Mosis Amt dazu gebrauchen, daß er schöne, herrliche Exempel und Legenden von den geliebten Heiligen GOttes geschrieben hat, wie sie GOttes Wort haben gehört, Glauben und Liebe und alle Tugend gehabt und geübt; welches man alles fein in Mose findet. Aber das ist noch viel tröstlicher und lieblicher, daß er schöne Sprüche von dem HErrn Christo setzt, wie er sei unser Heiland.

16. Das sind zwo Ursachen, warum man Mosen fleißig lesen soll; als, daß man erstlich drinnen findet gar herrliche Exempel der Altväter und Patriarchen, die da GOttes Wort gehabt, auch ihren Glauben und Liebe haben leuchten und sehen lassen. Item, man sieht auch, wie GOtt habe die Bösen und Gottlosen gestraft, als Sodoma und Gomorra. Zum andern, daß er auch Weissagung und Prophezeien von Christo führt. In diesen beiden Stücken soll Moses mein Lehrer und Prediger sein.

17. Aber wie er die Juden regiert, daß sie dieses essen und jenes trinken sollen, item, was sie für Kleider tragen, wie sie freien sollen, das geht mich nichts an. Ich will von Mose zwei Stücke, als, Prophezeien von Christo und Exempel eines göttlichen Lebens haben, und nicht Gesetze oder Gebote. Und was da anlangt die Prophezeien und die Exempel, welche Moses beschreibt, da steht er noch; aber was er gebeut und heißt thun nach Art und Gewohnheit des Gesetzes, da ist es mit Mose aus und gilt nicht mehr, gleichwie es mit dem jüdischen Volke aus ist, und mit dem Priesterthume zur Zeit des HErrn Christi auch ein Ende hatte.

18. Denn es sollte alles nicht länger währen, denn bis der andere Lehrer käme, der da eine andere Lehre bringen würde; wie denn der andere Psalm, V. 7., davon sagt: „er solle von einer solchen Weise predigen: Du bist mein Sohn, heute hab ich dich gezeuget“ etc. Das soll derselbige Lehrer ausstreichen und im Maul führen, das sonst Moses hat in der Feder gehabt. So nehmen wir nun Mosen an als einen Zeugen von Christo, und halten uns nach seinen Exempeln, auf daß wir ein Vorbild und Spiegel eines christlichen Glaubens und Wandels draus nehmen, und geht uns Mosis als ein Gesetzgeber nichts an.

19. So wollen wir nun hören, wie man das Osterlamm essen solle, und was wir für Exempel darinnen haben. Wir hören, wie um diese Zeit GOtt den Kindern von Israel das osterliche Fest habe eingesetzt, und ihnen geboten, daß sie das Osterlämmlein leiblich essen sollten, und er sie dieselbe Nacht aus Egypten in die Wüste geführt, und aus dem Diensthause Egypti erlöst habe.

20. Ich habe aber gesagt, wie man Mosen lesen und verstehen solle, daß wir in solchen Stücken daraus Exempel der göttlichen Werke nehmen, wie er mit seinem Volke sei umgegangen, und daß wir kein Gesetz draus machen, als müßten wir jetzt auch das Osterlämmlein essen. Denn alhier fähet sich Mosis Amt an, daß er dies Volk der Kinder von Israel regieren sollte. Aber viel weiter geht das Wort, so Abraham verheißen war, daß in seinem Samen gesegnet werden sollten alle Geschlechter auf Erden; in welchem Worte auch wir begriffen sind. Und was von Christo in Mose gesagt wird, das geht uns auch an; des Segens, das ist, der Predigt des Evangelii müssen wir uns annehmen. Aber des Gesetzes Mosis sollen wir uns nicht weiter annehmen, denn daß es uns Exempel des Lebens vorstelle, und Stärkung unseres Glaubens gebe, daß, gleichwie GOtt mit ihnen ist umgegangen, also wolle er uns auch thun.

\begin{schriftzitat}
V. 2. Dieser Mond soll bei euch der erste Mond sein.
\end{schriftzitat}

21. Die Juden haben diesen seinen Brauch gehabt, daß sie das Jahr nach den Monden zählen; wir thun das Gegenteil. Den neuen Mond heißen sie, wenn das Licht angeht, wie die Bauern und der gemeine Mann davon redet im Kalender, wenn der Mond angeht, neu ist, oder wenn er zur Hälfte, oder gar voll ist. Und haben die Juden für den ersten Mond des Jahres gehalten den Abib, das ist bei uns der April, wenn alle Dinge anfangen neu zu werden; als, in der Fasten, wenn das Gras und die Bäume ausschlagen; und das heißen sie auf hebräisch den Neuenjahrsmond. Wiewohl nichts daran liegt, wo du ansähest. Aber sie waren also im Gesetze verbunden, daß sie mit dem Monden mußten das Jahr anfangen. So war nun ihr erster Mond und Jahrstag mit dem März, oder wenn unser April anfähet, darnach der neue Mond anstehet, um die Zeit, wenn es grün wird, um die Fastenzeit.

\quellseite{425}
\begin{schriftzitat}
V. 3. Saget der ganzen Gemeinde Israel, und sprechet: Am zehnten Tage dieses Monden nehme ein jeglicher ein Schaf, wo ein Hausvater ist, je ein Schaf in Einem Hause.
\end{schriftzitat}

22. Zehn Tage sollen sie zählen von dem neuen Mond, das ist, da das erste Viertheil [des] Monds weg war; im selben Monden sollte ein jeglicher Hauswirth ein Schaf zu Einem Hause nehmen \&c. Die Juden haben das Haus voller gehabt denn wir; die Männer haben viel Weiber genommen, und die Weiber denn viel Mägde gebraucht, und die Mägde waren auch ehelich, und waren der Mägde Kinder bei der Frauen Kinder bei einander, also, daß ein Haus gleich wie ein eigen klein Königreich, oder wie eine Stadt gewesen, darinnen es krimmelte und wimmelte von Kindern und Menschen; wie sie noch thun, wo sie wohnen.

23. Das Schaf oder Lamm sollte keinen Wandel haben, nicht, daß es nur weiß wäre, sondern daß es nicht verwundet, nicht schäbig und sonst mangelhaftig wäre, noch ein Bein gebrochen hätte, sondern ganz heil und gesund; welches denn gar wohl also verdeutlicht ist. So soll es auch ein Männlein sein. Der lateinische Text hat allhier sehr geirrt, daß sie auch ein Responsorium davon gemacht haben.

24. So hat auch nicht die ganze Gemeinde ein Schaf opfern sollen, sondern, so weit Israel wohnte, so sollte ein jeglicher Hauswirth sein Schaf haben, und da er es allein nicht bezwingen konnte neben dem Hausgesinde, mußte er seinen Nachbar dazu nehmen, auf daß dies Osterlämmelein denselbigen Abend aufgegessen würde.

25. Und sollten es von der Heerde nehmen. Wo sollte man es sonst nehmen? Daran sind sie hart gebunden gewesen, daß sie im Hause kein Schaf zum Osterlämmelein sonderlich aufziehen haben dürfen, sondern zur Heerde gehen, und zu seiner Zeit das Lämmelein heraus heben. Am zehnten Tage sonderte man es also ab, und stellte es in das Haus, und behielten es daselbst vier Tage. So spricht nun GOtt: Ein jedermann solle sein Schaf schlachten und opfern, welches denn einerlei Ding ist. Denn alles, was sie opferten, das schlachteten sie. Darum braucht man in hebräischer Sprache immerdar dieser Worte eines für das andere. Als, zum Römern am zwölften Capitel, R. 1., spricht Paulus: Ihr sollet euren Leib opfern, das ist, schlachten, tödten. Also ward Isaak geopfert, das ist, es galt ihm das Leben, er sollte den Hals herhalten, getödtet, geschlachtet und geopfert werden.

26. Also hat GOtt das Fest bestimmt, daß es sollte angehen am Abend des vierzehnten Tages, auf daß am fünfzehnten Tage des ersten Monds der rechte Ostertag wäre. Das ist dies Gesetz, daher ist es ohne Zweifel geschehen, daß das wahre, rechtschaffene, liebliche und gnadenreiche Osterlämmlein, unser HErr Jesus Christus, am Palmtage ist zu Jerusalem eingeritten, hat vier Tage im Tempel gepredigt, darnach das Osterlämmlein mit seinen lieben Jüngern gegessen; ist bald darnach im Garten am Oelberge gefangen worden, verspottet, verspeiet, gegeißelt, mit Dornen gekrönt, zum Tode verurtheilt, gekreuziget und getödtet, und also recht geopfert worden. Das ist nicht also zugegangen, daß man ihn hätte auf einen Altar gelegt und darauf geschlachtet, sondern er ist auf dem Kreuze geopfert und daran gestorben.

27. So mußte nun ein jeder in seinem Hause ein Lämmelein schlachten (wie man sonst pflegt einem Lämmelein den Hals abzustrecken), und ließen es braten, wie ihr denn hören werdet.

28. Dies Osterlämmlein essen geht nun auf den HErrn Christum, mit dem trifft es alles so lustig überein, daß es Wunder ist. Sonst geht es uns nicht viel an, wie die Juden haben das Osterlämmlein gegessen. Sie mußten mit des Lämmleins Blut die Pfosten bestreichen, dazu sie ein Büschelein Ysops gebrauchten, wie der Prophet David solches berührt im 51. Psalm, R. 9., da er spricht: Besprenge mich, HErr, mit Ysopo; und mußte ein jegliches Haus also beschmiert sein am vierzehnten Tage desselbigen Monats. Darnach mußte man bei dem gebratenen Lämmlein Oblaten essen, und bei dem Gebrateniß Salßen oder bittere Lactuken haben. Item, sie mußten gegürtet stehen, als die, so wegfertig wären, die sich trollen und die Nacht davon ziehen wollten, als die alle Stunden bereit wären, aus dem Lande zu gehen, oder die da auf sein wollten, und davonreisen.

29. Er wollte nicht ein solch Mahl anrichten, da man sich toll und voll fressen und saufen sollte, sondern dabei man wacker und frisch auf wäre [Röm. 13,13.]. Er wollte nicht ein köstlich Mahl anrichten, sondern es sollte ein Zeichen und Bedeutung sein, damit GOtt sie üben und behalten wollte im wahrhaftigen Gottesdienste,
\quellseite{426}
auf daß sie nicht den heidnischen Abgöttern dieneten. Denn es soll des HErrn Gang (Passah) heißen, das Osterlamm oder der Ostertag, als man im Evangelisten Luca [Cap. 22, 1. f.] findet, der Tag, daran man das Osterlamm opferte, das ist, das Lamm gegessen hat, das man mußte auf Ostern essen. Der vierzehnte Tag ward genannt der erste Tag der süßen Brode.

30. Aber warum wird das Osterlamm ein Gang genannt, oder ein Uebergang, und das Fest auch ein Gang geheißen? Der Name ist daher genommen, daß in dieser Nacht GOtt durch ganz Egypten gegangen ist und alle Erstgeburt todtgeschlagen hat. Den Mord hat GOtt angerichtet in derselbigen Nacht.

31. Wohl an, das sind die Ceremonien, Weise,\footnote{„Weise“ fehlt in der Erlanger.} Sitten und Gebehrde, die man gehalten hat über dem Osterlamm, daß sie dieselbige Nacht mußten zu Hause bleiben und die Thüren zuthun, und sich nicht haben lassen merken, daß sie daheim wären; mußten in einem Hause essen, und nichts überbleiben lassen, alles aufessen oder verbrennen. Item, kein Bein an ihm brechen. Kein Fremder, oder Gekaufter, noch Hausgenoß und Miethling, mußte nicht davon essen. Das war die Weise der Juden, also mußten sie sich schicken, das Osterlammlein zu essen. Wenn ein Gast in der Herberge lag, oder ein fremder Landsaß, der da nicht einheimisch war, so durfte er nicht mitessen, sondern allein der Wirth, sein Weib, Söhne, Töchter, Knechte und Mägdge, denn GOtt hatte es seinem Volke gegeben. Wer aber ein gekaufter Knecht war, der mußte beschnitten werden, und alsdann mochte er auch vom Osterlamme essen.

32. Also war es beschrieben und gefasset, daß sie das Osterlämmlein daheim essen sollten, und nichts ließen überbleiben, nichts daran zerbrechen mußten. Mit diesem ihrem Gesetze waren nur die Juden beladen, das geht uns nichts an; es sollte auch nicht ewig bleiben, sondern nur stehen, bis daß Christus käme.

33. Das soll eine ewige Weise sein. Ewig nennen die Juden nicht alles, das kein Ende hat, sondern ewig heißen sie, das für und für währet.

34. Den vierzehnten Tag auf den Abend mußten sie anfangen zu essen süße Brod, und essen bis auf den ein und zwanzigsten Tag desselbigen Monds. Also sehen wir, daß das Osterfest acht Tage gewährt, und eine ganze Woche gestanden habe. Aber es hat nicht allezeit am Sabbath angefangen, sondern unterweilen am Montage, zuweilen auch am Mittwochen, nachdem der Mond angegangen und gelaufen; gleichwie noch unser Chrishtag oft so fällt. Auch wenn wir den Ostertag halten, so zählen wir nicht also von dem vierzehnten Tag an, denn er oft über acht Tage hernach gehalten wird, wie denn auch die Pfingsten. Aber es ist uns nicht geboten, und wir sind darum nicht Heiden, ob wir es gleich nicht also halten. Denn die Christen haben allezeit Ostertag, ihr Osterlamm währet für und für.

35. Dennoch hat man es in der christlichen Kirche behalten, wie es die Juden gelassen haben, gleichwie man andere Bräuche mehr behalten, als, daß die Weiber noch sechs Wochen haben nach der Kindesgeburt. Jedoch hält man es so ferne, daß man nicht uns lasse ein Gebot daraus machen, oder daß es unser Gewissen verbinde, oder daß wir es thun müßten. Denn wir halten die Zeit nicht wie die Juden, daß wir nach dem Mond rechneten; sondern, wenn er nur kommt der Ostertag, so feiern wir ihn. Jedoch halten wir ihn allezeit auf einen Sonntag, es sei am Monden wie es wolle, er sei im Anfange, Mittel oder Ende, darob uns denn die Juden schelten. Aber es liegt nichts daran. Am Abend fähet allezeit der Juden Fest an, wie 1 Mos. 1, 5. geschrieben steht: „Da ward aus Abend und Morgen der erste Tag“, und fortan. Wenn der Tag bei ihnen ein Ende hatte, so ging der folgende Tag an. Wie denn auch bei uns die folgende Nacht ist der Anfang des nachkommenden Tages. Die Nacht mit dem völligen Tage ward für Einen Tag gerechnet, und dieses nennen die Astronomi oder (wie man sie gemeiniglich heißt) Sternkinder noch diem naturalem. Das ist der Juden Fest und ihr Brauch gewesen.

36. Diese Historie sollen wir wohl fassen, auf daß wir sehen, was GOtt damals habe im Sinne gehabt. Dies Gesetz wird ihnen von GOtt geboten, gleichwie der Befehl von den ungegärten Broden, welcher auch gegeben ward von unserm HErrn GOtt. Darnach ist Moises zugefahren, und hat es ihnen vorgebracht.

37. Nachdem sie nun das Fest angefangen und das Osterlamm gegessen haben, da ist GOtt mitten in der Nacht gekommen, da der Schlaf
\quellseite{427}
am süßesten und besten gewesen, und hat alle Erstgeburt in Egypten erschlagen.

38. Das ist eine greuliche Schlacht gewesen, so ein jämmerlich Ansehen gehabt, daß GOtt mitten in der Nacht also plötzlich zufährt, und die Erstgeburt tödtet; da denn ein groß Geschrei und Weheklagen sich erhebt, daß man eins hier, das andere dort todt findet. Und man wird wohl in Einem Hause viel erstgeborne Söhne todt gefunden haben, als da viel Hausgenossen innen gewesen sind; da werden sie gedacht haben, daß das Land gar voller Teufel sei. Nun, GOtt hat damals also gethan, er habe dazu gebraucht gleich eines guten Engels, oder eines bösen Teufels. Gleich als wenn bei uns ein Feuer aufgeht, da erschrickt jedermann; aber was sollte für ein Schrecken werden, wenn eine solche Plage über ein ganzes Land kommt, wie alle Erstgeburt alhier erschlagen wird?

39. Aber GOtt hat Noth halben den Pharaonem also heimgesucht, daß er in der Nacht aufsteht, und die Kinder Israel aus dem Lande treibt. Vor drei Tagen wollte er sie nicht ausziehen lassen; nun hätte er Geld und Gut darum gegeben, daß er ihrer los wäre. Zuvor hat er sie gehalten, nun stößt er sie selbst aus dem Lande; er will sie nicht bei ihm bleiben lassen, sie müssen des Nachts davon.

40. Das ist die letzte Plage gewesen, dadurch GOtt das Volk Israel hat erlösen wollen, und es ist auch die grausamste Plage gewesen; noch hat sie nichts geholfen, daß Pharao oder die Seinen sich erkannt hätten. So kommt nun darauf das vierzehnte Capitel, die Letze mit Pharao und der Garaus, daß er gar verstocket bleibt, und verfolgt die Kinder von Israel, wird darüber in das rothe Meer gestürzt.

41. Aber wie demüthigt sich der Schalk? Hüte dich vor solchen, denn es kommt oft, daß die Gottlosen sich so geistlich stellen, und rühmen sich, als sind sie frömmer als keine Heiligen. Also, Pharao bittet jetzt, daß sie ihm gnädig seien, und weichen aus seinem Reiche, welche er zuvor nicht konnte sehen, noch wollte leiden. Denn die Egypter sprechen: Wir sind alle des Todes, der Tod ist in dem ganzen Lande. Diese Plage ist nicht Finsterniß, noch Geschwöhr, Blattern und Drüsen, sondern also heftig und herbe, daß sie sagen: Wir müssen alle sterben. Es ist diese Strafe nicht wohl auszudenken, sie wissen nicht, wo sie daheim gewesen sind, darum sagen sie: Je ehe je besser aus dem Lande geschafft.

42. Also sind die Israeliten eilends ausgetrieben, daß sie auch den Teig, davon sie wollten Brod backen, in Tüchern mit aus Egypten wegtragen. Sie hatten erst angefangen ungesäuert Brod zu essen, und hatten Teig eingeknetet, um eine halbe Nacht, konnten nicht fertig werden; da trieben sie die Egypter aus, daß sie also Wasser und Mehl nehmen und einen Teig kneten, und in Kleider, Mäntel und Böcke den ungesäuerten, rohen Teig fassen mußten. Und werden die Egypter darüber also bestürzt, daß sie den geliehenen Schmuck von den Juden nicht wieder fordern, werden darum gebraucht.

43. Denn, wie gesagt, GOtt gab den Juden diese Gnade vor den Egyptern, daß sie güldene und silberne Trinkgeschirre und Kleinode ihnen gerne borgten und liehen, und entwandten also ihren Wirthen ihr Gut. Daran ist recht geschehen, denn sie hatten zuvor den Kindern von Israel ihre Söhne erwürgt, und sie mit schweren Frohndiensten beladen, und kein Geld dafür gegeben; nun gibt ihnen GOtt jetzt den Lohn und den schönen Schatz.

44. Also ziehen sechsmal hunderttausend Mann aus Egypten, ohne was da ist der Pöbel gewesen; item, die Weiber und Kinder. Egypten ist reich und groß gewesen, darum haben die Juden von ihnen viel Güter entlehnt, ziehen hin und wieder auf dem Felde. Wo haben sie doch ihre Küchen gehabt? Also haben wir das Fest, wie es ist eingesetzt und gehalten worden.

45. Dieser Auszug der Kinder von Israel aus Egypten wird in der heiligen Schrift sehr gerühmt. Derwegen sollen wir diese Historie wohl lernen und wissen, um der folgenden Texte willen. Denn GOtt rückt dem jüdischen Volke diese Wohlthat für und für auf: daß sie ihn nicht dabei erkenneten und ehreten. Darum wir auch den GOtt anrufen, ehren und ihm dienen sollten, der seinen Sohn für uns Menschen am Kreuze hat sterben lassen, und geboren werden, der uns auch geboten hat zu beten. Auch richtet er das Priestertum auf, daß er dem jüdischen Volke diese Wohlthat einbildete, und täglich davon predigte. Er gab ihnen dasselbige äußerliche Werk, auf daß er es dabei ergriffen, erkannt und gefaßt werden möchte.

46. Denn GOtt hat wohl gewußt, daß Rottengeister kommen würden, und der eigenen Andacht
\quellseite{428}
so viel werden, daß keiner den andern ansehen würde, und daß einer würde diesen Gott haben wollen, dem eine graue Kappe gefiele; jener einen andern Gott suchen, in einer schwarzen Kappe, so doch Gott durch seine eigene Andacht wollte gesucht werden, sondern er bestimmt uns ein Werk vor, dabei wir ihn sollten kennen lernen. Darum, wenn du sprichst: Ich will GOtt dienen, dem ich gelobt, so viel Tage zu fasten etc., hörst du, GOtt will es nicht haben, denn es ist ein Werk, das du dir selbst auferwählst und bestimmest. Aber Gott spricht: Ich will dir ein Werk bestimmen, bei dem Werke sollst du mich erkennen, nämlich, du sollst ehren den GOtt, der das Volk Israel [hat] aus Egypten geführt.

47. Wiewohl hernachher die falschen Propheten Altäre aufgerichtet haben unter diesem Titel; denn sie haben wohl gewußt, daß man diese Zeichen und Werke müßte treiben; aber sie verführten dadurch das Volk. Gleichwie auch wir, die wir Christum haben, predigen ihn, daß er gekreuzigt worden, und unser Heiland sei. Aber die falschen Lehrer, Pfaffen und Mönche sagen darnach: Siehe, alhier im Kloster, in der Kappe, Platten und Stricke, da ist auch Christus. Also nehmen sie diesen Titel und Werk, und ziehen es auf sich. Aber wären sie rechte Propheten gewesen, so sollten sie das gepredigt und getrieben haben, daß die Leute den GOtt ehren sollten, der da nicht zu BethEl noch BethAlven ist, sondern der die Kinder von Israel hätte aus Egyptenland geführt. Also auch jetziger Zeit, wenn ich sagen wollte, daß ich das Werk dazu setze, als, ich diene dem GOtt, der meine Jungfrauschaft, Keuschheit, Armuth und Gehorsam ansieht, da ist es falsch und unrecht.

48. Sie sind nicht vierhundert und dreißig Jahre allzumal in Egypten gewesen, sondern es ist von der Zeit an zu rechnen, da Abraham, 1 Mos. 12, 1., Gottes Stimme hört; von der Verheißung an, da Abraham fünf und siebenzig Jahre alt ist gewesen, bis auf die Zeit, da Moses achtzig Jahre alt ist. Also, wenn man diese Jahre zusammen zeucht, so ist es wahr; sonst sind sie nur die Hälfte dieser Jahre oder ein wenig drüber in Egypten gewesen. Unter Pharaonis Tyrannei sind sie nicht über hundert Jahre gewesen. Also besteht die Schrift allenthalben mit Ehren. Auf diese Meinung redet auch St. Stephanus in den Geschichten der Apostel, da er [Cap. 7, 6.] spricht, daß sie vierhundert und dreißig Jahre sind in Egypten gewesen.

49. Und beschleußt auch also St. Paulus in der Epistel zum Galatern [Cap. 3, 11.], da er lehrt, Mosis Gesetz mache niemand gerecht. Das beweist er also: Denn Abraham ist verheißen viel eher zuvor, denn Moses gekommen ist, daß durch seinen Samen alle Geschlechter auf Erden sollten gesegnet werden, alle Heiden und Welt sollte von ihm Segen und alle Güter kriegen; daselbst ist das Gesetz noch nicht gewesen. Darum so kann es nicht das Gesetz thun, das Gesetz hat bei sich nicht Hülfe, einen fromm zu machen etc., sondern aus dem Evangelio bekommen wir solches.

50. Darum sollen wir es also zusammentragen und vergleichen, daß es nicht von der Wohnung geredet ist, die sie in Egypten haben gehabt, sondern von der ganzen Zeit an, da dem Abraham ist der verheißene Same zugesagt worden, und darnach zu ihm gesprochen, er sollte in ein ander Land ziehen. Es ist eine Synedoche, gleichwie Christus drei Tage und drei Nächte ist im Grabe gelegen, da er nur Einen Tag und zwo Nächte ist drinnen gewesen; aber er hat ergriffen und gefasset drei Tage. Den ersten, auf welchen er gekreuzigt, desselben Stück und die folgende Nacht zählt man für Einen Tag; darnach den Sonnabend (wie wir es nennen); Ausgangs\footnote{Gießebensche und Erlanger: Ausgangs.} der nachgehenden Nacht, am Morgen des dritten Tages, ist er erstanden; dann ist es Ein ganzer Tag, und zwo Nächte mit einem Stücke des Tages. Also nimmt er alhier auch totum pro parte, und partem pro toto, denn sie eine solche Zeit, und so viel Jahre haben ergriffen.

\subsection{Allegoria oder geistliche Deutung dieses zwölften Capitels}

51. Im zwölften Capitel haben wir gehört, wie die Kinder von Israel das Osterlamm gegessen haben, und Gott dasselbige Zeit der Ostern eingesetzt hat, daß man damit erhalten sollte das Gedächtnis der wunderbaren Erlösung der Kinder Israel aus Egypten; wie man es aber unter den Christen auch sollte essen, das ist daneben auch angezeigt. Wir werden aber durchaus gelehrt, daß wir an Gott glauben
\quellseite{429}
sollen, und daß unser Glaube also gestalt sein solle, daß wir wissen, GOtt sei in den höchsten und größten Nöthen nahe bei uns.

52. Nun haben wir noch vor uns die geistliche Auslegung und heimliche Deutung dieses Capitels und erzählter Geschichte zu handeln. Ich habe aber zuvor gesagt, und sage es noch allezeit, daß wer da suchen und forschen will den heimlichen Verstand oder heimliche Deutung, welche unter einer Geschichte verborgen steckt, daß er sie also führe, daß sie sich reime mit dem Glauben; denn die heilige Schrift ist sonst hell und klar, und darf keiner seiner Träume hieher tragen. Es hat solches im alten Testamente etwas bedeutet in dem Regimente, so der HErr Christus führt; wie es denn mit hellen, ausdrücklichen Worten und Thaten bezeugt wird, daß dies auf Christum geht; denn es ist alles um den Mann zu thun. Aber wir wollen so viel davon handeln, als wir können.

53. Erstlich so sollte ein jeglicher Hausvater ein Lämmlein nehmen für sein Haus, und so viel Menschen dazu ziehen, als es genug war, daß sie es ganz aufessen könnten, sonst sollten sie ihre Nachbarn auch dazu berufen. Dieses Osterlamm hat St. Paulus 1 Cor. 5, 7. 8. deutlich beschrieben und abgemalt, und herrlich herausgestrichen, was es bedeutet, da er spricht: „Wir haben auch ein Osterlamm, das ist Christus, für uns geopfert; darum lasset uns Ostern halten, nicht im alten Sauerteige“ etc. Als sollte St. Paulus sagen: Es ist nicht der Juden Osterlamm, die ein äußerlich und leiblich Lämmlein genießen haben, sondern ein geistliches, und das nicht gesehen wird. Denn ob er wohl vorzeiten von den Aposteln ist gesehen worden, die mit Christo gegessen und getrunken haben, sind mit ihm umgegangen, so wird er doch jetzt nicht gesehen, sondern er ist gen Himmel aufgefahren, und sitzt jetzt zur rechten Hand seines himmlischen Vaters [Joh. 20, 17.]. Darum so essen wir nun den HErrn Christum, unser rechtes Osterlamm, im neuen Testamente. Wie denn auch St. Johannes der Täufer sagt [Joh. 1, 29.]: „Siehe, das ist GOttes Lamm, so da trägt die Sünde der Welt.“

54. Dieses ist eine gewisse Deutung, daß man aus Christo ein Osterlamm mache, und auf ihn führe die heimliche Auslegung. Solches kann nicht fehlen, so wenig Christus fehlen kann, sonderlich wenn dieser Spruch Johannis des Täufers dazu kommt; wie man denn allhier St. Pauli Zeugniß auch hat, daß er das Osterlamm auf den Herrn Christum deutet und von ihm auslegt.

55. So ist nun das erste, daß man das Schaf nehmen soll von der Heerde. Es sind aber viel Osterlämmer gewesen, möchte derhalben einer sagen: Haben wir denn auch so viel Christos, als viel da Christen sind, gleichwie man allhier so viel Osterlämmer hat haben müssen, als viel Häuser gewesen sind? Da wisse, daß dort ein leiblich, äußerlich Wesen ist, da sie nicht alle in einem Hause wohnen konnten; aber dieselbige große Menge damals ist nun in ein einiges Wesen gezogen, daß, wiewohl sie jetzt an Einem Orte nicht alle sind, so müssen wir doch alle zumal nur Ein Osterlamm haben. Das ist das erste, daß Christus nicht allein soll bleiben unter den Juden, sondern unter den Heiden auch angenommen werden.

56. Das andere ist, daß man auch die nächsten Nachbarn zum Essen des Osterlamms ziehen und nehmen solle etc. Damit will er haben (wie es auch St. Paulus zum Römern lehrt), daß Christus desto mehr ausgebreitet werde. Denn es ist nicht genug dran, daß diese oder jene Christum erkennen, sondern wir sollen es ausbreiten und jedermann verkündigen, auf daß viel Leute zu diesem einigen Hause kommen möchten, ja, die ganze Welt zum Reiche Christi gebracht würde. Vielleicht hat man damals viel Häuser gehabt, darinnen so viel Personen gefunden sind worden, als da mußten bei dem Osterlämmlein sein, daß man es ganz und gar aufessen konnte. Aber allhier, im neuen Testament, fehlt's allezeit, da sind immermehr Christen genug [Röm. 12, 2.], daß wir aufhören möchten, Gäste zu diesem Osterlamm zu laden, sondern wir sollen immerdar weiter gehen und predigen, und uns auch zu denen finden, welchen Christus zuvor nicht ist gepredigt worden, und sie lehren, die da Christum nicht erkannt haben, daß sie zu dem geistlichen Reiche Christi auch gebracht werden. Ihr Osterlämmleinnehmen war leiblich, aber unseres ist geistlich. Und GOtt hat dasselbige Osterlämmlein genommen, wir haben es nicht erwählt; dieses Osterlämmleinnehmen geht geistlich zu.

57. Zum dritten so soll man ein solch Schaf nehmen, da kein Fehl an ist, als, daß da nicht hinkt, einäugig, schäbigicht oder räudig wäre. Das ist, man soll Christum also erkennen und essen,
\quellseite{430}
daß man ihn allein habe und halte als den, der ohne Fehl und gesund sei. Denn man müßte das Osterlamm rein erziehen, das ist, reine Lehre von Christo haben, wie es vor GOtt gilt. Ein Lämmlein, das vor den Leuten keinen Fehl hat, da fragt GOtt nicht nach, sondern darum ist es zu thun, daß wir Christum ganz rein erkennen, als das unschuldige, unbefleckte Lämmlein, das da keine Sünde gethan hat, wie St. Petrus in seiner [ersten] Epistel [Cap. 1, 19.] und auch die zum Hebräern davon redet [Cap. 7, 26.], daß wir einen solchen Bischof haben, der da unschuldig und heilig ist. Denn sonst ist kein Mensch auf Erden, der diesen Ruhm und Titel führen könnte, daß er nicht Sünde habe [Sprüchw. 20, 6.], ausgenommen der HErr Christus, der hat keinen Fehl.

58. Das ist nun unser christlicher Glaube, daß wir wissen und bekennen, alle Menschen haben Fehl, sind befleckt, schuldig und gebrechlich, auch die Christen, sie haben alle Sünde; allein der HErr Christus ist vollkommen, rein, heilig und gerecht. Darum so ist auch das die Ursache, daß wir ihn allezeit essen müssen. Denn diejenigen, so da Christum nicht kennen, noch ihm anhangen, die sind gänzlich todt und sind nichts. Aber die ihn erkennen und annehmen, das sind rechte, lebendige Leute [Joh. 6, 53. 54.]. Und dieselbigen sind etwas; aber dennoch sind sie nicht ganz und gar rein; sie sind wohl Christen, aber darum nicht der HErr Christus selbst. Es ist wohl etwas, aber doch sind sie allhier in diesem Leben noch nicht gar rein. Derhalben soll ihm diesen Titel niemand selbst zueignen oder zuschreiben, daß er gar ohne Fehl sei, auf daß er Ursache habe, Christum allezeit zu essen, dieweil er lebt.

59. Zum vierten so sollen sie ein Männlein nehmen. Das kann dahin gezogen werden, daß Christus ein Mann ist gewesen. Aber es ist noch nicht genug, denn es geht nicht auf den Glauben. Denn ein Weib ist ebensowohl ein Mensch als der Mann. Aber in der heiligen Schrift [1 Mos. 3, 17. 1 Cor. 11, 3.] steht also, daß der Mann gesetzt sei zum Haupte des Weibes, und daß das Weib genieße aller Güter und Ehre des Mannes. Denn Mann und Weib Ein Leib ist, daß eines des andern sich soll annehmen, und was Einem geschieht, das solle dem andern auch geschehen, denn es ist Ein Ding. Also soll Christus nicht für sich eine Person oder Mann sein, ohne Flecken und Fehl, voller Weisheit und Gerechtigkeit, unschuldig und heilig, das mich denn nicht hilffe; denn er, und nicht ich, hätte daran genug.

60. Also haben ihn die Sophisten gemalt, wie er Mensch und GOtt sei, zählen seine Beine und Arme, mischen seine beiden Naturen wunderlich in einander; welches denn nur eine sophistishe Erkenntniß des HErrn Christi ist. Denn Christus ist nicht darum Christus genannt, daß er zwo Naturen hat. Was geht mich dasselbige an? Sondern er trägt diesen herrlichen und tröstlichen Namen von dem Amte und Werke, so er auf sich genommen hat; dasselbige gibt ihm den Namen. Daß er von Natur Mensch und GOtt ist, das hat er für sich; aber daß er sein Amt dahin gemeinet und seine Liebe ausgeschüttet, und mein Heiland und Erlöser wird, das geschieht mir zu Troste und zugute; es gilt mir, darum, daß er sein Volk von Sünden los machen will. Matthäi am 1. Capitel, V. 21., wird angezeigt vom Engel Gabriel, daß er soll Jesus heißen; nicht darum, daß er GOtt und Mensch ist, sondern, daß er das Amt soll führen und in das Werk treten, den Leuten von Sünden und Tode zu helfen. Das macht ihn zu einem Manne. Dafür sollen wir ihn auch halten, daß er das Haupt und Erbe des Christenthums und aller Gottseligkeit ist. Der HErr Christus ist der Christenheit Mann, sie ist sein Weib; wie St. Paulus zum Ephesern am fünften Capitel, V. 25–27., solches auch meisterlich ausstreicht, und spricht: „Christus habe seine Gemeine geliebet, und sich selbst für sie gegeben, auf daß er sie heiligte, und hat sie gereiniget durch das Wasserbad im Wort, auf daß er ihm darstellte eine herrliche Gemeine, die nicht habe einen Flecken oder Runzel, oder deß etwas, sondern, daß sie heilig und unsträflich sei.“

61. Also, wenn ein Mann sein Weib liebt, so ist sie ihm die Schönste und Liebste, wenn sie allein fromm und ehrlich ist. Denn die Liebe theilt ihr mit alle Ehre, und gibt ihr dahin Leib, Gut und alles, den Namen und Titel, daß sie habe, was der Mann hat. Sie sitzt mit ihm in gemeinen Gütern und Ehre. Und wie ein Weib zu ihrem Manne kann sagen, also spricht auch ein Christ zu dem HErrn Christo. Denn das Weib hat alles Macht im Hause, allein, daß der Mann das Haupt und der Herr sei oder Regent; dennoch nimmt sich das Weib alles an, und
\quellseite{431}
greift zum Gute, als wäre es ihr eigen. Also stellt sich ein Christ auch gegen dem HErrn Christo, allein, daß er, der HErr Christus, das Haupt, der Mann und HErr ist, von dem man alles Gut und Ehre hat, welches allein der Unterschied ist. Denn Christi Unschuld, Leben und Gerechtigkeit ist mein, so ist das Himmelreich und der Heilige Geist auch mein, auch alles, was er hat, vermag und erworben hat, das ist mein, mein. Denn die Kirche ist sein Fleisch und Blut, es sollen zwo Ein Leib sein, sie sitzen in gemeinen Gütern. Also soll man Christum erkennen, daß er der Mann und das Haupt ist der Christenheit, und daß die Kirche seine Hausfrau und Braut sei.

62. Zum fünften, so soll das Lämmlein eines Jahres alt sein. Auf daß es nicht zu jung wäre, so muß [es] nicht unter Einem Jahre sein, sondern eines vollkommenen Alters. Also ist der HErr Christus auch ein vollkommener Mensch, er hat vollkommene Gewalt im Himmel und auf Erden [Matth. 28, 18.], er kann wohl regieren, ihm mangelt und gebricht nichts an dem, so zum Regiment gehört.

63. Er ist ein Mann, daß er die Unvermöglichen annimmt. Item, er ist eines Jahres alt, das ist, daß er es thun kann und thun will, was ihn gelüstet. Wer Christum also isset, der ist es wohl werth, daß GOtt ihn ansieht. Er ist alt genug, das ist, stark genug und vermag es. Es soll jung sein, das ist, Lust und Liebe hat GOtt zu ihm; er ist angenehm und angesehen, stark und gewaltig, und er kann und will helfen mit Lust und Liebe. Also hat GOtt vorzeiten den HErrn Christum abgemalt mit dem Osterlämmlein.

64. Zum sechsten, von den Lämmern und Ziegen soll man es nehmen. Den HErrn Christum soll man nehmen von den Menschen, denn er ist ihnen auch gleich, und soll uns dienen und helfen. Darum spricht GOtt: Von der Heerde oder vom Haufen sollt ihr das Osterlämmlein nehmen. Man soll es in der Heerde lassen gehen, und nicht ein Sonderliches im Winkel auferziehen. Gott bestellet es fleißig, auf daß er ja wohl Christum in uns menge und backe, und nicht aussondere, auf daß er des Fleisches und Blutes sei, deß wir sind. Welches denn über die Maßen sehr rühmt die Epistel zu den Hebräern [Cap. 2, 16.]: daß GOtt nicht die engelische Natur, sondern unser Fleisch und Blut erwischt hat, Abrahams Samen. Das ist, GOtt ist nicht ein Engel, sondern ein Mensch worden, und ein solcher Mensch, der da ist unter den andern Menschen umhergegangen [Apost. 10, 38.], gleichwie ein anderer Mensch, wie St. Paulus in der Epistel an die Philipper am andern Capitel, V. 6. 7., es gar meisterlich ausstreicht, da er spricht: „Als er GOtt gleich war, ließ er sich deß nicht dünken für einen Raub, sondern äußerte sich deß“, und stellte sich, gleich als wäre er nicht GOtt, und hat sich hernieder gelassen, gedemüthiget als ein anderer Mensch, und durch alles uns gleich befunden worden etc.

65. Das heißt Christum recht beschrieben. Er hat alle leibliche Nothdurft gelitten, die wir leiden; er hat geredet und gelacht wie wir; er hat sich gestellt, gleich als wäre er nicht GOtt, sondern hat sich der göttlichen Gewalt und Natur geäußert; das ist, er ist gleich ein Mensch gewesen, als wir sind [Hebr. 2, 14.]. Alhier haben wir nun Stärke und Trost an Christo, auf daß wir ihn für einen solchen Menschen erkennen, als wir sind, und nicht vor ihm fliehen mögen, oder Scheu vor ihm tragen, denn es ist keine lieblichere Creatur denn ein Mensch. Wie denn derjenige fühlt, der allein ist. Denn wenn er des Nachts wandert, so ist nicht so lieblich einen Hund oder Pferd hören, als wenn man einen Menschen hört. Denn zu dem Menschen versieht man sich mehr Gutes, denn wenn man einen Engel hörte, davor man erschrecken und sich entsetzen würde, wie die Exempel der Schrift vielmal bezeugen. Und obwohl zuweilen die Menschen unter sich arg und böse sind, so ist doch alhier die rechte Art und Natur des Menschen in Christo, daß wir in Anfechtungen und allen Nöthen zu ihm Zuflucht haben sollen, als zu dem, der da helfen kann. Und also ist Christus alhier auch abgemalt, daß er nicht sonst einem wilden Thiere verglichen, sondern als ein Lämmlein sei, das fein lieblich ist, keinen Fehl habe, GOtt so angenehm, daß er es nicht verwerfe, daß ein Männlein sei, und sich unser annehme. Item, daß eines Jahres alt sei, das Macht und Stärke habe.

66. Zum siebenten so sollte man das Lämmlein behalten bis auf den vierzehnten Tag des Monden, das ist, bis auf die Zeit der Zukunft Christi. Denn also ist es beschrieben und bestimmt vom Mose, daß Christus das Osterlämmllein
\quellseite{432}
für uns werden sollte. Das hat sich verzogen, und ist das Lämmlein behalten worden bis auf den vierzehnten Tag, daß ist, bis daß die Zeit kam, daß er sollte erscheinen [Gal. 4, 4.]. Oder man mag die vierzehn Tage also nehmen, daß sie anzeigten das Verlangen; das kleine Stündlein, da die Seele in Angst und Noth ist [Jes. 54, 8. Ps. 30, 6.], und begehrt Trost; daß das Gesetz vorher geht, daß es das Gewissen dringe zur Erkenntniß der Sünde; daß es dürste und hungere nach der Gnade GOttes [Jes. 66, 2.]. Denn der HErr Christus schmeckt niemand, denn einer hungrigen und durstigen Seele.

67. Darum, so wird es vierzehn Tage aufgeschoben, das ist, eine kurze Zeit, auf daß die Seelen zu der Erkenntniß ihrer Sünde, Noth und Gebrechen geführt werden, und im Hunger und Durst der Gnade liegen, daß dieselbigen vierzehn Tage die Zeit des Gesetzes sei, darinnen wir alle stecken und liegen müssen, das ist, hungern und dürsten. Denn mit Verlangen, Lust und Liebe will der HErr diese Speise angenommen haben. Diese Speise gehört nicht für eine satte Seele.

68. Alhier werden auch in und durch Christum alle guten Werke ausgeschlossen, so man thun kann, GOttes Gnade zu verdienen. Die vierzehn Tage sind die Zeit des Gesetzes, die Erkenntniß der Sünde, da uns das Gesetz zu Sündern macht. Derhalben so sollen wir nicht auf unsere guten Werke bauen, oder etwas vornehmen, das vor GOtt bestehen solle; sondern wir sollen dies Lämmlein allein haben. Es lautet, als sollte man ein Schaf nehmen; aber er redet auf die hebräische Art, daß, wer dies Osterlämmlein nicht isset, der sei verloren.

69. Zum achten: Und ein jegliches Häuslein im ganzen Israel sollte [es] schlachten zwischen Abends. Das ist, in allen Häusern soll man ein solch Osterlämmlein essen. Wer nun das nicht thut, und Christum, das rechte Osterlämmlein, nicht isset, dem ist nicht geholfen. Aber wie wird Christus von uns geschlachtet? Das geschieht, wenn wir erkennen, daß Christus für uns gestorben sei, und wenn wir, nach St. Pauli Meinung, Christum bekennen und predigen, da er [Röm. 15, 16.] spricht: Ich heilige das Evangelium, das ist, opfere das Evangelium. Da bedeutet er das Predigtamt, daß es solle ein Schlachtamt sein. Das Opfern zieht er allezeit auf das Predigtamt. Denn ich thue GOtt ein groß Opfer, wenn ich von Christo predige [Ps. 116, 17.]. Mit der Predigt von Christo opfere ich GOtt das höchste und allerliebste Opfer, und erfülle alle die Opfer, was sie bedeuten, und tödte den alten Menschen, und bekehre sie, daß sie neue Menschen werden.

70. Zum neunten so soll man es schlachten zwischen Abends. Man soll predigen und bekennen, wenn es halb Nacht und Abend ist; dieselbige Zeit wird der Abend genannt oder die letzte Stunde, wie der HErr Christus in den Parabeln anzeigt; und Johannis am fünften Capitel, V. 28., wird gesagt: „Es kommt die Zeit“ etc. Denn der jüngste Tag soll folgen auf diese Predigt, und nach dieser Predigt soll keine andere aufkommen von GOtt. Der Teufel wird wohl andere aufbringen, wie er denn für und für Ketzerei erweckt; aber GOtt hat erstlich das Gesetz gegeben; nunmals, am Ende der Welt, kommt Christus, und bringt die Predigt des Evangelii, welche Lehre der Heilige Geist gebracht hat. Denn dieses ist die letzte Predigt in der Welt, und wird genannt eine Predigt, die da geschieht auf den Abend [Matth. 20, 8.], das ist, am Ende der Welt. Dies Lämmlein soll man schlachten, das ist, von diesem Christo soll man allein predigen.

71. Zum zehnten: „Und sollt seines Blutes nehmen, und beide Pfosten an der Thüre, und die oberste Schwelle damit bestreichen, an den Häusern, da sie es innen essen.“ Solches kann man auf das Predigtamt ziehen. Denn das Aussprengen oder Bestreichen geschieht mit dem Munde oder mit der Zunge. Andere ziehen es auf menschliche Gedanken, auf das Gedächtniß und Verstand, oder auf die Vernunft; aber GOtt redet von den Sachen, so in die Gemeine hinein gehören. Denn in der ganzen Christenheit, in der Versammlung, sollen sie essen das Osterlamm, die da nicht hat meine oder deine Gedanken, sondern etwas Anderes. Derhalben so thue nach St. Petri Lehre, der da [1. Ep. 1, 19.] auch sagt: „Wir sind besprenget mit dem Blute Christi“ [1 Joh. 1, 7. Apost. 1, 5.]. Denn wenn wir durch den Heiligen Geist erleuchtet sind, so soll man allenthalben von dem Osterlamme predigen, auf daß man in der ganzen Welt wisse, wie das Lämmlein für uns gestorben sei, und sein Blut für uns vergossen habe. Wenn ich also predige von Christo, daß er für uns sein
\quellseite{433}
Blut vergossen hat, dann so bestreiche und besprenge ich die oberste Schwelle. Wenn ich predige, daß er mich gegen GOtt versöhnt habe, daß, wo ich mich hinkehre, da sehe ich das Blut allenthalben in dem Ausgange,\footnote{In diesem Satze ist die Construction gebrochen. Es sollte wohl so heißen: Wenn ich predige, daß er mich gegen GOtt versöhnt habe, dann sehe ich, wo ich mich hinkehre, daß Blut allenthalben in dem Ausgange etc.} und sonderlich an der Thür, da man aus- und eingehet; denn darum ist die Thür vom Zimmermann gemacht.

72. Die Thür heißt in der heiligen Schrift der Wandel und Leben eines Menschen, so er führt. Darum wird in den Geschichten der Apostel am ersten Capitel, V. 3., gesagt von St. Luca, daß der HErr Christus unter ihnen sei aus- und eingegangen, das ist, seinen Wandel und Wesen unter ihnen gehabt. Wie sonst auch in dem 121. Psalm, V. 8., gesagt wird: GOtt habe seinen Eingang und Ausgang bei uns, gleichwie es sich sonst in diesem Leben zuträgt. So soll nun das Leben also gefasset sein, daß durch die Predigt das Blut soll über sich gesprengt, und die Ueberschwellen damit bestrichen werden, nämlich, wenn ich von meinem Wesen und Leben rede, wie es vor GOtt gilt, daß mich da nichts anders helfe denn das Blut Christi.

73. Das ist denn ein recht Bestreichen, und da ist denn das Blut Christi zwischen GOtt und mir. Was ich dann thue, das gefällt GOtt, denn das Blut handelt zwischen mir und GOtt. Darum so gefällt mein Werk GOtt wohl. Denn also sollen alle Werke eines Christen eingelegt und gefasset werden in das Blut Christi, auf daß sie angenehm und GOtt wohlgefällig wären.

74. Denn das Blut Christi schützt mich zur rechten Hand, wider die Anfechtung der Güter, wenn es einem wohl geht; item, zur linken Seite, wenn der Satan mich anficht mit Verfolgung Leibes, Gutes und Ehre. Zur rechten Hand, wenn er uns ansicht mit Schein und List, als da geschieht mit den Rottengeistern. Auf der linken Seite, wenn er einher platzt und schreicht mit Verfolgung, Aengstigung, Schwermuth, und anderm, damit er uns linkwärts zusetzt, wie denn der ein und neunzigste Psalm, V. 7., sagt: „Ob gleich tausend fallen zu deiner Seite, und zehntausend zu deiner Rechten, so wird es doch dich nicht treffen.“

75. Da sollen nun die Prediger machen, wehren und predigen. Wenn du in Anfechtung liegst, daß sie dich vermahnen, an das Blut Christi zu denken und dich deß zu trösten, damit du wissest, daß es GOtt gefalle. Das muß man fleißig immerdar predigen, ob irgend falsche Propheten möchten aufstehen, daß man allezeit bei dem Blute Christi bleibe, und dazwischen gehe, daß unser Wandel gefasset sei in dem Blute Christi, es lebe oder sterbe einer, oder es gehe ihm sonst süß oder sauer. Dieses soll das Predigtamt treiben [Jes. 30, 8. 21. Ezech. 33, 2. ff.].

76. Nun kommt er darauf, wie man das Osterlämmlein essen soll. Bisher hat er vom Blute des Osterlämmleins gesagt, wie man es nehmen, und die oberste Schwelle damit besprengen und bestreichen soll, welches anders nichts ist, denn wie man predige, verkündige, abmale und zu erkennen gebe, wer Christus sei. Der Glaube ist das Essen, der uns enthält und stärkt. Gleichwie der Bauch, wenn er isset, trinkt, Speise zu sich nimmt, schluckt er es in sich, und gibt es wieder an den Ort, da es verkehrt wird in Blut, daß die Speise des Leibes Nahrung wird: also isset der Mensch auch geistlich, wenn er verschlingt und verdäuet Christum, bringt in sich Christum, und Christus nährt ihn.

77. Darum so ist dasselbige Essen anders nichts, denn der wahrhaftige, rechte Glaube des Herzens, wenn du mit rechtem Glauben annimmst Christum, und weißt,\footnote{Eislebensche: wissest.} daß er sein Blut für dich vergossen habe, und daß dich solches in Kreuz und Noth tröstet und stärkt, dieweil du es ohne alles Wanken des Herzens also glaubst, also issest du Christum und verdauest\footnote{Im Original: „däufest“.} ihn in dir, und er kommt in dich, daß ihr geistlicher Weise ein Ding werdet, und einerlei Gedanken und Sinn habet, auch einerlei Willen, Weisheit, Klugheit, Stärke und Gewinn, daß man ein neuer Mensch werde, der da täglich zunimmt, wächst, groß, fett und stark wird in dem Erkenntniß des HErrn Christi [Col. 3, 10.]. Wie denn davon der Herr Christus Johannes am 6. Capitel, V. 35., auch sagt: „Wer mich isset, den wird nicht hungern.“ Dasselb hast du auch das geistliche Essen des Herzens. Denn was der Christenmensch mit dem Munde zu sich nimmt, das hilft ihn nichts an seinem Christenthum [1 Cor. 8, 8.], aber wenn das Herz etwas durch den Glauben zu sich reißt, das hilft, und dadurch
\quellseite{434}
wird einer vor GOtt ein reicher, völliger Christ, daß alles GOtt wohlgefällt, was er thut.

78. Zum zwölften sagt Moses: „Ihr sollt also Fleisch essen in derselbigen Nacht.“ Er heißt „Fleisch“ den HErrn Christum, wiewohl er Christi Fleisch noch nicht deutlich ausdrückt; aber er will sagen: Ihr seid geneigt zu essen, das dem Leibe wohlthut und lieblich ist: wohlan, ich will euch einmal ein recht Fleisch zu essen geben, davon ihr lustig werdet zu essen, und das da eine liebliche Speise sei.

79. Moses hat uns bis anher vorgestellt das Osterlämmlein, dadurch der Heilige Geist uns hat bedeuten und zeigen wollen unsern lieben HErrn und Heiland Jesum Christum, und dasselbige auf zweierlei Weise; erstens, wie ich von ihm soll predigen, auch was ich von ihm halten solle. Als, daß das Lämmlein solle ein Männlein sein, und nicht ein Weiblein. Item, Eines Jahres alt, daran kein Fehl noch Gebrechen sei. Auch, daß man es des Abends schlachte und opfere, damit auch die zwo Säulen an der Thüre bestreiche etc. Darinnen das Predigtamt des göttlichen Worts uns eigentlich beschrieben ist, daß wir von Christo predigen sollen, wie er GOtt und uns mit einander versöhne, und in der Welt uns führe, beide, zur Rechten und zur Linken.

80. Darnach ist das andere Stück gewesen, wie man das Osterlamm essen solle. Nämlich, man solle dies Lämmlein nehmen, opfern, schlachten und braten, daß man es esse. Das ist nun, das Predigtamt führen, und Christum oder das mündliche Wort vortragen, und ist „essen“ hinein- und zu sich nehmen, daß es der Mensch verdaue und in sich verwandele, davon er fett und gesund wird, zunimmt und wächst [Ps. 107, 9.]. Denn darum wird es geschlachtet, daß man es essen soll. Also wird es darum auch gepredigt, daß wir von Christo errettet, fett und stark gemacht werden, welches denn anders nichts ist, denn glauben. Der Glaube ist das geistliche Essen und Verdauen, dadurch alles gestärkt, und Christus besser erkannt wird, was man an ihm habe und sich zu ihm versehen solle. Daher wird Johannes am sechsten Capitel, V. 54. 51., auch gesagt: „Wer mein Fleisch isset“, item: „Wer von dem Brode isset, der stirbet nimmermehr.“ Darnach, so gibt das Essen auch eine Kraft und Stärke, so da lebendig macht, und Christum in uns, und hinwieder uns in Christo erhält.

81. Sonst ist Fleisch essen ein böse Zeichen in der heiligen Schrift, wie denn GOtt auch im Mose verboten hat, Blut zu essen [1 Mos. 9, 4. 3 Mos. 7, 26. 5 Mos. 12, 16.]. Und also heißt anderswo „Fleisch essen“ glauben, und da hält man etwas vom Fleisch und Blute. Aber wenn man predigt vom vernünftigen Leben, wie denn die Heuchler gethan haben, da ist denn Fleisch und Blut, wie es von Adam geboren ist. Da sind denn die rechten Fleischfresser, welche glauben, daß sie durch ihre Kräfte und Vermögen noch etwas thun können, daß sie zu GOtt bekehrt werden. Das ist denn verboten, und werden solche Fleischfresser ungestraft nicht bleiben.

82. Aber allhier wird viel ein ander Fleisch zugerichtet, daß man es essen möge, und wird geboten, wie man es braten und mit ungesäuertem Brode und bittern Salzen essen soll. Denn es ist eine andere Speise. Also sollte er sagen: Bisher hat man gelehrt und gepredigt, wie man fromm werde, aber es ist Menschentand, und nur wie es die Vernunft gut dünkt. Wollt ihr euch wohl speisen und lehren lassen, so will ich auch einmal euch Fleisch zu essen geben, aber ein solch Fleisch, darinnen Christus wohnt. Also ist Christus unsere Speise, daß er uns gibt seinen Leib und Blut, wie er das natürlich von seiner Mutter empfangen hat; das ist unsere geistliche Speise und Trank, nicht um des Fleisches willen allein, sondern daß GOtt drin wohnt. Er spricht nicht: Das Fleisch ist die Speise, und das Blut ist der Trank, sondern: Mein Fleisch und mein Blut; mein, denn ich bin GOtt.

83. Wenn ich das glaube, daß Christus am Kreuz für mich gehangen hat, und wahrhaftig Fleisch und Blut worden, und dennoch GOtt ist, und diese Person, so GOtt ist, mag hinan für mich sein Fleisch und Blut (denn GOtt könnte sonst nicht leiden oder sterben; aber darum hat er das Fleisch an sich genommen, auf daß er leiden könnte, und hat sein Fleisch und Blut für mich gegeben, damit hat er mir geholfen) — wenn ich nun glaube, daß dies Fleisch für mich gegeben und geopfert sei am Kreuze, das heißt essen und trinken, das ist,
\quellseite{435}
glauben, Christi Blut und Leib sei für mich gegeben.

84. Dieses hätte aller Heiligen Fleisch nicht thun können, keiner hätte sagen dürfen: So du glaubst, daß dieser Leib für dich gegeben sei, so wirst du selig. Nein, das können sie nicht sagen; allein dieser Sohn GOttes, diese göttliche Person, wendet dies Stück, sein Leib und Blut, an mich. Darum gilt sein Fleisch und Blut so viel, als er gilt. Denn die Person ist GOtt, derhalben so hat dies Fleisch und Blut kein Ende, sondern bedeutet etwas Ewiges, und bleibt auch ewig. Sonst ist alles Fleisch eitel und nichts [Jes. 40, 8. Ps. 62, 10.], ausgenommen, wenn dieser Mann kommt und spricht: „Mein Fleisch ist eure Speise“ etc. [Joh. 6, 55.], denn ich bin GOtt. Diese Person ist zu hoch.

85. Dieses Essen ist, an Christum glauben; und sein Blut trinken, das ist, von Herzen halten, daß er für uns in den Tod gegeben sei. Das hat GOtt durch dies äußerliche Osterfest der Juden anzeigen wollen. Sie waren leiblich, und mußten leiblich vom Lämmlein essen, aber wir essen geistlich davon, das ist, wir glauben, daß des HErrn Christi Fleisch und Blut für uns gegeben sei. Aber dasselbe Essen geschieht durch den heiligen Geist.

86. Zum dreizehnten spricht er: „In derselbigen Nacht.“ Wenn der Abend anging, da schlachtete man das Lämmlein; bis daß dies Osterlämmlein zugerichtet ward, so ist es gut Nacht gewesen, und haben es müssen im Finstern essen. Also sollen wir das Osterlämmlein auch des Nachts essen. Wenn ich recht an Christum glaube, mein Herz und Seele sich nährt deß, daß er seinen Leib für mich gegeben hat, und zweifle nicht daran, dann so wird es Nacht; ich habe dann nichts mehr in der weiten Welt [Phil. 1, 21. Gal. 2, 20.], mein Leben ist in Christo verborgen, und wir sind der Welt abgestorben, wie St. Paulus zum Colossern [Cap. 3, 3.] sagt.

87. Sonst hat die Welt mit ihrer Vernunft ein Licht [2 Cor. 4, 3.], und leuchtet schön, spricht: Das sollst du an diesem Tag essen, und also sollst du dich auf jenen Tag kleiden, so und so leben, dies und jenes thun. Dieser Orden thut heute das, morgen thut jener Orden etwas Anderes, er betet, fastet, und thut dies und jenes; binden und fassen also die Gewissen mit Gesetzen, daß es alles solle vor GOtt gelten, und sehen nach diesen Dingen. Aber ein Christ schließt die Augen zu, und fragt nichts darnach, lebt in Tag hinein, spricht: Vor GOtt bin ich nichts besser, ich esse Fleisch oder Fische, oder was mir vorgehet werde [Röm. 14, 2. f.]. Eines Christen Gewissen kehrt sich nicht dran, sondern spricht: GOtt hat mir ein Osterlamm vorgesetzt, davon soll ich essen, und dabei will ich auch bleiben. Denn an seinem Worte da soll man allein hangen. Kommen denn äußerliche Dinge, so laß man sie gehen. Derhalben ist die Welt einem Christen eine lautere Nacht und Finsterniß, welches der Glaube macht, welcher allein an dem Blute Jesu Christi hangt, und sonst nichts ansieht noch achtet.

88. Zum vierzehnten so zeigt Moses oft und viel an, daß man das Osterlämmlein soll mit ungesäuertem Brode essen, oder süße Brode dazu gebrauchen, da kein Sauerteig zu gekommen sei, und acht Tage lang solche süßen Brode essen. Was dieses bedeutet, das legt der HErr Christus selbst aus im Evangelio Matthäi am sechzehnten, V. 11., da er spricht: „Hütet euch vor dem Sauerteig der Pharisiäer.“ Denn dieselbigen ihre Menschenaufsatzungen höher hielten denn GOttes Gebot; hoben GOttes Gebot auf, daß nur ihre Aufsätze blieben. Diese ihre heuchlerische Lehre heißt der HErr „Sauerteig“, damit die Schälke, die Pharisäer, das Volk verführen. So wird nun mit den süßen Broden angezeigt, daß man den Glauben rein behalten solle; denn er kann keinen Zusatz leiden. Man soll Menschenlehre nicht drein mengen, denn der Glaube leidet es nicht.

89. Eben dasselbe wird auch 1 Cor. 5, 7. von St. Paulo gesagt: „Feget aus den alten Sauerteig“; denn es ist alles Sauerteig, was da nicht christlich lebt; und befiehlt, daß man denjenigen, so seine Stiefmutter zur Ehe hatte genommen, sollte aus der Gemeinde stoßen und nicht leiden. So ist nun Sauerteig falsche Lehre und ein böse ergerlich Leben und Exempel. Denn ein räudig Schaf steckt sonst gerne eine ganze Heerde an, und die flechtenden\footnote{„flechtende“ wohl so viel als flechtenartige. (Diese Bedeutung fehlt bei Dietz.)} Seuchen machen auch andere Leute krank. Darum dies Wort „Sauerteig“ begreift allhier die Lehre, wiewohl auch oft das Leben damit verstanden wird. So sollen wir nun den Glauben also fassen, daß
\quellseite{436}
wir Christum, das Lämmlein, einnehmen, damit unser Gewissen rein, lauter und schlecht allein an Christo hange und bleibe. Aber das geht gar schwer zu, und wird denen sauer, die mit Gesetzen gefangen gehalten werden. Darum muß man allein bloß im Glauben leben. Der Glaube bringt uns zu das Leiden und Sterben unsers lieben HErrn Jesu Christi. Sonst will unsere Natur Ceremonien und allerlei verdienstliche Werke aufrichten, darauf das Gewissen sich steuern\footnote{steuern = stützen.} und verlassen soll, aber das ist Sauerteig. So will nun Moses uns ermuntern, daß wir den Glauben rein und unverfälscht behalten, und uns hüten sollen vor allen Zusäßlein und Menschensatzlein; denn Sauerteig ist gleich als ein Zusatz.

90. Zum fünfzehnten so soll man es mit bittern Salzen essen. Denn es schmeckt wohl, wenn ein wenig Essig dabei ist; das gibt die Natur, daß bittere Salzen und saurer Kraut neben gutem Gebraten\footnote{Gebraten oder Gebratnis = Braten.} wohl schmecke. Aber es bedeutet, daß man den alten Adam kreuzigen soll, denn soll es sauer werden. Der Glaube richtet das Herz auf, macht es stark und getrost in Christo, der für uns gestorben und in [den] Tod gegeben ist. Darauf gehört und soll folgen, daß der alte Adam nach dem äußerlichen Menschen muß gepanzerseget werden [Gal. 5, 13. 16. 17.]; Sauerkräutlein muß dabei sein, das Gewissen muß angegriffen und getroffen werden. Denn das Evangelium ist nicht eine Predigt für grobe, rohe, wüste Sünder, die ohne einige Andacht leben, sondern ist ein Trost für die betrübten Seelen [Matth. 11, 28.]. Denn es ist eine zarte Speise, die will haben eine hungrige Seele. Daher die liebe Jungfrau Maria in ihrem Magnificat [Luc. 1, 53.] auch singt: „Die Hungrigen hat er mit Gütern gefüllet.“

91. Sonst fällt der tolle Pöbel hinein, und wollen alle evangelisch und christliche Brüder sein, richten darnach Rotten und alles Unglück an. Den Teufel auf ihren Kopf sind sie. Ein Christ ist nicht frech, wild und roh, sondern sein Gewissen ist blöde, kleinmüthig und verzagt, die Sünde beißt sie, und fürchtet sich vor GOttes Zorn und vor dem Teufel und Tode. Einem solchen niedergeschlagenen und zertriebenen Herzen schmeckt der HErr Christus wohl. Item, die Erlösung von der Sünde, Tod, Teufel und Hölle schmeckt denen auch wohl, die in dem Tode stecken, und fühlen solche Noth, und wollten gerne Ruhe haben; die bekommen sie, wenn das Herz den Glauben hat; aber sie fühlen auch daneben, wie gebrechlich der alte Adam sei.

92. Zum siebzehnten so soll man das Osterlämmlein nicht roh essen, noch mit Wasser gesotten, sondern nur am Feuer gebraten. Es ist zweierlei Glaube: einer, da man das Evangelium hört und der Heilige Geist das Herz erleuchtet, rührt und anzündet, daß man es annehme und glaube. Der Glaube macht denn los von Sünden und selig; aber der andere Glaube wächst aus unserm Kopfe. Als, wenn man das Evangelium von Christo hört predigen, wie er für uns gestorben sei, und solches will verstehen und fassen mit der Vernunft; wie man der Leute allzumal findet, welche meinen, sie wissen es, und haben einen Glauben; das ist ein Glaube, aus unsern Kräften hergekommen, und ist ein Traum, ein Bild, so des Nachts einem vorkommt, aber kein rechter Glaube. Es ist gleich damit, als wenn einer im Schlafe Geld findet, oder eine Stadt gewinnt; da bäucht einem im Traume, daß es Ernst sei; aber wenn er aufmacht, so ist es nichts dahinter. Also ist es mit denen auch; sie meinen, sie haben es, und sind gute Christen, haben den Glauben und verstehen Christum sehr wohl. Aber dieser Glaube gewinnt bald ein Ende; wenn Verfolgung hergeht, daß man Gefahr an Gut, Ehre und Leib leiden soll, da werden die Augen aufgethan, und ist der Traum hinweg. Des HErrn Christi wird gar vergessen; denn es ist nur ein Schein gewesen und ein Wahn vom Glauben. Gleichwie man eines Menschen Bildniß und Angesicht in einem Wasser oder Spiegel sieht, das sieht dem ganz ähnlich, und dünkt ihn wohl, als sähe er denselbigen Menschen, aber es ist wahrhaftig nichts. Also, wenn es zum Treffen kommt, so fleugt dieser Glaube auch hinweg, läßt sich sehen, daß er nichts sei. Solches wird allhier gepredigt, daß man das Lämmlein mit Feuer braten solle.

93. Mit Wasser kochen ist, mit der Vernunft anzünden, da das Feuer nicht zum Lamme kommt, noch die Kälte des Wassers wärmen möge; die Hitze wird vielmehr durch das Wasser eingetrieben. Wasser heißt ein solcher Dünkel und Meinung, da man meint, man kenne Christum
\quellseite{437}
, und ist doch falsch. Sondern also soll es sein: Ich höre das Evangelium wohl, und es dünkt mich, ich verstehe es; aber ich traue nicht darauf, sondern ich bitte GOtt, daß er mir den Glauben gebe auf sein Wort [Luc. 17, 5.], welches der rechte und nicht erdichtete Glaube ist. Ich habe ihn im Traume nicht erdichtet, sondern der Heilige Geist hat ihn rechtschaffen gemacht in meinem Herzen, ohne meine Kräfte, durch das Wort, das gepredigt wird. Solch Feuer zündet er selbst an [Luc. 12, 49.], und [er]leuchtet mein Herz mit einem neuen Lichte und Brunst, daß ich mitten im Tode des HErrn Christi nicht vergesse.

94. Der wässerige oder der falsche Glaube verschwindet. Gleichwie ein Bildniß aus einem Spiegel verschwindet, also hält dieser Glaube den Puff und Stich nicht, sondern der feurige Glaube, den der Heilige Geist anzündet, derselbige besteht, ob er wohl gleich zuweilen erschreckt und betrübt wird. Denn da ist noch der alte Adam, der da zappelt, dennoch so hält er aus. Das Feuer ist der Heilige Geist, der das Osterlämmlein in uns gebraten macht, und einen rechten Glauben anrichtet, so Christum annimmt, und bei ihm bleibt, und wir auch gebraten und nicht gekocht werden, gleichwie er durch das Feuer gebraten ist. Zu diesem Glauben gehört nun Anfechtung und Trübsal, auf daß er flugs brenne und nicht verlösche. Und wer recht glaubt, der hat ein Feuer, das ist, er steht nicht auf seine Vernunft. Denn der Glaube ist nicht eine menschliche, sondern göttliche Kraft, wie St. Paulus zum Römern [Cap. 1, 16.] sagt.

95. Diese Kraft verstöret nicht, ist auch nicht müßig, sondern wird überschüttet mit Trübsal, auf daß der Mensch sehe, daß er nicht in seiner, sondern in göttlicher Kraft hergehe. Und der Glaube ist nimmermehr stärker und herrlicher, denn wenn die Trübsal und Anfechtung am größten ist; derhalben wird der Glaube feurig genannt. Und ich fürchte sehr, daß unser vieler Glaube auch wässerig sei, daß wir nur mit dem Maule viel davon reden; wenn aber die Püffe kommen, daß wir verfolgt werden, dann so kommt dies Feuer, und wird das Lämmlein in unserm Herzen gebraten. So ist nun jenes ein falscher Glaube, dieses aber ein rechter Glaube; jenes ist ein wässeriger Glaube, dieses aber ein feuriger Glaube.

96. Zum siebzehnten, man soll das Osterlämmlein gar essen, sein Haupt mit seinen Schenkeln und Eingeweide. Das Haupt haben sie gemacht zur Gottheit, die Schenkel aber die Menschheit Christi. Ich will es auf einen Haufen nehmen. Denn er will also viel sagen: Man soll das Osterlämmlein ganz und gar aufessen, und nichts davon überlassen. Ganz essen verstehe ich also, daß man Christum gar esse, und an ihm allein hange und nichts anders suche, denn ihn. Dieses wird dienen wider der Juden Glauben, die Christum nicht für den rechten Heiland halten, und vorgeben, Messias sei noch nicht gekommen, und auf einen andern Messiam, auf einen andern Glauben, und auf eine andere Lehre warten. Aber wisse, wer Christum hat, der hat alles auf einen Haufen.

97. Also, wenn ich an Christum glaube, so habe ich Vergebung der Sünde, Herrschaft über den Tod und Teufel, habe auch das ewige Leben. Allhier mangelt nichts, es fehlt nicht um eine Klaue, noch um ein Haar. GOtt will uns allezeit im reinen Glauben erhalten, da soll man keinen Zusatz machen, daß man ihn wolle bessern, oder zweifeln, als mangele etwas daran; wie denn St. Paulus zum Colossern [Cap. 1, 12.] nennt den Glauben das ganze Erbe. Als sollte er sagen: Der Glaube hat es rund in sich gefaßt, es sei oben, mitten oder am Ende: nimm es, wo du willst, so hast du es gar. Also sind wir sicher und gewiß, daß wir alles haben und alles essen.

98. Zum achtzehnten,\footnote{Die Erlanger hat hier und bei der folgenden Abtheilung: „Zum 19.“} daß man nichts davon überlasse bis morgen. Die Juden warten, gleich als sollte noch ein ander Osterlämmlein kommen, und wollen ihr Osterlämmlein erstreckt haben, bis ein anderes komme, und eine andere Lehre ihnen gebracht werde, oder der Tag anbreche; welche Lehre jenes Osterlämmlein verbrannt und aufgehoben, geht noch immer hin, obwohl der Tag vergangen ist und das Osterlämmlein kommen.\footnote{Der Sinn dieses Satzes wird sein: wiewohl das rechte Osterlamm, Christus, die Lehre von einem zukünftigen Osterlamme aufgehoben hat, wird diese Lehre von den Juden noch immer festgehalten, obwohl etc.} Das ist, daß alte Testament ist aus, es gilt nicht mehr; derhalben hat er geboten, sie sollten sich an dasjenige halten, und ob etwas anders vorhanden und übrig wäre, so sollte es mit Feuer, das ist, mit dem
\quellseite{438}
Heiligen Geiste verbrannt werden. Als sollte er sagen: „Wir haben alles in Christo, und ist nicht vonnöthen, daß wir allerlei Bräuche, Sitten, Weise und Ceremonien halten, oder die Werke des alten Testamentes thun, sondern wir haben alles in Christo. Also bleibt der Glaube an Christo, und hat seine Freude, Lust und Spiel an ihm. Dieses hat der Juden Osterlämmlein sie erinnern sollen, bis daß Christus, das rechte, wahre Lämmlein käme, welches jenes aufgehoben hat. Wie wir denn sehen, daß alles um seinetwillen gethan ist, und auf ihn alles geht, auf daß wir in der Predigt und Erkenntniß des HErrn Christi vollkommen und reich würden [Col. 3, 16.].

99. Also haben wir gehört, wie man das Osterlämmlein essen, das ist, von Christo predigen soll. Man soll predigen und lehren, daß man ihn ganz esse, das ist, man soll nichts predigen, denn allein den ganzen Glauben, der thut gegen GOtt alles. Nun wird er auf eine wunderbare Weise die Leute zubereiten, die das Osterlämmlein essen sollen, da wir doch sollten die Leute zuvor unterrichten, als, wie sie die Hände waschen sollten, und neue Kleider anziehen, ehe sie das Osterlämmlein essen. Aber GOtt kehrt allhier die Ordnung gar um. Denn die heilige Schrift wollte gerne wehren unsern Werken, auf daß wir zuvor glauben und durch das Wort GOttes Vertrauen an Christum haben, uns demüthigen, und alles in GOtt stellen.

100. Passah heißt, Eingang oder Uebergang, und ist daher also genannt, daß der Engel aus GOttes Befehl auf eine Nacht ging durch ganz Egyptenland, und würgte oder tödtete alle Erstgeburt der Egypter. Es ist aber hiermit bedeutet des HErrn Christi Leiden, Sterben und Auferstehung, dadurch er aus dieser Welt gegangen, und zu seinem himmlischen Vater gekommen ist. In demjenigen Durchgange hat er Sünde, Tod, Teufel und Hölle erlegt und geschlagen, und aus dieser egyptischen Gefängniß alle seine Christen und Gläubigen erlößt, und führt sie seinem himmlischen Vater zu. Dies ist unser, der Christen, und das rechte, tröstliche, freudenreiche Passah oder Ostern. Also sagt der Evangelist Johannes, am 16. Capitel, V. 28., daß Christus zu seinen Jüngern gesagt habe: „Wiederum verlasse ich die Welt und gehe zum Vater.“ Das ist der rechte Gang.

101. Zum neunzehnten: „Um eure Lenden sollt ihr gegürtet sein.“ Die heilige Schrift redet mäßig und züchtig und nicht frech. St. Petrus macht es geistlich, und will, daß wir nicht im Glauben mancherlei Secten haben sollen, sondern in Einem Glauben bleiben. Aber allhier will GOtt die äußerliche Zucht fassen, daß man Keuschheit halten solle, und gebeut, man solle die Lenden gürten. Denn die heilige Schrift schreibt die Keuschheit den Lenden zu. Also wird im ersten Buch Mosis auch gesagt, daß aus den Lenden Abrahä geboren werden sollen etc. Darum so bedeutet es nicht allein die Zucht oder Keuschheit, sondern den ganzen Leib. Gleichwie der ganze Mensch von Einem Weibe geboren wird, also bedeutet es allhier auch die ganze Zucht des äußerlichen ganzen Wandels, daß man ein ehrbar und züchtig Leben führe, nicht wild, roh oder bübisch lebe [Röm. 12, 2.]. Denn wenn man an GOtt glaubt, so isset man das Osterlämmlein. Aber darnach umgürte ich meine Lenden, das ist, halte den Leib in Zucht und im Zwange, damit ich selbst nicht in Sünden gerathe, auch andern mit bösem Leben kein Aergerniß vorstelle [1 Cor. 9, 27.].

102. Zum zwanzigsten: „Und eure Schuhe an euren Füßen haben.“ GOtt gibt einem jeglichen sein bescheiden Theil: diesem Weibe den Mann, und wiederum dem Manne auch jenes Weib, und daß einer also, der andere sonst lebt, thut und gebaret unter den Leuten. Darum können und sollen sie nicht alle auf einerlei Weise einhergehen. Solches bedeutet das Schuhanhaben. Denn wer da Schuhe anzieht, der bereitet sich zu gehen, und ist dahin gerichtet, daß er wandern will.

103. Man kann nicht alle Fälle und Exempel der Patriarchen oder ihre Werke nachahmen. Das ist eine böse Allegorie. Man soll der Heiligen Exempel nicht alleinthalben nachthun, sondern ihren Glauben und die Früchte des Glaubens ansehen, und dem folgen lernen. GOtt führt einen jeden seines Weges. Aber deß befleißige dich, daß du glauben mögest, wie sie geglaubt haben. Darum bedeutet dies Stück die Rüstung oder Bereitschaft des Evangelii, daß unsere Füße gestiefelt und angezogen sind mit dem Evangelio des Friedens. Denn darum zieht man Schuhe an, daß man gerüstet und geschickt sei zu gehen; das ist, man soll Uebung haben, das Evangelium zu verstehen, zu lesen, predigen, treiben, singen, immerdar damit umgehen
\quellseite{439}
; daß allhier der ganze Weg gefaßt ist, und geschicht sei, daß wir alles mit dem Evangelio, in dem Evangelio, und durch das Evangelium thun, was wir zu thun haben, wie es denn St. Paulus [Eph. 6, 15.] auch also auslegt.

104. Zum ein und zwanzigsten: „Und sollt Stäbe in euren Händen haben.“ Der Stab dient dazu, daß er dem Menschen hilft im Gehen, und er sich drauf lehnen könne, wenn er müde wird; item, daß er sich drauf verlassen kann, und einen Muth fassen. Solcher Stab ist nun auch das liebe Evangelium [Ps. 23, 4.]. Wenn wir in dieser Welt leben und Anstöße haben, und müde gemacht werden, da soll immer bei uns aus dem Evangelio Trost, Stärke, Vermahnung und Reizung sein. Die Predigt des Evangelii soll immerdar lehren, vermahnen und trösten, reizen, treiben und locken, daß man fortfahre und nicht faul oder laß werde. Darum so ist vonnöthen, daß man für und für mit dem Worte anhalte, und den alten, faulen Madensack aufwecke, daß er nicht schläfrig und faul werde. Solches heißt, den Stab in Händen und Schuhe an den Füßen haben, das ist, lehren und vermahnen aus GOttes Wort.

105. Zum zwei und zwanzigsten: „Und sollt essen, als die hinweg eilen.“ Sie sollen es nicht also eilends essen und behende käuen, und schnattern, wie die Gänse den Hafer essen, sondern sie sollen es eilends essen, als die da gerne davon wären, daß dies „eilends essen“ mehr das Herz betrifft denn den Mund, als die da in Angst sind. Darum so ist es eine gegenwärtige Eile, und nicht eine schlechte Eile. Das Essen möchte nicht lange währen, irgend zwo, drei oder vier Stunden. Es soll nur ein Rückbißlein\footnote{Rückbißlein = ein kleiner Imbiß.} oder Morgenbrod sein, daß man stehe und esse, da man nicht viel Gerichte gibt, und flugs sich auf den Weg macht. Denn GOtt spricht: Ihr sollt mit Eile ausziehen aus Egypten. Denn sie wurden gejagt, gleich als wären die Egypter mit Keulen hinter ihnen her. Derhalben haben sie das Osterlämmlein mit Eile, Zittern und Zagen verzehrt, daß sie nicht lange getischt haben, noch gespeiset, als wären sie zur Hochzeit geladen, sondern die da auf und davon wollten.

106. Solches bedeutet der Christen Hoffnung. Denn die an Christum glauben, die wissen, daß es mit diesem Leben muß aus sein, und wir allhier keine bleibende Statt haben, wie die Epistel zum Hebräern [Cap. 13, 14.] sagt. Es gilt allhier nicht, sich niederseßen, sondern kurzum stehen [Eph. 6, 14.], das ist, die das Evangelium hören und Christum erkennen, die haben keine Ursache allhier zu bleiben, sondern sie gedenken fort, und wollen sich davon in ein ander Land machen. Da gilt es auch hin. Gleichwie die Kinder von Israel aus Egypten eilten nach dem gelobten Lande, also sollen die Gläubigen aus diesem Leben nach einem andern und ewigen Leben gedenken. Darum so schmeckt einem Christen dies Leben nicht, er gedenkt mehr dorthin, denn hierher; er hat auch mehr Lust dorthin, denn zu diesem weltlichen Leben [Col. 3, 1. ff.]. Aber der alte Adam wehrt sich hart, und diese bittere Galle will ihm nicht schmecken; jedoch der Glaube muß fort. Denn in diesem Leben ist anders nichts denn Sünde, Plage, Unglück; so klebt das Fleisch und der Teufel uns noch am Halse; so straft GOtt auch die Sünde in dieser Welt [Ps. 90, 5. ff.]. Aber wie komme ich aus diesem Leben? Denn wer anhebt an die Wohllust dieser Welt, der will es nicht gerne entbehren, aber wer an Christum glaubt, und den Heiligen Geist hat, der verachtet denn dies Leben etc.

107. „Denn es ist des HErrn Passah“, das ist, Durchgang, er hat also in seinem Durchgange Egypten gestraft.

108. Das ist die Auslegung des Osterlämmleins, welches GOtt also befohlen, daß die Kinder von Israel essen sollten. Solch lieblich Osterlämmlein, das so lange Zeit zuvor geordnet gewesen, ist nun eingestellt bei uns, die wir das Evangelium und Christum haben, und von dem wahrhaftigen Osterlämmlein, so da Christus ist, essen, trinken und ernährt werden. In diesem Glauben haben die Kinder von Israel auch das Osterlämmlein gegessen, und sind in diesem Glauben selig worden, allein, daß sie eine Decke vor den Augen gehabt, und es ihnen etwas dunkel und verborgen gewesen ist, und nur das äußerliche, leibliche Lämmlein gesehen haben. Aber dieweil sie GOttes Wort gehabt, so haben sie Christum auch erkannt, und bei dem Osterlämmlein sich seiner Wohltat erinnert und getröstet, und in diesem unserm Glauben gestanden und selig worden [Röm. 1, 2. Eph. 4, 5.].

\quellseite{440}
109. Nun folgen am Ende noch etliche Stücke: als, zum drei und zwanzigsten, daß kein Fremdling noch Miethling, oder Hausgenoß, von dem Osterlämmlein essen sollte. Das ist der Unterschied unter den Figuren und der Erfüllung oder Deutung. Die Figuren nennen sie gemeiniglich ein Werk oder That, das in Maßen eines Gesetzes ist gestellt, daß es ein Werk bedeute, das die Menschen thun; und diese Bedeutung der Figur bleibt bei GOtt, daß er es thun sollte, oder Christus durch den Heiligen Geist. Jenes haben die Menschen gethan; aber dieses geschieht ohne Menschenwerk. Denn glauben an Christum geschieht nicht durch unsere Werke. Ein Christ wird nicht daher, daß man's ihn also heißt; ein Christ läßt sich nicht machen durch Gesetze, Gebot oder Werk, sondern es geschieht von oben herab, als ein göttlich und nicht menschlich Werk [Joh. 3, 3. Jac. 1, 17.]. Also werden die Figuren erfüllt; jenes haben die Leute gethan, dieses erfüllt GOtt.

110. Alhier wird geboten, daß ein ausländischer Mann oder Weib, ein\footnote{Im Original: kein.} Fremdling und Unbeschnittener nicht sollen vom Osterlämmlein essen, sondern allein die Kinder von Israel. Dies ist ein Werk und Gebot, das nicht die Menschen gethan haben; es ist nicht ein menschlich, sondern ein göttlich Werk, und bedeutet, es könne niemand das Osterlämmlein recht essen, er sei denn ein Glied Christi, und ein Bürger der christlichen Kirche [Eph. 2, 19.], denn es geschieht durch den Heiligen Geist, und nicht durch gute Werke. Derhalben so gebraucht dieses Lämmleins kein Fremdling noch Ungläubiger ohne den Heiligen Geist recht. Ob wohl bei den Israeliten ihrer oft zween oder mehr in einem Hause bei einander wohnten, der Herr und Hausgenoß, so konnte doch der Fremdling, Hausgenoß, Miethling oder Knecht davon nicht essen; es half ihm nichts, daß er bei dem Herrn in Einem Hause war und wohnte, er durfte des Osterlämmleins nicht mitgebrauchen und genießen. Also ist es auch, wenn gleich im Christenthume viel sich zur Kirche versammeln und Alle Christen genannt werden, sind unter den Christen, haben den Schein und Wert, empfangen das Sacrament des Altars, werden getauft, und stecken doch in Unglauben und gottlosem, sündlichem Wesen, so sind es falsche Christen und Miethlinge [Matth. 7, 21. ff.], es sei denn, daß sie weiter kommen; sonst bleiben sie unsere Hausgenossen, und genießen des Osterlämmleins nicht. Denn sie sind nur mit dem Namen, Schein und äußerlichen Wandel Christen, aber inwendig und im Grunde ist nichts dahinter. Darum so muß ein jeglicher Christum in sich bilden, und christlich werden, sich ganz und gar erneuern, sonst wird er im neuen Testamente nicht können von diesem Osterlämmlein essen.

111. Ein Knecht, das ist, ein getaufter Knecht aus den ausländischen Leuten, und nicht ein israelitischer Knecht, der doch beschnitten wäre, dieser sollte gezählt werden unter die Juden. Es ist aber diese Bedeutung, daß man den Unglauben und den alten Saft oder Adam soll vom Herzen abschneiden. Denn die Beschneidung geschieht durch das Evangelium und Heiligen Geist. Darum, ob er wohl ein Heide ist, so kommt er doch unter die Juden. Solches bedeutet die Heiden, welche sind Fremdlinge, aber erkauft durch das Blut Christi, und herzu gebracht, daß sie mit den Juden, das ist, mit den Aposteln und rechtschaffenen Gliedern Christi, sind auch gläubig, und mit ihnen gar Eins sind worden. Derhalben so sind sie geistlicher Weise beschnitten worden. Das geschieht nicht durch unser, sondern des Heiligen Geistes Werk.

112. Vorzeiten war ein getaufter Knecht nach dem jure civili nicht tüchtig, aber jetzt hat man Knechte, die man vermiethet; man kauft sie nicht mehr, sondern man mietet sie, und gibt ihnen einen Lohn. Dieses Gesetz, daß vorzeiten die Menschen gleich als [ein] Ochs oder Schaf gekauft worden, hebt GOtt hier nicht auf, sondern er bestätiget es. Ein armer Mann mußte seine Töchter oder Söhne verkaufen, wie denn im Evangelisten Matthäo [Cap. 18, 25.] der Herr, welcher mit seinem Knechte rechnete, hieß, ihn, sein Weib und Kinder verkaufen und bezahlen \&c. Solches ist vorzeiten die Weise und Brauch gewesen, und hindert an christlichem Leben nichts.

113. Darum ist es ein ungeräumtes Vorgeben, daß zu unserer Zeit etliche thun, die da sagen: Sind wir Christen, so müssen alle Dinge uns gemein sein, und machen aus dem Glauben ein menschlich Ding. Die Leute machen und ordnen nicht Einen christlichen Glauben; so wird man es auch nicht mit ihrer Vereinigung ausrichten; sondern es kommt von oben herab.

\quellseite{441}
Aus Rottensammlungen werden nicht Christen. Wenn es gälte, den Namen anzuschreiben, so wollten wir die besten Christen sein. Und also könnten die Türken sich auch versammeln. Vermöchte das christliche Wesen nicht; denn es kommt nicht von der Erde, sondern von oben herab.

114. Also hat der Seelmörder, der Papst, auch gethan, daß er geboten hat, man soll Wallfahrt gehen, die Heiligen anrufen, Ablaßbrief kaufen, und hat mit diesen Gesetzen wollen Christen machen. Aber noch nicht, lieber Geselle, es soll zum Glauben niemand gezwungen noch ein Christ genannt werden, sondern warten, bis GOtt ihn zum Christen mache. Thue den Mund und die Augen zu. GOtt läßt das Evangelium predigen in Haufen, darnach gibt er einem einen andern Muth und den Heiligen Geist, derselbige lebt darnach christlich. Jetzt will man in der Welt alle zu Christen machen, und sie zwingen zum Christenthume, und sagen: Lieber Bruder, willst du nicht ein Christ sein, so wollen wir dein Haus abbrennen.

115. Die Leibeigenschaft ist nicht wider das christliche Wesen, und wer es sagt, der leugt, sondern die christliche Freiheit erlöst die Seelen, und Christus ist ein Stifter derselbigen geistlichen Freiheit, die man nicht sieht. Was äußerlich ist, das läßt GOtt gehen, und fragt nicht so groß darnach.

116. Zum vier und zwanzigsten so soll ein jeglich Haus nur ein einig Osterlämmlein haben. Dies Gebot traf das leibliche Osterlämmlein an; da mußte man also thun, daß ein jeglich Haus nur Ein Osterlämmlein schlachtete. Aber GOtt spricht: Ich will es mit meinem geistlichen Osterlämmlein also machen, daß alle Welt nur ein einig Osterlämmlein haben soll. Denn zuvor aßen alle Welt nicht Ein Osterlämmlein allein, sondern ein jeglicher Hausvater hatte für sich ein Osterlamm. Darum durch das einige Osterlamm in einem jeglichen Hause wird bedeutet, daß in der Christenheit allein der einige Christus soll regieren, sonst soll alles andere ausgeschlagen werden [Apost. 4, 12.], das ist, allhier gilt allein der einträchtige Sinn und Verstand von Christo, auf daß nicht Rotten und Secten unter uns gemacht würden, und man Trennung anrichte, wie denn bei den Corinthiern [1 Cor. 1, 12.] geschah, da eines Theils sagten: Ich bin Apollisch; der andere war Petrisch, der dritte Paulisch, sondern Christen sein mögen, die alle Einen Mund, Sinn und Verstand hätten.

117. Diese Einigkeit unter den Christen ist nicht erzwungen durch das Gesetz, sondern eingegossen durch den Heiligen Geist. Das wird bedeutet, daß in Einem Hause nicht mehr denn Ein Osterlamm gegessen wird. Es hilft nichts zu unserer Seligkeit, denn allein der HErr Christus. Aeußerliche Dinge mögen helfen zu einer Ordnung, aber diesen Sinn und Einigkeit in der Christenheit erhält allein der HErr Christus durch den Heiligen Geist, daß alle in Einem Sinne, Herz und Glauben daher gehen, wie wir in unserm Symbolo bekennen mit diesen Worten: Der Heilige Geist hält in Einem Sinne gar eben \&c.

118. Darum sind das tolle Prediger, die durch gute Werke wollen Christen machen, denn sie zerrütten den einigen Sinn und das Werk, so von Christo gemacht ist, und wollen es verwandeln in Moslem; zerrütten und verderben also den einfältigen Sinn und Glauben; welches jetzt die tollen Propheten thun, und das tolle Volk hernach folgt, welches ohne das dazu sehr geneigt ist. Bedeutet also dies Stück die Einigkeit der Gemüther, des Geistes und Glaubens bei den Christen, daß man werde einerlei Sinn und Verstand haben.

119. Zum fünf und zwanzigsten, daß sie mußten das Osterlamm nicht außerhalb ihres Hauses essen; das ist, dies Osterlamm soll nicht unter Unreinen gegessen werden. Und geht solches wider die Ketzer und falschen Propheten, die Christum heraus ziehen; dort ist es geboten, alhier wird es ungeboten gehalten. Innerhalb des Hauses essen bedeutet dies Predigtamt, so man inwendig der Gemeinde GOttes treibt. Jedoch, ob ich gleich predige, so geht es doch nicht in das Herz, es sei denn bei denen, die da Christi Glieder sind, und da ein einfältiger Sinn ist. Darum, wenn ich wollte heraus fahren und sagen: Ich will diesen zum Christen machen mit diesem und jenem Werke, als, durch Bilderstürmen oder die Güter gemein machen, das geht nicht an.

120. Zum sechs und zwanzigsten, daß man das Lämmlein nicht zertheile. Wenn man das Lämmlein will austheilen unter die, so da nicht Christen sind, oder zertheilen, das ist unrecht. Zuvor ist dies auch geboten gewesen; nun aber wird es gehalten ohne Gebot, durch den Heiligen
\quellseite{442}
Geist, der läßt sich nicht zertheilen und schneiden in Stände, Erden, Werke und Seelen; sondern allein bei Einem Glauben bleibt es. St. Johannes der Evangelist [Cap. 19, 35. f.] hat solches auf Christum gezogen, da er am Kreuze gehangen ist. Solches hat er gethan aus dem übrigen Reichthume seines Verstandes, und zusammen geflochten das geistliche Lämmlein Christum und das leibliche Lämmlein. Zuvor sind die Gebeine leiblich gemein, und des Lämmleins Blut ist auch leiblich gewesen, gleichwie des HErrn Christi Fleisch und Blut auch leiblich ist, aber nun muß der Nutz und Brauch geistlich geschehen. Jener Brauch war gar leiblich, aber dieser Brauch ist geistlich, davon in demselben Evangelisten [Cap. 6, 54–56.] der HErr Christus sagt: „Wer da isset mein Fleisch“ etc.

121. Da mache nicht ein geistlich Blut von Christo, denn er hat einen geistlichen Leib und Blut nicht für uns gegeben, sondern ein natürlich Fleisch, so er von der Jungfrau Maria an sich genommen. Darum dies geistlich Essen steht im Wort „essen und trinken“. Wiewohl die Juden beides fleischlich verstanden; aber der HErr Christus spricht [Joh. 6, 63.]: „Meine Worte sind Geist und Leben“; zeigt damit an, solch Essen und Trinken müsse geistlich geschehen. Das ist nun, davon hier steht und gesagt wird: Lasset bleiben die leiblichen Beine.

122. Wir wollen aber jetzt bei dem Brechen bleiben, was geistlich brechen sei. Gebeine bedeuten die Stärke, so das Christenthum trägt. Denn Fleisch und Blut sonst das Christenthum nicht trägt, sondern der Heilige Geist gibt es, daß man Kraft, Muth und Stärke hat, und man nicht untergedrückt werde. Also verschlingt der Tod auch wohl den HErrn Christum, aber man muß ihm die Gebeine lassen. Er ist gestorben aus Kraft und Stärke des Fleisches, aber auferstanden durch Stärke des Geistes. Dies sind die Gebeine, die allhier in dem geistlichen Verstande bedeutet werden, als, daß man den Herrn Christum nicht hat können unterdrücken, wiewohl er sehr schwach gewesen und gestorben, wie ein anderer Mensch, Jes. 63, 5. Also werden allhier zusammen gesetzt die leiblichen und geistlichen Beine; das ist, das Osterlamm Christus, welchen jenes Osterlamm bedeutet hat, wird also sterben, daß er dennoch ewig bleibe. Und wird allhier der Tod und Auferstehung des HErrn Christi bedeutet, denn seine Gebeine sind nicht zerbrochen worden.

123. Jenes aber ist ein Werk und Gebot gewesen, daß man den hingerichteten oder erhenkten Menschen thun konnte, als, ihnen ihre Gebeine am Galgen oder Kreuz zerbrechen; aber diesem Menschen, Christo, sollten die Juden die Gebeine nicht brechen, das ist, in seiner höchsten Schwachheit und im Tode konnte er noch ganz behalten werden. Das hat müssen ein leiblich Zeichen sein, daß, gleichwie am Lämmlein die Gebeine nicht zerbrochen sind, also ist Christus auch ganz geblieben und vom Tode wieder auferstanden.

124. Also haben wir nun das Stück auch verrichtet mit dem Osterlämmlein des alten Testaments, und bedeutet solches das Predigen des Evangelii, daß man Christum erkenne und an ihn glaube, in welchem alles erfüllt wird. Wer das nun weiß, der sieht, wie der Heilige Geist abgemalt habe, was des Evangelii Amt sein solle, was man lehren solle, was Christus sei; nämlich, daß man predige, wie er sein Blut für uns vergossen und seinen Leib für uns gegeben habe, und daß man solches nicht haben noch empfangen könne, denn allein durch den Glauben.

125. Das ist nun der erste Theil, von der geistlichen Deutung des Osterlämmleins. Nun folgt das andere Stück, von den Erstgebornen, was mit denselben ferner geschehen sei, nachdem sie das Osterlämmlein gegessen und alles gethan hatten, was der HErr Mosi und Aaron geboten.

\begin{schriftzitat}
V. 29. Und zur Mitternacht schlug der HErr alle Erstgeburt in Egypten etc.
\end{schriftzitat}

126. Dieses ist das letzte Zeichen, und ein schrecklich, groß Wunderwerk, das GOtt in Egypten gethan hat, da er alle Erstgeburt der Egypter todtgeschlagen hat; aber der Israeliter Erstgeburt ist lebendig geblieben. Mit welchem Wunderwerke GOtt seine Allmacht beweist, daß er in Einer Nacht die Erstgeburt an Menschen und Vieh tödtet, welches uns billig zur Furcht GOttes reizen sollte; dazu seinen Ernst sehen läßt gegen die unbußfertigen Sünder, daß er mit der Strafe nicht ferne von ihnen sei. Damit will er jedermann zur Besserung Ursache geben, darnach auch, seine Gnade und Barmherzigkeit erzeigen, daß er in dieser grausamen Plage, so über die Egypter ergeht, dennoch der
\quellseite{443}
Israeliter Erstgeburt verschont,\footnote{Eislebensche: wird.} und im Zorne der Gnade und Barmherzigkeit noch eingedenkt ist [Hab. 3, 2.]. Darum wir uns alles Gutes zu GOtt zu versehen haben.

127. Letztlich bestätigt dies Wunderzeichen auch GOttes Wahrheit [Ps. 111, 7. ff.]. Denn er hatte verheißen, er wollte die Erstgeburt der Egypter todtschlagen, und die Kinder Israel aus dem Diensthause Egypti führen. Dieses geschieht jetzt beides, daß in des Königs Hause, und in des geringsten Kuhhirten Hause, Todte an Menschen und Vieh gefunden werden, und die Egypter mit Gewalt die Kinder Israel aus ihrem Lande selbst stoßen und treiben.

128. Diese Historie haben wir gehört, nun wollen wir auf die Allegorie oder heimliche Deutung auch kommen. Aber, so du willst folgen den Deutungen, und des Gewissen spielen von der Erlösung der Kinder Israel, so nimm endlich die Juden, darnach die Christen, an denen fein und herrlich gesehen wird diese Erlösung. Demnach mag man auch es auf diejenigen ziehen, die es gleich haben. Als, da im neuen Testamente viel Christen gemacht wurden, da wollten die Juden recht haben, so wollten die Apostel auch recht haben. Die Juden wollten den Preis allein und mit Gewalt behalten, daß sie GOttes Volk wären; sie hatten es aber versiehen, und sich selbst deß entsetzt und verlustig gemacht. Davon handelt nun der HErr Christus auch und spricht: „Die Ersten werden die Letzten werden, und sollen die Letzten die Ersten sein“ [Matth. 20, 16. 19, 30. Marc. 10, 31. Luc. 13, 30.]. Alhier steht zweierlei Volk: Die Juden, als die Ersten, hatten GOttes Wort und Zeugniß der Propheten, daß sein Volk auf Erden ihnen sollte vorgezogen werden; die stolziren auf ihre Erstgeburt, werden aber todtgeschlagen. Aber darnach ist ein Theil Volks, so da folgt den Aposteln, welche Christum predigen, und das sind die, so auf Christum getauft werden und selig sind; dieselbigen Christen und Apostel schweigen ihrer Erstgeburt und werden erhalten. Da erhebt sich über der Erstgeburt ein Kampf und Streit.

129. Also rühmt sich der Papst auch wider uns von seiner Erstgeburt, daß er von den Aposteln herkomme, und habe seine Lehre von den alten Vätern und von so viel Conciliien; aber wir stehen darauf und sagen: Unsere Lehre ist recht, eure ist falsch, denn sie ist Menschentand, und kann aus GOttes Wort nicht bewiesen werden. Darum wollen die Verheiligen auch allezeit einen Vorzug haben, sie wollen in der Welt die Ersten und Besten sein, die Vornehmsten, gehen empor. Und wenn sie denn deß gar gewiß sein wollen, so fähet der Heilige Geist ein neu Wesen an, das verdrießt sie denn; gleichwie jetzt auch geschiehet. Sie wollen die Ersten sein, und für die höchsten Heiligen gehalten werden; aber da bleibt es nicht lange nach (wie denn jetzt vor Augen), daß sie die größten Sünder und ärgsten Heuchler erfunden werden, und welche sie für Ketzer halten, werden alsdann für Heilige geachtet. Dies sind GOttes Werke, welche man an den Aposteln und Juden sieht, und auch an uns selbst noch auf diesen heutigen Tag und gegenwärtige Stunde es sich findet.

130. Durch den Pharaonem und die Egypter verstehe die Juden, welche zur Zeit der Apostel gelebt haben, da das Evangelium anging und hervorbrach. Und die Namen stimmen fein damit überein. Denn Pharao heißt barhäuptig, da eines Haupt bloß und nicht bedeckt ist, das ist, der da los steht, da Mißbrauch ist, der nichts zu schaffen hat. Auf welche Meinung der Apostel zu den Galatern am 5. Capitel, V. 2–4., ohne Zweifel mit diesen Worten sieht: Wenn ihr wollet durch die Werke selig werden, so seid ihr aus der Gnade entfallen, und Christus ist euch kein nütze. Als sollte er sagen: Er geht euer müßig, und ihr sein müßig. Denn er soll unser Haupt sein, und wir seine Glieder, und er soll in uns schweben; es soll nicht ein Leib oder Kloß sein, so kein Haupt habe. Die Christenheit soll nicht ein Pharao sein, das ist, der da ledig und bloß hergeht, und niemand unterthan ist.

131. Dieser Name reimt sich nun fein auf die Juden. Die wollten ein frei Volk und Sonderling sein, haben auswendig ein scheinbarliches Leben geführt, und sich gar herausgezogen aus dem Gehorsam des Evangelii, meinten, sie gefielen GOtt eben so wohl, wenn sie auswendig viel opferten, und den Glauben ließen anstehen inwendig, dadurch man doch GOtt allein ehrt. Darum geht GOtt ihrer müßig, und sie seiner auch, und ist eines dem andern entfallen. Sie sind frei ledig von GOtt,
\quellseite{444}
und bloß, in welchen er nicht herrscht durch den Glauben, welche doch waren die Erstgeburt, die lieben Junker, ein heilig Volk, und sollten das Erbe sein. Also ist es den Juden gegangen zur Zeit der Apostel.

132. Nun thut Christus einen Gang, er geht aus diesem Leben zum Vater, von welchem Gange er in Johanne dem Evangelisten [Cap. 16.], wie angezeigt, viel sagt. Durch diesen Gang hat er alle Erstgeburt todtgeschlagen, das ist, diese Ehre hat er ihnen genommen, daß sie nicht sollen die Ersten sein, sondern verdammt sein, und wissen, daß sie nicht GOttes Volk wären; diese Erstgeburt solle nun todt sein. Es stand also bis auf Christum; aber da diese Nacht kommt, da sind die Pharaones alle todt; da Christus von dem Tode aufersteht, mit demselbigen geht alle Herrlichkeit und Pracht der Juden zu Boden. Jetzt ist es alles todt und soll nicht gelten. Mosis Werk und alle äußerliche Heiligkeit liegt danieder, er will sie nicht haben. Es soll nichts helfen, denn allein glauben an Christum.

133. Solches bedeutet GOttes Durchgang in Egypten, da er alle Erstgeburt erschlagen hat; das ist, Christus ist von dem Tode auferstanden, durch diesen Gang hat er alles getödtet, das da jüdisch ist, auf daß es hinführet nichts gelte. Was jüdisch ist, das ist, alles, was an ihm hat eine äußerliche Weise, und führt solche Sitten, daran man hängt, das muß geistlich todtgeschlagen werden. Das geschieht nun durch den Heiligen Geist, der gibt die Kraft, daß diese äußerlichen Dinge nichts gelten.

134. Also sagen wir auch, der Papst habe nicht anders gewußt, läßt ihm auch noch anders nicht einreden, denn daß er neben den Seinen die Erstgeburt inne habe. Er hält sich für die Erstgeburt, dagegen andere arme Christen und Laien, oder der gemeine Mann, seien\footnote{Gießebensche: sei.} gar nichts, sondern sie müssen durch der Geistlichen Hülfe gen Himmel kommen. Daher ist der Papst Junker, gleichwie auch Pharao gewesen. Nun kommt das Evangelium gleich als ein Donnerschlag, und schlägt es alles zu Boden und tödtet es, sagt, daß es alles nichts sei [Ps. 144, 6.], item, daß es auch todt sei, und wer darauf vertrauet, der sei auch erschlagen und todt. Das ist die Schlacht und der Durchgang, den Christus mit seiner Auferstehung angerichtet hat. Und diese Schlacht thut er noch mit seinem lieben Evangelio. Denn, GOtt Lob und Dank, ihrer sind jetzt allenthalben gar viel, die des Papstes Irrthum erkennen.

135. Es finden sich aber auch solche, die da wollen alles unterdrücken. Das denn daher kommt, daß der Teufel etliche tolle Leute erweckt, die alles erschlagen wollen, Juden aus uns machen, alles ermorden; so man doch geistlich nur tödten soll, als, mit dem Munde predigen, daß das Gewissen erleuchtet werde. Wenn das Wort GOttes gepredigt wird, so hat man recht getödtet. Denn die Herzen wissen es sonst nicht; gleichwie man jetzt nicht weiß, warum man Klosterleute zu Christen machen will. Der Teufel hat gefühlt, daß der Papst solle erschlagen werden, und sterben müsse; solches hat er mit diesem Aufruhr hindern wollen, und macht damit dem Evangelio ein böse Geschrei, daß man jagete: Siehe, sind das die guten Evangelischen, so da jetzt die Klöster stürmen, die Schlösser verbrennen etc.? Aber wir sagen, daß man also lange predigen solle, bis daß GOtt sie tödte, auf daß es ein göttlich Werk sei, daß GOtt durch das Evangelium in das Herz komme, und erschlage alle Heiligkeit, Frömmigkeit und Gerechtigkeit solcher Menschen.

\section{Das dreizehnte Capitel}

\begin{schriftzitat}
V. 1–4. Und der HErr redete mit Mose und sprach: Heilige mir alle Erstgeburt, die allerlei Mutter bricht bei den Kindern Israel, beide unter den Menschen und dem Vieh, denn sie sind mein. Da sprach Mose zum Volk: Gedenket an diesen Tag, an dem ihr aus Egypten, aus dem Diensthause, gegangen seid, daß der HErr euch mit mächtiger Hand von hinnen hat ausgeführt; darum so sollst du nicht Sauerteig essen. Heute seid ihr ausgegangen, in dem Monden Abib.
\end{schriftzitat}

\quellseite{445}
1. Nächst haben wir gehört, zum Theil, was die Erstgeburt sei, nämlich, daß es die Juden sind, die mit ihrer äußerlichen Heiligkeit obenan saßen, und die Vornehmsten auf Erden waren; auch wie dieselbigen geistlich geschlagen werden müssen, durch das Schwert des Evangelii, sollen sie anders mit dem HErrn Christo gen Himmel fahren, damit all ihr Wesen und Gesetz aus sei, und müssen die Letzten werden, ja, die Allerverachtesten auf Erden; dagegen sind die Heiden die Ersten worden. Das andere Theil dieses Textes ist von der Erstgeburt, die GOtt hieß heiligen, und nicht tödten. Welches bedeutet diejenigen, so von den Juden und Heiden zum christlichen Glauben kommen würden. Also hat man die Erstgeburt getödtet; aber die andere Erstgeburt bleibt lebendig, und wird behalten. Und ist in dem das Evangelium erfüllt, da der HErr Christus [Matth. 20, 16.] spricht: Die Ersten werden die Letzten, und die Letzten die Ersten.

2. Dergleichen singt Moses [5 Mos. 32, 21.] in seinem Liede: „Ich will sie wieder reizen an dem, das nicht ein Volk ist, an einem närrischen Volke will ich sie erzürnen; denn sie haben mich erzürnet und zum Eifer gereizet.“ Als sollte er sagen: Gleichwie ihr mich, den wahrhaftigen GOtt, verlaßt, daß ich nicht muß euer GOtt sein, und erwählt andere Götter: also will ich euch, so ihr mein Volk gewesen seid, auch verlassen, und diejenigen annehmen und erwählen zu meinem Eigenthum, welche zuvor nicht mein Volk waren [Hos. 2, 23.], welches euch denn verdrießen wird. Wie denn auf diesen heutigen Tag die Juden noch darum zürnen; sie sind dahin mit aller ihrer Herrlichkeit. Den Hader hat GOtt angerichtet. Sie sind keinem Volke GOttes günstig, denn allein ihnen selbst; aber es ist aus, es gilt nicht mehr Juden sein, sondern Christen sein [Eph. 2, 13. 19. Col. 3, 11.].

3. Also mag man es auch noch auf alle deuten, die solcher ihrer Art sind. Denn die Papisten wollen auch noch heutiges Tages GOttes Volk sein. Um den Namen ist es zu thun, daß ein jeder ihm diesen Namen zuschreibt, und wollen dieses von sich gerühmt haben: Invocatum est nomen tuum super nos [wir sind ja nach deinem Namen genannt]. Dann müssen die wahrhaftigen Christen und Volk GOttes, oder die Erstlinge und Erstgebornen, verachtet werden und unterliegen, auch Ketzer gescholten werden, aber die andern müssen es gar sein, und die christliche Kirche heißen, und den Heiligen Geist haben. Dieses kann niemand scheiden, denn allein GOtt mit der Offenbarung seines Evangelii, der schlägt unter sie, und urtheilet, welches sind die wahrhaftigen oder falschen Christen.

4. Solches wird wohl also bleiben bis an das Ende der Welt. Denn die Welt will schön und heilig sein, und gerühmt werden, daß sie des Himmels und der ewigen Seligkeit würdig sind, und verfolgen drüber andere Leute, daß sie wollen die Kirche sein; müssen also die wahrhaftigen Christen die Letzten heißen. Aber unser lieber HErr GOtt hält dieselbigen, so in der Welt also verachtet sind, für die Ersten. Und dieses ist uns zum Exempel vorgeschrieben, auf daß wir wissen, es müsse also gehen alle denen, die da wollen die Vornehmsten und Ersten sein, und diejenigen auch getröstet werden möchten, die da eine Zeitlang geplagt und verfolgt werden um des göttlichen Worts und christlichen Namens willen. Wie ist aber das Heiligen zugegangen? Der Text spricht:

\begin{schriftzitat}
V. 11–13. Wenn dich nun der HErr in das Land der Cananiter gebracht hat, wie er dir und deinen Vätern geschworen hat, und dir es gegeben; so sollst du aussondern dem HErrn alles, was die Mutter bricht, und Erstgeburt unter dem Vieh, das ein Männlein ist. Die Erstgeburt vom Esel sollst du lösen mit einem Schafe; wo du es aber nicht lösest, so brich ihm das Genick. Aber alle erste Menschgeburt unter deinen Kindern sollst du lösen.
\end{schriftzitat}

5. Droben, im 12. Capitel, haben wir gehört, wie GOtt ein Gesetz und Gebot vom Osterlämmlin jährlich zu essen gegeben und gestiftet hat, auf daß man der herrlichen Liberation der Kinder Israel aus Egypten nicht vergesse.

6. Nun folgt das andere Gesetz, von der Opferung oder Heiligung der ersten Geburt, zum Gedächtniß, daß der Israeliten erste Geburt in Egypten wäre erhalten worden. Das Geheiligte ward dem Priester gegeben, es war vom Vieh oder sonst von Früchten; das nahm unser HErr GOtt zu sich, und es ward den Priestern zugestellt. Dafür, daß sie das Gesetz lehrten und GOttes Wort treiben mußten, sollte die Erstgeburt ihre Nahrung sein; denn sonst hatte der Stamm Levi in der Austheilung des cananäischen Landes keine Landgüter bekommen, wie
\quellseite{446}
andere Stämme. Wenn ein Knäblein geboren ward, so gehörte es auch dem HErrn; aber er ließ zu, daß man es lösete mit einem Ort des Guldens.\footnote{Ort des Guldens = ein Viertelgulden.} Und das war zu thun um die Prediger des göttlichen Worts, damit sie das Volk mit dem göttlichen Worte versorgten. Aber wie man sonst viel Gesetze und Ordnungen GOttes unterließ, also hielten sie dies Gebot GOttes auch nicht, und wurden die Priester verachtet; man gab ihnen nichts, daß sie mußten Hunger leiden. Darum so suchten auch die Priester falsche Gottesdienste, und gingen auch an Abgötterei [Mal. 1, 7. 8.], da gab denn das Volk Geldes und Guts genug zu.

7. Gleichwie man vorzeiten im Papstthum auch viel Almosen gegeben, und die Priester reich gemacht hat, und das Geld für den Ablaß weggeschleudert. Denn wenn man falsche Propheten in der Welt hat gehabt, welche die Leute mit der Nase umgeführt, und uns um Leib und Seele, auch um das Gut gebracht, denen hat man Geldes genug gegeben. Die Königin Isabel konnte bei 800 Propheten von ihrem Tische herrlich speisen und ernähren, da Elias dagegen in der Flucht und Exilio sein mußte, und andere hundert Propheten von Abdia [Obadia], des Königes Kämmerer, in einer Höhle versteckt, mit Wasser und Brod gespeist worden [1 Kön. 18, 4.]. Dieses läßt sich ärgerlich und thöricht vor der Vernunft ansehen, aber ich urtheile wohl selbst also, daß GOtt daran närrisch thäte, daß er sein Evangelium also dahin schlendert, und seine Prediger in der Welt die höchste Armuth dabei leiden müssen.

8. Das geistliche Heiligen oder Opfern bedeutet, daß man soll im Amte des Predigens sein. Denn wer da will Deutung führen, der ziehe sie den mehreren Theil auf das Predigtamt, denn darum ist es gar zu thun. Denn GOtt will durch das Evangelium alles ausrichten; [aus] der Ursache so soll sich auch alles dahin richten und gezogen werden. Wenn ich predige, so soll ich anders nichts denn den Glauben an Christum predigen; denn der gilt allein vor GOtt. Denn so schlage ich die alte Erstgeburt todt, und die andere Erstgeburt heilige ich GOtt; das ist, diese Erstgeburt ist ihm angenehm, die nimmt er an, sie ist GOtt zugeeignet, und unterrichtet die Leute von GOtt. Darnach [soll ich] sagen, daß die GOttes Kinder sind, die glauben von Herzen an Christum, der ist auch der Erstgeborne vor allen Creaturen,\footnote{Eislebensche: aller Creaturen.} wie St. Paulus [Col. 1, 15.] sagt, und macht uns auch zu Erstgebornen und Kindern GOttes. Und ob diese wohl die Letzten in der Welt sind, und andere für die Obersten und Vornehmsten gehalten werden, so ist doch GOttes Urtheil da, und hält sie für selig, die andern aber für Verdammte.

9. Wer nun das Evangelium annimmt und ihm folgt, derselbige opfert mit, daß man von Herzen und mit dem Munde bekenne, daß diese GOttes Kinder sind, die da glauben an Christum. Also sollen alle Opfer dahin gerichtet sein, daß man GOtt lobe und ehre. Es soll alles Lobopfer sein. Denn im 50. Psalm, V. 23., wird gesagt: „Wer mir Dankopfer bringt, der preiset mich.“ Das ist das allerbeste Opfer, das GOtt geschehen mag. Darum gebet Gott hin und wieder in den Psalmen, man solle ihm Lobopfer bringen. Das reucht ihm sehr wohl an. Wenn wir predigen und bekennen, daß unsere Werke nichts sind, sondern wir müssen unseres Herrn Gottes Gnadenwerke, seine Gunst, Güte und Barmherzigkeit haben, das ist opfern. Wenn ich dies mit dem Munde bekenne, dann opfere ich mit. Das geht nicht mit Werke an, sondern daß man glaubt mit dem Herzen, daß unsere Werke nichts sind, aber GOttes Barmherzigkeit und Gnade, in Christo uns geschenkt, die thue es. Wenn ich nun GOtt in der Welt also preise, und bekenne ihn vor den Leuten, so bin ich ein Priester und opfere das schönste Opfer.

10. Das ist nun die Erstgeburt opfern und heiligen. Welche nun auf ihre Werke trauen, und auf ihr gutes, heiliges Leben sich verlassen, dieselbigen preisen und loben sich selbst, und richten einen Abgott auf, gleich als wollte GOtt durch ihre Werke gelobt sein, und schreiben den Werken zu, gleich als könnten dieselbigen ihnen helfen, welche Ehre allein dem HErrn Christo gebührt. Nun ehren sie mit dem Munde den HErrn Christum wohl, aber inwendig im Herzen lästern und schänden sie ihn [Jes. 29, 13. Matth. 15, 8. Marc. 7, 6.]. Das heißt nun, GOtt die Erstgeburt zueignen, heiligen oder opfern.

11. Was ist aber das Recht der Erstgeburt? Sonst gehörten dem erstgebornen Sohne zwei Theile des väterlichen Erbes, zweimal so viel,
\quellseite{447}
als einem andern Kinde. Danach gebührte ihm auch das Regiment, die Herrschaft und das Priesterthum, daß er im Hause war Priester und Prediger; daß meinte Gut, die größte Macht und größte Ehre; er war König und Priester, Kaiser und Papst.

12. Solche Güter haben wir Christen allzumal geistlich. Wir haben auch zwei Theile. Erstlich sind wir Priester, daß wir genug haben an der Seele, inwendig, daß wir der Seelen Heil und Seligkeit erlangt durch Christum, predigen und bekennen ihn, loben, preisen und rufen ihn an in allen Nöthen. Danach haben wir die Verheißung, daß wir nicht sollen Hungers sterben; und ob gleich zuweilen nicht Baarschaft oder Vorrath vorhanden ist, so müßten ehe die Raben geflogen kommen und Speise bringen, wie dem Elia geschah [1 Kön. 17, 6.], oder vom Himmel Brod regnen, als in der Wüste Manna gefiel [2 Mos. 16, 4.], wenn wir allein glaubten.

13. Ein christlich Herz hat viel zu schaffen und zu leiden [Ps. 34, 20.], aber es ist ohne Sorge, also gewiß und sicher, gleich als wenn es die ganze Welt voller Güter hätte [Ps. 73, 25. 26.]. Denn es gedenkt: Mein GOtt kann allezeit mehr schaffen, denn in der Welt ist. Wenn man ihm gleich die Güter wegnimmt, so kann man ihm doch nicht GOtt nehmen. Bedarf er nachmals etwas, so kann ihm GOtt wohl mehr geben, denn aller Welt Güter sind sein. Aber es plagt und martert uns Christen so gar sehr diese Bauchsorge, oder die zeitliche Nahrung, daß wir diesen Trost nicht fest genug fassen. Darum so ist ein Christ auch ein König, das ist, ein Christ ist ein Herr über alles, und alle Creaturen müssen ihm gehorsam sein; aber geistlich.

14. Auswendig kann ihm das Leben und die Güter genommen werden, aber er arbeitet und thut, was ihm befohlen ist. Und ob ihm wohl alles genommen wird, so bessert er sich doch allezeit davon, und sein Glaube wächst immerdar und mehret sich, und regieret also im Herzen, daß ihn weder Reichthum noch Armuth bekümmern, weder trauriger noch fröhlicher machen. Wenn er reich ist, Geld und Gut hat, so ist er nichts fröhlicher; er ist auch nichts trauriger, wenn er arm ist [Phil. 4, 12.]. Schilt oder lobt man ihn, so gilt es ihm gleich viel. Das macht der Muth, den er hat; das ist der freudige, sichere Glaube, damit er GOtt anhängt.

15. Nun sage mir einer, welche Herrschaft besser sei, daß einer mit dem Schwerte fechte und streite, oder daß einer habe einen solchen Muth? Ob du gleich mit dem Schwerte einen überwindest, so thut er, was du willst; aber ein christlich Herz ist ein solcher Herr, welchen niemand ändern noch betrüben kann; er bleibt allezeit für sich. Solche Herren macht er aus uns (wie St. Paulus Röm. 8, 35–39. auch sagt: „Wer will uns scheiden von der Liebe GOttes? Weder Hohes oder Niedriges, Fürstenthum oder Gewaltigen“ etc.), daß uns nichts kann bekümmern. Aber wo werden diese Erstgebornen gefunden, und wo sind solche Christen? St. Paulus ist derselbigen einer gewesen.

16. Zum dritten so ist er auch ein Priester; er hat die Gewalt zu predigen GOttes Wort, und vor GOtt zu beten, das unverständige Volk zu lehren, geht hinauf zu GOtt durch das Gebet. Denn wir sind deß gewiß, wenn wir getauft sind, daß unser Gebet angenehm sei; und wenn wir alle beten, so sind wir alle Priester, können alle GOttes Wort lehren, so vom Himmel uns verkündigt ist. Das ist Gottes Wort, daß ein jeglicher predige, wenn es noth ist. Das ist nun die Figur und Bedeutung derselbigen Herrlichkeit. Die Christen sind Erstgeborne im Geiste; jene haben die leibliche Herrlichkeit der Erstgeburt gehabt.

17. Nun erhebt sich die Frage, davon wir zuvor auch etwas gesagt: Weil gleichwohl weder hier im Texte, oder anderswo nichts gemeldet wird vom Wiedergeben des entlehnten Guts von den Egyptern, ob die Israeliter auch recht gethan haben, daß sie den Egyptern ihr Silbergeräth genommen? Haben sie hiermit nicht einen Diebstahl begangen? Sie haben es zu ihrem Feste geborgt und entlehnt, und da man es ihnen geliehen hat, fahren sie zu in der Nacht, und nehmen es hinweg. Haben sie Recht dazu, oder ist es ein Diebstahl, oder haben sie es geraubt? Hat GOtt nicht in den zehen Geboten befohlen, man soll nicht stehlen? 2 Mos. [Cap. 11, § 4 f.] ist auch geantwortet: Solches heißt nicht gestohlen, sondern ihnen von GOtt gegeben, aus dieser Ursache: Unseres HErrn GOttes ist alles. Nun hat er ihnen vom Himmel befohlen, sie sollten es von den Egyptern entlehnen, und mit sich nehmen. Stehlen ist, wenn ich einem andern das Seine nehme, wenn ich einem andern etwas nehme, das er mir nicht gibt, sondern ich entwende oder entfremde es
\quellseite{448}
heimlich. GOtt hat den Egyptiern das Silbergeschirr gegeben; nun fährt er zu und nimmt es ihnen wieder, und gibt es den Israeliten; er borgt es ihnen ab, und täuscht sie und spricht: Also will ich es haben. Vielleicht darum, daß sie schwere Arbeit in Egypten gethan hatten, und gar keinen Lohn dafür bekommen. Aber GOtt hat es an ein Kerbholz geschnitten, und wollte ihre Arbeit bezahlen.

18. Das ist also geschehen nach der Historie. Aber was bedeutet es im neuen Testamente? Es soll eine Anzeigung sein der Welt Wohlbered(en)heit und Weisheit, daß kluge Leute mit Sprüchen nehmen und ziehen die Christen zu sich etc. Nun, GOtt gebe uns nicht viel weltliche Weisheit. Wir sollen lassen bleiben im äußerlichen Weltreiche die schönen Sprüche und Lehren der heidnischen Scribenten, und scheiden von einander das Reich GOttes, darinnen Christen sind, und das äußerliche Weltreich, und sagen: Zu dem Reiche GOttes gehört nichts von der weltlichen Weisheit. Es regiert und herrscht allhier anders nichts, denn allein das reine, lautere Wort GOttes.

19. Sind sonst andere weise und kluge Sprüche in der Vernunft, dieselbigen zeuch heraus in dies äußerliche Leben; da geht es hin, daß einer mehr Vernunft und Verstand habe denn der andere. Darum so haben wir viel hübsche Sprüche in allen Zungen und Sprachen, zu diesem äußerlichen Leben und Wesen. Also ist einer auch besser und geschickter denn der andere, ein Regiment zu führen, dieweil er hat eine höhere Vernunft und Verstand, welches ist die weltliche Weisheit, die denn GOtt auch schenkt und gibt. Davon mag ein anderer auch ein Exempel nehmen, Lehre und Sprüche von ihm borgen; nach welchem Exempel seine Lande und Leute auch im Frieden bleiben mögen, und sich daraus beßern; wie denn die Heiden zur Besserung der Polizei oder des Weltregiments gethan haben, davon die Historien zeugen.

20. Auf solche Vernunft oder Verstand reimen sich fein die Historien, denn sie haben alles aus der Erfahrung; wie denn auch die Deutschen im Sprüchworte sagen: Guter Muth ist halber Leib. Dieses ist ein feiner Spruch, aus der Vernunft gezogen, und dient dazu, wenn einer betrübt ist, daß man zu ihm sage: Ei, nimm dir einen guten Muth vor, so kommst du herdurch. Aber es ist eine äußerliche Tröstung, da aus der Vernunft einer den andern tröstet; welches denn auch eine Gabe und Geschenk GOttes ist, und man soll es erkennen, daß solches von GOtt kommt; wie er sonst auch Reichthum gibt. Also ist es auch eine sonderliche Kunst, Krieg zu führen; item, Pferde zu bereiten; es sind alles Gaben GOttes, die er in die Rapuse geworfen hat, damit wir uns regieren sollen. Aber in den Dingen, die über uns sind, da die Conscienzien soll regiert werden, da gelten diese Sprüche ganz und gar nichts; sondern da muß allein regieren das lautere, helle und reine Wort GOttes.

21. Derhalben so gilt diese Bedeutung nichts, daß man mit heidnischen Sprüchen die christliche Kirche beschweret hat, ja, dadurch die Kirche gar zerstört, wie denn die Papisten auch gethan haben. Deß dancke ihnen aber der Teufel; denn das heißt mit der menschlichen Weisheit die Christenheit zu Boden gestoßen. Denn Weltweisheit kann und weiß nicht mehr, denn wie man mag Frieden auf Erden haben; aber GOtt will haben, daß du ein Christ seiest, und wissest, wie du einen gnädigen GOtt und das ewige Leben haben mögest, welches du aus der Vernunft nicht lernen wirst, sondern GOttes Wort lehrt solches.

22. Ich habe oft davon gesagt, und wollte es ja mächtig gerne, daß wir von einander scheideten diese zwei Reiche. Denn die Vernunft, wie schön und herrlich sie auch ist, so gehört sie doch in das Weltreich allein, da hat sie ihre Herrschaft und Gebiet. Aber im Reiche Christi, da hat allein GOttes Wort die Oberhand. Gleichwohl will der Teufel immerdar mit der Vernunft in der Kirche regieren, die Hand im Sode haben, und mit heidnischen schönen Sprüchen und Anschlägen herrschen. Aber das Gold und Silber, die schönen Kleinode, die in Egypten geraubt sind, ist die heilige Schrift. Denn alsbald die Apostel sich von den Juden absonderten, und sie gar neue Leute wurden und GOttes Volk, da hatten die Juden verloren die heilige Schrift; den Buchstaben und das Pergament haben sie wohl behalten, aber den rechten Verstand und Meinung oder Sinn haben sie nicht [2 Cor. 3, 3. f.].

23. Darum, bei den wahrhaftigen Israeliten da bleibt die heilige Schrift, und jene sind es ihnen schuldig gewesen. Da haben wir keine Sprüche, hübsche Historien und Exempel, die wir gebrauchen zum ewigen Leben; da werden
\quellseite{449}
wir durch Sprüche sein unterrichtet; item, sehen, wie GOtt den Glauben in den Christen gewirkt hat. Denn an dem Glauben ist alles gelegen, darauf steht das ganze christliche Leben. Dies ist das Gold und die köstlichen Geschmeide, so man den Egyptern geraubt und entwendet hat, gleichwie wir jetzt den hohen Schulen die heilige Schrift auch gestohlen haben. Ob sie sich wohl mit vielen Büchern tragen, so haben wir doch den reinen Verstand der Schrift, und schmücken uns damit, das ist, stärken und bessern uns im Glauben durch solche Sprüche. Im Propheten Ezechiel [Cap. 16, 10. 13.] spricht GOtt auch: „Ich habe dir mein Gold und Seide gegeben“, das ist, meine heilige Schrift, und den rechten Verstand, den ich gegeben habe durch die Propheten, den hast du verkehrt, und eine mächtigere Maie daran gemacht, nach alle deinem Willen.

\begin{schriftzitat}
V. 18. Und die Kinder Israel zogen gewappnet aus Egyptenland.
\end{schriftzitat}

24. Sie hatten auch Harnisch an. Was war es vermöthen, da sie an Kriegsrüstung und Macht viel geringer und schwächer waren denn der König in Egypten, und GOtt wollte sie auch durch seine Kraft ausführen, daß er allein den Ruhm davon hätte? Noch heißt er sie, sich rüsten, und ihren Kriegshaufen zurichten, daß denn ein Ansehen hat, als wollten sie sich mit den Heiden schlagen; aber es ist darum geschehen, daß wir GOtt nicht versuchen. Die Mittel sind GOttes Larven und Mummerei, darinnen er läuft auf Erden. Er will mir helfen in allen Dingen, an Leib und Seele, und daß ich allein auf sein Wort vertraue; dennoch will er auch, daß ich das Meine dazu thun soll, Roß, Spieß, Schwert und Wagen haben, Heerhaufen machen, gleich als wollte man durch das Schwert überwinden, so es doch GOttes Meinung nicht war. David im 20. Psalm, V. 8., sagt: Ich verlasse mich nicht auf mein Roß, Schwert oder Bogen \&c.

25. Wozu dient denn dieses alles? Ei, unter dieser Rüstung und Kriegswehre, da will GOtt bei dir sein und sich darunter verbergen, daß andere Leute denken möchten, du würdest es ausrichten mit deiner Kriegsrüstung und eigener Macht, so es doch GOtt allein thut [Psalm 60, 14. Ps. 144, 1. ff.]. Also thut GOtt auch mit andern Dingen. Er heißt uns beten, arbeiten \&c., welches alles nur ein lauter Spiegelfechten ist. Denn wo GOtt nicht segnete, so möchte nicht ein Haar, noch ein einiger Strohhalm, sondern es wäre mit allem Dinge aus. Dennoch will er haben, daß ich mich also stelle: dann, wo ich nicht pflügte oder säete, daß ich gar nichts würde haben.\footnote{Der Sinn des letzten Satzes wird sein: GOtt will, daß ich mich so stelle, als ob ich dann gar nichts haben würde, wenn ich nicht pflügte oder säete.}

26. Denn ohne meine Arbeit will er es nicht lassen geschehen, und dennoch soll es nicht durch meine Arbeit gethan sein. Er will nicht haben, daß ich daheim sitze, faulenze und es dem lieben GOtt befehle, und warte, bis mir ein gebraten Huhn in das Maul fliege. Das ist GOtt versuchen, daß er mir sollte lassen etwas wachsen, da sonst nichts wächst; das will er nicht haben. Darum ist es wohl geredet, daß die Alten gesagt: Dat Deus omne bonum, sed non per cornua taurum. Und also sollen wir uns halten im Herzen, daß wir wissen, es werde nichts daraus, wenn er es nicht thue, auf daß GOtt Raum habe, und er eine Hütte bei dir kriege, darinnen er sich verberge. Also haben sie ihrer Kriegsrüstung gebraucht, gleich als wären sie gar wohl geschickt zum Streit, und gar kriegerisch. Als wollte GOtt sagen: Ich will euch gebrauchen zum Schein allein in dieser Sache.

\begin{schriftzitat}
V. 20–22. Also zogen sie aus von Suchoth, und schlugen Gezelt[e] auf in Etham, vorne an der Wüste. Und der HErr zog vor ihnen her, des Tages in einer Wolkensäule, daß er sie den rechten Weg führete, und des Nachts in einer Feuersäule, daß er ihnen leuchtete, zu wandeln Tag und Nacht. Die Wolkensäule und Feuersäule wichen nimmer von dem Volke.
\end{schriftzitat}

27. Dieses ist auch ein herrliches Wunderwerk, daß ihnen GOtt des Tages eine weiße Wolkensäule, und des Nachts eine Feuersäule am Himmel gestellt hat, darnach die Kinder Israel auf der Reise sich richten sollten. Dieses Mirakels gedenkt die heilige Schrift sonst oft, als, 4 Mose am vierzehnten Capitel, V. 14., auch Nehemia am neunten Capitel, V. 12. Denn GOtt wollte damit zu verstehen geben, daß er gegenwärtig bei seinem Volke sein wollte, sie vor ihren Widerwärtigen schützen, und mit Rath und That nicht verlassen in allerlei Gefährlichkeiten, sondern ein Auskommen in der Trübsal
\quellseite{450}
geben, daß man es ertragen könne [1 Cor. 10, 13.]. 28. Wir Christen wandeln auch aus Egypten, aus dem Reich des Teufels und der Sünden, nach dem gelobten Lande, das ist, nach dem ewigen Leben; und kommen in die Wüste, das ist, in allerlei Trübsal, Noth und Anliegen; da mußten wir den Weg nicht zu treffen, daß wir heraußer kämen, und dieselbigen überwinden könnten. Aber GOtt gibt uns Geleitsleute, als, des Tages die Wolke, und des Nachts die Feuersäule, das ist, die Predigt des göttlichen Worts und den Brauch der hochwürdigen Sacramente; die leuchten und scheinen uns vor, daß wir den Weg in der Wüste treffen und finden mögen zum ewigen Leben.

\section{Das vierzehnte Capitel}

\begin{schriftzitat}
V. 1–4. Und der HErr redete mit Mose, und sprach: Rede mit den Kindern Israel, und sprich, daß sie sich herum kehren, und ihr Gezelt aufschlagen gegen dem Thal Hiroth, zwischen Migdol und dem Meer, gegen Baal-Zephon, und daselbst gegenüber das Gezelt aufschlagen an das Meer. Denn Pharao wird sagen von den Kindern von Israel: Sie wissen nicht, wo aus im Lande, die Wüste hat sie beschlossen. Und ich will sein Herz verstocken, daß er ihnen nachjage, und will an Pharao und an aller seiner Macht Ehre einlegen; und die Egypter sollen inne werden, daß ich der HErr bin. Und sie thaten also etc.
\end{schriftzitat}

1. Bisher sind erzählt worden etliche Plagen, so GOtt über Egypten hat ergehen lassen um der Kinder von Israel willen. Darauf denn gefolgt ist, im 12. Capitel, die zehnte Plage, welche trifft die Erstgebornen, und ist die letzte Plage, darauf alsbald die Kinder von Israel sind aus Egypten gezogen. Aber, auf daß dieselbige, neben den andern Plagen, und auch die Erlösung aus Egypten, in der Welt nicht vergessen, sondern ewiglich gepredigt und jedermann bekannt würde, so hat GOtt das Osterfest eingesetzt, und allerlei äußerliche Weise und Bräuche oder Ceremonien an diesem Feste zu halten geboten. Aber wenn wir dieselbigen auswendig ansehen, so scheint es als eitel Thorheit. Denn es lautet seltsam, daß sie das Osterlamm mit ungesäuertem Brode essen mußten, und aufgeschürzt dazu stehen, gleich als Wandersleute. Man kann die Ursache nicht ausrechnen, warum es GOtt also befohlen habe. Es hat aber ohne Zweifel GOtt durch diese sonderliche Weise binden wollen der Nachkömmlinge Augen und Ohren an diesen Auszug der Kinder von Israel aus Egypten.

2. Gleichwie wir Christen auch den Befehl haben, daß wir das Gedächtniß unsers lieben HErrn und Heilands Jesu Christi predigen, davon rühmen, singen und sagen, und von dem herrlichen Werke immerdar und ewiglich predigen und reden, dichten und schreiben sollen, daß er der Sohn GOttes, von Maria geboren, gekreuziget und gestorben sei, und von dem Tode wieder auferstanden am dritten Tage, auf daß er uns von dem Tode, Sünde, Hölle und des Teufels Gewalt erlösete [1 Cor. 1, 23. ff.], und daß dieses Gedächtniß bleibe in der Welt, so hat der HErr Christus das hochwürdige Sacrament des Altars eingesetzt, eine äußerliche Ceremonie und Weise, daß man im Brod und Wein den Leib und Blut des HErrn Christi empfahen, daran wir uns üben, davon predigen, singen und sagen sollten, und seines Leidens, Sterbens und Auferstehung dabei gedenken. Bei den Sacramenten will er ergriffen und gefunden werden, und bei uns sein.

3. Also ist auch die Taufe ein gewiß Zeichen, dabei man den wahrhaftigen GOtt, der Himmel und Erde geschaffen hat, und des Sohn für uns gestorben ist, und uns den Heiligen Geist in unsere Herzen schenkt, ergreift. Gleicher Weise hat demselbigen Volke Israel das Osterlämmlein ein Gedächtniß sein sollen, dabei sie gewiß wären, GOtt wäre bei ihnen gegenwärtig, und [daß] sie ihm dieneten und ihn anrufeten, die weil er sie durch so viel Mirakel aus Egypten geführet hätte. Und zu Stärkung dieses Gedächtnisses wird dies äußerliche Zeichen des Osterlämmleins eingesetzt. Denn GOtt hat
\quellseite{451}
sich allezeit der Welt offenbart durch sein Wort und äußerliche Zeichen, und das hat er darum gethan, daß er es von uns nicht leiden kann, daß wir eine eigene Brücke in Himmel bauen und ihn da suchen wollen.

4. Wie denn auch vor dieser Zeit gethan haben, die mit ihren Werken gen Himmel klettern wollten, als wir Mönche vorgenommen. Wir haben gepredigt: So ich beschoren bin, und habe eine Platte, trage eine schwarze Kappe, dann so gefalle ich GOtt. Ja, du gefällst dem Teufel auf deinem Kopf. Solches ist alles von ihnen eingesetzt und erdichtet. Ich will ihn mit meinem Zeichen holen.\footnote{Hier ist die Relation wieder sehr unklar. Mit diesem letzten Satze soll vielleicht gesagt sein: Ich will ihn [GOtt] mit den von mir selbst aufgeworfenen, äußerlichen Zeichen zu erlangen suchen. Vgl. § 7 f.} Solche von mir erdichtete Gebärde und Wahrzeichen gelten nichts; es ist der Teufel. Ich soll wohl ein Zeichen haben, und zwar ein äußerlich Zeichen; aber nicht das, das ich eingesetzt, gestiftet und erdacht habe, sondern das GOtt hat eingesetzt. Ich soll ein solch Zeichen haben, davon GOtt selbst spreche: Das setze ich ein, daß es ein solch äußerlich Stück sei, damit du mich fühlest, daß ich ein solcher GOtt, und gegenwärtig bei dir sei.

5. Dazu dient dies nach der Historie, auf daß sie gewiß sageten: Allhier wohnt GOtt wahrhaftig. Also hat er die Sacramente auch zum Zeichen gegeben, daß man sagen könnte: Da wohnt GOtt wahrhaftig. Ich habe es nicht eingesetzt, so hat es auch sonst kein Mensch erdacht, sondern es ist vom Himmel gekommen; auf daß ich nicht durch mich zuführe, und nicht eine Weise aufwürfe, die dazu diente, daß es sollte GOtt gefallen.

6. Gleichwie Jeroboam zwei Kälber aufrichtete, eines zu Dan, und das andere zu Bethel, und predigte: „Siehe, Israel, allhier ist dein GOtt, der dich aus Egypten geführet“ [1 Kön. 12, 28. ff.], und das Volk fuhr zu, und opferte daselbst, meinte, es wäre also der rechte Gottesdienst, und GOtt wäre daselbst gegenwärtig; und daß sie den wahrhaftigen GOtt anrufeten, ließen sie ihnen nicht ausreden. Darüber haben sich die Propheten mit den falschen Lehrern zerschlagen und zerpredigt, und gehen auch frisch drob zu Grunde, daß sie alle um dieser Strafpredigt willen erschlagen.

7. Also haben wir auch gethan, und mit des Pabsts und Mönche erdichteten Gottesdienste uns zerrauft und zerfatzbalgt, die da vorgeben, daß ihr Ablaß, Messe, Vigilien, Wallfahrt, Klosterleben, Fasten, Beten und Almosengeben sollte ihnen machen einen gnädigen GOtt, und solch ihr Leben wäre der rechte Gottesdienst.

8. Aber GOtt spricht: Allhier findest du mich, auf dem Predigtstuhl, in der Taufe, im Abendmahle, denn das ist meine äußerliche Ordnung. Aber wir, im Papstthum, haben vorgegeben und gesagt: Willst du in den Himmel kommen, so mußt du aus der Welt in ein Kloster laufen, in eine Kappe kriechen; da, da findest du unsern HErrn GOtt; also sind wir hineingefallen, wie die Säue zum Thore hinein.

9. Nun, obwohl solches gar närrisch und thöricht scheint, dennoch so soll es in der christlichen Kirche allein gelten. Denn GOtt hat seinem Volke wollen äußerliche Zeichen geben, äußerliche Stücke und Ordnung, dabei sie ihn antreffen und finden möchten, und hat damit ihnen steuern und wehren wollen, auf daß sie nicht eigene Gottesdienste erfinden sollten.

10. Sie mußten nicht Schuhe anhaben; item, Stäbe in Händen tragen, und andere Ceremonien dazu gebrauchen. Es scheint nichts herrlich daselbst, sondern es hat alles ein närrisch Ansehen. Aber zur selbigen Zeit, ehe denn die Taufe und das Sacrament des Altars geordnet ward, da ist es ein groß Ding gewesen, daß GOtt mit denselbigen äußerlichen Stücken die Leute hat an sich ziehen und binden wollen, auch Augen und Ohren an diese Ordnung knüpfen, daß sie recht und wahrhaftig wären.

11. Unsere erdichtete Andacht und Werke sind nicht also, wenn wir auch der Kappe und Klosterlebens gleich recht gebrauchten. Denn der Menschen Tand gilt nicht, und ob man es gleich wohl gebraucht, so hilft mich doch meine Kappe nichts. Die Beschneidung hilft auch nicht, wie St. Paulus [1 Cor. 7, 19.] solches lehrt; aber die göttlichen Ordnungen gehen alle auf den Glauben. Also bringt mich das Sacrament des Altars auf den Glauben, denn die Worte: „Nehmet hin, esset und trinket“ \&c., die fordern den Glauben. Derselbige Glaube fehlt nicht. Ob ich gleich ein Schalk bin, so ist doch GOttes Wort und Ordnung nichtsdestoweniger in seiner Kraft, und bleibt bei den Gläubigen. Also haben sie auch gesungen, und sich erinnert, wenn sie das Osterlämmlein gegessen haben: „Ich bin der
\quellseite{452}
„HErr, dein GOtt, der dich aus Egypten geführet hat.“ Das mußt du nun glauben.

12. Die es nun mißbraucht haben, die haben es mit ihrem Schaden gethan, die Ordnung ist für sich köstlich, und das Wort recht, wahrhaftig, gut und heilsam geblieben. Aber wenn ich sage: Ich will in ein Kloster gehen, so werde ich selig werden, so ist dasselbige Wort nicht recht, denn ich thue es ohne Glauben. Darum, was GOtt ordnet, wie närrisch es scheint, dieweil es GOttes Wort für sich hat, das gilt dem Glauben. So viel nun diese gehandelte Geschichte anlangt, hat zu unserer Zeit solches alles aufgehört, und gilt nun etwas Anderes, als, das Sacrament der Taufe und des Altars.

13. Dieses sei auf die Historie gesagt. Hernacher sind die Kinder von Israel ausgezogen aus Egypten, und haben alle entlehnten und geborgten Güter mit sich genommen, und alle Erstgeburt todt hinter ihnen gelassen, und müssen also bei Nacht ausziehen, werden ausgetrieben in der Eile, daß sie auch keinen Teig zu Brod zubereiten können.

\begin{schriftzitat}
V. 9–12. Und die Egypter jagten ihnen nach, und ereileten sie, da sie sich gelagert hatten am rothen Meer mit Rossen und Wagen, und Reitern, und allem Heer des Pharao im Thal Hiroth gegen BaalZephon. Und da Pharao nahe zu ihnen kam, huben die Kinder von Israel ihre Augen auf, und siehe, die Egypter zogen hinter ihnen her, und sie fürchteten sich sehr, und schrieen zu dem HErrn. Und sprachen zu Mose: Waren nicht Gräber in Egypten, daß du uns mußtest wegführen, daß wir in der Wüste sterben: Warum hast du uns das gethan, daß du uns aus Egypten geführt hast? Ist es nicht das, daß wir dir sagten in Egypten: Höre auf, und laß uns den Egyptern dienen. Denn es wäre uns je besser, den Egyptern dienen, denn in der Wüste sterben?
\end{schriftzitat}

14. Da werden wir bald hören, wie sie errettet sind von dem Könige Pharaone; und haben daran ein schönes Exempel des Glaubens. Sie sind je auf GOttes Geheiß, auf seine Weisung und Anleitung, nach dem rothen Meer gezogen, und allda ihre Gezelt aufgeschlagen, haben für sich GOttes Befehl und Beruf, kommen aber darüber in große Fährlichkeit und Noth. Denn, da sie in die Wüste kommen, sind sie also gelegen, daß sie das rothe Meer vor der Nase gehabt, und sind ihrer sechsmal hunderttausend streitbarer Mann, ohne was sonst Weiber, Kinder, Knechte und Mägde mit gewesen, die konnten über das Meer nicht springen. Darnach waren auf beiden Seiten hohe Berge, daß sie also gleich als in der Kluft liegen, zwischen zweien großen Gebirgen, und vorne hatten sie das Meer. Hinterwärts kommt der Feind Pharao, und dringt mit aller Macht auf sie, und gedenkt der Feind: Ei, es ist gut, dort werden sie mir über das Meer nicht fliehen, so werden sie mir über das hohe Gebirge nicht klettern können. Fällt also in Verblendung, daß er meint, er habe sie nun, wie er selbst wolle. Und wahrlich, er hatte sie auch, gleichwie man eine Maus in der Falle hat, und Rebhühner in einem Garne haben mag. Er gedachte aber nicht, daß GOtt ihnen helfen würde. Aber GOtt verkündigt den Kindern von Israel zuvor diese Gefährlichkeit, und verheißt ihnen auch Hülfe und Erlösung; denn er wolle Ehre an ihm einlegen, und beweisen, daß er allein GOtt sei.

15. Daraus man lernen mag, wie es denen geht, die im göttlichen Berufe wandeln, daß Anfechtung, Noth und Gefahr ihnen unter Augen stößt. Denn der Teufel feiert nicht, sondern schickt denselbigen allerlei Anstöße zu, ob sie des Berufs überdrüssig und müde werden möchten. Aber GOtt weiß wohl um solche Gefährlichkeit, und will ihnen wieder draus helfen. Denn er schickt solche Noth den Seinen nicht darum zu, daß er sie verderben wollte, sondern seine Gnade und Barmherzigkeit gegen ihnen zu erzeigen, und seinen Ernst, Zorn und Ungnade an seinen Feinden sehen zu lassen. Derwegen spricht er auch an diesem Orte, er wolle an Pharao und aller seiner Macht Ehre einlegen, und die Egypter sollen inne werden, daß er der HErr sei.

16. Aber sie, die Kinder Israel, können solches weder sehen noch erkennen; ja, die Todesangst und Noth tritt ihnen in das Angesicht, wie sie mit ihren eigenen Worten bezeugen, da sie sagen zu Mose: Ja, „waren nicht Gräber genug in Egypten, daß wir daselbst gestorben wären?“ etc. Fallen wir über die Berge, so brechen mir die Fälle entzwei; springen wir in das Meer, so ersaufen wir; laufen wir denn zurück, so fallen wir in des Feindes Schwert. Und hat sich also ein groß Zetergeschrei und
\quellseite{453}
Klage unter dem Volke erhoben, denn sie nicht alle gläubig gewesen sind und GOtt vertraut haben. So sind ihrer auch viel mit aus Egypten gezogen um fleischlicher Freiheit willen, daß sie nach der großen Arbeit, Beschwerung und Frohndiensten nun wollten Junker sein. Aber da das Kreuz und Todesgefahr daher geht, da läßt ihr Unglaube sich merken, da ist in ihrem Munde eitel Lästerung wider Mose, und an ihren Gebärden große Ungeduld und Verzweiflung gegen GOtt zu spüren.

17. Mosis Herz hat allhier auch gezittert und gezagt. Denn er war ihr Herzog und Heerführer, ihr Oberster und Hauptmann, daß er gedacht hat: Du hast sie ausgeführt, und bringst nun so viel Leute um den Hals. Siehe, wie sie dir zusprechen. Solche innerliche Wehmuth und Herzengegeschrei Mosis hört GOtt leise und bald, obgleich Moses nicht viel schreit. Denn das ist die Ursache, daß ihm viel zu angst und bange dazu war.

18. Das ist eine rechte, schöne Art und Stück des Glaubens. Da fühlt man des göttlichen Wortes Kraft, und was der Christen Glaube sei. Der Christ vertrauet GOtt, und ist wohl zufrieden, wie es GOtt mit ihm macht. Er hätte die Kinder von Israel wohl bald können durch das Meer führen; aber er will sie noch besser versuchen, und läßt sie die ganze Nacht liegen, daß sie hinter ihnen der Egypter Gezelt sehen. Doch, der Engel des HErrn geht in einer Wolke vor den Kindern von Israel her, und da jetzt die Egypter ihnen nachfolgen, so läßt sich die Wolke hernieder, und setzt sich zwischen die Egypter und Israeliter. Nun meinen die Egypter nicht anders, denn es sei nur ein Wetter, und trübe, finstere Wolken; aber GOtt spricht zu Mose: Schlage in das Meer mit deinem Stabe; alsbald kommt ein Wind, und reißt das Meer von einander, daß es stand zu beiden Seiten, wie zwo Mauern.

19. Zuvor waren sie gar traurig, klagen, sie sind des Todes, und sind dieses ihre Worte: „Waren nicht Gräber in Egypten?“ etc. Wollen sagen: Pfui dich an, Moses, du verzweifelter Bösewicht, du wärest werth, daß man dich mit den Zähnen oder Zangen zerrisse. Jedoch bedenken wir, wenn uns solche Geschichte wird vorgetragen, sie haben den Sachen zu viel gethan; aber wären wir allda gewesen, so hätten wir eben dasselbige auch geklagt. Wenn noch auf den heutigen Tag irgend ein Kriegsvolk käme, und uns ungewarneter Sache überfiele, so verzweifelten wir alle. Aber allhier müssen wir ihren Glauben ansehen, und ihnen in das Herz hinein schauen. Ihr Glaube schließt die Augen zu, und sieht das Meer, noch die Berge, noch den Feind nicht, sondern hängt allein an dem Wort Gottes, das GOtt gesagt hatte: Ich will es thun, und aus Pharaonis Diensthause die Kinder von Israel ausführen.

20. Gleichwohl betrachte man, wie wunderlich er es angreift. Da er sie will los und ledig machen, führt er sie erst recht in Tod hinein. GOtt spricht: Ich will es thun. Aber einer möchte sagen: Ei, wie willst du es thun? ist doch allhier des Pharaonis Schwert; item, das hohe Gebirge und das tiefe Meer, und die große Macht und Gewalt der Feinde. Wo sollen wir hin? Wenn man einen will lebendig machen, so ihm der Kopf ist abgeschlagen, das geht langsam an. Aber GOtt spricht: Laß mich sorgen. Sagst du dann: Ei, es ist unmöglich, die Vernunft kann es nicht glauben! Ja, antwortet GOtt, das will ich auch haben; siehe auf mein Wort. Ich muß glauben, daß GOtt könne das Meer wegsprühen als ein Stäublein, und daß er könne die Berge versetzen, und das Schwert Pharaonis auch wegtreiben, und auch die Erde wegthun. Denn, hat er doch Himmel und Erde gemacht.

21. Dieses haben ihrer viel geglaubt. Denn GOtt kann die Kunst, daß er das Meer also trocken machen kann, gleich als ginge eine Brücke drüber, und das weiche Wasser muß also hart werden, als eine Mauer; dagegen müssen die Berge so weich werden, als ein Wasser oder Fluß, und muß das Meer ein feiner, trockener Weg sein. Da zuvor einer hat ein wenig hinein getreten, da hat es gequitschet.\footnote{quitschen wohl so viel als: quatschen, ein plätscherndes Geräusch machen.} Warum, wer diese Mauern des Meeres gesehen und geschauet, daß dies Meer ist zur Mauer worden, der hat mögen gedenken: Wer wollte so kühn sein, daß er hinein ginge? Und wenn jemand so keck wäre, so müßte er doch sich fürchten, daß nicht das Wasser über ihn fiele, und ihn ersäufte. Denn das Meer stund aufgehebt auf beiden Seiten, gleich als wenn es mitten von einander geschnitten wäre. Da mußte man fürchten, daß
\quellseite{454}
nicht das Wasser einem auf den Kopf fallen möchte. Denn niemand hält es, es ist keine Wand drunter gezogen.

22. Aber der Glaube, der hält es, der weiß, daß es GOtt hält. Das Meer wird sich weiter aufgethan haben, denn diese Stadt Wittenberg lang oder breit ist; sechs, sieben oder acht Meilen Wegs lang. Das ist eine große Scheidung und eine herrliche Erlösung, und ist das Mirakel so groß, daß man sich darob nicht genug verwundern kann. So lange steht den Kindern Israel das Meer, bis sie alle mit ihrer Habe und Zugehör hinüber sind. Gedenke, wie viel Zeit sie dazu haben müssen nehmen, daß sie mit einem so großen Kriegsvolke haben können hindurch kommen. Denn es will Maße und Raum haben, daß dies Volk hindurch gezogen ist; sie werden Einen Tag, zween oder drei gezogen sein, und im Gehen mit gegessen und getrunken haben, denn sie sonst mitten im Meere nicht viel Speise noch Futter gehabt.

23. Es ist ein trefflich Ding gewesen, wiewohl es hier mit kurzen Worten beschrieben ist; aber wenn man ihm nachdenkt, und rechnet es aus, so ist es ein groß, wunderbarlich Ding, und ist der Boden oder Grund im Meere so tief gewesen, und das Wasser so hoch gestanden, als zwei hohe, große Gebirge. Das hat gewährt sechs oder sieben Meilen, daß im Meere ist ein trockener Weg und Fort gewesen, gleichwie sonst eine Straße auf einem Lande.

24. Nun sind es in der Wahrheit kühne Leute gewesen, die sich hinein gewagt und begeben haben. Ist es nicht ein Mirakel? Aber es ist die rechte Art des Glaubens, welche man allhier an den Kindern von Israel sieht. Der Glaube thut die Augen zu, wiewohl nicht gar, denn das Herz zappelt. Aber sie vertrauen dem Worte, und glauben GOtt, daß er könne entweder das Meer pfastern, oder die Berge dem Felde gleich und eben machen, oder auch die Egypter wegblasen, oder auf einmal sie alle zu Tode schlagen.

25. Wenn wir doch auch also glaubten, so oft wir in Noth und Gefahr Leibes und Lebens kämen, wenn der Tod daher geschlichen käme, oder wir in unserer Feinde Hände fielen, daß wir sprächen: Ich habe einen GOtt, der mir zugesagt hat, er wolle mich nicht verlassen; da ist sein Wort, er wolle mich erretten, und am Leben erhalten. Wenn ich das glaube, so sehe ich nicht auf die Feinde, auf den Tod, Schwert, Pfeilen, Hunger, Gefängniß, oder was sonst die Noth ist, so mir zuhanden kommt, sondern ich gehe hindurch. Und also müssen zuletzt die Feinde zu Pulver und Asche werden, und die Spieße zu Strohhalm gemacht werden; wie denn Pharao mit aller seiner Kriegsrüstung vor ihren Augen ersäuft und verschwindet.

26. Wer aber nicht glaubt, mit dem kann GOtt kein Mirakel noch Wunderzeichen thun, sondern sie müssen untergehen, und können dem Tode nicht entlaufen. Aber die da GOtt vertrauen, die werden errettet, sollte eher Himmel und Erde vergehen. Denn er hat es also verheißen, und gesagt: Wirst du glauben, so soll dir geholfen werden [Ps. 125, 1. ff.]. Sein Wort gilt ihm mehr, denn Himmel und Erde. Denn er hat mit seinem Worte alles erschaffen; darum so will er auch, man soll seinem Worte und Verheißungen mehr glauben. Wenn gleich alle Welt, der Himmel, Erde und alle Sterne wider dich wären und als deine Feinde sich erzeigten, doch solltest du meinem Worte vertrauen. Denn, ist seine mündliche Stimme nicht also mächtig, daß sie könnte hinweg reißen den Himmel und auch die Berge? Man sieht es wohl nicht, aber wer es glaubt, der hat es.

27. Der Glaube ist ein recht Mirakel, der macht in einem Menschen, der sonst eine schwache, arme Creatur ist, einen solchen großen Muth, daß er dadurch also stolz wird, daß er sagen kann: So alle Teufel auf mich fielen, ja, alle Könige, Kaiser, Himmel und Erde wider mich wären, dennoch glaube ich und weiß, daß ich werde erhalten werden. Der Glaube ist ein groß Ding, man kann seine Macht, Kraft, Stärke und Gewalt nimmermehr genug verstehen.

28. Darngegen ist dieses Exempel des Glaubens gar herrlich. Es hat das göttliche Wort und der Glaube große Kraft; so ist dem Glauben nichts unmöglich, noch zu schwer, oder zu bitter [Marc. 9, 23.]. Der Glaube ist ein groß, herrlich Werk. Wer glaubt, der ist ein Herr, und ob er gleich stirbt, so muß er doch wieder leben. Ist einer arm, so muß er doch reich sein; ist einer krank, so muß er doch wieder gesund werden. Wenn man durch den Glauben also den Karren\footnote{Eislebensche: Karn.} geschmiert hat, so geht das Fuhrwerk fort. Aber wir üben den Glauben nicht,
\quellseite{455}
und erkennen noch lernen nicht, daß GOtt uns in allen Nöthen helfen könne; item, daß der Glaube stärker sei denn Himmel und Erde, oder alle Creaturen. Denn GOtt spricht selber [Luc. 21, 33.]: „Himmel und Erde müssen vergehen; aber sein Wort, das vergehe nicht.“

29. Also sieht man allhier ausdrücklich und eigentlich, daß die Kinder Israel durch ihre Kriegsrüstung ihnen selbst nicht geholfen, wiewohl sie Schwert und andere Rüstung auch gehabt haben. Denn die Faust oder menschliche Vernunft kann allhier nicht helfen; das Wasser im rothen Meere kann man nicht ausgießen, so kann man die hohen Berge auch nicht zermalmen oder zerschmettern, ob sie gleich ewiglich gestähet hätten. Aber es ist keine andere Hülfe allhier, denn GOttes Wort und Verheißung, daß GOtt spricht: Ich bin bei dir; ich will dein GOtt sein. Dies Wort ist nicht unser Werk. Nein, daß er zu Mose sagt [V. 26.]: „Recke deine Hand aus über das Meer, daß das Wasser herfalle über die Egypter, über ihre Wagen und Reiter.“ Diesem Worte sollen sie allein folgen und anhangen; und wir auch zu ihm, als zu dem rechten Asilo, Zuflucht und Trost haben.

30. Es ist der Vernunft ein lauter Scherz und Spott, daß Moses mit dem Stabe in das Meer schlägt, und das Meer sich fürchtet, und von einander läuft, und steht auf beiden Seiten als eine Mauer, gleich als nähme Moses einen Löffel oder ein Schüffen,\footnote{„Schüffen“ ist wohl so viel als: Schüppe, Schöpfer.} und wollte das Wasser ausschöpfen, so es doch die ganze Welt nicht könnte ausschöpfen. Aber das ist es, daß GOtt spricht: Ich will euch erretten. So spricht Moses und die Kinder von Israel: Das glauben wir, lieber HErr. Und sagt Moses denn ferner: Wohlan, Meer, thue dich auf. Also bringt der Glaube hindurch, und theilt die Wasser und das Meer von einander. Daher sollen wir auch lernen glauben, und diese schönen, tröstlichen Exempel nicht verachten.

31. Da nun die Kinder von Israel herdurch sind, gedenkt der unsinnige, thörichte und närrische Pharao: Ich will auch hindurch. Denn er sieht nicht, daß ein Mirakel und Wunderwerk GOttes ist; sondern ist verblendet und verstockt, sagt in seinem Herzen: Wir sind auch fromm, wir wollen hindurch gehen, gleich als die Kinder von Israel; und sind doch verzweifelte, böse Buben, so GOttes Volk verfolgen, und gedenken, nun wollen sie erst recht mit ihnen handeln und sie plagen. Aber da Pharao und die Egypter eine Zagereise hinein kommen in das Meer, und die Kinder von Israel hindurch sind, so spricht GOtt zu Mose: Schlage in das Meer, dann fällt das Wasser (so als Mauern auf beiden Seiten stand) wieder in einander.

32. Ei, hat sich denn der mächtige König Pharao nicht gewehret? Die Räder fließen von den Wagen hinweg, und schwimmt alle seine Macht im rothen Meere. Denn GOtt hat ein Wetter, einen Wind, Brausen, Schrecken und Flucht unter sie kommen lassen, daß sie schreien: O lasset uns fliehen, nun ist Zeit Fliehens. Aber sie waren zu weit in das Loch gekommen, und steckten alle mit einander im Wasser, etliche hunderttausend Mann. Da plumpt das Wasser gar zusammen, und müssen alle ersaufen, daß nicht Einer davongekommen wäre.

33. Also kann GOtt sein Volk erretten und ihnen helfen. Wenn Pharao den Glauben hätte gehabt, so wären die Wassermauern auch stehen geblieben; aber da er nicht den Glauben hat, so werden die Steine wieder zu Wasser. Der Glaube macht aus Wasser eitel Steine, auch aus Feuer macht er Wasser, und aus Wasser kann er Feuer zurichten, und GOtt macht es, wie derjenige will, so da glaubt, wie der 145. Psalm, V. 19., auch sagt: Deus facit voluntatem timentium se, er thut den Willen derer, die ihn fürchten. Das will GOtt von ihm wissen lassen, und solches sollen wir uns zu ihm versehen, daß er den Gläubigen helfen will. Darum hat er diese Wunderwerke gethan, daß wir auch lernen glauben und GOttes Hülfe erwarten. Wie denn im Propheten Jesaia am 30. Capitel, V. 15., GOtt uns zum Glauben und Hoffnung der göttlichen Hülfe vermahnt, da er spricht: „In Stille-sein und Hoffnung werdet ihr stark sein.“ Wie allhier Moses auch spricht: „Fürchtet euch nicht, stehet fest, und sehet zu, was für ein Heil der HErr heute an euch thun wird.“

34. Also haben wir bis anher gehört von dem Auszuge der Kinder von Israel aus Egypten, gleichwie zuvor von dem Sündflüßlein und von der Erstgeburt, was der Heilige Geist durch diese beiden Stücke bedeutet und vorgebildet habe, daß im neuen Testamente geschehen sollte. Nun ist das dritte und letzte noch hinterstellig zu handeln, als, wie die Kinder von Israel durch
\quellseite{456}
das rothe Meer gegangen sind. Da denn GOtt gar ein groß Wunderwerk gethan hat, daß er das Volk Israel mit trockenem Füßen durch das rothe Meer geführt hat, da sie sonst auf beiden Seiten gefangen waren. Denn vor ihren Augen hatten sie das rothe Meer; auf dem Rücken folgte ihnen der Feind nach, der König von Egypten; zu beiden Seiten war ein groß Gebirge und hohe Felsen, darüber sie nicht konnten steigen; waren also in großen Aengsten und Nöthen, und war der Tod da vor Augen. Darum sagten sie auch zu Mose: „Waren nicht Gräber in Egypten?“ Solche Hindernisse alle ungeachtet, richtet GOtt diesen Zug oder Durchgang durch das rothe Meer gleichwohl aus. Damit er hat anzeigen wollen, daß er uns helfen könne und wolle in allen unsern Nöthen und Anliegen, daß er auch wisse Wege und Raum zu finden, da gar kein Weg oder Raum ist, und in Summa aus Nichts Alles machen könne.

\subsection{Deutung dieser Geschichte}

35. Nun wollen wir auch deuten, was das rothe Meer sei. In griechischer Sprache wird es das rothe Meer genannt, nicht, daß das Wasser roth sei, denn es ist einerlei Meer, sondern darum, daß es ein roth Ufer gehabt; die Berge am Ufer oder Rande haben rothe Erde gehabt. In hebräischer Sprache wird es das Schilfmeer geheißen, darum, daß viel Schilf und Röhricht am Ufer hin und wieder gewachsen ist.

36. Wir haben aber drei Dinge gehandelt, ehe denn es zum Treffen kommt. Erstlich, so erschrecken die Israelitter, und sind gleich als wären sie todt. Darnach, daß GOtt verheißt, er wolle die Egypter verstocken, daß sie sollen den Israeliten nachfolgen bis an das Meer, und spricht zu Mose: Schlage du mit deinem Stabe in das Meer, so soll sich das Wasser zertheilen. Zum dritten, daß er einen Engel läßt vor dem Volke herziehen in der Luft, und derselbige muß sich hinter das Volk lagern und zurückziehen, und sich niederlassen zwischen die Israeliter und Egypter. Da nimmt sich unser HErr GOtt dieses seines Volks nicht anders an, denn wie sich ein Vater oder Mutter ihrer Kinder annimmt, und für sie sorgt, und ihnen hilft [Ps. 103, 13.]. Denn GOtt nimmt sich mit solcher Sorge und Fleiß ihrer an, daß er sie in Nöthen nicht verläßt, sondern hilft ihnen aus, und errettet die Seinen. Aber er kommt nicht eher, es sei denn alles aus, und ganz und gar desperirt, daß die Noth am größten ist [Ps. 9, 10.]. Es muß zuvor alles zu Trümmern gehen, und die Ochsen am Berge stehen, und die Saite auf das höchste gezogen sein, daß sie jetzt zerspringen will; das Wasser muß über Berge und Thal gehen, ehe denn sie von den Egypten erlöst werden.

37. Darum so hat sich die Hülfe also lange verzogen, daß sie darüber gar in Verzweiflung kommen müssen, und denken: Ach, wer gestorben und schon begraben wäre! Murren derhalben und sind ungeduldig. Damit wird angezeigt unsere Schwachheit und Sünde. Wenn das Gewissen erschreckt wird, so fürchtet es GOttes Zorn; das macht denn einen verzagt und blöde.

38. Da ist denn Pharao auch hinter einem her, das ist, die Gesetzgeber, die einem den Himmel zu enge und die Hölle gar zu weit machen, und dahin treiben, daß einer gleich verzweifeln möchte, und fürchten, er müsse ewig verloren sein. Also fürchteten sich die Kinder von Israel, daß sie müßten Haare lassen, und in dieser Noth des Todes verloren sein. Dieses ist nun eine elende Marter, darüber in den Psalmen oft geklagt wird, daß die Heiligen sagen: HErr, du hast mich gar verworfen von deinem Angesichte [Ps. 51, 13.]. Item, im 6. Psalm, V. 1. f.: „Ach HErr, strafe mich nicht in deinem Zorn, und züchtige mich nicht in deinem Grimm. HErr, sei mir gnädig, denn ich bin schwach; heile mich, HErr, denn meine Gebeine sind erschrocken, und meine Seele ist sehr erschrocken. Ach du HErr, wie lange!“ Wo ich mich umsehe, da sterbe ich unter meinen Feinden. Ich werde allenthalben geängstigt, Himmel und Erde ist mir zu enge worden, ich konnte nirgendes bleiben, allenthalben waren meine Feinde.

39. Da werden wir gelehret, daß wir in Verzweiflung nicht sollen fallen, noch verzagen, und wie man in Anliegen soll aus der Noth kommen. Denn wenn es zu solchem harten Stande und Zügen nicht kommt, so schmeckt uns GOttes Gnade und Hülfe nicht.

40. Darnach, daß Mosi gesagt wird, er solle mit dem Stabe in das Meer schlagen, so solle sich das Wasser zertheilen, gar trocken werden, und stehen wie zwo Mauern. Da muß das Wasser, das sonst andere ersäufen sollte, der Schutz werden, und sie bei dem Leben behalten. Dasjenige, das würgen sollte, und Schaden anrichten
\quellseite{457}
, daß muß Frommen thun und lebendig machen. Das auch zur Hölle führen sollte, das muß gen Himmel helfen.

41. Viel Propheten und Psalmen haben sich mit dem rothen Meere bekümmert, daß es dem Volke Israel hat müssen zum Leben dienen, und doch der Egypter Tod war. Also wunderbarlich ist GOtt, daß er mir mit dem zur Seligkeit hilft, das sonst meine Hölle ist, und daß er mich durch Schande zu Ehren bringt. Es wäre das rothe Meer den Kindern von Israel der rechte bittere Tod gewesen; sie hätten müssen alle drinnen ersaufen, wie es denn den Egyptern widerfuhr; dennoch ist dieser Tod den Israeliten ein Leben. Und wird allhier gesagt, das Meer sei gestanden als zwo Mauern, gleich als wäre gar kein Wasser allda.

42. Aber wodurch geschieht solches? Daher, daß Moses mit einer Ruthe, Stecken oder Stabe, aus GOttes Geheiß drein geschlagen hat. Der Schlag, der thut es. Es ist wohl ein närrisch Ding, daß er mit einem Stabe in das Wasser schlagen soll, und dieser Schlag soll solche Kraft haben, daß er das Meer theilen soll, daß es stehen muß, wie zwo Mauern, und der Israeliter Schutz und Schirm werden. Wenn es in der heiligen Schrift nicht geschrieben stünde, so würde man es für die höchste Thorheit und Narrenwerk halten, daß ein Holz oder Stecken solche großen Dinge solle thun. Aber das ist es, daß es GOtt ebensowohl ohne den Stab Mosis hätte können thun, und es wäre des Stabes nicht vonnöthen gewesen. Denn GOtt hat es gethan, und nicht der Stab, wenn man es will recht ansehen. Denn der Text sagt: Da das Meer mit dem Stabe geschlagen war, da kam ein Wind, der das Meer weghob, daß es trocken ward. Der Wind hat es weggenommen, und nicht der Stab.

43. Der Stab ist das heilige Evangelium; wie denn allenthalben in der heiligen Schrift wird GOttes Wort ein Stab genannt, als im drei und zwanzigsten Psalm, V. 4.: „Dein Stecken und Stab trösten mich“ etc., und Jesaia am elften Capitel, V. 15., wird gesagt: „Er wird seine Hand gehen lassen über das Wasser mit seinem starken Winde, und die sieben Ströme schlagen, daß man mit Schuhen dadurch gehen mag.“ Das geschieht mit dem Stabe oder Stecken, daß es muß gehen und geschlagen werden, und mit Freuden getrieben werden zwischen den Sündern; das macht denn die Gewissen getrost. Wenn der Schlag geschieht, so weicht Sünde, Teufel, Tod und alles Unglück. Nach dem Schlage hebt der Wind das Meer hinweg, daß es Raum wird; das ist, wenn GOttes Wort getrieben wird, dann ist der Stab in das Wasser geschlagen, und nach dem Worte folgt der Heilige Geist, der bläst das Gewässer hinweg und macht einen Muth, erneuert die Herzen; dann ist hinweg, verschwindet, zerstäubt und zerfliegt der Tod, die Sünde und alles Unglück.

44. Dieses ist nun eine wunderbare Kunst GOttes, daß er durch das mündliche Wort, wenn es gepredigt wird, mit uns handelt, den Heiligen Geist gibt und schenkt, welcher schafft, daß einem Christenmenschen nichts schaden muß, sondern das Gegenspiel muß darauf erfolgen, daß ihm auch hülflich und nützlich sei, was an ihm selbst und seiner Natur nach schädlich und verderblich wäre [Sprichw. 8, 34. f.]. Das Stabschlagen in das Meer ist, das Evangelium frei schlagen stoßen, predigen und treiben, in Anfechtungen leiblich und geistlich. Dann kommt der Heilige Geist und macht Raum. Also geht es, wenn mich die Sünden drücken, mein böse Gewissen beschweret und zur Verzweiflung treibt; je mehr mich dieselbigen kränken und beschweren, je mehr ich getrieben werde, daß ich auf GOtt hoffen soll.

45. Das ist ein lustiger, feiner Krieg, daß ich GOttes Wort habe und höre, daß Christus gepredigt wird, er sei mein HErr, und stehe für mich [Eph. 6, 13. Ps. 144, 1.]. Dann ist mein voriges Leben und Wesen, auch meine Sünde da, welches ist gleich als das rothe Meer, da spricht Pharao: Dies und das hast du nicht gehalten: aber je mehr mich das Gesetz, Sünde und Teufel martert, je mehr ich hoffe, und wird mir Raum gemacht, wenn ich dem Worte GOttes anhange. Also müssen auch die Sünden einem Christen helfen; und wenn er nicht Sünde hätte, so käme er nicht so wohl hindurch. Denn so ich nicht die Sünde, das böse Leben und Gewissen fühlte, so schmeckte mir nimmermehr die Kraft des göttlichen Wortes alle wohl. Aber dieweil sie mich martern, und die Hölle heiß und schrecklich machen, darum so werde ich der Kraft des göttlichen Worts desto mehr gewahr, daran ich hange, und bringt mich hindurch.

46. Das hat GOtt in dieser Sache also geordnet; und was ich von der Sünde rede, das
\quellseite{458}
wird von andern auch gesagt. Du kommst nicht davon, denn allein durch das Wort GOttes, wenn man mit dem Stabe in das Meer schlägt; das ist, wenn du GOttes Wort annimmst, dann muß dir kein Unglück schaden. Das Wort GOttes theilt es gewaltiglich, daß dieselbige Sünde dich fördere und dir helfe, daß du viel besser herdurch kommst, denn sonst. Also macht das Meer lebendig, welches sonst tödtet. Das geschieht durch den Stab, und daß der Wind kommt, der alles dürre macht und trocknet. Der Stab hat es nicht gethan, sondern der Stab hat nur geschlagen; der Wind aber hat es gethan. Denn GOtt will den Heiligen Geist geben durch das Wort; ohne Wort will er es nicht thun [Joh. 3, 8. Gal. 3, 2.]. Er hätte wohl ohne den Stab das Wasser zertheilen können, aber er hat es nicht wollen thun. Also wird der Wind dir nicht zertheilen und vertreiben die Anfechtung und das Unglück, ohne den Stab, oder ohne das Wort GOttes. Halte dich allein zum Worte, nimm es in die Faust, und schlage frisch drein mitten in die Sünde und in den Tod; hier ist GOttes Wort, dann wird es alles beiseits weggehen, und Raum machen.

47. Der Teufel hat sonst die Welt erfüllt mit Geistern, die vom Winde und Geist predigen, aber ohne Mosis Stab, und die heilige Schrift verlassen. Ich warne euch davor, daß ihr euch wohl vorsehet, denn ich fürchte sehr, wir werden GOttes Wort wiederum verlieren um unserer Undankbarkeit willen. Der Wind theilt das Meer nicht, es sei denn, daß man habe mit dem Stabe drein geschlagen. GOtt erleuchtet dich nicht ohne sein göttliches Wort.

48. Der Stab ist gegen dem Meere nichts, wenn man es mit der Vernunft will ausrechnen. Denn der Stab ist irgend zwo Ellen lang, so ist das Meer in die acht oder zehn Meilen breit; noch thut es Moses und schlägt mit seinem Stabe drein, dieweil GOtt solches ihn geheißen, er sollte drein schlagen; darauf folgte bald der Wind, der es theilte und trocknete. Also ist auch das Wort GOttes ein schlechter Buchstabe, schwarze Tinte, ein Wort wie ein ander Wort. Da sagt man denn: Ei, was sollte das erlösen von Sünden, Tod, Teufel und Hölle? Aber siehe du zu, und gehe mit dem göttlichen Worte fleißig um, so wird der Heilige Geist dazu kommen, und mit dem Worte folgen, und alle dein Unglück ausschlagen und dir herdurch helfen.

Der andere Haufe Pharaonis, und zwar er Pharao selbst, die ersaufen im rothen Meer; das ist, wo GOttes Wort nicht ist, sondern Menschenlehren und Menschensatzungen sind, oder auch gleich das Gesetz GOttes, und Vertrauen auf gute Werke, da ersäuft alles.

49. Pharao ist ein Lehrer ohne GOtt und ohne den Heiligen Geist, ein müßiger Prediger, der da macht, daß man aus der Gnade GOttes fällt, und in eigener Heiligkeit und Vertrauen auf gute Werke einhergeht. Wer diesem Lehrer folgt, der muß zu Boden gehen, er kann es nicht erwehren. Wer aber GOttes Wort hat, der fühlt die Sünde, das Gesetz und den Tod, und schaden ihm nichts. Wer den Stab aber nicht hat, sondern hört andere Geister, denen er folgt, der wird in seinem Gewissen ersaufen.

\begin{schriftzitat}
V. 19. 20. Da erhub sich der Engel GOttes, der vor den Gezelten herzog, und machte sich hinter sie, und die Wolkensäule machte sich auch von ihrem Angesichte, und trat hinter sie, und kam zwischen die Gezelt(e) der Egypter und Israel. Es war aber eine finstere Wolke, und erleuchtete die Nacht, daß sie die ganze Nacht, diese und jene, nicht zusammen kommen konnten.
\end{schriftzitat}

50. Das ist auch ein Mirakel, daß GOtt heißt einen Engel, sich zwischen die Kinder von Israel und Egypter lagern. Das ist ein herrlicher Trost, daß GOtt väterlich für die Seinen sorgt. Denn er läßt seine Hülfe allhier erscheinen, und gibt uns die Engel zu Wächtern und Geleitsleuten. Gleich als wollte GOtt allhier sagen: Ich will es also machen, daß Pharao euch nachjagen soll; da richtet euch nach. So hat uns GOtt in seinem Wort, was zukünftig ist und uns übergehen würde, auch zuvor verkündigt. Das heißt ja väterlich und getreulich gewarnt und vermahnt, auch daneben zugesagt, daß er uns nicht verlassen will. Denn da ist GOttes Wort, das sagt: Die dich verfolgen, sollen dir dienlich und förderlich sein. Wie denn noch heutiges Tages der Pabst und Tyrannen wider ihren Willen das Evangelium gefördert haben, ob sie gleich sehr dawider getobt und gewüthet, und es ganz und gar haben dämpfen wollen.

51. Darum handelt GOtt väterlich genug mit den Israeliten, er verkündigt ihnen, er wolle sie nicht verlassen. Daher haben die Heiligen viel Trost genommen, daß GOtt seinen Heiligen also nahe sei, daß ihnen auch die Engel
\quellseite{459}
dienen müssen, wie der vier und dreißigste Psalm, V. 8., auch sagt: „Der Engel des HErrn lagert sich um die her, die den HErrn fürchten, und hilft ihnen aus“; und der 91. Psalm, V. 11. 12., sagt: „Er hat seinen Engeln befohlen über dir, daß sie dich behüten auf allen deinen Wegen, daß sie dich auf den Händen tragen, und du deinen Fuß nicht an einen Stein stoßest.“ Wenn wir das glauben, so sollten wir desto fröhlicher sein. Denn die lieben Engel sind unsere Wächter und Geleitsleute, ja, unsere Knechte und Diener, so auf die Christen warten müssen, daß ihnen kein Leid widerfahre, wie die Epistel zum Hebräern [Cap. 1, 14.] zeuget, daß die Engel dienstbare Geister sind, ausgesandt zum Dienst derer, die da selig werden sollen. Also sehen wir alhier auch, daß die Kinder von Israel nicht irgend eine Mauer oder Wall um sich haben, sondern gleich als große Heerscharen, dadurch sie vor dem Pharaone geschützt werden, daß sie nun sicher und GOttes Hülfe gewiß sind.

52. Aber da gehört Glaube zu, daß man wisse, die Engel GOttes sind um uns her, gleichwie ein Kriegsheer sich in das Feld lagert und niederläßt mit seinem Harnisch, Spieß, Büchsen, Roß und Wagen. Dies ist ungezweifelt wahr, wer es allein glauben könnte, und GOtt vertraute, der wäre denn gewiß, daß die lieben Engel bei ihm stehen. Derhalben so sollten wir dieses Schutzes der Engel uns trösten und sicher sein, daß es in Gefährlichkeiten nicht werde Noth mit uns haben oder gewinnen. Ich selbst wollte lieber einen Engel um mich haben, denn vier und zwanzig türkische Kaiser mit aller ihrer Macht und Gewalt; wenn sie gleich hundertmal tausend Büchsen bei sich hätten, so ist es doch alles gegen Einem Engel gar nichts [2 Kön. 19, 35.].

53. Und wenn wir sonst keinen andern Trost hätten, denn diesen einigen, so wäre es genug. Sage mir, ist es nicht ein groß Ding, daß GOtt für uns also sorgfältig ist, und so genau auf uns sieht? Was soll er doch mehr thun? Ja, sprichst du, ich möchte gerne einen Engel sehen. Lieber, da wäre denn des Glaubens nicht vonnöthen; aber man muß es glauben. Deß haben wir ein fein, herrlich Exempel in der Historie des Propheten Elisä [2 Kön. 6, 14–16.]. Sein Knecht sah, daß ein groß Kriegsvolk der Syrer um die Stadt Dothan kam, darinnen Elisäus war, und belagerte sie. Das mußte nun der Prophet nicht; darum, da es ihm von dem Knechte angezeigt ward, sprach er: „Fürchte dich nicht, ihrer sind mehr bei uns, denn bei ihnen“, und bat GOtt, daß er dem Knechte die Augen aufthäte; da sahe er, daß die Berge um Dothan alle voller feuriger Rosse und Wagen waren. Ja, wenn man einem die Augen also aufthut, so möchte er sehen. Darum, wer allein GOtt vertrauen kann, so sieht man nicht allein hunderttausend Engel, sondern alle Berge voll Engel, ja, man sieht GOtt selbst um sich her.

54. Also wird alhier auch angezeigt, daß der Engel sich des Volks Israel mußte annehmen wider die Egypter, auf daß wir wissen, daß wir einen solchen HErrn haben, der uns behütet und bewacht, daß uns ja kein Leid widerfahre [Ps. 121, 4.]. Wahrlich, es ist eine Sünde und Schande, daß wir davon so viel schöner Historien und Sprüche der heiligen Schrift haben, und dennoch dem lieben GOtt nicht vertrauen können, der so starke, gewaltige und tröstliche Verheißungen uns diesfalls gegeben hat; noch gehen wir dahin, und achten es nichts. Das macht alles unser leidiger Unglaube.

\begin{schriftzitat}
V. 21. 22. Da nun Mose seine Hand ausreckte über das Meer, ließ es der HErr hinweg fahren durch einen starken Ostwind die ganze Nacht, und machte das Meer trocken, und die Wasser theilten sich von einander. Und die Kinder Israel gingen hinein mitten in das Meer auf dem Trockenen, und das Wasser war ihnen für Mauern, zur Rechten und zur Linken.
\end{schriftzitat}

55. Das ist das andere Mirakel, so alhier geschiehet, daß Gott das Wasser als Mauern stehen läßt, und die Tiefe des Meeres muß austrocknen, auf daß die Kinder Israel könnten hindurch gehen; damit GOttes rechte Hand, Kraft und Allmacht gespürt wird, daß er den Seinen in der Noth, als ein gnädiger GOtt, wohl helfen könne. So lasset uns lernen an GOtt glauben. Denn was der Glaube vermöge, was er für Hülfe erlange, und was er für einen Nachdruck habe, das zeigt uns dies Exempel an. Davon rühmt und predigt die Epistel zum Hebräern am elften Capitel, V. 29., mit diesen Worten: „Durch den Glauben sind sie durch das rothe Meer gegangen, als durch ein trocken Land; und da es ihnen die Egypter wollten nachthun, gingen sie alle unter, versanken und ertranken.“ [Nun folgt] das dritte Mirakel:

\quellseite{460}
\begin{schriftzitat}
V. 24–26. Als nun die Morgenwache kam, schaute der HErr auf der Egypter Gezelt[e] aus der Feuersäule und Wolke, und machte ein Schrecken in ihren Gezelt[en]. Und stieß die Räder von ihren Wagen, führte sie mit Ungestüm. Da sprachen die Egypter: Lasset uns fliehen von Israel, der HErr streitet für sie wider die Egypter. Aber der HErr sprach zu Mose: Recke deine Hand aus über das Meer, daß das Wasser wieder herfalle über die Egypter, über ihre Wagen und Reiter etc.
\end{schriftzitat}

56. Als die Egypter nachjagten den Kindern von Israel, und sie ereilten, schreckte sie Gott durch einen Engel, daß sie fliehen wollen; aber das Meer schlägt wieder zusammen, behält und ersäuft sie alle.

57. Dies Mirakel soll in uns erwecken GOttes Furcht, daß wir uns daran spiegeln und schauen, wie GOtt die Gottlosen und Unbußfertigen strafen und stürzen könne, wie er denn allhier dem Pharao und Egyptern also thut. Daher Josua am 4. Capitel, V. 23. 24., davon sagt: „GOtt habe das Schilfmeer darum vertrocknet, und sie hindurch gehen lassen, auf daß alle Völker auf Erden erkenneten die Hand des HErrn, wie mächtig sie sei, und daß man GOtt den HErrn allezeit fürchtete“ etc.

58. Nun müssen wir das letzte Stück auch handeln, als den Spruch St. Pauli 1 Cor. 10, 1. 2., da er sagt: „Unsere Väter sind alle unter der\footnote{Eislebensche: den Wolken.} Wolke gewesen, und sind alle durch das Meer gegangen, und sind alle unter Mosen getauft, mit der Wolke und mit dem Meer“, welches am selben Orte nicht ist ein Geheimniß oder Allegoria, sondern eine Historie. Wenn aber jemand Lust hat zu heimlichen Auslegungen, will eine Historie deuten, und in eine Allegoriam ziehen, der soll es auf das göttliche Wort ziehen, und sehen, wie man mit dem Glauben und Gewissen umgehe. St. Paulus zeucht es an als ein Exempel und Historie für sich; als sollte er sagen: Unsere Väter sind im Meer getauft, gleich als wir; das ist, sie haben damals GOttes Wort auch gehabt, daß sie geheißen wurden, sie sollen sterben und durch Christum eingehen in das ewige Leben; daß sie also durch das Wort und den Glauben an Christum erlöst sind. Derhalben so haben sie durch den Glauben auf das Wort dieselbige Gerechtigkeit erlangt, die wir auch haben. Abraham ist durch das Wort GOttes an Christum auch in das ewige Leben gegangen und getauft worden.

59. Das Getauftwerden an diesem Ort und sonst mehrmals in der heiligen Schrift (als, da Christus Matth. 20, 22. den zweien Söhnen Zebedäi zur Antwort anzeigt: „Könnet ihr euch taufen lassen mit der Taufe, damit ich getauft werde?“) bedeutet und begreift in sich eines jeglichen Fährlichkeit, Anfechtung oder Werk, von GOtt ihm befohlen und auferlegt. Das ist seine Taufe. Denn er hat GOttes Wort, und wird geführt in mancherlei Widerwärtigkeit, daraus ihm denn GOtt hilft durch das Wort, und wird ihm alle sein Leiden und Versuchung gleich als eine Taufe.

60. Daß aber St. Paulus [1 Cor. 10, 6.] sagt: „Solches alles widerfuhr ihnen zum Vorbilde“, haben etliche ausgelegt, als heiße das Vorbild eine Allegorie oder geistliche Deutung, aber es ist unrecht; sondern also heißt es: den Juden ist solches alles widerfahren, und diese That, als der Durchgang durch das rothe Meer, ist zur Figur, zum Exempel und Beispiel geschehen, anzuzeigen, daß [es] uns auch also gehen werde. Denn gleichwie sie ihre Noth und Fährlichkeit gehabt, daraus GOtt sie errettet hat, also haben wir unsere Noth und Taufe auch, daraus wir erlöst werden. Als sollte er sagen: Wer da will ein rechter Christ sein und selig werden, der muß getauft, das ist, gerollt, gewanzert und geplagt werden, Apost. 14, 22.; deß haben wir viel Exempel in der heiligen Schrift. Und es kostet mehr, die heilige Schrift recht lernen, denn auslernen, wie die Nottirer und Irrgeister ihnen selbst träumen, darum sie der Schrift bald satt haben.

61. Vorbild wird allhier nicht genannt ein Geheimniß, es ist nicht eine Figur gewesen, sondern ein großer Ernst; GOttes Wort, so da lebendig macht [Ebr. 4, 12.], und der rechte Glaube ist allda gewesen; darum so ist es ihnen nicht im Schein geschehen, sondern es ist die That selbst da gewesen. Die Figur oder Vorbild muß man nicht auf die geistliche Deutung allein ziehen und deuten; es ist ein Exempel gewesen, daß, wie sie gethan haben, also müssen wir auch thun.

62. Wiewohl alle Historien ihre Auslegung auch haben, damit sie Christum bedeuten, dennoch so sind sie für sich auch ein Ernst, und sind
\quellseite{461}
nicht allein die Hülsen, sondern auch der Kern. St. Paulus 1 Cor. 10, 4. spricht: „Sie haben einerlei Trank mit uns gehabt“, das ist, einerlei Glauben an Christum, wie wir auch haben. Wie sollten sie denn allein eine Bedeutung geführt haben, und nicht erschnappt haben Christum selbst? Ja wohl. Also lerne, daß es nicht eine Bedeutung oder Figur sei, sondern ein Exempel.

63. Ein jeglicher muß in seinem rothen Meer gehen, wenn er herdurch haben soll. Da lerne, wie du das rothe Meer zertheilen sollest, und dir eine sichere Straße machen. Siehe dich nicht um nach deinen guten Werken, sondern nimm den Stab, das ist, das göttliche Wort. Wenn du das treibst, dann kommt der Wind, das ist, der Heilige Geist kommt mit dem Worte und macht Bahn, daß du denn sagen kannst: Ei, wenn dies Unglück nicht gewesen wäre, so wäre ich dahin, und nicht also hindurch gekommen, wie nun. Das Wasser muß den Kindern von Israel helfen und ihre Mauer sein; also, was uns sonst hindern will und schädlich sein, das muß uns zu Nutz und allem Guten gedeihen.

\section{Das fünfzehnte Capitel}

\begin{schriftzitat}
V. 1. Da sang Mose und die Kinder Israel dieses Lied dem HErrn, und sprachen: Ich will dem HErrn singen, denn er hat eine herrliche That gethan, Roß und Wagen hat er in das Meer gestürzt.
\end{schriftzitat}

1. Wir haben nächst gehört von einem großen Mirakel und Wunderwerke, so GOtt an dem Volke Israel gethan, da er sie durch das rothe Meer geführt, und den König Pharao mit seinem Kriegspolke darinnen ersäuft hat. Welches eine sonderliche Wunderthat und erschrecklich Werk GOttes gewesen, und darum geschehen, auf daß dies Volk Israel an GOtt glauben, ihm vertrauen, und für solche unaussprechliche, große Wohlthat danken und dienen sollte. Denn er sie dem Tode gleich aus dem Rachen und Zähnen gerissen. Es ist aber solche Geschichte auch darum beschrieben, daß wir aus dieser wunderbaren That erkennen sollten, wer GOtt wäre, nämlich, der auch mitten im Tode helfen könne. Wie denn der 68. Psalm, V. 21., auch sagt: „Wir haben einen GOtt, der da hilft, und den HErrn HErrn, der vom Tode errettet.“ Darum so soll man ihm vertrauen, daß er alle Sachen gut machen werde.

2. Dazu sollen wir dieser Historie recht gebrauchen. Denn wie die Kinder von Israel haben vom HErrn Hülfe bekommen, also will er unser Nothhelfer auch sein [Röm. 15, 12.]. Die Kinder Israel sehen die Egypter am Ufer des Meeres alle todt und ersoffen; darum so fürchtet das Volk den HErrn, und glauben an ihn, und von der Zeit an erkennen sie, was sie für einen GOtt haben, und lernen ihn fürchten und ihm gehorsam sein. Solche Andacht war bei den Kindern von Israel groß, wenn sie lange geblieben wäre; aber sie währt nicht länger, denn ein Tanz und hohe Messe.

3. Nun folgt Mosis und der Kinder Israel Freudenlied oder Lobgesang, da sie GOtt für diese großen Wohlthaten danken; welchen Psalm man in Israel hernach für und für alle Jahre gesungen hat. Denn es ist die Weise und Gebrauch gewesen, daß man von GOttes Wunderwerk gesungen, oder, wie man zu sagen pflegt, ein Spiel angerichtet hat, davon man singen und sagen sollte. Denn GOtt handelt also mit uns, daß wir beide Ohren voll haben, und überall Ursache und Anreizung genug, ihn zu preisen, loben und ehren. Deß haben wir allhier an Mose ein Exempel, welchem David auch hat nachgefolgt und viel herrlicher Psalmen und Gesänge gemacht. Darnach haben andere Väter auch so gethan, wenn GOtt ein sonderlich Wunder an ihnen bewies, als, wenn irgend eine redliche Schlacht geschehen ist, und GOtt Sieg gegeben, daß man darnach davon gesungen hat. Und das ist billig, daß man GOttes Werk treibe mit Predigen und Singen, daß aller Welt Ohren davon voll werden.

\subsection{Ich will dem HErrn singen}

4. Das ist die Summa dieses Psalms, daß Moses und das Volk wollen GOtt singen und loben. Es soll ein Lobpsalm sein, da sie singen
\quellseite{462}
und sagen wollen von GOttes Allmacht, Ernst, Gewalt, und auch von seiner Gnade und Güte, daß er ein Wunderwerck an ihnen gethan hat. Ihrer Person halben haben sie kein Lied zu singen, denn es ist nichts in dieser That, davon sie sich rühmen könnten; sondern sie sind erschrocken und gleich[sam] halb todt gewesen, von wegen der grausamen Leibes- und Lebensgefahr, in welcher sie waren, als die jetzt sollten in das Grab gelegt werden. Daß aber etwas Gutes geschehen ist, haben sie unserm HErrn GOtt zu danken; dem singen sie auch Preis und Lob dafür, und danken ihm.

5. Darum sind dies eitel fröhliche Worte, die aus dem Herzen heraus brechen, welches GOtt erkannt hat, und will ihn nun preisen und loben. Denn es ist ein groß, rühmlich Werck gewesen, und darf sich diese That wohl sehen lassen und an Tag kommen. GOtt hat ein Meisterstück alhier bewiesen; es ist ein köstlicher, hochtrefflicher und großer Sieg gewesen; daß die Egypter im rothen Meer ersäuft sind, daran mag man GOttes Gewalt und Macht wider seine Feinde spüren und sehen.

\begin{schriftzitat}
V. 2. Der HErr ist meine Stärke und Lobgesang etc.
\end{schriftzitat}

6. Wir könnten es nicht thun; die sechsmal hunderttausend Mann wären so stark nicht gewesen, Pharaonem mit seiner Kriegsmacht zu tilgen etc. Wer nun zu GOtt also sagen kann: Du bist meine Stärke, derselbige ist anders gerüstet, denn wir faule Tropfen sind. Wie wird aber GOtt unsere Stärke? In den Sachen, die GOtt antrifft, und nicht allein in diesen, sondern in allen Dingen müssen wir verzweifeln an uns, daß wir nichts vermögen. Also geben sie ihm die Ehre und bekennen, daß sie keine Stärke noch Kräfte gehabt, und was geschehen sei, das sei durch seine Macht und Gewalt ausgerichtet, es sei nicht ihr Werck, sondern GOttes Gnade und Segen, er sei der rechte Kriegsmann, der ihre Feinde gestürzt habe.

7. Dieses ist ein hochmüthiger Trotz, daß man also sagen kann: „Der HErr ist meine Stärke.“ Und haben die Worte ihrer gar viel gesungen, aber wenig verstanden. Sie wollen aber sagen: Dieses ist GOttes Werck; Trotz sei jedermann geboten, der uns ein Leid thue, oder ein Härlein krümme [Röm. 8, 31. f.]. Wenn demnach GOtt meine Stärke und Kraft ist, welches Stärke wird mir doch können etwas abbrechen? Denn ich kann sprechen: Ob ich gleich ein armes Würmlein bin, so habe ich doch GOttes Stärke bei mir [Ps. 22, 7. Jes. 41, 10.]. Item, ob ich gleich nichts bin, und auf Erden schwach und kränk, und so matt, daß ich auch nicht den Fliegen wehren könnte, so mich stechen, dennoch so bin ich stark, denn GOtt ist meine Stärke. Das haben die Kinder von Israel aus der Erfahrung lernen müssen, daß sie an ihnen verzweifeln, und sagen: Wir wissen nirgend mit die Egypter zu schlagen, denn mit dem Glauben und Vertrauen auf GOtt. Und in diesem Vertrauen fasse ich in mich GOttes Stärke, und GOttes Stärke faßt mich darnach in sich, sintemal GOtt diejenigen nicht verläßt, die auf ihn vertrauen.

\subsection{Mein Lobgesang}

8. Das ist, mein Ruhm. Als wollte er sagen: Ich weiß nichts zu preisen, rühmen und zu singen denn von dem.

\subsection{Und ist mein Heil}

9. Mein Sieg, mein Heil in dieser Sache, mit dem Könige Pharaone; mein Heil, so den Feind überwältigt, der mich keck macht durch das Vertrauen auf sich, und tilgt den Feind. Das ist gesagt von der Gewalt, die da sitzt wider Noth und Tod. Wollte GOtt, daß wir diesen Vers auch treffen könnten, wenn wir unterliegen, daß wir wider die Sünde, den Tod und Teufel, und sonst alles, was uns bei will und zusetzt, könnten sagen: Ich habe einen Sieg, Stärke und Heil, so größer ist, denn du bist; Trotz, der mir ein Leid thue. Also überwindet man den Tod und den Teufel und alles Widerwärtige. Das ist ein Großes, daß ein armer, schwacher Mensch fühlt, er habe eine solche Stärke bei sich, daß er auch den Tod und Teufel überwinden möge. Die Israeliter haben mit dem Schwert Pharaonem in das Meer nicht gestürzt und ersäuft, sondern dieweil sie an GOtt hingen durch den Glauben und Vertrauen ihres Herzens. Darum so ist der Glaube eine göttliche Stärke, ein göttlich Werck und Ding etc.

\section[Das ist mein GOtt, ich will ihn preisen, er]{Das ist mein GOtt, ich will ihn preisen, er ist meines Vaters GOtt}

10. Vor großer Freude des Geistes, so die Christen aus den Wercken GOttes haben, reden sie von einerlei Ding auf mancherlei Weise; da
\quellseite{463}
läuft der Mund gar über mit ausbündigen, sonderlichen Worten. Wenn das Herz voller Freuden ist, und ein Ding wohl bedacht hat, dann folgt darauf, daß der Mund viel davon redet [Matth. 12, 34.]. Denn das Herz lehret den Mund, daß er spricht: Allhier ist GOtt, an den ich mich mit dem Glauben halte, bei dem will ich bleiben.

\subsection{Ich will ihn erheben}

11. Das ist, ich will ihn zieren. Eines Theils haben es gedeutet, als wollte er ihm einen köstlichen Tempel bauen. Also legen es die Juden aus, als wollte er GOtt eine Hütte, das ist, eine Wohnung zurichten, da GOtt nahe bei ihnen wäre. Aber das Wörtlein „erheben“ heißt hier zieren, schmücken. Denn die Juden haben GOtt den Tabernakel oder die Hütte nicht ehe gebaut, GOtt wollte es denn von ihnen haben, und hieß es sie, wie er es hernach den David hieß, daß sein Sohn Salomo sollte den Tempel bauen [2 Sam. 7, 12. 13.].

\subsection{Er ist meines Vaters GOtt}

12. Als sollte er sagen: Er ist mein GOtt, und meiner Väter GOtt, ja, wir meinen allein diesen GOtt, der mit unsern Vätern, Abraham, Isaak und Jakob, geredet, und ihnen sich offenbart hat, und verheißet ihrem Samen und Nachkommen, er wollte sie nicht verlassen, sondern sie mehren als den Sand am Meer und als die Sterne am Himmel. Aber wir gedachten, daß wir arm wären und steckten dem Tode und Teufel im Halse, meinten, wir hätten keinen GOtt, aber nun, weil wir hindurch sind durch das rothe Meer, so haben wir seine Stärke erfahren. Da will Moses sagen: Das möchte der GOtt sein, der mit meinen Vätern geredet hat, und ihnen Verheißung gegeben, daß er ihre Kinder aus dem Diensthause Egypten führen wollte; und das hat er nun mit der That redlich bewiesen, darum will ich ihn hoch preisen und erheben. Von diesem Erheben, wie singt Maria auch in ihren Lobgesange [Luc. 1, 46.]: „Meine Seele erhebt den HErrn“! Moses macht es wunderlich; er gibt GOtt viel Namen, damit er ihn will preisen, hoch loben, welches denn die Hebräi heißen „erheben“.

13. Das sind die rechten Werke, die GOtt gefallen, damit wir GOtt ehren, schmücken, zieren und schön machen, daß er schöner wird, denn die Sonne ist, und gegen ihm sonst nichts schön ist [Ps. 104, 1.]. Er darf es zwar nicht, daß man ihn erhebe, denn er sitzt über alle Himmel, im Thron seiner Herrlichkeit [Jes. 6, 1.]. Aber wenn ich mit meiner Predigt und mit meinem Lobe ihn ausbreite und von ihm predige, wie ein mächtiger und feiner GOtt er sei, alsdann wird er erhoben und bekannt gemacht [Ps. 145, 1. 2.]. Sonst ist GOtt unbekannt bei den Menschen, er taugt nichts bei ihnen, und stinkt bei ihnen. Darum, wenn er erhoben und geziert wird, daß durch meinen Mund, meine Predigt und Bekenntniß sein Lob unter die Leute kommt, daß der GOtt, so zuvor stank, nun heilig gepredigt und gehalten wird, dann ist es der schönste Kranz, den man ihm kann aufsetzen, und der hübscheste Rock, den man ihm mag anziehen, ja, das hübscheste Haus, so man GOtt bauen mag. Dieser Mund richtet viel aus, wenn ich GOtt also schmücke, daß die Leute viel von ihm halten, daß er den Leuten in das Herz kommt, und seinen andern bösen Gedanken von GOtt kriegen, oder aus Eingeben des bösen Geistes ihn verdächtig halten [Jac. 3, 9. Ps. 34, 2. 9.].

14. Allhier ist Moses ein Meister zu reden, und da stammelt er nicht, wie wir sonst droben [Cap. 4, 10.] gehört haben. Andere Narren mögen immer hinfahren, die ihm wollen Tempel bauen, Bilder aufrichten und schöne Kleider schenken; aber das ist der rechte Schmuck, wenn er in seiner Weisheit erkannt wird, daß man viel von ihm hält; als, daß GOtt getreu, wahrhaftig, allmächtig, gerecht und weise sei, das ist ein schöner Schmuck. Wie man sonst von einem Menschen sagt: Das ist ein ehrbarer Mann, denn er führt ein fromm, redlich Leben. Und im Sprüchworte sagt man, daß Zucht ein Weibsbild mehr ziert denn alle Kleider. Davon Salomon auch sagt im Buche der Sprüche [Cap. 11, 22.]: „Ein unzüchtig Weib, das schön geputzt ist, die ist nicht anders, denn wie eine Sau, die einen güldenen Ring in der Nase hat“; aber welch Weib züchtig und redlich ist, die hat den höchsten Schmuck. Denn ob wohl eine Hure mit eitel Perlen und Edelgesteinen, von den Fußsohlen-an bis auf das Haupt, geschmückt wäre, so ist sie dennoch eine Hure, da dagegen ein ander Weib in einen Kittel, das fromm und züchtig ist, alle ihren Hurenschmuck überwiegt und verstellet.\footnote{verstellen = in den Schatten stellen.}

\quellseite{464}
15. Davon, spricht Moses, will ich predigen, das soll mein Werck sein. Man kann GOtt kein wohlgefälliger Werck thun, denn von ihm predigen und ihn loben. Du darfst ihm nicht fasten, Kirchen bauen, oder dich zu Tode martern; damit kann ich geschmückt und schön gemacht werden; aber GOttes Schmuck ist, wenn ich das Lobopfer gebe, daß ich von ihm predige und unter die Leute bringe, und seine Macht und Stärke bekannt mache. Das soll ich thun; das heißt recht erheben.

\begin{schriftzitat}
V. 3. Der HErr ist der rechte Kriegsmann, HErr ist sein Name.
\end{schriftzitat}

16. Ich meine, der könne kriegen. Wie närrisch thut doch derjenige, so auf Harnisch, Spieß, Schwert und Büchsen seinen Trost setzt, und auf Festungen sich verläßt, dadurch Ehre und Ruhm sucht? [1 Sam. 17, 45. 47.] Willt du Kriege führen, so suche den rechten, billigen Krieg, der mit gutem Gewissen geschieht. Andere Leute kriegen, gleichwie die Säue mit den Hunden. GOtt ist allein der mächtige HErr, der da Glück und Sieg in Kriegen gibt, wie im 144. Psalm, V. 10., gesagt wird: Qui dat salutem regibus [der den Königen Sieg gibt]. Er hat befohlen den Kindern von Israel, daß sie mit geharnischter, gewappneter Hand sollten ausziehen aus Egypten, wie sie denn sechsmal hunderttausend streitbarer Mann haben. Nun ist die Macht und Gewalt der Egypter auch groß, welche sich alle wider den Kriegshaufen der Kinder von Israel legen. Aber was geschieht? Die Kriegsrüstung der Egypter wird in einem Hui alle vertilgt.

17. Die Israeliter hätten in acht Tagen nicht so viel Kälber (wenn sie auch gleich wären gebunden gewesen) schlachten und würgen können, als der Egypter gewappneter Kriegsleute allhier umgebracht werden; denn es ist gar ein gewaltiger König gewesen, der sich allhier wider dies Volk gelegt hat. Aber GOtt erregt und bringt ihn darum zu Felde, daß bewiesen werde, er sei ein rechter Kriegsmann; der weiß wohl, wo er den Geharnischten treffen solle. Wenn das Stündlein kommt, so stößt er es alles über einen Haufen, daß nicht Einer überbleibt von denen, die sich wider ihn oder seine Kirche gesetzet und aufgelehnt haben. Solches beweist mächtiglich diese That am Pharao und viel, viel andere Geschichten mehr bezeugen es.

18. Diesem HErrn soll man vertrauen, der eine solche Gewalt hat; und wer wollte an ihm zweifeln? Moses breitet diese Schlacht mit vielen, langen Worten aus, und sagt, wie es zugegangen sei. Zwar es ist nicht Kinderspiel gewesen, daß man hat sehen sollen, daß Pharao so viel schöner herrlicher Leute mit ihm gehabt, so viel Fürsten, Herren und weise, thätige, streitbare Leute und Helden, damit er gedenkt wohl die ganze Welt zu fressen; aber da es zum Treffen kommt, so gehen sie alle unter und versinken wie Blei im Wasser, müssen ersaufen, auf daß man sehe, GOtt sei der rechte Kriegsmann.

19. Dieses thäte GOtt noch heutiges Tages, wenn wir einen rechten Glauben hätten. Stände ein König auf, wie im 7. Psalm, V. 2., gesagt wird, so spräche ich: GOtt, auf dich will ich vertrauen. Wollen sie denn mein Fleisch umgeben und tödten, so müssen sie mir nicht schaden. Denn ich will mehr auf GOtt vertrauen, denn auf meine Stärke und Vermögen; denn prächtig und mächtig ist seine Hand. Moses sagt allhier: Du bist mein Heil und Stärke, du, du hast es gethan [Apost. 9, 5.]. Wie rühmen sich doch allhier die Kinder von Israel! Sie geben den Egyptern Schuld, daß sie wider GOtt gethan haben, und wider GOtt gefochten und gestritten haben.

20. So sollten wir auch lernen und uns dazu gewöhnen, GOtt die Sache anheim zu geben, daß man sagte: Man hat nicht wider mich, sondern wider dich, GOtt, gestritten. Ich bin deine Schwachheit, du bist meine Stärke; ich bin eitel Sünde, du lauter Gerechtigkeit; ich bin dein Unflath, du bist meine Zierde und Schmuck, da wird es wohl besser um uns stehen. Der Teufel sieht uns wohl für schwach, kraftlos und matt an; aber wenn es zum Treffen geht, so spricht GOtt: Ich bin allhier. Dann geht es gleich wie mit einem armen Würmlein, das an einem Fischangel klebt, welches der Fisch will fressen und verzehren, aber er wird darüber gefangen. Also läßt uns GOtt auch schwach sein, und ist dennoch bei uns mit seiner Stärke, und erhält uns. Dann trotzen die Feinde einher und wollen uns gar vertilgen und fressen, aber bald liegen sie zu Boden; ob sie wohl meinen, sie wollen uns auf einen Bissen fressen, so fehlt es ihnen doch. Denn GOttes Stärke ist in unserer Schwachheit allda verborgen gelegen, und überwindet alle Gefährlichkeit.

\quellseite{465}
21. Ich soll sagen zu GOtt: Ich bin dein Knecht, dein Glaube und dein Wort ist in mir gewesen, und weil ich dein Wort gehabt habe, so wollen sie mich unterdrücken; aber siehe zu, wie sie sein sollen anlaufen, die da GOtt wollen unterdrücken; es soll ihnen gehen, wie es allhier den Egyptern geht. Sie sind her, und wollen mich fressen, wissen aber nicht, daß sie dich, GOtt, fressen wollen; das werden sie nicht hinausführen. GOtt gebe, daß wir es auch einmal also rühmen können.

\begin{schriftzitat}
V. 7. Denn da du deinen Grimm ausließest, verzehrete er sie wie Stoppeln.
\end{schriftzitat}

22. In diesen Worten erzählt Moses den großen Ernst GOttes gegen seine Feinde, wie die Verfolger seiner Kirche und Christen sollen alsobald untergehen. Denn gleich als sich das Stroh oder Stoppeln wider ein groß Feuer und Glut nicht aufhalten noch wehren kann, also ist es auch, wenn Menschen wider unsern HErrn GOtt streiten und fechten [Jes. 40, 24. 41, 1. 2.]. Es könnten hochmüthigere und trotzligere Worte nicht sein, denn Moses und die Kinder von Israel allhier gebrauchen, da sie vom Pharao und den Egyptern errettet sind, daß sie sagen, GOttes Zorn sei ein verzehrend Feuer, und GOttes Feinde sind Stroh oder Stoppeln.

23. Hier haben sie viel einen andern Muth und Sinn, auch eine andere Zunge und Sprache, denn sie zuvor gehabt. Erstlich gedachten die Kinder von Israel, sie wären eitel Wasserblasen gegen Pharao und den Egyptern. Aber nun sehen sie, daß, wenn einer unserm HErrn GOtt vertrauet, so müssen seine Feinde sein als Stroh und Stoppeln gegen einem Feuer, da denn die Stoppeln einem Feuer nicht Schaden thun mögen, sondern das Feuer je länger je größer machen, und zuletzt vom Feuer verzehrt, zu Staub und Asche gemacht werden.

24. Man könnte von der großen Gewalt Pharaonis und der Egypter nicht verächtlicher noch schimpflicher reden, denn allhier Moses und die Kinder von Israel thun, daß sie diesen mächtigen, gewaltigen, prächtigen König, so viel Fürsten, den Herrn und das höchste Vermögen seines ganzen Reichs, Stoppeln und Stroh heißen. Wenn jetzt einer den römischen Kaiser, die Churfürsten und Fürsten, das ganze Reich Stoppeln nennete: würde man nicht sagen, er wäre toll und thöricht? Aber Mosis Herz ist gar voll von göttlicher Hülfe und Gewalt, deren er sich allhier freut und GOtt dafür dankt.

25. Das soll uns nun zum Troste dienen, auf daß alle dasjenige, so uns anficht und beschweret, ob es wohl einen Schein hat, als sei es groß und wolle uns gar fressen und verschlingen, nicht erschrecke. Denn wenn wir mit dem Glauben solches ansehen, und unsere Schwachheit fühlen, auch die große Macht unserer Widersacher anschauen, so können wir doch GOttes Stärke uns trösten, und der Widersacher Macht und Gewalt verachten. Denn sie sind gegen GOtt gleich als Nichts. Ja, die Welt, der Himmel und Erde und alle Creaturen sind gegen GOtt gleichwie ein Tröpflein Wassers gegen einem Eimer voll Wassers [Jes. 40, 15.].

26. Dieses erkennt man aus vorgenommener Historie, da Moses zu verstehen gibt, wie nach ihrem Durchgange das Meer wieder durch einander gefallen sei, gleich als wollte Himmel und Erde untergehen. Denn es ist das Wasser so weit von einander gestanden, weiter denn eine große Stadt lang oder breit ist. Da es nun plötzlich wieder zusammen schlägt, wie meinst du, was es muß für ein Brausen und Gepraffel gewesen sein? Es wird nicht anders gethan haben, denn als wollte alles über einen Klumpen fallen. Erschrickt doch einer vom Brausen und Hauchen eines Wassers; was sollte denn allhier nicht geschehen sein? Zuvor hieß GOtt das Wasser stille stehen, wie Moses allhier [V. 8.] sagt: „Durch dein Blasen thaten sich die Wasser auf, und die Fluthen stunden auf Haufen.“ Geist und Wind heißt ein Ding bei den Hebräern. Er ließ einen Grimm kommen, oder er schnaubte einmal, da sahe ich deine Kraft und Gewalt, so alle die Stärke und Macht der Egypter (gleich als das Feuer die Stoppeln verzehrt) hat vertilgt.

27. Moses beschreibt nicht allein, wie sie gesinnt sind gewesen, und wie sie sich gefühlt haben, daß sie gar kleinmüthig und verzagt gewesen; denn da sind hohe Berge, haben sie gedacht, wo sollen wir hinfliehen? da ist nichts, denn eitel Tod vor unsern Augen; sondern meldet auch der Feinde Sicherheit und Ruhm, da sie die Kinder Israel beschlossen haben, gleich als wenn man die Säue zusammen in einen Stall bringt. Darum schreien sie: Gewonnen, gewonnen! und ist bei ihnen eitel Kraft und Stärke und gewisse Hoffnung des Sieges. Dagegen der Israeliter
\quellseite{466}
Haufe muß seufzen und sagen: O wir sind des Todes, und alle verloren etc. Da läßt sich GOtt nicht sehen mit seiner Stärke und Gewalt, sondern er ist in großer Schwachheit bei den Kindern von Israel; noch erhält er sie, und stürzt den Pharaonem, und bläst den Kindern von Israel einen Muth ein.

28. Dieses sind treffliche Exempel und Worte eines aufrichtigen Glaubens, daß wenn ein Feind tobt, wüthet und stolziret, und läßt sich dünken, der Himmel hänge alles voller Geigen, du aber bist dagegen kraftlos und ohnmächtig, daß du dann schließen und sagen könnest: Lieber HErr GOtt, du bist mein, und ich bin dein. Und das ist GOttes Werk, dazu hat er Lust, wie auch der 73. Psalm, V. 18., singt, daß GOtt etliche hoch erhebt, auf daß er sie wieder niedrige und stürze, und die Jungfrau Maria singt es auch in ihrem Lobgesange [Luc. 1, 52.]: Deus deposuit potentes de sede etc. [, er stößet die Gewaltigen vom Stuhl, und erhebet die Elenden]. Denn wenn ich auch ein Ding hart niederwerfen will, so hebe ich es hoch empor; auf harte Büsse sparet man hoch. Darum hebt GOtt die Feinde hoch, daß sie hart fallen sollen. Also ließ er den Pharao und die Egypter auch hoch herfahren, gleich als hätten sie die Kinder von Israel gar in der Hand; aber wenn man zu sicher wird, so thut das Wasser seinen Schlund zu und frißt sie; da liegt denn Pharao im rothen Meer, und Sennacherib verliert vor Jerusalem sein Kriegsvolk [2 Kön. 19, 35.]. Antiochus, der tyrannische König, kommt auch jämmerlich um in Persienlande [1 Macc. 6, 8. ff.].

29. Daraus sollen wir lernen, daß GOtt ein solcher Mann ist: wen er erhebt, mit dem steht es gefährlich; was er aber sinken läßt, das bedeutet etwas Gutes [Ps. 113, 7. Sir. 3, 19.]. Das Hochheben ist schrecklich, denn er zerbricht es auch gerne, und macht das daraus, daß es nicht war. Solches steht allein dem Schöpfer zu. Das thut er; nämlich, was hoch und groß ist, das stürzt er zu Boden, und was nichts ist, das erhöht er. Wie auch der HErr Christus selbst [Luc. 16, 15.] zu den Pharisaern sagt, es sei vor GOtt ein Greuel, was vor den Menschen hoch und groß ist. Darum, wenn sie meinen, sie sind alles, und stehen auf festen Beinen, so liegen sie zu Boden. Wiederum, was verachtet und verworfen ist, des nimmt sich GOtt an [Ps. 10, 12.]. Das ist also unsers HErrn GOttes Art und Weise.

\begin{schriftzitat}
V. 10. Und sanken unter wie Blei im mächtigen Wasser.
\end{schriftzitat}

30. Gleich als wenn man ein Blei oder Kloß in das Wasser wirft, welches bald zu Grunde sinkt, also sind in einem Hui die Egypter auch ersoffen, daß in einem Augenblicke weder Mann, Roß noch Wagen da ist, und sie auch nicht mehr schwimmen noch einen Finger regen können [1 Thess. 5, 3.].

\begin{schriftzitat}
V. 11. HErr, wer ist dir gleich unter den Göttern?
\end{schriftzitat}

31. Der Heilige Geist ist gar voller Wörter, er geht gar über mit Rühmen, gleichwie ein Roß, darinnen ein Most gieret,\footnote{gieren = gähren.} also schäumet Mosis von übriger Freude und Gedanken. Es sind mancherlei Götter auf Erden gewesen; aber kein Gott hat so eine herrliche That gethan, als der rechte GOtt. GOtt verhängt und läßt zu, daß andere Götter auch aufgeworfen werden; wie St. Paulus [1 Cor. 8, 4. 5.] davon sagt, daß sich die Teufel für GOtt ausgeben und wollen angebetet sein, aber es ist nur ein einiger GOtt über die Bösen und über die Gottesfürchtigen. Aber das ist der Unterschied, daß der Dienst und Amt ungleich sind. Denn dieser will auf eine solche Weise, jener auf eine andere Art GOtt dienen, und fehlen alle GOttes. Gleichwie aber nur ein einiger GOtt ist, also ist auch nur einerlei Weise GOtt zu dienen. Denn die andern wollen unserm HErrn GOtt auch dienen, aber sie dienen dem Teufel.

32. Darum, von wegen so mancherlei Abgötterei, grausamer Sünden und Mißhandlungen, dazu der Teufel gar hurtig, unverdrossen und willig ist, verhängt GOtt, daß ein Volk das andere dämpft und tilgt. Denn der Satan ist der Welt Gott, er hat mancherlei Engel, die ihm dienen. Also, in diesem Capitel wird der König Pharao und seine Fürsten erschlagen, aber GOtt thut es durch den Teufel. Wiewohl die Gottlosen oft die Frommen auch unter sich zwingen und eine Schlacht gewinnen; denn der Satan hilft und schüzt zuweilen seine Diener und Volk auch. Diese Gewalt hat er wohl, aber es ist nichts gegen der Stärke des, so einen Himmel über den andern gebauet hat; darum so ist er ein GOtt über alle Götter, er hat einen Himmel
\quellseite{467}
über alle Himmel gebauet. Wir heißen das einen Himmel, so weit wir ihn sehen; über diesen sichtbaren hat er andere Himmel mehr.

33. Also ist nur Ein HErr und GOtt, ob wohl die Menschen ihnen selbst andere Götter mehr träumen und erdichten. In Hispania haben sie St. Jakob zum Patron, daselbst unter dem Namen St. Jakob ehren und dienen sie dem Teufel. Zu Rom haben sie St. Peter; anderswo hat man St. Catharina, St. Barbara und Nicolaus als Nothhelfer angeruffen. Der Teufel ist ein Gott und Fürst der Welt; darum so ist er mächtig und gewaltig. Er kann zuweilen helfen, und GOtt verhängt es also; denn GOttes Rathschläge sind wunderbarlich. Darum will Moses sagen: Es ist nichts mit den Strohpoppen\footnote{Strohpoppen = Popanz von Stroh.} und Göttern, gegen dir. Er bekennt, daß sie Götter heißen; aber er spricht: Wer ist dir gleich unter den Göttern, der so hehr und heilig, schrecklich, löblich und wunderthätig sei?

34. Das ist gleich der Beschluß des ersten Theils, da in diesem Lobpsalm Moses GOttes Gewalt und Ernst gegen die Verfolger und Feinde seines Worts gepreiset hat, da er will sagen: Da ist keiner. Es sind viel Götter auf Erden; aber es ist keiner so trefflich in Heiligkeit, als du. Denn alle anderen Götter, die man erfunden und sich ausgeben, die schänden sich selbst; aber wer diesen GOtt ehrt, der wird heilig, und er heiligt auch GOtt. Die Mönche haben St. Bernhardum, Benedictum und Franciscum geehrt; aber dem Teufel unter dem Namen der Heiligen gedient, und unter dieser Heiligen Namen nur Schande angerichtet; aber dies Volk hat den rechten GOtt und den rechten Geist; und wie GOtt heilig ist, also gehen sie auch in Heiligkeit. Andere Götter sind unwendig unflätig und unrein; es ist keiner so schrecklich, löblich und wunderthätig, als du.

\begin{schriftzitat}
V. 12. Da du deine rechte Hand ausstrecktest, verschlang sie die Erde.
\end{schriftzitat}

35. Das ist, sie waren unter dem Himmel; aber da das Meer und Wasser zusammen fiel, da tauchten sie zu Grunde und fielen in das Erdreich; beide Theile schlugen zusammen, und sind die Egypter mitten im Wasser, darum so sind sie in Abgrund der Erde gesunken.

\begin{schriftzitat}
V. 13. Du hast geleitet durch Barmherzigkeit dein Volk, das du erlöset hast, und hast sie geführet durch deine Stärke zu deiner herrlichen Wohnung.
\end{schriftzitat}

36. Alhier dankt und preist er GOtt über der herrlichen Wohlthat, Gnade und Barmherzigkeit, daß er nicht allein sein Volk, die Kinder von Israel, aus Egypten geführet und erlöst hat, sondern auch ihr Geleitsmann gewesen ist, und sie regieret auf dem Wege aus Egypten nach dem rothen Meer, und sie auch durch das rothe Meer gebracht. Der Engel des HErrn ist vor ihnen hergegangen des Nachts in einer feurigen Säule, und des Tages in einer weißen Wolke etc. Welche Wolke sich setzte vor der Egypter (Samptstraße\footnote{So von uns herübergenommen aus der alten Ausgabe Walchs. (Eislebensche: Straf.)}) zwischen die Egypter und Kinder von Israel, auf daß sie nicht konnten zusammen kommen.

37. Das meint er nun alhier, daß es den Egyptern nach ihrem Vornehmen nicht gegangen ist, sondern GOtt ist bei den Israeliten gestanden, auf daß sie gewiß wären, sie sollten nicht ausziehen oder fortdrücken, es gefiele denn GOtt wohl, und er wollte bei ihnen sein. Davon sagt Moses: Es war deiner Güte Schuld, nicht unseres Verdienstes oder guten Werke. Denn wenn du nach unserm Verdienste handeln solltest, so hättest du uns wohl ungekeitet und ungeführet gelassen.

38. Ihr könnt denken, daß aus der großen und erschrecklichen That ein groß Geschrei erfolgt sei in die umliegenden Länder, so von der Wundterthat gehört, daß die ganze Macht des Königreichs Egypten im rothen Meer ersoffen sei, und daß die Kinder von Israel das Volk wären, welches trocken durch das Meer gegangen wäre, und alle Länder gedacht haben: Hilf GOtt, wer will vor diesem Volke bleiben? Denn wenn der Türke in diese Lande herein fiele, und mit trocknen Füßen durch die Elbe zöge, welch ein groß Geschrei würde durch alle Lande davon werden? Wie würde jedermann das Herz entfallen? Viel mehr ist dies alhier geschehen, als sich dies herrliche Wundterwerck hat zugetragen. Aber dieselbigen Lande sind dennoch zu GOtt nicht bekehrt worden, sondern gedachten, es wäre irgend durch Zauberei geschehen. Denn es ist damals viel Schwarzkünstelei und Zauberei im Brauche gewesen, dazu die Abgötterei weiblich geholfen. Ob nun wohl Israel
\quellseite{468}
dem rechten GOtt gedient, so dachten doch die Heiden: Ein Gott ist stärker, denn der andere. Item, ob gleich der Israeliten GOtt Eines gelungen hat, so muß darum der Egypter Gott nicht unrecht sein.

\begin{schriftzitat}
V. 14. Da das die Völker hörten, erbebeten sie, Angst kam die Philister an.
\end{schriftzitat}

39. Das ist, die umliegenden Länder, als, die Philister, Edomiter, Moabiter, Cananiter und andere fürchteten sich, sie taumelten, liefen zusammen, es erhob sich ein Rumor, da dies geschehen ist. Die Philister waren die nächsten Nachbarn, darum gedenkt er ihrer, und waren die Philister ein stark, mächtig Volk, wohnten gegen dem Abend des Landes, da sie hinziehen sollten. Die Kinder von Israel sind von Egypten herein gezogen, als von Abend gegen Morgen, als wollten sie in das Land gegen Mitternacht ziehen; welches die rechte Bahn und Straße war auf das Land, da die Philister wohnten, welches Land viel großer Städte, und Fürsten, und Riesen hatte; aber GOtt führte sie lange umher in der Irre, lenkte sie um die Berge herum, um ihrer Sünde willen. Diesen Philistern ward auch bange vor dem Geschrei. Das Volk Israel hat nie ganz und gar der Philister mächtig werden können, daß sie ihnen wären unterthänig gewesen; zinsbar sind sie den Kindern von Israel gewesen, weiter hat es auch David nicht bringen können [2 Sam. 5, 17. ff. 1 Chron. 15, 8. ff.].

\begin{schriftzitat}
V. 15. Da erschraken die Fürsten Edom, Zittern kam die Gewaltigen Moab an, alle Einwohner Canaan wurden feige.
\end{schriftzitat}

40. Die Edomiter wohnten auf der linken Hand, und kamen von Esau, Jacobs Bruder, her, und die Kinder von Israel zogen fast bei vierzig Jahren um ihr Land her. Moab lag etwas weiter gegen Morgen. Edom war hierwärts besser gelegen. Ueber diesen Völkern lag Canaan. Diese Länder allzumal, die sie vor sich hatten, haben sich über diesem Wunderwerke entsetzt, wiewohl es gewaltige, große, mächtige und starke Völker waren. Dennoch dachten sie, wie sie wollten die Kinder von Israel schlagen und dämpfen. Das ist die That, wie es mit den Kindern von Israel ergangen, und wie sie erlöst sind. Nun schließt er den Lobgesang mit einem Gebet, und spricht:

\begin{schriftzitat}
V. 16. Laß über sie fallen Schrecken und Furcht, durch deinen großen Arm, daß sie erstarren wie die Steine.
\end{schriftzitat}

41. Als sollte er sagen: Aus dem Geschrei dieses großen Wunderwerks laß sie feige werden; drücke nach, lieber HErr Gott, drücke nach, sie haben Sorge, Furcht und Erschrecken, und wir müssen herdurch und sollen herdurch. Das rothe Meer haben wir überwunden, ei, so wollen wir diesen Haufen der Philister, Edomiter, Moabiter und Canaaniter auch schlagen und überwinden. O hilf, daß wir es bald thun, und thue du es; denn du bist allein mächtig. Du hast neulich den Pharao und die Egypter im rothen Meer ersäuft, darum hoffen wir, du werdest andere Völker mehr vor uns her vertilgen, und uns in das gelobte Land bringen.

42. Wie soll er es aber thun? Nimm ihnen das Herz! Denn diese Kunst hat GOtt, er hauet einem nicht eine Faust oder Bein ab, sondern er nimmt die Mannheit und den Muth, daß einem das Herz entfällt, daß kein Muth oder Mannheit da ist. Wie der 76. Psalm, V. 12. 13., auch sagt: „Bringet Geschenke dem Schrecklichen, der den Fürsten den Muth nimmt, und schrecklich ist unter den Königen auf Erden.“ Wenn der Muth hinweg ist, dann steht man wie ein Narr, ja, es kann einer kaum auf den Beinen stehen, Arme und Beine werden zitternd und machtlos, daß einer da liegt wie ein Kloß, läßt auf sich hauen und stechen, als auf einen Kloß oder Holz. Ein Knabe kann einen Mann erwürgen. Also läßt unser HErr GOtt einem die Fäuste ganz, den Harnisch und Spieß und das Messer unentzogen. Aber er gibt ein verzagt Herz, wenn man mit dem Feinde schlagen soll. Darum liegt es nicht an der Faust, noch am Schwerdte, sondern der Muth, der thut es. Wenn einer ein Ding in Sinn nimmt, daß er es thun darf, so ist es schon halb geschehen. Denn hüte dich vor dem, der dich mit Ernst meint; denn wenn es im Sinne also beschlossen ist, so geht die Faust bald hernach.

43. Dieses ist unseres HErrn Gottes Behendigkeit, seine Kriegskunst, Vortheil und bester Harnisch, nämlich, daß er den Scharrhansen und Leutefressern den Muth nimmt; dann wird einer geringer denn ein Kind, und blöder denn ein Weib. Darum spricht Moses allhier: Nimm ihnen das Herz und den Muth, und gib es uns; laß über sie fallen Schrecken und Furcht, daß
\quellseite{469}
sie erstarren, wie die Steine. Dann wollen wir sie überwinden; wenn ihnen die Mannheit entgeht, dann werden sie matt und halb todt sein, daß sie nicht einen Finger werden regen können.

\subsection{Bis das Volk hindurch komme, das du erworben hast}

44. Als sollte er sagen: Es ist doch dein Volk, es liegt dir dran, auf daß du geehrt und gepreiset werdest; darum nimm ihnen den Muth, und gib uns ein Herz. Also streitet unser HErr GOtt wider seine Feinde; da mag denn danach streiten, wer da will, ich will mit diesem Krieger zufrieden sein [Ps. 108, 12. 13.].

\subsection{Bringe sie hinein, und pflanze sie auf dem Berge deines Erbtheils}

45. Das ist, bringe sie in das Land Canaan, treibe die Heiden aus. „Dein Erbe.“ Denn das Land ist dein, so sind wir dein Erbtheil; darum wirst du, GOtt, allda eine Wohnung, Tempel und Sitz aufrichten, da du wirst gefunden werden.

46. Denn GOtt hat allezeit also gethan, daß er auf Erden gegeben hat ein leiblich Zeichen, eine Person, Ort und Stätte, da er gewißlich hat wollen gefunden werden. Denn wo wir nicht durch ein leiblich, äußerlich Zeichen gebunden und gefangen werden, so wird ein jeglicher GOtt suchen, wo es ihn gelüstet. Darum haben die heiligen Propheten viel geschrieben von dem Tabernakel, von der Wohnung und Hütte, da er gegenwärtig sein wollte. Also hat GOtt stets gethan. Dergleichen hat er uns Christen auch einen Tempel gebaut, da er wohnen will, nämlich, das mündliche Wort, die Taufe und das Abendmahl, welches da sind leibliche Dinge. Aber unsere falschen Propheten, Rottengeister und Schwärmer verachten es, und werfen es hinweg, gleich als tauge es nichts, und sagen: Ja, ich will sitzen und warten, bis mir ein fliegender Geist und Offenbarung vom Himmel komme.

47. Aber hüte dich davor. Wir wissen es auch wohl, daß Wasser, Brod- und Wein uns nicht selig machen; aber wie gefällt dir das, daß im Abendmahl nicht schlecht Brod und Wein, oder auch in der Taufe pur lauter Wasser ist; sondern GOtt spricht, daß er in der Taufe sein will, sie soll uns von Sünden reinigen und waschen? Und im Abendmahl, unter dem Brod und Wein, wird der Leib und Blut des HErrn Christi gegeben. Willst du nun allhier GOtt und sein Zeichen verachten, und das Wasser in der Taufe ansehen und halten, gleich als das Wasser, so in der Elbe fließt, oder damit du kochst? Oder, willst du das Wort des Evangelii gleich achten dem Worte oder Reden, so Bauern in einem Kretzschmar oder Tabern reden? Denn GOtt hat gesagt: Wenn das Wort von Christo gepredigt wird, dann bin ich in deinem Munde, und ich gehe mit dem Worte durch deine Ohren in das Herz.

48. Darum so haben wir ein gewiß Zeichen, und wissen, wenn das Evangelium gepredigt wird, so ist GOtt gegenwärtig da, er will sich daselbst finden lassen; daselbst habe ich ein leiblich Zeichen, dabei ich GOtt erkennen und finden möge. Also ist er auch bei der Taufe und Abendmahl; denn er hat sich verbunden allda zu sein. Laufe ich aber zu St. Jakob, oder in das Grimmenthal,\footnote{„in Grymthal zu Franken“, Jäger, Carlstadt, S. 188, Anm.} gehe in ein Kloster und suche GOtt anderswo, da werde ich seiner fehlen. Und wenn jetzt die Rottengeister also predigten: Gleichwie das Klosterleben, Anrufung der Heiligen, Messe und Wallfahrt nichts ist, also ist die Taufe und Abendmahl auch nichts; das klappt noch lange nicht. Denn es ist ein großer Unterschied, wenn GOtt etwas ordnet und einsetzet, oder wenn Menschen etwas stiften. Ja, du sollst GOttes Ordnungen und Stiftungen glauben, sie anbeten und in großen Ehren halten; also hat er es Mosi auch befohlen.

49. „Bringe sie in das Land“, das ist, ordne und mache namhaftig einen gewissen Ort, auf daß, wer nicht persönlich daselbst dich anbeten kann, daß derselbige seinen Leib hierher kehre, und sein Angesicht dahin wende und bete. Also habe ich GOtt auch an einem gewissen Ort, nämlich, allhier im Worte und Sacramenten, daß, wenn gleich einer zu Rom ist, oder wo er sonst sein mag, wenn er sein Angesicht zum Worte und Sacramenten nur kehrt und anbetet, so findet er allda unsern HErrn GOtt, und wenn er sich auch gleich in einem Strohhaufen wollte finden lassen, so sollte man ihn daselbst suchen und ehren. Darum so haben die Propheten gesagt: Wir wollen zu seinem Fußschemel kommen, denn er ist heilig; daselbst wollen wir ihn anbeten.

\quellseite{470}
50. Aber unsere Propheten (so viel klüger sind denn jene; in welchen der Geist GOttes sich thätlich sehen ließ) sagen: Meinst du, daß dies Holz, Wasser, Brod und Wein heilig sei? Das weiß ich selbst wohl, daß es ohne das Wort und seiner Natur halben die Heiligkeit nicht habe. Aber wenn GOtt will diesen Dingen durch sein Wort Heiligkeit und Kraft geben, so wird es nicht Mangel haben, sondern GOtt, der darauf sitzt und diese Dinge heiligt, der ist in seinem Worte. Wenn er nur sagt: Allhier bin ich, und da ist mein Wort, da findet man GOtt; von dem wird es geheiligt. Darum, daß er allda will sein, derhalben so will ich ihn, oder das äußerliche Zeichen nicht verachten.

51. So will nun Moses sagen: Laß sie einsitzen und im Lande einzuwurzeln, und beständiglich daselbst einnehmen, und nicht da bleiben gleich als Gäste, die da essen und trinken, und nur eine Nachtherberge daselbst hätten, und wieder von dannen ausreißeten; sondern laß sie zunehmen und wachsen, daß sie drinnen gepflanzet werden, und wesentlich oder beharrlich bleiben mögen.

\subsection{Auf dem Berge}

52. Das ganze Land nennt er Gebirge; und redet Moses davon, gleich als wäre es schon geschehen, da es doch noch nicht geschehen war, sondern GOtt hat es im Sinne, seinen Tabernakel und Hütte dahin zu bauen. Das Land war sonst nicht so köstlich seiner Früchte halben, sondern davon, daß GOtt da wohnen wollte. Und GOttes Wohnung und Hütte soll man nicht rechnen nach Holz und Stein, sondern daß GOtt daselbst redet, wie der 60. Psalm, V. 8., auch sagt: „GOtt redet in seinem Heiligthum, des bin ich fröhlich.“ Und also ist es auch, wo GOtt redet in seinem Worte und in der Predigt, das rühmt Moses hoch, daß es ein köstlich Land sei, und spricht: Pflanze sie an den Ort, da du deine Wohnung gemacht hast.

53. Daher sagt auch ein anderer Psalm [Ps. 122, 1.]: „Ich freue mich, daß zu mir gesagt worden ist: Wir wollen in das Haus des HErrn gehen“; nicht da Orgeln und Pfeifen sind, sondern da der HErr unser GOtt redet; daselbst geht an Leben, Seligkeit und Barmherzigkeit. So haben die Propheten GOtt gehabt, daß sie an einem leiblichen Zeichen gewiß gewesen sind, daß GOtt da wäre; sonst laufen wir hierher und dorthin, und wissen nichts Gewisses von GOtt. Daher ist dieser in ein Barfüßlerkloster gelaufen, jener ein Carthäuser worden; einer tappt hier, der andere da. Denn GOtt sitzt nicht daselbst, redet auch nicht an demselbigen Orte, lockt auch nicht zu sich; wie er sonst thut, wo sein Wort gepredigt wird.

54. O! das wollen die großen Geister nicht leiden. Aber wohlan, laß toben und wüthen alle Heiden; er will ein Reich anrichten, das in Ewigkeit bleiben soll, wie er denn diesen Lobgesang beschließt und spricht [V. 18.]: „Der HErr wird König sein immer und ewig.“ Denn er ist auch allein ein wahrhaftiger, ewiger, allmächtiger, ernstlicher und barmherziger, gnädiger König und HErr, der Pharao im rothen Meere hat Mores gelehrt und ersäuft, uns Israeliter aber mit trockenen Füßen durch das Meer gehen lassen und erlöst. Nun folgt ein neu Capitel, das man von den andern unterscheiden sollte.

\begin{schriftzitat}
V. 20. 21. Und Mirjam, die Prophetin, Aarons Schwester, nahm eine Pauke in ihre Hand, und alle Weiber folgten ihr nach hinaus mit Pauken und Reigen. Und Mirjam sang ihnen vor: Lasset uns dem HErrn singen, denn er hat eine herrliche That gethan, Mann und Roß in das Meer gestürzt.
\end{schriftzitat}

55. Mirjam ist Mosis und Aarons Schwester gewesen, und nach den Sprachen werden die Namen verändert, als, Hans auf deutsch, wird sonst graece Johannes genannt, Claus wird graece Nicolaus genannt. Also müssen wir ihnen nachrechnen; die neuen Geister geben darauf nicht groß Achtung. Maria, des HErrn Christi Mutter, wird auch Mirjam genannt c. Diese Mirjam, Mosis Schwester, nimmt ein Exempel von ihrem Bruder Mose, und von den andern Männern, und singt dem HErrn auch ein Danklied mit Pauken und Reigen, zur Dankbarkeit, daß sie und die andern Weiber GOtt auch preisen und loben um seiner Wohlthat und Wunderwerke willen, die er an ihnen geübt hatte, wollen dadurch andere zu GOttes Erkenntniß und Dienst auch führen und locken.

\begin{schriftzitat}
V. 22–25. Moses ließ die Kinder Israel ziehen vom Schilfmeer hinaus zu der Wüste Sur; und sie wanderten drei Tage in der Wüste, daß sie kein Wasser fanden. Da kamen sie gen Mara; aber
\end{schriftzitat}

\quellseite{471}
Sie konnten des Wassers zu Mara nicht trinken, denn es war fast bitter; daher hieß man den Ort Mara. Da murrete das Volk wider Mose und sprach: Was sollen wir trinken? Er schrie zu dem HErrn, und der HErr wies ihm einen Baum, den that er in das Wasser, da ward es süß.

56. Bisher haben wir gehört das Geheimniß großer, herrlicher Thaten, an den Kindern von Israel bewiesen, da sie aus Egypten geführt worden, auch wie sie GOtt gelobt und gedankt haben für diese großen Wunderwerke. Nun wollen wir wieder zu der Historie kommen, und hören, wie es ihnen weiter gegangen sei; und ist dieser Text das Ende des fünfzehnten Capitels. GOtt hat sie durch ein groß Mirakel aus Egypten geführt; darum so bleibt er noch bei den Wundern, und hört nicht auf Mirakel zu thun, die ganze Reise über, bei vierzig Jahren, dieweil sie in der Wüste sind, bis daß sie in das gelobte Land kommen, und thut diese lange Zeit ein Mirakel und Wunder nach dem andern.

57. Aber allhier ist uns auch ein Exempel vorgeschrieben und abgemalt derjenigen, so GOttes Wunderwerke sehen, groß halten, loben und rühmen, und dennoch bald wieder zurückfallen und aller Wohlthaten GOttes vergessen, in dem, daß sie wider GOtt murreten und ungeduldig wurden, daß sie in dreien Tagen nicht Wasser zu trinken hatten. Kurz zuvor hatten sie gesehen und erfahren das unaussprechliche Mirakel, wie sie aus Egypten geführt, und durch das rothe Meer gebracht waren. Das sollte in ihren Augen billig ein groß Ding sein, denn das rothe Meer ist nicht so schmal gewesen, als die Elbe oder der Rhein ist, sondern acht oder zehn Meilen Weges breit. Daraus sind sie wunderbarlich errettet, und der Feind Pharao darinnen mit alle seinem Kriegsvolke versenkt und ertränkt.

58. Was geschieht aber? Ueber drei Tage vergessen sie alles, verachten die vorige göttliche Hülfe und Beistand, schlagen es in Wind, haben sich gegen GOtt gestellt, gleich als sei er nicht mehr bei ihnen, und sie keinen sein nimmer. Ist das nicht Sünde und Schande? Da sie kommen zur Wüste hinaus, auf eine Haide und Gehölze, da keine Leute wohnen, noch Länder und Städte sind, da nicht ein flach Feld ist, und gleichwohl diese Wüste stracks zum gelobten Lande zugeht, daselbst finden sie kein Wasser. Wiewohl, als sie kamen gen Mara, daselbst Wasser fanden, konnten aber desselbigen Wassers nicht trinken, denn es war fast bitter; daher hieß man den Ort Mara, das ist, bitter oder Bitterkeit; da murrete das Volk wider Mosen und sprachen: „Was sollen wir trinken?“

59. Hier können sie nimmer halten, die Geduld wird zu kurz. Sie haben irgend einen Tag Mangel an Wasser gehabt; denn sonst werden sie vom rothen Meer Wasser mit sich getragen haben; aber da sie einen oder zween Tage Durst leiden, und nicht bald Wasser haben, da murren sie, sind nicht länger denn drei Tage fromm. Es hat aber damit GOtt ihren Glauben versucht und probirt, daß er eine neue Anfechtung, Noth und Trübsal unter sie schickt, als den Durst. Da gedenken sie nicht: Siehe, hat uns GOtt aus dem Tode und aus dem rothen Meer heraus geführt, so wird er uns ja auch zu trinken geben, daß sie also ihm geglaubt und vertraut, daß er diesem Durst auch steuern würde. Darnach, so sahen sie auch die Wolkensäule und das Feuer Tag und Nacht vor ihnen hergehen, dem sie nachfolgten, als einem Panier, welches ein öffentlich Zeichen gewesen ist, daß GOtt als ein Herzog und oberster Feldherr vor ihnen sichtiglichen herzog, am Tage und auch des Nachts. Welches ist ein herrlich Exempel der Gegenwärtigkeit GOttes bei seinem Volk. Dennoch so bewegt es sie nichts.

60. Sie fragen nichts nach dem vorigen Mirakel und Erlösungen, welche sie doch sollten gestärkt und getröstet haben, daß der GOtt, so ihnen zuvor geholfen hatte, noch lebte und noch aushelfen würde, wie 1 Sam. 17, 37. David auch sagt: „Der GOtt, der mir von dem Löwen und Bären half, wird mich auch von dem unbeschnittenen Philister erretten.“ Und 2 Cor. 1, 8–10. sagt St. Paulus, GOtt habe ihm aufgeholfen in Asia aus seiner Leibesschwachheit und Todesnoth durch ihre Fürbitte, und werde ihm ferner helfen, und spricht also: „Wir wollen euch nicht verhalten, lieben Brüder, unsere Trübsal, die\footnote{Eislebensche: unser Trübsal, der.} uns in Asia widerfahren ist, da wir über die Maßen beschwert waren, und über Macht, also, daß wir uns des Lebens erwogen, und beschlossen hatten, wir müßten sterben. Das geschah aber darum, daß wir unser Vertrauen nicht auf uns selbst stellten, sondern auf GOtt, der die Todten auferwecket, welcher
\quellseite{472}
ung von solchem Tode erlöset hat, und noch täglich erlöset, und hoffen, er werde uns auch hinfort erlösen, durch Hülfe eurer Fürbitte für uns.“

61. So sehen nun die Kinder Israel nur auf den Bauch, wie der gorret\footnote{gorren = murren.} und zu trinken haben will. Freßlinge sind es, die mehr auf den Bauch und Mangel des Wassers Achtung geben, denn daß sie auf GOttes Mirakel, Wort und Verheißung sehen. Fallen also herunter und sorgen nur für den Leib. Gleich als wenn sie nicht Wasser haben, so haben sie auch keinen GOtt mehr. Sie wollen den Bauch versorgt haben, es bleibe GOtt im Himmel, oder sei, wo er wolle. Das mögen wir uns lassen eine Warnung sein, und unsere Augen nicht in die gegenwärtige Noth und Gefahr stecken, sondern sehen auf GOttes Wort, auf seine vorigen Mirakel und Exempel, wie er andern geholfen habe; denn dadurch wird der Glaube durch den Heiligen Geist gestärkt.

62. Nun, GOtt schweigt dazu etwas stille, und läßt sie sein wie sie sind. Und merke allhier, daß sie nicht alle zugleich gemurrt haben; wie denn hernacher im Texte folgt. Denn etliche unter ihnen sind noch gestanden, und [haben] auf GOtt, der unter ihnen Wunderzeichen thue, gewartet, und der vorigen Mirakel auch recht gebraucht, und gehofft, daß GOtt die Seinen nicht verlassen werde. Denn wenn GOtt Wunderzeichen thut, schickt er zuvor etliche, oder Einen aufs wenigste, der den Glauben hat. Denn er thut nicht Mirakel, es sei denn zuvor der Glaube da, um deswillen er Wunder übt. Darum sind allhier noch etliche gewesen, die geglaubt haben, als Moses, Aaron und andere. Denn Moses schrieet zum HErrn um Wasser; wie denn hernacher auch noch zwei fromme Männer gefunden worden, als, der Josua und Caleb, die in das gelobte Land kamen, da die andern alle sterben mußten. Also sind allhier auch noch etliche, wenige Fromme, um derer willen er Wasser gibt, und die bösen Schälke der Frommen Wohlthat genießen; und die Frommen müssen es allhier entgelten, darum, daß die bösen Schälke und Buben gemurrt haben, dieweil sie unter ihnen sind, gleichwie sonst die Bösen oft der Frommen genießen.

63. GOtt hätte eine Bornquelle machen können, aber er will es nicht thun. Denn er handelt wunderbarlich; gleichwie hier, daß im Felsen unnatürlicher Weise Wasser steht, und es noch täglich also geht, daß alle Wasserquellen aus den Bergen und Steinen gehen. Zudem so ist es nicht weniger ein Wunder, daß die Wolken Wasser geben. Wo nimmt er es doch? Es kann nichts Dünneres und Weicheres sein denn eine Wolke, denn es ist nichts anders denn eine Luft; wie kann nun aus einer Luft Wasser und Regen werden; oder, wie kann er aus Wolken Steine machen, als, wenn es graupet und hagelt? Ist sich deß nicht auch zu verwundern, daß in einer halben Stunde der ganze Himmel mit Regen überzogen wird? Aber niemand denkt ihm nach, es wird alles verachtet, weil es gemein ist.

64. Es geschieht aber alles dahin und ergeht darum also, daß GOtt in uns den Glauben gerne erweckt und baut, auf daß wir ihm vertraueten. Denn er will nicht lassen mangeln. Und wenn wir auch gleich in steinernen Klüften säßen, dennoch könnte er daraus Wasser und alle Nothdurft machen. Item, was bitter ist, das kann er wohl süß machen, und wiederum, bald süß auch sauer und bitter machen. Also, was weich ist, das kann er auch hart machen, die Creaturen müssen ihm allein gehorsam sein [Ps. 135, 6.].

65. Also ist allhier von Natur das Wasser bitter gewesen, und also geschaffen; aber, da sie es jetzt trinken sollen, heißt der HErr einen Baum oder Holz hinein werfen, da wird es süß. Nicht, daß dies Holz solche große Kraft hätte; sondern es ist ein Mirakel gewesen, daß Gott durch sein Wort, ohne einiges Zuthun des Mosis, hat thun wollen, und war bald das Wasser nicht bitter, wie zuvor; ja, er kann es lassen bitter bleiben, und lassen süß schmecken etc.

66. Das ist alles darum geschrieben und geschehen, daß wir den Glauben, seine Anfechtung und Uebung wohl lernen, und wissen, daß wir einen solchen GOtt haben, der Alles aus Nichts macht [Röm. 4, 17.]. Ihm ist nichts zu klein, auch keines zu groß; und, die da glauben, sollen alles haben.

\begin{schriftzitat}
V. 26. Daselbst stellte er ihnen ein Gesetz, und ein Recht, und versuchte sie und sprach: Wirst du der Stimme des HErrn deines GOttes gehorchen, und thun, was recht ist vor ihm, und zu
\end{schriftzitat}

\quellseite{473}
Ohren fassen seine Gebote, und halten alle seine Gesetze, so will ich der Krankheiten keine auf dich legen, die ich auf Egypten gelegt habe; denn ich bin der HErr, dein Arzt.

67. GOtt hat ihnen oft Gesetze gegeben, als, droben im 12. Capitel, da sie das Osterlämmlein essen sollten; item, wie sie abborgen sollten gülden und silbern Geräthe von den Egyptern. Also gibt er nun allhier durch den Menschen Mosen, und nicht durch sich, Gesetze, Ordnungen und Rechte; wie er denn hernacher, auf dem Berge Sinai, durch sich selbst wiederum die zehn Gebote durch große Mirakel geben wird.

68. Zu der Stunde, als das bittere Wasser durch den Baum süß gemacht war, da ließ er eine Predigt gehen auf das Wunderwerk, und hat Mosis dies Mirakel zur Ursache genommen, eine starke, gute Predigt darauf zu thun; als sollte er sagen: Ihr habt gesehen GOttes Wunderwerk wie denn auch zuvor; derhalben so sage ich euch, lasset euer Murren aufstehen, seid nicht so ungeduldig nach diesem Mirakel, so ihr neulich gesehen habt; gebet euch drein, gehorchet der Stimme des HErrn eures GOttes, und thut, was recht ist. Das Land ist dir verheißen, er wird dich nicht lassen. Wirst du folgen seiner Stimme, so will er dein Arzt sein, dein Helfer, es soll dir wohlgehen, er wird dich noch wohl erhalten; wie du denn jetzt erfahren hast an dem Mirakel mit dem bittern Wasser, das da ist süß worden.

69. Dieses ist die Ursache und Meinung der Predigt gewesen, und anderer Reden Mosis mehr zu diesem Volke. Denn er vermahnt sie, und richtet sie auf zum Glauben, daß sie Gott gehorsamten und ließen ihn recht haben, er würde sie nicht lassen, ob sie gleich ein wenig etwas leiden müßten. Denn GOtt versuchte sie, und ließ sie schwach und krank werden, darum, daß er wollte bei ihnen sein, und sie wieder heilen: Er will ihr Helfer und Arzt sein. Das ist dies Theil, da sie drei Tage in der Wüste gezogen sind.

\subsection{Allegorie dieser Geschichte}

70. Nun folgt das Mysterium, oder die Allegorie über diesen Text; denn viel haben ihn wunderbarlich gedeutet. Aber, was du willst geistlich deuten, das sollst du immerdar ziehen auf den Glauben an Christum und auf das Evangelium. Eines Theils haben es also ausgelegt, daß Mara, das bittere Wasser, bedeute allerlei Kreuz, Leiden und Unglück, so einem Menschen widerfahren mag. Aber daß Moses ein Holz drein wirft und es süß macht, sei, daß man das Sacrament des Abendmahls, den Leib und Blut Christi drein werfe, und man gedenke an das Leiden und Sterben; dann schmecke alles Kreuz und Leiden wohl; dann sei kein Leiden so herbe und bitter, daß nicht süß und lieblich werde. Dieses lassen wir christliche und gute Gedanken sein; wiewohl sie kindisch genug sind, dennoch so werden sie auf das Leiden Christi gezogen; aber recht deuten geht hinein in Geist, auf das Gewissen und Glauben.

71. „Mara“ ist das Gesetz GOttes, und „wandeln in der Wüste“ ist die Schwermuth des Lebens, so geführt wird unter der Last des Gesetzes und der guten Werke, da nimmermehr keine Ruhe noch Friede ist, sondern der Geist wird gar getödtet, und wird durch die Gesetze alles je länger je ärger. Denn das Gesetz macht böse Gewissen, und richtet nur Zorn an, und mehrt die Sünde, Röm. 4, 15. Derhalben, je mehr für das Gewissen das Gesetz gepredigt wird, je mehr der Mensch erschreckt werden muß durch das bittere Wasser; das macht alle betrübt, es richtet zu schwermüthige Leute; wie denn die erfahren, die unter dem Gesetze sind.

72. Nun, da erhebt sich ein Murren und Ungeduld wider GOtt, welches denn natürlich geschieht [Matth. 20, 11.]. Denn da kann das Herz anders nicht thun, denn daß es GOtt muß feind werden. Ich rede jetzt geistlich vom Gesetze. Denn äußerliche, grobe Menschen muß man mit dem äußerlichen Gesetze zwingen [1 Tim. 1, 9.], aber die zarten Seelen (da rede ich jetzt von), die zu dem Evangelio gehören, diese führt GOtt ein drei Tagereisen und zermartert sie, und bringt sie in den großen Durst, Angst und Noth, zum bittern Wasser, das sie trinken sollen, das ist, daß sie ein bitter Gewissen haben. Das Murren, welches folgt, ist des Gewissens, das heimlich wider GOtt zürnt; mit dem Munde spricht es, es habe GOtt lieb, aber das Herz leugt. Denn es will nicht haben, daß es so hart versucht werde, und ist allen Strafen feind, so über die Sünde gehen.

73. Die Strafen machen ein bös Gewissen, thun demselben wehe und sind dem alten Adam ganz verdrießlich; aber wenn nicht Strafen wären, o so wäre ein sündlich Leben gar ein
\quellseite{474}
sein Ding. Dieweil sie nun dem Gesetze und Zorn GOttes feind sind, so sind sie auch GOtt feind, und solches aus nachfolgender Ursache, daß unwidersprechlich folgt: Wer die Strafen nicht billiget, der will GOttes Gerechtigkeit nicht haben, und erdichtet ihm selbst allein einen solchen GOtt, der barmherzig sei (wie jetzt die Welt pflegt zu thun, und weiß sich meisterlich darauf zu behelfen), darum so will derselbigen auch GOtt selbst nicht haben. Derhalben so ist inwendig im Gewissen ein heimlicher Groll, Lästerung und Murren wider GOtt, da das Herz gerne wollte, daß gar kein GOtt noch Strafe wäre. Daher muß das Gewissen immerdar zittern und beben, auch vor einem rauschenden Blatt sich fürchten und entsetzen.

74. Dies ist nun eine heimliche Sünde, daß man solches zudeckt, davon nicht viel Leute wissen [Jer. 17, 9.]. Sollten sie GOtt lieben und ihm günstig sein? Sie wollten gerne GOtt fliehen und meiden, und wollten, daß kein GOtt wäre; wie man denn an denen sieht, die da sterben sollen, oder etwas um ihrer Mißthat willen erdulden. Und am jüngsten Tage werden sie es viel kenntlicher beweisen, da GOtt von sich treiben wird die Herzen, mit Sünde der Ungeduld und Zorns beladen, eben darum, daß sie ungeduldig gewesen, und wider GOtt gemurrt; welches denn GOtt sehr wohl hört, so sieht es auch ein jeglicher, der geistlich urtheilen und richten kann; aber die es thun, die erkennen es nicht also. Es thun's aber alle, die erschreckt werden durch das Erkenntniß der Hölle und des Gesetzes, ehe denn sie den Heiligen Geist bekommen und ihre Herzen gemildert werden. Darum, daß wir dem Gesetze günstig und hold werden, solches macht der Baum, welcher von GOtt dem Mosi gezeigt wird; was es auch muß für ein Baum gewesen sein, denn hier wird er mit Namen nicht ausgedrückt.

75. Es wird aber zweierlei allhier angezeigt: erstlich, daß das Wasser, das ist, das Gesetz, nicht süß gemacht wird ohne Zuthun Mosis, welcher den Menschen durch Gesetzschrecken gar und mürbe macht, und mit Bitterkeit also ängstiget, daß er nach Hülfe verlangend wird; alsdann, so der Heilige Geist kommt, bald wird es süß. Aber unsere Rottengeister wollen den Heiligen Geist haben ohne das göttliche Wort. So ist nun dieser Baum das liebe Evangelium; das bittere Wasser ist das Gesetz, oder die Erkenntniß der Sünde. Der Baum des Lebens [1 Mos. 2, 9.] ist das liebe Evangelium, das Wort von GOttes Gnade, Barmherzigkeit und Güte; wenn das Evangelium in das Gesetz und Erkenntniß der Sünde getaucht wird und rührt das Herz an, darinnen das Gesetz Traurigkeit, Angst, Schrecken und Betrübniß anrichtet, da schwemmt es. Mosis thut seinen äußerlichen Dienst; also hier auch. Es thut nichts ohne das äußerliche Evangelium, denn durch dasselbige folgt Süßigkeit und Lust zum Gesetze. Man erkennt daraus, daß Christus für uns gestorben sei, und empfähet den Heiligen Geist, und kriegt Lust und Liebe zu GOtt, dem man sonst zuvor feind war. Wenn denn GOtt ihm gefällt, und diese Lust ist angegangen, dann ist Friede, und geschieht, was GOtt heißt und gebeut, und auch das Gesetz haben will. Es ist ihm lieblich und lustig, und hört gerne davon, dieweil das Gesetz ist von einem andern erfüllt; er ist gar ein anderer Mensch worden.

76. Also sollte man es auf das Predigtamt ziehen; und darum treibe ich es, daß ich die jungen Theologen, so in der heiligen Schrift studiren, übe, und frisch leite und treibe zum Wort GOttes, und den Glauben auf GOttes Wort gründe. Die andere, vorige Auslegung [§ 70] ist zu schwach und zu kindisch, erlangt den Grund und Geist nicht.

77. Wenn denn nun die Bitterkeit weggenommen wird, und das Herz gereinigt ist, so hat es nicht Betrübniß. Dann wird erkannt, daß GOtt ihnen das Gesetz gegeben habe, daß sie der Stimme des HErrn gehorchen. So wollte er auch keine Krankheit auf sie legen, wie er auf Egypten gelegt hatte, sondern er wolle ihr Arzt sein, da sie gleich mit Krankheit überfallen würden; das ist, wenn das Evangelium gehet wird, das denn auch den Tod süß, ja, zu einem sanften Schlafe macht. Dann ist der Grund gelegt, und das ist die rechte Lehre, da muß man anhalten, treiben und vermahnen, daß der Mensch sich zu Christo kehre, und lerne, wie er uns helfe von Sünden, Gesetz, Tod und Teufel. Dabei bleibe, und lerne GOtt vertrauen. Wenn du nun ferner versucht wirst mit Bitterkeit, so bleibe allezeit bei dem Worte, so wird GOtt dein Medicus sein, und wird dich schützen vor aller Betrübniß.

78. Denn in diesen zweien Stücken steht das Predigtamt: erstlich lehren, daß man nicht weiß;
\quellseite{475}
darnach, vermahnen und anhalten, daß man wohl einbilde, was man gelernt und erkannt hat, und nicht laß werde oder dasjenige vergesse [2 Tim. 4, 2.]. Wir sind fleisch und Blut, darum so kann es nicht genug gepredigt werden. Solches sieht man auch in der heiligen Schrift, die predigt immerdar mit andern Worten von einerlei Lehre. GOtt weiß, daß der alte Schalk, der alte Adam, faul und träge ist; er läßt ihm bald die reine und rechtschaffene Lehre nehmen, darnach sich irriger, verführerischer Menschenfündlein bereden.

\section{Das sechzehnte Capitel}

Cap. 15, 27. Und sie kamen in Elim, da waren zwölf Wasserbrunnen, und siebenzig Palmbäume, und lagerten sich daselbst an das Wasser etc.

Cap. 16, 2. 3. Und es murrete die ganze Gemeine der Kinder Israel wider Mosen und Aaron, in der Wüste, und sprachen zu ihnen: Wollte GOtt, wir wären in Egyptenland gestorben, durch des HErrn Hand, da wir bei den Fleischtöpfen saßen, und hatten die Fülle Brod zu essen. Denn ihr habt uns darum ausgeführt in diese Wüste, daß ihr diese ganze Gemeine Hungers sterben lasset.

1. Wir haben den Text des 16. Capitels im andern Buch Mosis gehört, wie die Kinder von Israel von dem rothen Meer zu der Wüste Sin gezogen sind, und wie sie erstlich in Elim gekommen, aber nicht lange allda geblieben, da sie denn gefunden haben zwölf Wasserbrunnen und siebenzig Palmbäume. Von demselben Orte sind sie kommen in die Wüste Sin, da denn diese große That geschehen ist, so hernach folgt. Darnach sind sie gekommen in Raphidim, da hatte das Volk abermals nichts zu trinken; von dannen sind sie auf den Berg Sinai gekommen.

2. Nun wird alhier gemeldet, daß sie zwölf Brunnen gefunden haben, daß sie des Wassers, so GOtt durch ein Mirakel ihnen gegeben hatte, mehr nicht bedurften; darnach haben sie auch siebenzig Palmbäume gefunden. Und kann wohl sein, daß sich durch sonderliche Gottes Schickung dieses Brunnenwasser also in zwölf unterschiedliche Quellen getheilt hat, damit das Volk Israel Wassers genug haben möchte; wie GOtt zuvor auch gethan hatte.

3. Also versorgt GOtt sie mit Speise und Trank, daß sie leibliche Nahrung haben müssen, im Exilio und in der Wüste zu trinken und zu essen finden; er bringt sie allhier gar an einen lustigen Ort, da sie sich wieder erquicken und erlustieren mögen. Palmen haben eine süße Frucht, die Datteln; davon haben sie gezehret, gegessen und des Brunnenwassers getrunken; vielleicht werden sie auch vom rohen Zeige noch etwas übrig behalten haben. Dieweil sie nur fünfzehn Tage-reisen aus Egypten zur Wüste Sin gezogen sind, so wird ohne Zweifel jedermann für sein Haus einen Scheffel oder zween Mehl zu Brode mit sich auf den Weg genommen, und etlichermaßen auf Vorrath gedacht haben. Daraus haben auch wir zu lernen kürzlich, daß GOtt unser Hirte und Pfleger sei, der uns weide, speise und genug gebe; wie der 23. Psalm, V. 1., solches auch lehrt.

4. Aber wenn das Mehl aus dem Sacke ist und kein Teig mehr vorhanden, dann regt sich der Abgott, Junker Bauch, und läßt sich merken, er könne unseres HErrn GOttes Verzug nicht auswarten. Denn dies Volk murret wider GOtt, daß sie nicht Brod zu essen haben; gleichwie sie droben murrten, da Wasser zu trinken mangelte. Es sind in der Wahrheit die Kinder von Israel gar ein böse Volk gewesen, die GOttes Mirakel und Gutthat bald vergessen und ihres Berufs überdrüssig werden, verzweifeln an GOttes Gnade und Barmherzigkeit. Mosen und Aaron lästern sie, als die da wollten das Volk Hungers sterben lassen; ja, GOtt selbst schänden [sie], als habe er sie durch Mosen nicht lassen aus Egypten führen. Darum urtheile, was für ein Kränklein dies Volk sei, wie sie zweifeln an GOttes Befehl, daß Moses sie aus Egypten in das Land Canaan führen sollte, item, verzagen an GOttes Hülfe.

5. Und wir sind gleich der Art; solch gottlos Wesen sieht uns auch aus den Augen. Wenn
\quellseite{476}
nimmer Vorrath vorhanden ist, so hört unser Glaube auch auf. Sie sind nicht mehr denn fünfzehn Tagereisen gegangen, da sagen sie: Moses hat es wohl ausgerichtet, wir können nicht zurück gehen, wir wollten denn Hungers sterben; auch können wir nicht weiter gehen, noch zur Seite, von wegen der Einöde und Wildniß,\footnote{Eißlebensche: Gewildniß.} ausbrechen; der Teufel hat uns in die Wüste geführt.

6. Nach der Vernunft haben sie Ursache genug wider Mosen, dieweil sie nichts zu essen oder zu trinken haben. Denket ihr selbst, wenn du fünfzehn Tage und Nächte in der Wüste wärest und könntest nicht heraus, da gedächtest du selbst auch, daß du müßtest Hungers sterben. Diese Ursache haben sie gehabt, daß das ganze Volk murrt wider Mosen, und spricht: Gib uns Speise, warum hast du uns lassen aus Egypten ziehen, daß du uns, unsere Kinder und Vieh Hungers sterben ließest?

7. Welche spöttische, höhnische und lästerische Worte sind doch das? Es ist ein tiefer Unglaube, daß sie also zurück fallen, GOttes Wort und Verheißung fahren lassen, an die vorigen Wunderwerke und Hülfe GOttes nicht mehr gedenken; und wollen die Juden dennoch fromme Leute sein, dazu von GOtt ungestraft sein, und wollen doch gleichwohl, daß alle Teufel Aaron und Mosen wegführten. Denn sie sagen allhier: Der HErr weiß nichts drum. Sie sprechen: Dort, in Egypten, wären wir selige Leute gewesen, aber allhier in der Wüste müssen wir Hungers sterben. Hundert Teufel haben uns hieher gebracht. Habt ihr uns darum ausgeführt, daß das ganze Volk Hungers sterbe?

8. Und nachdem sie alle mit einander also gemurrt und geredet haben wider Mosen und Aaron, sollten den beiden billig die Hosen gestunken haben; aber sie stehen wie die Helden, ihr Herz ist voller GOttes, schlagen diese Lästerung aus, so sie müssen hören, nämlich, daß sie sagen: Ihr seid Gesellen, die ihr uns ausgeführt habt, und in dem gehandelt als die Bösewichter; sollte GOtt solches euch geheißen haben? Da sind die Wunder alle hinweg, die sie zuvor gesehen und empfunden, ja, gerühmt und gelobt haben im vorigen Lobgesange, da sie gänzlich überwunden und überwiesen gewesen sind, daß diese Ausführung aus Egypten kein Menschenwerk oder Betrug sei. Wohlan, es geht nicht anders zu; wo der Abgott Bauch kommt, und der Magen beginnt zu platzen, da ist es alles weg. Eine edle, zarte Frucht ist der Mensch, wenn er ihm selbst gelassen wird.

9. Aber wir wollen es besser ausstreichen, auf daß wir sehen mögen, wer wir doch sind gegen GOtt. Denn wir verfluchen, lästern, ja, verfolgen auch seine Boten und Knechte, und machen seine Wunder zu Schanden. Allhier fällt eine ganze Gemeinde dahin, bei sechsmal hunderttausend Mann, viel trefflicher, großer Leute. Josua, Caleb und andere gehen jetzt alle dahin mit dem Haufen und Obersten des Volks, und bleiben\footnote{Eißlebensche: blieben.} allein diese zween beständig, die halten noch Farbe und treten nicht zurück. Was wollen wir uns doch verwundern, daß jetzt etliche dahin vom Evangelio taumeln und purzeln, dieweil so viele Leute allhier dahinfallen, welche der Teufel alle frißt? Und das richtet allein die Bauchsorge aus, und bringt sie zu diesem greulichen Abfall [Matth. 6, 25.].

10. Es ist ein wunderbärlich Ding um einen Christen, der da steht, wenn es noth thut. Dennoch pflegt man solche Leute in der Welt nicht hoch zu achten und drängt sich nicht sehr um sie. Aber unsere Natur ist also geartet, wie Salomon [Sprüchw. 20, 14.] davon sagt: „Wenn man ein Ding hat, so ist es böse, und wird nicht geachtet und taugt nirgend zu, aber wenn es hinweg ist, so preiset man es denn.“ Also geht es auch: Wenn GOtt genug gibt, so haben wir dennoch nicht satt und ist noch nirgend gut genug, gaffen immerdar vor sich auf den Mangel, und man verläßt die Fülle und Bescherung GOttes, so vorhanden ist, bis man drum kommt. O! dann ist’s sehr gut gewesen.

11. Ich will es noch besser auslegen, auf daß es vernommen werde. Nimm vor dich einen Ehemann, und einen, der außerhalb der Ehe lebt; oder gleich einen armen und reichen Mann, von welchen das Sprüchwort lautet: Gut macht Muth, Armuth wehe thut. Wer Armuth leidet im Ehestand, der sieht nicht, was es will mit ihm werden, und duldet die Armuth mit Unwillen; alle seine Augen, Gedanken und Sinne sehen auf den Mangel; daß er aber nicht stirbt im Hunger, und einen sichern, fröhlichen Muth dazu hat, auch ohne Sorgen lebt, das gibt ihm
\quellseite{477}
Gott, und ist des Armen sonderlicher Vortheil, und Zuvoraus.\footnote{Zuvoraus = Vorzug.} Er bekümmert sich nicht, dieweil er weiß, man stiehlt ihm nichts, er kann mit Frieden schlafen; dagegen, daß er das Haus voll Kinder hat, das ist in seinen Augen nichts, der Güter rechnet er keines nicht [Ps. 128, 3. ff. Ps. 127, 3. 4.]. Wenn man dieses auf einer Wage wägen würde, so würdest du sehen, wie weit solche Güter alle Armuth übertreffen. [Er] geht also dahin und braucht des Himmels und der Erde.

12. Aber dagegen ein Reicher, wenn er gleich viel Güter hat, so muß er auch viel Mäuler ernähren. Da muß er sorgen, daß ihm das Gut und Reichthum nicht gestohlen werde, und denkt immerdar, wie er am Gut wachse und zunehme; er hat nicht genug, sondern gedenkt immerdar: Ei, hätte ich noch das dazu. Also hat er nimmermehr einen sichern und fröhlichen Muth, das Gut macht ihm Unmuth; zudem weder Ruhe noch Frieden. Wiederum, der Arme sieht nur des Reichen Gut an; aber er betrachtet nicht, daß er einen Vortheil habe vor ihm, wie kurz zuvor angezeigt, der dem Reichen gänzlich abgeht. Ursache, es heißt, wie Augustinus sagt: Divitias invenisti, requiem perdidisti [Reichthum hast du gefunden, die Ruhe hast du verloren.].

13. Also fehlt jetzt den evangelischen Bauern auch nichts, allein, daß sie aufstehen, wie sie Zins und Geschos, auch den Zehnten und andere Gerechtigkeit der Obrigkeit geben sollen, die daran nicht Unrecht thut, daß sie ihren Tribut fordert [Röm. 13, 1. ff.]; da wenden sie allein die Augen her. Aber dagegen sehen sie nicht, daß ihre Aecker und Wiesen, Haus und Hof im guten Frieden ist; diese Sicherheit macht ihnen die Obrigkeit, sonst könnten sie nicht eine Stunde lang sicher schlafen in ihren Häusern. Diese Sicherheit und solch groß Gut sieht man nicht. Ein Trunk Brunnens schmeckt einem solchen besser, denn Gesottenes und Gebratenes\footnote{Eißlebensche: gesotten und gebraten.}: noch gaffet und sieht er sich um nach den großen Hansen, die auf schönen Hengsten daher traben und in güldenen Ketten prangen. Ach! gedenkt dieser Narr, wollte GOtt, daß ich an seiner Stätte allda sitzen sollte! und muß ihm alles stinken, was in seinem Hause ist; sieht aber und weiß nicht, was für Beschwerung, Gefahr und Unlust diesem Stande anklebt und nachfolgt.

14. Also thut auch einer, der außerhalb dem Ehestande lebt; der hält sein Wesen für eitel Ungemach, und läßt sich dünken, der Ehestand sei nur ein Rosengarten, und sieht nicht den Vortheil, daß er nur den Bauch, oder Ein Maul habe zu versorgen. Widerum, der im ehelichen Stande ist, der sieht auch nicht die Güter, die ihm gegeben sind, nämlich, daß der Mann und das Weib von GOtt geschaffen und zusammen verordnet sind; item, daß sie Kinder haben und Güter ihnen beschert werden; dieser läßt ihm träumen, daß der ledige Stand ein Paradies sei.

15. Es sind die Güter und Gaben GOttes in einem jeglichen Stande zehnmal größer, denn der Schade und Mangel, so einer darinne fühlt. Wenn du des Nächsten Güter rechnetest mit deinem Mangel, und dagegen deinen Mangel mit deinen Gütern, so würdest du auch thun dasjenige, davon die Weisen geschrieben haben: Wenn es könnte sein, daß einer seine Güter und sein Unglück auf einen Haufen trüge, und austheilte zugleich, so würde ein jeglicher sagen: Lieber, gib mir deine Güter; und des Unglücks daneben nicht haben wollen.

16. Also hat der heidnische Poet Horatius gesagt, eben dies anzuzeigen, es lasse ihm niemand an seinem Stande genügen: Optat ephippia bos piger, optat arare caballus, der Ochse wollte gerne ein Pferd sein, und das Pferd wollte gerne pflügen wie der Ochs \&c., so es doch niemand besser machen noch anders ordnen sollte, denn wie es GOtt gemacht hat.

17. Also wird es den Bauern auch gehen. Wenn es nicht wider GOtt wäre, so wollte ich aus einem Korn den aufrührischen Bauern wünschen, daß sie nur Ein Jahr sollten Fürsten sein, Briefe lesen, Klage hören, Händel verrichten und entscheiden, Briefe schreiben, und anderes mehr thun, was in diesem Stande will ausgerichtet und beschickt sein; da würde einer desselbigen so bald müde werden, daß er spräche: Behalte dir deine geelen Sporen, silbern Schwerdt, güldene Ketten, sammeten Rock, Spieß und Harnisch, schöne Rosse und behangene Wagen; ich will bleiben, der ich vorhin gewesen bin.

18. Also sind wir alle mit einander gesinnt, und unsere Brüder, die Kinder von Israel in der Wüste, gedachten nicht lange an das Manna oder Himmelbrod. Gleichwie sie in Egypten auch nicht gedachten an das Fleisch und Brod, so sie überflüssig hatten, sondern sie heulen und
\quellseite{478}
schreien allein darüber, daß sie mit den Frohndiensten und andern beschwert und bedrängt waren. Sie haben nicht betrachtet ihr Gutes, so ihnen GOtt gegeben hatte, sondern allein das Unglück und das Böse gesehen. Aber an diesem Orte, da sie aus dem Mangel gekommen sind, und das Gute erlangt, so laufen sie zurück. Also ist unser Auge ein Schalk, und kann ihm GOtt selbst nicht recht thun.

19. Ei, wollen wir denn gar nichts leiden? soll uns denn allhier nichts gebrechen? Im Himmel wird uns nichts mangeln noch gebrechen; in diesem Leben aber müssen wir (wir wollen oder wollen nicht) Gebrechen und Mangel mit unter haben. Bist du arm, hast du kein Haus oder andere Güter, so hast du auch Gebrechen: warum gedenkst du aber nicht, du habest einen gesunden Leib, gesunde Augen und andere Sinne, du hast deine Stärke, Kinder und anderes? \&c. Dagegen ist dein Gebrechen gar geringe und klein, den du daneben hast.

20. Aber man muß anheben zu lästern und zu schänden um eines kleinen und geringen Gebrechens willen. Gleichwie allhier die lieben Söhnlein, die Kinder von Israel, auch thun; da sie satt sind und mit Fleisch und Brod erfüllt, sahen sie zu und gedachten das Himmelsbrod, darum sie doch zuvor sehr zu GOtt schrieen. Viele sagen, GOtt sei zornig, stäupe und schicke uns Krankheit zu; aber wir schweigen davon kein stille, daß wir es auch wohl verdienen. Er greife es an, wie er wolle, doch taugt es uns nirgendhin; sollte er denn nicht stäupen? Wenn wir GOtt wären, wir hielten so lange nicht an uns, daß wir nicht flugs mit Keulen drein schlügen; aber er ist viel zu geduldig. Da es gegenwärtig ist, was er uns gibt, so verachtet man es; wenn es hinweg kommt, so lobt man es, und will keinesweges erstlich bedenken, wie gut es sei. Weil einem Mann sein Weib noch lebt, so sieht er nicht, wozu sie ihm nütze sei; wenn sie aber stirbt, so gedenkt er: Ach GOtt, hätte ich mein Weib noch!

21. Also, wenn die Sonne Einen Tag außen bliebe und nicht schiene, so würden wir alle betrübt werden; aber weil sie täglich sich einstellt, sind wir der Güter GOttes gewöhnt, meinen nicht anders, denn gleich als müßte sie uns scheinen; weil sie nun immerdar scheint, o so danken wir unserm HErrn GOtt auch nicht dafür. Da die Kinder von Israel in Egypten gefangen und dienstbar waren, verlangte sie nach der Freiheit; und da sie nun die Freiheit haben, so verachten sie die Freiheit. Also geht es: Wie es GOtt mit uns macht und thut, so wollen wir es nicht ansehen; noch trägt er es. Daher St. Paulus in den Geschichten der Apostel [Cap. 13, 18.] spricht: GOtt habe ihre Weise und Sitten in der Wüste geduldet und getragen; und das ist auch wahrlich also, er kann wohl leiden und tragen [2 Mos. 34, 6. Ps. 103, 8.].

22. Eine solche schändliche Unart ist an uns, daß wir für die gegenwärtigen Gaben und Güter undankbar sind, und allein einen kleinen Gebrechen ansehen. Und ein jeder mag heimgehen und rechnen die Güter, die er habe, so wird er viel mehr Güter denn Mangel finden, und danke GOtt dafür. Was willst du mehr haben, denn gesunden Leib? Item, deine Habe und Güter, und daß du schöne Kinder hast?

23. Aber wie sollen wir ihm helfen? Salomo sagt [Pred. 7, 15.]: In die malorum memor esto bonorum, et econtra: In die bonorum, memor sis malorum [am guten Tage sei guter Dinge, und den bösen Tag nimm auch für gut]. Geht dir es übel und fühlst das Böse, so kehre deine Augen zum Guten, und gedenke, wie gottesfürchtige Herzen thun, daß du die Augen abkehrst vom Unglück, und wendest sie über sich zu einem größern Glücke. Also sollen sie sich trösten, sonst machen sie ihnen ihr Leben sauer, und richten mit dem Grämen lauter nichts aus.

24. Von St. Martino, einem Bürger zu Rom, liest man, da ihm alle seine Güter genommen sind, und er in das Elend getrieben ward um des christlichen Glaubens willen, daß er gesagt habe:\footnote{Eislebenische: hatte.} Ob man mir gleich alles nimmt, so haben sie mir doch Christum nicht weggenommen. Geht hinweg mein Weib und Kind, mein Gut, Haus und Hof, meine guten Freunde, meine Ehre, Würde und Herrlichkeit, so habe ich doch den HErrn Christum noch. Also sieht ein Christ mit seinen Augen noch mehr auf das gegenwärtige Glück und Gut, denn auf das Unglück und Uebel, so einem begegnet.

25. Aber unsere Natur kann es nicht thun; sie thut, wie allhier diese elenden Leute thun. Dennoch stellt sich unser HErr GOtt als ein GOtt, der große Geduld habe, und es ihnen zugut könne halten, und spricht: Wohlan, ich will
\quellseite{479}
euch eure Lust büßen; wiewohl es unrecht ist, daß sich GOtt neigen sollte nach dem Willen dieses Volks. GOtt hat ihnen droben [Cap. 15, 26.] eine Verordnung gegeben, sie sollten seiner Stimme gehorchen, so wollte er ihr Arzt sein; darum gibt er ihnen Himmelsbrod, auf daß sie nicht sollten Hungers sterben; aber sie gedenken hier an die Verheißung nicht. Nun versucht er das Allerletzte, ob sie ihm wollten gehorsam sein, daß sie sechs Tage sollen Manna eintragen, aber am Sabbathtage sollten sie in ihren Hütten bleiben und nicht heraus gehen, darum so gab GOtt ihnen zweifach Manna, für den sechsten und auch für den siebenten Tag. Ein feiner Haushalter ist GOtt, er gibt das tägliche Brod, und gibt nicht mehr einem jeden, denn er essen kann.

26. Da heben sie wieder an zu glauben, und gedenken: Ei, ist dies nicht ein feiner Kornboden oder Kornkasten, daß alle Morgen die Luft und Wolken so viel Brods oder Manna gibt, als viel wir essen können? [Ps. 145, 15. 16.] Als der erste Abend kommt, da des Morgens das Manna gefallen, da war nichts mehr übrig; da werden ihrer viel gedacht haben, was werden wir den Morgen haben? Also haben sie gedacht: HErr, gib uns unser Brod heute, für das morgende sorgen wir nicht; wie er es denn auch damals beweiset hat.

27. Also versorget uns GOtt, auf daß er dem Bauche sein Wundern nehme. Denn der Bauch schlägt unserm HErrn GOtt immerdar sein Wort und Wunder nieder; aber man muß einmal, nach der Lehre St. Pauli, den Bauch und Speise unter die Erde graben. Darum lebe man also, und steure\footnote{steuern = stützen.} sich auf die väterliche Treue GOttes, daß man ungewiß sei, was man morgen eintragen möge. Und wenn auch dieser große HErr mit seinen Worten nicht soll gelten, o! so mag man dem Teufel glauben. Wie es die Erfahrung zeugt, daß ihrer, leider, allzuviel an GOtt treulos und dem leidigen Teufel, nur des garstigen Bauches halben, anhängig werden.

\begin{schriftzitat}
V. 4–8.
\end{schriftzitat}

\begin{quote}
Da sprach der HErr zu Mose: Siehe, ich will euch Brod vom Himmel regnen lassen, und das Volk soll hinaus gehen und sammeln täglich, was es bedarf, daß ich es versuche, ob es in meinem Gesetze wandle, oder nicht; des sechsten Tages aber sollen sie sich schicken, daß sie zwiefältig eintragen über das sie täglich sammeln. Moses und Aaron sprachen zu allen Kindern Israel: Am Abend sollt ihr inne werden, daß euch der HErr aus Egyptenland geführet hat, und des Morgens werdet ihr des HErrn Herrlichkeit sehen. Denn er hat euer Murren gehört. Was sind wir, daß ihr wider uns murret? Weiter sprach Mose: Der HErr wird euch am Abend Fleisch zu essen geben, und am Morgen Brod die Fülle, darum, daß der HErr euer Murren gehöret hat, daß ihr wider ihn gemurret habt. Denn was sind wir? Euer Murren ist nicht wider uns, sondern wider den HErrn.
\end{quote}

28. Hier antworten Moses und Aaron auf die Lästerung der Israeliten, da sie sprechen: Ach, daß wir in Egypten durch GOttes Hand gestorben wären; damit sie so viel wollen zu verstehen geben: Dort haben wir einen gnädigen GOtt gehabt, der Teufel hat euch zu uns geführt; und sagen: Nun, ihr lästert GOtt und murrt wider denselben; morgen sollt ihr ein Wunderwerk sehen, damit je GOtt zu preisen und herrlich zu halten sei. Wir wollen es euch hiermit angezeigt haben, ihr sollt es fühlen und greifen, auch erfahren, daß es GOtt thun werde, und dabei gewahr werden, daß nicht wir, sondern er euch aus Egypten geführt habe; denn sonst sollte morgen kein Wunder noch Mirakel kommen; GOtt wird aber seine Ehre damit retten und beweisen.

29. Mit dem tollen Pöbel soll man nach solcher Gnade und Barmherzigkeit handeln. Es ist wohl sticherlich,\footnote{sticherlich = in spitziger Weise geredet.} und sie hätten lieber etwas Besseres denn das, nämlich, wenn sie sagten und ihnen mit zu behuelften: Meinet ihr, daß wir die Leute sind, die euch konnten ausführen aus Egypten, und bitten, ihr wollet euch an uns nicht rächen? Aber sie haben sich so wohl und recht entschuldigt. Das hat Moses gesagt und Aaron befohlen, ihnen in das Angesicht zu reden, auch andern zu verkündigen, auf daß sie es unter das Volk brächten. Aaron nimmt es auch aus Moses Munde, daß er es unerschrocken vortrage, und sage ihnen eben die Meinung, wider wen sie gemurret und gemurrt haben.

\begin{schriftzitat}
V. 9.
\end{schriftzitat}

\begin{quote}
Und Mose sprach zu Aaron: Sage der ganzen Gemeinde der Kinder Israel: Kommet herbei vor den HErrn, denn er hat euer Murren gehöret.
\end{quote}

\quellseite{480}
30. Zum HErrn kommen ist, wenn man sich versammelt und GOttes Wort hört; denn wo sein Wort ist, da ist auch der HErr. Es ist damals noch nicht der Tempel gewesen. Diesen seinen Amt richtet Aaron redlich aus. Er straft sie erstlich, und bringt sie in Erkenntniß der Sünde, daß sie wider GOtt gemurrt haben, und seine vorigen Mirakel alle in den Wind geschlagen. Nun geht die Predigt an, und mitten in der Predigt wandten sie sich um zur Wüste; denn er hat diese Predigt nicht in Einem Tage ausgerichtet. Und dieses ist unter der Predigt des Aaron geschehen, daß GOtt sich sonderlich sehen läßt in einer Wolke, er scheint, daß man seiner gewahr wird in einer Wolke, aus welcher Mosis allein diese Stimme gehört hat, und nicht andere Leute. Das heißt, die Leute können tragen. Und wie er gepredigt hat, also geht es auch; aber er ließ das Wort vorhergehen, wie er denn pflegt zu thun, und darnach das Werk hernach folgen. Durch diese Predigt sind sie abermals bekehrt, und hat Moses und Aaron große Mühe und Arbeit mit ihnen gehabt, damit, daß sie ihren Glauben aufrichteten und stärkten. Darum so thut GOtt auch dies Mirakel.

31. Was man Wachteln heiße, darüber streiten noch unter sich die Lehrer: eines Theils meinen, es seien Rebhühner, andere, es seien Brachvögel; aber gewiß ist es, daß es Fleisch gewesen sei. Ueber ein Jahr hernach hat GOtt fast dergleichen Mirakel gethan, wie man liest 4 Mos. 11, 31. 34., da denn angezeigt wird, wie hoch sie geflogen sind in der Luft, nämlich, daß man sie mit den Händen hat greifen mögen, und man keines Netzes dazu bedurft habe. Daß aber derselbige Ort in der Wüste Lustgräber genannt ward, geschah darum, dieweil man daselbst begrub das lüsterne Volk. Nun darf man nicht gedenken, daß es einerlei Geschichte oder Historien seien, sondern es sind zwo Historien. Denn allhier haben sie auch Manna oder Himmelsbrod mit dem Fleische; dort haben sie allein Fleisch. Aber also kam diese Speise fein mit Haufen, daß alles voll lag; daß man ja nicht lange es durfte zusammenraffen und auffassen, so lag es alles fein bei einander.

32. Manna ist fast gewesen wie Coriander; wiewohl Coriander ist noch ein kleinerer Gestörn, gleichwie bei uns Reis oder Hirse; allein Manna ist weißer und runder gewesen, und schmeckte wie Honig mit Semmel vermengt und durchzogen, hat einen feinen, süßen, lieblichen Geschmack gehabt. Dasselbige Brod oder Speise wird Manna genannt, sein schön und klar, wie ein weißgeformtes Mehl.

\begin{schriftzitat}
V. 13–15. Und am Morgen lag der Thau um die Gezelt her, und als der Thau gefallen war, siehe, da lag etwas in der Wüste dünne und klein, wie der Reif auf dem Lande ist. Und da es die Kinder von Israel sahen, sprachen sie untereinander: Das ist Man, denn sie wußten nicht, was es war.
\end{schriftzitat}

33. Sie sprechen: Es ist etwas; denn es ist ein neu Wunderwerk gewesen, darum so hat es keinen Namen, dieweil es noch ist eine neue Speise gewesen. Man heißt auf hebräisch, Barschaf, Geschenk oder Gabe. Als sollten sie sagen: Es ist GOttes Gabe, welche GOtt unter uns austheilt. Gleich als wenn wir ein fremd, unbekannt Ding bekämen, das doch gut und scheinbarlich wäre, und wir hießen es Manna. Wir Deutschen heißen es aus dem 6. Capitel Johannis, V. 31., Himmelsbrod, aber sie haben es GOttes Gabe genannt.

\begin{quote}
Mose aber sprach zu ihnen: Es ist das Brod, das euch der HErr zu essen gegeben hat.
\end{quote}

34. Als sollte er sagen: Da habt ihr es: GOtt hält seine Zusage; nun esset und sättiget euch davon.

\begin{schriftzitat}
V. 19. 20. Und Mose sprach zu ihnen: Niemand lasse etwas davon über bis morgen. Aber sie gehorchten Mose nicht. Und etliche ließen davon über bis morgen. Da wuchsen Würmer drinnen und wards stinkend. Und Mose ward zornig auf sie.\footnote{Diese Textesworte fehlen in den Ausgaben, wiewohl das Folgende ohne dieselben völlig unverständlich ist.}
\end{schriftzitat}

35. Das ist ja gar hoch versucht und GOtt ein Ziel gesteckt, da sie GOtt gleich getappt, gegriffen und gefühlt haben, und er hat sich lassen sehen in einer Wolke. Das ist wahrlich GOtt versuchen, welches er doch hart verboten hat, daß man nicht soll sagen: Gib mir es also, oder auf diese Stunde; daß ich GOtt eine Person, Ort, Weise und Raum vorstellte, durch welche, wo, wann, und wie er mir helfen solle; sondern ich soll es ihm heimstellen, und allein bitten um die Nothdurft; daß sie denn allhier nicht gethan haben.

\quellseite{481}
36. Also haben wir von einem großen Mirakel gehört, so GOtt mit dem Manna gethan hat, damit er die Kinder von Israel vierzig Jahre gespeist, und täglich neu Manna gegeben, und nicht mehr gereicht hat, denn sie haben aufessen können. Das dünkt uns eine schlechte That sein; aber es ist ein groß Werk gewesen. Wie denn solche Wohlthat GOttes, mit dem Manna, allenthalben in der heiligen Schrift gerühmt wird.

37. Und GOtt hat es erstlich darum gethan, daß er seine Herrlichkeit wider die Ungläubigen und Lästerer wollte sehen lassen, daß er allein der rechte GOtt wäre. Wie denn Moses sagt: „Im Abend sollet ihr inne werden, daß euch der HErr aus Egyptenland geführet habe, und des Morgens werdet ihr des HErrn Herrlichkeit sehen“ etc., und also beweisen, daß aus GOttes Befehl und Geheiß Moses und Aaron die Kinder von Israel hätten aus Egypten geführt. Darnach so hat GOtt auch damit wollen anzeigen, daß er seine Gläubigen, so in ihrem Berufe wandeln, speisen und ernähren wolle; wie denn im 5. Buch Moses am 8. Capitel, V. 3. 16., gesagt wird: „GOtt hat dir Manna gegeben, da du, noch deine Väter nicht wußten, was es für eine Speise war, auf daß er dir anzeigete, daß der Mensch nicht allein von dem Brode lebe, sondern von einem jeglichen Wort, so aus dem Munde GOttes gehet.“ Denn GOtt könne die Seinen auch ernähren, wo gleich kein Brod ist; als allhier in der Wüste, da weder Acker, Korn, Mühle, Brod, Bäcker noch anderer Vorrath war; dennoch, ehe sie sollen Hungers sterben, so muß ihnen vom Himmel Manna gegeben werden. Gleichwie Elias, zur Zeit der Theurung in Israel, am Bache Crith [1 Kön. 17, 5. 6.] auch wunderlich ernährt und gespeist war. Und diese Verheißung gibt auch der HErr Christus [Matth. 6, 33.], so man „erstlich sucht das Reich GOttes und seine Gerechtigkeit, alsdann solle das andere alles uns zufallen“. Es sagt auch der 37. Psalm, V. 25.: „Ich bin jung gewesen, und alt worden; ich habe aber nie gesehen den Gerechten verlassen, noch seinen Samen nach Brod gehen.“

38. Aber dies Mirakel mit dem Manna hat bei ihnen wenig geholfen; dieweil das Manna täglich gefallen, so ist es bei ihnen gar gemein worden, und haben es nicht geachtet. Gleich als daß uns jetzt die Sonne täglich aufgeht, welches denn auch ein groß Wunderwerk ist, daß sie nicht Einen Tag, oder zween Tage, sondern viel Monden, ja, alle Jahre so scheint; aber es ist jetzt also in die Gewohnheit gekommen, daß wir gedenken, es könne nicht anders sein. Vergleichen ist uns auch kein Wunder, daß jährlich Wein und Korn wächst. Es sollte billig GOtt, wie mit diesen, also auch mit andern täglichen Mirakeln, unsern Glauben ermuntern; denn, daß aus der Erde ein Weizenkorn oder anderlei Gewächs kommen soll, ist so ein groß Mirakel, als wenn GOtt aus dem Himmel das Manna auf diesen heutigen Tag noch gäbe; aber es hat nicht das Ansehen, so macht es das Gewöhnliche verächtlich.

39. So haben sie nun für eine jede Person oder Haupt des Tages einen Gomor\footnote{2 Mos. 16, 36.: „Ein Gomor ist der zehende Theil eines Epha.“ Ein Epha rechnet man zu 15 Berliner Metzen. Sonach wäre der Gomor 1½ Berliner Metze. (Karl von Raumer, „Der Zug der Israeliten aus Aegypten nach Canaan“, S. 27.)} gesammelt. Wie viel deß gewesen sei, weiß man nicht eigentlich; denn das Pfund, die Ellen und Maß verändern sich; wo man nur über zwei oder drei Meil kommt, bald stimmen sie nicht mehr überein. Gleichwie der Scheffel, die Sprache, die Sitten, und äußerliche Ceremonien mit einander auch in der Welt ungleich sind. Denn so manch Land, so manche Sitten und manches Maß ist. Der Gomor ist so groß gewesen, daß darein gegangen ist, als viel sie haben auf Einen Tag essen können. Nun ist es gleichwohl ein sonderlich Mirakel gewesen, daß der Kleine nicht mehr denn der Große, und der Große nicht mehr denn der Kleine hat verzehren mögen; einer hat so viel empfangen, als der andere, der dritte oder vierte eben also.

40. Es kann auch also verstanden werden, daß ein jeder mit Haufen zugetragen hat; der eine Hand voll, jener eine Schüssel, der einen Hut voll, und haben es hernach zu Hause ausgetheilt durch dies gemeine Gemäß. Darnach ist einer gekommen, der wenig gebracht, der kriegt auch aus dem Gomor, so viel als ein anderer, daß er genug hatte; gleichwie der Text auch also lautet. Es kann einerlei Verstand sein, daß sie es durch dies gemeine Maß ausgemessen haben. Item, es kann auch also sein, daß ein jeder viel oder wenig gesammelt habe, und wenn er beim gekommen ist, so hat er nicht mehr noch weniger bekommen, es hat mit dem Messen zugetroffen.

\quellseite{482}
41. Man spürt auch beineben diesem, so jetzt vom Austheilen gesagt ist, daß sie sind geizig gewesen und nicht zu erfüllen, und haben GOtt nicht vertrauet, daß er ihnen Brod geben würde, in dem, daß sie, wider das Verbot GOttes, über das gesetzte Maß bei sich brachten. Aus welcher Uebermaß Würmer worden sind, auf daß GOttes Wort erfüllt würde. Und darob, daß es sonst auch mit ihnen nicht fortging, mußte die edele Speise verderben. Es ist ein großes Heer und ein großer Haufe Volkes gewesen, darum haben sie auf den Abend alles rein aufgegessen. Da hat denn GOtt des Morgens abermal ihrer so viel zu speisen gehabt. Es ist von der Hitze des Tages nicht zerschmolzen, man hat es gebacken, gekocht und gebraten, und sind nicht zufrieden damit, daß es geschmeckt hat als Semmelmehl, und wie ein jeder selbst gewollt, sondern sie haben auch anderes draus machen und suchen dürften.

42. Zu dem ist ihnen dies ein Mirakel gewesen, daß sie am sechsten Tage zween Gomer gefunden; darum gingen sie hinein und fragten Mosen, und sprachen: Wie meinst du es, daß wir des Tages nur Einen Gomer sammeln sollen, oder es würden Würmer draus machen; haben wir doch heute zween Gomer? Was antwortet er darauf?

\begin{schriftzitat}
V. 22–24. Alle Obersten der Gemeine kamen hinein und verkündigten es Mose. Und er sprach zu ihnen: Das ist es, daß der HErr gesagt hat: Morgen ist der Sabbath der heiligen Ruhe des HErrn; was ihr backen wollt, das backet, und was ihr kochen wollt, das kochet; was aber übrig ist, das lasset bleiben, daß es behalten werde bis morgen. Und sie ließen es bleiben bis morgen, wie Moses geboten hatte; da ward es nicht stinkend, und war auch kein Wurm drinne.
\end{schriftzitat}

43. Wer hielt es da, daß es nicht verdarb? GOtt hat es also geboten, da hat er sein Wort drüber gegeben, er erhielt es auch, daß [es] nicht madig würde noch stinkend. Aus diesem siehst du, daß der Sabbath gewesen sei, ehe denn das Gesetz Mosis gekommen ist, und ist auch wohl von der Welt Anfang her gewesen [1 Mos. 2, 2.], sonderlich daß die Frommen, die den wahrhaftigen Gottesdienst gehabt, an diesem Tage zusammen gekommen sind, und GOtt angerufen haben; davon ist sint der Zeit her derselbige Tag genannt worden der Sabbath, und ist ein geheiligter Tag gewesen, abgesondert von andern Tagen, auf daß alba Menschen und Vieh ruheten. Aber darnach ist aus dem Himmel dieses Gebot durch Mosen öffentlich gegeben, daß sie am siebenten Tage ruhen, und mit Arbeiten und andern Unheiligen inne halten sollten.

44. Zuletzt wollten sie es zu viel machen, da stäupt er auch unter sie, wie ihr denn hören werdet; darum so mußten sie den Gomer auch behalten bis auf ihre Nachkommen [V. 33.]. Die Juden ziehen diese That an Johannis am sechsten Capitel, V. 31., des HErrn Christi Wort damit zu verlegen, und rühmen: „Unsere Väter haben Manna gegessen in der Wüste, das vom Himmel kam.“ Darauf antwortet Christus meisterlich, und spricht, V. 32.: „Moses hat es euch nicht gegeben, sondern mein Vater gibt euch das rechte Brod vom Himmel.“ Und sagt weiter, V. 49. 50.: „Sie sind alle gestorben, die es gegessen haben; wer aber dieses isset, der stirbt nicht.“

45. Also hört ihr sonst mehrmal: Alle GOttes Gaben und Wunder, wenn sie nicht gespeiset sind mit oder in dem Worte GOttes, so helfen sie nichts. Das Wasser in der Taufe, wenn nicht GOttes Wort dabei wäre, ist nichts. Also, wenn bei dem Manna nicht GOttes Wort wäre, so hätte es die Kraft nicht. Das ist es nun, das der HErr Christus in vorhergehenden Worten meint: Es haben’s ihrer viel gegessen, und sind gleichwohl gestorben, denn sie haben diesen Anhang des Worts nicht mitgeessen, darum sind sie auch nicht erhalten. Dergleichen sind auch alle Creaturen GOttes, wenn sie ohne GOttes Wort empfangen werden, so ist kein Leben, keine Macht noch Kraft in ihnen, sondern es ist ein leiblich Ding.

\subsection{Allegorie des Manna}

46. So hat nun GOtt durch das Manna wollen bedeuten das rechte Manna, welches er, Christus, selbst ist; wie er denn im Texte Johannis [Cap. 6, 50. 58.] es also deutet und spricht: „Dies ist das Brod“, und zeigt auf das große Manna, auf sich selbst; als sollte er sagen: Da werdet ihr finden, was das Manna sei, das eure Väter gegessen haben.

47. Ihr mißt aber, wie und welcher Gestalt Christus das Brod sei. Man muß ihn nicht also essen, wie es die Juden verstanden, da sie
\quellseite{483}
sagten [Joh. 6, 52.]: „Wie kann dieser uns sein Fleisch geben?“ Denn sie meinten, sie müßten ihn mit den Zähnen zerreißen, und mit dem Munde trinken. Aber wie er eine geistliche Speise ist, also isset man ihn auch geistlich. Essen ist allhier glauben; mit dem Glauben muß man die Speise empfangen, daß Christi Fleisch für uns getödtet, und sein Blut für uns vergossen ist. Christus will sagen: Wer da glaubt, daß ich vom Vater gesandt sei, daß ich mein Fleisch und Blut für euch, die ihr selig werden sollt, gebe, derselbige hat das rechte Manna, und wird nicht sterben. Denn das ist unser Glaube, daß, der da glaubt an Christum, gestorben für unsere Sünde, der soll selig werden; man muß glauben an Christum, gestorben und auferstanden von den Todten. Wer isset und trinket Christi Blut, der hat das ewige Leben. Es sind andere Worte, aber es ist einerlei Meinung. Also bleibt es wahr, wer das geistliche Manna isset, der stirbt nicht. Eure Väter haben dies geistliche Manna nicht gegessen, das ist, an Christum nicht geglaubt, darum sind sie gestorben.

48. Diesen Text haben unsere Widersacher, die Sacramentischwärmer, zu ihrem Behelf genommen, und haben ihn fälschlich auf das Sacrament gedeutet. Wiewohl es wahr ist, daß Christi Blut und Fleisch allda gegenwärtig ist, so wird es auch mit dem Munde empfangen, leiblich isset und trinket man es. Aber diesem Essen rechnet man nicht zu, daß es lebendig mache. Denn da gehen viele dazu, essen und trinken, und sterben doch, sie gehen dazu zur Strafe und zum Gerichte, welches denn von St. Paulo hart gestraft wird; derhalben so haben sie nicht das Leben. Und der Sequenz, den man singt am Tage corporis Christi: Manna datur patribus, ist nicht recht, es wird falsch also gesungen, und ist ein närrischer Sequenz, sondern du sollst es auf den Glauben und auf das göttliche Wort ziehen. Das Wort wird mit dem Munde nicht empfangen und gegessen, wie das Sacrament des Altars, sondern das Wort wird mit dem Herzen ergriffen; wenn man glaubt, und wenn man es hinein in das Herz und Glauben zeucht, so ist es recht. So will nun der HErr Christus im Johanne sagen: Der Juden Manna hat mich, Christum, bedeutet.

49. Will man nun sagen: Wie kannst du das Brod sein, so man isset und trinkt? Also geschieht es, wenn man mich austheilt. Mit was Meßer? Daß man durch den Mund und Zunge von Christo predige, und vermeldet, daß er für uns sein Blut vergossen habe und gestorben sei; da lege ich dir Christum vor, nicht auf Tellern, sondern in das Herz lege ich dir sein Fleisch und Blut, das für dich gegeben ist. Das heißt in das Herz gelegt. Wenn du es also empfängst, issest und trinkst, das ist, wenn du es also glaubst, so bist du selig [Offenb. 19, 9.]. Das ist es, das der HErr Christus sagt: „Wer mein Fleisch isset, der wird nimmermehr sterben.“ Sonst geht es alles dahin. Kurzum, ihr seid alle verloren, wo ihr nicht mein Fleisch esset, das ist, mit dem Herzen mich esset und an mich glaubt, und ich euch durch das Wort vorgetragen werde.

50. Also habt ihr das rechte Manna, das GOtt durch jenes bedeutet hat. Und kürzlich zu reden: Manna ist uns jetziger Zeit das Sacrament und das Evangelium, nämlich, diese Predigt, darinnen gehandelt wird, daß der Leib und das Blut Christi für uns gegeben sei zur Vergebung der Sünde. Darum, wenn du das Evangelium hörst, so regnet es Manna, und da sammelst du Manna, das wird dir dann vom Himmel gegeben. Wenn du es nun nimmst, das ist, glaubst, so isset du auch, und bleibst ewig, gleichwie das Wort GOttes auch ewig bleibt.

51. So ist auch das Manna weiß gewesen, gleichwie ein Thau, wie ein Schnee oder Reif; das ist, das Evangelium ist nicht allein eine säuberliche, feine, reine Predigt, sondern auch eine liebliche, süße Predigt. Moses ist ein blutiger, rother Prediger, er führt roth im Schild, er ist ein scharfer, zorniger Prediger, er tödtet einen Menschen; aber das Evangelium ist weiß, rein, und hat keinen Flecken, es hat die schönste Farbe, und ist dem Lichte am nächsten. Denn das Licht ist die allertröstlichste Farbe; aber Finsterniß sind schrecklich. Gleichwie auch in Morgenländern Könige, Fürsten und große Herren haben pflegen weiße Kleider anzuziehen, und weiße Fächeln\footnote{Fächel (Fächel) — Binde, Schnur um Hüfte, Stiefeln und Kleider. Vgl. Col. 558.} um die Hüfte zu tragen. Weiß ist eine fröhliche Farbe, lieblich und tröstlich. Und ist die weiße Farbe die beste Farbe, von aller Schwärze abgesondert, neigt sich am nächsten zum Lichte.

52. Darum so geht sie nirgendshin, denn auf Christum, sie will nicht, daß man ihr etwas zueigen soll; das ist: Wer da predigt, daß dich
\quellseite{484}
Christus allein selig mache, der läßt das Manna weiß und bei seiner rechtschaffenen Farbe bleiben. Item, es ist das Evangelium also süß und lieblich, es verkündiget eitel Trost, Frieden und Vergebung der Sünde, so durch den Tod Christi uns widerfahren ist. Wie zu dem Gichtbrüchigen der HErr Christus [Matth. 9, 2.] sagt: „Sei getrost, mein Sohn, dir sind deine Sünden vergeben“; item [Matth. 11, 28.]: „Kommet zu mir alle, die ihr mühselig und beladen seid, ich will euch erquicken“ etc.

53. Darnach, so ist es gewesen wie Coriander, klein gestoßen, gleich als Hirsien oder Reis; das ist, das Evangelium ist eine geringe, schlechte Predigt. Wie denn der HErr Christus selbst [Matth. 13, 31.] das Evangelium auch einem kleinen Senfkorn vergleicht. Also zerstoßen und gering ist das Evangelium anzusehen, es wird von der Welt verachtet, gelästert und geschändet. Dennoch so hat das evangelische Wort, so wir predigen, so große Kraft, daß es selig macht; es hat eine göttliche Kraft bei sich. Solch Wort GOttes, so aus des Menschen Munde geredet wird und erschallt, das wird mit der Vernunft nicht empfangen, es hat keinen äußerlichen Schmuck noch Ansehen; noch bringet es mit sich den Sieg wider den Tod und Teufel, Hölle und Verdammniß, auch Vergebung der Sünden, und macht zu GOttes Kindern, zu Herren über Himmel und Erde [Röm. 1, 16. 2 Tim. 1, 8. 10.]. Solche Gewalt haben wir im Worte GOttes, welches Wort doch ein schlecht Ding ist; denn es ist um einen Augenblick zu thun, so ist es vorüber [Ps. 147, 15.]. Darum so ist es nicht ein ansehnlich oder prächtig Ding, es hat keinen Schein, und ist doch so kräftig.

54. Darnach so schmeckt es wie Honig und Semmel, das ist, es tröstet und erquickt die Seele. Andere haben allhier phantasirt, daß Honig und Semmel so wohl schmeckt, sei die Gottheit und Menschheit; die Gottheit muß erst ins Brot kommen, soll sie uns zu Nutzen gedeihen und schmecken. Aber sie fallen hinein mit dem vita contemplativa [mit dem beschaulichen Leben]; der Satan stürzt sie also herab. Es wird niemand die Gottheit schmecken, denn wie sie will geschmeckt sein, nämlich, daß sie in der Menschheit Christi betrachtet werde; und wenn du nicht also die Gottheit findest, so wirst du nimmermehr Ruhe haben. Darum laß sie immer anhin speculiren und reden von der Beschaulichkeit, wie alles mit GOtt buhle, und einen Vorschmack des ewigen Lebens gebe, und wie die geistlichen Seelen ein beschaulich Leben anfangen.

55. Aber lerne du mir GOtt nicht also kennen, laß ihn sein, wie groß und mächtig er ist; aber hebe allhier an, und rühre ihn an, und sage erstlich: Ich weiß von keinem GOtt, denn der für mich gegeben ist. Also wird dir seine Majestät süß sein, wenn du weißt,\footnote{Eißlebensche: wisset.} du habest einen solchen GOtt, der für dich gegeben sei, welcher Himmel und Erde in seiner Hand hat; dann wird dich keine Creatur, noch diese deine Majestät nicht erschrecken. Denn daß er für dich gegeben ist, das ist seiner Barmherzigkeit und Liebe Schuld; allhier fasse ihn, und da ist er auch weich und süß [1 Mos. 32, 30. 5 Mos. 33, 2. 3. Joh. 20, 20.].

56. Sonst ist GOtt zweierlei; als GOtt für sich, da er ist ein HErr und GOtt in seinem unbegreiflichen Wesen. Darnach ein GOtt, der für dich gegeben ist, daß GOtt ein HErr sei über dich und alle Menschen. Davon predigt Moses, und er will es auch von allen haben, daß man ihn für einen HErrn halten soll. Aber allhier fleucht jedermann vor diesem GOtt; denn er ist gar zu herbe und schrecklich. Aber da ist es nun ein ander Ding, wenn dir das Evangelium von demselbigen GOtt predigt, und ihn dir recht vorbildet, und macht, daß du Lust und Freude an ihm habest; nämlich also, wenn es bezeuget und weist, wie er für dich gestorben und gegeben sei; da ist er süß und wie Honig.

57. Zum sechsten so ist das Manna nicht von Menschenhänden gemacht, noch aus der Erde gewachsen, sondern vom Himmel herabgekommen, gleichwie der Thau herab fällt. Dahin haben die Propheten gesehen, wenn sie von dem heiligen Evangelio geweissagt oder gepredigt haben. Denn es ist das Evangelium eine solche Predigt, die von keinem Menschen erdichtet ist, wie St. Petrus in seiner [zweiten] Epistel [Cap. 1, 21.] auch sagt: „Keine Weissagung kommt aus menschlicher Wahl oder Andacht, sondern die Heiligen GOttes haben geredet, getrieben durch den Heiligen Geist.“ Item, der HErr Christus spricht [Matth. 10, 20.]: „Ihr seid es nicht, die da reden, sondern der Geist meines Vaters, der in euch ist“ etc. Das Evangelium wird auch darum vom Himmel kommend genannt
\quellseite{485}
, denn die Vernunft versteht es nicht, sie weiß nicht, daß GOtt seinen Sohn für mich gegeben, daß der Sohn in die Welt gesandt sei, und für mich Mensch ist worden und gestorben. Davon kann niemand predigen, niemand kann solche Lehre erfinden, sie komme denn von oben herab wunderbarer Weise, ohne Menschen Zuthun [Joh. 3, 27.]. Das ist so viel gesagt: Das Evangelium ist eine himmlische Predigt, nicht ein menschlich Wort, man kann es auf dem Erdreiche nicht finden. Darum so sind alle andern Lehren erdichtet, falsch und Teufelslehren, denn sie kommen nicht vom Himmel.

58. Es kommt das Manna auch des Morgens, wenn die Sonne aufgeht; das ist, das Evangelium ist aufgegangen, und kommt zur Zeit der Gnaden, da die Christenheit anging, und die Sonne der Gerechtigkeit, das ist der HErr Christus, aufgegangen oder auferstanden ist. Denn nach des HErrn Christi Auferstehung und Himmelfahrt kommt der Heilige Geist, und füllt dieser Thau in die ganze Welt. Darum ist das Evangelium eine Speise, von GOtt zubereitet ohne unser Zuthun. Also haben wir auch Manna, die Himmelsspeise, so wohl zugerichtet und uns vorgelegt wird. So hat er uns das Evangelium geschenkt, es ist unsere Gabe; wir haben nicht darum gebeten, noch darnach getrachtet, sondern es kommt aus lauter Gnaden, und offenbart sich uns.

59. Daß aber das Manna also aufgelesen wird, daß ein jeder so viel kriegt als der andere, und daß man es nach dem Maß Gomer mißt, das hat diese Bedeutung, daß ein Christ dem andern gleich ist, und keiner mehr oder weniger hat denn der andere, sie haben alle Einen GOtt, Einen Christum, Eine Taufe, Geist und Glauben etc. [Eph. 4, 4–6.] Das mag wohl sein, daß die Gomor ungleich gewesen sind, eines steinern, das ander hölzern, das dritte kupfern oder gläsern. Also haben auch unter uns etliche einen schwachen Glauben, die andern einen starken Glauben, mehr oder wenigern Unterricht vom Glauben, so ist doch solcher aller Christus. St. Petrus hat nicht einen bessern Christum, denn ich habe, und wie viel St. Petrus hat, also viel habe ich auch).

60. Nach dem Namen Christi soll sich keiner des andern schämen, noch besser denn der andere sich dünken lassen, denn der einige Christus ist für mich und dich gegeben. Ich habe ihn ja so viel gestanden, als ihn Maria, die Propheten und alle Heiligen gehabt haben; und dieweil ich an denselbigen glaube, so habe ich so viel an ihm, als ein anderer. Wir sind gleich, mein Maß hat so viel in sich, als jener; es ist gleich und so groß, als jener; obwohl, die da schwach sind in Widerwärtigkeit und in Anfechtung, ein Geringfügigeres haben, und die Bekenner\footnote{„Bekenner“ wird hier wohl so viel sein als Märtyrer (confessores). Diese Bedeutung fehlt bei Diez.} haben ein scheinbarlicher Maß, so sind sie doch, nach Christo zu rechnen, alle durchaus gleich, und haben einerlei Gut des Glaubens, das ist, sie haben den ganzen Christum. Dieses ist gar ein tröstlich Stück, daß, gegen Christo gerechnet, alle Heiligen und Gläubigen gleich sind.

61. Dieses sollen wir wissen wider die Lumpenprediger, die uns einen Unterschied gemacht haben durch die Heiligen, und haben uns aus den Heiligen Abgötter gemacht. Ich rede aber allhier von den Christen, wo deren einer ist, ob er sich wohl allen unterwirft und gegen jedermann demüthigt, so ist er doch vor GOtt so groß als Petrus; ob er gleich nicht so große Gaben hat, und so viel gesammelt als andere, so hat er doch in dasselbige Maß gesammelt, und hat auch ein gleich Maß, als St. Petrus, das ist, eben denselbigen Christum; darum so darf er sich nichts rühmen. Aber er mache nicht eine äußerliche Gleichheit draus, er ziehe es nicht in äußerliche Sachen. Denn darum heißt man nicht Christen.

62. Der HErr Christus hat im Evangelio einerlei Samen, welcher doch an viererlei ungleichen Orten wird ausgestreut. Also heißt alles Christen, und haben das Evangelium alle; aber das vierte Theil des Samens bleibt nur gut, und bringt seine Frucht. Solches Christenvolk habe ich nie auf Erden gesehen, noch erkannt; wenn man sich unter einander rottirt und selbst hoch hält, so ist es falsch. Es soll eine Stadt sich nicht nennen eine christliche Stadt, sie habe denn den rechten Beweis des Christenthums. Darum so gib das vierte Theil heraus, oder das andere Theil, und rechne, so viel du willst; die Personen kann ich nicht zählen, aber dies kann ich sagen: Wo das Evangelium ist, daselbst sind Christen. Item, ich kann sagen: Inwendig sind sie alle gleich, auswendig aber sehr ungleich, da hat einer mehr, und der andere weniger gesammelt.

\quellseite{486}
63. Aber was wollen wir zu dem Spruche St. Pauli sagen, 2 Cor. 8, 13–15., da er dieses auf die Almosen und Wohlthun zeucht, und mit wunderlichen Worten bettelt und vermahnt, daß man Armen, Dürftigen helfen soll? jedoch spricht er, man soll nicht geben, daß sie vollauf hätten, und andere darben und mangeln müßten. Wie denn faule Tropfen pflegen herauszunehmen, und lassen andere hinein geben; die lassen es sich viel kosten, jene wollen es mit Müßiggang und Faulenzen verschwenden. Nun dies sind seine Worte [2 Cor. 8, 13–15.]: „Nicht geschiehet das der Meinung, daß die andern Ruhe haben und ihr Trübsal, sondern daß es gleich sei. So diene euer Ueberfluß ihrem Mangel, diese theure Zeit lang, auf daß auch ihr Ueberfluß hernach diene eurem Mangel, und geschehe, daß gleich ist. Wie geschrieben stehet: Wer viel sammelte, hatte nicht Ueberfluß, und der da wenig sammelte, hatte nicht Mangel.“ Wie reimt sich doch das daher auf das Almosen?

64. Antwort: Ein anderes ist es, Exempel führen, und ein anderes, geistliche Deutung machen. Das Exempel wird auf das Werk gezogen, und die Deutung geht auf den Glauben. Als: Christus ist für uns gestorben, und [hat] uns ein Exempel gegeben, daß wir hernachher auch sterben sollen; wenn ich denn sterbe, so thue ich, wie Christus gethan hat. Dieses ist nicht eine Bedeutung, sondern ein Exempel und Nachfolge. Derhalben so soll man es auch ein Exempel, und nicht eine Allegorie heißen. So will nun St. Paulus so viel sagen, und sie unterrichten: Ich wollte es gerne also machen, daß es dem Exempel nach ginge, welches dort unter der Israeliter Zeit in der Wüste gewesen ist, als da, der wenig Manna aufliest, bekommt doch seinen Theil gleich sowohl, wie der, so viel gesammelt hat.

65. Also sollte es auch mit den zeitlichen Gütern zugehen. Ihr habt der geistlichen Güter reichlich, und nicht weniger als jene; ihr habt das Evangelium von GOtt, den reichen Schatz, der bei euch ausgetheilt. Darum so seid ihr von GOttes wegen nun schuldig, das Zeitliche andern mitzutheilen. Und nun ist es gleich, daß, wer da wenig hat, der hat so viel, als der da gar viel gesammlet; wer den geistlichen Schatz hat, der wird den andern gleich. Jetzt habt ihr nicht weniger denn ich, oder die andern Apostel, wiewohl wir es erst gehabt haben, und ihr habt es von uns empfangen. Darum, wer unter euch an der leiblichen Nahrung vermögend und reich ist, der theile aus denen, so nothdürftig sind, oder [denen,] die euch das Geistliche mitgetheilt haben, die da predigen, und im Kampfe sind, und euren geistlichen Reichthum mehren, auf daß ihr reich werdet und krieget den Schatz; so gebet ihnen wiederum, und nähret sie, damit sie genießen eures zeitlichen Guts, wie ihr ihres geistlichen genießet und mit theilhaftig werdet.

66. Daher sagt er auch in der ersten Epistel zum Corinthhern Cap. 9, 11.: „Wenn sie das Geistliche euch aussäen, ist es ein Großes, wenn sie das Zeitliche wieder einernten?“ Sie speisen eure Seelen; ihr aber ernährt nur unsern Leib; wir dienen euch mit dem Schatz, daß wir das Evangelium euch predigen, darum sollt ihr auch Unterhaltung geben, und uns ernähren. Also trägt eines das andere. Darum so ist dies nicht eine Allegorie, sondern ein Exempel.

67. Nun sagt der Text weiter: Was überblieb, da wurden Würmer draus. Denn es war von GOtt geboten, man sollte nichts auf den Morgen lassen überbleiben, gleichwie man an dem Osterlamme nichts mußte lassen überbleiben. Nach der Historie ist es das, daß GOtt den Unglauben nicht haben will, daß man ihm nicht vertraut noch glaubt; darum so straft er sie, daß alles, das übrig ist geblieben, ihnen nicht hat müssen zu Nutz oder zugute kommen. Das ist, der Unglaube macht es madig und stinkend, und daß es ein Greuel wird; wie sie denn auch inwendig gewesen sind, ob sie wohl außerhalb dem Manna vor ihren Augen sich fromm und heilig vor GOttes Augen dünken ließen; darum will er sie nicht lassen bleiben\footnote{„bleiben“ von uns gesetzt statt: „leiden“ in den Ausgaben.} um des Unglaubens willen.

68. Es ist aber die geistliche Bedeutung, daß die Juden durch das alte Manna das Gesetz verstanden haben, dahin haben sie es deuten und ziehen wollen, daß über einen Tag nicht währt. Sie haben den gestrigen Tag, wir den heutigen Tag; ihr Gesetz ist aus, das Evangelium ist nun aufgegangen; darum so wird nun das Gesetz stinkend und madig, es wird greulich und verwerflich. Moses wird zornig über die Juden, das ist, das Gesetz verdammt sie; aber sie fragen nichts darnach. Also habt ihr die Historie von dem Manna.

\quellseite{487}
\section{Das siebenzehnte Capitel}

\begin{schriftzitat}
V. 1. Und die ganze Gemeinde der Kinder Israel zog aus der Wüste Sin, ihre Tagreisen, wie ihnen der HErr befahl, und lagerten sich in Raphidim. Da hatte das Volk kein Wasser zu trinken.
\end{schriftzitat}

1. Aus dem vorhergehenden Capitel ist geredet von dem großen Wunderwerk, so GOtt an den Kindern Israel gethan hat, daß er sie in der Wüste vierzig Jahre mit dem Himmelsbrod ernährt; wiewohl sie dennoch für und für ungläubig geblieben, und wider GOtt gemurret haben, und dieses Wunderwerk (wie denn auch alle anderen Mirakel) in Wind geschlagen; wie wir denn dergleichen auch noch auf den heutigen Tag thun. Aber GOtt hat mit demselbigen Wunderzeichen beweisen wollen seine Macht auf Erden unter den Menschen; aber wir lassen uns dadurch nicht bewegen zum Glauben.

2. Nun fährt Moses fort und beschreibt weiter andere Wunderthaten, die GOtt an ihnen gethan habe. Denn wir haben allhier zwei Zeichen, so er diesem Volke zugute bewiesen: Eines, daß er ihnen Wasser gibt zu trinken in der Wüste. Darnach, daß er sie durch Mosis Hand wunderbarer Weise in dem Kriege wider die Amalekiter errettet und erhält.

3. Das eine Stück wollen wir erst besehen. Diese Historien gehen nach einander, gleichwie die Gebirge in der Wüste, da ein Berg auf den andern folgt; wie denn die Gelehrten wissen, wie diese Berge nach einander gelegen sind. Mara und Etham ist ein groß Gebirge in Arabia; also sind allhier die Berge Horeb und Raphidim nahe bei einander gelegen, und haben doch andere Namen; denn es wird hernach der Berg Horeb auch Sinai genannt, und Raphidim liegt hart dran. Sie sind einen krummen und langen Weg in der Wüste gezogen durch das Gebirge.

4. Da folgt nun ein Mirakel, nachdem GOtt das Volk Israel versucht, und läßt sie dürsten. In Elim hatten sie Wassers genug; jetzt sind sie in der Wüste, da ist gar kein Wasser, und scheint wieder, gleich als wäre kein GOtt, der nach ihnen fragte. Und es macht’s der Mangel, daß sie der vorigen Mirakel alle vergessen, als, da das bittere Wasser süß ward; item, er ihnen auch das Himmelsbrod gegeben, und daß der Engel des Tages in der Wolkensäule und des Nachts in der Feuersäule ihnen vorgegangen. Das ist nun eine schändliche Klage, daß uns der Bauch regieren solle, und hindern, daß wir die gegenwärtigen Mirakel nicht sehen. Uns dünkt jetzt wohl: O wären wir damals vorhanden gewesen, und hätten die Wunderwerke gesehen, so wollten wir gerne geglaubt haben. Ja wohl!

5. Es sagt der HErr Christus Matthäi am 6. Capitel, V. 28. 29.: „Sehet die Blumen oder Lilien auf dem Felde an, wie sie wachsen; sie arbeiten nicht, auch nähen sie nicht, noch ist Salomo in aller seiner Herrlichkeit nicht bekleidet gewesen, als derselbigen eines.“ Item, V. 26.: „Sehet die Vögel unter dem Himmel an, sie säen nicht, sie ernten nicht, sie sammeln nicht in die Scheuern, und euer himmlischer Vater nähret sie doch.“ Dieses ist ein täglich Mirakel, so wir sehen und erfahren, wie GOtt unzählig viel Vögel ernährt, und allen Blumen ihre lieblichen, schönen Farben und Schmuck gibt, zum Zeugniß eines großen Wunderwerks: noch geht es im gemeinen Haufen niemand zu Herzen.

6. Und redet der HErr Christus noch gröber und deutlicher davon, und spricht [Matth. 6, 27.]: „Welcher ist unter euch, der seiner Länge eine Elle möchte zulegen?“ Als sollte er sagen: Ihr seid die nicht, die ihr euch selbst hättet geschaffen oder gemacht, ihr habt euch selbst nichts gegeben: noch habt ihr nicht so viel Glauben, daß ihr unserm HErrn GOtt vertrauen könntet. Und wir müssen alle sagen, daß wir nicht einen Augenblick unseres Lebens sicher sind, und müssen dazu bekennen, daß eines andern Macht und Gewalt sei, daß wir ernährt werden, nämlich unseres lieben GOttes. Dieses Mirakel haben wir täglich alle im Hause, tragen es an unserem eigenen Leibe, noch sollen wir nichts darnach fragen. Der HErr Christus spricht noch dazu [Matth. 6, 25.]: „Ist nicht der Leib mehr denn die Kleider?“ Als sollte er sagen: Ihr sorgt für die Kleider, so GOtt doch euer Leib und Leben erhält, und auch dieselbigen Kleider euch verschafft.

\quellseite{488}
7. Es mag nun der Juden spotten und lachen, wer da will, und sagen: Pfui über Juden, daß sie solches nicht betrachtet und GOtt geglaubt haben. Aber höre du, greif in deinen eigenen Busen, so wirst du darinnen auch einen solchen ungläubigen Juden finden; und pfui dich auch an, solltest\footnote{Eislebensche: sollst.} du dich nicht auch schämen, daß man so einen groben Esel in dir finden sollte, der es für kein Mirakel hält, daß GOtt dir dein Brod, Kleider, Leib und Leben noch täglich nothdürftiglich verleiht, Speise und Nahrung gibt? Aber niemand denkt daran, und dankt ihm dafür, daß er das Leben erhält. Es ist ein gewöhnlich Ding, alt sein; niemand dankt GOtt dafür, wenn er des Morgens aufsteht, daß ihn GOtt die Nacht über gnädiglich behütet hat. Also ist auch das Himmelsbrod ein täglich Ding gewesen. Aber wenn es GOtt wieder entzeucht, als, daß sie allhier nicht Wasser haben, da wollte man es denn gerne wieder haben; ja, es wäre sonst kein Wunderwerk gewesen. Also, wenn er uns täglich alles gäbe, so achteten wir es ganz und gar nichts (denn siehe es an den Reichen, die alles genug haben, die fragen nach nichts), damit käme man nimmermehr zu GOtt.

8. Derhalben so müssen wir versucht werden, müssen Noth und Mangel haben, und Fehl an Essen und Trinken, an Kleidern, Geld und Gut, an Ehre, Gunst und Förderung bekommen, auf daß wir Ursache gewinnen, GOtt zu suchen, auf daß wir GOttes nicht vergessen, wenn wir satt und voll wären. Denn das ist sonst unsere Art und Natur, es ist bei uns Mangel gar genug; noch sind wir nicht zu zähmen und irre zu machen. Es muß der heilige Moses im 5. Buch [Cap. 32, 15.] auch darüber klagen: Dilectus meus incrassatus, dilatatus et impinguatus, oblitus est Dei Creatoris sui [er ist fett und dick und stark worden, und hat den GOtt fahren lassen, der ihn gemacht hat]. Also bleibt es wohl. Den Leuten ist doch wie dem Esel, der löckt\footnote{löcken = hüpfen, springen; hier: hinten ausschlagen.} hinter sich, das Futter sticht ihn, wird geil, geht auf das Eis tanzen, und bricht ein Bein. Denn, gibt GOtt alles genug, so werden wir muthwillig und sicher, und vergessen unseres HErrn GOttes. Will aber GOtt ein wenig Lob, Preis, Ehre, Anrufung oder Danksagung von uns haben, so muß er uns lassen Mangel und Noth

leiden, das Futter etlichermaßen entziehen, und das Leibliche lassen in die Schanze schlagen: so fragt man nichts nach den ewigen Gütern unseres HErrn GOttes.

9. Die Unachtsamkeit und der Undank haben so gar Überhand bekommen, daß niemand zufrieden ist mit seinem Weibe und Kinde, Haus und Hof, Hand, Maul, Nase, mit seiner Stärke und Gesundheit. Dieses achtet man alles nicht. Und was an diesen Stücken ein anderer hat, beliebet uns viel mehr, denn das Unsere; doch GOtt wirft es in die Rappuse. Aber wenn ein wenig kommt eine kleine Armuth, oder daß GOtt schickt eine Krankheit, oder sonst Gefährlichkeit, alsbald so schreiet man: O hilf GOtt, hilf GOtt. Ja, kommt du nun getrollt? Warum, weil unser lieber GOtt und Vater so hartherzig ist, und wir so blind sind, daß wir seine Güter und Wohlthaten nicht sehen, noch ihm dafür danken, so behält er oft einen kleinen Zuspruch zu uns, auf daß wir Ursache haben, ihn anzurufen.

10. Also seht ihr, was wir für Leute sind, und wer GOtt auch sei; item, was die Welt, Fleisch und Blut thue gegen GOtt. Die Welt ist blind, halsstarrig, und ganz böse: GOtt aber dagegen ist freundlich und gnädig, wohlthätig und von großer Güte, wie ihn also die heilige Schrift, und sein eigenes Wort commendirt.

\begin{schriftzitat}
V. 2. 3. Und sie zankten mit Mose, und sprachen: Gib uns Wasser, daß wir trinken. Mose sprach zu ihnen: Was zanket ihr mit mir? Warum versucht ihr den HErrn? Da aber das Volk daselbst dürstete nach Wasser, murreten sie wider Mose und sprachen: Warum hast du uns lassen aus Egypten ziehen, daß du uns, unsere Kinder und Vieh, Durst sterben ließest?
\end{schriftzitat}

11. Der Bauch zappelt abermals. Es hat der liebe Moses ein beschwerlich, verdrießlich Amt gehabt, daß er so ein unwillig, hartnäckig, halsstarrig Volk hat regieren sollen; es hat ihn auch mürbe genug gemacht, ist daneben in großer Fahr Leibes und Lebens schier alle Stunden gestanden, denn er hat immer warten müssen, wenn sie kämen und ihn erwürgeten.

12. Nun wollen sie wiederum sagen: Wir müssen nicht, ob wir von GOtt aus Egypten geführt sind, sondern du hast es gethan. Das ist ein greulicher Unglaube und erschreckliche Versuchung. Da muß Moses Undank und Lästerung hören, sie schänden GOtt und alle seine
\quellseite{489}
Wunder, dazu auch seine Boten und Propheten; daß muß Moses allhier anhören und leiden.

13. Ein solch fröhlich Ding ist es, den gemeinen Pöbel regieren. Wer das nicht weiß, der hebe nur an zu regieren, und versuche es ein wenig, er wird es bald erfahren. Ich wollte, daß alle seltsamen, wunderlichen, und aufrührischen Köpfe nur zwei Jahre regieren sollten, sie würden die Hörner bald abstoßen. Die Heiden haben gesagt: Magistratus ostendit virum, das ist, man kennt einen Mann nicht, er habe denn zuvor andere Leute regiert, und sei ihnen vorgestanden. Denn an dem Orte pflegt es gewöhnlich also zuzugehen, daß, thut er etwas, so gefällt es einem also, dem andern mit nichten; und wenn es Hunderten gefällt, so sind ihrer Tausend, denen es nicht gefällt. Wenn GOtt einen will zum Märtyrer machen, den setzt er an einen Ort, daß er irgend regieren muß; will er denn sein Regiment mit gutem Gewissen führen, und seine Seele verwahren, dem mag man nicht den Teufel über die Thüre malen, noch zu Gewatter bitten; sondern der gemeine Pöbel wird es wohl ausrichten, daß er so geleimt werde, daß man ihn möchte um einen Finger wickeln.

14. Wenn Moses wäre ein ungeduldiger und unleidlicher Mann gewesen, so hätte er gesagt: In des Teufels Namen, was habe ich von euch? Ich habe nichts für alle meine Mühe, Arbeit und Gutthat, denn Murren, Zanken, item, Undankbarkeit. Das ist wahrlich ein schöner, niedrlicher Lohn, den man ihm dafür gibt. Er theile ihnen das Herz im Leibe mit, und setzte für sie Leib und Leben, Ehre und Gut; aber dagegen hat er dieses Trankgeld davon, daß sie ihn schalten als einen Bösenwicht.

15. Aber also geht es den Christen: sie sollen für Gutes und Wohlthaten einnehmen Schande, Schaden, Nachtheil und Undankbarkeit. In der Welt geht es nicht anders zu. Wer den andern aus dem Kothe zeucht, der stößt ihn darnach wieder hinein; und wer den andern vom Galgen erlöst, der knüpft ihn darnach zur Dankbarkeit wieder hinan.

16. Also ist es auch dem HErrn aller Herren selbst gegangen. Denn da Christus Blinde sehend, Taube hörend, Stumme redend machte, und von allerlei Krankheiten die Leute heilte, und jedermann half, zuletzt schlugen sie ihn an das Kreuz. Eben also ist es allhier auch Mosi begegnet in der Wüste, daß er für seine saure Mühe und Arbeit Gefahr des Lebens muß gewärtig sein. Wer sich nun deß nicht erwägen will, daß er in der Welt Undank einnehme, der denke nur, daß er sich bald niederlege und sterbe. Denn, wenn er will aufrichtig hindurch gehen und regieren, so wird er den Teufel und die Welt bald verdrießen.

17. Aber dies ist die rechte Schule, die uns lehrt, dieses Lebens müde und überdrüssig werden. Denn wenn man in der Welt für Wohlthat und Liebe nichts denn Undankbarkeit, Herzleid und Haß einnehmen soll, so wird man des Lebens bald müde, daß man spricht: Ach, wer todt wäre! Komm, lieber GOtt, und laß mich nur sterben, komm Tod und hole mich! Wie denn Moses hernach in diesem Buche [Cap. 32, 32.] wird sagen: Ach HErr, nimm mich nur hinweg. Ja, so muß man einem das Regiment heiß machen. Es dünkt einen wohl fein sein, daß man über andere herrsche; aber versucht es, so wirst du sagen: Ach, wer ein Bauer auf einem Dorfe dafür wäre, und unter einem grünen Baume mit Frieden säße!

18. Darum so sagt Moses allhier auch: „Was zanket ihr mit mir?“ Ich habe es nicht gethan. Gehet ab zurück! Seid ihr denn nun rasend und unsinnig? Ist es doch unmöglich, daß ein solch herrlich Werk ein Mensch thun sollte, als, Himmelsbrod geben, und Wasser aus einem Felsen fließen lassen. Und weist sie Moses von sich auf unsern Herrn Gott. Es ist Herr Omnes ein toller Teufel, er gebricht nicht, bis daß er auf das Maul geschlagen wird. Mit Wohlthaten verwöhnt man sie, aber nur mit Besen muß man sie stäupen. Es ist Moses alle Tage, Stunden und Augenblicke in Gefahr des Todes. Das ist der Verdienst. Ei, du freundliche, holdselige Welt, du feines Kräutlein! man sollte dich lieben; du bezahlst wohl getreue Dienste; du erkennst kein eines blutsauere Arbeit. Also schändlichen Lohn gibt sie, daß einem wohl gelüsten möchte, daß er mit Fäusten drein schlüge, oder daß man diese schöne Frucht mit Maien besteckte, und dem Teufel zum neuen Jahre brächte.

19. Aber man muß diesen Undank verbeißen, und mit Geduld überwinden und ausharren, auf daß nur der Gottesdienst rein bleibe. Wenn Moses nicht auf den Befehl Gottes gegangen hätte, so hätte er vorlängst vom Predigtamte und Regimente ablassen müssen. Und hätte ich
\quellseite{490}
mich auch an der Leute Haß, Feindschaft, Verfolgung und Undankbarkeit stoßen sollen, so hätte ich mein Theil auch längst hinweg, würde nicht viel mehr gelesen, geschrieben und gepredigt haben. Aber ich habe es nicht wollen thun. Denn wir sollen der Welt dienen und ihr Gutes thun, ob sie wohl dagegen das Gute mit Bösem vergilt. GOtt wird es wohl annehmen und reichlich vergelten, was die Welt nicht bezahlt und erstattet; er wird unsern getreuen Dienst ansehen. GOtt ist gütig, gnädig, langmüthig und von großer Geduld, wie im Propheten Hosea [Cap. 11, 8. 9.] auch gesagt wird: „Was soll ich aus dir machen, Ephraim? Soll ich dich zurichten wie Adama und Zeboim? Aber meine Liebe ist zu brünstig, daß ich nicht thue nach meinem Grimm“ etc. Er sieht durch die Finger, daß er eine Zeitlang solche Undankbarkeit vertragen kann; er wird es aber zu seiner Zeit ungerochen nicht lassen [2 Mos. 32, 34.]. Über Menschen können es nicht thun, die Undankbarkeit verbittert sie zu hart, darum gehört dazu eine göttliche Geduld.

\begin{schriftzitat}
V. 4–6. Moses schrie zum HErrn und sprach:
\end{schriftzitat}

Wie soll ich mit dem Volke thun? es fehlet nicht weit, sie werden mich noch steinigen. Der HErr sprach zu ihm: Gehe vorhin vor dem Volke, und nimm etliche Aeltesten von Israel mit dir, und nimm deinen Stab in deine Hand, damit du das Wasser schlagest, und gehe hin. Siehe, ich will daselbst stehen vor dir auf einem Fels, Horeb, da sollst du den Fels schlagen, so wird Wasser heraus laufen, daß das Volk trinke. Moses that also vor den Aeltesten von Israel.

20. Horeb ist gewesen ein dürre, rauh, hoch Gebirge, ein kahler Berg, da gar nichts gewachsen ist.

21. Alhier sind abermal zwei Stücke: eines, daß GOtt einen Befehl thut, und redet mit Mose, er solle mit etlichen Aeltesten in Israel hingehen an den Berg Horeb, und mit seinem Stabe an den Fels schlagen, so werde Wasser heraus laufen, daß das Volk trinke. Sonst hätte Moses wohl an alle Bäume im Walde schlagen mögen, so hätte er doch kein Wasser daraus gebracht. Darnach, so hat er das äußerliche Werk auch haben wollen.

22. Aber es ist ein rühmlich Mirakel, da wir mit starrenden Augen vorüber gehen, und nicht ansehen, auch so dicke Ohren haben, daß wir es nicht vernehmen, auf daß wir daran sehen könnten, wie gewaltig die Creaturen in GOttes Hand sind. Was kann unnatürlicher sein, denn ein Fels, ein harter Stein, daß darinnen Wasser solle sein? Wenn aus Schnee oder aus Wolken Wasser würde, und regnete, das wäre glaublicher; denn der Wolken Art ist, daß sie weich und wässerig sind; aber ein Fels, ein dürre, hart, trocken Ding, ein kahler Berg, daß der soll Wasser von sich geben, das hat eben das Ansehen, als wenn man spräche, daß aus Feuer solle Wasser kommen, und wiederum aus Wasser müsse Feuer werden. Droben [Cap. 16, 4.] haben wir dergleichen auch gehört, da in der Wüste kein Acker, kein Pflug noch Egge, weder Same, Korn, Weizen, Roggen noch Gerste war, davon die Kinder von Israel hätten essen mögen und ihre Nahrung haben: dennoch wuchs Brod in der Luft, es fiel täglich Manna, Himmelsbrod, in das Lager.

23. Wie reimt sich das? Die Luft gibt Speise und Nahrung, und alhier, die Steine oder Felsen fließen mit Wasser. Es ist eine wunderbare Gabe. Wie es denn auch seltsam und wunderberlich ist, daß Körner aus der Erde wachsen. Wer kann diese Kunst, und wer hat diese Gewalt? GOtt hat sie, der kann solche unnatürlichen Dinge thun, auf daß wir daraus uns einbilden mögen, was er für ein GOtt sei, und was er für Gewalt habe, auf daß wir an ihm nicht verzagten oder verzweifelten, sondern festiglich glaubten, und ihm vertrauten, daß er auch könne das Leder an der Tasche zu Golde machen, und aus Staub eitel Korn auf dem Boden machen, und die Luft mir zum Keller voll Weins machen. Das soll man ihm vertrauen, daß er eine solche große Gewalt habe, und wir wissen mögen, wir haben einen solchen GOtt, der diese Kunst könne, und daß es um ihn alles regne und schneie mit Wunderwerken.

24. Aber es hilft nicht. Ach GOtt! welch ein greulich, erschrecklich Gericht wird über die Welt gehen, daß sie sich nicht kehret an diese Mirakel! Lerne alhier erkennen, was die Welt heiße, wie ein teuflisch Ding es um die Welt sei: sie ist verstockt und verblendet, und sieht nicht GOttes Mirakel, noch was GOtt sei, der sich so reichlich ausgeschüttet mit Wohlthaten, und beweist mit Mirakeln, und mildiglich hilft jedermann.

\quellseite{491}
\begin{schriftzitat}
V. 7. Da hieß man den Ort Massa Meriba, um des Zanks willen der Kinder Israel, und daß sie den HErrn versucht, und gesagt hatten: Ist der HErr unter uns oder nicht?
\end{schriftzitat}

25. Sie wollen sagen: Haben wir einen GOtt, oder nicht? Pfui euch an, euch Lästerer, daß ihr so unverschämt GOtt versucht, und an GOtt verzweifelt. In Henkers\footnote{In der Eislebenschen: Hängers.} Namen, habt ihr nicht Manna gehabt? Oder seid ihr nicht aus Egypten gereist, und mit trockenen Füßen durch das Meer gegangen? Pfui uns auch an, wenn wir dergleichen sagen, so GOtt uns bald nicht gibt, wie wir ein Ding haben wollen! Sie stehen im Zweifel, ob GOtt unter ihnen sei oder nicht. Es ist eine solche Sünde gewesen, damit sie GOtt versucht haben, ob es wohl die heilige Schrift allhier verschweigt; und ist ein Wunder, daß es GOtt also lange gelitten hat, daß sie ihn immerdar versuchten.

26. Es ist eine grausame, große Sünde, also Gott mißbieten\footnote{mißbieten = ungebührlich sein.} und mißtrauen. Und läßt die ungläubige Rotte ihre Art und Tücke allhier redlich sehen; darum geben sie dem Orte auch davon einen Namen, auf daß man ja dieser Sünde des Volks nicht vergesse; heißen es Massa, das ist, tentatio; und Meriba, das ist, jurgium, ein Zank. Denn sie haben GOtt versucht, sie haben mit Mose oder mit GOtt gehadert und gezankt an diesem Orte. Zwar GOtt hat diese Sünde hart verdrossen, wie man das aus dem fünf und neunzigsten Psalm, V. 7–11., beweisen kann, da gesagt wird: „Heute, so ihr seine Stimme höret, so verstocket eure Herzen nicht, wie zu Meriba geschah, wie zu Massa in der Wüste, da mich eure Väter versuchten, fühlten und sahen mein Werk, daß ich vierzig Jahre Mühe hatte mit diesem Volke etc.; und sprach: Es sind Leute, deren Herz immerdar den Irrweg will, und die meine Wege nicht lernen wollen, daß ich schwur in meinem Zorn: Sie sollen nicht zu meiner Ruhe kommen.“ Da zeucht der Prophet diese Sünde an, daß dieser Zank, Hader, oder Versuchung, GOtt sonderlich verdrossen habe. Und wird solches auch von Mose im 5. Buch am 6. Capitel, V. 16., ihnen aufgerückt: „Versuchet GOtt euren HErrn nicht, wie ihr ihn versucht in Massa Meriba.“ Item, der 78. Psalm, V. 41. 56., und 106. Psalm,

\begin{schriftzitat}
V. 21. ff., gedenken sein auch; wiewohl nicht gemeldet wird in der Schrift, daß sie darüber eine Strafe erlitten hätten.
\end{schriftzitat}

27. Das ist nun das herrliche Mirakel, daß GOtt die Kinder Israel von dem Durst erlöst hat, und durch ein großes Wunder Wasser aus einem Steine oder Felsen verschafft, wie er ihnen zuvor Manna vom Himmel gegeben hat. Damit fürwahr uns ein tröstlich Exempel vor die Augen gestellt, daß GOtt uns erhalten und erquicken wolle, ob gleich auch die höchste und äußerste Noth uns übergehen sollte. Darum, so sollen wir ihm auch alle leibliche Noth vertrauen. Denn, ehe es mangeln sollte, so müßten ehe Steine zu Wasser werden, und Brod vom Himmel fallen, welches zuvor nie erfahren, und über alle Natur ist. Aber hiervon zu diesemmal genug. Lasset uns weiter zum Texte greifen.

\begin{schriftzitat}
V. 8–14. Da kam Amalek, und stritt wider Israel in Raphidim. Und Mose sprach zu Josua: Erwähle uns Männer, zeuch aus, und streite wider Amalek. Morgen will ich auf des Hügels Spitze stehen, und den Stab GOttes in meiner Hand haben. Und Josua that, wie Mose ihm sagte, daß er wider Amalek stritte. Mose aber, und Aaron, und Hur, gingen auf die Spitze des Hügels. Und dieweil Mose seine Hände empor hielt, siegte Israel; wenn er aber seine Hände niederließ, siegte Amalek. Aber die Hände Mosis waren schwer; darum nahmen sie einen Stein, und legten ihn\footnote{„ihn“ fehlt in der Eislebenschen.} unter ihn, daß er sich darauf setzte. Aaron aber und Hur unterhielten ihm seine Hände, auf jeglicher Seite einer. Also blieben seine Hände steif, bis die Sonne unterging. Und Josua dämpfte den Amalek und sein Volk durch des Schwerts Schärfe. Und der HErr sprach zu Mose: Schreibe das zum Gedächtnis in ein Buch, und befiehl es in die Ohren Josua, denn ich will den Amalek unter dem Himmel austilgen.
\end{schriftzitat}

28. Nun folgt das andere Wunderwerk, welches wahrlich in meinen Augen auch seltsam ist, daß GOtt die Amalekiter erschlägt, und die Israeliter errettet. Droben [Cap. 12, 37.] habt ihr angehört, daß die Kinder von Israel, sechsmal hunderttausend streitbare Mann, wohl geharnischt und gerüstet, aus dem Königreiche Egypten ausziehen, gar ein gewaltiges, großes
\quellseite{492}
Kriegsvolk. Hier aber werden diese Eisenfresser und Scharrhansen,\footnote{Scharrhansen = trotzige, hochfahrende Leute.} dies mächtige Volk, zu einem solchen losen Volke, daß sie sich nicht schützen können wider das geringe Volk der Amalekiter. Ich hätte gedacht, ein solch Kriegsvolk sollte im Kreuz,\footnote{In den Ausgaben: ein Kreuz. „im Kreuz“ = kreuzweise} durch die ganze Welt gereist sein. Aber ihre große Macht und Gewalt läßt GOtt so schändlich zu Schanden werden, und wird ein so groß Volk von wenigen Amalekitern übel geschlagen, welche Amalekiter nur ein Stück Landes inne hatten um das Gebirge Horeb.

29. Alhier geht eine neue Anfechtung und Trübsal daher, daß die Amalekiter zu Felde ziehen, und das Volk Israel, die fremden Gäste, schlagen wollen. Nun waren die Israeliter des grausamen Tyrannen und Feindes, des Pharaonis und der Egypter, nur vor wenig Tagen los worden, welche denn vor ihren Augen im rothen Meere ersoffen und umgebracht sind, daß sie meinten, sie wären jetzt gar über den Berg hinüber, es hätte nunmehr keine Noth oder Gefahr mit ihnen. Denn GOtt hat ihnen auch in der Wüste Brod vom Himmel, und Wasser aus den Steinfelsen gegeben. Aber es will noch nicht ablassen und aus sein; ja, ein neues Unglück plagt einher, nämlich dies, daß Amalek, der Feind, erstehet seinen Vortheil, und meint, das Volk Israel ist müde von der großen Reise, da sie aus Egypten gezogen waren, hatten auch in der Wüste Hunger und Durst erlitten, stecken noch in der Wüste und in dem Gebirge, da ihnen nichts zukommen kann, und sie zum Streite nicht wohl gefasset sind; darum gedenkt er, sie zu überfallen, zu schlagen und zu vertilgen.

30. Solches ist ein recht Bild und Spiegel dieses Christenlebens in dieser Welt, da immerdar eine Verfolgung, Kreuz und Noth auf die andere folgt; da heißt es: Nulla calamitas sola; und, wie man im Sprichwort sagt: Auf einen Sonnenschein folgt gemeiniglich gerne ein Platzregen. Es ist um dies zeitliche Leben eines Christen gethan, gleich wie um das Aprilwetter. Denn im April das Wetter nicht stetig ist, und nicht für und für die Sonne scheint, sondern jetzt ist der Himmel hell und klar, bald so regnet es, schneiet es, schlosset es, und hagelt es wieder drauf; flugs vergeht solches trübe und nasse Wetter wieder, und kommt darauf ein Sonnenschein. Also wechselt GOtt auch um mit den Christen; jetzt haben sie Glück, bald Unglück; jetzt ist Freude, bald Leid; jetzt Leben, bald Tod. Aber GOtt hilft immerdar in solcher Trübsal, und gibt darinnen ein Auskommen, daß man die Anfechtung ertragen könne [Ps. 68, 20.].

31. Da erhebt sich eine Frage von diesen Orten, da sie eingezogen sind: Wie Amalek habe können wider ganz Israel streiten (denn des Volks Israel viel mehr gewesen ist, denn jener, und daß Mosis mit dem Gebete sie allein geschlagen, und sie es nicht mit der Faust gethan haben), dieweil sie hernachmals die beiden viel gewaltigeren Könige, als Og und Sihon, nicht bestreiten konnten? Item: Warum Moses daselbst nicht auch gebetet? Nun, wenn es in der Schrift nicht stünde, so glaubte ich es selbst nicht; ich spräche, es wäre erlogen. Aber also ist es zugegangen: Josua erwählte etliche. Es ist nicht ein Krieg gewesen wider das ganze Volk; denn dies große Volk hat auch einen großen Raum mit seinem Lager eingenommen, sie sind weit von einander gelegen, und breit gezogen; da hat Amalek ein tausend oder zwei angegriffen von dem großen Haufen, und also sein Müthlein gekühlt.

32. Es ist der gemeine Brauch auch also zu reden, wenn ein Volk an einem Orte geschlagen wird; ob wohl das wenigste Theil die Schlappe einnimmt, so sagt man doch, das ganze Volk sei geschlagen. Gleichwie man noch sagt, wenn eine Stadt gestürmt, oder ein Volk geschlagen wird, so geht das Geschrei, der Fürst sei geschlagen und überwunden. Also ist in diesen Worten die Synedoche auch, daß Moses redet vom ganzen Haufen, ob gleich nur ein Theil gute Klappen bekommen hat. Gleichwie man auch sagt: Der hat einen Hund geschlagen, wenn er ihn kaum an ein Bein getroffen hat; denn niemand kann einen Hund gar an allen Orten schlagen. Also spricht man auch: Der hat einen Mann verwundet,\footnote{In der Eislebenschen: gewundet.} ob er wohl nur ihm ein Bein, Arm, oder sonst am Leibe verwundet hat. Item, man sagt: Der hat ein Haus gepocht,\footnote{pochen = zertrümmern.} ob er wohl nur den Ofen in der Stube eingeworfen, oder die Fenster ausgeschlagen, oder nur in die Thüre gehauen hat. Wohlan, das ist die gemeine Weise, zu reden von einem Stücke, da
\quellseite{493}
man doch das ganze Ding mit versteht. Also hat alhier Amalek nur einen Ort [des] Lagers, oder ein Theil vom Volke angegriffen; das lautet also, als hätte er das ganze Volk Israel überfallen.

33. Das Stücklein hat Moses geklagt. Nun wird er Josua befehlen, daß er wider Amalek ausziehen solle. Zwar GOtt selbst befiehlt es, daß sie sich rüsten sollen zum Streit, und also stellen, gleich als wollten sie Amalek mit dem Schwerte gar tilgen und schlagen; und sie sollten ihn doch nicht mit dem Schwerdt alle machen. Was gaukelt denn ihr Schwerdt und Rüstung, so es nichts thun soll, und GOtt heißt sie es gleichwohl? Denn drunten [V. 9.] wird Moses sagen: „Ich will droben auf dem Berge sein, und den Amalek schlagen.“ Josua, ob er wohl kriegete, schlug ihn nicht, sondern Moses auf dem Berge, mit seinem Gebet. Josua gehet von Leder; so hebt Moses zu GOtt seine Hände auf, mit einigen Anrufen. Also geht es beides neben einander.

34. Daraus wir auch lernen sollen, wie GOtt uns schützen und vertheidigen wolle, und solches auf zweierlei Weise: durch Mittel, und, wenn es die Noth erfordert, in andere Wege; gleichwie Josua des Schwerdts gebraucht wider den Amalek. Aber Moses, der gegen [den] Himmel die Hände aufhob, und betete, richtet mehr aus wider die Amalekiter, denn Josua mit der Schärfe des Schwerdts. Derhalben, so soll man mit nichten verzweifeln, ob gleich Mangel da ist; auch nicht stolz, sicher und aufgeblasen werden, noch trotzen, wenn irgend ein Vortheil vorhanden ist, sondern auf GOtt vertrauen. Es ist alhier ein Ueberfluß und großer Vorrath der Kriegsrüstung; aber aus ihren Kräften thun sie es nicht, und er muß ihnen nicht zu Nutze kommen, sondern Moses mit seinem Gebet, oder Händeaufheben, schlägt den Amalek.

35. Also wehrt die heilige Schrift überall, daß man nicht zur rechten oder zur linken Hand zu weit ausweiche, das ist, daß man nicht hoffärtig und vermessen sei, wenn etwas Gutes uns zusteht; auch nicht verzage, wenn nichts da ist, oder daß Mangel und Noth erscheint; sondern auf der Mittelbahn soll man bleiben. Denn wenn es da ist, so kann es GOtt wohl entziehen und wegnehmen; wiederum, wenn es nicht allda ist, so kann und weiß er es auch vom Himmel wohl zu verschaffen.

36. Dazu, so haben wir in diesem Texte ein herrlich Exempel eines ernsten Gebets, und was dasselbige für Kraft und Wirkung habe. Denn da Moses seine Hände aufhebt und betet neben Aaron und Hur, da siegen die Israeliter, und liegen ob wider die Amalekiter. Darum ist es wahr, wie die heilige Schrift sagt (Sir. 35, 21.): „Das Gebet der Elenden durchdringet die Wolken, und lässet nicht ab von dem Allerhöchsten, bis es erlanget, was es begehret.“ Item, das im 50. Psalm, V. 15., gesagt wird: „Rufe mich an in der Zeit der Noth, so will ich dich erretten, und du sollst mich preisen.“ Und das GOtt im Propheten Jesaia Cap. 65, 24. spricht: „Es wird sein, ehe sie schreien werden, will ich sie erhören“ etc. Denn wenn ein christlich Herz ernstlich zu GOtt betet, schreiet, und spricht: Hilf, lieber GOtt, nun hilf GOtt, laß dich das erbarmen im Himmel! seufzt, fleht, und hält an, so ist es unmöglich, daß ein solch Gebet von GOtt nicht sollte erhört werden. Es muß und soll alles da sein; derhalben schlägt solch Gebet diesen zornigen und grimmigen Feind, den Amalek, auch.

37. Nun wollen wir handeln die Allegorieen dieses 17. Capitels, und sagen, was GOtt habe anzeigen wollen. Damit aber, daß [das] Wasser aus einem steinernen Felsen entsprünge, solches stellen wir an seinen Ort, dieweil diese Allegorie fast gemein ist, und St. Paulus in der 1. Epistel zu den Corinthiern am 10. Capitel, V. 4., sie auch anrühret etlichermaßen, und spricht: „Unsere Väter haben daselbst einen geistlichen Trunk gethan von dem Fels in der Wüste, der ihnen nachgefolget ist,\footnote{Eißlebenische: sei.} welches war Christus.“ Und deutet daselbst St. Paulus den Fels auf Christum, und will, daß sie eben allda beiderlei haben getrunken, den geistlichen und leiblichen Trunk. Den leiblichen aus dem natürlichen Fels, damit der Leib erquicket und der Durst gestillet wurde. Aber diese Wohltat wäre ihnen nicht widerfahren, und das Wasser wäre ihnen nicht gegeben worden, hätten sie nicht den Glauben gehabt. Demewegen haben sie geistlich auch getrunken, das ist, sie haben geglaubt; wie wir denn auch noch an denselbigen Christum glauben. Allein, das ist zwischen uns und den Kindern von Israel der Unterschied, daß sie an Christum glaubten, der noch nicht gekommen
\quellseite{494}
war, sondern er kam hernach; wir aber glauben an den Christum, der gekommen ist.

38. So ist nun durch den Fels der HErr Christus geistlich bedeutet worden, darum, daß er der rechte Grund ist, darauf wir bauen und vertrauen sollen, ja, darauf wir uns auch stählen und verlassen mögen. Darum der HErr Christus auch sagt zu Petro [Matth. 16, 16. 18.], da er bekannte, Christus wäre der Sohn des lebendigen GOttes: „Du bist Petrus, und auf diesen Fels will ich meine Kirche bauen.“ Und St. Peter wird davon genannt, daß er heiße ein Felsen. Aber die Kirche ist auf Christum gegründet, das ist, sie glaubt an Christum, daß er sei unsere Gerechtigkeit, Stärke, Weisheit, Heiligkeit und Leben [1 Cor. 1, 30.]. Wer dies ihm vertrauet, der ist durch diesen Glauben und Vertrauen gegründet auf den Fels, und bleibt wohl unumgestoßen.

39. Daß nun das Wasser aus dem Felsen heraus läuft, das ist ein unglaublich Ding, und wird von der Vernunft nicht begriffen. Aber GOtt sprach zu Mose: „Schlage an den Fels, so wird Wasser herausfließen“ etc. Wenn nun Moses nicht hätte den Glauben gehabt, so wäre kein Wasser aus dem Steinfels gesprungen. Sollte Moses seiner Vernunft hierin gefolgt haben, so hätte er müssen sagen: GOtt wäre ein Narr, oder der Teufel hätte solches geheißen. Ja, wenn es ein Schneeballen wäre, der möchte Wasser geben; aber dies war zu hoch, daß Wasser aus dem Steine sollte fließen. Darum, der Glaube Mosis und der andern brachte das Wasser heraus, damit, daß sie dem Worte Gottes getraueten.

40. Die Bedeutung ist aber, daß noch aus dem gekreuzigten und verachteten Christo Wasser der Gnaden des Heiligen Geistes entspringt und herfließt; und dies ist denn auch unglaublich, daß da Wasser des Lebens und Heiligen Geistes sein solle. Wie denn dies Wasser also ausgelegt wird Joh. 4, 13. 14., da der HErr Christus zu dem samaritanischen Weibe sagt, wer des Wassers trinke, das er gebe, den werde\footnote{Gießlebensche: wird.} ewiglich nicht dürsten; sondern das Wasser, das er geben werde, das werde in ihm ein Brunn des Wassers werden, das in das ewige Leben quillet. Hier sehe man auch, wie sich das reimt, oder wie man dies fassen könne. Ein gekreuzigter Mensch, der da Fleisch und Blut hat, und läßt sich tödten (denn er ist Fleisch, und stirbt), aus deß Tode, aus seinem Fleische und Blute soll ein Geist, der lebendig mache, herauskommen? Fleisch und Geist, Tod und Leben, wie reimen sie sich doch mit einander? Wer will doch alhier den Heiligen Geist heraus saugen? Das müßte ein geschickter Meister sein, der diese Kunst könnte.

41. Darum dünkt es die Welt überaus närrisch sein, daß aus dem verdammten und erhenkten Menschen Christo sollte der Heilige Geist kommen, und lebendig Wasser quellen, welches erquicke und selig mache an Leib und Seele. Das ist übernatürlich, und keines Menschen Verstande auf Erden begreiflich, er sei gleich, wer er wolle. Wer es nun nicht glaubt, der hat es nicht; es gehört Glaube hierzu, daß du es gewiß dafür hältst, daß Christus der lebendige Brunnquell und himmlische Fels sei, so für uns gestorben, und sein Leben für uns gelassen habe, und also durchgraben und eröffnet; wie auch der heilige Johannes in der Geschichte seines Leidens daher sieht Cap. 19. 34.; alles darum, daß ich durch ihn den Heiligen Geist überkäme.

42. Daß aber Moses mit dem Stabe an den Fels schlägt, bedeutet, daß, ob wohl hundertmal und aber hundertmal Christus gekreuzigt wäre, so käme doch der Heilige Geist nicht, wenn es nicht in das Predigtamt gefasset wäre. Da muß man pochen und klopfen an den Fels, das ist, man muß von ihm predigen, man muß die Ruthe des Mundes haben [Jes. 11, 4.]. Das Predigtamt ist Mosis Ruthe und Stab; wer den Leuten aus Vorhaltung des Gesetzes ihre Gewissen pocht und schlägt, und darauf predigt, daß Christus für uns gestorben sei, der schlägt auf den Fels, und aus derselbigen Predigt, oder aus diesem Schlagen und Klopfen, da kommt der Heilige Geist. Siehe, ein solch gering, verächtlich Ding ist es, es ist nicht dafür anzusehen. Ein Fels und Stab ist ein dürre Ding; so ist es auch um das mündliche und leibliche Wort ein gering Ding, und soll doch daraus fließen der Heilige Geist in die Herzen der Zuhörer, daß sie davon trinken, das ist, einen lebendigen Geist empfangen.

43. Aber die Vernunft spricht: Wie kann der Heilige Geist aus dem mündlichen Worte empfangen werden, ist er doch unmeßlich? Ja, er reinigt uns von Sünden, errettet uns vom Tode, und bringt uns in das Leben, macht uns ewig
\quellseite{495}
selig, segnet uns mit Gütern zeitlich und ewiglich. Das scheint ein närrisch Ding sein, daß durch eines Menschen Wort solches geschehen solle. Gleichwie es auch ein seltsam und närrisch Ansehen gehabt, daß Moses an den Fels geschlagen. Aber Moses hat geistlich anzeigen wollen, daß man mit dem leiblichen Predigtamte müsse auf diesen Fels, und nicht sonst auf die Erde schlagen, das ist, man soll nichts predigen, noch auf nichts schlagen oder treffen, denn allein auf Christum. Ich soll mit meiner Zunge immerdar schlagen und dreischen auf diesen Fels, dann fließt bald Wasser heraus, das ist, dann wird mir und andern der Heilige Geist gegeben, ja, er kommt zu dir auch, auf daß wir erneuert und wieder lebendig gemacht werden.

44. Aber dies geschieht auf dem Berge Horeb, welcher gar ein dürrer Ort ist, es liegt gar ein dürre Land umher, und ist ein ganz groß Gebirge gewesen, gleichwie im Böhmer- und Thüringerwalde ist. Dies Gebirge hat viel Stücke oder Theile gehabt, und darum auch viel Namen bekommen, daß es hat geheißen der Berg Horeb, item, der Berg Sinai; hat in einer dürren Einöde oder Wüste gelegen. Dieses bedeutet, das Evangelium werde gepredigt, und Wasser aus einem harten Felsen gegeben, anderswo nirgends, denn wo dürre, trockene Herzen sind, das ist, elende, betrübte Sünder. Wie ihr denn zuvor oft gehört habt, daß das Evangelium nirgends reichen Trost predigt, oder zu Herzen gezogen wird, denn bei den dürftigen, hungrigen und dürftigen Seelen. Wie denn das Magnificat [Luc. 1, 53.] auch also singt: „Die Hungrigen hat er mit Gütern erfüllet.“ Und sonst wird von dem HErrn Christo gesagt: Das Evangelium wird den Armen gepredigt [Matth. 11, 5. Luc. 7, 22.].

45. Unseres HErrn GOttes Art und Weise ist diese, daß er zu schaffen hat und wirkt allein da, wo nichts ist. Denn mit dem, so hoch und groß ist, handelt GOtt anders nicht, denn daß er es zerbreche. Wiederum geht er mit den Armen also um, daß er sie tröste, und daß er die Todten lebendig mache, und predigt durch das Evangelium ihnen Vergebung der Sünde, macht die Seelen und Gewissen fröhlich und muthig. Denn wo die Seele sündlich, und das Gewissen erschrocken ist, da bekommt es am Berge Horeb wieder Kraft und Stärke. Der Prophet Elias hat an diesem Berge Horeb auch gewohnt [1 Kön. 19, 8.], auf daß wir sehen, daß die Propheten und das Gesetz mit diesem Berge viel zu thun haben; das ist, das Gesetz treibt die Menschen dahin, daß sie ihre Sünde erkennen. Denn so kommt Moses, und schmeißt mit einem Stabe auf den Fels, da geht Wasser heraus; das ist, da wird ihnen der Heilige Geist gegeben, daß sie wieder erquickt und gelabt werden.

46. Das ist die Versuchung oder der Zank zu Massa Meriba. Denn in Horeb geht es also zu: Wenn das Gesetz gepredigt wird, so demüthiget es die Gewissen, es offenbart die Sünde, da bleibt denn Hader und Zank nicht außen. Denn es fallen etliche Herzen zu, die mit GOtt durch ihre Werke handeln wollen, und streben wider das Evangelium, nehmen es nicht an, als denn die Juden thun, Apost. 13, 46., andere aber, die nehmen es mit Glauben an. Dieser Kampf und Hader bleibt wohl, daß diese wollen Werke haben, jene aber den Glauben; diese wollen das Evangelium nicht leiden, andere, die wollen es haben. Und währt solcher Zank auch noch auf den heutigen Tag, welcher zu Massa Meriba sich zuträgt. Wie denn auch im 95. Psalm, V. 7. 8., gesagt wird, da David von dem neuen Predigtamte redet: „Heute, wenn ihr werdet hören meine Stimme, so verhärtet eure Herzen nicht, wie ihr thatet zu Meriba“, sondern nehmet ohne Zank die Stimme an, und lasset eure guten Werke immer hinfallen.

47. Nun wollen wir auch hören von Amalek, was er bedeutet. Josua ward geboten, er sollte streitbare Männer auserlesen, und wider Amalek kriegen; aber Moses blieb auf dem Berge, hebt seine Hände auf, und betet. Und dieweil er die Hände aufhob, so lag Israel oben; wenn er aber die Arme sinken ließ, so verlor Israel. Denn Mosis Hände wurden schwer,\footnote{Eißlebenische: schwerer.} darum mußte Moses niedersitzen, und wurden Steine ihm unter die Arme gelegt. Es standen ihm aber zween auf beiden Seiten, so die Hände Mosis empor hielten.

48. Allhier sind nun zwo Historien. Das Volk Amalek nennt die heilige Schrift die Erstlinge unter den Heiden, wie Bileam im vierten Buche Mosis [24, 20.] sagt: Primitiae gentium Amalek, das ist, da Israel aus Egypten zog, da waren die Amalekiter die ersten unter den
\quellseite{496}
Heiden, da sie durchziehen mußten, und die sich auch wider die Israeliter legten. Darum ward geschrieben, 5 Mos. 25, 17–19., daß man sie vertilgen sollte, denn sie hatten sich am ersten wider Israel gesetzt; wie denn allda gesagt wird: „Gedenke, was dir die Amalekiter thaten auf dem Wege, da ihr aus Egypten zoget, wie sie dich angriffen auf dem Wege, und schlugen deine Hintersten, alle die Schwachen, die da hinten nachzogen, da du müde und matt warest, und fürchteten GOtt nicht. Wenn nun der HErr, dein GOtt, dich zur Ruhe bringt, von allen deinen Feinden umher, im Lande, das dir der HErr, dein GOtt, gibt zum Erbe einzunehmen, so sollst du das Gedächtniß der Amalekiter austilgen unter dem Himmel; das vergiß nicht.“ Und dies ward durch Saul und Samuel ausgerichtet, welche diese Feinde vertrieben und vertilgten.

49. Dies Volk Amalek bedeutet die Juden, oder das Volk Israel. Denn als der HErr Christus auf Erden kam, legten sie sich mit aller Macht und Gewalt wider Christum. Dies Volk ist der rechte Amalek worden. Denn da das Evangelium durch Johannen den Täufer, durch den HErrn Christum selbst und die lieben Apostel gepredigt ward, waren sie die allerersten, die sich wider das göttliche Wort legten, dasselbige mit großem Ernste verfolgten, und viel Christen tödteten. Und sie achteten auch noch nicht das göttliche Wort; denn das ist ihr Amt, wider Christum mit dem Schwert und mit der Zunge streiten. Es ist ein geistlicher Krieg und Streit. Aber es geht den Juden, wie von dem Amalek geschrieben steht, daß ein ewiger Krieg zwischen Amalek und den Kindern Israel gewesen sei, bis die Amalekiter sind vertilgt worden. So hat GOtt auch seine Strafen für und für über die Amalekiter ergehen lassen, und sie bis in das dritte und vierte Geschlecht gestraft.

50. Durch die Apostel und ihre Predigt ward damals den jüdischen Amalekitern solches zur Warnung treulich genugsam angezeigt; dieweil es aber nichts half, kam ihnen der Glaube in die Hand mit den Römern, durch welche das jüdische Volk gedämpfet ist, und wird auch noch heutzutage durch das Evangelium verstöret und ausgetilgt, daß sie nichts mehr thun können. Und ob sie, die Juden, wohl einzeln in die Länder sich hin und wieder verstecken, so können sie doch nicht wieder zu Kräften kommen, daß sie ein eigen Volk würden, daß sie einen eigenen König und Haupt hätten. Gleichwie auch Amalek nie zu einem eigenen Königreiche wieder gemacht ist; sondern, da es erstlich so ein großes, herrliches Reich war, und, mit allerlei Vorrath versehen, bei einander stand, wie ein gewaltiger, großer See, oder ein tieffließend Wasser, so ward es darnach so gering und klein, als ein Bach, da man mit den Füßen hindurch waten möchte, wie also Jesaias der Prophet [Cap. 11, 15.] davon redet.

51. Josua erwählt aber Männer, so wider Amalek streiten sollen; das ist, GOtt erweckt Prediger und Apostel, die sich wider die Juden legen. Und diese fechten und kämpfen mit dem Schwerte des Geistes, das ist, mit dem Evangelio.

52. Aber damit richtet man wenig aus, wenn nicht unser Moses, das ist, der HErr Christus, droben auf der Spitze des Berges wäre, und unser Fürst und Herr sei. Denn wir reden nur das Wort, und predigen es; aber wo er nicht das Herz rührt, und den Heiligen Geist gibt, so ist es geschehen. Das wird in dem angezeigt, daß Moses mit seinem Händenausheben mehr streitet, denn Josua mit seinem Schwerte; das ist, wir predigen, aber er gibt das Gedeihen und den Nachdruck dazu, wie solches St. Paulus [1 Cor. 3, 6.] auch bezeugt, da er spricht: „Paulus hat gepflanzet, und Apollo begossen, aber GOtt hat das Gedeihen dazu gegeben.“ Also ist Moses allhier ein Bild und Figur des HErrn Christi, denn auf Christum soll man auch alle Figuren ziehen.

53. Daß aber Moses schwere Hände hat c., bedeutet Mosis Predigtamt, oder die Lehre des Gesetzes. Denn wenn Mosis Amt im Schwange geht, und das Gesetz getrieben wird, dann sind seine Hände schwer, das ist, die Werke des Gesetzes sind schwer zu thun. Aber also haben wir Mosen im neuen Testamente nicht, als die Juden ihn gehabt; wir haben nur den halben Mosen, er thut bei uns Christen mehr nicht, denn daß wir anfänglich durch ihn zur Erkenntniß der Sünden gebracht werden; wie denn Paulus [Röm. 3, 20.] sagt, es werde die Sünde durch das Gesetz erkannt. Darnach haben wir Christum, der den innerlichen Menschen regiert durch den Heiligen Geist; aber auswendig sind seine Hände noch schwer. Denn er hat noch
\quellseite{497}
viel schwacher Christen, die nicht hernach können mit den Werken; wie denn St. Paulus zum Römern [Cap. 7, 18. 19.] auch darüber klagt; der alte Schalk will nicht hernach. Nach dem innerlichen Gesetze ist es alles erfüllt, aber nach dem alten Adam sind die Hände noch schwer, es will nicht hernach, da ist das Gesetz eine schwere Last und Bürde, so da drückt und beschweret. „Denn das Fleisch gelüstet wider den Geist“, sagt St. Paulus (Gal. 5, 17.). Und der HErr Christus (Matth. 26, 41.) spricht: „Der Geist ist willig, aber das Fleisch ist schwach.“ Das Fleisch zeucht uns immerdar zurück vom Geiste.

54. Aber die zween Diener, so Moses Hände empor heben, bedeuten auch das Predigtamt in beiden Testamenten. Man muß wehren zur rechten Hand und zur linken Hand mit Vermahnen, Anhalten, Strafen, Träuen, Verheißen und Trösten. Man muß widerhalten, daß man nicht geil, faul und laß werde, wie denn unser Fleisch und Blut dazu seine Lust hätte. Also will das Wort getrieben sein, auf daß wir nach dem Fleisch und dem Geist keck seien. Wenn der Geist frisch regiert, so hat es mit dem Fleische nicht Noth, wenn wir allein das Wort getrost treiben. Mit dem Predigtamte des göttlichen Worts muß man das Fleisch im Zaume behalten.

\begin{schriftzitat}
V. 15. Und Mose bauet einen Altar, und hieß ihn: Der HErr Nissi (das ist, ein Heerzeichen, ein Panier).
\end{schriftzitat}

55. Dieser Altar ist Jesus Christus, welcher ist unser Fähnlein oder Panier [Jes. 5, 26.], auf welchem Altar wir opfern, das ist, GOtt danken und loben. Denn wir haben das Evangelium und den Heiligen Geist, und erkennen Christum, und thun auch christliche Werke allhier im Fleische. Das ist unser Zeichen, daß wir in Christo und durch Christum GOtt danken für sein heiliges Wort. Sonst haben die Prediger allhier eitel Moses gemachet. Nun ist es gut, daß eines Theiles Moses sei nach dem äußerlichen Menschen, aber inwendig da soll Moses nichts zu thun haben, sondern Christus allein regieren.

\section{Das achtzehnte Capitel}

\begin{schriftzitat}
V. 1–8. Und da Jethro, der Priester in Midian, Moses Schwäher, hörte alles, was GOtt gethan hatte mit Mose und seinem Volke Israel, daß der HErr Israel hätte aus Egypten geführet, nahm er Zipora, Moses Weib, die er hatte zurück gesandt, sammt ihren zween Söhnen. Der einer hieß Gerson, denn er sprach: Ich bin ein Gast worden in fremden Landen; und der andere Elieser, denn er sprach: GOtt meines Vaters ist meine Hülfe gewesen, und hat mich errettet von dem Schwerte Pharao. Da nun Jethro, Moses Schwäher, und seine Söhne und sein Weib zu ihm kam in die Wüste, an den Berg GOttes, da er das Gezelt aufgeschlagen hatte, ließ er Mose sagen: Ich, Jethro, dein Schwäher, bin zu dir gekommen, und dein Weib und ihre beiden Söhne mit ihr. Da ging ihm Mose entgegen hinaus, und neigte sich vor ihm, und küssete ihn. Und da sie sich unter einander gegrüßet hatten, gingen sie in die Hütte. Da erzählte Mose seinem Schwäher alles, was der HErr Pharao und den Egyptern gethan hatte, Israels halben, und alle die Mühe, die ihnen auf dem Wege begegnet war, und daß sie der HErr errettet hätte.
\end{schriftzitat}

1. Ehe denn dies geistliche Regiment bestellet wird, so ordnet Moses das weltliche Regiment; denn man sonst das Evangelium nicht wohl predigen kann, es sei denn zuvor in einem Lande stille, und werde guter Friede gemacht; darum setzt Moses dieses Capitel vorher. Denn wo es nicht Friede oder stille ist, da soll das Evangelium wohl ungepredigt bleiben.

2. So beschreibt er nun in diesem Capitel, wie sein Schwäher Jethro, der Priester in Midian, zu ihm in die Wüste kommt, und bringt Ziporam, Moses Weib, und seine zween Söhne mit. Denn droben im 2. Capitel, V. 15. 16., haben wir gehört, wie Moses aus Egypten entwich vor Pharaone, und verhielt sich bei dem Priester in Midian, hütete allda der Schafe, bis daß er ihm seine Tochter zum Weibe gab, mit welcher er zween Söhne zeugte. Und als er wieder nach Egypten zog, sandte er das Weib mit den Kindern
\quellseite{498}
zurück zu ihrem Vater, der droben [2 Mos. 2, 18.] Reguel genannt wird, allhier aber wird er Jethro geheißen. Vielleicht ist der Jethro der Frau Stiefvater gewesen, und wird Reguel, der rechte Vater, nun gestorben sein. Da nun dieser Jethro hört, daß die Kinder von Israel aus Egypten nach der Wüste\footnote{Die Worte: „nach der Wüste“ fehlen in der Erlanger.} gezogen waren, so kommt er zu seiner Tochter Mann, Mosen, und bringt ihm das Weib und seine Kinder.

3. Aber was geht dieses uns an? Warum ist es geschrieben, daß man Mosi sein Weib und zween Söhne bringt, und Jethro darnach wahrhaftig bekehrt wird, den wahren GOtt erkennt, und opfert hernach demselben? Wir wollen daraus drei Stücke oder drei Lehren nehmen. Erstlich: Moses ist in einem Amte gewesen, von GOtt dazu berufen und erfordert, daß er das Reich Egypten und den König Pharao plagen, und das Volk Israel ausführen sollte aus dem Diensthause Egypti. Derwegen verläßt er darüber sein Haus und Hof, Weib und Kind in Midian, befiehlt sie unserm HErrn GOtt; der wartet allhier seines Hauses, seiner Habe und Güterlein, Weibes und Kindes. Derselbige versorgt solches alles viel besser, denn wenn Moses selbst wäre gegenwärtig gewesen. Aber dieweil Moses GOtt ist gehorsam gewesen, und hat dem Beruf mit Fleiß nachgesetzt, so Gott ihm befohlen, derhalben so thut GOtt ihm wieder die Ehre, daß er sein Weib und Kind nothdürftiglich versorgt, und zu gelegener Zeit dem Mosi herrlicher zuführen läßt, denn vielleicht sonst geschehen wäre.

4. Daraus wir denn diese Lehre nehmen sollen, daß, wenn wir im Gehorsam Gottes daher gehen, und wir um seinetwillen Haus und Hof und andere Güter lassen müssen, daß wir darum nicht trauern, oder uns zu hart bekümmern. Es soll mit den Gütern viel besser zugehen denn zuvor; man hat sie alsdann dem rechten Hausvater befohlen. Zu jenemmal, als Moses sein Weib und Kind aus Midian mit in Egypten nehmen wollte, da ward er todtkrank; darum mußte er sie wieder zurückschicken. Wie er sie nun nicht bei ihm hat, geht es ihm viel glückseliger, denn da sie um und neben ihm waren. Also will GOtt uns lehren, daß wir ihm vertrauen sollen, und alles um seinetwillen verlassen; denn er kann es wohl erhalten, bewahren, und viel reichlicher wieder geben, denn wir wünschen oder begehren möchten. Daher wird vom HErrn Christo auch im Evangelio [Matth. 19, 29.] gesagt: Wer da verläßt Acker, Haus, Hof, Weib und Kind, um meines Namens willen, der soll es hundertfältig wieder haben. Wären sie sonst mit einander fortgezogen, so wären vielleicht sie oder er in Egypten gestorben. Aber nun sind sie alle am Leben geblieben, kommen frisch und gesund wieder zusammen.

5. Zum andern, daß dies große Mirakel und Werk, als die Erlösung der Kinder Israel aus Egypten, nicht unfruchtbar bliebe, so gibt er allhier zu verstehen, daß solches Auszugs auch die Heiden genossen haben, auf daß sich die Juden nicht rühmten, daß sie allein Gott kennen, sondern die Heiden kennen ihn nun auch; GOtt hat unter ihnen auch seine Heiligen und Auserwählten. Er will nicht allein der Juden, sondern auch der Heiden GOtt sein [Röm. 3, 29.], ja, aller Creaturen HErr und GOtt. Derhalben, da er die Kinder von Israel ausführte aus dem Diensthause Egypti, und ihrer viel murrten und sich nicht besserten, da fährt GOtt zu, und nimmt auch die Heiden an zu seinem Volke. Denn dieser Heide, Jethro, bessert sich und bekehrt sich zu Gott, und freut sich über dem Erkenntniß des wahrhaftigen GOttes, und über der herrlichen Wohlthat, so GOtt an den Kindern von Israel bewiesen hatte, daß er sie mit gewaltiger Hand aus dem Diensthause Egypti los gemacht hatte. Und wiewohl sich dieses für ein schlecht und gering Ding läßt ansehen, und scheint, als sei Jethro allein in die Wüste zu Mose gekommen, so hat er doch genießlich viel Leute mit sich gebracht, denn er\footnote{Erlanger: es.} ist ein Regent, ja, das Haupt in Midian gewesen.

6. Die Midianiter sind von Abraham hergekommen, von der Ketura, Abrahams Weibe, wie im ersten Buch Mose am 25. Capitel, V. 2., geschrieben steht. Damit hat GOtt anzeigen wollen, daß das Volk Israel von andern Menschen abgesondert sei mit äußerlichen Satzungen, allhier auf Erden, um dieser Ursache willen: denn Christus hat von diesem Volke seine Menschheit an sich nehmen sollen. Aber nichtsdestoweniger hat er ihm auch die Heiden vorbehalten, welchen er inwendig den Glauben und den Heiligen Geist gegeben, ob er gleich nicht hat eine
\quellseite{499}
äußerliche Sonderung mit ihnen gemacht, als mit den Juden, und Israels Samen. Sie haben auch nicht gehabt Mosen und die Propheten, er hat sie nicht abgetheilt von andern, wie die Juden; doch inwendig hält er sie so gut als die Juden, er gibt ihnen eben denselbigen Glauben, Erkenntniß GOttes und den Heiligen Geist. Hat also GOtt seine Kirche und große Zahl der Gläubigen auch unter den Heiden; wie solches das Exempel Melchisedecks, Hiobs, Naaman Syri, der Niniviten und anderer mehr ausweisen.

7. So ist er ein Priester und Regent gewesen, wie die alten Väter auch mit regiert haben, als jetzt die Könige und Fürsten thun, hat gleichwohl auch mitzupredigt, und das Volk gelehrt als ein Priester. Denn das weltliche Reich ist nicht also vom Predigtamte damals abgesondert gewesen, als es jetzt ist. Und dieweil er nun in dem heiligen Predigtamte gewesen ist, und geglaubt, so hat er ihm diesen Schatz und Erkenntniß GOttes allein nicht behalten, sondern er sollte es anlegen, und damit wuchern, wie der HErr Christus Matthäi am 25. Capitel, V. 14. ff., bezeuget. Darum wird ohne Zweifel durch ihn das Land Midian sein bekehrt worden, und zum Erkenntniß des Worts GOttes gekommen; wie die Kinder von Israel dergleichen auch gehabt haben. Daß also GOttes Erkenntniß und der Glaube weit ist ausgebreitet worden, und nicht allein bei den Juden geblieben, sondern die Heiden haben es auch mit genossen. Denn GOttes Werke haben sollen in aller Welt bekannt werden, auf daß GOtt nicht allein von den Juden, sondern von allen Menschen in der Welt gepreist und gelobt würde; wie denn folgt im Text.

\begin{schriftzitat}
V. 10. 11. Und Jethro sprach: Gelobet sei der HErr, der euch errettet hat von der Egypter und Pharao Hand, der weiß sein Volk von der Egypter Hand zu erretten. Nun weiß ich, daß der HErr größer ist, denn alle Götter; darum, daß sie Hochmuth an ihnen geübt haben.
\end{schriftzitat}

8. Jethro lobt GOtt, und bekennt ihn vor jedermänniglich, da thut er seine Bekenntniß. Das ist das rechte Opfer, das er GOtt leisten kann, daß man nach dem Erkenntniß GOttes sich freut über dem, daß GOtt ein solcher HErr sei, der Israel errettete vom Pharaone; darum sei er auch allein der rechte GOtt, vor allen andern, der den großen Bäumen steuern kann, daß sie nicht in Himmel wachsen. So sind die Egypter sicher und vermessen gewesen, haben Hochmuth an den armen Israeliten geübt. Wer also GOtt erkennt, der hat ihn wohl und recht erkannt; und also will er gelobt und erkannt sein, auf daß man wisse, er sei ein solcher GOtt, der allen helfen könne, die da in Nöthen stecken, und seiner Hülfe begehren; der auch demüthige alle hoffärtigen und hohen Leute [1 Petr. 5, 5.].

9. Nun, das rühmt er von GOtt, wie er über alle Götter sei, und keinesweges die Hochtrabenden und Uebermüthigen neben sich dulden könne; darum er auch mit den Egyptern nicht eins war. Denn die Egypter waren hoffärtig, und gar zu muthig, und drückten das Volk Israel unter sich; daraus hilft ihnen niemand, denn allein dieser GOtt. Und mit diesem GOtt wollen wir es auch halten, der ein solcher GOtt ist, daß er kann helfen und aus aller Noth erretten, und demüthige die Hoffärtigen. Aber wo sind die, die beides glauben? Man verachtet beide Stücke. Denn, die da hoffärtig sind, Uebermuth, Frevel und Stolz treiben, die halten es nicht dafür, daß es GOtt sehe oder rächen werde; sie wollen wissen, daß GOtt gnädig ist, und mit der Strafe verziehe, damit sie Raum haben zu ihrer Bosheit, und daß ihnen oft ihre Bubenstücke eine Zeitlang hingehen [Ps. 64, 6.]. Denn GOtt sieht dazu durch die Finger, das macht sie so verwegen; darum glauben sie nicht, daß ein GOtt sei. Denn sie führen ihr gottlos Wesen also hinaus, wie sie es nur im Sinne haben [Ps. 14, 1. 53, 2.].

10. Was geschieht nun auf dieser Seite bei uns, die wir glauben und leiden sollen? Hier will es auch in uns nicht, daß GOtt könne aus großen Herren und Hansen kleine oder gar keine, und aus reichen Leuten arme und Bettler machen, oder aus Klugen, Weisen und Hochgelehrten Narren machen, oder Heilige zu Schanden setzen werde. Ja, wir glauben auch nicht, daß GOtt auf uns ein Auge habe, und nach uns etwas frage. Wenn mir es fehlt an der Nahrung, oder stehe in Gefahr des Todes, komme etwa in Schande und Schmach, von Stund an denke ich, es sei gar mit mir aus, und niemand mehr vorhanden, der da helfen könne. Oder, wenn es sonst kommt, daß das Herz betrübt ist, und man in Unglück steckt, da glaubt es niemand, daß GOtt bei uns sei; so es doch GOtt von sich läßt predigen und sagen, daß er gegenwärtig bei uns in Nöthen sei; als Ps. 91, 15.: „Ich
\quellseite{500}
bin bei ihm in der Noth, ich will ihm heraußer helfen und erretten.“ Er beweist es auch redlich, als, an dem Pharaone, den er kann in das rothe Meer stürzen, und die Kinder von Israel aus seiner Tyrannei ledig machen. Wer wollte sonst diesem Volke ausgeholfen haben? Es war je niemand, der sich ihrer annahm, oder gedachte sie zu retten.

11. Darum gedenke du auch also: Es gehe mir, wie es wolle, so ist doch GOtt daheim. Wenn ich nun GOtt kenne, und weiß, daß ein GOtt sei, dann denke und sei gewiß, er schaue auf dich; denn das ist seine Art, daß er auf uns sieht, nach uns fragt, und Acht darauf hat, und ist bedacht darauf, daß er uns helfen möge. Denn er ist ein solcher GOtt, der dieses von sich läßt singen, sagen und schreiben, daß er erhöhe die Untergedrückten und Niedrigen, und wiederum stürze die Gewaltigen von dem Stuhl, oder drücke die Hoffärtigen zu Boden [Luc. 1, 51. 52.]. Aber dieweil wir nicht glauben, darum widerfährt uns auch keine Rettung; aber die es glauben, denen hilft er getreulich, ob es sich gleich eine Weile verzieht, doch kommt er zuletzt [Hab. 2, 3.], und läßt es nicht fehlen, wer allein ihm auszuhalten, ausharren könnte, wie die heilige Schrift sehr fleißig hierzu vermahnt.

12. Ein solch Licht hat er im Herzen gehabt, daß er also mit Freuden singt; als sollte er sagen: Ich lobe GOtt, der den Elenden hilft, und stürzt alle Hohen und Stolzen. Dies Liedlein oder Gesang singe ihm nach, wer da will sich als einen Christen erzeigen.

13. Nun fällt eine Frage vor: Ob mehr Götter sind denn Einer? Droben [Cap. 12, 12.] habt ihr gehört, da die Erstgeburten alle erschlagen worden sind, habe GOtt ein Gericht gehalten an den Göttern in Egyptenlande. Es sind aber die Egypter ein närrisch Volk gewesen, die voller Götter gesteckt haben; wiewohl es auch das klügste Volk unter der Sonne gewesen ist, dafür ich es gänzlich halte. Denn Joseph und die Kinder Israel brachten ein groß Licht hinein, denn in Israel war allein der rechte GOtt bekannt [Ps. 76, 3.], und denselbigen haben sie auch gepredigt und bekannt gemacht. Wie denn der 105. Psalm, V. 22., spricht: Joseph habe die Weisen in Egypten gelehrt, und habe sie von GOtt unterrichtet.

14. Aber hernach sind sie von diesem Erkenntniß GOttes und dem hellen Lichte wieder abgefallen, und viel blinder worden, denn irgend andere; wie es denn pflegt also zuzugehen, daß nach großem, hellem Lichte auch pflegen große Finsternisse zu folgen. Darum haben sie in Egypten hernach Kühe und Ochsen, Kälber, Hunde, Katzen und Vögel, dazu Sonne und Mond angebetet. Und in den folgenden Capiteln wird man hören, wie die Kinder von Israel solches von den Egyptern gelernt, und hernach auch ein gülden Kalb aufgerichtet haben, und darum getanzt. Aber was sind doch dieselbigen Götter? Uns geht es auch also. Wenn wir auf einer Seite des rechten GOttes gefehlt haben, dann haben wir seiner gar gefehlt, und wer ihn nicht ganz ergreift, der hat nichts von ihm; allein, daß wir nicht solche abgöttischen Bilder aufrichten, wie sie thaten.

15. Wie, wenn ich aber glaube, daß GOtt durch die Messen versöhnt und zufrieden gestellt werde, und nehme nicht an den Mann Jesum Christum, der von der Jungfrau Maria geboren ist, und für mich gelitten hat und gestorben ist, ist das nicht GOttes gefehlt? Ich nenne wohl mit dem Munde den rechten Namen GOttes und des HErrn Christi, aber im Herzen habe ich seiner vergessen. Denn meine Gedanken sind nicht rechtschaffen in meinem Herzen, und ich nehme meinen eigenen Dünkel und folge ihm, davon doch die heilige Schrift nichts lehrt, ja, vielmehr auf das heftigste wehrt und verbeut; und das nicht bei Halsabhauen (wie man pflegt zu reden), sondern bei dem höchsten Verlust, nämlich der Seelen Seligkeit; und damit habe ich den göttlichen Namen gelästert und geschändet, und mir einen andern Gott gemacht. Denn es soll dem HErrn Christo gefallen, was mir geliebet, da mir doch vielmehr gefallen sollte, was Christo geliebet; denn er soll mich meistern, und ich oder andere ihn nicht. Aber also gibt mir der Teufel den Dünkel ein, und übervortheilt mich, daß er mir den Namen Christi und GOttes daran schmiert; so wird man betrogen. Dagegen lehrt mich die heilige Schrift, daß ihm dies gefalle, und solches GOttes ewiger Wille sei, an Christum glauben, und auf sein Leiden, Sterben und Blutvergießen sich verlassen, wie Johannis im 6., V. 29., gesagt wird: „Das ist GOttes Werk, daß ihr an den glaubet, den er gesandt hat“, und ja nicht irgend ein Werk aufrichten, darauf ich meinen Trost setzte; alsdann so treffe ich ihn.

\quellseite{501}
16. Und also ist die ganze Welt zu unsern Zeiten in lauter Abgötterei gewesen. Darein ist sie also geraten und gefallen, daß man erstlich ein wenig gewichen ist und gestrauchelt hat, darnach sind sie je mehr und mehr in größere Irrthümer gediehen, bis daß alle Lande voller Abgötter worden, wie der Prophet Zephanja am ersten Capitel, V. 4. f., sagt.

17. Also ging es auch Jerobeam, da er zwei Kälber anrichtete, eines zu Dan, das andere zu Bethel, stiftete allda Gottesdienst und opferte, und war eben daselbst der Gottesdienst, so zu Jerusalem war; er schmierte auch den Namen daran, daß er sprach: „Israel, da ist dein GOtt, der dich aus Egypten geführt hat“ [1 Kön. 12, 28. 29.]. Dieser König Jerobeam wußte von dem rechten GOtt wohl, der im Himmel ist, und der den Kindern von Israel hatte aus dem Diensthause Egypti geholfen; noch nimmt er diesen Namen, und schmiert ihn auf seine Abgötterei, und spricht: Derselbige Gott wird alhier bei diesen Kälbern geehrt, man dient ihm allda, und sollte ihm der Dienst gefallen, den sie thaten bei diesen Kälbern, so doch GOtt in seinem Worte gesagt hatte: Zu Jerusalem will ich geehrt und angerufen sein, da der Gnadenstuhl ist [2 Kön. 21, 4.]. Aber dieses mußte alles nicht gelten; darum war man her, fuhr weiter zu, und richtete hier in einer Stadt und dort in einer andern besondere Gottesdienste an, und folgten diesem Exempel Jerobeams nach, daß man in allen lustigen Thälern und grünen Auen Altäre baute, und das Land all voll Gottesdienstes machte, daß man schier deren nicht mehr erdenken konnte. Doch hieß es alles: Siehe, das ist der GOtt, der uns aus Egypten geführt hat, und unsern Vätern die Verheißung gegeben hat. Da ist der Name wohl geblieben, aber das Herz und die Meinung war falsch. Dieses ist nun den Namen GOttes gelästert und geschändet, denn der leidige Teufel ist unter dem Titel und Namen GOttes gewesen.

18. Also kommt und fließt alle Abgötterei aus dem Unglauben und Blindheit her, daß wir GOtt nicht erkennen. So ist unserer Mönche Abgötterei erstlich in der Wüste angefangen, daß sie nicht wollten mit den Leuten umgehen, flohen in die Wildniß, und wollte einer dies und jenes nicht essen, mit andern Leuten keine Gemeinschaft haben, andere wollten diese und jene Kleidung tragen, und mit solchen Dingen dem wahren Gott dienen. Darnach, so sind wir also schändlich herunter gekommen, daß wir glauben, wir könnten durch Ablassbriefe, oder auf Papier gen Himmel kommen, haben darnach Lichtlein und Kerzlein vor die Abgötter gestellt; dadurch wollten wir selig werden.

19. Ist das nicht ein grober Spott? Wenn einer wollte der Egypter spotten, die da Kälber, Ochsen, Raben, Knoblauch und Zwiebeln angebetet haben, so haben wir es ja so arg gemacht als sie, ja, wohl ärger. Denn da gedachte man anders nicht, denn, dies oder jenes Werk will ich thun, will eine Kirche bauen, so wird mir GOtt gnädig sein. Welcher Teufel hat dir das gesagt oder befohlen? Ja, sprichst du, es dünkt mich also gut. Ja, Lieber, dünkt es dich also? Aber du solltest es von GOtt lernen, was gut wäre, so willst du es ihm lehren. Es sind Gedanken, die der Mensch lernt von seinem Herrn und Gott, dem Teufel. Denn es ist kein Gott weder im Himmel oder auf Erden, der sich ließe versöhnen mit Glockengießen, mit Lichtern und Butterbriefen oder Ablassbriefen. Aber du träumst dir selbst aus Eingeben des Teufels von einem solchen Gott, und brauchst solches dem rechten,\footnote{„rechten“ fehlt in der Erlanger.} wahrhaftigen Gott zu Schanden und Unehren.

20. Der wahrhaftige GOtt ist nicht also gesinnt, daß er sich ließe bezahlen mit Lichtlein und Wachs, sondern er schickt seinen Sohn vom Himmel, der wird Mensch, vergießt sein Blut für mich am Kreuze, und stirbt; das wagt er dran, daß ich selig werden möchte. Dieses ist der rechte, wahrhaftige GOtt, die andern Götter alle sind nur leidige Teufel, wiewohl du den Namen GOttes und seines Wortes daran schmirst, und muß darüber GOttes Name unter dem Teufel zu Unehren werden.

21. Ein Idolum, oder Abgötterei, ist anders nichts, denn ein menschlicher Wahn und Gedanke, vom Teufel in das Herz eingebildet, und führt doch den Namen des wahren GOttes; dadurch verleitet denn der Teufel. Wenn ein Mensch dies hört, so denkt er: O, dies Werk gefällt GOtt wohl, der Himmel und Erde gemacht hat. Derhalben so geht man hinan; aber man weiß nicht, daß der Teufel darunter steckt. Also sind des Papsts Briefe voller GOttes und seines Worts gesteckt; aber was ist darunter geweisen
\quellseite{502}
? Anderes nichts, denn der Teufel. Denn alhier hat man gar geschwiegen des HErrn Christi Todes und Blutvergießens, und hat ein jeglicher sich seines Gehorsams und guten Werke gerühmt.

22. Danach, so ist die Abgötterei auch heraußer geschlagen in die äußerlichen Werke, daß wir Stricke, Herzen und Kappen vor uns genommen haben, gleichwie die Egypter Hunde und Katzen angebetet. Und soll die Welt noch eine Weile stehen, so werdet ihr sehen, daß die reine, wahrhaftige Erkenntniß Christi noch gar wird wieder untergehen. Wiewohl der Name Christi wird bleiben; aber es werden darunter Rottengeister und Secten aufstehen, und das Volk mit dem, so in der Schrift nirgends gefunden wird, verführen; als, daß sie die Leute auf die guten Werke weisen werden, und werden dann viel ärger werden, denn sie zuvor gewesen sind.

23. Unserm HErrn GOtt gefällt anders nichts, denn allein das einige Werk Christi. Derhalben, gleich wie es uns zuvor unter dem Babstthum ist gegangen, also wird es uns wieder gehen; wenn dies helle Licht des Evangelii wird wieder verbleichen und verlöschen, so werden wir wieder anbeten die Platten, abgeschornene Haare und Stricke der Barfüßermönche, und sagen: O, dies gefällt GOtt wohl. Aber es ist Abgötterei, und gehört zum Teufel.

24. Nun, Jethro saget alhier: Alle Götter könnten das nicht thun, das dieser GOtt thut; das ist, so du alle Gottesdienste der andern Götter zusammen brächtest, und alle Gutdünkel auf einen Haufen und Klumpen schmelzest, so hülfen sie doch den Menschen gar nicht, wenn man sie gleich alle zusammen auf einen Haufen nähme.

25. Die Mönche sprechen, wenn sie einen in Nöthen trösten: Habe Geduld. Aber es können Heiden auch Geduld haben, und ist ihre Weise auch gewesen, also zu trösten. Darum tröstet solcher Werktheiligen Lehre kein Gewissen, das in Egypten ist, und sie können alle über einem Haufen dem beschwerten Gewissen in Armuth, in Sündenangst oder Todesnoth nicht rathen. Sie können ihren Gott anrufen und ihm helfen [1 Kön. 18, 24. ff.], aber er kann ihnen nicht helfen aus einer einigen Sünde, wenn sie auch gleich zu ihm kommen.

26. So wisse nun, lieber Mensch, wenn du in Sünden liegst, was du thun sollst, nämlich,

daß du wissest, du habest einen solchen GOtt, der dir helfen wolle; denn also ist er gegen dir gesinnt. Du wirst dir selbst nicht helfen können. Da liegt viel an, daß man dies wisse. Die Mönche wissen das nicht, sondern sagen: Habe Reue und Buße für deine Sünde; und lehren unnöthige Dinge, nicht das Erkenntniß des HErrn Christi. Aber wie wird hiermit der Seele geholfen, wenn man mir gleich hundertmal sagte: Hast du gesündigt, so denke und habe Reue und Leid über deine Sünde?

27. Aber es hilft nicht, das Gewissen wird dadurch nicht stille, sicher und zufrieden, sondern je mehr verzagt und blöde, und meint nicht anders, denn GOtt sei zornig, und stehe mit der Keule hinter uns, und wolle uns gar todt haben, bis so lange daß der kommt, der also spricht: Willst du in deinem Gewissen sicher sein, so thue ihm also: Laß fallen deine und aller Menschen gute Werke, und lerne GOtt erkennen, und zweifle nicht daran, dies sei GOttes Art und Natur, daß er denen helfen wolle, die an seinen Sohn Jesum Christum glauben. Das nimm mit festem, gewissem Vertrauen an. Da geht durch Zuthun des Heiligen Geistes der Mensch daran, und befindet es gewißlich also.

28. Derhalben so ist nun unser GOtt größer, denn alle anderen Götter; es gilt nicht mehr, denn Christum erkennen. Und wenn denn das Gewissen recht unterrichtet ist, so folgt auch darauf ein recht gut Werk. Wo der Verstand im Herzen rechtschaffen ist, was GOtt anbetrifft, daß ich weiß, was ich mich zu GOtt versehen solle, dann so bekenne ich ihn auch recht mit dem Munde, und diene ihm auch recht. Sonst, wenn das Herz falsch ist, so sind auch die Werke falsch. Aber hier richtet ein Christ seine Worte und Werke dahin, daß er damit vielen diene, auf daß es nicht ein heimlicher Schatz sei, sondern jedermann gemein werde. Nun folgt das dritte Stück.

\begin{schriftzitat}
V. 13–16. Des andern Morgens setzte sich Mose, das Volk zu richten, und das Volk stund um Mose her vom Morgen an bis zu Abend. Da aber sein Schwäher sahe alles, was er mit dem Volke that, sprach er: Was ist, das du thust mit dem Volke? Warum sitzest du allein, und alles Volk stehet um dich her, von Morgen an bis zu Abend? Mose antwortete ihm: Das Volk kommt zu mir, und fragt GOtt um Rath. Denn wo sie
\end{schriftzitat}

\quellseite{503}
was zu schaffen haben, kommen sie zu mir, daß ich richte zwischen einem jeglichen und seinem Nächsten, und zeige ihnen GOttes Recht und seine Gesetze.

29. Moses lernt von Jethro, seinem Schwäher, wie er die Kinder von Israel regieren solle; denn er sieht, wie er so viel Volks zu Unterthanen habe, nämlich sechsmal hunderttausend streitbare Männer, die mit Mose aus Egypten gezogen waren. Unter so vielen hat sich wahrlich gar viel Zank und Haders erregt und erhoben. Nun ist die Last, alles zu verrichten, einem allein zu schwer; darum, da Moses will alle Sachen selbst verhören, und auf sich nehmen zu entscheiden, eben damit wird das Volk versäumt, und hindert einer den andern; darum sagt er zu ihm:

\begin{schriftzitat}
V. 17. 18. Es ist nicht gut, daß du thust. Du thust närrisch, denn das Volk auch, das mit dir ist. Das Geschäfte ist dir zu schwer, du kannst [es] allein nicht ausrichten.
\end{schriftzitat}

30. Er will sagen: Du, Mose, handelst nicht weislich, daß du alle Sachen lässest auf dich legen, du richtest nichts aus; so werden die Leute ungeduldig darüber, und wird nichts zum Ende gebracht; sondern thue ihm also:

\begin{schriftzitat}
V. 19–22. Gehorche meiner Stimme, ich will dir rathen, und GOtt wird mit dir sein. Pflege du des Volks vor GOtt, und bringe die Geschäfte vor GOtt, und stelle ihnen Rechte und Gesetze, daß du sie lehrest den Weg, darinnen sie wandeln, und die Werke, die sie thun sollen. Siehe dich aber um unter allem Volke nach redlichen Leuten, die GOtt fürchten, wahrhaftig, und dem Geize feind sind; die setze über sie, etliche über tausend, über hundert, über fünfzig, und über zehn, daß sie das Volk allzeit richten.
\end{schriftzitat}

31. Er will, daß, nach Gelegenheit und Vermögen eines jeden Verstandes und Geschicklichkeit, die Leute im Regimente versehen, und ihnen Personen vorgesetzt werden, die richten in ihren Sachen. Was nun dem zu viel ist, der über zehn regiert, das bringe man vor den, so über hundert Mann regiert; und wo es demselbigen auch zu schwer wäre, so bringe man es vor den, so über tausend gesetzt ist. Und allhier ist das Volk erstlich mit Tausenden gezählt worden, davon der Prophet Micheas [Cap. 5, 1.] auch sagt; daß also ein jeglich Tausend zu seinem Obersten und Hauptmann ist gerechnet gewesen. Was aber unter diesen gar keiner hat vertragen oder beilegen können, das hat man sollen gelangen lassen an Mosen, als an die hohe Obrigkeit und obersten Richter.

32. Dies Stück gehört uns auch an. Es ist bis anher das Volk Israel ausgeführt aus Egypten, und ist weder das geistliche oder weltliche Regiment bestellt oder verordnet gewesen; aber da nun dies Volk in der Wüste liegt, und sicher ist vor ihrem Feinde, da geht das Regiment an, und nimmt Moses vor, dies Volk zu ordnen mit weltlichen und geistlichen Gesetzen und Geboten. Und der Heide Jethro greift Mosi vor, gibt ihm Rath und Anleitung, wie er regieren und es recht treffen solle.

33. Ihr habt aber oft gehört, was das weltliche und geistliche Regiment sei. Im geistlichen Reiche, da regiert allein unser HErr GOtt, da ist Christus das Haupt der Gläubigen. Diese Gläubigen werden nicht gesehen; wie denn der HErr Christus auch nicht gesehen wird [Joh. 18, 36.]. Aber des weltlichen Regimentes maßt sich der HErr Christus nicht an, denn da hat er sonst Häupter und Leute zugegeben, die es inne haben, und verwalten nach Maß und Acht der Billigkeit. So sind nun im geistlichen Reiche Diener und Amtleute die Prediger, die da nicht regieren, sondern das Wort Gottes ist alda Regent, und das Predigen von dem Haupte Christo [Apost. 20, 28. 1 Cor. 3, 11.].

34. So wird nun allhier beschrieben, wie das Volk Israel sei äußerlich in eine Ordnung gebracht, und wie das äußerliche weltliche Reich sei bestellt. Er faßt nicht erstlich das geistliche Regiment, sondern das leibliche Reich, regnum animale, wie es St. Paulus [1 Cor. 2, 6.] nennt, davon auch der HErr Christus [Matth. 22, 21.] sagt: „Gebet dem Kaiser, was des Kaisers ist, und Gott, was Gottes ist.“ Also setzt der HErr Christus auch eher das weltliche Regiment denn das geistliche. Und solches darum, daß Noth halben das weltliche Schwert erst muß geordnet werden. Denn es muß zuvor Friede auf Erden sein, sonst kann man nicht predigen, man habe denn Raum und Zeit dazu, daß es stille und guter Friede sei. Wenn Unfriede ist, daß man Spieße, Büchsen und Schwerter in Händen haben soll, da kann man mit der Predigt des göttlichen Worts nicht viel umgehen.

35. Derhalben so hat GOtt auch sechs Tage
\quellseite{504}
gegeben, an denen man arbeiten solle, aber am siebenten Tage solle man allein GOttes Wort hören. Diesen Tag hat er sonderlich zum Frieden geordnet, auf daß man an demselbigen Tage GOttes Wort allein hörete. Und darum will GOtt noch das Regiment in der Welt haben, damit es alles in der Stille und beträglich\footnote{beträglich = verträglich.} daher gehe, und man friedlich leben, des Seinen warten, Kinder auferziehen möge, und vor allen Dingen auch GOttes Wort hören und lernen könne. Dieser Ursache halben wird das Schwert erstlich eingesetzt, und das Volk in eine weltliche Ordnung gebracht, und bestellet, wie man regieren und der Unterthänigkeit sich halten solle.

36. Solch Regiment wird von einem schlechten Menschen, der da kein Zeugniß von GOtt hatte, daß er es thun sollte, vorgeschlagen; wiewohl GOtt es hernach bestätigt hat; sondern Jethro fährt frei heraus, und lehrt Mosen, der da gar voll des Heiligen Geistes war, wie er regieren möge, da doch das Gegenspiel geschehen sollte, und Moses den Heiden lehren, wie er regieren müßte. Aber es geschieht darum, anzuzeigen, wie GOtt das Weltreich in die Vernunft gefasset habe; und da hat er Witz genug gegeben, leibliche Sachen zu regieren. Die Vernunft und die Erfahrung lehren, wie man Weib und Kind regiere, Kühe aus- und eintreiben solle, und was sonst die leibliche Nahrung betrifft. Dieses ist alles der Vernunft Gabe und Geschenk, ihr von GOtt mitgetheilt und verliehen; davon darf man nicht die heilige Schrift um Rath fragen, sondern GOtt hat auch unter alle Heiden solche Gabe in die Rappuse geworfen.

37. Und will GOtt sagen: O, das weltliche Regiment habe ich schon gemacht und bestellt. Wo? Genesis, Cap. 1, 28., da gesagt ward: „Wachset und mehret euch, und füllet die Erde, und herrschet über Thiere, Fische und Vögel“ etc., bringet die Erde unter euch. Da hat GOtt zum weltlichen Regimente Macht und allen Vorrath gegeben, alsbald als er nur den Menschen geschaffen hatte. Dazu darf ich keines Geistes, daß ich ein Pferd anders regieren muß, denn eine Sau oder Kuh; item, daß ich ein Vierfaß anders mache, denn ein Weinfäß; das lehrt die Vernunft und Erfahrung. Solches ist unsern fünf Sinnen unterworfen, daß man es durch die Vernunft beschicken kann; das weltliche Regiment ist in den Stücken gefasset, die da unter uns sind.

38. Es reicht aber nicht dahin in diese Dinge, die GOtt uns nicht unterworfen hat, als, über das Gewissen. Wie denn viel thörichter Fürsten sind, die ihre Macht und Gewalt über den Himmel führen wollen, und die Gewissen regieren, auch was man glauben oder nicht glauben solle; da doch das weltliche Reich nur mit dem umzugehen soll, was die Vernunft fassen kann. Die Vernunft geht mit demjenigen um, was unter uns ist, und nicht, was über uns ist. Denn ich kann einen Ochsen regieren, daß er gehen muß, wie ich es haben will. Ein Haus muß auch gebaut werden und stehen, wie mir es gefällt, oder ich breche es wieder ab, und baue es anders. Darum, so geht das weltliche Regiment nur mit den Sachen und Gütern um, die äußerlichen, leiblichen Dingen unterworfen sind. Und da können Weltleute mit ihrer Vernunft klüger sein in leiblichen Dingen denn geistliche Leute.

39. Heiden sind viel weiter erfunden worden denn Christen; sie haben viel läustiger, ausrichtiger\footnote{läustig = geläufig, hurtig; ausrichtig = zweckmäßig.} und geschickter Weltsachen ordnen und zu ihrer Einigkeit bringen können denn die Heiligen GOttes. Wie denn Christus auch im Evangelio sagt [Luc. 16, 8.]: „Die Kinder der Welt sind klüger, denn die Kinder des Lichts, in ihrer Art.“ Sie wissen besser äußerliche Sachen zu regieren denn St. Paulus oder andere Heiligen. Daher haben auch die Römer so herrliche Gesetze und Rechte gehabt. Denn die Vernunft lehrte sie, daß man die Mörder strafte, die Diebe henken sollte, und wie man sonst Erbgüter austheilen möchte; das haben sie alles gewußt, und fein ordentlich gethan, ohne Rath und Unterricht der heiligen Schrift oder der Apostel. Wie denn St. Paulus in der Epistel, so er ihnen geschrieben, diesfalls nichts gebeut, oder vorschreibt; allein erinnert er sie, daß sie wohlgeordnetem und von GOtt gegebenem Regimente Folge thun [Röm. 13, 1. ff.].

40. So sollen wir nun lernen, daß GOtt weltliche Obrigkeit bestätigt und haben wolle, ob er die gleich nach unsern Gedanken nicht eingesetzt, oder, daß er dazu gleich der Heiden gebraucht. Denn er will die Obrigkeit dazu geordnet haben, nicht, daß man sie allein fürchtete und ehre (wie man ihnen denn auch Ehre schuldig
\quellseite{505}
ist [Röm. 13, 1. ff.], sondern, daß man stille und friedlich lebe, und von GOttes Wort, göttlichem Namen und Reiche predigen könne. Darum, wo eine Obrigkeit ist, da gibt oftmals unser lieber GOtt Gnade, Frieden und Raum dazu, daß man sein Wort predigen möge.

41. Als, zu der Apostel Zeit gab er das römische Kaiserthum. Wiewohl es ein gottloses Reich war, und sich hart wider die Christen legte, doch regierten sie durch die Vernunft, und wurden von jedermanniglich gefürchtet, hielten guten Frieden; es war auch zu ihrer Zeit allenthalben Friede, die Welt stand gar offen. Dies war ein irdisch, vernünftig Reich. Aber wie weltgescheidet und weise sie immer waren, sahen sie dennoch nicht, wozu GOtt sie gebrauchte; nämlich, daß seine lieben Apostel möchten auf und nieder in diesem Kaiserthume fahren, und das Evangelium sicher predigen. Da lief auch das Evangelium schnell, in einem Hui, ungefähr in etlichen, wenigen Jahren, durch das ganze römische Reich, wie denn GOttes Wort schnell läuft [Ps. 147, 15.]. Es ist ein eilendes Wort, in einem Nu reißt und rauscht es durchaus; und wenn es hinweg ist, dann kommt der Teufel, und will auch predigen.

42. Das ist ein Stück dieses Capitels, da das weltliche Regiment mit Geboten, Rechten und Gesetzen geordnet wird, wie er weise, gottesfürchtige Leute über Aemter setzen solle. Und haben gehört, daß GOtt erstlich weltliche Obrigkeit einsetzen läßt bei diesem Volke, ohne sonderlichen Befehl vom Himmel, und wie er auch sonst läßt die Menschen dazu gebrauchen ihrer natürlichen Vernunft. Denn er hat auch desshalben die Vernunft und fünf Sinne den Menschen gegeben. Und was Obrigkeit nach weltlichen Rechten allhier thut, urtheilt und richtet, das hat GOtt gethan und geurtheilt. Wie denn Mosis hiervon spricht: „Das Volk kommt zu mir, und fragt (nicht mich, sondern) GOtt um Rath, alsdann zeige ich ihnen GOttes Recht und Gesetze.“

43. Nun zeigt ihm Jethro auch an, und wir wollen es kürzlich besehen, wie diejenigen sollen geschickt sein, die regieren und in Aemter gesetzt werden sollen, welches die Conditiones oder Zugehöre einer frommen Obrigkeit oder Richters sein sollen. Davon sagt Jethro: „Siehe dich um, unter allem Volke, nach redlichen Leuten, die GOtt fürchten, wahrhaftig, und dem Geize feind sind, die sehe über sie.“ Das sollen die Eigenschaften, Art und Tugenden einer frommen Obrigkeit sein. Wie sie denn im ersten Capitel des fünften Buchs Mosis, V. 16. 17., auch wiederholt werden, da Moses gebeut den Richtern, sie sollen ihre Brüder verhören, und recht richten zwischen ihnen, keine Person im Gerichte ansehen, und niemandes Person scheuen; denn das Gerichtamt ist GOttes. Und 2 Chron. am neunzehnten, V. 6. 7., spricht der König Josaphat zu den Richtern: „Sehet zu, was ihr thut; denn ihr haltet das Gericht nicht den Menschen, sondern dem HErrn, und er ist mit euch in dem Gerichte. Darum lasset die Furcht des HErrn bei euch sein, und hütet euch, und thut es. Denn bei dem HErrn ist kein Unrecht, noch Ansehen der Person, noch Annehmen des Geschenks.“

44. So läßt nun in diesem achtzehnten Capitel die heilige Schrift nichts außen, was da zur Sache dient, oder das allhier fehlen möchte, und malt die Personen mit rechten Farben ab, die da sollen dienen zum weltlichen Reiche. Davon man denn wohl möchte ein Muster nehmen, was man für Leute gebrauchen solle zum geistlichen Amte, oder zur Seelsorge [1 Tim. 3, 1. ff.]. Denn, dieweil GOtt will haben, man soll solche Leute auslesen, die nur über Leib und Gut, und andere äußerliche Sachen zu gebieten und zu herrschen haben sollen: wie viel mehr will er auch solche Leute, und wohl viel bessere haben, die in geistlichen Sachen richten, und der Kirche, mit seinem Blute erworben, vorstehen möchten [1 Tim. 4, 6. ff. Tit. 1, 5. ff. 2 Tim. 1, 8. ff. 2, 1. ff.].

45. Aber, wie tauft er sie, wie malt er sie ab, die da sollen Oberste des Volkes sein in zeitlichen und leiblichen Sachen? Er beschreibt aber nicht allein die amtstücklichen Personen in leiblicher Herrschaft, sondern er sticht auch heimlich mit an, und rührt die Laster und Sünden, so gemeiniglich unter diesen Leuten eingerissen sind, und zu Zoten gehen.\footnote{zu Zoten gehen = einwurzeln.}

46. Erstlich so sollen sie GOtt fürchten. Denn wenn ein Regent, Prälat, Fürst und Herr, Bürgermeister oder Richter nicht GOtt fürchtet, so wird er sein Amt nicht wohl ausrichten; er wird nicht daran gedenken, wenn er jemand Unrecht thäte, daß er Gott damit erzürnen
\quellseite{506}
würde. Denn wer sonst in einem Amte sitzt, und andere richten und regieren soll, dann kommt mancherlei Zank und Hader den Menschen zuhanden, daß sich die Leute zweien und uneinig werden, und geben dem Richter oder Bürgermeister zu schaffen genug. Soll er nun urtheilen und richten, so wird er allezeit einen Theil erzürnen, und auf sich laden; und da wird es gemeiniglich kommen, daß er die großen Hansen und die reichen Junker, seine guten Freunde und Schwäger, oder die sonst vor andern angesehen sind, wird vor den Kopf stoßen und erzürnen müssen.

47. Denn dieselbigen Scharrhansen und große Herren wollen immerdar recht haben, wollen auch, daß man schließen und urtheilen solle, wie es ihnen gefalle. Thut man es nicht, so werden sie des Richters Feinde, und muß denn ein Regent daran wagen sein Gut, Gunst, Weib, Ehre, Leib und Leben, daß ihm daran Schaden widerfahre. Wer da nicht einen Muth hat, und darf es wagen, der gehört hieher nicht. Denn, wenn er sich will fürchten, so ist es um ihn geschehen. Denn, hilft er einem zum Rechte, so hat er den andern zum Feinde; so geht ihm dies und jenes ab. Aber hinweg mit diesem Regenten, er kann nicht fromm sein. Wenn er aber dem Buben hilft, daß er bleibt, und drückt den Armen unter, dann ist einer in der Welt ein guter Fürst, Herr oder Richter.

48. Darum so soll er das Gottfürchten wohl lernen und halten, und gedenken: Wenn gleich dieser reiche, mächtige und starke Hans oder Herr mein Feind wird, da schlage Glück zu, ich habe bei mir in meinem Amte und Berufe noch einen, der viel stärker, ansehnlicher und gewaltiger ist denn er, und wenn dieser gleich alle Teufel, Fürsten und Könige auf seiner Seite hätte, die alle ärger wären, als er ist, was frage ich darnach, wenn der bei mir ist, der droben im Himmel sitzt? [Ps. 56, 4. 5.] Also soll man die Händel beschließen, daß man sage: Lieber GOtt, ich schreibe es dir zu, ob es mir auch darüber mein Leben kostete. Dann spricht GOtt: Ei, halte fest; ich will auch fest halten. Also geht es hindurch, oder muß brechen, und fürchtet einer unsern HErrn GOtt mehr denn die Menschen. Aber wo findest du solche Regenten? Wo sind sie?

49. Moses, im fünften Buche am ersten Capitel, V. 15–17., sagt auch wohl, daß Richter und Obrigkeiten solche Leute sein sollen, da er spricht: er habe sie gesetzt zu Häuptern über tausend, über hundert, über fünfzig, und über zehn, und Amtleute unter die Stämme, und den Richtern geboten, und gesprochen: Sehet zu, richtet recht; denn das Gericht ist GOttes; sehet im Gerichte keine Person an. Ja wohl, es ist ein harter Zusatz; es thut's auch niemand, er habe denn eine sonderliche Gnade und Gabe. Denn es findet sich allerlei, das einen Richter hindert, als Furcht, Gunst, Geld, Gewalt; da muß man denn haben einen festen Muth, gleich als einen steinernen Fels.

50. Davon haben auch die Heiden gesagt: Magistratus ostendit virum, willst du wissen, ob einer ein herzhaftiger, getroster und unverzagter Mann sei, so befehl ihm ein Amt, da wirst du bald sehen, was er sei. Der gemeine Mann weiß nicht, was regieren ist. Wenn einer aber ein Heer führen soll, das Fähnlein in der Hand haben, und vorne an der Spitze stehen soll, so wird ihm alle Kunst, Vernunft und Muth zu kurz, die Hosen werden ihm schlottern, und das Herz zittern; und ist er nicht mit dieser Tugend gefasset, so wird er GOttes gar vergessen.

51. Wenn wünschen helfen sollte (nach der Welt Art zu reden), so sollte man wünschen, daß ein jeder möchte ein Jahr regieren; und wenn einer einem feind wäre, so sollte er ihm gönnen, daß er müßte ein Regent werden. Junge Leute meinen immerdar, sie sind klüger, denn andere Leute, und denken, andere haben närrisch regiert, oder machen es ja noch nicht, wie es sein sollte; sie aber wollten es viel besser machen: wollen elf Regel treffen und umschießen auf einen Wurf, da ihrer nur neune auf dem Bosseleid\footnote{Bosseleid = Regelbahn.} stehen. Kommt es ihnen aber einmal zum Versuchen, so machen es dieselbigen oft am allerärgsten. Ein solch Ding ist es um das Regieren, daß kein Reich rechtschaffen ist, es sei denn die Furcht GOttes dabei; und wenn man einen findet, der diese Gabe hat, der regiert wohl, und es ist eine Gabe GOttes; gleich als ein hübsch Angesicht auch GOttes Gabe ist, welches er in die Kappuse wirft.

52. Sonst, wenn einer wüßte, was Regieren wäre, der ließe davon, oder flöhe in einen Winkel, denn es muß sich einer erwägen, daß er jedermann zum Feinde haben werde. Darum spricht
\quellseite{507}
Jethro allhier, daß Amtleute sollen diese Tugend haben, daß sie GOtt fürchten, und nicht vor Menschen sich entsetzen. Denn, wenn ich GOtt fürchten soll, was soll ich mir denn vor einem Menschen grauen lassen? Zähle mir aber einer die Fürsten und Regenten alle über einen Haufen, die mehr GOtt denn den Menschen fürchten, wie viel, meinst du, wird man ihrer zusammen bringen? Ich wollte alle ihre Namen auf einen Finger schreiben; oder, wie jener sagt, er wollte aller frommen Fürsten Namen auf einen Pettschaftsring graben.

53. Nun, also sollten die Regenten geartet sein, und darum ist es auch geschrieben, daß über diesen Leuten die Obrigkeit sollte geschlagen und gezogen werden. Zwar, wer ein Stücklein davon hat, der danke unserm lieben HErrn GOtt dafür. Denn die Welt ist sonst nicht werth, daß sie solche Amtsverwalter haben sollte. Es gebeut's GOtt wohl, aber niemand hält es; gleichwie die andern Gesetze GOttes auch nicht gethan und erfüllt werden. Es ist wohl ein Gebot, daß man solche Obrigkeiten haben sollte, die da GOtt fürchteten; aber niemand trachtet und richtet sich darnach. David ist ein solcher gewesen, und etliche wenige fromme Könige im Reiche Israel. Denn David fürchtete GOtt, und ist seines Gleichen nicht gewesen, wird auch nicht kommen auf Erden; er ist die Krone aller Fürsten. Noch waren Leute an seinem Hofe, ja, sein eigener Sohn Absalom, die ihm feind waren, und jagten ihn aus dem Königreich, daß David allezeit gleich als auf der Schudel saß [2 Sam. 15, 13. f.].

54. Denn wer da will ein solcher Regent sein, der muß aller Ungunst auf sich laden, auch mancherlei Gefahr gewärtig sein. Wer das nicht ertragen kann, der bleibe davon. Sonst sieht der gemeine Pöbel, daß ein groß Ding sei, obenan sitzen; aber nimm es in die Hand, und siehe, was es sei, sonderlich, wenn du dem nach das Regiment führen willst, daß du deinen Menschen fürchtest. Wenn einem Christen, der es versteht, ein solch Amt angeboten würde, so würde er davor fliehen, und das Amt nicht annehmen. Wenn einer sich will kehren an des Reichen, Gewaltigen Murren oder Zürnen, so werde er kein Amtmann. Denn das wird dir gewißlich begegnen, bist du ein solcher Regent, so wirst du Unlust genug davon haben.

55. Darum will auch GOtt haben, daß man für die Obrigkeit bitten solle, wie auch St. Paulus in der ersten [Epistel] zu Timotheo am andern, V. 2., sagt; item, daß man sie auch in Ehren halte. Denn es kostet viel, Oberherren sein, beide in geistlichem und weltlichem Regimente. Ich rede aber von den rechtschaffenen. Wer da will ein frommer Regent sein, dem wird so viel Unlust, Sorge und Fahr unter Augen stoßen, daß er lieber dafür möchte Steine tragen, oder ein Bauersmann sein. So aber einer in der Regierung ein Schalk und Bube sein will, der mag seinen Lohn allhier himweg nehmen; dort wird er es wohl finden, wie er es haben und bekommen soll [Luc. 16, 25.].

56. Das ist eine Tugend, nämlich GOtt fürchten, dies sollen Regenten thun. Ihrer werden aber wenig befunden, die solches thun, achten und betrachten.

57. Die andere Tugend ist wahrhaftig sein, daß er die Wahrheit liebe. Denn das folgt auf die Gottesfurcht, wo man GOtt vor Augen hat, pocht und trotzt auf den droben, da ist man auch wahrhaftig, da wird man mit herzlicher Liebe und Lust thun, was recht ist, und wird es getreulich ausrichten. Sonst finden sich allerlei Tücke, Schein und Gleisnerei in der Welt, daß man die Leute mit der Nase umführt; item, den Rechten eine wächserne Nase andreht, und ehe man sich umsieht, so hat man aus Gutem Böses gemacht; und wiederum, auch was Unrecht ist, zu Rechte verkehrt [Jes. 5, 20.]. Davon redet Salomo viel im Buche der Sprüche, daß er, als selbst ein großer König, aus vieler Erfahrung gelernt hat. So soll nun ein Regent dem Rechten und der Wahrheit geneigt sein, und helfen dem, das billig und dem Rechten ähnlich ist, und hinwieder verdammen, was unrecht wird befunden, hintan gesetzt alles Ansehen der Personen.

58. Die dritte Tugend ist, daß er dem Geiz feind sei. Die erste Tugend, die geht gegen GOtt; aber die andern zwo Tugenden ziehen sich mehr herunter gegen die Leute. Aber wo findet man solche Herren? Ich halte es dafür, wenn man die Geizigen in der Welt abscheiden und absondern sollte, so würde man wenig Fürsten finden, die nicht geizig wären, und nicht ihre Unterthanen schatzten, schindeten und aussaugten. An den Höfen regiert jetzt Untreu, Finanzerei,\footnote{Finanzerei = Ränke, Betrug.} Eigenmuß, und der Geiz in den
\quellseite{508}
Fürsten und in ihren Räthen. Denn sie haben Raum und Ursache dazu, und beschönen es mit diesem Deckel und Vorwenden; dieweil in der Obrigkeit Amt sitzen gebieret und trägt vielerlei Anfechtungen. Rechnen aber dieses gar nicht, daß sie werden geehrt und hochgehalten. Man gibt ihnen ihren Geschoß, Steuer, Zins und Rente, dazu so werden sie auch sehr gefürchtet. Daher kommt es denn, daß sie, um der Ehre, Guts und Gewalt willen, ein stolz, hoffärtig Herz bekommen, daß sie sagen: Man muß thun, was ich haben will, und gedenken nicht, daß sie über sich auch einen GOtt im Himmel sitzen haben [Eph. 6, 9.].

59. Erstlich wird einer aus Uebermuth frech und kühn, bricht und beugt das Recht nach alle seinem Gefallen, darum, daß er sich läßt dünken, er sei gewaltig; darnach, auf daß ihm die Hände geschmiert und gefülbert, das ist, Geld und Geschenke ihm zugebracht und zugetragen werden. Daher wird auch im 5. Buch Mosis am sechzehnten, V. 19., gesagt, daß die Geschenke der Weisen Augen verblenden. Es muß ein ehrbar Herz sein, und recht gottesfürchtig, daß sich nicht mit der Welt Gütern sollte betrügen, und Sachen nicht verderben, sondern dem Armen seine gute Sache und Gerechtigkeit des Reichen Unrecht vorziehen\footnote{Eislebensche: ersürziehen.} [Dan. 5, 17.].

60. Es ist sonst eine große Gelegenheit und Occasion im Regimente, Unrecht zu thun, sonderlich wenn einer geizig ist, und GOttes nicht achtet. Wiewohl noch andere Ursachen mehr sind, einem Unrecht zu thun, als, Zorn, Neid und Haß, wenn einer einem Regenten ist zu nahe gewesen, oder etwas zuwider gethan, da rächt man sich gerne bald wieder. Darum sind dies die vornehmsten Ursachen, daß man in der Obrigkeit Stande das Recht verkehret und nicht gleich zu handelt; als: wenn man GOtt nicht fürchtet, nicht wahrhaftig ist, und Geld oder Geschenke nimmt. Aber soll es recht zugehen, so gehört hierher ein solcher Mann, der da nicht steht nach Ehre, Gewalt, Gut oder Geld.

61. Das ist eine feine Regel, was für Leute die Regenten sein sollen. Er spricht nicht allein, daß sie nicht geizig sein sollen, sondern auch dem Geiz feind sein; das ist, also redlich sollen sie sein, daß sie nicht allein für ihre Person milde sind, sondern auch andern wehren, die sich den Geiz wollen reiten lassen. Sie sollen nicht Gift\footnote{Gift = Geschenk, Gabe.} und Gaben nehmen, sondern recht handeln, GOtt, der Wahrheit zu Dienste und zu Wohlgefallen. Aber siehe dich um, du wirst solcher Regenten wenig finden; wie denn Mosis ihrer auch wenig gefunden hat.

62. Nichtsdestoweniger muß man dieses den Regenten predigen, auf daß sie wissen, wie sie sollen geschickt sein, sonderlich, daß sie GOtt fürchten. Wer diese Gnade nicht hat, der lasse das Amt stehen. Denn man wird allezeit Mangel an den Leuten sehen und befinden. Das ist aber gewißlich wahr, wenn sie sich erkennen, wer sie wären, und ihr Vermögen erwägen, man würde einem ein Fürstenthum nachtragen, und er würde es abschlagen und nicht annehmen wollen. Man würde ihn dazu bitten und flehen müssen; noch würde er sich des Regiments weigern. Aber wenn man selbst darnach rennt, läuft, steht und bittet, so steht es nicht wohl, ist gar ein böse Zeichen; es gebiehet auch nimmermehr den Leuten zur Besserung, wie die Erfahrung gibt und lehrt.
